% Using report class because it's mostly like article but has chapters. https://tex.stackexchange.com/questions/36988/regarding-the-book-report-and-article-document-classes-what-are-the-mai
\documentclass[11pt, letterpaper, landscape]{report}
\usepackage{multicol}
\usepackage[inline]{enumitem}

\usepackage[letterpaper, left=.4in, right=.4in, top=.25in, bottom=.5in]{geometry}

\usepackage{graphicx}
\graphicspath{{/home/ben/nxc/dont_delete_images/}}

\usepackage{xcolor}
\usepackage{fontspec}
\usepackage{amsmath}
\usepackage{amssymb}
\usepackage{accents}
\usepackage{unicode-math}
\usepackage{libertinus}
\setmainfont{Libertinus Serif}
%\setmathfont{Libertinus Math}
\setmonofont{Source Code Pro}[Scale=MatchLowercase]

\newif\iflecturenotes
\lecturenotestrue  % comment out with % for textbook
%\lecturenotesfalse % comment out with % for lecture notes

\iflecturenotes
  % Lecture notes version: show notes, hide book content
  \newenvironment{notes}{}{}
  \usepackage{comment}
  \excludecomment{book}
\else
  % Textbook version: show book content, hide notes
  \newenvironment{book}{}{}
  \usepackage{comment}
  \excludecomment{notes}
\fi

% Gets rid extra space caused by \left and \right
% https://tex.stackexchange.com/questions/2607/spacing-around-left-and-right
\let\originalleft\left
\let\originalright\right
\renewcommand{\left}{\mathopen{}\mathclose\bgroup\originalleft}
\renewcommand{\right}{\aftergroup\egroup\originalright}

% This did the same thing as above, but requires use of mleft and mright which breaks use of cdlatex commands like lr( TAB
%\usepackage{mleftright}

%%%%%%%%%%%%%%%%%%%%%%%%%%%%%%%%%%%%%%%%%%%%%%%%%%%%%%%%%%%%
% page numbering
%%%%%%%%%%%%%%%%%%%%%%%%%%%%%%%%%%%%%%%%%%%%%%%%%%%%%%%%%%%%
\usepackage{fancyhdr}
\pagestyle{fancy}
\fancyhf{} % clear existing header/footer entries
% Place Page X of Y on the right-hand
% side of the footer
\fancyfoot[C]{p\thepage~\S\thesection}

%%%%%%%%%%%%%%%%%%%%%%%%%%%%%%%%%%%%%%%%%%%%%%%%%%%%%%%%%%%%
% Adjusting spacing; maybe more adjusting needed later
%%%%%%%%%%%%%%%%%%%%%%%%%%%%%%%%%%%%%%%%%%%%%%%%%%%%%%%%%%%%
\setlength\parskip{1.4ex plus 0.2ex minus 0.2ex}
\setlength\parindent{0pt}
\setlength\columnsep{18pt}
%\setlist{labelwidth=2em}
\setlist{leftmargin=24pt, topsep=0.3ex,
         partopsep=0pt, itemsep=0.7ex plus 0.2ex minus 0.2ex, parsep=0pt}

%%%%%%%%%%%%%%%%%%%%%%%%%%%%%%%%%%%%%%%%%%%%%%%%%%%%%%%%%%%%
% Title Headings
%%%%%%%%%%%%%%%%%%%%%%%%%%%%%%%%%%%%%%%%%%%%%%%%%%%%%%%%%%%%
\usepackage{titlesec}
\titleformat*{\section}{\Large\bfseries}
\titleformat*{\subsection}{\large\bfseries}
\titleformat*{\subsubsection}{\normalsize\bfseries}
\titleformat*{\paragraph}{\normalsize\bfseries}
\titleformat*{\subparagraph}{\normalsize\bfseries}
\titlespacing*{\section}{0pt}{0pt}{0pt}
\titlespacing*{\subsection}{0pt}{0pt}{0pt}

%%%%%%%%%%%%%%%%%%%%%%%%%%%%%%%%%%%%%%%%%%%%%%%%%%%%%%%%%%%%
% Counters and environments for Examples, Thms, etc.
%%%%%%%%%%%%%%%%%%%%%%%%%%%%%%%%%%%%%%%%%%%%%%%%%%%%%%%%%%%%
% let's just make warm-ups exercises like so:
% \exercise (\emph warm-up)

% in-section exercises
\newcounter{exercise}[section]
\newcommand{\exercise}{\refstepcounter{exercise}\par%
\noindent\textbf{Exercise~\theexercise.} }

% end-of-section exercises
\newcommand{\exatend}{\refstepcounter{exercise}\par%
\noindent\textbf{\theexercise.} }

\newcounter{example}[section]
\newcommand{\example}{\refstepcounter{example}%
\par\noindent\textbf{Example~\theexample.} }

\newenvironment{definition}{\begin{bbox}\par\noindent\textbf{Definition. }}
  {\end{bbox}}

\newcounter{theorem}[chapter]
\counterwithin{theorem}{chapter}
\newenvironment{theorem}[1][]{\refstepcounter{theorem}\begin{bbox}\par
    \noindent\textbf{Theorem~\thetheorem: #1}}{\end{bbox}}

\newenvironment{lemma}[1][]{\refstepcounter{theorem}\begin{bbox}\par
    \noindent\textbf{Lemma~\thetheorem: #1}}{\end{bbox}}

\newcommand{\exerciseheading}{\noindent\textbf{\Large{\thesection~Exercises}}}

%\usepackage{footnote}
\usepackage{dblfnote}
\newcommand{\hintfootnote}[1]{%
\begingroup
\renewcommand{\thefootnote}{H}%
\footnote{\theexercise. Hint: #1}%
\endgroup}

\newcommand{\creditfootnote}[1]{%
\begingroup
\renewcommand{\thefootnote}{C}%
\footnote{\theexercise. Credit: #1}%
\endgroup}

%%%%%%%%%%%%%%%%%%%%%%%%%%%%%%%%%%%%%%%%%%%%%%%%%%%%%%%%%%%%%%% 
% This creates two columns out of minipages
%%%%%%%%%%%%%%%%%%%%%%%%%%%%%%%%%%%%%%%%%%%%%%%%%%%%%%%%%%%%%%% 
\newcommand{\btwocol}[2]{%~\\ just put new line yourself
  \begin{minipage}[t]{0.48\linewidth} \vspace{0em} #1 \end{minipage}
  \hfill
  \begin{minipage}[t]{0.48\linewidth} \vspace{0em} #2 \end{minipage}
}

\newcommand{\bbtwocol}[4]{%#1 and #2 should be decimals which sum to at most 1. #3 and #4 are content.
  \begin{minipage}[t]{#1\linewidth} \vspace{0em} #3 \end{minipage}
  \hfill
  \begin{minipage}[t]{#2\linewidth} \vspace{0em} #4 \end{minipage}
}

\newcommand{\bbthreecol}[6]{%#1 - #3 should be decimals which sum to at most 1. #4 - #6 are content.
  \begin{minipage}[t]{#1\linewidth} \vspace{0em} #4 \end{minipage}
  \hfill
  \begin{minipage}[t]{#2\linewidth} \vspace{0em} #5 \end{minipage}
  \hfill
  \begin{minipage}[t]{#3\linewidth} \vspace{0em} #6 \end{minipage}
}

\newcommand{\bbfourcol}[8]{%#1 - #4 should be decimals which sum to at most 1. #5 - #8 are content.
  \begin{minipage}[t]{#1\linewidth} \vspace{0em} #5 \end{minipage}
  \hfill
  \begin{minipage}[t]{#2\linewidth} \vspace{0em} #6 \end{minipage}
  \hfill
  \begin{minipage}[t]{#3\linewidth} \vspace{0em} #7 \end{minipage}
  \hfill
  \begin{minipage}[t]{#4\linewidth} \vspace{0em} #8 \end{minipage}
}

\newlength\oldfboxsep%
\newlength\myminiwidth%

\newsavebox{\bboxbox}
\newenvironment{bbox}{%
  \setlength\oldfboxsep\fboxsep
  \setlength{\fboxsep}{.1in}%
  \setlength\myminiwidth\linewidth%
  \addtolength\myminiwidth{-0.2in}%
  \vspace{0.1in}%
  \begin{lrbox}{\bboxbox}\begin{minipage}{\myminiwidth}%
    }{%
    \end{minipage}\end{lrbox}%
  \noindent\fbox{\usebox{\bboxbox}}\vspace{0.1in}\setlength\fboxsep\oldfboxsep%
}

%%%%%%%%%%%%%%%%%%%%%%%%%%%%%%%%%%%%%%%%%%%%%%%%%%%%%%%%%%%% 
% Math Style
%%%%%%%%%%%%%%%%%%%%%%%%%%%%%%%%%%%%%%%%%%%%%%%%%%%%%%%%%%%% 
\newcommand{\R}{\mathbb{R}}
\newcommand{\Poly}{\mathbb{P}}

\newcommand{\Lim}[1]{\raisebox{0.5ex}{\scalebox{0.8}{$\displaystyle \lim_{#1}\;$}}}

\DeclareMathOperator{\col}{col}
\DeclareMathOperator{\row}{row}
\DeclareMathOperator{\nul}{null}
\DeclareMathOperator{\rank}{rank}
\DeclareMathOperator{\spn}{span}
\DeclareMathOperator{\proj}{proj}
\DeclareMathOperator{\range}{range}

\newcommand{\bbvec}[2]{
  \begin{bmatrix} #1 \\ #2 \end{bmatrix}
}

\newcommand{\bbbvec}[3]{
  \begin{bmatrix} #1 \\ #2 \\ #3 \end{bmatrix}
}

\newcommand\nm[1]{\left\lVert #1 \right\rVert}

%%%%%%%%%%%%%%%%%%%%%%%%%%%%%%%%%%%%%%%%%%%%%%%%%%%%%%%%%%%% 
% Lines in augmented matrices: https://latex.org/forum/viewtopic.php?t=22393
%%%%%%%%%%%%%%%%%%%%%%%%%%%%%%%%%%%%%%%%%%%%%%%%%%%%%%%%%%%% 
\makeatletter
\renewcommand*\env@matrix[1][*\c@MaxMatrixCols c]{%
  \hskip -\arraycolsep
  \let\@ifnextchar\new@ifnextchar
  \array{#1}}
\makeatother


%%%%%%%%%%%%%%%%%%%%%%%%%%%%%%%%%%%%%%%%%%%%%%%%%%%%%%%%%%%% 
% Horizontal, tabbed question lists
% Argument to horiparts is the number of columns.
%%%%%%%%%%%%%%%%%%%%%%%%%%%%%%%%%%%%%%%%%%%%%%%%%%%%%%%%%%%% 

\RequirePackage{tabto}

% This complicated one added extra space above for no reason. Simplified to below.
% \def\changemargin#1{\list{}{\leftmargin#1}\item[]}
% \let\endchangemargin=\endlist

% \newenvironment{horiparts}[1]{\begin{changemargin}{1.4em}\noindent %
%     \NumTabs{#1} \begin{enumerate*}[label=(\alph*) \hspace{1ex}, itemjoin={\tab}]}
%     {\end{enumerate*}\end{changemargin}}

\newenvironment{horiparts}[1]{\noindent %
    \NumTabs{#1} \begin{enumerate*}[label=(\alph*) \hspace{1ex}, itemjoin={\tab}]}
    {\end{enumerate*}}

\title{Yet Another First Course in Linear Algebra}
\author{Ben Smith}
\date{}

\definecolor{darkblue}{rgb}{0,0,.7}
\usepackage[bookmarksnumbered=true,colorlinks=true,linkcolor=darkblue]{hyperref}

%%%%%%%%%%%%%%%%%%%%%%%%%%%%%%%%%%%%%%%%%%%%%%%%%%%%%%%%%%%% 
% vectorgrid stuff begins here
%%%%%%%%%%%%%%%%%%%%%%%%%%%%%%%%%%%%%%%%%%%%%%%%%%%%%%%%%%%% 
\usepackage{tikz}
\usepackage{tkz-euclide} % for labelling angles
\usepackage{pgfplots}
\usetikzlibrary{arrows.meta,calc,positioning, decorations.pathreplacing}%,external}
% speeds up compiling
%\tikzexternalize[prefix=figures-ext/]

% Define pgfkeys for the vector grid
\pgfkeys{
  /vectorgrid/.cd,
  xmin/.store in=\vgxmin,
  xmax/.store in=\vgxmax,
  ymin/.store in=\vgymin,
  ymax/.store in=\vgymax,
  step/.store in=\vgstep,
  labint/.store in=\vglabint,
  xu/.store in=\vgxu,
  yu/.store in=\vgyu,
  xv/.store in=\vgxv,
  yv/.store in=\vgyv,
  % scale/.store in=\vgscale,
  % Default values
  xmin=-10,
  xmax=10,
  ymin=-10,
  ymax=10,
  step=1,
  labint=2,
  xu=1,
  yu=3,
  xv=-2,
  yv=4,
  % scale=0.7,
}

% The main command
\newcommand{\vectorgrid}[1][]{%
  \pgfkeys{/vectorgrid/.cd, #1}%
  %   Clip to viewing area
  \clip (\vgxmin,\vgymin) rectangle (\vgxmax,\vgymax);

  % Standard grid
  \draw[step=\vgstep cm, gray, very thin] (\vgxmin,\vgymin) grid (\vgxmax,\vgymax);

  % Axes
  \draw[-] (\vgxmin,0) -- (\vgxmax,0);% node[right] {$x$};
  \draw[-] (0,\vgymin) -- (0,\vgymax);% node[above] {$y$};

  \pgfmathsetmacro{\bignum}{50}
  % Main vector lines through origin (thicker)
  % \draw[thick, blue] ($-\bignum*(\vgxu,\vgyu)$) -- ($\bignum*(\vgxu,\vgyu)$);
  % \draw[thick, red] ($-\bignum*(\vgxv,\vgyv)$) -- ($\bignum*(\vgxv,\vgyv)$);

  % Lines parallel to u vector, passing through multiples of v
  \foreach \k in {-20,...,20} {
    \pgfmathsetmacro{\px}{\k*\vgxv}
    \pgfmathsetmacro{\py}{\k*\vgyv}
    % Only draw if the base point is reasonably close to viewing area
    \pgfmathparse{abs(\px) <= abs(\vgxmax)*2+abs(\vgxmin)*2 && abs(\py) <= abs(\vgymax)*2+abs(\vgymin)*2 ? 1 : 0}
    \ifnum\pgfmathresult=1
      \draw[blue, very thin] ($(\px,\py) - \bignum*(\vgxu,\vgyu)$) -- ($(\px,\py) + \bignum*(\vgxu,\vgyu)$);
    \fi
  }

  % Lines parallel to v vector, passing through multiples of u
  \foreach \k in {-20,...,20} {
    \pgfmathsetmacro{\px}{\k*\vgxu}
    \pgfmathsetmacro{\py}{\k*\vgyu}
    % Only draw if the base point is reasonably close to viewing area
    \pgfmathparse{abs(\px) <= abs(\vgxmax)*2+abs(\vgxmin)*2 && abs(\py) <= abs(\vgymax)*2+abs(\vgymin)*2 ? 1 : 0}
    \ifnum\pgfmathresult=1
      \draw[blue, very thin] ($(\px,\py) - \bignum*(\vgxv,\vgyv)$) -- ($(\px,\py) + \bignum*(\vgxv,\vgyv)$);
    \fi
  }

  % Draw vectors u and v from origin (optional - for visualization)
  \draw[-Stealth, very thick, blue] (0,0) -- (\vgxu,\vgyu);% node[above right] {$\vec{u}$};
  \draw[-Stealth, very thick, blue] (0,0) -- (\vgxv,\vgyv);% node[above left] {$\vec{v}$};

  % Axis labels
  \pgfmathtruncatemacro{\intxmin}{ceil(\vgxmin)}
  \pgfmathtruncatemacro{\intxmax}{floor(\vgxmax)}
  \foreach \x in {\intxmin,...,\intxmax} {
    \pgfmathparse{mod(\x,\vglabint) == 0 && \x != 0 && \x != \vgxmin && \x != \vgxmax ? 1 : 0}
    \ifnum\pgfmathresult=1
      \draw (\x,0.1) -- (\x,-0.1) node[anchor=north,fill=white,inner sep=1pt] {\small$\x$};
    \fi
  }

  \pgfmathtruncatemacro{\intymin}{ceil(\vgymin)}
  \pgfmathtruncatemacro{\intymax}{floor(\vgymax)}
  \foreach \y in {\intymin,...,\intymax} {
    \pgfmathparse{mod(\y,\vglabint) == 0 && \y != 0 && \y != \vgymin && \y != \vgymax ? 1 : 0}
    \ifnum\pgfmathresult=1
      \draw (0.1,\y) -- (-0.1,\y) node[anchor=east,fill=white,inner sep=1pt] {\small$\y$};
    \fi
  }
  % \end{tikzpicture}
}

%%%%%%%%%%%%%%%%%%%%%%%%%%%%%%%%%%%%%%%%%%%%%%%%%%%%%%%%%%% 
% vectorgrid stuff ends here
%%%%%%%%%%%%%%%%%%%%%%%%%%%%%%%%%%%%%%%%%%%%%%%%%%%%%%%%%%% 

% parallelogram symbol!
\newcommand*{\pgram}[1][]{%
  \pgfpicture\pgfsetroundjoin\pgfsetlinewidth{0.9pt}
    \pgftransformxslant{.6}%
    \pgfpathrectangle{\pgfpointorigin}{\pgfpoint{.60em}{.65em}}%
    \pgfusepath{stroke,#1}%
  \endpgfpicture}

\begin{document}
\begin{multicols*}{2}

\raggedcolumns


\hspace{1pt}

\vspace{2in}

\begin{center}
\begin{Large}
Yet Another First Course in Linear Algebra
\end{Large}

\vspace{2em}
\begin{large}
Ben Smith

LA Pierce College
\end{large}

© 2026 Licensed under CC BY-SA 4.0. To view a copy of this license, visit \href{https://creativecommons.org/licenses/by-sa/4.0/}{https://creativecommons.org/licenses/by-sa/4.0/}
\end{center}

\newpage
%Trick to get page numbers to line up with PDF page numbers:
%\addtocounter{page}{1} % at start of chapter 1.
\tableofcontents

\section{Preface for Instructors}
\subsection{Features}
There are many free elementary linear algebra textbooks to choose from. Here are some features of this one:
\begin{itemize}
\item \textbf{Lots of exercises.} Most sections have 20 to 50, including some true/false, basic calculations, easy proofs, and a few hard proofs. There's an emphasis on making connections between a graphical view of linear combinations, systems of linear equations, linear transformations, and matrix-vector multiplication.
\item \textbf{Open source.} You are free to modify your own copy of the \LaTeX{} code as much or as little as you like, distribute your version freely, as long as you use the CC BY-SA 4.0 license and give me credit (eg, ``by Ben Smith and Your Name'').
\item \textbf{Lecture notes/textbook combo.} The file \texttt{yafcla\_notes.pdf} includes big blank spaces to work out examples and some proofs in class. The file \texttt{yafcla\_book.pdf} includes the same content but with example solutions and all proofs included. In class, the instructor can either annotate the PDF with tablet/projector, or write on the printed notes with a document camera/projector, or record videos for a flipped classroom. Students can write on the printed notes or annotate their copy of the PDF notes. The landscape orientation was chosen because it's the natural orientation of most screens. Instructors probably want to modify the lecture notes to fit their class, by expanding/contracting blank space, or by adding/removing examples or space for proofs.
\item \textbf{Flexible rigor.}  In my experience, expecting linear algebra students to write some simple proofs in HW and tests shakes them out of their recipe-following habits from earlier math classes, and results in much better understanding of the concepts. The subject matter is ripe for lots of simple proof exercises. It's also a good place to learn \emph{about} mathematical rigor, and one principle of mathematical rigor is that every theorem should have a proof, based on previous definitions and theorem; so, in the textbook, every theorem (with one exception in chapter 3) is proven. That said, depending on your students, it's probably not a good idea to spend class time proving everything. Some of the proofs are much harder than anything I'd expect a student to do, and some are tedious and not instructive. For the proofs I teach in my class, there's a blank space in the lecture notes under the theorem, and other proofs are relegated to the textbook. In your class, you may prefer to cover more or fewer proofs. You can even skip assigning proofs for HW altogether. 
\item \textbf{Motivation first.} Where possible, I have tried to include concrete examples to illustrate definitions and theorems \emph{before} the definition/theorem. Sometimes these give an argument why a theorem should true in general, which I've found is very helpful for students at this level, but shouldn't be confused with a ``proof'' by example. The exposition in chapters \ref{chap:determinants} and \ref{chap:eigenstuff} are unconventional, in a motivation-first kind of way. Chapter \ref{chap:determinants} introduces determinants as \emph{signed volume} of parallelograms/parallelopipeds, and shows how this leads to the traditional cofactor expansion, an approach inspired by Sergei Treil's excellent \emph{Linear Algebra Done Wrong}. Chapter \ref{chap:eigenstuff} covers linear difference equations and diagonalizability before eigenstuff, which arises naturally as a way to describe linear difference equations.
\item \textbf{Brisk pace.} At my school, linear algebra is sadly only a 3-unit class which meets twice a week. This makes it hard to finish the material with many textbooks, which are often divided into too many sections, or repeat too much material. This book covers a standard outline of topics, plus an optional proof-writing section, in 25 sections (or, it \emph{will} have 25 sections, after I write chapter 6). The longer sections are divided into subsections, in case your school has shorter class periods.
\item \textbf{No calculator needed.} Every exercise and example in the book can be done with a reasonable amount of computation by hand. Fractions and big numbers are avoided. Instead of asking questions like ``\emph{diagonalize the following $3\times 3$ matrix}'' which takes many steps and may require a computer, we'll have questions like ``\emph{write the characteristic equation of the following matrix},'' or ``\emph{given that the eigenvalues of the following matrix are 8, 9, and 10, diagonalize it.}'' I have a few reasons for doing this. One is that I can give exams with no calculators allowed, which eliminates the advantage that students students with a pricey graphing calculators have. Another reason is that I believe that doing mental math gives students a better \emph{feel} for things like linear combinations and matrix multiplication. 
\end{itemize}

\subsection{Possible downsides}
Here are some things that you may find are downsides of the book:
\begin{itemize}
\item \textbf{No technology instruction.} If you want to teach your students to use sympy, sage, or MATLAB, this book won't help you do it. And all of the exercises are doable by hand, so students won't feel motivated to use technology. If you have time, and want to make technology part of your class, I recommend David Austin's excellent free online book \emph{Understanding Linear Algebra}, which teaches sage with fancy embedded code widgets throughout the book. It also has some excellent application sections which you could use to supplement this book.
\item \textbf{Not a ton of applications.} This book does have some application exercises, and I do try hard to give a sense of how linear algebra is useful for anyone studying or working in STEM. But it doesn't include the wide range of application problems that some books have. Again, you may want to supplement with Austin. And if you have problems you think I should include, feel free to email them to me at smithbt@laccd.edu or with a github pull request.
\item \textbf{It is a work in progress.} Here is a list of things, from highest to lowest priority, that I still need to do, as of January 8, 2026:
\begin{itemize}
\item A mini-section about the Graham-Schmidt process.
\item Chapter 6 (Symmetric matrices and quadratic forms).
\item More synthesis/review type questions in the last few chapters.
\item Solutions to exercises.
\item More application exercises throughout.
\end{itemize}
\end{itemize}

\subsection{How to compile your own version}
There are many good reasons to edit the \LaTeX{} code and compile your own version. Maybe you want your own version of the lecture notes with different examples, or more/less blank space. Maybe you want to add your own exercises. Maybe you want the book to be in portrait orientation. Maybe you want little pictures of cute cats on every page. Go for it! Here's what you need to do:
\begin{enumerate}
\item Install a distribution of \LaTeX{} on your computer. Windows users usually use MiKTeX, mac users usually use MacTex, and linux users usually use TeXLive.
%% Commented out the next step because adding the line \usepackage{libertinus} installs the fonts automatically.
%\item Install the free Libertinus font family on your computer. It's easy to find on the internet, and installation instructions depend on your operating system. Alternatively, you can use any font you have installed on your computer, as long as it comes with math support, by changing the following lines in yafcla.tex:
%\begin{verbatim}
%\setmainfont{MyFont Serif}
%\setmathfont{MyFont Math}
%\end{verbatim}
%You could also comment out these lines (with \texttt{\%} at the beginning) and it will use the default Computer Modern font. Note that changing the font may make blank spaces and column/page breaks move to awkward places.
\item Make your changes to the yafcla.tex file. If you want to make a version of the lecture notes, make sure the second of the following two lines is commented out like so:
\begin{verbatim}
\lecturenotestrue  % comment out with % for textbook
%\lecturenotesfalse % comment out with % for lecture notes
\end{verbatim}
This will make it so anything in a \texttt{book} environment will be ignored when you compile. If you want to make a version of the textbook, move the \texttt{\%} to the other line and anything inside a \texttt{notes} environment will be ignored.
\item Compile with the \texttt{lualatex} command. In a terminal (aka console or command prompt), use \texttt{cd} to  navigate to the directory where you have your \texttt{yafcla.tex} file, and run the command:

\texttt{lualatex --file-line-error --synctex=1 \\ -interaction=nonstopmode yafcla.tex}

This may prompt you to install a bunch of latex packages, and may take some time, but it should eventually create \texttt{yafcla.pdf}. Then, run the above command again because the first time, labels and references won't be linked up. Then you should rename \texttt{yafcla.pdf} to something else, like \texttt{yafcla\_book.pdf} or \texttt{yafcla\_notes.pdf}, so that doesn't get overwritten when you compile the other version.
\end{enumerate}

\columnbreak
\section{Preface for Students}

\subsubsection{What is linear algebra about?}

One of the first kinds of equations you learned about in a basic algebra class was linear equations---${y=mx+b}$. And in science classes, many fundamental laws and models can be written as linear equations, specifically with ${b =0}$:
\begin{align*}
\begin{split}
\text{Force} &= \text{(Mass)}\text{(Acceleration)}, \\
\text{Voltage} &= \text{(Current)}\text{(Resistance)}, \\
\text{Force of spring} &= \text{(Spring constant)}\text{(Displacement)}, \\
\text{Part} &= \text{(Proportion)}\text{(Whole)},
\end{split}
\end{align*}

and so on. But the world is often more complicated, with many $y$ variables that depend on many $x$ variables. And some (but not all) of the above equations can have multiple variables, like $x_1, x_2, x_3$ and $y_1, y_2, y_3$ instead of just $x$ and $y$. So instead of just ${y=mx}$ you might have
\begin{align} \label{eq:basic system}
  \begin{split}
    y_1 &= m_{11}x_1+m_{12}x_2+m_{13}x_{3} \\
    y_2 &= m_{21}x_1+m_{22}x_2+m_{23}x_{3} \\
    y_3 &= m_{31}x_1+m_{32}x_2+m_{33}x_{3}
  \end{split}
\end{align}

Linear algebra is fundamentally about \emph{systems of linear equations} like this, which show up in every STEM discipline when you get beyond 100 and 200-level courses. We're going to start by learning how to solve them, which means finding the $x$'s that make all the equations work, given values of the $y$'s and $m$'s, and determine when the system has no solution, one solution, or an infinite number of solutions.

Then, we'll see how we can write the system as a \emph{vector equation}, like
\[ \bbbvec{y_1}{y_2}{y_3} = x_1\bbbvec{m_{11}}{m_{21}}{m_{31}} +
x_2\bbbvec{m_{12}}{m_{22}}{m_{32}} +
x_3\bbbvec{m_{13}}{m_{23}}{m_{33}}\]

or as a \emph{matrix equation}, like
\[ \bbbvec{y_1}{y_2}{y_3} = M \bbbvec{x_1}{x_2}{x_3} \text{ where } M =
\begin{bmatrix}
  m_{11} & m_{12} & m_{13} \\
  m_{21} & m_{22} & m_{23} \\
  m_{31} & m_{32} & m_{33}
\end{bmatrix}
\]

In Chapter \ref{chap:matrices-and-transformations} we'll see how matrices themselves are really interesting mathematical objects, which can be combined with addition and multiplication just like numbers. And just like ${F=ma}$ can be combined with ${a=\frac{\Delta v}{t}}$ to make ${F = m\frac{\Delta v}{t}}$, we can combine matrix equations $\vec{y} = M \vec{x}$ and $\vec{x} = R \vec{z}$ to make $\vec{y} = (MR)\vec{z}$. We'll see how the properties of matrices $M$ and $R$ are related to properties of $MR$ and what that all says about existence and uniqueness of solutions to the related systems of equations like equation \ref{eq:basic system}.

In Chapter \ref{chap:determinants} we'll look at \emph{determinants} which sort of describe how ``big'' a matrix is.

Chapter \ref{chap:eigenstuff} is about eigenvectors and eigenvalues, which are important for describing what happens when matrices are raised to powers: ${\vec{y} = M^t \vec{x}}$. This is kind of like a multivariable version of exponential functions like ${y = ae^{kt}}$.

A lot of the time, a solution to a system like equation \eqref{eq:basic system} doesn't exist, but you can still find a \emph{least-squares solution}, which basically means getting the right sides of the equations as close as possible to the left sides. In Chapter \ref{chap:dot-products-etc} we'll learn how to find least-squares solutions using dot products and related theory.

In Chapter \ref{chap:symmetric-matrices} we'll learn about symmetric matrices and principal component analysis. Results from this chapter are widely used in many kinds of data analysis, from signal processing to computer vision to machine learning to genetics.

Finally, Chapter \ref{chap:abstract-vector-spaces} is about how most everything we've learned about vectors actually applies to many kinds of mathematical objects such as functions, polynomials, matrices, and solutions to differential equations.

\chapter{Existence and Uniqueness of Solutions}
\label{chap:existence-and-uniqueness-of-solutions}

\section{Systems of Linear Equations}
\label{sec:systems-of-linear-eqns}

In this class, we're going to study:

\begin{itemize}
\item two-dimensional space, aka the $xy$-plane, aka $\R^2$,
\item three-dimensional space, aka $xyz$-space, aka $\R^3$,
\item and in general, $n$-dimensional space, aka $\R^n$, where we write points in the form $(x_{1}, x_2, \ldots, x_n)$.
\end{itemize}

\begin{definition} A \textbf{system of linear equations} can be written in the form:
\begin{align*}
  \begin{split}
    a_{11}x_1+a_{12}x_2+ \cdots + a_{1c}x_{c} &= y_1 \\
                                              &\vdots \\
    a_{r1}x_1+a_{r2}x_2+ \cdots + a_{rc}x_{c} &= y_r
  \end{split}
\end{align*}

where all of the \textbf{coefficients} $a_{11}, \ldots, a_{rc}$ and \textbf{constant coefficients} $y_1, \ldots, y_r$ are real numbers. A \textbf{solution} of the system is a list of numbers $(x_1, \ldots, x_c)$ that makes all of the equations true when plugged in. The set of all solutions of the system is called the \textbf{solution set} and two systems are \textbf{equivalent} if they have the same solution set.
\end{definition}

\columnbreak
\example Solve the system $\begin{aligned}
  x_1 + x_{2} &=3 \\
  -x_{1} + 2x_{2} &= 0
\end{aligned}$ by graphing both equations, and also by the Gaussian elimination method that you may remember from previous classes.

\begin{book}
\textbf{Solution.} Think of $x_1$ and $x_2$ as $x$ and $y$, respectively. To graph them, it helps to first convert them into slope-intercept form:
\begin{align*}
x_2 &= -x_1+3 \\
x_{2} &= \frac{1}{2}x_{1}
\end{align*}
\btwocol{
The first one has $x_2$-intercept 3 and slope $-1$, and the second equation has $x_2$-intercept 0 and slope $\frac{1}{2}$ as shown.}
{
\begin{center}
\begin{tikzpicture}[line cap=round,line join=round]
\begin{axis}[
x=25pt,y=25pt,
axis lines=middle,
axis line style={-},
ymajorgrids=true,
xmajorgrids=true,
xlabel=$x_1$,
ylabel=$x_2$,
xmin=-0.5, xmax=4.5,
ymin=-0.5, ymax=4.5,
xtick={-4,-3,...,5},
ytick={-3,-2,...,4},]
\addplot[red,very thick, dotted] {-x+3};
\node[right,red] at (axis cs:0,3){$x_2=-x_1+3$};
\addplot[blue,thick] {0.5*x};
\node[blue, below] at (axis cs:3.5,1.6) {$x_2=\frac{1}{2}x_1$};
\end{axis}
\end{tikzpicture}
\end{center}
}

  Each line has an infinite number of points which are solutions of that individual equation, but the only solution of the \emph{system} is $(2,1)$.

  The elimination method means adding or subtracting multiples of both sides of one equation to another equation, with the goal of eliminating some variables. Let's add equation 1 to equation 2 to replace equation 2:
\begin{align*}
    x_1 + x_{2} &=3 \\
    3x_2&= 3
\end{align*}
Then we can multiply both sides of equation 2 by $\frac{1}{3}$:
\begin{align*}
    x_1 + x_{2} &=3 \\
    x_2&= 1
\end{align*}
And finally subtract both sides of equation 2 from equation 1:
\begin{align*}
   x_1 &=2 \\
  x_2&= 1
\end{align*}
\blacksquare
\end{book}

\begin{notes}
\vfill
\end{notes}

Even though you've probably had experience with the elimination method in previous classes, you might not have thought very hard about \emph{why} it works. The proof of the next theorem goes into this in (maybe too much) depth. \emph{Fun Fact:} The first written record of Gaussian elimination comes from China, circa 150 BC. It was independently invented in several cultures. In Europe, it was first widely published by Isaac Newton and well-known before Gauss was born.

\begin{notes}
\columnbreak
\end{notes}

\begin{theorem}[Gaussian elimination doesn't change solutions.]
  \label{theorem:gauss}

  If one system of linear equations is transformed into another system using one of the operations below, then the two systems are equivalent.

  \begin{itemize}
  \item \textbf{Swapping:} Switch the positions of two equations. (aka \emph{interchange}.)
  \item \textbf{Scaling:} Multiply both sides of an equation by a non-zero real number.
  \item \textbf{Multaddiplacement:} Re\emph{place} each side of one equation by the \emph{sum} of itself and a \emph{mult}iple of another equation's sides. (aka simply \emph{replacement} in most textbooks.)
  \end{itemize}
\end{theorem}

\begin{notes}
\vfill
\columnbreak
\end{notes}

\begin{book}
\textbf{Proof:}  It's clear that swapping doesn't change the solution set, because the same equations are just written in a different order. Scaling only affects one equation, and doesn't change the solution set of that equation, just like in basic algebra class.

Multaddiplacement requires some explaining. Call the first system $S$ and the transformed system $T$, and suppose you're replacing equation 2 with the sum of itself and $k$ times equation 1. And let's label the left sides as $L$ and right sides as $R$ like so:

  \[ S: \left\{
      \begin{aligned}
        L_1 &= R_1 \\
        L_2 &= R_2 \\
            & \vdots
      \end{aligned} \right.
    \text{ transforms to } T: \left\{
      \begin{aligned}
        L_1 &= R_1 \\
        L_2 + k(L_1) &= R_2 + k(R_1) \\
            & \vdots
      \end{aligned}\right. \]

Suppose that $(s_1,\ldots,s_n)$ is a solution of $S$. This means when $(s_1,\ldots,s_n)$ is plugged in to the variables,

\[L_1\text{ is the same number as }R_1 \text{, and}\]
\[L_2\text{ is the same number as }R_2.\] So you're really just adding the same number to both sides of equation 2:
  \[L_2 + k(L_1)\text{ is the same number as }R_2 + k(R_1).\] Thus $(s_1,\ldots,s_n)$ is also solution of $T$.

Now suppose that $(t_1,\ldots,t_n)$ is a solution of $T$. This means that when $(t_1,\ldots,t_n)$ is plugged in,
  \[L_1\text{ is the same number as }R_1\text{ and}\]
  \[L_2 + k(L_1)\text{ is the same number as }R_2 + k(R_1).\]
And this also means that
  \[-k(L_1)\text{ is the same number as }-k(R_1),\]
and you can add it to both sides of equation 2:
  \[L_2+k(L_1) - k(L_1)\text{ is the same number as  }R_2+k(R_1) - k(R_1),\]
which means
  \[L_2\text{ is the same number as }R_2.\] Thus $(t_1,\ldots,t_n)$ is also a solution of $S.$

Putting it together, any solution of $S$ is also a solution of $T$ and any solution of $T$ is also a solution of $S$, which means they have the same solution set. \blacksquare
\end{book}

\example Write and sketch a system of two equations in two variables with:

\begin{enumerate}[label=(\alph*)]
\item No solution
\item An infinite number of sol'ns
\item Exactly two sol'ns
\end{enumerate}

\begin{notes}
\vfill
\end{notes}

\begin{book}
\noindent \textbf{Solution.}
\begin{enumerate}[label=(\alph*)]
\item \bbtwocol{0.4}{0.5}{\vspace*{-1ex}Any parallel, non-intersecting lines will do, such as $x_1+x_2=4$ and $x_1+x_2=2$:}
{\begin{center}
\begin{tikzpicture}[line cap=round,line join=round]
\begin{axis}[
x=25pt,y=25pt,
axis lines=middle,
axis line style={-},
ymajorgrids=true,
xmajorgrids=true,
xlabel=$x_1$,
ylabel=$x_2$,
xmin=-0.5, xmax=5.5,
ymin=-0.5, ymax=5.5,
xtick={-4,-3,...,5},
ytick={-3,-2,...,5},]
\addplot[red,very thick, dotted] {4-x};
\node[red,above,rotate=-45] at (axis cs:1,3){$x_1+x_2=4$};
\addplot[blue,thick] {2-x};
\node[blue, above, rotate=-45] at (axis cs:1,1) {$x_1+x_2=2$};
\end{axis}
\end{tikzpicture}\end{center}}

\item \bbtwocol{0.4}{0.5}{\vspace*{-1ex}The two equations must be equivalent, with the same slope and intercept, but maybe written in different forms, such as $-2x_1+x_2=-1$ and $2x_1-x_2=1$:
}{\begin{center}
\begin{tikzpicture}[line cap=round,line join=round]
\begin{axis}[
x=25pt,y=25pt,
axis lines=middle,
axis line style={-},
ymajorgrids=true,
xmajorgrids=true,
xlabel=$x_1$,
ylabel=$x_2$,
xmin=-0.5, xmax=4.5,
ymin=-0.5, ymax=4.5,
xtick={-4,-3,...,5},
ytick={-3,-2,...,5},]
\addplot[blue,thick] {-1+2*x};
\node[blue, below right, rotate=64] at (axis cs:1.2,1.5) {$2x_1-x_2=1$};
\addplot[red,very thick, dotted] {-1+2*x};
\node[red,above,rotate=64] at (axis cs:1.5,2){$-2x_1+x_2=-1$};
\end{axis}
\end{tikzpicture}
\end{center}}

\item \bbtwocol{0.4}{0.5}{\vspace*{-1ex}To intersect at exactly two points, they can't both be linear. $x_2=(x_1-4)^2$ and $x_1+x_2=6$ works.}{
\begin{center}
\begin{tikzpicture}[line cap=round,line join=round]
\begin{axis}[
x=20pt,y=20pt,
axis lines=middle,
axis line style={-},
ymajorgrids=true,
xmajorgrids=true,
xlabel=$x_1$,
ylabel=$x_2$,
xmin=-0.5, xmax=6.5,
ymin=-0.5, ymax=6.5,
xtick={-4,-3,...,6},
ytick={-3,-2,...,6},
domain=-0.5:6.5]
\addplot[red,very thick, dotted] {(x-4)*(x-4)};
\node[red,left] at (axis cs:6.2,5){$x_2=(x_1-4)^2$};
\addplot[blue,thick] {6-x};
\node[blue, above, rotate=-45] at (axis cs:3.5,2.5) {$x_1+x_2=6$};
\end{axis}
\end{tikzpicture}
\end{center}}

\end{enumerate}
\end{book}

The above example shows why any system of two linear equations with two variables has either zero solutions, one solution, or an infinite number of solutions. It turns out this is also true for any number of equations with any number of variables, which is hard to visualize, but we'll see why it's true by analyzing the systems themselves.

\begin{definition}
A system is called \textbf{consistent} if it has at least one solution; if it has exactly one solution, the solution is \textbf{unique}. A system with no solutions is called \textbf{inconsistent}.
\end{definition}

\begin{notes}
\columnbreak
\end{notes}

\example \label{example:three}
Is each system consistent? If so, is the solution unique?

\begin{horiparts}{3}
\item $\begin{aligned}
x_1 &= 5 \\
x_2 + x_3 &= 13 \\
-2x_2 - x_3 &= -19
\end{aligned}$

\item $\begin{aligned}
x_1 + 2x_3 &= 3 \\
x_2 &= 5 \\
-5x_1 - 10x_3 &= 22
\end{aligned}$

\item $\begin{aligned}
x_1 + 2x_3 &= 3 \\
x_2 &= 5 \\
-5x_1 - 10x_3 &= -15
\end{aligned}$
\end{horiparts}

\begin{notes}
\vfill \hspace{1pt}
\columnbreak
\end{notes}

\begin{book}
\textbf{Solution.} (a) The first equation already tells us what we want to know about $x_1$; it must equal 5. The second two equations have two variables $x_2$ and $x_3$. Let's use multaddiplacement to clear the $x_2$ from equation 3. The shorthand notation on the right means ``replace equation 3 with the sum of itself and 2 times equation 2.''
\[\begin{aligned}
x_1 &= 5 \\
x_2 + x_3 &= 13 \\
  x_3 &= 7
\end{aligned} \quad \quad
\begin{matrix}
\\ \\ e_3\leftarrow e_3+2e_2
\end{matrix}
\]
Now we know what $x_3$ must be, and you \emph{could} substitute 7 into $x_3$ in the second equation to solve for $x_2$. Instead, let's practice using the standard method we'll use throughout the class, and clear the $x_3$ from equation 2 using multaddiplacement:
\[\begin{aligned}
x_1 &= 5 \\
x_2 &= 6 \\
  x_3 &= 7
\end{aligned} \quad \quad
\begin{matrix}
\rule{1em}{0em}\\ e_2\leftarrow e_2-1e_3 \\ \rule{1em}{0em}
\end{matrix}
\]
This system has only one solution $(5,6,7)$, and by theorem \ref{theorem:gauss}, so does the original equation. In other words, the original system is consistent with a unique solution.

(b) There's no need to mess with the second equation ($x_2=5$) but the first and third equations make a system of two equations with two variables, so let's solve them the same way. We can clear the $-5x_1$ from equation 3 by adding 5 times equation 1 to it:
\[
\begin{aligned}
x_1 + 2x_3 &= 3 \\
x_2 &= 5 \\
0 &= 37
\end{aligned} \quad \quad 
\begin{matrix}
\rule{1em}{0em}\\ \rule{1em}{0em} \\ e_3\leftarrow e_3+5e_3 
\end{matrix}
\]
Now, no matter what values of the variables you plug in, there is no way to make it true that $0=37$. So this system (and thus the original system) has no solution; it is inconsistent. You may also notice that the first and third equations in the original system, have the same slopes but different intercepts if you graph them as lines in the $x_1x_3$-plane.

(c) This one is almost the same as the last system, except for the $-15$ instead of $22$. Doing the same step to clear the $x_1$ from equation 3 yields:
\[
\begin{aligned}
x_1 + 2x_3 &= 3 \\
x_2 &= 5 \\
0 &= 0
\end{aligned} \quad \quad 
\begin{matrix}
\rule{1em}{0em}\\ \rule{1em}{0em} \\ e_3\leftarrow e_3+5e_3 
\end{matrix}
\]
Just like no values of the variables could make $0=37$ \emph{true}, no values of the variable can make $0=0$ \emph{false}. So any solution of the first two equations solves the whole system. We know that $x_2$ must be $5$, but $(x_1,x_3)$ can be any point on the line $x_1+2x_3=3$. So this system (and the original system) has an infinite number of solutions such as $(3,5,0)$, $(1,5,1)$, $(-1,5,2)$, and $(-3,5,3)$. In other words, the system is consistent but the solution is \emph{not unique}. \blacksquare
\end{book}

Here's a rough strategy for Gaussian elimination that we'll be using throughout the class:

\begin{enumerate}
\item If the top equation doesn't have $x_1$, use swaps to get it there first. Scale the top equation so the coefficient of $x_1$ is $1$. 
\item Use multaddiplacements to clear all the $x_1$ terms from the other equations.
\item If the second equation doesn't have $x_2$, use swaps to get it there. Scale the second equation so the coefficient of $x_2$ is $1$. (If $x_2$ isn't in any of the equations below equation 1, skip this and the next step.)
\item Use multaddiplacements to clear all the $x_2$ terms from the other equations.
\item Repeat with the rest of the variables until each variable shows up in only one equation, if possible.
\end{enumerate}

Note that in Example \ref{example:three} (b) and (c), it's not possible to get $x_3$ in an equation by itself because $x_3$ got canceled out along with $x_2$. And then it's not possible to clear out the $x_3$ from the first equation. In general, when there's an infinite number of solutions or no solution, you can't get each equation to have a single variable, but the strategy above will get you as close as possible.

\example
Solve the system:
$\begin{aligned}
x_1 - 2x_2 &= -7 \\
-2x_1 + 4x_2 + x_3 &= 18 \\
x_1 - 2x_2 + x_3 &= -3
\end{aligned}$

\begin{book}
\textbf{Solution.} Let's follow the above steps. The first equation already has $x_1$ so we can proceed to step 2, and use multaddiplacements to clear the $x_1$'s from equations 2 and 3. 
\begin{align*}
&\begin{aligned}
x_1 - 2x_2 &= -7 \\
   x_3 &= 4 \\
 x_3 &= 4
\end{aligned} &
\begin{aligned}
&\\
e_2 &\leftarrow e_2+2e_1 \\
e_3&\leftarrow e_3-1e_1
\end{aligned}
\end{align*}
Now that $x_2$ isn't in any of the equations below the first, we skip steps 3 and 4, and move on to clear $x_3$ out of all equations except the second using a multaddiplacement:
\begin{align*}
&\begin{aligned}
x_1 - 2x_2 &= -7 \\
   x_3 &= 4 \\
 0 &= 0
\end{aligned} &
\begin{aligned}
&\\
&\\
e_3&\leftarrow e_3-2e_1
\end{aligned}
\end{align*}

Now, we can see there is only one value that $x_3$ can be, but there are an infinite number of $(x_1,x_2)$ solutions to the first equation. So the whole system has an infinite number of solutions, such as $(-7,0,3)$, $(-5,1,3)$, and $(-3,2,3)$. \blacksquare
\end{book}


\columnbreak
\exerciseheading

\noindent For exercises \ref{exercise:bob} through \ref{exercise:fred}, use Gaussian elimination to determine if the system is consistent or not, and if so, if the solution is unique.

\exatend \label{exercise:bob} $\begin{aligned}[t]
2x_1 + 6x_2 &= 18 \\
-2x_1 - 5x_2 &= -14
\end{aligned}$

\exatend  $\begin{aligned}[t]
x_1 + x_2 &= 4 \\
-2x_1 - 2x_2 &= 0
\end{aligned}$

\exatend $\begin{aligned}[t]
x_1 + 2x_2 &= 0 \\
3x_1 + 4x_2 &= 0
\end{aligned}$

\exatend $\begin{aligned}[t]
x_1 + 2x_2 &= 0 \\
3x_1 + 6x_2 &= 0
\end{aligned}$

\exatend $\begin{aligned}[t]
-2x_1 +0x_2 - 7x_3 &= 1 \\
-2x_1 +0x_2 - 6x_3 &= -2
\end{aligned}$

\exatend \label{exercise:fred} $\begin{aligned}[t]
2x_1 + x_2 &= -3 \\
2x_1 - x_2 &= 11 \\
x_1 - x_2 &= 9
\end{aligned}$

\noindent For exercises \ref{exercise:trish} through \ref{exercise:flo}, Is the statement true or false? Justify your answer. If true, either reference the definition/theorem/paragraph where it's stated in the text, or explain why it's true. If false, maybe the statement is an incorrect version of something in the text which you can reference; or you can give a counterexample that shows it's not always true. (These instructions apply to T/F problems throughout the book.)

\exatend \label{exercise:trish} (T/F?) Suppose you use Gaussian elimination to reduce a system of two equations and get $\begin{aligned}
x_1 &= 1 \\
x_2 &= 2
\end{aligned}$. This equation has two solutions.

\exatend (T/F?) A \emph{solution} of a system of equations is a list of numbers $(x_1, \ldots, x_n)$ that makes at least one of the equations true when plugged into the variables.

\exatend (T/F?) Each of the elimination operations (swapping, scaling, multaddiplacement), can be reversed using an elimination operation.


\exatend (T/F?) If two systems of linear equations have the same solution sets, then they're consistent with each other.

\exatend  (T/F?) If two linear systems have the same solution set, then they're equivalent to each other.

\exatend \label{exercise:flo} (T/F?) If you scale an equation in a system by multiplying both sides of the equation by zero, the resulting system is equivalent to the first one.


\noindent For exercises \ref{exercise:harderelim} through \ref{exercise:harderelimend}, use Gaussian elimination to determine if the system is consistent or not, and if so, determine if the solution is unique. If the system has a unique solution, find it.

\exatend \label{exercise:harderelim} $\begin{aligned}[t] %unique soln
x_1 - 3x_2 + 5x_3 &= 23 \\
-3x_2 + 6x_3 &= 21 \\
x_3 &= 8
\end{aligned}$

\exatend $\begin{aligned}[t] %no soln
4x_1 - x_2 + 3x_3 &= 4\pi \\
-3x_1 + x_2 - 3x_3 &= -3\pi \\
-4x_1 + 2x_2 - 6x_3 &= 2 - 4\pi
\end{aligned}$

\exatend $\begin{aligned}[t] %infinite number of solns
x_1 + 2x_3 &= 3 \\
x_2 - 4x_4 &= 0 \\
x_3 + 2x_4 &= -5 \\
2x_1 - 3x_2 + 5x_3 + 14x_4 &= 1
\end{aligned}$

\exatend $\begin{aligned}[t] %infinite number of solns
x_1 + 2x_2 &= 3 \\
4x_1 + 8x_2 &= 12 \\
-3x_1 - 6x_2 &= -9
\end{aligned}$ \label{exercise:harderelimend}

\exatend (T/F?) A system of three linear equations with two variables is always inconsistent. \emph{(Remember to always justify T/F answers.)}

\exatend (T/F?) A system of two linear equations with three variables is always consistent with an infinite number of solutions.

\exatend (T/F?) If a system of two linear equations with three variables is consistent, then it has an infinite number of solutions.

\exatend (T/F?) If a system of linear equations has all of the constant coefficients (on the right sides of the equal sign) equal to zero, then the system is consistent.

\begin{minipage}{\linewidth}
\exatend Find the value(s) of $p$ that make the following system consistent: $\begin{aligned}
x_1 + x_2 + 5x_3 &= -4 \\
-3x_1 - 2x_2 - 13x_3 &= 8 \\
-7x_1 - 5x_2 - 31x_3 &= p
\end{aligned}$
\end{minipage}

\exatend What can you conclude about $c$ and $d$ if the following system has an infinite number of solutions:
$\begin{aligned}[t]
x_1 + 2x_2 &= 3 \\
cx_1 + dx_2 &= 6
\end{aligned}$

\exatend What can you conclude about $c$ and $d$ if the following system has a unique solution?
$\begin{aligned}[t]
x_1 + 2x_2 &= 3 \\
cx_1 + dx_2 &= 6
\end{aligned}$

\exatend What can you conclude about $c$ and $d$ if the following system is inconsistent?
$\begin{aligned}[t]
x_1 + 2x_2 &= 3 \\
cx_1 + dx_2 &= 6
\end{aligned}$

\exatend What can you conclude about $a,b,c$ and $d$ if the following system is inconsistent?
$\begin{aligned}[t]
ax_1 + bx_2 &= 0 \\
cx_1 + dx_2 &= 5
\end{aligned}$

\columnbreak
\section{Row Echelon Form and Pivots}
\label{sec:row-echelon-form-and-pivots}

\exercise (\emph{Warm-up}) Use Gaussian elimination to determine if the system is consistent or not, and if so, determine if the solution is unique. If the system has a unique solution, find it.

\begin{enumerate}[label=(\alph*)]
\item $\begin{aligned}[t]
2x_1 + 4x_2 + 6x_3 &= 8 \\
x_1 + 2x_2 + 3x_3 &= -5
\end{aligned}$

\vspace{3em}
\item $\begin{aligned}[t]
2x_1 + 4x_2 + 6x_3 &= 8 \\
x_1 + 2x_2 + 3x_3 &= 4
\end{aligned}$
\vspace{3em}
\item $\begin{aligned}[t]
x_1 - 2x_2 + x_3 &= 0 \\
x_2 - x_3 &= -1 \\
x_3 &= 6
\end{aligned}$
\end{enumerate}

\vspace{3em}

\begin{definition}
The \textbf{coefficient matrix} of a system of $r$ linear equations with $c$ variables is an array with $r$ rows and $c$ columns, where each entry in the matrix corresponds with the coefficient of a variable in the equation. Each row of the matrix represents a row of the equation and each column of the matrix represents a variable.


So a system of linear equations
\begin{align*}
  \begin{split}
    a_{11}x_1+a_{12}x_2+ \cdots + a_{1c}x_{c} &= y_1 \\
                                              &\vdots \\
    a_{r1}x_1+a_{r2}x_2+ \cdots + a_{rc}x_{c} &= y_r
  \end{split}
\end{align*}
has coefficient matrix
\[
\begin{bmatrix}
a_{11} & a_{12} & \cdots & a_{1c} \\
\vdots & \vdots & \vdots & \vdots \\
a_{r1} & a_{r2} & \cdots & a_{rc}
\end{bmatrix}
\]
The \textbf{augmented matrix} of a system has an extra column on the right with the constant coefficients, like so:
\[
\begin{bmatrix}[cccc|c]
a_{11} & a_{12} & \cdots & a_{1c} & y_1 \\
\vdots & \vdots & \vdots & \vdots & \vdots \\
a_{r1} & a_{r2} & \cdots & a_{rc} & y_{r}
\end{bmatrix}
\]
If a matrix has $r$ rows and $c$ columns, we say its \textbf{size} is $r \times c$. So the size of this augmented matrix is $r \times (c+1)$.
\end{definition}

\example For the system in exercise 1(c), write the augmented matrix of the original system, and the augmented matrix of the reduced system that tells you the solution. What is the size of the augmented matrix?

\begin{book} \textbf{Solution.} The original system was 
$\begin{aligned}
x_1 - 2x_2 + x_3 &= 0 \\
0x_1+ x_2 - x_3 &= -1 \\
0x_1+0x_2+x_3 &= 6
\end{aligned}$ so its augmented matrix is $\begin{bmatrix}[ccc|c]
1 & -2 & 1 & 0\\
0 & 1 & -1 & -1 \\
0 & 0 & 1 & 6
\end{bmatrix}$. 
The system showing the solution was $\begin{aligned}
x_1 + 0x_2 + 0x_3 &= 4 \\
0x_1+ x_2 + 0x_3 &= 5 \\
0x_1+0x_2+x_3 &= 6
\end{aligned}$ so its augmented matrix is $\begin{bmatrix}[ccc|c]
1 & 0 & 0 & 4\\
0 & 1 & 0 & 5 \\
0 & 0 & 1 & 6
\end{bmatrix}$. Both augmented matrices have 3 rows and 4 columns, so their size is $3\times 4$. \blacksquare
\end{book}

\begin{notes}
\columnbreak
\end{notes}

\begin{definition}
A \textbf{leading entry} of a row in a matrix is the leftmost non-zero entry in the row.

We say a matrix is in \textbf{row echelon form} (or simply \textbf{echelon form}) if it satisfies:
\begin{enumerate}[label=(\roman*)]
\item Any rows with all zeros are below rows with non-zero entries.
\item Every leading entry is to the right of the leading entry in row above. (This doesn't apply to the top row.)
\end{enumerate}
A matrix is in \textbf{reduced row echelon form} if it's in row echelon form and also satisfies:
\begin{enumerate}[label=(\roman*)] \setcounter{enumi}{2}
\item Every leading entry is the only non-zero entry in its column.
\item Every leading entry is 1.
\end{enumerate}

The Gaussian elimination operations from theorem \ref{theorem:gauss} also have matching \textbf{row operations}:
\begin{itemize}
\item \textbf{Swapping:} Switch the positions of two rows. (aka \emph{interchange}.)
\item \textbf{Scaling:} Multiply all entries in a row by a non-zero real number.
\item \textbf{Multaddiplacement:} Re\emph{place} a row by the \emph{sum} of itself and a \emph{mult}iple of another row. (aka simply \emph{replacement} in most textbooks.)
\end{itemize}
We say that two matrices are \textbf{row-equivalent} if one can be transformed to the other by a series of these row operations.
\end{definition}
By theorem \ref{theorem:gauss}, if the augmented matrices of two systems are row-equivalent, then the two systems are equivalent (have the same solution sets).

\begin{notes}
\example Circle which of the following matrices are in row echelon form, and double-circle the ones that are in reduced row echelon form.

$\begin{bmatrix}
-5 & 7 \\
0 & -3 \\
0 & 0
\end{bmatrix} \quad \quad\quad
\begin{bmatrix}
1 & 0 & 2  \\
0 & 1 & -7 
\end{bmatrix} \quad \quad \quad
\begin{bmatrix}
1 & 0 & 0 \\
1 & 0 & 0  
\end{bmatrix} \quad \quad \quad
\begin{bmatrix}
0 & 7 & 12 \\
0 & 0 & -3 \\
0 & 0 & 0 \\
0 & 0 & 0
\end{bmatrix} $
\vspace{2em}

$\begin{bmatrix}
1 & 0 & 0 \\
0 & 1 & 0 \\
0 & 0 & 1
\end{bmatrix} \quad \quad\quad
\begin{bmatrix}
1 & 4 & 0 & 0 \\
0 & 0 & 1 & 3 
\end{bmatrix} \quad \quad\quad
\begin{bmatrix}
1 & 2 & 3
\end{bmatrix} \quad \quad\quad
\begin{bmatrix}
0 & 0 & 1 \\
0 & 1 & 0 \\
1 & 0 & 0
\end{bmatrix}$


\vfill
\columnbreak
\end{notes}

\begin{book}
\example Here are some matrices that are not in row echelon form:
\[\begin{bmatrix}
1 & 0 & 0 \\
1 & 0 & 0  
\end{bmatrix} \quad \quad \begin{bmatrix}
0 & 0 & 1 \\
0 & 1 & 0 \\
1 & 0 & 0
\end{bmatrix}\]
Here are some matrices in row echelon form:
\[\begin{bmatrix}
-5 & 7 \\
0 & -3 \\
0 & 0
\end{bmatrix} \quad
\begin{bmatrix}
1 & 0 & 2  \\
0 & 1 & -7 
\end{bmatrix} \quad
\begin{bmatrix}
0 & 7 & 12 \\
0 & 0 & -3 \\
0 & 0 & 0 \\
0 & 0 & 0
\end{bmatrix} \quad
\begin{bmatrix}
1 & 0 & 0 \\
0 & 1 & 0 \\
0 & 0 & 1
\end{bmatrix} \quad
\begin{bmatrix}
1 & 4 & 0 & 0 \\
0 & 0 & 1 & 3 
\end{bmatrix} \quad
\begin{bmatrix}
1 & 2 & 3
\end{bmatrix} 
\]
And the following are in reduced row echelon form:
\[\begin{bmatrix}
1 & 0 & 2  \\
0 & 1 & -7 
\end{bmatrix} \quad
\begin{bmatrix}
1 & 0 & 0 \\
0 & 1 & 0 \\
0 & 0 & 1
\end{bmatrix} \quad
\begin{bmatrix}
1 & 4 & 0 & 0 \\
0 & 0 & 1 & 3 
\end{bmatrix} \quad
\begin{bmatrix}
1 & 2 & 3
\end{bmatrix} \quad \blacksquare
\]
\end{book}

\begin{definition}
If there is a reduced row echelon matrix equivalent to $A$ with a leading entry in row $i$, column $j$, we say $i,j$ is a \textbf{pivot position} of $A$, and column $j$ is a \textbf{pivot column}.
\end{definition}
\emph{Note:} As we'll see in theorem \ref{theorem:uniqueness-of-rref}, every matrix is row-equivalent to a \emph{unique} reduced row echelon matrix, which means we don't have to worry about situations where, say $2,2$ is a leading entry in one reduced row echelon form of $A$, but not a leading entry in another reduced row echelon form of $A$.

\example Suppose each of the following are augmented matrices for a system. Is each consistent, and if so, how many solutions does it have?

\begin{horiparts}{3}
\item $\begin{bmatrix}[cc|c]
1 & 0 & 2  \\
0 & 1 & -7 
\end{bmatrix}$
\item $\begin{bmatrix}[c|c]
-5 & 7 \\
0 & -3 \\
0 & 0
\end{bmatrix}$
\item $\begin{bmatrix}[ccc|c]
1 & 4 & 0 & 0 \\
0 & 0 & 1 & 3 
\end{bmatrix}$
\end{horiparts}

\begin{book}
\textbf{Solution.} (a) This augmented matrix represents the system $\begin{aligned}
x_1&=2 \\
x_2&=-7 
\end{aligned}$, which is consistent with the unique solution $(2,-7)$.

(b) The system for this augmented matrix is $\begin{aligned}
-5x_1&=7 \\
0 &= -3 \\
0 &= 0
\end{aligned}$ which is inconsistent because no value of $x_1$ can make the equation $0=-3$ true.

(c) This represents the system $\begin{aligned}
x_1+4x_2&=0 \\
x_3&=3
\end{aligned}$. There's only one value of $x_3$ that works, but there are an infinite number of $(x_1,x_2)$ pairs that make the first equation true. Some solutions of the system are $(0,0,3),(-4,1,3)$, and $(-8,2,3)$. So it's consistent with an infinite number of solutions. \blacksquare
\end{book}

\begin{notes}
\vfill \hspace{1pt}
\columnbreak
\end{notes}

\begin{definition}
Variables which correspond to pivot columns are called \textbf{basic variables}. Any non-basic variables are called \textbf{free variables}. When you solve for the basic variables in terms of the free variables, you have a \emph{parametric} description of the solution set. Any choice of the free variables determines a different solution.
\end{definition}

\example Reduce the following augmented matrix for a system to reduced row echelon form, and describe the solution set.
\[\begin{bmatrix}[ccc|c]
-3 & 12 & 1 & -5 \\
-4 & 16 & 1 & -5 \\
-1 & 4 & 1 & -5
\end{bmatrix}\]

\begin{book}
\textbf{Solution.} To avoid having to multiply by $\frac{1}{3}$, let's swap rows 1 and 3:
\[\begin{bmatrix}
-1 & 4 & 1 & -5 \\
-4 & 16 & 1 & -5 \\
-3 & 12 & 1 & -5 
\end{bmatrix}\]
Now we can use easier multaddiplacements to clear out the first column and proceed from there:
\begin{align*}
&\begin{bmatrix}
-1 & 4 & 1 & -5 \\
0 & 0 & -3 & 15 \\
0 & 0 & -2 & 10 
\end{bmatrix} &\begin{matrix}
\rule{0em}{1em} \\ r_2 \leftarrow r_2 -4r_1 \\ r_3 \leftarrow r_3-3r_1
\end{matrix} \\
&\begin{bmatrix}
1 & -4 & -1 & 5 \\
0 & 0 & 1 & -5 \\
0 & 0 & 1 & -5 
\end{bmatrix} &\begin{matrix}
r_1 \leftarrow -1r_{1} \\ r_2 \leftarrow \frac{-1}{3}r_2 \\ r_3 \leftarrow \frac{-1}{2}r_3
\end{matrix} \\
&\begin{bmatrix}
1 & -4 & 0 & 0 \\
0 & 0 & 1 & -5 \\
0 & 0 & 0 & 0 
\end{bmatrix} &\begin{matrix}
r_1 \leftarrow r_1+1r_2 \\ \rule{0em}{1em} \\ r_3 \leftarrow r_3 -1r_2
\end{matrix} 
\end{align*}
Now that we're in reduced row echelon form, and see that columns 1 and 3 have pivots, so $x_1$ and $x_3$ are the basic variables, and $x_2$ is a free variable. So we can write all the variables in terms of $x_2$. Let's start by noticing that row 1, which is the equation $x_1-4x_2=0$, can be written $x_1=4x_2$. So any solution satisfies:
\[\begin{aligned}
x_1&=4x_2 \\
x_2&=x_2 \\
x_3&=-5
\end{aligned}\]
Some examples of solutions are $(0,0,-5), (4,1,-5),$ or $(8,2,-5)$. Equivalently we could write that the solution set is:
\[\left\{ (4x_2,x_2,-5) : x_2 \text{ is in } \R\right\} \]
This \emph{set-builder notation} stands for ``the set of all points of the form $(4x_2,x_2,-5)$ such that $x_2$ is in $\R$.'' \blacksquare
\end{book}

\begin{notes}
\vfill
\end{notes}

\begin{theorem}[Existence and uniqueness basics.]
\label{theorem:existence-and-uniqueness-basics}
\begin{enumerate}[label=(\roman*)]
\item A system is consistent if and only if the augmented matrix does \emph{not} have a  pivot in the rightmost column.

\item If the system is consistent and there is at least one non-pivot column in the coefficient matrix, then there are an infinite number of solutions.

\item If the system is consistent and there are no non-pivot columns in the coefficient matrix, then the solution is unique.
\end{enumerate}
\end{theorem}

\begin{notes}
\columnbreak
\end{notes}

\begin{book}
\textbf{Proof:} (i) (\rightarrow) Suppose a system is consistent. Then there's no pivot in the rightmost column, because if there were, the system would be equivalent to a system with a ``$0 = \text{nonzero}$'' equation and it wouldn't be consistent.

(\leftarrow) Suppose a system has no pivot in the rightmost column. Then, using the Gaussian elimination strategy in section \ref{sec:systems-of-linear-eqns} the system's augmented matrix can be reduced to a row-equivalent reduced echelon matrix with no leading entry in the rightmost column. This equivalent system has a solution, since you can solve for the basic variables in terms of any free variables (if they exist), and plug any real number into the free variables to get a solution.

(ii) Suppose the system is consistent. If there is at least one non-pivot column in the coefficient matrix, then the system is equivalent to a reduced system with free variable(s). You can solve for the basic variables in terms of the free variables. There are an infinite number of values of the free variables that each lead to a different solution, so there are an infinite number of solutions.

% This is worded carefully to avoid non-uniqueness or rref weirdness. If anyone asks, define "non-pivot column" as a column that is not a pivot in SOME rref form.
(iii) If there are no non-pivot columns in the coefficient matrix, then any reduced echelon matrix equivalent to the coefficient matrix has a pivot in every column, so every column of a reduced coefficient matrix has only a single 1 in it, and zeros everywhere else. So every non-zero row of the reduced echelon \emph{augmented} matrix represents an equation in the form ``$x_i= \text{a number}$'' which means the system has only one solution.
\end{book}

\exerciseheading

For exercises \ref{exercise:1.2.0b} through \ref{exercise:1.2.0e}, is the matrix in row echelon form? What about reduced row echelon form? If it's not in reduced row echelon form, use Gaussian elimination to make it so.

\exatend \label{exercise:1.2.0b} $
\begin{bmatrix}
1 & 2 & 3 \\
0 & 2 & 0
\end{bmatrix}
$

\exatend $
\begin{bmatrix}
1 & 8 & \pi \\
0 & 0 & 0 
\end{bmatrix}
$

\exatend $
\begin{bmatrix}
1 & 0 & 0 & 8 \\
0 & 0 & 1 & 6 \\
0 & 1 & 0 & 7
\end{bmatrix}
$

\exatend $
\begin{bmatrix}
2 & 3 \\
6 & 9 \\
-2 & -3 \\
\end{bmatrix}
$

\exatend $\begin{bmatrix}
0 & 1 & 2 & 4 \\
-2 & 3 & 0 & 26
\end{bmatrix}$

\exatend \label{exercise:1.2.0e}
$\begin{bmatrix}
1 & 2 & 0 & 5 \\
1 & 3 & 0 & 9 \\
-2 & -4 & 0 & -7
\end{bmatrix}$

For exercises \ref{exercise:1.2.1b} through \ref{exercise:1.2.1e}, an augmented matrix of a system is given. Solve the system (either write \emph{no solution}; or if has a unique solution, write it; or if it has an infinite number of solutions, write the parametric form).

\exatend \label{exercise:1.2.1b} $
\begin{bmatrix}[ccc|c]
1 & 0 & 0 & 8 \\
0 & 0 & 1 & 6 \\
0 & 1 & 0 & 7
\end{bmatrix}
$

\exatend $
\begin{bmatrix}[cc|c]
1 & 0 & 0 \\
0 & 0 & 1  \\
0 & 1 & 0 
\end{bmatrix}
$

\exatend  $
\begin{bmatrix}[ccc|c]
1 & 2 & 3 & 0
\end{bmatrix}
$

\exatend $\begin{bmatrix}[ccc|c]
-2 & 4 & -4 & -6 \\
2 & -4 & 5 & 11 \\
-3 & 6 & -6 & -11
\end{bmatrix}$

\exatend $\begin{bmatrix}[ccc|c]
-2 & 4 & -4 & -6 \\
2 & -4 & 5 & 11 \\
-3 & 6 & -6 & -9
\end{bmatrix}$

\exatend \label{exercise:1.2.1e} $\begin{bmatrix}[cccc|c]
1 & 1 & 5 & 1 & 0 \\
-2 & -1 & -7 & -2 & 0 \\
-1 & -1 & -5 & 2 & 0
\end{bmatrix}$

Remember to give some justification for all T/F answers.

\exatend (T/F?) A system of 2 equations with 4 variables has a $4 \times 2$ coefficient matrix.

\exatend (T/F?) A system of 2 equations with 4 variables has a $2 \times 5$ augmented matrix.

\exatend (T/F?) Every matrix in reduced row echelon form is also in row echelon form.

\exatend (T/F?) One example of a valid row operation is replacing row 2 with row 2 minus row 2, resulting in a row of all zeros.

\exatend (T/F?) If there's a leading entry in column 3 of an echelon form coefficient matrix, then $x_3$ is a basic variable.

\exatend (T/F?) If an augmented matrix has a row equal to $
\begin{bmatrix}
0 & 0 & 0 & -2
\end{bmatrix}$, then the system is inconsistent.

\exatend (T/F?) If a coefficient matrix has non-pivot columns, then the corresponding system has an infinite number of solutions.

\exatend (T/F?) If a coefficient matrix has non-pivot columns, and the rightmost column of the corresponding augmented matrix is not a pivot column, then the system has an infinite number of solutions.

\exatend Find all the values of $y_{1}$ and $y_2$ so that the system with this augmented matrix is inconsistent, if possible: $
\begin{bmatrix}[cc|c]
-1 & 2 & y_1 \\
3 & -6 & y_2
\end{bmatrix}
$

\exatend Find all the values of $y_1, y_2,$ and $y_3$ so that the system with this augmented matrix is inconsistent, if possible: $\begin{bmatrix}[ccc|c]
-1 & -3 & -4 & y_1\\
2 & 5 & 7 & y_2\\
0 & -1 & -2 & y_3
\end{bmatrix}
$

\exatend  Find all the values of $y_1, y_2,$ and $y_3$ so that the system with this augmented matrix is inconsistent, if possible: $\begin{bmatrix}[ccc|c]
-1 & -3 & -4 & y_1\\
2 & 5 & 7 & y_2\\
0 & -1 & -1 & y_3
\end{bmatrix}
$

\exatend Show that scaling (multiplying) a row by zero can change the solution set using an example.

\exatend Give a $3 \times 3$ coefficient matrix $A$, not in echelon form, so that the system with augmented matrix $
\begin{bmatrix}[ccc|c]
 & & & y_1 \\
 & A & & y_2 \\
 & & & y_3 \\
\end{bmatrix}$ is consistent for any values of $y_1,y_2,$ and $y_3$.

\exatend Give a $3 \times 3$ coefficent matrix $A$, not in echelon form, so that the system with augmented matrix $
\begin{bmatrix}[ccc|c]
 & & & y_1 \\
 & A & & y_2 \\
 & & & y_3 \\
\end{bmatrix}$ is consistent for some values of $y_1,y_2,$ and $y_3$, and inconsistent for others. Give an example of values $(y_1,y_2,y_3)$ that makes it consistent and also an example $(y_1,y_2,y_3)$ that makes it inconsistent.

\exatend Find values of $a,b$, and $c$ so that $f(x) = c + bx + ax^2$ has $f(0) = -3, f(1) = -2,$ and $f(2) = -3$.

\columnbreak
\section{Vectors, Linear Combinations, and Span}
\label{sec:vectors-linear-combinations-span}

Suppose $(x_1, \ldots, x_c)$ is any point in $\R^c$. We will often represent such a point as a $c \times 1$ matrix $\bbbvec{x_1}{\vdots}{x_c}$, call it a \textbf{vector}, and write it as $\vec{x}$. In fact, we'll use the terms \emph{point} and \emph{vector} interchangeably, and the notations $(x_1, \ldots, x_c)$ and $\bbbvec{x_1}{\vdots}{x_c}$ interchangeably. Sometimes it will be helpful to think of a vector as any arrow with the same magnitude and direction as the arrow from $(0, \ldots, 0)$ to the point $(x_1, \ldots, x_c)$.

\begin{tikzpicture}
\begin{axis}[
x=30pt,y=30pt,
axis lines=middle,
axis line style={-},
ymajorgrids=true,
%minor y tick num=3,
%yminorgrids=true,
xmajorgrids=true,
%minor x tick num=3,
%xminorgrids=true,
%xlabel=$x_1$,
%ylabel=$x_2$,
xmin=-4, xmax=4,
ymin=-1, ymax=4,
xtick={-4,-3,...,4},
ytick={-1,0,...,4},
xticklabels={},
yticklabels={},
]
\draw[-Stealth, very thick] (axis cs:0,0) -- (axis cs:1,3)
node[right] {$(1,3) \text{ aka } \bbvec{1}{3}$};
\addplot[only marks, mark=*, mark size=2pt, very thick] coordinates {(1,3)};
\draw[-Stealth, very thick] (axis cs:-3.5,1) -- (axis cs:-2.5,4);
\draw[right] (axis cs:-3,2.5) node {$(1,3) \text{ aka } \bbvec{1}{3}$};
\draw[right] (axis cs:0.5,1.5) node {$(1,3) \text{ aka } \bbvec{1}{3}$};
\end{axis}
\end{tikzpicture}

But, don't go overboard: we will never write a vector/point as a single-row matrix like $
\begin{bmatrix}
x_1 & \ldots & x_c
\end{bmatrix}$.

Vectors can be scaled (multiplied) by any real number, or added together like so:

\[k\bbbvec{u_1}{\vdots}{u_c} = \bbbvec{ku_1}{\vdots}{ku_c} \quad \text{ and } \quad \bbbvec{u_1}{\vdots}{u_c} + \bbbvec{v_1}{\vdots}{v_c} = \bbbvec{u_1+v_1}{\vdots}{u_c+v_c} \]

It's important to understand what these operations look like with points/arrows:

\columnbreak
\example Suppose $\vec{w}=\bbvec{-1}{2}$ and $\vec{z}=\bbvec{3}{1}$. Find and graph:
\begin{enumerate}[label=(\alph*)]
\item $\vec{w}$ and $\vec{z}$
\item $\vec{w}+\vec{z}$
\item $k \vec{w}$ for $k=-2,-1,0,0.5,2$
\end{enumerate}

\begin{notes}
\begin{tikzpicture}
\begin{axis}[
x=30pt,y=30pt,
axis lines=middle,
axis line style={-},
ymajorgrids=true,
%minor y tick num=3,
%yminorgrids=true,
xmajorgrids=true,
%minor x tick num=3,
%xminorgrids=true,
%xlabel=$x_1$,
%ylabel=$x_2$,
xmin=-2, xmax=4,
ymin=-4, ymax=4,
xtick={-2,-1,...,4},
ytick={-4,-3,...,4},
%xticklabels={},
%yticklabels={},
]
\end{axis}
\end{tikzpicture}
\end{notes}

\begin{book}\bbtwocol{.45}{.55}{
\textbf{Solution.} (a) See the graph.

(b) $\bbvec{-1}{2}+\bbvec{3}{1} = \bbvec{-1+3}{2+1}=\bbvec{2}{3}$

(c)
\begin{align*}
-2 \bbvec{-1}{2} &= \bbvec{2}{-4}\\
-1 \bbvec{-1}{2} &= \bbvec{1}{-2}\\
0 \bbvec{-1}{2} &= \bbvec{0}{0}\\
0.5 \bbvec{-1}{2} &= \bbvec{-0.5}{1}\\
2 \bbvec{-1}{2} &= \bbvec{-2}{4}
\end{align*}
}{
\begin{tikzpicture}
\begin{axis}[
x=30pt,y=30pt,
axis lines=middle,
axis line style={-},
ymajorgrids=true,
%minor y tick num=3,
%yminorgrids=true,
xmajorgrids=true,
%minor x tick num=3,
%xminorgrids=true,
%xlabel=$x_1$,
%ylabel=$x_2$,
xmin=-2, xmax=4,
ymin=-4, ymax=4,
xtick={-2,-1,...,4},
ytick={-4,-3,...,4},
%xticklabels={},
%yticklabels={},
]
\draw[-Stealth, thick] (axis cs:0,0) -- (axis cs:-1,2) node[left]{$\vec{w}$};
\draw[-Stealth, thick] (axis cs:0,0) -- (axis cs:3,1) node[right]{$\vec{z}$};
\draw[-Stealth, thick, dashed] (axis cs:-1,2) -- (axis cs:2,3) node[midway,above]{$+\vec{z}$};
\draw[-Stealth, thick, dashed] (axis cs:3,1) -- (axis cs:2,3) node[midway,right]{$+\vec{w}$};
\addplot[only marks, mark=*, mark size=2pt, very thick] coordinates {(2,3)} node[right]{$\vec{w}+\vec{z}$};
\addplot[only marks, mark=*, mark size=2pt, very thick] coordinates {(-2,4)} node[below right]{$2\vec{w}$};
\addplot[only marks, mark=*, mark size=2pt, very thick] coordinates {(-0.5,1)} node[left]{$0.5\vec{w}$};
\addplot[only marks, mark=*, mark size=2pt, very thick] coordinates {(0,0)} node[below left]{$0\vec{w}$};
\addplot[only marks, mark=*, mark size=2pt, very thick] coordinates {(1,-2)} node[above]{$-1\vec{w}$};
\addplot[only marks, mark=*, mark size=2pt, very thick] coordinates {(2,-4)} node[above]{$-2\vec{w}$};
\end{axis}
\end{tikzpicture}
}
\end{book}

\begin{theorem}[Properties of vector operations.]
\label{properties-of-vector-operations}
For all vectors $\vec{x}, \vec{y}$, and $\vec{z}$ in $\R^c$ and all scalars $r$ and $t$:
\begin{enumerate}[label=(\roman*)]
\item $\vec{x}+ \vec{y} = \vec{y} + \vec{x}$ \quad (vector addition is \textbf{commutative})
\item $(\vec{x} + \vec{y}) + \vec{z} = \vec{x} + (\vec{y} + \vec{z})$ \quad (vector addition is \textbf{associative})
\item $\vec{x} + \vec{0} = \vec{0} + \vec{x} = \vec{x}$
\item $\vec{x} + (- \vec{x}) = \vec{0}$
\item $(r+t) \vec{x} = r \vec{x} + t \vec{x}$
\item $r(\vec{x} + \vec{y}) = r\vec{x} + r\vec{y}$
\item $r(t\vec{x}) = (rt)\vec{x}$
\end{enumerate}
\end{theorem}

\begin{notes}
\columnbreak
\end{notes}
Each of these is true because of a corresponding property of real numbers which can be applied to each entry of the vectors. Note that we will study mathematical objects in this class where some of these rules are not true.

Let's prove (i) and leave some others for exercises.

\begin{notes}
\vfill
\end{notes}

\begin{book}
\textbf{Proof.} (i) Suppose $\vec{x}=\bbbvec{x_1}{\vdots}{x_c}$ and $\vec{y}=\bbbvec{y_1}{\vdots}{y_c}$ are any two vectors in $\R^c$. Then:
\begin{align*}\vec{x}+\vec{y} &= \bbbvec{x_1}{\vdots}{x_c} + \bbbvec{y_1}{\vdots}{y_c} \\
&= \bbbvec{x_1+y_1}{\vdots}{x_c+y_c} &\text{(Def'n of vector addition)} \\
&= \bbbvec{y_1+x_1}{\vdots}{y_c+x_c} &\text{(Commutative property of real number addition)}\\
&= \bbbvec{y_1}{\vdots}{y_c} + \bbbvec{x_1}{\vdots}{x_c} &\text{(Def'n of vector addition)}\\
&= \vec{y}+\vec{x} & \blacksquare
\end{align*}

\end{book}

\begin{definition} A \textbf{list} of vectors is a finite set of vectors, that all have the same number of entries, written in some order. Sometimes we'll write a list with parentheses, sometimes without. So $\vec{a}_{1}, \vec{a}_2, \vec{a}_3$ means the same thing as $(\vec{a}_{1}, \vec{a}_2, \vec{a}_c)$ but not the same thing as $\vec{a}_3,\vec{a}_2,\vec{a}_1$.

If $\vec{a}_1,\vec{a}_2, \ldots, \vec{a}_c$ is a list of vectors in $\R^n$ and $x_{1}, x_2, \ldots, x_c$ are in $\R$, then 
\[ x_1\vec{a}_1 + x_2\vec{a}_2 + \cdots + x_c\vec{a}_c\]
is called a \textbf{linear combination} of $\vec{a}_1,\vec{a}_2, \ldots, \vec{a}_c$ with coefficients (aka weights) $x_{1}, x_2, \ldots, x_c$.
\end{definition}

\begin{notes}
\columnbreak
\end{notes}

\example Given $\vec{u} = \bbvec{-1}{2}$ and $\vec{v}=\bbvec{2}{1}$, graph some vectors in the set:
\[ \left\{ t\vec{u} + r\vec{v} : t,r \text{ are integers} \right\}\]
(Read this as ``the set of all vectors of the form  $t\vec{u} + r\vec{v}$ such that $t,r$ are any integers.'')

\begin{notes}
\begin{tikzpicture}[line cap=round,line join=round]%,x=2.0cm,y=1.0cm]
\begin{axis}[
x=25pt,y=25pt,
axis lines=middle,
axis line style={-},
ymajorgrids=true,
xmajorgrids=true,
xlabel=$x_1$,
ylabel=$x_2$,
xmin=-6, xmax=6,
ymin=-4, ymax=4,
xtick={-6,-5,...,6},
ytick={-4,-3,...,4},]
\end{axis}
\end{tikzpicture}

\columnbreak
\end{notes}

\begin{book}
\textbf{Solution.} Below are some vectors in the set, color coded to make it easier to find them on the graph. First lets note that $t$ can be $0$ which gives multiples of $\vec{v}$:
\[\textcolor{red}{0\vec{u}+0\vec{v} = \bbvec{0}{0}\quad 0\vec{u}+1\vec{v} = \bbvec{2}{1} \quad 0\vec{u}+2\vec{v} = \bbvec{4}{2} \quad 0\vec{u}-1\vec{v} = \bbvec{-2}{-1}} \]
Likewise $r$ can be zero giving:
\[ \textcolor{blue}{-1\vec{u}+0\vec{v} = \bbvec{1}{-2} \quad 1\vec{u}+0\vec{v} = \bbvec{-1}{2} \quad }\]
Some other examples of vectors in the set:
\[\textcolor{orange}{-1\vec{u}+2\vec{v}= \bbvec{5}{0}} \quad \textcolor{violet}{1\vec{u}-1\vec{v} = \bbvec{-3}{1}}\]
All other vectors in the set that fit in the graph are labeled as black dots.

\begin{tikzpicture}[line cap=round,line join=round]%,x=2.0cm,y=1.0cm]
\begin{axis}[
x=25pt,y=25pt,
axis lines=middle,
axis line style={-},
ymajorgrids=true,
xmajorgrids=true,
xlabel=$x_1$,
ylabel=$x_2$,
xmin=-6, xmax=6,
ymin=-4, ymax=4,
xtick={-6,-5,...,6},
ytick={-4,-3,...,4},]
\addplot[red, only marks, mark=*, mark size=2pt, very thick] coordinates {(0,0)};
\addplot[red, only marks, mark=*, mark size=2pt, very thick] coordinates {(2,1)};
\addplot[red, only marks, mark=*, mark size=2pt, very thick] coordinates {(4,2)};
\addplot[red, only marks, mark=*, mark size=2pt, very thick] coordinates {(-2,-1)};
\addplot[red, only marks, mark=*, mark size=2pt, very thick] coordinates {(-4,-2)};
\draw[red,-Stealth, thick] (axis cs:0,0) -- (axis cs:2,1) node[above]{$\vec{v}$};
\addplot[blue, only marks, mark=*, mark size=2pt, very thick] coordinates {(-1,2)};
\addplot[blue, only marks, mark=*, mark size=2pt, very thick] coordinates {(1,-2)};
\draw[blue,-Stealth, thick] (axis cs:0,0) -- (axis cs:-1,2) node[above]{$\vec{u}$};
\addplot[orange, only marks, mark=*, mark size=2pt, very thick] coordinates {(5,0)};
\addplot[purple, only marks, mark=*, mark size=2pt, very thick] coordinates {(-3,1)};
\addplot[only marks, mark=*, mark size=2pt, very thick] coordinates {(-4,3)};
\addplot[only marks, mark=*, mark size=2pt, very thick] coordinates {(-1,-3)};
\addplot[only marks, mark=*, mark size=2pt, very thick] coordinates {(3,-1)};
\addplot[only marks, mark=*, mark size=2pt, very thick] coordinates {(4,-3)};
\addplot[only marks, mark=*, mark size=2pt, very thick] coordinates {(1,3)};
\addplot[only marks, mark=*, mark size=2pt, very thick] coordinates {(-5,0)};
\addplot[only marks, mark=*, mark size=2pt, very thick] coordinates {(-6,2)};
\addplot[only marks, mark=*, mark size=2pt, very thick] coordinates {(-3,-4)};
\addplot[red, only marks, mark=*, mark size=2pt, very thick] coordinates {(-6,-3)};
\addplot[blue, only marks, mark=*, mark size=2pt, very thick] coordinates {(-2,4)};
\addplot[only marks, mark=*, mark size=2pt, very thick] coordinates {(3,4)};
\addplot[red, only marks, mark=*, mark size=2pt, very thick] coordinates {(6,3)};
\addplot[blue, only marks, mark=*, mark size=2pt, very thick] coordinates {(2,-4)};
\addplot[only marks, mark=*, mark size=2pt, very thick] coordinates {(6,-2)};
\end{axis}
\end{tikzpicture}

\end{book}

\example Find coefficients that make $\bbvec{2}{2}$ a linear combination of $\vec{u} = \bbvec{-1}{2}$ and $\vec{v}=\bbvec{2}{1}$. First use the graph to make an estimate, then set up and solve a system of equations.

\begin{notes}
\vfill
\end{notes}

\begin{book}
\textbf{Solution.} It looks like you can get to $(2,2)$ with a little more than $1$ times $\vec{v}$ and a little less than half of $\vec{u}$. So a decent guess is that $0.4\vec{u} + 1.2\vec{v} = \bbvec{2}{2}$. This turns out to be exactly right, but we won't always be so lucky. Let's see how to do it analytically. We are looking for $t$ and $r$ that make the following vector equation true:
\[t \bbvec{-1}{2} + r \bbvec{2}{1} = \bbvec{2}{2}\]
Doing the scalar multiplication and vector addition on the left gives:
\[\bbvec{-1t +2r}{2t+1r} = \bbvec{2}{2}\]
These two vectors are equal if and only if both of their components are equal. So the vector equation is equivalent to the system of equations
\[\begin{aligned}
-1t+2r &= 2 \\
2t +1r &= 2
\end{aligned}\]
which we can solve by doing row operations with the augmented matrix:
\begin{align*}
&\begin{bmatrix}[cc|c]
-1 & 2 & 2 \\
2 & 1 & 2
\end{bmatrix} \\
&\begin{bmatrix}[cc|c]
-1 & 2 & 2 \\
0 & 5 & 6
\end{bmatrix} &\begin{matrix}
\rule{0em}{1em} \\ r_2 \leftarrow r_2+2r_1 \end{matrix} \\
&\begin{bmatrix}[cc|c]
1 & -2 & -2 \\
0 & 1 & 1.2
\end{bmatrix} &\begin{matrix}r_1\leftarrow -1r_1 \\ r_2 \leftarrow \frac{1}{5}r_2 \end{matrix} \\
&\begin{bmatrix}[cc|c]
1 & 0 & 0.4 \\
0 & 1 & 1.2
\end{bmatrix} &\begin{matrix} r_1\leftarrow r_1+2r_2 \\ \rule{1em}{0pt} \end{matrix} 
\end{align*}
Translating this into a system of equations, we have the solution: $t=0.4, r=1.2$. \blacksquare
\end{book}

\begin{definition}
The \textbf{span} of $\vec{a}_1, \ldots, \vec{a}_c$, written $\spn ( \vec{a}_1, \ldots, \vec{a}_c )$, is the set of \emph{all} linear combinations of $\vec{a}_1, \ldots, \vec{a}_c$.

The span of the empty list $()$ is defined to be $\left\{ \vec{0} \right\}$.
\end{definition}

\begin{notes}
\columnbreak
\end{notes}

\example Let $\vec{u}=\bbvec{3}{1}$ and $\vec{v}=\bbvec{-6}{-2}$.
\begin{enumerate}[label=(\alph*)]
\item Graph $\spn(\vec{u},\vec{v})$.
\begin{notes}
\vspace{12em}
\end{notes}
\item How many vectors are in the list $( \vec{u},\vec{v} )$?
\item How many vectors are in $\spn(\vec{u},\vec{v})$?
\item Is $\spn(\vec{u},\vec{v}) = \R^2$?
\begin{notes}
\vspace{2em}
\end{notes}
\item For what value(s) of $k$ is $\bbvec{k}{5}$ in $\spn(\vec{u},\vec{v})$?
\begin{notes}
\vspace{8em}
\end{notes}
\item How many solutions does $x_1\vec{u}+x_2\vec{v}=\vec{0}$ have?
\end{enumerate}

\begin{notes}
\columnbreak
\end{notes}

\begin{book}
\textbf{Solution.} (a) $\spn(\vec{u},\vec{v})$ includes vectors like:
\[-2\vec{u}+0\vec{v}=\bbvec{-6}{-2}, \quad 0\vec{u}+\frac{-1}{2}\vec{v}=\bbvec{3}{1}, \quad 0\vec{u}+0\vec{v}=\vec{0}, \quad 0.123\vec{u}+0.2\vec{v}=\bbvec{-0.831}{-.277}\]
All of these are points on the line that includes $\vec{u},\vec{v},$ and $\vec{0}$. And, any point on this line can be written as a multiple of $\vec{u}$ or $\vec{v}$. So $\spn(\vec{u},\vec{v}) = \left\{ (x_1,x_2) : x_2=\frac{1}{3}x_1 \right\}$.

\begin{tikzpicture}[line cap=round,line join=round]%,x=2.0cm,y=1.0cm]
\begin{axis}[
x=20pt,y=20pt,
axis lines=middle,
axis line style={-},
ymajorgrids=true,
xmajorgrids=true,
xmin=-7, xmax=4,
ymin=-3, ymax=2,
xtick={-7,-6,...,6},
ytick={-4,-3,...,4},]
\draw[thick] (axis cs:-9,-3) -- (axis cs:6,2);
\draw[-Stealth,thick] (axis cs:0,0) -- (axis cs:-6,-2) node[above]{$\vec{v}$};
\draw[-Stealth,thick] (axis cs:0,0) -- (axis cs:3,1) node[above]{$\vec{u}$};
\end{axis}
\end{tikzpicture}

(b) The list $(\vec{u},\vec{v})$ contains two vectors: $\vec{u},$  and $\vec{v}$.

(c) The set $\spn(\vec{u},\vec{v})$ contains an infinite number of vectors, like $\vec{u},\vec{v}, \vec{0}, \pi \vec{v}, 17\vec{u},$ etc.

(d) No, there are points in $\R^2$, such as $(0,1)$, that can't be written as a linear combination of $\vec{u}$ and $\vec{v}$.

(e) We want to see what values of $k$ make it so $x_1\vec{u}+x_2\vec{v} = \bbvec{k}{5}$ has a solution. We can write this vector equation several ways:
\begin{align*}
x_1 \bbvec{3}{1} + x_2 \bbvec{-6}{-2} &= \bbvec{k}{5} \\
\bbvec{3x_1}{x_1} + \bbvec{-6x_2}{-2x_2} &= \bbvec{k}{5} \\
\bbvec{3x_1-6x_2}{x_1-2x_2} &= \bbvec{k}{5}
\end{align*}
These vectors are equal if and only if both components are equal, so we have the system with the following augmented matrix which we can reduce:
\begin{align*}
&\begin{bmatrix}[cc|c]
3 & -6 & k \\
1 & -2 & 5
\end{bmatrix} & \hspace{1pt} \\
&\begin{bmatrix}[cc|c]
1 & -2 & 5 \\
3 & -6 & k 
\end{bmatrix} &r_1 \leftrightarrow r_2 \\
&\begin{bmatrix}[cc|c]
1 & -2 & 5 \\
0 & 0 & k-15
\end{bmatrix} &r_2 \leftarrow r_2-3r_1
\end{align*}
This has a solution if and only if $k-15$ is zero. In other words, $\bbvec{k}{5}$ can be written as a linear combination of $\vec{u}$ and $\vec{v}$ if and only if $k=15$. This agrees with the graph; the point (15,5) is on the line through $(3,1)$ and $(0,0)$.

(f) Looking at the above graph, you can see that $(0,0)$ can be written as $2\vec{u}+\vec{v}$, or as $\vec{u}+0.5\vec{v}$, or as $0\vec{u}+0\vec{v}$. So intuitively, it looks like $x_1\vec{u}+x_2\vec{v}=\vec{0}$ has an infinite number of solutions, but let's check, writing it as a system of equations:
\begin{align*}
x_1 \bbvec{3}{1} + x_2 \bbvec{-6}{-2} &= \bbvec{0}{0} \\
\bbvec{3x_1-6x_2}{x_1-2x_2} &= \bbvec{0}{0} 
\end{align*}
This is equivalent to the system of equations with augmented matrix
$\begin{bmatrix} [cc|c]
3 & -6 & 0 \\
1 & -2 & 0
\end{bmatrix}$ which can handily be reduced to $\begin{bmatrix} [cc|c]
1 & -2 & 0 \\
0 & 0 & 0 \\
\end{bmatrix}$. There is no pivot in the rightmost column, so the system is consistent. Since $x_2$ is a free variable, there are an infinite number of solutions. \blacksquare

\end{book}

\example Let $\vec{u}=\bbvec{2}{-1}$ and $\vec{v}=\bbvec{4}{1}$. 
\begin{enumerate}[label=(\alph*)]
\item Show that $\spn(\vec{u},\vec{v}) = \R^2$. (In this case we say the list $(\vec{u},\vec{v})$ \emph{spans} $\R^2$.)
\begin{notes}
\vspace{20em}
\end{notes}
\item How many solutions does $x_1\vec{u}+x_2\vec{v}=\vec{0}$ have?
\end{enumerate}

\begin{notes}
\vspace{12em}
\columnbreak
\end{notes}

\begin{book}
\textbf{Solution.} (a) We want to show that \emph{every} vector in $\R^2$ can be written $x_1\vec{u}+x_2 \vec{v}$ for some $x_1,x_2$. In other words, we want to show that for any $(y_1,y_2)$, the following equation has a $(x_1,x_2)$ solution:
\[x_1 \bbvec{2}{-1} + x_2 \bbvec{4}{1} = \bbvec{y_1}{y_2}\]
We can write this as
\[\bbvec{2x_1+4x_2}{-1x_1+1x_2} = \bbvec{y_1}{y_2}\]
which is equivalent to the system of equations:
\[\begin{aligned}
2x_1+4x_2&=y_1 \\
-1x_1+1x_2&=y_{2}
\end{aligned}\]
So let's row-reduce the augmented matrix:
\begin{align*}
&\begin{bmatrix}[cc|c]
2 & 4 & y_1 \\
-1 & 1 & y_2
\end{bmatrix} \\
&\begin{bmatrix}[cc|c]
-1 & 1 & y_2 \\
0 & 6 & y_1 +2y_{2}
\end{bmatrix} &\begin{matrix} r_1 \leftrightarrow r_2 \\ r_2 \leftarrow r_2+2r_1 \end{matrix}
\end{align*}
This is now in echelon form, and there's no pivot in the rightmost column, no matter what $y_1$ and $y_2$ are. Theorem \ref{theorem:existence-and-uniqueness-basics} says the system is always consistent; thus any $(y_1,y_2)$ vector can be written $x_1\vec{u}+x_2\vec{v}$.

(b) Just like with (a), we can write $x_1\vec{u}+x_2\vec{v}= \vec{0}$ in the equivalent forms below, and row reduce the augmented matrix.
\begin{align*}
x_1 \bbvec{2}{-1} + x_2 \bbvec{4}{1} = \bbvec{0}{0} \\
\bbvec{2x_1+4x_2}{-1x_1+1x_2} = \bbvec{0}{0} \\
\begin{bmatrix}[cc|c]
2 & 4 & 0 \\
-1 & 1 & 0
\end{bmatrix} \\
\begin{bmatrix}[cc|c]
-1 & 1 & 0 \\
0 & 6 & 0 
\end{bmatrix} \\
\begin{bmatrix}[cc|c]
1 & 0 & 0 \\
0 & 1 & 0 
\end{bmatrix}
\end{align*}
So the original vector equation $x_1\vec{u}+x_2\vec{v}= \vec{0}$ is equivalent to $x_1=0, x_2=0$, which has a unique solution. \blacksquare
\end{book}

\example Let $\vec{u} = \bbbvec{1}{3}{0}$ and $\vec{v}=\bbbvec{0}{0}{2}$.
\begin{enumerate}[label=(\alph*)]
\item Graph $\spn ( \vec{u},\vec{v} )$.
\begin{notes}
\vspace{12em}
\end{notes}
\item For what value(s) of $k$ is $\vec{y}=\bbbvec{h}{2}{4}$ in $\spn( \vec{u},\vec{v} )$?
\begin{notes}
\vspace{12em}
\end{notes}
\item How many solutions does $x_1\vec{u} + x_2\vec{v} = \vec{0}$ have?
\end{enumerate}

\begin{book}
\btwocol{\textbf{Solution.} (a) The span of $(\vec{u},\vec{v})$ consists of all points in $\R^3$ that you can get to by moving any distance in the $\vec{u}$ direction (including backwards), plus any distance in the $\vec{v}$ direction. All these points make up a \emph{plane}, which includes the points $\vec{u},\vec{v}$, and $\vec{0}$, as shown.}{
\begin{tikzpicture}[line cap = round, scale=0.7, x={(-0.7cm,-0.7cm)}, y={(1cm,0cm)}, z={(0cm,1cm)}]
\draw (0,0,0) -- (3.5,0,0) node[left]{$x_1$};
\draw (0,0,0) -- (0,4.5,0) node[anchor=north west]{$x_2$};
\draw (0,0,0) -- (0,0,3.5) node[right]{$x_3$};
% Draw tick marks on x-axis
\foreach \x in {,1,2,3} {\draw (\x,0,0) -- (\x,0.1,0) node[right] {\scriptsize $\x$};}
% Draw tick marks on y-axis
\foreach \y in {1,2,3,4} {\draw (0,\y,0) -- (0.1,\y,0) node[below left] {\scriptsize $\y$};}
% Draw tick marks on z-axis
\foreach \z in {1,2,3} {\draw (0,0,\z) -- (0,0.1,\z) node[right] {\scriptsize $\z$};}
\node[below left,rotate=-20,blue] at (1.5,4.5,3) {$\spn(\vec{u},\vec{v})$};
\fill[blue!30,opacity=0.5] (-1.5,-4.5,-3) -- (-1.5,-4.5,3) -- (1.5,4.5,3) -- (1.5,4.5,-3) -- cycle;
\draw[-Stealth,thick] (0,0,0) -- (1,3,0) node[below] {$\vec{u}$};
\draw[-Stealth,thick] (0,0,0) -- (0,0,2) node[left] {$\vec{v}$};
\end{tikzpicture}}

(b) The vector $\vec{y}$ is in $\spn(\vec{u},\vec{v})$ if and only if there exists an $x_1,x_2$ solution of the vector equation
\[x_1 \bbbvec{1}{3}{0} + x_2 \bbbvec{0}{0}{2} = \bbbvec{h}{2}{4}\]
and like the previous examples, we can write this in the form
\[\bbbvec{1x_1+0x_2}{3x_1+0x_2}{0x_1+2x_2} = \bbbvec{h}{2}{4}\]
which is equivalent to a system with an augmented matrix that we can row-reduce:
\begin{align*}
&\begin{bmatrix}[cc|c]
1 & 0 & h \\
3 & 0 & 2 \\
0 & 2 & 4
\end{bmatrix} \\
&\begin{bmatrix}[cc|c]
1 & 0 & h \\
0 & 0 & 2-3h \\
0 & 1 & 2
\end{bmatrix} &\begin{matrix} \rule{1em}{0pt} \\ r_2 \leftarrow r_2-3r_1 \\ r_3\leftarrow 0.5r_3\end{matrix} \\
&\begin{bmatrix}[cc|c]
1 & 0 & h \\
0 & 1 & 2 \\
0 & 0 & 2-3h
\end{bmatrix} &\begin{matrix} \rule{1em}{0pt} \\ r_2 \leftrightarrow r_3 \\ \rule{1em}{0pt}\end{matrix} 
\end{align*}
This is now in echelon form, and represents a consistent system if and only if $2-3h=0$. Thus $\vec{y}$ is in $\spn(\vec{u},\vec{v})$ if and only if $h=\frac{2}{3}$.

(c) Like before, the equation $x_1\vec{u}+x_2\vec{v}=\vec{0}$ is equivalent to the system of equations with the following augmented matrix which can be reduced like so:
\begin{align*}
&\begin{bmatrix}[cc|c]
1 & 0 & 0 \\
3 & 0 & 0 \\
0 & 2 & 0
\end{bmatrix} \\
&\begin{bmatrix}[cc|c]
1 & 0 & 0 \\
0 & 0 & 0 \\
0 & 1 & 0
\end{bmatrix} &\begin{matrix} \rule{1em}{0pt} \\ r_2 \leftarrow r_2-3r_1 \\ r_3\leftarrow 0.5r_3\end{matrix} \\
&\begin{bmatrix}[cc|c]
1 & 0 & 0 \\
0 & 1 & 0 \\
0 & 0 & 0
\end{bmatrix} &\begin{matrix} \rule{1em}{0pt} \\ r_2 \leftrightarrow r_3 \\ \rule{1em}{0pt}\end{matrix} 
\end{align*}
This represents the system of equations $\begin{aligned}
x_1=0 \\ x_2=0
\end{aligned}$ which has only one solution. You could also cite theorem \ref{theorem:existence-and-uniqueness-basics} since there are no free variables: there's a pivot position in each of the $x_1$ and $x_2$ columns. \blacksquare
\end{book}


\exerciseheading

\exatend Simplify $-3 \bbvec{0}{-7}-2 \bbvec{4}{3}$.

\exatend Simplify $3 \bbbvec{3}{2}{1} -2 \bbbvec{0}{-5}{2} - \bbbvec{9}{16}{-1}$.

\btwocol{\exatend \label{exercise:1.3.1} Let $\vec{u}=\bbvec{1}{2}$ and $\vec{v}=\bbvec{3}{-4}$. Find and graph $\vec{u}, \vec{v}, \vec{u}+ \vec{v}, \vec{u}-\vec{v}, \vec{v}- \vec{u}$, and $k \vec{u}$ for $k=-2,-1,0,1,2$.}{
\begin{tikzpicture}
\begin{axis}[
x=23pt,y=23pt,
axis lines=middle,
axis line style={-},
ymajorgrids=true,
%minor y tick num=3,
%yminorgrids=true,
xmajorgrids=true,
%minor x tick num=3,
%xminorgrids=true,
%xlabel=$x_1$,
%ylabel=$x_2$,
xmin=-3, xmax=4,
ymin=-4, ymax=6,
xtick={-3,-2,...,3},
ytick={-3,-2,...,5},
%xticklabels={},
%yticklabels={},
]
\end{axis}
\end{tikzpicture}}

\exatend Using $\vec{u}$ and $\vec{v}$ from exercise \ref{exercise:1.3.1}, use the graph to estimate solutions $(x_1,x_2)$ to each vector equation:
\begin{horiparts}{2}
\item $x_1 \vec{u}+ x_2 \vec{v} = \bbvec{2.5}{0}$
\item  $x_1 \vec{u}+ x_2 \vec{v} = \bbvec{2}{-1}$
\end{horiparts}

\exatend  Using $\vec{u}$ and $\vec{v}$ from exercise \ref{exercise:1.3.1}, translate each vector equation into a system of equations and solve:

\begin{horiparts}{2}
\item $x_1 \vec{u}+ x_2 \vec{v} = \bbvec{2.5}{0}$
\item  $x_1 \vec{u}+ x_2 \vec{v} = \bbvec{2}{-1}$
\end{horiparts}

%\columnbreak
\exatend The graph below shows $\vec{u}, \vec{v}$, and other points. Use it to solve each vector equation for $(x_1, x_2)$:

\begin{horiparts}{2}
\item $x_1\vec{u}+x_2\vec{v} = \vec{a}$
\item $x_1\vec{u}+x_2\vec{v} = \vec{b}$
\item $x_1\vec{u}+x_2\vec{v} = \vec{c}$
\item $x_1\vec{u}+x_2\vec{v} = \vec{d}$
\item $x_1\vec{u}+x_2\vec{v} = \vec{e}$
\end{horiparts}

\begin{tikzpicture}[scale=0.7]
\vectorgrid[xmin=-3,xmax=9,ymin=-3,ymax=7,step=1,labint=1,xu=1,yu=3,xv=-1,yv=2]
\fill (-1,2) circle (3pt) node[left] {$\vec{u }$};

\fill (1,3) circle (3pt) node[right] {$\vec{v}$};
\fill (-2,-1) circle (3pt) node[right] {$\vec{a}$};
\fill (2,6) circle (3pt) node[right] {$\vec{b}$};
\fill (0,0) circle (3pt) node[right] {$\vec{c}$};
\fill (5,0) circle (3pt) node[right] {$\vec{d}$};
\fill (4,4.5) circle (3pt) node[right] {$\vec{e}$};
\end{tikzpicture}

\exatend \label{exercise:1.3estimate3} Use the graph below (and no calculation) to estimate three different solutions of $x_1\vec{u}+x_2\vec{v}+x_3\vec{w} = \vec{0}$.

\begin{tikzpicture}
\begin{axis}[
x=25pt,y=25pt,
axis lines=middle,
axis line style={-},
ymajorgrids=true,
%minor y tick num=3,
%yminorgrids=true,
xmajorgrids=true,
%minor x tick num=3,
%xminorgrids=true,
%xlabel=$x_1$,
%ylabel=$x_2$,
xmin=-4, xmax=4,
ymin=-3, ymax=3,
xtick={-4,-3,...,4},
ytick={-3,-2,...,3},
xticklabels={},
yticklabels={},
]
\draw[-Stealth, very thick] (axis cs:0,0) -- (axis cs:-2,2)
node[above] {$\vec{u}$};
\draw[-Stealth, very thick] (axis cs:0,0) -- (axis cs:1,1)
node[right] {$\vec{v}$};
\draw[-Stealth, very thick] (axis cs:0,0) -- (axis cs:2,0)
node[below] {$\vec{w}$};
\end{axis}
\end{tikzpicture}

%\columnbreak
For exercises \ref{exercise:1.3solveb} to \ref{exercise:1.3solvee}, convert the vector equation to a system of equations, write its augmented matrix, and solve it (either write \emph{no solution}; or if has a unique solution, write it; or if it has an infinite number of solutions, write the parametric form).

\exatend \label{exercise:1.3solveb} $ x_1\bbvec{5}{-3} + x_2\bbvec{-11}{7} = \bbvec{-4}{4}$
\exatend $x_1 \bbvec{-2}{4} + x_2 \bbvec{-6}{13} + x_3 \bbvec{2}{0} = \bbvec{0}{7}$

\exatend $x_1 \bbvec{1}{2} + x_2 \bbvec{4}{8} = \bbvec{1}{7}$
\exatend $x_1 \bbvec{1}{2} + x_2 \bbvec{4}{8} = \bbvec{1}{2}$
\exatend $x_1 \bbvec{-2}{2} + x_2 \bbvec{1}{1} + x_3 \bbvec{2}{0} = \bbvec{0}{0}$ (Note: these are the vectors from exercise \ref{exercise:1.3estimate3}. Does your answer agree with that answer?)
\exatend \label{exercise:1.3frodo} $ x_1 \bbbvec{1}{2}{-2} + x_2\bbbvec{0.5}{1}{-1} + x_3 \bbbvec{0}{1}{1} = \bbbvec{-2}{1}{9} $
\exatend \label{exercise:1.3bilbo} $ x_1 \bbbvec{1}{2}{-2} + x_2\bbbvec{0.5}{1}{-1} + x_3 \bbbvec{0}{1}{1} = \bbbvec{-2}{1}{12} $

\exatend $ x_1 \bbbvec{-2}{5}{-1} +x_2 \bbbvec{6}{-14}{3} +x_3 \bbbvec{-3}{11}{-2} = \bbbvec{8}{-25}{4} $
\exatend \label{exercise:1.3solvee} $ x_1 \bbbvec{-2}{5}{-1} +x_2 \bbbvec{6}{-14}{3} +x_3 \bbbvec{-3}{11}{-2} = \bbbvec{0}{0}{0} $

\exatend Write four vectors that are in $\spn ( \bbbvec{1}{2}{3}, \bbbvec{0}{10}{0} )$.


\exatend Let $M$ be the list $\bbvec{5}{-3},\bbvec{-11}{7}$. (Hint: see your answer to exercise \ref{exercise:1.3solveb}.)
\begin{enumerate}[label=(\alph*)]
\item Is $\vec{y} = \bbvec{-4}{4}$ in $M$? How many vectors are in $M$?
\item Is $\vec{y} = \bbvec{-4}{4}$ in $\spn ( M )$? How many vectors are in $\spn ( M )$?
\end{enumerate}

\exatend Let $W$ be the list  $\left( \bbbvec{1}{2}{-2}, \bbbvec{0.5}{1}{-1}, \bbbvec{0}{1}{1}  \right) $. 
\begin{enumerate}[label=(\alph*)]
\item Does $W$ span $\R^3$?  (ie, is $\spn ( W ) = \R^3$?) Justify your answer. (Hint: see exercise \ref{exercise:1.3bilbo}.)
\item For what values of $k$ is $\bbbvec{k}{3}{3}$ in $\spn ( W )$?
\end{enumerate}

Remember to justify any true/false answers, maybe by citing something from the section, or writing your own reasoning, or giving a counterexample.
\exatend (T/F?) If $\vec{a} = \bbvec{1}{2}$ and $\vec{b}=\bbvec{3}{4}$ then $2\vec{a}+\vec{b}=
\begin{bmatrix}
2 & 3 \\
4 & 4
\end{bmatrix}
$
\exatend (T/F?) $\begin{bmatrix}1 & 2 & 3\end{bmatrix}$ is another way to write  the vector $\bbbvec{1}{2}{3}$.
\exatend (T/F?) If $\vec{u} \ne \vec{0}$, then scalar multiples of $\vec{u}$ are always points on the line through the origin and $\vec{u}$.
\exatend (T/F?) The system of linear equations with augmented matrix $
\begin{bmatrix}[cc|c]
\vec{a}_1 & \vec{a}_2 & \vec{y}
\end{bmatrix}$ is consistent if and only if $\vec{y}$ is in $\spn ( \vec{a}_1, \vec{a}_2 )$.
\exatend (T/F?) If $\vec{u} \ne \vec{0}$, then $\spn ( \vec{u} )$ contains an infinite number of vectors.
\exatend (T/F?) If $\vec{u} = \vec{0}$, then $\spn ( \vec{u} )$ has an infinite number of vectors.
\exatend (T/F?) If $\vec{u},\vec{v}$ are any two vectors in $\R^2$, then the list $( \vec{u}, \vec{v} )$ spans $\R^2$.
\exatend (T/F?) For any vectors $\vec{a}_1, \ldots, \vec{a}_c$, the origin is always contained in $\spn ( \vec{a}_1, \ldots, \vec{a}_{c} )$.


\exatend Prove theorem \ref{properties-of-vector-operations} (ii).
\exatend Prove theorem \ref{properties-of-vector-operations} (v).
\exatend Prove theorem \ref{properties-of-vector-operations} (vii).

\exatend Give an example of two vectors $\vec{u}$ and $\vec{v}$ in $\R^3$, where $\spn ( \vec{u},\vec{v} )$ is \emph{not} a plane.

\exatend Suppose $\bbvec{y_1}{y_2}$ is in $\spn \left( \bbvec{2}{1}, \bbvec{6}{3} \right) $. What can you conclude about $y_1$ and $y_2$?

\exatend Suppose $\spn \left( \bbvec{2}{1}, \bbvec{6}{k} \right) = \R^2$. What can you conclude about $k$?

%inspired by LADR 2.A#1 -- difficult
\exatend Suppose $\vec{u}_1, \vec{u}_2, \vec{u}_3$ spans $\R^{3}$. Prove that the list $\vec{u}_1-\vec{u}_2, \vec{u}_2-\vec{u}_3, \vec{u}_3 $ also spans $\R^3$.

\columnbreak
\section{Homogeneous Systems and Linear Independence}
\label{sec:homogeneous-systems-and-linear-independence}

\exercise (\emph{Warm-up}) Write a vector equation that is equivalent to the system. (Don't solve, just change the notation.)

$\begin{aligned}[t]
4x_1 - x_2 + 3x_3 &= 4\ \\
-3x_1 + x_2 - 3x_3 &= -3 \\
-4x_1 + 2x_2 - 6x_3 &= 2 
\end{aligned}$

\exercise (\emph{Warm-up}) Choose the correct geometric description of $\spn (  \vec{v}_1,\vec{v}_2 )$, where $\vec{v}_1 = \bbbvec{3}{0}{2}$ and $\vec{v}_2=\bbbvec{-2}{0}{3}$.
\begin{enumerate}[label=(\Alph*)]
  \item The two points $(3,0,2)$ and $(-2,0,3)$.
  \item The line through the two points $(3,0,2)$ and $(-2,0,3)$.
  \item The plane through the three points $(0,0,0)$, $(3,0,2)$, and $(-2,0,3)$.
  \item All of $\R^3$.
\end{enumerate}

\exercise (\emph{Warm-up}) Choose the correct geometric description of $\spn ( \vec{v}_1 )$, where $\vec{v}_1=\bbvec{1}{2}$.
\begin{enumerate}[label=(\Alph*)]
  \item The point $(1,2)$.
  \item The line through $(0,0)$ and $(1,2)$.
  \item All of $\R^2$.
\end{enumerate}

\exercise (\emph{Warm-up}) Choose the correct geometric description of $\spn (\vec{v}_1,\vec{v}_2)$, where $\vec{v}_1=\bbvec{1}{2}$ and $\vec{v}_2=\bbvec{3}{6}$.
\begin{enumerate}[label=(\Alph*)]
  \item The points $(1,2)$ and $(3,6)$.
  \item The line through $(0,0)$ and $(1,2)$.
  \item All of $\R^2$.
\end{enumerate}

\columnbreak
\exercise (\emph{Warm-up}) \label{exercise:1.4dependence1} Suppose that $x_1 \vec{a}_1 + x_2 \vec{a}_2 + x_3 \vec{a}_3 = \vec{0}$ and that $x_3 \ne 0$. Write $\vec{a}_3$ as a linear combination of $\vec{a}_1$ and $\vec{a}_2$.

\begin{notes}
\vfill
\end{notes}
\exercise (\emph{Warm-up}) \label{exercise:1.4dependence2} Suppose that $y_1 \vec{a}_1 + y_2 \vec{a}_2 = \vec{a}_3$. Write $\vec{0}$ as a linear combination of $\vec{a}_1, \vec{a}_2$, and $\vec{a_3}$.

\begin{notes}
\vfill
\end{notes}
The last two exercises show that for any vectors $\vec{a}_1, \vec{a}_2$, and $\vec{a_3}$, the following two statements are equivalent:
\begin{itemize}
  \item It is possible to write $\vec{a}_3$ as a linear combination of the other two vectors.
  \item There exist numbers $x_1,x_2,x_3$, with $x_3 \ne 0$, so that $x_1 \vec{a}_1 + x_2 \vec{a}_2 + x_3 \vec{a}_3 = \vec{0}$.
\end{itemize}

\subsection{Linear Independence: Main Ideas}\label{subsec:linear-independence-main-ideas}

\begin{definition}
A linear system of equations that can be written in the form 
\[ x_1 \vec{a}_1 + \cdots + x_c\vec{a}_c = \vec{0}\]
is called \textbf{homogeneous}. Homogeneous systems are always consistent since $\vec{x} = ( x_1, \ldots, x_c )=\vec{0}$ is always a solution, called the \textbf{trivial solution.} A \textbf{non-trivial} solution is any solution $\vec{x} \ne \vec{0}$. 

When a non-trivial solution exists, we say the list $\vec{a}_1, \ldots, \vec{a}_c$ is \textbf{linearly dependent}. When the only solution is the trivial solution, we say the list is \textbf{linearly independent.} The empty list $()$ is defined to be linearly independent as well.

If $A =
\begin{bmatrix}
\vec{a}_1 & \cdots & \vec{a}_c
\end{bmatrix}
$, the solution set of $ x_1 \vec{a}_1 + \cdots + x_c\vec{a}_c = \vec{0}$ is called the \textbf{null space} of $A$ and is written $\nul A$.
\end{definition}

\begin{notes}
\columnbreak
\end{notes}

\example \label{example:1.4first-nul} Let $A =\begin{bmatrix}
\vec{a}_1 & \vec{a}_2
\end{bmatrix} =
\begin{bmatrix}
3 & 4 \\
6 & 8
\end{bmatrix}$. 
\begin{enumerate}[label=(\alph*)]
\item Is the list of vectors $\vec{a}_1, \vec{a}_2$ linearly dependent or independent?
\item Can you write $\nul A$ as the span of some vector(s)?
\item Can you write column 2 as a scalar multiple of column 1?
\item Is $\spn ( \vec{a}_1, \vec{a}_2 ) = \R^2$?
\end{enumerate}

\begin{notes}
\columnbreak
\end{notes}

\begin{book}
\textbf{Solution.} (a) We want to know if the equation $x_1\vec{a}_1+x_2\vec{a}_2=\vec{0}$ has a non-trivial solution: with $(x_1,x_2)$ not equal to $(0,0)$. Like many examples in the last section, $x_1 \bbvec{3}{6} + x_2 \bbvec{4}{8} = \bbvec{0}{0}$ is equivalent to the system of equations with augmented matrix:
\[\begin{bmatrix}[cc|c]
3 & 4 & 0 \\
6 & 8 & 0
\end{bmatrix}\]
Replacing row 2 with the sum of itself and $-2$ times row $1$, we get the equivalent matrix:
\[\begin{bmatrix}[cc|c]
3 & 4 & 0 \\
0 & 0 & 0
\end{bmatrix}\]
This has a pivot position in column 1, so $x_1$ is a basic variable, and $x_2$ is a free variable since there is no pivot in the second column. By theorem \ref{theorem:existence-and-uniqueness-basics} this has an infinite number of solutions, not just $(x_1,x_2)=(0,0)$. So, the list $\vec{a}_1,\vec{a}_2$ is linearly dependent.

(b) In part (a) we showed that $x_1\vec{a}_1+x_2\vec{a}_2=\vec{0}$ is equivalent to the system with augmented matrix $\begin{bmatrix}[cc|c]
3 & 4 & 0 \\
0 & 0 & 0
\end{bmatrix}$. Let's get this all the way to row reduced echelon form by scaling row 1 by $\frac{1}{3}$:
$\begin{bmatrix}[cc|c]
1 & 4/3 & 0 \\
0 & 0 & 0
\end{bmatrix}$
This is equivalent to the system:
\[\begin{aligned}
x_1+4/3x_2 &=0 \\
0&=0
\end{aligned}\]
The trick now is to write \emph{all} the variables ($x_1$ and $x_2$) in terms of the free variable $x_2$:
\begin{align*}
x_1&=-4/3x_2 \\
x_2&=x_2
\end{align*}
Any real number can be plugged in here for $x_2$ to get a different solution of $x_1\vec{a}_1+x_2\vec{a}_2=\vec{0}$. To see how this represents a span, write the solutions as:
\[\bbvec{x_1}{x_2} = \bbvec{-4/3x_2}{x_2} = x_2 \bbvec{-4/3}{1} \]
The set of all vectors of this form is $\spn \left( \bbvec{-4/3}{1} \right)$.

\bbtwocol{0.68}{0.3}{
(c) Yes: $\bbvec{4}{8}= \frac{4}{3}\bbvec{3}{6}$.

(d) No: $\vec{a}_2$ is \emph{in} $\spn(\vec{a}_1)$, since $\vec{a}_2$ is a scalar multiple of $\vec{a}_1$. So, adding $\vec{a}_2$ into the mix doesn't get you anywhere: $\spn(\vec{a}_1,\vec{a}_2)$ is the line through $\vec{a}_1,\vec{a}_2,$ and $\vec{0}$. \blacksquare}
{
\begin{tikzpicture}[line cap=round,line join=round]%,x=2.0cm,y=1.0cm]
\begin{axis}[
x=15pt,y=15pt,
axis lines=middle,
axis line style={-},
ymajorgrids=true,
xmajorgrids=true,
xmin=-1.5, xmax=5.5,
ymin=-1.5, ymax=9.5,
xtick={-7,-6,...,9},
ytick={-4,-3,...,8},]
\draw[-Stealth,thick] (axis cs:0,0) -- (axis cs:4,8) node[below right]{$\vec{a}_2$};
\draw[-Stealth,thick] (axis cs:0,0) -- (axis cs:3,6) node[below right]{$\vec{a}_1$};
\draw[thick,blue] (axis cs:-5,-10) -- (axis cs:5,10); 
\node[rotate=63.4,above,blue] at (axis cs:2,4){$\spn(\vec{a}_1,\vec{a}_2)$};
\end{axis}
\end{tikzpicture}
}
\end{book}

\example Let $A =\begin{bmatrix}
\vec{a}_1 & \vec{a}_2 & \vec{a}_3
\end{bmatrix} =\begin{bmatrix}
-1 & -3 & 1 \\
2 & 5 & 2 \\
-1 & -2 & -1
\end{bmatrix}$.
\begin{enumerate}[label=(\alph*)]
\item Is the list of vectors $( \vec{a}_1, \vec{a}_2, \vec{a}_3 )$ linearly dependent or independent?
\item Can you write $\nul A$ as the span of some vector(s)?
\item Can you write one of the columns as a linear combination of the preceding columns?
\item Is $\spn ( \vec{a}_1, \vec{a}_2, \vec{a}_3 ) = \R^3$?
\end{enumerate}

\begin{notes}
\columnbreak
\end{notes}

\begin{book}
\textbf{Solution.} (a) We want to know if the equation $x_1\vec{a}_1+x_2\vec{a}_2+x_3\vec{a}_3=\vec{0}$ has a nontrivial solution, with $(x_1,x_2,x_3) \ne (0,0,0)$. This vector equation is equivalent to the system of equations with the following augmented matrix, which we'll row-reduce:
\begin{align*}
&\begin{bmatrix}[ccc|c]
-1 & -3 & 1 & 0 \\
2 & 5 & 2 & 0 \\
-1 & -2 & -1 & 0
\end{bmatrix} \\
&\begin{bmatrix}[ccc|c]
-1 & -3 & 1 & 0 \\
0 & -1 & 4 & 0 \\
0 & 1 & -2 & 0
\end{bmatrix} &\begin{matrix}
\rule{1em}{0pt} \\ r_2 \leftarrow r_2+2r_1 \\ r_3 \leftarrow r_3-1r_1
\end{matrix} \\
&\begin{bmatrix}[ccc|c]
-1 & -3 & 1 & 0 \\
0 & -1 & 4 & 0\\
0 & 0 & 2 & 0
\end{bmatrix} &\begin{matrix}
\rule{1em}{0pt} \\ \rule{1em}{0pt} \\ r_3 \leftarrow r_3+r_2
\end{matrix}
\end{align*}
At this point, we have a pivot in each of the three variable columns, so there are no free variables and the equation $x_1\vec{a}_1+x_2\vec{a}_2+x_3\vec{a}_3=\vec{0}$ has a unique solution, and it must be the trivial solution $(x_1,x_2,x_3)=(0,0,0)$. Thus the list $(\vec{a}_1,\vec{a}_2,\vec{a}_3)$ is linearly independent.

(b) From part (a), the null space contains only one vector, $\vec{0}$. Any scalar multiple of $\vec{0}$ is also $\vec{0}$, so $\nul A = \spn(\vec{0})$.

(c) If it was possible to write $\vec{a}_2$ as a linear combination of $\vec{a}_1$, it would mean that there exists $x_1$ so that $x_1\vec{a}_1=\vec{a}_2$. But this would lead to a non-trivial solution of the homogeneous equation
\[x_1\vec{a}_1-1\vec{a}_2+0\vec{a}_3=\vec{0}\]
and we showed in part (a) that this equation doesn't have a non-trivial solution. Likewise, if it was possible to write $x_1\vec{a}_1+x_2\vec{a}_2=\vec{a}_3$, we'd have a non-trivial solution of
\[x_1\vec{a}_1+x_2\vec{a}_2-1\vec{a}_3=\vec{0}\]
which we know doesn't exist. So it is not possible to write any of the vectors as a linear combination of the previous vectors.

(d) If $\vec{y}=(y_1,y_2,y_3)$ is any vector in $\R^3$, we can write $x_1\vec{a}_1+x_2\vec{a}_2+x_3\vec{a}_3=\vec{y}$ as an augmented matrix and row-reduce it using the same steps in part (a):
\[\begin{bmatrix}[ccc|c]
-1 & -3 & 1 & y_1 \\
2 & 5 & 2 & y_2 \\
-1 & -2 & -1 & y_3
\end{bmatrix} \sim \begin{bmatrix}[ccc|c]
-1 & -3 & 1 & y_1 \\
0 & -1 & 4 & y_2+2y_1\\
0 & 0 & 2 & y_3+y_2+y_1
\end{bmatrix}\]
The pivot positions are in the first three columns no matter what the values of $y_1,y_2,y_3$ are. So there is no pivot in the rightmost column, and by theorem \ref{theorem:existence-and-uniqueness-basics}, the equation $x_1\vec{a}_1+x_2\vec{a}_2+x_3\vec{a}_3=\vec{y}$ has a solution $(x_1,x_2,x_3)$ for any $\vec{y}$ in $\R^3$. In other words, every vector $\vec{y}$ in $\R^3$ is in $\spn(\vec{a}_1,\vec{a}_2,\vec{a}_3)$, so $\spn(\vec{a}_1,\vec{a}_2,\vec{a}_3)= \R^3$. \blacksquare
\end{book}

\example \label{example:single-non-zero}Let $A =\begin{bmatrix}
\vec{a}_1
\end{bmatrix} =\bbvec{-3}{0}$.
\begin{enumerate}[label=(\alph*)]
\item Is the list of one vector $( \vec{a}_1 )$ linearly dependent or independent?
\item Can you write $\nul A$ as the span of some vector(s)?
\item Does $( \vec{a}_1 )$ span $\R^2$?
\end{enumerate}

\begin{notes}
\vfill
\end{notes}

\begin{book}
\textbf{Solution.} (a) Does the equation $x_1 \bbvec{-3}{0} = \bbvec{0}{0}$ have a non-zero solution $x_1$? No; the only solution is $x_1=0$, so the list $( \vec{a}_1 )$ is linearly independent.

(b) The solution set of $x_1 \bbvec{-3}{0} = \bbvec{0}{0}$ is the set of only one vector $\left\{ 0 \right\}$, which we can write as $\spn(0)$, since any multiple of $0$ is itself.

(c) No, $\spn ( \vec{a}_1 )$ is the line through the origin and $(-3,0)$, aka the $x_1$-axis. It doesn't include points like $(17,5)$, so $\spn ( \vec{a}_1 )$ is not $\R^2$. \blacksquare
\end{book}

\example \label{example:set-contains-zero} Let $A =\begin{bmatrix}
\vec{a}_1 & \vec{a}_2
\end{bmatrix} =
\begin{bmatrix}
0 & 3 \\
0 & 2
\end{bmatrix}.$
\begin{enumerate}[label=(\alph*)]
\item Is the list of vectors $\vec{a}_1, \vec{a}_2$ linearly dependent or independent?
\item Can you write $\nul A$ as the span of some vector(s)?
\item Can you write column 2 as a scalar multiple of column 1?
\item Does $( \vec{a}_1, \vec{a}_2 )$ span $\R^2$?
\end{enumerate}

\begin{notes}
\vfill \hspace{2pt}
\columnbreak
\end{notes}

\begin{book}
\textbf{Solution.} (a) Does the equation $x_1 \bbvec{0}{0} + x_2 \bbvec{3}{2} = \bbvec{0}{0}$ have non-trivial solutions, with $(x_1,x_2) \ne (0,0)$? Yes; for example $(57,0)$ is a non-trivial solution since $57 \bbvec{0}{0} + 0 \bbvec{3}{2} = \bbvec{0}{0}$. So the list $(\vec{a}_1,\vec{a}_2)$ is linearly dependent.

(b) The homogeneous equation can be written with the augmented matrix
\[\begin{bmatrix}[cc|c]
0 & 3 & 0 \\
0 & 2 & 0
\end{bmatrix}\]
which is row-equivalent to:
\[\begin{bmatrix}[cc|c]
0 & 1 & 0 \\
0 & 0 & 0
\end{bmatrix}\]
Here, $x_1$ is a free variable because there's no pivot in column 1. Let's do the trick of writing both $x_1$ and $x_2$ in terms of the free variable $x_1$:
\[\begin{aligned}
x_1=x_1 \\
x_2=0
\end{aligned}\quad \text{ or equivalently }\quad \bbvec{x_1}{x_2} = \bbvec{x_1}{0} = x_1 \bbvec{1}{0}\]
So, the solution set of $x_1 \bbvec{0}{0} + x_2 \bbvec{3}{2} = \bbvec{0}{0}$ is $\spn \left( \bbvec{1}{0} \right)$.

(c) No; There is no $x_1$ such that $x_1 \bbvec{0}{0} = \bbvec{3}{2}$. (Note, however, that column 1 can be written as a multiple of column 2.)

\bbtwocol{.6}{.35}{(d) We know $\spn \left( \bbvec{3}{2} \right)$ is the line through the origin and $(3,2)$. Adding multiples of $\bbvec{0}{0}$ doesn't change anything. So $\spn(\vec{a}_1,\vec{a}_2) = \spn(\vec{a}_2)$ which is a line, not all of $\R^2$. \blacksquare
}
{
\begin{tikzpicture}[line cap=round,line join=round]
\begin{axis}[
x=20pt,y=20pt,
axis lines=middle,
axis line style={-},
ymajorgrids=true,
xmajorgrids=true,
xmin=-1.5, xmax=4.5,
ymin=-1.5, ymax=3.5,
xtick={-7,-6,...,9},
ytick={-4,-3,...,8},]
\addplot[only marks, mark=*, mark size=2pt, very thick] coordinates {(0,0)};
\draw[-Stealth,thick] (axis cs:0,0) -- (axis cs:3,2) node[below right]{$\vec{a}_2$};
\draw[thick,blue] (axis cs:-6,-4) -- (axis cs:6,4); 
\node[rotate=33.4,above,blue] at (axis cs:1.5,1){$\spn(\vec{a}_1,\vec{a}_2)$};
\node[below right] at (axis cs:0,0){$\vec{a}_1$};
\end{axis}
\end{tikzpicture}
}
\end{book}

\begin{theorem}
\label{theorem:one-vector-is-independent}
If a list contains only a single vector and it's not $\vec{0}$, then the list is linearly independent.
\end{theorem}

\begin{theorem}
\label{theorem:zero-vector-is-dependent}
If a list of vectors contains $\vec{0}$, then it's linearly dependent.
\end{theorem}

Examples \ref{example:single-non-zero} and \ref{example:set-contains-zero} should help you understand why these theorems are true, but they aren't rigorous proof. The proofs are left as exercises.

\textbf{Logic Fun Corner:}  A statement of the form ``If $P$, then $Q$'' is false if there exists a \emph{counterexample}, which makes $P$ true and $Q$ false. The \emph{converse} of ``If $P$, then $Q$'' is ``If $Q$, then $P$''. If both are true then we write ``$P$ if and only if $Q$'' or abbreviate it ``$P$ iff $Q$.''

For example, the statement ``If $x^2 = 9$, then $x=3$'' is false because there's a counterexample: \rule{5em}{0.8pt}.

\example Is the converse of theorem \ref{theorem:zero-vector-is-dependent} true?

\begin{notes}
\vfill
\end{notes}

\begin{book}
\textbf{Solution.} The converse of theorem 1.5 is the statement: \emph{If a list of vectors is linearly dependent, then it contains $\vec{0}$.} A counterexample to this would be a list of vectors that is linearly independent but does not contain $\vec{0}$. The list of vectors from example \ref{example:1.4first-nul} does the trick:
\[( \vec{a}_1, \vec{a}_2 ) = \left( \bbvec{4}{8}, \bbvec{3}{6} \right)\]
\blacksquare
\end{book}

\textbf{Strategy to write nul \emph{A} as a span:} This is the strategy we used in examples \ref{example:1.4first-nul} through \ref{example:set-contains-zero}. After reducing $
\begin{bmatrix}[c|c]
A & \vec{0}
\end{bmatrix}
$ to reduced echelon form, and solving for \emph{all} the (basic \emph{and} free) variables in terms of the free variables, we can write
\[ \vec{x} = \begin{bmatrix} 
x_1 \\ x_2 \\ \vdots \\ x_c
\end{bmatrix} =
\begin{bmatrix}
0 + \textrm{maybe some combination of free variable terms} \\
0 + \textrm{maybe some combination of free variable terms} \\
\vdots \\
0 + \textrm{maybe some combination of free variable terms} \\
\end{bmatrix}
\]
without constants on the right sides. This means that any solution of a homogeneous system can be written as a linear combination of some vectors, with the free variables as weights. So the solution set is the span of these vectors. You can also write the solution set in \textbf{parametric vector form}, eg $\left\{ x_3 \bbvec{5}{6} : x_3 \text{ is in } \R \right\}$ instead of $\spn \left( \bbvec{5}{6} \right)$.

\columnbreak
\begin{theorem}[Characterization of linear dependence.]
\label{theorem:characterization-of-linear-dependence}

A list of vectors $W = \vec{a}_1, \ldots, \vec{a}_c $ is linearly dependent iff one of the vectors can be written as a linear combination of the others.

Specifically, if $W$ is linearly dependent, and $\vec{a}_1 \ne \vec{0}$, then some $\vec{a}_i$ with $i>1$ is a linear combination of the preceding vectors: \[x_1\vec{a}_1+\cdots+x_{i-1}\vec{a}_{i-1}=\vec{a}_i\]
\end{theorem}
% Prove the above theorem in the book only.
\begin{notes}
Warm-up exercises \ref{exercise:1.4dependence1} and \ref{exercise:1.4dependence2} as well as examples \ref{example:1.4first-nul} through \ref{example:set-contains-zero} give some justification. We'll skip proving this in class, but there's a proof in the textbook.
\end{notes}

\begin{book}
\textbf{Proof:} For the first statement, assume one of the vectors, say $\vec{a}_j$, can be written as a linear combo of the others: \[x_1\vec{a}_1 + \cdots + x_{j-1}\vec{a}_{j-1} + x_{j+1}\vec{a}_{j+1} + \cdots + x_{c}\vec{a}_{c} = \vec{a}_j\]
Subtracting $\vec{a}_j$ from both sides gives:
\[x_1\vec{a}_1 + \cdots + x_{j-1}\vec{a}_{j-1} + (-1)\vec{a}_j+ x_{j+1}\vec{a}_{j+1} + \cdots + x_{c}\vec{a}_{c} = \vec{0}\]
Since $-1 \ne 0$, we have $( x_1, \ldots, x_{j-1}, -1, x_{j+1}, \ldots ,x_{c} ) \ne \vec{0}$,  a non-trivial solution to the homogeneous equation, so $W$ is linearly dependent.

Now for the second statement and the ($\rightarrow$) direction of the first statement, suppose $W$ is dependent. Then there is some non-trivial $(x_1,\ldots, x_c)$ solution to the homogeneous equation, with some of the $x$-value(s) not equal to $0$. Let $x_i$ be the right-most of these (with the greatest $i$ such that $x_i \ne 0$):
\[x_1\vec{a}_1 + \cdots + x_{i-1}\vec{a}_{i-1} + x_i\vec{a}_i + 0\vec{a}_{i+1} + \cdots + 0\vec{a}_c = \vec{0} \]

Since $x_i\ne 0$, we can multiply both sides by $\frac{1}{-x_i}$ as we solve for $\vec{a}_{i}$:
\begin{align}
\label{eq:1.4characterization-thm}
x_1\vec{a}_1 + \cdots + x_{i-1}\vec{a}_{i-1}  + 0\vec{a}_{i+1} + \cdots + 0\vec{a}_c &= -x_i\vec{a}_i \nonumber \\
\frac{x_1}{-x_i}\vec{a}_1 + \cdots + \frac{x_{i-1}}{-x_i}\vec{a}_{i-1}  + 0\vec{a}_{i+1} + \cdots + 0\vec{a}_c &= \vec{a}_i
\end{align}
Thus $\vec{a}_i$ can be written as a linear combination of the other vectors, proving the first statement. Note that if $\vec{a}_1=\vec{0}$, it may be that $i=1$ and the only way to write a vector in $W$ as a linear combo of the others is $0\vec{a}_{2} + \cdots + 0\vec{a}_c = \vec{a}_{1}$. But if we assume that $\vec{a}_1 \ne \vec{0}$, then $i$ can't be $1$, and equation (1.1) shows $\vec{a}_i$ with $i>1$ is a linear combination of the preceding vectors, proving the second statement. \blacksquare

\end{book}

\begin{notes}
\vfill
\end{notes}
\subsection{Theorems About Pivot Positions}\label{subsec:theorems-about-pivot-positions}

Maybe we've been a little sloppy talking about ``the'' basic variables and free variables, and ``the'' pivot positions of a matrix. But is it possible you could row-reduce a matrix using two different sequences of row operations, resulting in different reduced echelon matrices, with different pivot columns? Could a single column be both a pivot column and a non-pivot column, depending on the row operations used? Fortunately, no:

\begin{notes}
\columnbreak
\end{notes}

%TODO: double check for circular logic 
\begin{theorem}[Uniqueness of reduced row echelon form.]
\label{theorem:uniqueness-of-rref}

Every matrix is row-equivalent to a unique reduced row echelon matrix.
\end{theorem}
\begin{book}
Before we can prove Theorem \ref{theorem:uniqueness-of-rref} we need to prove:
\end{book}

\textbf{Lemma. } Suppose $A$ and $B$ are augmented matrices for two systems of homogeneous equations, with each matrix the same size, and each in reduced row echelon form, and $A \ne B$. Then the two systems don't have the same solution sets.
\begin{notes}

The lemma might seem intuitive but the proof is pretty technical and has been relegated to the textbook. Let's use it to prove the theorem.

\vfill
\columnbreak
\end{notes}

\begin{book}
\textbf{Proof of Lemma:} There are two cases to consider:

\begin{itemize}
\item Case 1: The two systems have the same basic variables.
\item Case 2: The two systems don't have the same basic variables.
\end{itemize}

In case 1, the matrices have the same pivot positions, with the leading entry 1's in the same places. Everything left of the leading 1's are zero, and the right-most column is all zeros, and everything above and below the leading 1's in the pivot positions are zeros, so the only place $A$ can be different from $B$ is in a free variable column to the right of a leading 1. Let column $i$ denote the column with the leading 1 in this row, and say the entry that is different is $a$ in $A$ and $b$ in B, both in column $j$. Since $x_j$ is a free variable, there's a solution of each system with $x_j=1$ and all the other free variables equal to zero. In $A$'s system, this gives a solution with $x_i=-a$ and in $B$'s system, it gives $x_i=-b$. So the systems don't have the same solution sets.

In case 2, there is a variable $x_i$ which is a basic variable for one matrix, say $A$, and is a free variable in $B$. In $B$, you can find a solution where $x_i=1$ and all the other free variables are zero. Since every basic variable can be written in terms of only the free variables to the right of it,
% (see TODO: REFERENCE EQUATION/EXAMPLES SHOWING THIS)
this means that any \emph{basic} variables to the right of of $x_i$ in this solution of $B$ will also be zero. So we have a solution for system $B$ where $x_i=1$ and $x_j=0$ for all $j>i$. This cannot be a solution of $A$, since in system $A$, if all the free variables to the right of $x_i$ are zero, then $x_i$ must be zero too. \blacksquare

\textbf{Proof of Thm \ref{theorem:uniqueness-of-rref}.} First of all, the Gaussian elimination strategy in \ref{sec:systems-of-linear-eqns} can be used to reduce any matrix to reduced row echelon form; the best way to see why is to get some experience using it rather than a formal proof. 

To show the reduced echelon form is unique, let $A$ be any matrix, and let $B$ and $C$ be reduced row echelon matrices that are each row-equivalent to $A$. We need to show $B = C$. Augment each matrix with a column of zeros, and denote the augmented matrices as $A^+, B^+$, and $C^+$. These now represent augmented matrices of homogeneous systems, and since row operations don't affect the column of zeros, the augmented matrices are all row-equivalent to each other, and by theorem \ref{theorem:gauss}, the systems all have the same solution set. By the above lemma, $B^+$ cannot be \emph{not} equal to $C^+$, so $B^+=C^+$, and ignoring the columns of zeros, $B=C$. \blacksquare

\emph{Credit:} This proof was inspired by ``A First Course in Linear Algebra'' by Ken Kuttler.
\end{book}
% For Class: Show the lemma (put it in the lecture notes but not the proof.) Then do the proof of  uniqueness-of-rref theorem.

\example  Let $A =\begin{bmatrix}
\vec{a}_1 & \vec{a}_2 & \vec{a}_3 & \vec{a}_4
\end{bmatrix} =
\begin{bmatrix}
1 & -5 & 4 & 2 \\
2 & -9 & 11 & 4 \\
-1 & 5 & -3 & 2
\end{bmatrix}$.

\begin{enumerate}[label=(\alph*)]
\item Is the list of vectors $\vec{a}_1, \vec{a}_2, \vec{a}_3, \vec{a}_4$ linearly dependent or independent?
\item Can you write $\nul A$ as the span of some vector(s)?
\item Can you write one of the columns as a linear combo of the preceding columns?
\item Is $\spn ( \vec{a}_1, \vec{a}_2, \vec{a}_3, \vec{a}_4 ) = \R^3$?
\end{enumerate}

\begin{notes}
\columnbreak
\end{notes}

\begin{book}
\textbf{Solution.} (a) The homogeneous equation $x_1\vec{a}_1 + x_2\vec{a}_2+x_3\vec{a}_3+x_4\vec{a}_4=\vec{0}$ can be represented with an augmented matrix we can row-reduce:
\begin{align*}
&\begin{bmatrix}[cccc|c]
1 & -5 & 4 & 2 & 0 \\
2 & -9 & 11 & 4 & 0 \\
-1 & 5 & -3 & 2 & 0
\end{bmatrix} \\
&\begin{bmatrix}[cccc|c]
1 & -5 & 4 & 2 & 0 \\
0 & 1 & 3 & 0 & 0 \\
0 & 0 & 1 & 4 & 0
\end{bmatrix} 
&\begin{matrix}
\rule{1em}{0pt} \\ r_2 \leftarrow r_2-2r_1 \\ r_3 \leftarrow r_3+r_1
\end{matrix} \\
&\begin{bmatrix}[cccc|c]
1 & -5 & 0 & -14 & 0 \\
0 & 1 & 0 & -12 & 0 \\
0 & 0 & 1 & 4 & 0
\end{bmatrix} 
&\begin{matrix}
r_1 \leftarrow r_1-4r_3 \\r_2\leftarrow r_2-3r_3 \\ \rule{1em}{0pt} 
\end{matrix}\\
&\begin{bmatrix}[cccc|c]
1 & 0 & 0 & -74 & 0 \\
0 & 1 & 0 & -12 & 0 \\
0 & 0 & 1 & 4 & 0
\end{bmatrix} 
&\begin{matrix}
r_1 \leftarrow r_1+5r_2 \\ \rule{1em}{0pt} \\ \rule{1em}{0pt} 
\end{matrix}\\
\end{align*}
The fourth column is a non-pivot column, so $x_4$ is a free variable; the rightmost column is not a pivot column so the homogeneous equation has an infinite number of solutions and the list of vectors is linearly dependent. 

(b) The above augmented matrix shows the homogeneous equation is equivalent to the system:
\[\begin{aligned}
x_1 -74x_4&=0 \\
x_2-12x_4&=0 \\
x_3+4x_4&=0
\end{aligned} \quad \text{ or equivalently } \quad \begin{aligned}
x_1 &=74x_4 \\
x_2 &=12x_4 \\
x_3 &=-4x_4
\end{aligned}\]
And using the strategy outlined above, we should write \emph{all} the variables in terms of the free variable $x_4$:
\[\begin{bmatrix}
x_1 \\ x_2 \\ x_3 \\ x_4
\end{bmatrix} = \begin{bmatrix}
74x_4 \\ 12x_4 \\ -4x_4 \\ x_4
\end{bmatrix} = x_4\begin{bmatrix}
74 \\ 12 \\ -4 \\ 1
\end{bmatrix}\]
Thus the solution set of the homogeneous equation, $\nul A$, is equal to $\spn \left( \begin{bmatrix}
74 \\ 12 \\ -4 \\ 1
\end{bmatrix} \right)$.

(c) Since the list is dependent, theorem \ref{theorem:characterization-of-linear-dependence} says yes. But let's do it. We know one solution of the homogeneous equation is $(74,12,-4,1)$, which means $74\vec{a}_1+12\vec{a}_2-4\vec{a}_3+1\vec{a}_4=\vec{0}$. This gives us a way to write $\vec{a}_4$ as a linear combination of the preceding vectors: $\vec{a}_4=-74\vec{a}_1-12\vec{a}_2+4\vec{a}_3$.

(d) Suppose $\vec{y}=(y_1,y_2,y_3)$. Consider the augmented matrix in part (a) but imagine it has $y_1,y_2,y_3$ in the rightmost column instead of $0,0,0$. The same sequence of row operations would reduce it to reduced row echelon form and the pivot positions would still be in each row in the first three columns; theorem \ref{theorem:uniqueness-of-rref} guarantees there is only one reduced row echelon matrix equivalent to $A$. So there is no way to get a pivot position in the rightmost column of $\begin{bmatrix}[c|c]
A & \vec{y}
\end{bmatrix}$, and the equation $x_1\vec{a}_1+x_2\vec{a}_2+x_3\vec{a}_3+x_4\vec{a}_4=\vec{y}$ is always consistent. Thus every vector $\vec{y}$ in $\R^3$ is in $\spn(\vec{a}_1,\vec{a}_2,\vec{a}_3,\vec{a}_4)$. \blacksquare
\end{book}

\begin{theorem}[The existence theorem.]
\label{theorem:1.4existence}
Suppose $A =
\begin{bmatrix}
\vec{a}_1 & \cdots & \vec{a}_c
\end{bmatrix}
$ is an $r \times c$ matrix. The following are equivalent:
\begin{enumerate}[label=(\roman*)]
\item $A$ has a pivot position in every row.
\item Every system of equations that has $A$ as its coefficient matrix is consistent.
\item $\spn(\vec{a}_1,\ldots,\vec{a}_c) = \R^r$.
\end{enumerate}
% proof (put in textbook version): if A has a pivot in every row, there can't be a pivot in the rightmost column of an augmented matrix, so the above "fact" or thm says it's consistent. If A doesn't have a pivot in every row, then consider the sequence of row operations that reduce it to row echelon form with all zeros in the bottom. Let row k be the row that ends up as the last row with all the swaps. The augmented matrix with constant coefficients y1 = 0, ... yk = 1, ..., ym = 0 when you reduce it, will swap the 1 into the bottom row, maybe with some non-zero scaling, or zeros added to it along the way. This will result in a row echelon matrix with a pivot in the rightmost column.
% NOTE: This proof relies on the uniqueness of rref. Because if rref isn't unique, there may be one rref matrix equivalent to A with pivots in every row and another without pivots in every row.

% In class/notes, skip proof and justify with several examples.
\end{theorem}

\begin{book}
\textbf{Proof.} Parts (ii) and (iii) are saying the same thing with different notation: that for any $\vec{y}$ in $\R^r$, there exists a solution $( x_1,\ldots,x_c )$ of the equation $x_1\vec{a}_1 + \cdots + \cdots x_c\vec{a}_c = \vec{y}$. So we'll prove (i) is equivalent to (ii).

Suppose $A$ has a pivot position in every row, and $\begin{bmatrix}[c|c]
A & \vec{y}
\end{bmatrix}$ is the augmented matrix for a system that has $A$ as its coefficient matrix. There's no way to reduce $\begin{bmatrix}[c|c]
A & \vec{y}
\end{bmatrix}$ to a reduced row echelon matrix that \emph{doesn't} have a pivot in every row left of $\vec{y}$, because the reduced echelon form of $A$ is unique by theorem \ref{theorem:uniqueness-of-rref}. So there's no way to get a row of the form $\begin{bmatrix}[ccc|c]
0 & \cdots & 0 & \text{k}
\end{bmatrix}$ with $k$ non-zero. Thus theorem \ref{theorem:existence-and-uniqueness-basics} part (i) says the system is consistent. 

Now suppose that $A$ does not have a pivot position in every row, and denote $B$ as the row reduced echelon matrix equivalent to $A$. All zero rows of $B$ must be on the bottom, so if we let $\vec{z}=(0, \ldots, 0, 1)$, then the system $\begin{bmatrix}[c|c]
B & \vec{z}
\end{bmatrix}$ is inconsistent. Now row operations that reduced $A$ to $B$ can be reversed, making $\begin{bmatrix}[c|c]
B & \vec{z}
\end{bmatrix}$ row equivalent to $\begin{bmatrix}[c|c]
A & \vec{y}
\end{bmatrix}$ for some vector $\vec{y}$; so $\begin{bmatrix}[c|c]
A & \vec{y}
\end{bmatrix}$ is the augmented matrix for a system with no solution.

\end{book}

\begin{notes}
\vfill
\end{notes}

\begin{theorem}[The uniqueness theorem.]
\label{theorem:1.4uniqueness}
Suppose $A =
\begin{bmatrix}
\vec{a}_1 & \cdots & \vec{a}_c
\end{bmatrix}
$ is an $r \times c$ matrix. The following are equivalent:
\begin{enumerate}[label=(\roman*)]
\item $A$ has a pivot position in every column.
\item $\nul A = \left\{ \vec{0} \right\}$.
\item The columns $\vec{a}_1, \ldots, \vec{a}_c$ are a linearly independent list of vectors.
\end{enumerate}
\end{theorem}

\begin{book}
\textbf{Proof.} Parts (ii) and (iii) are saying the same thing with different notation: that the only solution to the equation $x_1\vec{a}_1+ \cdots + x_c\vec{a}_c = \vec{0}$ is $(x_1,\ldots,x_c) = (0, \ldots, 0) = \vec{0}$. So we can just prove that (i) is equivalent to (iii).

Suppose that $A$ has a pivot position in every column. By the uniqueness of reduced row-echelon form, there are no non-pivot columns of $A$, which is the coefficient matrix of the homogeneous system with augmented matrix $\begin{bmatrix}[c|c]
A & \vec{0}
\end{bmatrix}$. And since homogeneous systems are always consistent, theorem \ref{theorem:existence-and-uniqueness-basics} part (iii) says the solution is unique: $(x_1,\ldots,x_c)=\vec{0}$.

Now suppose that $\vec{a}_1, \ldots, \vec{a}_c$ are a linearly independent list of vectors; that is, the homogeneous equation with augmented matrix $\begin{bmatrix}[c|c]
A & \vec{0}
\end{bmatrix}$ and coefficient matrix $A$ has only one solution. By theorem  \ref{theorem:existence-and-uniqueness-basics} part (ii), $A$ cannot have non-pivot columns, because if it did, there would be an infinite number of solutions.
\end{book}


\begin{theorem}
\label{theorem:more-vectors-than-entries-dependent}
If a list of vectors has more vectors than there are entries in the vectors, then it's linearly dependent.
\end{theorem}

\begin{book}
\textbf{Proof.} Suppose we have $c$ vectors, $\vec{a}_1,\ldots,\vec{a}_c$, in $\R^r$, with $c>r$. Then the matrix $A =
\begin{bmatrix}
\vec{a}_1 & \cdots & \vec{a}_c
\end{bmatrix}$ has more columns then rows. Since you can't have more than one pivot per row, there must be at least one non-pivot column in $A$, which corresponds to a free variable in $x_1\vec{a}_1 + \cdots + x_c\vec{a}_c= \vec{0}$. The free variable implies an infinite number of solutions by theorem \ref{theorem:existence-and-uniqueness-basics}.
\end{book}

\begin{notes}
\vfill \hspace{2pt}
\end{notes}

\exerciseheading

\textbf{\ref{subsec:linear-independence-main-ideas} Linear Independence: Main Ideas}

\exatend \label{exercise:1.4estimate0} Use the graph of $\R^2$ below (and no calculation) to answer:
\begin{enumerate}[label=(\alph*)]
\item Give one or two non-trivial solutions of $x_1\vec{u}+x_2\vec{v}+x_3\vec{w} = \vec{0}$ if possible.
\item Write $\vec{w}$ as a linear combination of $\vec{u},\vec{v}$ if possible.
\end{enumerate}

\begin{tikzpicture}
\begin{axis}[
x=25pt,y=25pt,
axis lines=middle,
axis line style={-},
ymajorgrids=true,
%minor y tick num=3,
%yminorgrids=true,
xmajorgrids=true,
%minor x tick num=3,
%xminorgrids=true,
%xlabel=$x_1$,
%ylabel=$x_2$,
xmin=-3, xmax=3,
ymin=-2, ymax=2,
xtick={-3,-2,...,3},
ytick={-2,-1,...,2},
xticklabels={},
yticklabels={},
]
\draw[-Stealth, very thick] (axis cs:0,0) -- (axis cs:2,-1)
node[above] {$\vec{u}$};
\draw[-Stealth, very thick] (axis cs:0,0) -- (axis cs:1,1)
node[right] {$\vec{v}$};
\draw[-Stealth, very thick] (axis cs:0,0) -- (axis cs:-1,0)
node[below] {$\vec{w}$};
\end{axis}
\end{tikzpicture}

\exatend Using the same $\vec{u},\vec{v}$ in the graph of $\R^2$ in the previous exercise:
\begin{enumerate}[label=(\alph*)]
\item Give one or two non-trivial solutions of $x_1\vec{u}+x_2\vec{v} = \vec{0}$ if possible.
\item Write $\vec{v}$ as a linear combination of $\vec{u}$ if possible.
\end{enumerate}

\exatend Each of the following are augmented matrices for linear systems. Which system(s) are consistent?

\begin{horiparts}{2}
\item $\begin{bmatrix}[ccc|c]
1 & 0 & 0 & 8 \\
0 & 1 & 0 & 9 \\
0 & 0 & 1 & 10
\end{bmatrix}$
\item $\begin{bmatrix}[ccc|c]
1 & 0 & 3 & -20 \\
0 & 1 & 17 & 0
\end{bmatrix}$
\item $\begin{bmatrix}[ccc|c]
1 & 0 & 0 & 8 \\
0 & 1 & 0 & 9 \\
0 & 0 & 0 & 10
\end{bmatrix}$
\item $\begin{bmatrix}[ccc|c]
1 & 0 & 0 & 8 \\
0 & 1 & 0 & 9 \\
0 & 0 & 0 & 0
\end{bmatrix}$
\end{horiparts}

\exatend Each of the following are augmented matrices for linear systems. Out of the system(s) that are consistent, which have an infinite number of solutions?

\begin{horiparts}{2}
\item $\begin{bmatrix}[ccc|c]
1 & 0 & 0 & 8 \\
0 & 1 & 0 & 9 \\
0 & 0 & 1 & 10
\end{bmatrix}$
\item $\begin{bmatrix}[ccc|c]
1 & 0 & 3 & -20 \\
0 & 1 & 17 & 0
\end{bmatrix}  $
\item $\begin{bmatrix}[ccc|c]
1 & 0 & 0 & 8 \\
0 & 1 & 0 & 9 \\
0 & 0 & 0 & 10
\end{bmatrix}$
\item $\begin{bmatrix}[ccc|c]
1 & 0 & 0 & 8 \\
0 & 1 & 0 & 9 \\
0 & 0 & 0 & 0
\end{bmatrix}$
\end{horiparts}

\exatend Each of the following are \emph{coefficient} matrices for linear systems. In other words, each is a matrix $A$ for a matrix equation $A \vec{x} = \vec{b}$. Which system(s) are always consistent, no matter vector $\vec{b}$ is on the other side of the equation?

\begin{horiparts}{3}
\item $\begin{bmatrix}
1 & 0 & 0 \\
0 & 1 & 0  \\
0 & 0 & 1
\end{bmatrix}$
\item $\begin{bmatrix}
1 & 0 & 3  \\
0 & 1 & 17
\end{bmatrix}$
\item $\begin{bmatrix}
1 & 0 & 0 \\
0 & 1 & 0 \\
0 & 0 & 0
\end{bmatrix}$
\end{horiparts}

\exatend Write your own example of a matrix $A$ that makes $A\vec{x}=\vec{0}$ have an infinite number of solutions. Justify your answer.

\exatend Which of the following is the solution set of $\begin{aligned}
x_1 + 2x_2 &= 0 \\
3x_1+6x_2 &= 0
\end{aligned}$?

(A) $\spn \left( \bbvec{1}{2} \right)$ \quad \quad
(B) $\spn \left( \bbvec{-2}{1} \right)$ \quad \quad
(C) $\R^2$ \quad \quad
(D) $\left\{ \vec{0} \right\}$

\exatend Describe the solution set of $x_1+3x_2=0$ as the span of some vectors. \hintfootnote{
$\begin{bmatrix}[cc|c]
1 & 3 & 0
\end{bmatrix}$ is already in reduced row echelon form.}

\vspace{1ex}
Remember to give some justification for all True/False answers:

\exatend (T/F?) The system of equations $\begin{aligned}
  x_1 + x_{2} &=3 \\
  -x_{1} + 2x_{2} &= 0
\end{aligned}$ is homogeneous because of the number 0 on the right side.


\exatend (T/F?) Every homogeneous linear system of equations is consistent.

\exatend (T/F?) The solution set of every homogeneous system can be written as the span of some vectors.

\exatend (T/F?) If $\vec{a}_1,\vec{a}_2,\vec{a}_3$ are in $\R^3$, and $\vec{a}_1, \vec{a}_2$ is a linearly dependent list, then the list $\vec{a}_1, \vec{a}_2, \vec{a}_3$ is linearly dependent.

\exatend (T/F?) If $\vec{a}_1,\vec{a}_2,\vec{a}_3$ are in $\R^3$, and $\vec{a}_1, \vec{a}_2$ is a linearly \emph{independent} list, then the list $\vec{a}_1, \vec{a}_2, \vec{a}_3$ is linearly \emph{independent}.

\exatend (T/F?) The list of vectors $\vec{u}, \vec{v}, 3\vec{u}-2\vec{v}$ is linearly dependent.

\exatend Let $A =\begin{bmatrix}
\vec{a}_1 & \vec{a}_2
\end{bmatrix} =\begin{bmatrix}
3 & 1 \\
2 & 4 \\
1 & 0 
\end{bmatrix}$.
\begin{enumerate}[label=(\alph*)]
\item Is the list of vectors $\vec{a}_1, \vec{a}_2$ linearly dependent or independent?
\item Write $\nul A$ as the span of some vector(s).
\item Can you write one of the columns as a linear combination of the preceding columns?
\item Does the list $( \vec{a}_1, \vec{a}_2 )$ span $\R^3$?
\end{enumerate}

\exatend Let $A =\begin{bmatrix}
\vec{a}_1 & \vec{a}_2 & \vec{a}_{3}
\end{bmatrix} =\begin{bmatrix}
3 & 1 & -5\\
2 & 4 & 1\\
1 & 0 & 0
\end{bmatrix}$.
\begin{enumerate}[label=(\alph*)]
\item Is the list of vectors $\vec{a}_1, \vec{a}_2, \vec{a}_3$ linearly dependent or independent?
\item Write $\nul A$ as the span of some vector(s).
\item Can you write one of the columns as a linear combination of the preceding columns?
\item Does the list $( \vec{a}_1, \vec{a}_2, \vec{a}_3 )$ span $\R^3$?
\end{enumerate}

\exatend Let $A =\begin{bmatrix}
\vec{a}_1 & \vec{a}_2 & \vec{a}_{3}
\end{bmatrix} =\begin{bmatrix}
2 & 4 & 1\\
1 & 0 & 0
\end{bmatrix}$.
\begin{enumerate}[label=(\alph*)]
\item Is the list of vectors $\vec{a}_1, \vec{a}_2, \vec{a}_3$ linearly dependent or independent?
\item Write $\nul A$ as the span of some vector(s).
\item Can you write one of the columns as a linear combination of the preceding columns?
\item Is $\spn ( \vec{a}_1, \vec{a}_2, \vec{a}_3 ) = \R^2$?
\end{enumerate}

\exatend Let $A =\begin{bmatrix}
\vec{a}_1 & \vec{a}_2 & \vec{a}_{3} & \vec{a}_4
\end{bmatrix} =\begin{bmatrix}
4 & 8 & -18 & 6 \\
-1 & -2 & 5 & -3 \\
-1 & -2 & 4 & 0
\end{bmatrix}$.
\begin{enumerate}[label=(\alph*)]
\item Is the list of vectors $ \vec{a}_1, \vec{a}_2, \vec{a}_3, \vec{a}_4 $ linearly dependent or independent?
\item Write $\nul A$ as the span of some vector(s).
\item Can you write one of the columns as a linear combination of the preceding columns?
\item Is $\spn ( \vec{a}_1, \vec{a}_2, \vec{a}_3, \vec{a}_4 ) = \R^3$?
\end{enumerate}

\exatend Let $A =
\begin{bmatrix}
10 & 20 & 30 \\
-9 & -18 & -27 \\
8 & 16 & 24
\end{bmatrix}$
\begin{enumerate}[label=(\alph*)]
\item Row reducing $A$ and using theorem \ref{theorem:1.4existence}, is the span of the columns of $A$ equal to $\R^3$?
\item Note that the second and third columns are scalar multiples of the first. Is the span of the columns a single point, a line, a plane, or all of $\R^{3}$?
\item Note that the third column is the sum of the first two columns. Use this fact to write a solution of $x_1\vec{a}_1+x_2\vec{a}_2+x_3\vec{a}_3 = \vec{0}$.
\end{enumerate}

\exatend Find a counterexample that proves the converse of theorem \ref{theorem:one-vector-is-independent} is false.

% The problems at the end of this subsection SHOULD be at the end of the section, but I moved them here in case teachers divide the section into two.
\exatend Prove or give a counterexample: If $\vec{a}_1, \ldots, \vec{a}_c$ is a linearly independent list of at least two vectors, and you remove one vector, then the remaining list is still independent.

\exatend Prove or give a counterexample: If $\vec{a}_1, \ldots, \vec{a}_c$ is a linearly dependent list of at least two vectors, and you remove one vector, then the remaining list is still dependent.

%inspired by LADR 2.A # 8
\exatend Prove or give a counterexample: If $\vec{a}_1, \ldots, \vec{a}_c$ is a linearly independent list of vectors in $\R^n$ and if $k$ is a real number with $k\ne 0$, then the list $k\vec{a}_1, \ldots, k\vec{a}_c$ is also linearly independent.

% inspired by LADR 2.A # 9
\exatend Prove or give a counterexample: If $\vec{a}_1, \ldots, \vec{a}_c$ and $\vec{b}_1,\ldots, \vec{b}_c$ are each linearly independent lists in $\R^n$, then the list $\vec{a}_1 + \vec{b}_1, \ldots, \vec{a}_c + \vec{b}_c$ is also linearly independent.

% inspired by LADR 2.A # 6 -- pretty hard
\exatend Suppose $\vec{u}_1, \vec{u}_2, \vec{u}_3$ is linearly independent in $\R^{c}$. Prove that the list $\vec{u}_1-\vec{u}_2, \vec{u}_2-\vec{u}_3, \vec{u}_3 $ is also linearly independent.

%inspired by LADR 2.A # 10 -- Challenge!
\exatend (\emph{Challenge}) Suppose $\vec{u}_1, \ldots, \vec{u}_c$ is a linearly independent list of vectors in $\R^n$, and $\vec{w}$ is in $\R^n$. Prove that if $\vec{u}_1 + \vec{w}, \ldots, \vec{u}_c+\vec{w}$ is linearly dependent, then $\vec{w}$ is in $\spn ( \vec{u}_1, \ldots, \vec{u}_c )$.

\textbf{\ref{subsec:theorems-about-pivot-positions} Theorems About Pivot Positions}

\exatend (T/F?) If there is a pivot position in every column of an $r \times c$ matrix $A$, then the columns of $A$ span $\R^r$.

\exatend Let $\vec{v}_1 =
\begin{bmatrix}
0 \\ -1 \\ 1 \\ -1
\end{bmatrix}$, $\vec{v}_2 =
\begin{bmatrix}
0 \\ 1 \\ 0 \\ 0
\end{bmatrix}$, and $\vec{v}_3 =
\begin{bmatrix}
-1 \\ 0 \\ 1 \\ 0
\end{bmatrix}$.
\begin{enumerate}[label=(\alph*)]
\item Before doing any calculations, what is the maximum possible number of pivots that $\begin{bmatrix}
\vec{v}_1 & \vec{v}_2 & \vec{v}_{3}
\end{bmatrix}$ could have? What does theorem \ref{theorem:1.4existence} say about whether $\spn ( \vec{v}_1,\vec{v}_2,\vec{v}_3 ) = \R^4$ or not?
\item What must be true about $\vec{y} = (y_1,y_2,y_3,y_4)$ in order for $\vec{y}$ to be in $\spn ( \vec{v}_1,\vec{v}_2,\vec{v}_3 )$? (Row-reduce $\begin{bmatrix}[ccc|c]
\vec{v}_1 & \vec{v}_2 & \vec{v}_{3} & \vec{y}
\end{bmatrix}$.)
\end{enumerate}

\exatend Is the list $\bbbvec{1}{2}{3}, \bbbvec{4}{5}{6}, \bbbvec{7}{8}{9}$ linearly dependent or independent?
\exatend Let $A =
\begin{bmatrix}
\vec{a}_1 & \vec{a}_2 & \vec{a}_3
\end{bmatrix} =
\begin{bmatrix}
2 & 4 & -6 \\
1 & 4 & -1 \\
1 & 3 & -2
\end{bmatrix}$ and $\vec{b}= \bbbvec{b_1}{b_2}{b_3}$. Does $x_1\vec{a}_1 + x_2\vec{a}_2 + x_3\vec{a}_3=\vec{b}$ have a solution for all possible $\vec{b}$? If not, describe the set of all $\vec{b}$ for which it does have a solution.

\exatend Let $\begin{bmatrix}
2 & 4 & -5 \\
1 & 4 & -1 \\
1 & 3 & -2
\end{bmatrix}$ and $\vec{b}= \bbbvec{b_1}{b_2}{b_3}$. Does $x_1\vec{a}_1 + x_2\vec{a}_2 + x_3\vec{a}_3=\vec{b}$ have a solution for all possible $\vec{b}$? If not, describe the set of all $\vec{b}$ for which it does have a solution.

\exatend (T/F?) A system of equations with more equations than variables cannot have an infinite number of solutions.

\exatend (T/F?) If $\vec{u}$ and $\vec{v}$ are in $\R^c$, and $\vec{v}$ can't be written as a multiple of $\vec{u}$, then the list $\vec{u}, \vec{v}$ is linearly independent. \hintfootnote{Be careful.}

\exatend (T/F?) If a $10 \times 12$ matrix has 10 pivot columns, then the span of the columns equals $\R^{10}$.

\exatend Write a list of two $\R^3$ vectors $\vec{a}_1, \vec{a}_2$ so that $x_1\vec{a}_1 + x_2\vec{a}_2 = \vec{0}$ has a unique solution. Then write a list of two $\R^3$ vectors $\vec{b}_1, \vec{b}_2$ so that $x_1\vec{b}_1 + x_2\vec{b}_2 = \vec{0}$ has an infinite number of solutions. 

\exatend Prove theorem \ref{theorem:one-vector-is-independent}. Remember that your proof should cover \emph{any} list containing a single non-zero vector, not just one example. So start with something like ``Suppose $W=(\vec{a})$ is a list with just one vector, and $\vec{a}\ne \vec{0}$.''

\exatend Prove theorem \ref{theorem:zero-vector-is-dependent}.

\exatend Explain why a system of linear equations with more variables than equations cannot have a unique solution.

\exatend Explain, using theorem \ref{theorem:1.4existence}, why a set of fewer than $r$ vectors cannot span $\R^r$.

\exatend Suppose $A$ is an $r \times c$ matrix, and for all $\vec{y}$ in $\R^r$, the equation $A\vec{x}=\vec{y}$ has at most one solution. Show that the columns of $A$ are linearly independent, using the definition. 

\exatend Suppose $A$ is an $r \times c$ matrix with $c$ pivot columns. Show that the columns of $A$ are linearly independent.

\exatend Find a number $k$ so that $\left\{ \bbbvec{2}{-4}{3}, \bbbvec{-1}{3}{-6}, \bbbvec{-3}{9}{k} \right\}$ is linearly independent.

\columnbreak
\section{Subspaces of $\R^c$}
\label{sec:subspaces-of-Rc}

\exercise (\emph{Warm-up}) The graph shows the unit disk $D = \left\{ (x_1,x_2) : x_1^2+x_2^2 \le 1 \right\}$. 

\begin{enumerate}[label=(\alph*)]
\item (T/F?) $\vec{0}$ is in $D$.
\item (T/F?) For any two points $\vec{x}$ and $\vec{y}$ in $D$, the vector $\vec{x}+ \vec{y}$ is in $D$. (If you're stuck, a good approach to this kind of problem is to try to find a counterexample: $\vec{x},\vec{y}$ so that $\vec{x}+\vec{y}$ is \emph{not} in $D$.)
\item (T/F?) For any point $\vec{x}$ in $D$ and any scalar $k$, $k\vec{x}$ is in $D$.
\end{enumerate}
\begin{tikzpicture}[scale=2]
% Fill the unit disk with light blue
\fill[blue!20] (0,0) circle (1);
\draw[blue] (0,0) circle (1);
% Add the label "D" in the center
\node[blue] at (0.5,0.5) {$D$};
% Optional: Add coordinate axes for reference
\draw (-1.2,0) -- (1.2,0);
\draw (0,-1.2) -- (0,1.2);
\draw[black,thick] (-1,0.06) -- (-1,-0.06) node[below] {$-1$};
\draw[black,thick] (1,0.06) -- (1,-0.06) node[below] {$1$};
\draw[black,thick] (0.06,1) -- (-0.06,1) node[left] {$1$};
\draw[black,thick] (0.06,-1) -- (-0.06,-1) node[left] {$-1$};
\end{tikzpicture}

\begin{notes}
\vfill
\end{notes}

\exercise (\emph{Warm-up}) The graph below shows $W = \spn ( (2,-1) )$
\begin{enumerate}[label=(\alph*)]
\item (T/F?) $\vec{0}$ is in $W$.
\item (T/F?) For any two points $\vec{x}$ and $\vec{y}$ in $W$, the vector $\vec{x}+ \vec{y}$ is in $W$. 
\item (T/F?) For any point $\vec{x}$ in $W$ and any scalar $k$, $k\vec{x}$ is in $W$.
\end{enumerate}
\begin{tikzpicture}[scale=0.8]
% Draw the axes
\draw (-3,0) -- (3,0);  % x-axis
\draw (0,-2) -- (0,2);  % y-axis
\foreach \x in {-3,-2,-1,0,1,2,3} {
\draw[black] (\x,0.1) -- (\x,-0.1);}
\foreach \y in {-2,-1,0,1,2} {
\draw[black] (0.1,\y) -- (-0.1,\y);}
\draw[black] (2,0.1) -- (2,-0.1) node[below] {$2$};
\draw[black] (0.1,-1) -- (-0.1,-1) node[left] {$-1$};
\draw[blue, thick] (-3,1.5) -- (3,-1.5);
\node[blue, above] at (-2,1) {$W$};
\end{tikzpicture}

\begin{notes}
\vfill \hspace{2pt}
\columnbreak
\end{notes}

\begin{definition}
A \textbf{subspace} $U$ of $\R^n$ is a subset of $\R^n$ for which:
\begin{enumerate}[label=(\roman*)]
\item The zero vector of $\R^n$ is in $U$.
\item If $\vec{p}$ and $\vec{q}$ are in $U$, then $\vec{p}+\vec{q}$ is in $U$. ($U$ is \textbf{closed under vector addition}.)
\item If $\vec{p}$ is in $U$ and $c$ is in $\R$, then $c\vec{p}$ is in $U$. ($U$ is \textbf{closed under scalar multiplication}.)
\end{enumerate}
\end{definition}

\example Show that the set $U = \left\{ \vec{0} \right\}$ is a subspace of $\R^n$. This is called the \emph{trivial subspace}.

\begin{notes}
\vfill
\end{notes}

\begin{book}
\textbf{Solution.} We need to show that it meets all three conditions of the definition:
\begin{enumerate}[label=(\roman*)]
\item $\vec{0}$ is certainly in $\left\{ \vec{0} \right\}$.
\item If $\vec{p}$ and $\vec{q}$ are in $\left\{ \vec{0} \right\}$, then $\vec{p}=\vec{q}=\vec{0}$ and $\vec{0}+\vec{0}=\vec{0}$ is in $\left\{ \vec{0} \right\}$.
\item If $\vec{p}$ is in $\left\{ \vec{0} \right\}$, then $\vec{p}=\vec{0}$, and for any $c$ in $\R$, we have $c\vec{p}=c\vec{0}=\vec{0}$ which is in $\left\{ \vec{0} \right\}$. \blacksquare
\end{enumerate}
\end{book}

\textbf{Note:} We say set $V$ is a \emph{subset} of $W$ if everything in $V$ is also in $W$. This means that $W$ is a subset of itself. We say $V$ is a \emph{proper subset} of $W$ if it's a subset of $W$ and there's something in $W$ but not $V$. So $W$ is not a proper subset of itself.
\example Show that $U = \R^n$ is a subspace of $\R^n$.

\begin{notes}
\vfill \hspace{2pt}
\columnbreak
\end{notes}

\begin{book}
\textbf{Solution.} All three of the conditions in the definition are clearly satisfied:
\begin{enumerate}[label=(\roman*)]
\item $\vec{0}$ is in $\R^n$.
\item If $\vec{p}$ and $\vec{q}$ are in $\R^n$, then $\vec{p}+\vec{q}$ is in $\R^n$.
\item if $\vec{p}$ is in $\R^n$ and $c$ is a real number, then $c\vec{p}$ is in $\R^n$. \blacksquare
\end{enumerate}
\end{book}

\example Is $M = \left\{ (a,b,c) : abc = 0 \right\}$ a subspace of $\R^3$?

\begin{notes}
\vfill
\end{notes}

\begin{book}
\textbf{Solution.} Let's think about each of the conditions:
\begin{enumerate}[label=(\roman*)]
\item Is $\vec{0}$ in $M$? Well, $0\cdot 0 \cdot 0 = 0$, so yes.
\item Suppose $\vec{p}=(p_1,p_2,p_3)$ and $\vec{q}=(q_1,q_2,q_3)$ are in $M$. In other words, suppose $p_1p_2p_3=0$ and $q_1q_2q_3=0$. Does this guarantee that $(p_1+q_1)(p_2+q_2)(p_3+q_3) = 0$? No it does not: one counterexample is $\vec{p}=(0,2,3)$ and $\vec{q}=(1,0,3)$, which makes $\vec{p}+\vec{q}=(1,2,6)$ which is not in $M$. This is enough to show that $M$ is not a subspace. \blacksquare
\end{enumerate}
\end{book}

\example Is $P = \left\{ (a,b,c) : a+b+c = 0 \right\}$ a subspace of $\R^3$?

\begin{notes}
\vfill \hspace{2pt}
\columnbreak
\end{notes}

\begin{book}
\textbf{Solution.} Let's think about each of the conditions:
\begin{enumerate}[label=(\roman*)]
\item Is $\vec{0}$ in $P$? Well, $0+0+0=0$, so yes.
\item Suppose $\vec{p}=(p_1,p_2,p_3)$ and $\vec{q}=(q_1,q_2,q_3)$ are in $P$. In other words, suppose $p_1+p_2+p_3=0$ and $q_1+q_2+q_3=0$. Does this guarantee that $(p_1+q_1) + (p_2+q_2) + (p_3+q_3) = 0$? Yes it does, because you can rewrite it as $(p_1+p_2+p_3) + (q_1+q_2+q_3)= 0 + 0 = 0$.
\item Suppose $\vec{p}=(p_1,p_2,p_3)$ is in $P$ and $c$ is in $\R$. Does this guarantee that $c\vec{p}$ is in $P$? Yes it does, since $cp_1+cp_2+cp_3=c(p_1+p_2+p_3) = c(0) = 0$. Thus $P$ satisfies all three conditions and is a vector space. \blacksquare
\end{enumerate}
\end{book}

\begin{theorem}[Null spaces are subspaces.]
\label{theorem:null-spaces-are-subspaces}

For any $r \times c$ vector $A =\begin{bmatrix}
\vec{a}_1 & \cdots & \vec{a}_c\end{bmatrix}$
, the null space of $A$ is a subspace of $\R^c$.
\end{theorem}

\begin{notes}
\vfill
\end{notes}

\begin{book}
\textbf{Proof.} (i) Since $0\vec{a}_1 + \cdots + 0\vec{a}_c = \vec{0}$, we have $\vec{0}$ is a solution of the homogeneous equation, so $\vec{0}$ is in $\nul A$.

(ii) Suppose $\vec{p}= ( p_1, \ldots, p_c )$ and $\vec{q} = ( q_1, \ldots, q_c )$ are in $\nul A$. In other words, suppose $p_1\vec{a}_1 + \cdots + p_c\vec{a}_c = \vec{0}$ and $q_1\vec{a}_1 + \cdots + q_c\vec{a}_c = \vec{0}$. Plugging in $\vec{p}+\vec{q} = (p_1+q_1,\ldots, p_c+q_c)$ into the homogeneous equation gives
\begin{align*}
( p_1+q_1 )\vec{a}_1 + \cdots + (p_c+q_c)\vec{a}_c &= \\
p_1\vec{a}_1+q_1\vec{a}_1 + \cdots + p_c\vec{a}_c+q_c\vec{a}_c &= \\
p_1\vec{a}_1+ \cdots + p_c\vec{a}_c + q_1\vec{a}_1 + \cdots + q_c\vec{a}_c &= \vec{0} + \vec{0} = \vec{0}
\end{align*}
so $\vec{p}+\vec{q}$ is also in $\nul A$.

(iii) Suppose $\vec{p} = (p_1, \ldots, p_c)$ is in $\nul A$, and suppose $k$ is a real number. Plugging in $k\vec{p} = (kp_1,\ldots,kp_c)$ into the homogeneous equation gives
\[
kp_1\vec{a}_1+\cdots+kp_c\vec{a}_c = k(p_1\vec{a}_1+ \cdots + p_c\vec{a}_c) = k\vec{0} = \vec{0}
\]
so $k\vec{p}$ is also in $\nul A$. \blacksquare
\end{book}

\begin{theorem}[Spans are subspaces.]
\label{theorem:spans-are-subspaces}

If $\vec{a}_1, \ldots, \vec{a}_c$ is a list of vectors in $\R^n$, then $\spn (\vec{a}_1,\ldots, \vec{a}_c)$ is a subspace of $\R^n$.
\end{theorem}

\begin{notes}
\vfill \hspace{2pt}
\columnbreak
\end{notes}

\begin{book}
\textbf{Proof.} (i) $\vec{0}=0\vec{a}_1+\cdots+0\vec{a}_c$ is in $\spn(\vec{a}_1,\ldots,\vec{a}_2)$.

(ii) Suppose $\vec{p}$ and $\vec{q}$ are in $\spn(\vec{a}_1,\ldots,\vec{a}_2)$. Then $\vec{p}= x_1\vec{a}_1+\cdots+x_c\vec{a}_c$ and $\vec{q}=y_1\vec{a}_1+\cdots+y_c\vec{a}_c$ for some $(x_1,\ldots,x_c)$ and $(y_1,\ldots,y_c)$, so
\begin{align*} \vec{p}+\vec{q} &= x_1\vec{a}_1+\cdots+x_c\vec{a}_c + y_1\vec{a}_1+\cdots+y_c\vec{a}_c \\
&= (x_1+y_1)\vec{a}_1+\cdots+(x_c+y_c)\vec{a}_c
\end{align*}
is in  $\spn(\vec{a}_1,\ldots,\vec{a}_2)$ as well.

(iii) Suppose $\vec{p}$ is in  $\spn(\vec{a}_1,\ldots,\vec{a}_2)$ and $k$ is a real number. Then $\vec{p}= x_1\vec{a}_1+\cdots+x_c\vec{a}_c$ for some $(x_1,\ldots,x_c)$, so 
\[ k\vec{p}= k( x_1\vec{a}_1+\cdots+x_c\vec{a}_c) = kx_1\vec{a}_1 + \cdots + kx_c\vec{a}_p\]
is in  $\spn(\vec{a}_1,\ldots,\vec{a}_2)$ as well.
\end{book}

\begin{definition}
The \textbf{column space} of a matrix $A$, written $\col A$, is the span of the columns of $A$.

The \textbf{row space} of $A$, written $\row A$, is the span of the rows of $A$, written as vectors.
\end{definition}

\example Let $A =\begin{bmatrix}
1 & 2 & 3 \\
4 & 5 & 6
\end{bmatrix}$. Find a few vectors in each of the following:

\begin{enumerate}[label=(\alph*)]
\item $\col A$
\begin{notes}
\vfill
\end{notes}
\item $\row A$
\begin{notes}
\vfill
\end{notes}
\item $\nul A$
\end{enumerate}

\begin{notes}
\vfill \hspace{2pt}
\columnbreak
\end{notes}

\begin{book}
\textbf{Solution.} (a) Any linear combination of the columns of $A$ is in $\col A$, such as:
\[0 \bbvec{1}{4} + 1 \bbvec{2}{5} + 0 \bbvec{3}{6} = \bbvec{2}{5} \quad \text{or} \quad 
\pi \bbvec{1}{4} + 0 \bbvec{2}{5} - 1\bbvec{3}{6} = \bbvec{\pi-3}{4\pi-6} \]

(b) Any linear combination of the rows of $A$ is in $\row A$, such as:
\[0 \bbbvec{1}{2}{3} + 2 \bbbvec{4}{5}{6} = \bbbvec{8}{10}{12} \quad \text{or} \quad
1.5 \bbbvec{1}{2}{3} -7 \bbbvec{4}{5}{6} = \bbbvec{-26.5}{-32}{-37.5}\]

(c) Vectors in $\nul A$ are solutions of the equation $x_1 \bbvec{1}{4} + x_2 \bbvec{2}{5} + x_3 \bbvec{3}{6}= \vec{0}$ which we can solve by row-reducing the corresponding augmented matrix:
\begin{align*}
\begin{bmatrix}[ccc|c]
1 & 2 & 3 & 0 \\
4 & 5 & 6 & 0
\end{bmatrix} \sim \begin{bmatrix}[ccc|c]
1 & 2 & 3 & 0 \\
0 & -3 & -6 & 0
\end{bmatrix} \sim \begin{bmatrix}[ccc|c]
1 & 0 & -1 & 0 \\
0 & 1 & 2 & 0
\end{bmatrix} \sim
\begin{aligned}
x_1-x_3&=0 \\
x_2+2x_3&=0
\end{aligned}
\end{align*}
Now that this is reduced row echelon form, we can write all the variables in terms of the free variable $x_3$:
\[\bbbvec{x_1}{x_2}{x_3} = \bbbvec{x_3}{-2x_3}{x_3} = x_3\bbbvec{1}{-2}{1}\]
Thus some vectors in the null space are $(1,-2,1), (0,0,0),$ and $(-3,6,-3)$. This problem shows that it's easy to find vectors in $\col A$ and $\row A$, but finding vectors in $\nul A$ involves solving an equation by row-reducing. \blacksquare
\end{book}

\example Let $A =\begin{bmatrix}
1 & 2 & 3 \\
4 & 5 & 6
\end{bmatrix}$. 

\begin{enumerate}[label=(\alph*)]
\item Is $\bbvec{12}{21}$ in $\col A$?
\begin{notes}
\vfill
\end{notes}
\item Is $\bbbvec{5}{7}{10}$ in $\row A$?
\begin{notes}
\vfill
\end{notes}
\item Is $\bbbvec{-\sqrt{7}}{2\sqrt{7}}{-\sqrt{7}}$ in $\nul A$?
\end{enumerate}

\begin{notes}
\vfill \hspace{2pt}
\columnbreak
\end{notes}

\begin{book}
\textbf{Solution.} (a) Too see if $(12,21)$ is in $\col A$, we need to see if there exists $(x_1,x_2,x_3)$ such that:
\[x_1 \bbvec{1}{4} + x_2 \bbvec{2}{5} + x_3 \bbvec{3}{6}= \bbvec{12}{21}\]
Let's row-reduce the corresponding augmented matrix:
\[\begin{bmatrix}[ccc|c]
1 & 2 & 3 & 12 \\
4 & 5 & 6 & 21
\end{bmatrix} \sim \begin{bmatrix}[ccc|c]
1 & 2 & 3 & 12 \\
0 & -3 & -6 & -27
\end{bmatrix} \sim \begin{bmatrix}[ccc|c]
1 & 0 & -1 & -6 \\
0 & 1 & 2 & 9
\end{bmatrix}\]
This system is consistent because there's no pivot in the rightmost column. Also, since there's a pivot in every row of $A$, we know \emph{every} $\R^2$ vector is in $\col A$ because there's no way to get a pivot in the rightmost column.

(b) To see if $(5,7,10)$ is in $\row A$, we need to see if there exists $(x_1,x_2)$ such that:
\[x_1 \bbbvec{1}{2}{3} + x_2 \bbbvec{4}{5}{6} = \bbbvec{5}{7}{10}\]
Let's row-reduce the corresponding augmented matrix:
\[\begin{bmatrix}[cc|c]
1 & 4 & 5 \\
2 & 5 & 7 \\
3 & 6 & 10
\end{bmatrix} \sim \begin{bmatrix}[cc|c]
1 & 4 & 5 \\
0 & -3 & -3 \\
0 & -6 & -5
\end{bmatrix} \sim \begin{bmatrix}[cc|c]
1 & 4 & 5 \\
0 & -3 & -3 \\
0 & 0 & 1
\end{bmatrix}\]
This is inconsistent because the last row says $0=1$, so $(5,7,10)$ is not in $\row A$.

(c) To check if $(-\sqrt{7}, 2\sqrt{7}, -\sqrt{7})$ is in $\nul A$, we just need to see if it's a solution of the homogeneous equation:
\begin{align*}
-\sqrt{7} \bbvec{1}{4} +2\sqrt{7} \bbvec{2}{5} -\sqrt{7} \bbvec{3}{6} &= \bbvec{0}{0} \\
 \bbvec{-\sqrt{7}+4\sqrt{7}-3\sqrt{7}}{-4\sqrt{7}+10\sqrt{7} -6\sqrt{7}} &= \bbvec{0}{0} \\
\bbvec{0}{0} &= \bbvec{0}{0}
\end{align*}
So it is in $\nul A$. \blacksquare
\end{book}

\begin{theorem}[Subspaces are spans.]
\label{theorem:subspaces-are-spans}

If $W$ is a subspace of $\R^n$, then $W=\spn (\vec{v}_1,\ldots, \vec{v}_c)$ for some linearly independent list of vectors $\vec{v}_1,\ldots,\vec{v}_c$ in $\R^n$, with $c \le n$.
\end{theorem}

\begin{notes}
\vfill 
\end{notes}

\begin{book}
\textbf{Proof.} Let $W$ be a subspace of $\R^n$. The following multi-step process shows how $\vec{v}_1,\ldots,\vec{v}_c$ can be constructed.

\textbf{Step 0.} If $W = \left\{\vec{0}\right\}$, then $W = \spn(\vec{0})$ and we're done. Otherwise, there exists some non-zero vector in $W$; call it $\vec{v}_1$.

\textbf{Step 1.} The list $\vec{v}_1$ of a single nonzero vector is linearly independent. Since $W$ is closed under scalar multiplication, we know $\spn ( \vec{v}_1 )$ is a subset of $W$. If $W = \spn ( \vec{v}_1 )$, we're done. Otherwise, there exists some vector $\vec{v}_2$ in $W$ but not in $\spn ( \vec{v}_1 )$.

\textbf{Step 2.} The list $\vec{v}_1,\vec{v}_2$ is linearly independent by theorem \ref{theorem:one-vector-is-independent}, since $\vec{v}_1 \ne \vec{0}$ and $\vec{v}_2$ is not a scalar multiple of $\vec{v}_1$.  Since $W$ is closed under scalar multiplication and vector addition, we know $\spn ( \vec{v}_1, \vec{v}_2 )$ is a subset of $W$. If $W = \spn ( \vec{v}_1, \vec{v}_2 )$, we're done. Otherwise, there exists some vector $\vec{v}_3$ in $W$ but not in $\spn ( \vec{v}_1, \vec{v}_2 )$.
\[\vdots\]
\textbf{Step $i$.} The list $\vec{v}_1,\vec{v}_2, \ldots, \vec{v}_i$ is linearly independent by theorem \ref{theorem:one-vector-is-independent}, since $\vec{v}_1 \ne \vec{0}$ and none of the other vectors are a linear combination of the preceding vectors.  Since $W$ is closed under scalar multiplication and vector addition, we know $\spn ( \vec{v}_1, \vec{v}_2, \ldots, \vec{v}_i )$ is a subset of $W$. If $W = \spn ( \vec{v}_1, \vec{v}_2, \ldots, \vec{v}_i )$, we're done. Otherwise, there exists some vector $\vec{v}_{i+1}$ in $W$ but not in $\spn ( \vec{v}_1, \vec{v}_2, \ldots, \vec{v}_i )$.
\[\vdots\]
This process must end before step $n+1$ because it's impossible to have a linearly independent list $\vec{v}_1, \vec{v}_2, \ldots, \vec{v}_{n+1}$ with more vectors than entries in each vector by theorem \ref{theorem:more-vectors-than-entries-dependent}.
\end{book}

\begin{notes}
\vfill
\end{notes}

Combining theorems \ref{theorem:spans-are-subspaces} and \ref{theorem:subspaces-are-spans} we get:

\begin{bbox}\textbf{Corollary.} A subset $W$ of $\R^n$ is a subspace if and only if it can be written as the span of a list of vectors.
\end{bbox}

\begin{notes}
\columnbreak
\end{notes}

\example Let $\vec{u}=(3,1)$ and $\vec{v}=(1,-1)$.
\begin{enumerate}[label=(\alph*)]
\begin{notes}
\item Draw each of the following points:

\bbtwocol{.3}{.7}{
  $\vec{u} + (0)\vec{v}$

  $\vec{u} + \vec{v}$

  $\vec{u} + 2\vec{v}$

  $\vec{u} + 2.5 \vec{v}$

  $\vec{u} + (-1)\vec{v}$

  $\vec{u} + (-2)\vec{v}$
  }{
\hspace{1pt}
\begin{tikzpicture}
\begin{axis}[
x=27pt,y=27pt,
axis lines=middle,
axis line style={-},
ymajorgrids=true,
%minor y tick num=3,
%yminorgrids=true,
xmajorgrids=true,
%minor x tick num=3,
%xminorgrids=true,
%xlabel=$x_1$,
%ylabel=$x_2$,
xmin=-1, xmax=6,
ymin=-2, ymax=3,
xtick={-1,0,...,6},
ytick={-2,-1,...,3},
]
\draw[-Stealth, very thick] (axis cs:0,0) -- (axis cs:3,1)
node[right] {$\vec{u}$};
\draw[-Stealth, very thick] (axis cs:0,0) -- (axis cs:1,-1)
node[right] {$\vec{v}$};
\end{axis}
\end{tikzpicture}
}
\end{notes}

\begin{book}
\item Draw each of the following points:

  $\vec{u} + (0)\vec{v}$

  $\vec{u} + \vec{v}$

  $\vec{u} + 2\vec{v}$

  $\vec{u} + 2.5 \vec{v}$

  $\vec{u} + (-1)\vec{v}$

  $\vec{u} + (-2)\vec{v}$
\end{book}

\item Is the set $W = \left\{ \vec{u}+t\vec{v} : t \text{ is in } \R \right\}$ a subspace of $\R^2$?
\end{enumerate}

\begin{notes}
\vfill
\end{notes}

\begin{book}
\bbtwocol{0.45}{0.52}{
\textbf{Solution.} (a) See the graph on the right.

(b) No, $W$ fails every one of the three conditions to be a subspace. (Can you explain why?)
}{
\begin{tikzpicture}
\begin{axis}[
x=24pt,y=24pt,
axis lines=middle,
axis line style={-},
ymajorgrids=true,
%minor y tick num=3,
%yminorgrids=true,
xmajorgrids=true,
%minor x tick num=3,
%xminorgrids=true,
%xlabel=$x_1$,
%ylabel=$x_2$,
xmin=-0.5, xmax=7.2,
ymin=-2.1, ymax=3.5,
xtick={-2,0,...,6},
ytick={-2,-1,...,3},
]
\fill (axis cs:3,1) circle (3pt) node[right] {$\vec{u} +0\vec{v}$};
\fill (axis cs:4,0) circle (3pt) node[above right] {$\vec{u} +\vec{v}$};
\fill (axis cs:5,-1) circle (3pt) node[right] {$\vec{u} +2\vec{v}$};
\fill (axis cs:5.5,-1.5) circle (3pt) node[right] {$\vec{u} +2.5\vec{v}$};
\fill (axis cs:2,2) circle (3pt) node[right] {$\vec{u} +(-1)\vec{v}$};
\fill (axis cs:1,3) circle (3pt) node[right] {$\vec{u} +(-2)\vec{v}$};
\draw[-Stealth, very thick] (axis cs:0,0) -- (axis cs:1,-1)
node[right] {$\vec{v}$};
\end{axis}
\end{tikzpicture}
}
\end{book}

\begin{bbox}
\textbf{Note:} In general, the set of points $\{\vec{u} + t \vec{v} : t \textrm{ is any real number}\}$ describes a line including the point $\vec{u}$, parallel to the vector $\vec{v}$, as an arrow.
\end{bbox}

\columnbreak
\exerciseheading

\exatend Suppose $A$ is an $r \times c$ matrix, $\vec{u}$ is in $\nul A$, and $\vec{v}$ is in $\col A$. Fill in each blank with $\vec{u}$ or $\vec{c}$.
\begin{enumerate}[label=(\alph*)]
\item The entries of \rule{2em}{0.8pt} are weights of a linear combination of columns of $A$.
\item \rule{2em}{0.8pt} is a linear combination of columns of $A$.
\item \rule{2em}{0.8pt} has $r$ entries.
\item \rule{2em}{0.8pt} has $c$ entries.
\item If \rule{2em}{0.8pt} is not $\vec{0}$, it means the columns of $A$ are linearly dependent vectors.
\end{enumerate}

\exatend Is the set $W=\left\{ ( x_1,x_2,x_3 ) : 7x_1-3x_2+\pi x_3= 0 \right\}$ a subspace of $\R^3$?

\exatend Is the set $Z=\left\{ ( x_1,x_2,x_3 ) : 7x_1-3x_2+\pi x_3= 8 \right\}$ a subspace of $\R^3$?

\exatend Is the set $U=\left\{ ( x_1,x_2,x_3 ) : x_1 = 3x_3 \right\}$ a subspace of $\R^3$?

\exatend Is the set $R = \left\{ (x+3y, z+x, y+2z, \pi x) : x,y, \text{ and } z \text{ are in } \R \right\}$ a subspace of $\R^4$?

\exatend Is the set $X=\left\{ ( x_1,x_2,x_3 ) : x_1 = x_3^2 \right\}$ a subspace of $\R^3$?

\exatend Is the unit square, $S = \left\{ (x_1,x_2) : 0 \le x_1 \le 1, 0 \le x_2 \le 1 \right\}$, a subspace of $\R^2$?

\begin{tikzpicture}[scale=3]
\draw[-] (-0.2,0) -- (1.3,0);
\draw[-] (0,-0.2) -- (0,1.3);
\fill[blue!20] (0,0) rectangle (1,1);
\draw[thick, blue] (0,0) rectangle (1,1);
\node at (0.5, 0.5) {$S$};
\draw[black,thick] (1,0.06) -- (1,-0.06) node[below] {$1$};
\draw[black,thick] (0.06,1) -- (-0.06,1) node[left] {$1$};
\end{tikzpicture}

\exatend Is Quadrant II, $Q_{II} = \left\{ (x_1, x_2) : x_1\le 0, x_2\ge 0 \right\}$, a subspace of $\R^2$?

\begin{tikzpicture}[scale=2.5]
\draw[-] (-1.5,0) -- (0.5,0);
\draw[-] (0,-0.5) -- (0,1.5);
\fill[blue!20] (0,0) rectangle (-1.5,1.5);
\node at (-1, 1) {$Q_{II}$};
\draw[blue,thick] (-1.5,0) -- (0,0);
\draw[blue,thick] (0,0) -- (0,1.5);
\draw[black,thick] (-1,0.06) -- (-1,-0.06) node[below] {$-1$};
%\draw[black,thick] (0.06,-1) -- (-0.06,-1) node[left] {$-1$};
%\draw[black,thick] (1,0.06) -- (1,-0.06) node[below] {$1$};
\draw[black,thick] (0.06,1) -- (-0.06,1) node[left] {$1$};
\end{tikzpicture}

\exatend Let $A = \begin{bmatrix}
1 & 3 & -7 \\
-3 & -9 & 23
\end{bmatrix}$. Find a few vectors in each of these subspaces:
\begin{enumerate}[label=(\alph*)]
\item $\col A$
\item $\row A$
\item $\nul A$
\end{enumerate}

\exatend Let $A = \begin{bmatrix}
1 & 3 & -7 \\
-3 & -9 & 23
\end{bmatrix}$.
\begin{enumerate}[label=(\alph*)]
\item Is $(-6,20)$ in $\col A$?
\item Is $(0,5,0)$ in $\row A$?
\item Is $(6,-2,0)$ in $\nul A$?
\end{enumerate}

\exatend Let $\vec{a}=(1,2)$ and $\vec{b}=(-2,1)$.
\begin{enumerate}[label=(\alph*)]
\item Draw the following points.

\bbtwocol{.3}{.7}{
$\vec{a} + 0\vec{b}$

$\vec{a} + 1\vec{b}$

$\vec{a} + 2\vec{b}$

$\vec{a} - \vec{b}$

$\vec{a} -1.8\vec{b}$
}{\begin{tikzpicture}
\begin{axis}[
x=22pt,y=22pt,
axis lines=middle,
axis line style={-},
ymajorgrids=true,
%minor y tick num=3,
%yminorgrids=true,
xmajorgrids=true,
%minor x tick num=3,
%xminorgrids=true,
%xlabel=$x_1$,
%ylabel=$x_2$,
xmin=-4, xmax=5,
ymin=0, ymax=4,
xtick={-4,-3,...,5},
ytick={0,1,...,4},
]
\end{axis}
\end{tikzpicture}}
\item Is the set of points $\left\{ \vec{a}+t\vec{b}: t \text{ is in } \R \right\}$ a subspace of $\R^2$?
\end{enumerate}


\exatend (T/F?) The set $\left\{ \vec{a}+t\vec{b}: t \text{ is in } \R \right\}$ describes a line including the point $\vec{b}$ parallel to the vector $\vec{a}$ (as an arrow).

\exatend (T/F?) The null space of $\begin{bmatrix}
\vec{a}_1 & \cdots & \vec{a}_c
\end{bmatrix}$ is the set of solutions $(x_1,\ldots,x_c)$ of the equation $x_1\vec{a}_1+\cdots+x_c\vec{a}_c = \vec{0}$.

\exatend (T/F?) The column space of $\begin{bmatrix}
\vec{a}_1 & \cdots & \vec{a}_c
\end{bmatrix}$ is the set of vectors $\vec{y}$ for which the equation $x_1\vec{a}_1+\cdots+x_c\vec{a}_c = \vec{y}$ is consistent.

\exatend (T/F?) If $A$ is an $m \times n$ matrix, then the row space of $A$ is a subspace of $\R^m$

\exatend (T/F?) If $A$ is a matrix, $\vec{u}$ is in the null space of $A$ and $\vec{x}$ is in the row space of $A$, then $\vec{u}$ and $\vec{x}$ have the same number of entries.

%bit of a challenge
\exatend Let $k$ be a real number. Prove that $U=\left\{ (x_1,x_2) : x_1=3x_2 + k\right\}$ is a subspace of $\R^2$ if and only if $k=0$.

% Trick: 0 in W means homogeneous eqn xu + yv = 0 has a non-trivial soln (1,t).
\exatend Prove that if the set $W = \left\{ \vec{u}+t\vec{v}: t \text{ is in } \R \right\}$ is a subspace then the list of vectors $\vec{u}, \vec{v}$ is linearly dependent.

% a bit of a challenge. Trick is to show it's a span of a vector, in either case.
\exatend Prove that if the list of vectors $\vec{u}, \vec{v}$ is linearly dependent, then the set $W = \left\{ \vec{u}+t\vec{v}: t \text{ is in } \R \right\}$ is a subspace. \hintfootnote{Consider the two separate cases: $\vec{u} = \vec{0}$ or $\vec{u} \ne \vec{0}$.}

% challenge. union of x and y-axes works: (x,y) : x=0 or y=0
\exatend Give an example of a non-empty subset of $\R^2$ that is closed under scalar multiplication, but is not a subspace of $\R^2$.

% challenge. Quadrant 1 works.
\exatend Give an example of a non-empty subset of $\R^2$ that is closed under vector addition but is not a subspace of $\R^2$.

% TODO: this might be onerous. Try it out.
\exatend Suppose $A$ is an $r \times c$ matrix, $\vec{u} = ( u_1, u_2, \ldots, u_c )$ is in the null space of $A$ and $(b_1,b_2,\ldots,b_c)$ is one of the rows of $A$.
\begin{enumerate}[label=(\alph*)]
\item Explain why $u_1b_1 + u_2b_2 + \cdots + u_cb_c=0$.
\item Suppose  $\vec{x} = ( x_1, x_2, \ldots, x_c )$ is in the row space of $A$. Explain why $u_1x_1+u_2x_2 + \cdots + u_cx_c = 0$.
\end{enumerate}


\exatend The \textbf{intersection} of two sets $P$ and $Q$ is defined as $P \cap Q = \left\{ x : x \text{ is in } P \text{ and } x \text{ is in } Q \right\}$. For example, $\spn(1,2) \cap \spn(3,0) = \left\{ \vec{0} \right\}$. And $\left\{ (x_1,x_2) : x_1 > 0 \right\} \cap \left\{ (x_1,x_2) : x_2>0 \right\}$ is quadrant I. Prove that if $P$ and $Q$ are subspaces of $\R^c$, then $P \cap Q$ is a subspace too.

\exatend \label{exercise:1.5-col-nul} \emph{(Challenge)} Give an example of a $4\times 4$ matrix $A$ so that $\col A = \nul A$.

\columnbreak
\section{Bases and Coordinates}
\label{sec:bases-and-coordinates}

\exercise (\emph{Warm-up}) Let $\vec{a}_1, \vec{a}_2$ be the columns of $A = \begin{bmatrix}
  1 & 3 \\
  0 & 4 \\
  0 & 2
\end{bmatrix}$.
\begin{enumerate}[label=(\alph*)]
\item Is the list $\vec{a}_1, \vec{a}_2 $ linearly independent? Why or why not?
\item Is $\spn (\vec{a}_1, \vec{a}_2) = \R^3$? In other words, does $(\vec{a}_1, \vec{a}_2)$ span $\R^3$? Why or why not?
\end{enumerate}

\begin{notes}
\vfill
\end{notes}

\exercise (\emph{Warm-up}) Let $\vec{b}_1, \vec{b}_2, \vec{b}_3$ be the columns of $B = \begin{bmatrix}
  2 & 4 & 6 \\
  0 & 2 & 2
\end{bmatrix}$.
\begin{enumerate}[label=(\alph*)]
\item Is $\vec{b}_1, \vec{b}_2, \vec{b}_3$ linearly independent? Why or why not?
\item Is $\spn (\vec{b}_1, \vec{b}_2, \vec{b}_3) = \R^2$? In other words, does $(\vec{b}_1, \vec{b}_2, \vec{b}_3)$ span $\R^2$? Why or why not?
\end{enumerate}

\begin{notes}
\vfill
\end{notes}

\exercise (\emph{Warm-up}) Let $\vec{e}_1,\vec{e}_2, \vec{e}_3$ be the columns of $I =\begin{bmatrix}
1 & 0 & 0 \\
0 & 1 & 0 \\
0 & 0 & 1
\end{bmatrix}$.
\begin{enumerate}[label=(\alph*)]
\item Is the list $\vec{e}_1,\vec{e}_2,\vec{e}_3$ linearly independent? Why or why not?
\item Is $\spn (\vec{e}_1, \vec{e}_2, \vec{e}_3) = \R^3$? In other words, does $(\vec{e}_1, \vec{e}_2, \vec{e}_3)$ span $\R^3$? Why or why not?
\end{enumerate}

\begin{notes}
\vfill \hspace{2pt}
\columnbreak
\end{notes}

\subsection{Bases and Coordinates: Main Ideas}\label{subsec:bases-and-coordinates-main-ideas}

\begin{definition}
Let $U$ be a subspace of $\R^n$. (Note that since $\R^n$ is a subspace of itself, this definition also applies to $U = \R^n$.) A list of vectors $\mathcal{B} = ( \vec{b}_1,\ldots, \vec{b}_c )$ is called a \textbf{basis} for $U$ if:
\begin{enumerate}[label=(\roman*)]
\item $\mathcal{B}$ is linearly independent.
\item $\spn \mathcal{B} = \R^n$.
\end{enumerate}
\end{definition}
\example Is $\mathcal{B} = \left( \bbbvec{0}{\pi}{0}, \bbbvec{3}{4}{0} \right)$ a basis of $\R^3$?

\begin{notes}
\vfill
\end{notes}

\begin{book}
\textbf{Solution.} Neither of the two vectors is a multiple of the other, so $\mathcal{B}$ is linearly independent. However, $\spn \mathcal{B} \ne \R^3$; for example $(0,0,1)$ is not in $\spn\mathcal{B}$. So $\mathcal{B}$ is not a basis of $\R^3$. \blacksquare
\end{book}

\example Is $\mathcal{B} = \left( \bbbvec{0}{\pi}{0}, \bbbvec{3}{4}{0} \right)$ a basis of $\spn \mathcal{B}$?

\begin{notes}
\vfill
\end{notes}

\begin{book}
\textbf{Solution.} This is the same $\mathcal{B}$ as the last example, and it is still linearly independent. This time, however, the second condition is not that $\spn \mathcal{B}=\R^3$, but that $\spn \mathcal{B}= \spn \mathcal{B}$, which is certainly true. So $\mathcal{B}$ is a basis of $\spn \mathcal{B}$. \blacksquare
\end{book}

\example Is $\mathcal{E} = \left( \bbbvec{1}{0}{0}, \bbbvec{0}{1}{0}, \bbbvec{0}{0}{1} \right)$ a basis of $\R^3$? (This is called the \textbf{standard basis} of $R^3$ and $I_3 =
\begin{bmatrix}
\vec{e}_1 & \vec{e}_2 & \vec{e}_3
\end{bmatrix}$ 
is called the $3\times 3$ \textbf{identity matrix}.)

\begin{notes}
\vfill \hspace{2pt}
\columnbreak
\end{notes}

\begin{book}
\textbf{Solution.} There are many ways to show that $\mathcal{E}$ is linearly independent and spans $\R^3$. One way is to notice that $\begin{bmatrix}
\vec{e}_1 & \vec{e}_2 & \vec{e}_3
\end{bmatrix}$ has a pivot in every column, so $\vec{E}$ is linearly independent by the uniqueness theorem (\ref{theorem:1.4uniqueness}). And Since $\begin{bmatrix}
\vec{e}_1 & \vec{e}_2 & \vec{e}_3
\end{bmatrix}$ has a pivot in every row, the existence theorem (\ref{theorem:1.4existence}) says $\mathcal{E}$ spans $\R^3$. \blacksquare
\end{book}

\example Is $\mathcal{B} = \left( \bbvec{2}{3}, \bbvec{7}{-1}, \bbvec{0}{10} \right)$ a basis of $\R^2$?

\begin{notes}
\vfill
\end{notes}

\begin{book}
\textbf{Solution.} Since $\mathcal{B}$ contains more vectors than entries in each vector, it is linearly dependent (theorem \ref{theorem:more-vectors-than-entries-dependent}), so it's not a basis. \blacksquare
\end{book}

\example Is $\mathcal{B} = \left( \bbvec{2}{3}, \bbvec{7}{-1}, \bbvec{0}{10} \right)$ a basis of $\spn \mathcal{B}$?

\begin{notes}
\vfill
\end{notes}

\begin{book}
\textbf{Solution.} It's still a linearly dependent list, so it's still not a basis. \blacksquare
\end{book}

\begin{definition}
Suppose $\mathcal{B} = (\vec{b}_1,\ldots,\vec{b}_c)$ is a basis for a subspace $U$ of $\R^n$, and that $\vec{y}$ is a vector in $U$. If 
\[ x_1\vec{b}_1+\cdots+x_c\vec{b}_c = \vec{y}\]
then the vector $(x_1,\ldots, x_c)$ is called the \textbf{$\mathcal{B}$-coordinate vector of $\vec{y}$} and we write $(x_1,\ldots,x_c) =\left[ \vec{y} \right]_{\mathcal{B}}$.
\end{definition}

\begin{notes}
\columnbreak
\end{notes}
\example $\mathcal B = \{\vec{b}_1, \vec{b}_2\} = \{(3,1), (-4,2)\}$ is a basis of $\R^2$. Use the graph to:
\begin{enumerate}[label=(\alph*)]
\item Find $\vec{y}$ if $[\vec{y}]_{\mathcal B} = (1,2)$
\begin{notes}
\vspace{1.5em}
\end{notes}
\item Find $[(-3,-1)]_{\mathcal B}$.
\begin{notes}
\vspace{1.5em}
\end{notes}
\item Find $[(10,0)]_{\mathcal{B}}$.
\begin{notes}
\vspace{1.5em}
\end{notes}
\item Find $[(7,-1)]_{\mathcal{B}}$.
\begin{notes}
\vspace{1.5em}
\end{notes}
\item Find $\left[  \vec{b}_{1}\right]_{\mathcal{B}}$
\begin{notes}
\vspace{1.5em}
\end{notes}
\end{enumerate}

\begin{notes}
\vfill
\end{notes}
\begin{tikzpicture}[scale=0.8]
\vectorgrid[xmin=-5,xmax=10,ymin=-3,ymax=6,step=1,labint=2,xu=3,yu=1,xv=-4,yv=2]
\node[above,blue] at (1.5,0.5){$b_1$};
\node[above,blue] at (-2,1){$b_2$};
\end{tikzpicture}

\begin{notes}
\vfill \hspace{2pt}
\columnbreak
\end{notes}

\begin{book}
\textbf{Solution.} (a) If $\left[ \vec{y} \right]_{\mathcal{B}}=(1,2)$, then $\vec{y}=1\vec{b}_1+2\vec{b}_2 = 1 \bbvec{3}{1} +2 \bbvec{-4}{2} = \bbvec{-5}{5}$.

(b) We want $x_1,x_2$ such that $x_1\vec{b}_1 + x_2\vec{b}_2=\bbvec{-3}{-1}$. But since $\bbvec{-3}{-1} = -\vec{b}_1$ (see what this looks like on the graph), we have $(x_1,x_2) = [(-3,-1)]_{\mathcal{B}} = (-1,0)$.

(c) On the grid, you can see that $(10,0) = 2\vec{b}_1-1\vec{b}_2$, so $[(10,0)]_{\mathcal{B}} = (2,-1)$.

(d) On the grid, you can see that $(7,-1) = \vec{b}_1-1b_2$, so $[(7,-1)]_{\mathcal{B}}=(1,-1)$.

(e) No need for the grid for this one. Since $\vec{b}_1=1\vec{b}_1+0\vec{b}_2$, we have $[\vec{b}_1]_{\mathcal{B}}=(1,0)$. \blacksquare
\end{book}

\begin{theorem}[Unique representation theorem.]
\label{theorem:unique-representation}

Let $\mathcal{B} = ( \vec{b}_1,\ldots,\vec{b}_c )$ be a basis for a subspace $U$ of $\R^n$. Then for each $\vec{y}$ in $U$, there is only one $\left[ \vec{y} \right]_{\mathcal{B}}$.
\end{theorem}

\textbf{Partial proof.}  Suppose $\vec{x}= (x_1,\ldots,x_c)=\left[ \vec{y} \right]_{\mathcal{B}}$ and $\vec{z}=(z_1,\ldots,z_c) = \left[ \vec{y} \right]_{\mathcal{B}}$. In other words, suppose $x_1\vec{b}_1+\cdots+x_c\vec{b}_c = \vec{y}$ and $z_1\vec{b}_1+\cdots+z_c\vec{b}_c = \vec{y}$.

Exercise \ref{exercise:prove-unique-representation} asks you to fill in the gaps below.

Since $\vec{y}-\vec{y}=\vec{0}$, we can write:

\begin{align*}
\text{(missing steps here!) } &= \vec{0} \\
\rule{4em}{0.8pt}\vec{b}_1+\cdots + \rule{4em}{0.8pt}\vec{b}_{c} &= \vec{0}
\end{align*}

Since the list $( \vec{b}_1,\ldots,\vec{b}_c )$ is linearly independent, we know all of the coefficients are zero: $\rule{4em}{0.8pt}=0, \ldots, \rule{4em}{0.8pt} =0$. Thus $\vec{x}=\vec{z}$ and $\left[ \vec{y} \right]_{\mathcal{B}}$ is unique. \blacksquare

\begin{theorem}[Removing dependent vectors doesn't change span.]
\label{theorem:removing-dependent-vectors}

Suppose $L = ( a_1,\ldots,a_n )$ is a list of vectors, and $U=\spn ( a_1,\ldots,a_n )$. If one of the vectors $\vec{a}_j$ is a linear combination of the other vectors, then removing $\vec{a}_j$ from the list doesn't change the span. In other words, $U=\spn ( \vec{a}_1, \ldots, \vec{a}_{j-1}, \vec{a}_{j+1}, \ldots, \vec{a}_n )$.
\end{theorem}

\textbf{Partial proof.} Suppose $\vec{a}_j$ is a linear combination of the other vectors: 
\[\vec{a}_j=x_1\vec{a}_1+\cdots+x_{j-1}\vec{a}_{j-1}+x_{j+1}\vec{a}_{j+1} + \cdots +x_n\vec{a}_n\]

To show that the span doesn't change, suppose that $\vec{u}= z_1\vec{a}_1+\cdots+z_n\vec{a}_n$ is in $U$. Exercise \ref{exercise:prove-removing-dependent} is to finish rest of the proof: show that $\vec{u}$ is in $\spn ( \vec{a}_1, \ldots, \vec{a}_{j-1}, \vec{a}_{j+1}, \ldots, \vec{a}_n )$.

\begin{notes}
\vfill \hspace{2pt}
\columnbreak
\end{notes}

\subsection{Bases for Column and Row Spaces}\label{subsec:bases-for-column-and-row-spaces}

The next two examples will help justify theorem \ref{theorem:basis-for-column-space} below.

\example \label{example:1.6columnsF}Let $F =
\begin{bmatrix}
\vec{f}_1 & \vec{f}_2 & \vec{f}_3 & \vec{f}_4 & \vec{f}_{5}
\end{bmatrix} = \begin{bmatrix}
1 & -5 & 0 & 0 & 6 \\
0 & 0 & 1 & 0 & 7 \\
0 & 0 & 0 & 1 & 8 \\
0 & 0 & 0 & 0 & 0
\end{bmatrix}$.
\begin{enumerate}[label=(\alph*)]
\item Explain why the list $( \vec{f}_1, \vec{f}_3, \vec{f}_4 )$ is linearly independent.
\item Explain why $\col F = \spn ( \vec{f}_1,\vec{f}_3,\vec{f}_4 )$.
\end{enumerate}

\begin{book}
\textbf{Solution.} (a) $\vec{f}_3$ is not a multiple of $\vec{f}_1$ because in the second position, the pivot entry 1 is not a multiple of 0. Likewise $\vec{f}_4$ is not a linear combo of $\vec{f}_1,\vec{f}_3$, because of the pivot 1 in the third position. Since $\vec{f}_1\ne \vec{0}$, theorem \ref{theorem:characterization-of-linear-dependence} implies $\vec{f}_1, \vec{f}_3, \vec{f}_4$ is a linearly independent list.

(b) Since $\vec{f}_2=-5\vec{f}_1$, and $\vec{f}_5=6\vec{f}_1+7\vec{f}_3+8\vec{f}_4$, we can use theorem \ref{theorem:removing-dependent-vectors} to remove $\vec{f}_2$ and $\vec{f}_5$ without changing the span: $\spn ( \vec{f}_1,\vec{f}_2,\vec{f}_3,\vec{f}_4,\vec{f}_5 ) = \spn ( \vec{f}_1,\vec{f}_3,\vec{f}_4 )$. \hfill \blacksquare
\end{book}

\begin{notes}
\vfill \hspace{2pt}
\columnbreak
\end{notes}

\example \label{example:1.6columnsA}Let $A =
\begin{bmatrix}
\vec{a}_1 & \vec{a}_2 & \vec{a}_3 & \vec{a}_4 & \vec{a}_{5}\end{bmatrix} =
\begin{bmatrix}
1 & -5 & 0 & -1 & -2 \\
1 & -5 & 0 & -1 & -2 \\
-2 & 10 & 0 & 3 & 12 \\
0 & 0 & 1 & 0 & 7
\end{bmatrix}$. 

$A$ is row-equivalent to $F$ from the previous exercise, which you can check if you want.
\begin{enumerate}[label=(\alph*)]
\item Explain why the list $( \vec{a}_1, \vec{a}_3, \vec{a}_4 )$ is linearly independent.
\item Explain why $\col A = \spn ( \vec{a}_1,\vec{a}_3,\vec{a}_4 )$.
\end{enumerate}

\begin{book}
\textbf{Solution.} (a) Suppose $x_1\vec{a}_1+x_3\vec{a}_3+x_4\vec{a}_4=\vec{0}$. The goal is to show this implies $(x_1,x_2,x_3) = (0,0,0)$. Since $A$ is row-equivalent to $F$, the same sequence of row operations converts $\begin{bmatrix}[ccc|c]
\vec{a}_1 & \vec{a}_3 & \vec{a}_4 & \vec{0}
\end{bmatrix}$ to $\begin{bmatrix}[ccc|c]
\vec{f}_1 & \vec{f}_3 & \vec{f}_4 & \vec{0}
\end{bmatrix}$, and so $(x_1,x_2,x_3)$ is also a solution of $x_1\vec{f}_1+x_3\vec{f}_3+x_4\vec{f}_4=\vec{0}$. Since we already showed the list $\vec{f}_1,\vec{f}_3,\vec{f}_4$ is independent, we can conclude $(x_1,x_2,x_3) = (0,0,0)$ and the list $\vec{a}_1,\vec{a}_3,\vec{a}_4$ is also independent.

(b) Here we'll use the fact that, since $A$ is row-equivalent to $F$, the following equations have the same solution set:
\begin{align}
x_1\vec{a}_1+x_2\vec{a}_2+x_3\vec{a}_3+x_4\vec{a}_4+x_5\vec{a}_5 &= \vec{0} \text{ and } \label{1.6eq1}\\
x_1\vec{f}_1+x_2\vec{f}_2+x_3\vec{f}_3+x_4\vec{f}_4+x_5\vec{f}_5 &= \vec{0} \label{1.6eq2}
\end{align}
Specifically, since 
\begin{equation}
\label{eq:1.6eq3}
\vec{f}_2=-5\vec{f}_1,\text{ and }\vec{f}_5=6\vec{f}_1+7\vec{f}_3+8\vec{f}_4
\end{equation}
we have two solutions of equation \ref{1.6eq2}: $(5,1,0,0,0)$ and $(-6,0,-7,-8,1)$, which are also solutions of \ref{1.6eq1} and so:
\begin{equation}
\label{eq:1.64}
\vec{a}_2=-5\vec{a}_1,\text{ and }\vec{a}_5=6\vec{a}_1+7\vec{a}_3+8\vec{a}_4\end{equation}
So, by theorem \ref{theorem:removing-dependent-vectors}, we can remove $\vec{a}_2$ and $\vec{a}_5$ without changing the span. \hfill \blacksquare
\end{book}

\begin{notes}
\vfill \hspace{2pt}
\columnbreak
\end{notes}

\begin{theorem}[Basis for column space.]
\label{theorem:basis-for-column-space}

The pivot columns of $A$ are a basis for the column space of $A$. (\emph{Note:} This doesn't say anything about ``pivot columns of an echelon matrix equivalent to $A$.'')
\end{theorem}

\begin{book}
\textbf{Proof.} It's cumbersome to write notation for this proof when we don't know what specific columns are pivot positions, so we'll be a little vague and refer to examples \ref{example:1.6columnsF} and \ref{example:1.6columnsA} for specifics. Suppose $A$ is row-equivalent to a row echelon matrix $F$. 
First let's show the pivot columns of $F$ and $A$ each makes and independent list of vectors. Each of the pivot columns of $F$ has a non-zero entry to the right of a row of zeros (like example \ref{example:1.6columnsF}(a)). So none of the pivot columns of $F$ is a linear combination of the pivot columns left of it, and since the first pivot column isn't $\vec{0}$, theorem \ref{theorem:characterization-of-linear-dependence} says the pivot columns of $F$ are linearly independent. Now, like in example \ref{example:1.6columnsA}(a), any linear combination of the pivot columns of $A$ that results in $\vec{0}$ corresponds with a linear combination of the pivot columns of $F$ that results in $\vec{0}$, and the only way to do that is with coefficients that are all $0$. So the pivot columns of $A$ are also linearly independent.

Now, let's show the pivot columns of $A$ span $\col A$, by showing that removing the non-pivot columns of $A$ doesn't change the span. The only non-zero entries in non-pivot columns of $F$  must be in positions right of pivot positions in a pivot column. So, like in example \ref{example:1.6columnsF}(b), the non-pivot columns of $F$ are linear combinations of pivot columns. Thus each non-pivot column corresponds with an equation in the form of \ref{eq:1.6eq3} above, which corresponds with a solution of
\[x_1\vec{f}_1+\cdots+x_c\vec{f}_c = \vec{0}\]
which is also a solution of 
\[x_1\vec{a}_1+\cdots+x_c\vec{a}_c = \vec{0}\]
because $A$ and $F$ are row-equivalent. Each of these solutions means the non-pivot columns of $A$ can be written as a linear combination of the pivot columns, like equation \ref{eq:1.64} above. Thus you can remove the non-pivot columns of $A$ without changing the span. \hfill \blacksquare
\end{book}

\begin{notes}
\vfill \hspace{2pt}
\columnbreak
\end{notes}

\example \label{example:1.6row-space} Label the rows of $F$ from example \ref{example:1.6columnsF} as:
\[\vec{\phi}_1, \vec{\phi}_2,\vec{\phi}_3,\vec{\phi}_4 = 
\begin{bmatrix}
1 \\ 5 \\ 0 \\ 0 \\ 6
\end{bmatrix},
\begin{bmatrix}
0 \\ 0 \\ 1 \\ 0 \\ 7
\end{bmatrix},
\begin{bmatrix}
0 \\ 0 \\ 0 \\ 1 \\ 8
\end{bmatrix},
\begin{bmatrix}
0 \\ 0 \\ 0 \\ 0 \\ 0
\end{bmatrix}
\]

And label the rows of $A$ from example \ref{example:1.6columnsA} as:
\[ \vec{\alpha}_1, \vec{\alpha}_2,\vec{\alpha}_3, \vec{\alpha}_4 =
\begin{bmatrix}
1 \\ -5 \\ 0 \\ -1 \\ -2
\end{bmatrix},
\begin{bmatrix}
1 \\ -5 \\ 0 \\ -1 \\ -2
\end{bmatrix},
\begin{bmatrix}
-2 \\ 10 \\ 0 \\ 3 \\ 12
\end{bmatrix},
\begin{bmatrix}
0 \\ 0 \\ 1 \\ 0 \\ 7
\end{bmatrix}\]
\begin{enumerate}[label=(\alph*)]
\item Explain why $\row F$ is a subset of $\row A$.
\item Explain why $\row A$ is a subset of $\row F$.
\item Explain why $\vec{\phi}_1, \vec{\phi}_2, \vec{\phi}_3$ is a basis for $\row F$ and thus $\row A$.
% Save this for an exercise! \item Is $\vec{b}_1, \vec{b}_2, \vec{b}_3$ a basis for $\row A$?
\end{enumerate}

\begin{book}
\textbf{Solution.} (a) When you do a multaddiplacement operation when reducing $A$ to $F$, you're replacing one a row with a linear combination of itself and another row. When you do a scaling operation, you're replacing a row with a linear combination of itself. So every row of $F$ is a linear combination of the rows of $A$. And every vector in $\row F$ is a linear combination of linear combinations of the rows of $A$, which is just a messier linear combination of the rows of $A$. So every vector in $\row F$ is in $\row A$.

(b) Every row operation is reversible, so there is a sequence of row operations that converts $F$ to $A$. The same logic as part (a) applies, and every vector in $\row A$ is in $\row F$. This and part (a) mean that $\row A = \row F$.

(c) Let's write them in reverse order: $\vec{\phi}_3,\vec{\phi}_2,\vec{\phi}_1$. Since $\vec{\phi}_2$ has a pivot position in position 3 and $\vec{\phi}_3$'s 3rd entry is the $0$ below that in $F$, $\vec{\phi}_2$ is not a multiple of $\vec{\phi}_3$. Likewise $\vec{\phi}_1$ has a pivot in position 1 which corresponds to zeros in the rows under it, so $\vec{\phi}_1$ is not a linear combination of $\vec{\phi}_2,\vec{\phi}_3$. Since $\vec{\phi}_3\ne \vec{0}$, theorem \ref{theorem:characterization-of-linear-dependence} shows $\vec{\phi}_3,\vec{\phi}_2,\vec{\phi}_1$ is linearly independent.

Also, removing $\vec{\phi}_4=\vec{0}$ clearly doesn't change the span, so:
\[\row A = \row F = \spn ( \vec{\phi}_1,\vec{\phi}_2,\vec{\phi}_3,\vec{\phi}_4 ) =\spn ( \vec{\phi}_1,\vec{\phi}_2,\vec{\phi}_3 ) \]
Thus the list of vectors $\vec{\phi}_1,\vec{\phi}_2,\vec{\phi}_3$ is a basis of $\row A$. \hfill \blacksquare
\end{book}

\begin{notes}
\vfill \hspace{2pt}
\columnbreak
\end{notes}

\begin{theorem}[Basis for row space.]
\label{theorem:basis-for-row-space}

If a matrix $A$ is row-equivalent to a row echelon matrix $F$, then $\row A = \row F$, and the pivot rows of $F$ make a basis for both $\row A$ and $\row F$.
\end{theorem}

\begin{book}
\textbf{Proof.} This proof is basically just an abstraction of the argument in example \ref{example:1.6row-space}, and that solution is easier to read than this proof. But this is included for completeness.

At every step in the process of row-reducing $A$ to $F$, the rows are a linear combination of the rows from the previous step. So every row of $F$ is a linear combination of the rows of $A$, and every vector in $\row F$ is a linear combination of linear combinations of rows of $A$. So $\row F$ is a subset of $\row A$. Likewise, these row operations can be reversed to transform $F$ to $A$. By the same argument, $\row A$ is a subset of $\row F$, and so $\row A = \row F$.

Now, let's say the top $k$ rows of $F$ are pivot rows, and write them in reverse order: $\vec{\phi}_k,\ldots,\vec{\phi}_1$, where $\vec{\phi}_k$ is the bottom-most pivot row and $\vec{\phi}_1$ is the top row of $F$. Since $\vec{\phi}_k$ is a pivot row, it's not $\vec{0}$. And since $\vec{\phi}_{k-1}$ has a pivot position above a 0 entry in $\vec{\phi}_k$, we know $\vec{\phi}_{k-1}$ is not a multiple of $\vec{\phi}_k$. Likewise each $\vec{\phi}_j$ is not a linear combination of $\vec{\phi}_k, \ldots, \vec{\phi}_{j+1}$, because row $\vec{\phi}_j$ has a pivot position above the 0 entries of $\vec{\phi}_k, \ldots, \vec{\phi}_{j+1}$. So theorem \ref{theorem:characterization-of-linear-dependence} shows that $\vec{\phi}_k, \ldots, \vec{\phi}_1$, the pivot rows of $F$, are linearly independent. And all of the non-pivot rows of $F$ are zero, so removing them doesn't change the span, and: 
\[ \row A = \row F = \spn ( \vec{\phi}_1,\ldots,\vec{\phi}_{k} )\]
\blacksquare
\end{book}

\begin{notes}
\vfill \hspace{2pt}
\columnbreak
\end{notes}

\exerciseheading

\textbf{\ref{subsec:bases-and-coordinates-main-ideas} Bases and Coordinates: Main Ideas}

\exatend Is the list of vectors $\bbvec{1}{2}, \bbvec{-2}{-4}, \bbvec{3}{0}$ a basis of $\R^2$?

\exatend Is the list of vectors $\bbvec{7}{1}, \bbvec{3}{0}$ a basis of $\R^2$?

\exatend Is the list of one vector $\bbvec{2}{5}$ a basis of $\R^2$? Is it a basis of a subspace of $\R^2$?

\exatend Is the list of vectors $\bbbvec{1}{2}{3}, \bbbvec{0}{-4}{5}, \bbbvec{-2}{0}{-1}$ a basis of $\R^3$? Is it a basis of a subspace of $\R^3$?

\exatend Is the list of vectors $\bbbvec{1}{2}{3}, \bbbvec{0}{-4}{5}, \bbbvec{-2}{0}{0}$ a basis of $\R^3$? Is it a basis of a subspace of $\R^3$?

\exatend Write your own list of vectors that spans $\R^4$ but isn't a basis of $\R^4$.

\exatend Write your own list of vectors in $\R^4$ that is linearly independent, but isn't a basis of $\R^4$.

\columnbreak
\exatend The grid below shows $\mathcal{B} = ( \vec{b}_1,\vec{b}_2 )$. Find:

\begin{horiparts}{2}
\item $\left[ (1,4) \right]_{\mathcal{B}}$
\item $\vec{y}$ if $\left[ \vec{y} \right]_{\mathcal{B}} = (1,4)$
\item $\left[ (5,-1) \right]_{\mathcal{B}}$
\item $\vec{y}$ if $\left[ \vec{y} \right]_{\mathcal{B}} = (5,-1)$
\end{horiparts}

\begin{tikzpicture}
\vectorgrid[xmin=-3,xmax=6,ymin=-3,ymax=5,step=1,labint=1,xu=1,yu=1,xv=-1,yv=2]
\node[right] at (1, 1) {$\vec{b}_1$};
\node[left] at (-1, 2) {$\vec{b}_2$};
\end{tikzpicture}

\exatend Recall that $\vec{e}_1$ means $(1,0)$ and $\vec{e}_2$ means $(0,1)$. In the grid above, with $\mathcal{B} = ( \vec{b}_1,\vec{b}_2 )$:
\begin{enumerate}[label=(\alph*)]
\item Note that $2\vec{b}_1-\vec{b}_2 = (3,0)$. Use this to find $\left[ \vec{e}_1 \right]_{\mathcal{B}}$.
\item Note that $\vec{b}_1+\vec{b}_2 = (0,3)$. Use this to find $\left[ \vec{e}_2 \right]_{\mathcal{B}}$.
\item Use your results from (a) and (b) to find $\left[ (4,2) \right]_{\mathcal{B}}$.
\end{enumerate}
% BEN: Use this problem and the one below for matrix inverses.

\columnbreak
\exatend The grid below shows $\mathcal{B} = ( \vec{b}_1,\vec{b}_2 )$. Find:

\begin{horiparts}{2}
\item $\left[ (-1,1) \right]_{\mathcal{B}}$
\item $\vec{y}$ if $\left[ \vec{y} \right]_{\mathcal{B}} = (-1,1)$
\item $\left[ (5,3) \right]_{\mathcal{B}}$
\item $\vec{y}$ if $\left[ \vec{y} \right]_{\mathcal{B}} = (5,3)$
\end{horiparts}

\begin{tikzpicture}
\vectorgrid[xmin=-3,xmax=6,ymin=-2,ymax=4,step=1,labint=1,xu=1,yu=1,xv=3,yv=1]
\node[right] at (1, 1) {$\vec{b}_1$};
\node[right] at (3, 1) {$\vec{b}_2$};
\end{tikzpicture}

\exatend Recall that $\vec{e}_1$ means $(1,0)$ and $\vec{e}_2$ means $(0,1)$. In the grid above, with $\mathcal{B} = ( \vec{b}_1,\vec{b}_2 )$:
\begin{enumerate}[label=(\alph*)]
\item Note that $-\vec{b}_1+\vec{b}_2 = (2,0)$. Use this to find $\left[ \vec{e}_1 \right]_{\mathcal{B}}$.
\item Note that $3\vec{b}_1-\vec{b}_2 = (0,2)$. Use this to find $\left[ \vec{e}_2 \right]_{\mathcal{B}}$.
\item Use your results from (a) and (b) to find $\left[ (\pi,\sqrt 2) \right]_{\mathcal{B}}$.
\end{enumerate}

\exatend Let $\mathcal{B}=( \vec{b}_1, \vec{b}_2 ) = \left( \bbvec{0}{2}, \bbvec{3}{-2} \right)$. Find:
\begin{enumerate}[label=(\alph*)]
\item $\left[ (-3,5) \right]_{\mathcal{B}}$
\item $\vec{y}$ if $\left[ \vec{y} \right]_{\mathcal{B}} = (-3,5)$
\item $\left[ (6,0) \right]_{\mathcal{B}}$
\item $\vec{y}$ if $\left[ \vec{y} \right]_{\mathcal{B}} = (6,0)$
\end{enumerate}

\exatend Let $\mathcal{B}=( \vec{b}_1, \vec{b}_2 ) = \left( \bbbvec{1}{2}{0}, \bbbvec{0}{-2}{4} \right)$. 
\begin{enumerate}[label=(\alph*)]
\item Explain why $\mathcal{B}$ is a basis of a subspace of $\R^3$. 
\item Find $\left[ (2,-2,12) \right]_{\mathcal{B}}$.
\item Find $\vec{y}$ if $\left[ \vec{y} \right]_{\mathcal{B}} = (0,1)$.
\item Find $\left[ (-3,-5,2) \right]_{\mathcal{B}}$.
\item Find $\vec{y}$ if $\left[ \vec{y} \right]_{\mathcal{B}} = (\pi,\sqrt 2)$.
\end{enumerate}

\exatend (T/F?) A basis of $U$ is a list of vectors that spans $U$ and is as small as possible.

\exatend (T/F?) A basis of $U$ is a linear independent list of vectors in $U$ that is as small as possible.

\exatend (T/F?) Every basis must contain $\vec{0}$.

\exatend (T/F?) Every linearly independent list of vectors is the basis of some subspace.

\exatend (T/F?) If $\mathcal{B} = ( \vec{b}_1, \ldots, \vec{b}_p )$ is a basis of a subspace $U$ of $\R^q$, and $\vec{y}$ is in $U$, then $\left[ \vec{y} \right]_{\mathcal{B}}$ is in $\R^q$.


\exatend (T/F?) If a list of vectors spans a subspace $H$, then some subset of the list is a basis of $H$. %True

\exatend (T/F?) If a list of vectors has \emph{more} vectors than entries in each vector, then it cannot be a basis of a subspace of $\R^n$.

\exatend (T/F?) If a  list of vectors has \emph{fewer} vectors than entries in each vector, then it cannot be a basis of a subspace of $\R^n$.

% problems from here to the end of the subsection are harder and should be at the end of the section.
\exatend Prove that every basis of $\R^n$ must have exactly $n$ vectors.

\exatend Prove or find a counterexample: Every subspace $U$ of $\R^n$ has a basis. %True, since (0) is a basis of {0}.

\exatend The $n \times n$ \textbf{identity matrix} is 
\[ I_n=
\begin{bmatrix}
1 & 0 & \cdots & 0 \\
0 & 1 & \cdots & 0 \\
\vdots & \vdots & \ddots & \vdots \\
0 & 0 & \cdots & 1
\end{bmatrix}\]
with $1$'s in the main diagonal and 0's everywhere else, and it's columns, $\mathcal{E} =  (\vec{e}_1, \vec{e}_2, \ldots, \vec{e}_n)$ are called the standard basis of $\R^n$. Suppose $\mathcal{B} = (  \vec{b}_1, \ldots, \vec{b}_n)$ is another basis of $\R^n$.
\begin{enumerate}[label=(\alph*)]
\item Explain why $\left[ \vec{b}_1 \right]_{\mathcal{B}} = \vec{e}_1, \ldots, \left[ \vec{b}_n \right]_{\mathcal{B}} = \vec{e}_n$.
\item Explain why $\left[ \vec{b}_1 \right]_{\mathcal{E}} = \vec{b}_1, \ldots, \left[ \vec{b}_n \right]_{\mathcal{E}} = \vec{b}_n$.
\end{enumerate}

\exatend \label{exercise:prove-unique-representation} Write the proof of theorem \ref{theorem:unique-representation} with the gaps filled in. 

\exatend \label{exercise:prove-removing-dependent} Write the finished proof of \ref{theorem:removing-dependent-vectors}.

\textbf{\ref{subsec:bases-for-column-and-row-spaces} Bases for Column and Row Spaces}

For problems \ref{exercise:1.6bases-start} through \ref{exercise:1.6bases-end}, find bases for $\nul A$, $\col A$, and $\row A$ for the given $A$.

\exatend \label{exercise:1.6bases-start}
$ A = \begin{bmatrix}
3 & 6 & 9 \\
-1 & -2 & -2
\end{bmatrix}$

\exatend $A =
\begin{bmatrix}
1 & 3 & -2 \\
-1 & -2 & 7 \\
-2 & -4 & 14
\end{bmatrix}$

\exatend $A =
\begin{bmatrix}
0 & 0 \\
\pi & -2\pi \\
3 & -6
\end{bmatrix}$

\exatend $A = \begin{bmatrix}
1 & 3 & -2 \\
-1 & -2 & 7 \\
-2 & -4 & 12
\end{bmatrix}$

\exatend $A = \begin{bmatrix}
2 & -6 & 4 \\
1 & -3 & 2 \\
-1 & 4 & 1 \\
-2 & 7 & -1
\end{bmatrix}$ 

\exatend \label{exercise:1.6bases-end} $A = \begin{bmatrix}
2 & 1 & -1 & -2 \\
-6 & -3 & 4 & 7 \\
4 & 2 & 1 & -1
\end{bmatrix}$


\exatend Consider the matrices $F$ and $A$ which are row-equivalent in examples  \ref{example:1.6columnsF} and \ref{example:1.6columnsA}. Are the pivot columns of $F$,  $(  \vec{f}_1,\vec{f}_3,\vec{f}_4)$, a basis of $\col A$? Explain.

\exatend Consider the matrices $F$ and $A$ which are row-equivalent in examples  \ref{example:1.6columnsF} and \ref{example:1.6columnsA}. Are the first three rows of $A$, the pivot rows, a basis of $\row A$? Explain.

\exatend (T/F?) If a matrix $A$ is row equivalent to a row echelon matrix $M$, then the pivot columns of $M$ are a basis of $\col A$.

\exatend (T/F?) If a matrix $A$ is row equivalent to a row echelon matrix $M$, then the pivot rows of $M$ are a basis of $\col A$.

\columnbreak

\section{Dimension}
\label{sec:dimension}

\exercise (\emph{Warm-up}) Explain why $\mathcal B = \left(\bbvec{2}{1}\right)$ is a not basis of $\R^2$.

\begin{notes}
\vfill
\end{notes}

\exercise (\emph{Warm-up}) Explain why $\mathcal B = \left(\bbvec{2}{1}, \bbvec{3}{-7}, \bbvec{0}{5}\right)$ is a not basis of $\R^2$.

\begin{notes}
\vfill
\end{notes}

\exercise (\emph{Warm-up}) Let $\mathcal B = \left( \bbbvec{1}{0}{1}, \bbbvec{0}{2}{2}\right)$. Note that since $\mathcal{B}$ is linearly independent, $\mathcal{B}$ is a basis of $\spn(\mathcal{B})$.

\begin{horiparts}{2}
\item Find $\vec{y}$ if $[\vec{y}]_{\mathcal B} = \bbvec{3}{-1}$.
\item Find $[\vec{y}]_{\mathcal B}$ if $\vec{y} = \bbbvec{1}{6}{7}$.
\item What do you think: is it possible to have a linearly independent list of 3 vectors in $\spn ( \mathcal{B} )$?
\end{horiparts}

\begin{notes}
\vfill \vfill \hspace{2pt}
\columnbreak
\end{notes}

\begin{theorem}[Properties of coordinates.]
\label{theorem:properties-of-coordinates}

Suppose $\mathcal{B} = ( \vec{b}_1, \ldots, \vec{b}_n )$ is a basis of $U$, a subspace of $\R^c$. Then, for any $\vec{y},\vec{z}$ in $U$ and any real number $k$:
\begin{enumerate}[label=(\roman*)]
\item $\left[ k\vec{y} \right]_{\mathcal{B}} = k \left[ \vec{y} \right]_{\mathcal{B}}$
\item $\left[ \vec{y} + \vec{z} \right]_{\mathcal{B}} =  \left[ \vec{y} \right]_{\mathcal{B}}+\left[ \vec{z} \right]_{\mathcal{B}}$
\item If $\left[ \vec{y} \right]_{\mathcal{B}} = \left[ \vec{z} \right]_{\mathcal{B}}$, then $\vec{y}=\vec{z}$ \quad (We say $f(\vec{y}) =\left[ \vec{y} \right]_{\mathcal{B}}$ is an \emph{injective} or \emph{one-to-one} function.)
\item If $\vec{x}$ is any vector in $\R^n$, then there is a $\vec{y}$ in $U$ so that $\left[ \vec{y} \right]_{\mathcal{B}}=\vec{x}$ \quad (We say $f(\vec{y}) =\left[ \vec{y} \right]_{\mathcal{B}}$ is a \emph{surjective} or \emph{onto} function.)
\end{enumerate}
\end{theorem}

\begin{book}
\textbf{Proof.} (i) Exercise \ref{exercise:prove-coordinate-scaling}.

(ii) Suppose $\vec{y}$ and $\vec{z}$ are in $U$, with $\left[ \vec{y} \right]_{\mathcal{B}}= (x_1,\ldots,x_n)$ and $\left[ \vec{z} \right]_{\mathcal{B}}= (w_1,\ldots,w_n)$. In other words, assume: \[\vec{y}=x_1\vec{b}_1 + \cdots x_n\vec{b}_n \quad \text{ and }\quad \vec{z}=w_1\vec{b}_1 + \cdots w_n\vec{b}_n\]
Then
\begin{align*}
\vec{y}+\vec{z} &= x_1\vec{b}_1 + \cdots x_n\vec{b}_n + w_1\vec{b}_1 + \cdots w_n\vec{b}_n \\
&=(x_1+w_1)\vec{b}_1 + \cdots + (x_n+w_n)\vec{b}_n
\end{align*}
which means: \[\left[ \vec{y}+\vec{z} \right]_{\mathcal{B}} = \bbbvec{x_1+w_1}{\vdots}{x_n+w_n} = \bbbvec{x_1}{\vdots}{x_n} + \bbbvec{w_1}{\vdots}{w_n}  = \left[ \vec{y} \right]_{\mathcal{B}}+\left[ \vec{z} \right]_{\mathcal{B}}\]

(iii) Suppose $\left[ \vec{y} \right]_{\mathcal{B}} = (x_1,\ldots,x_n)$ and $\left[ \vec{z} \right]_{\mathcal{B}}  = (x_1,\ldots,x_n)$. In other words, suppose:
\[\vec{y}=x_1\vec{b}_1+\cdots + x_n\vec{b}_n \quad \text{ and } \quad \vec{z}=x_1\vec{b}_1+\cdots + x_n\vec{b}_n\]
Then, of course, $\vec{y}=\vec{z}$.

(iv) Exercise \ref{exercise:prove-coordinate-surjective}. \blacksquare
\end{book}

\begin{notes}
Proofs of (i) and (iv) are exercises, so let's prove (ii) and (iii).
\vfill \hspace{2pt}
\columnbreak
\end{notes}

\begin{theorem}[More vectors than a basis are dependent.]
\label{theorem:more-vectors-than-a-basis}
If a subspace $U$ of $\R^c$ has a basis with $n$ vectors in it, then any list with more than $n$ vectors in $U$ is linearly dependent.
\end{theorem}

\begin{notes}
\vfill \vfill
\end{notes}

\begin{book}
\textbf{Proof.} Let's call the basis $\mathcal{B}=(  \vec{b}_1, \ldots ,\vec{b}_n)$, and suppose $\vec{z}_1, \ldots, \vec{z}_j$ is a list of vectors in $U$ with $j>n$. Recall the definition of $\left[ \vec{w} \right]_{\mathcal{B}}$ is the vector $(x_1,\ldots,x_n)$ such that $x_1\vec{b}_1+\cdots+x_nb_n = \vec{w}$. So, by definition, $\left[ \vec{z}_1 \right]_{\mathcal{B}},\ldots,\left[ \vec{z}_j \right]_{\mathcal{B}}$ is a list of $j$ vectors in $\R^n$, and since $j>n$, theorem \ref{theorem:more-vectors-than-entries-dependent} says the list is dependent in $\R^n$. So there exists $(y_1,\ldots,y_j)\ne \vec{0}$, such that:
\[y_1 \left[ \vec{z}_1 \right] + \cdots + y_j \left[ \vec{z}_j \right] = \vec{0}_n \]
(Here, $\vec{0}_n$ denotes the zero vector in $\R^n$, with $n$ zeros.) Using theorem \ref{theorem:properties-of-coordinates} parts (i) and (ii), we can rewrite the left side:
\begin{align*}
y_1 \left[ \vec{z}_1 \right]_{\mathcal{B}} + \cdots + y_j \left[ \vec{z}_j \right]_{\mathcal{B}} &= \vec{0}_n \\
\left[ y_1\vec{z}_1 \right]_{\mathcal{B}} + \cdots + \left[ y_j\vec{z}_j \right]_{\mathcal{B}} &= \vec{0}_n \\
\left[ y_1\vec{z}_1  + \cdots +  y_j\vec{z}_j \right]_{\mathcal{B}} &= \vec{0}_n 
\end{align*}
Now, since $\left[ \vec{0}_c \right]_{\mathcal{B}} = \vec{0}_n$ theorem \ref{theorem:properties-of-coordinates} part (iii) says
\[y_1\vec{z}_1  + \cdots +  y_j\vec{z}_j = \vec{0}_c\]
which means $(y_1,\ldots,y_j)$ is a non-trivial solution to the homogeneous equation and the list $( \vec{z}_1,\ldots,\vec{z}_j )$ is linearly dependent. \blacksquare
\end{book}

\begin{theorem}
\label{theorem:all-bases-same-size}
If $U$ is a subspace of $\R^c$, then every basis of $U$ has the same number of vectors.
\end{theorem}

\begin{notes}
\vfill 
\end{notes}

\begin{book}
\textbf{Proof.} We'll do a ``proof by contradiction,'' which means we'll show that it's impossible that the theorem is false. Assume that $\mathcal{A}$ and $\mathcal{B}$ are each bases of $U$. Then, if one of them (say, $\mathcal{A}$) has more vectors than the other (say, $\mathcal{B}$), then $\mathcal{A}$ has more vectors than a basis and theorem \ref{theorem:more-vectors-than-a-basis} says $\mathcal{A}$ is linearly dependent and thus not a basis. So, it's impossible to have bases with two different sizes.
\end{book}

\begin{definition}
The \textbf{dimension} of a subspace $U$ of $\R^c$, written $\dim U$, is the number of vectors in a basis of $U$. If $\dim U = n$, we say $U$ is $n$-dimensional.
\end{definition}

\begin{notes}
\columnbreak
\end{notes}
\example Find the dimension of $\R^3$.

\begin{notes}
\vfill
\end{notes}

\begin{book}
\textbf{Solution.} Recall that $(\vec{e}_1,\vec{e}_2,\vec{e}_3) = \left( \bbbvec{1}{0}{0}, \bbbvec{0}{1}{0}, \bbbvec{0}{0}{1} \right)$ is a basis of $\R^3$, called the standard basis. Since it has three vectors in it, $\dim \R^3=3$. \blacksquare
\end{book}

\example Find the dimension of the solution set of $x_{1}+8x_2-x_3=0$ in $\R^3$.

\begin{notes}
\vfill \vfill
\end{notes}

\begin{book}
\textbf{Solution.} Let's use the strategy of section $\ref{sec:homogeneous-systems-and-linear-independence}$ to describe the solution set as the span of some vectors. The augmented matrix $\begin{bmatrix}[ccc|c]
1 & 8 & -1 & 0
\end{bmatrix}$ is already in reduced row echelon form, so the next step is to write all of the variables in terms of the free variables $x_2$ and $x_3$:
\[\bbbvec{x_1}{x_2}{x_3}= \bbbvec{-8x_2+x_3}{x_2}{x_3} = x_2 \bbbvec{-8}{1}{0} +x_3 \bbbvec{1}{0}{1}\]
So the solution set is the span of $(-8,1,0)$ and $(1,0,1)$. And since each of these vectors has a 1 in the position the other has a 0, they're not multiples of each other, and form a linearly independent list, and thus a basis of the solution set. So the dimension of the solution set is $2$. \blacksquare
\end{book}

\example Find the dimension of $U = \left\{
\begin{bmatrix}
2p+t \\
-7p +4q-2r \\
6q-3r+5t \\
8t
\end{bmatrix} : p,q,r,t \text{ in } \R
\right\}$.

\begin{notes}
\vfill \vfill \hspace{2pt}
\columnbreak
\end{notes}

\begin{book}
\textbf{Solution.} Any vector in $U$ can be written:
\[p \begin{bmatrix}
2 \\ -7 \\ 0 \\ 0
\end{bmatrix} + q\begin{bmatrix}
0 \\ 4 \\ 6 \\ 0
\end{bmatrix} + r\begin{bmatrix}
0 \\ -2 \\ -3 \\ 0
\end{bmatrix}+t\begin{bmatrix}
1 \\ 0 \\ 5 \\ 8 
\end{bmatrix}\]
So we have a list of four vectors which spans $U$:
\[\vec{b}_1=\begin{bmatrix}
2 \\ -7 \\ 0 \\ 0
\end{bmatrix}, \quad \vec{b}_2=\begin{bmatrix}
0 \\ 4 \\ 6 \\ 0
\end{bmatrix}, \quad \vec{b}_3=\begin{bmatrix}
0 \\ -2 \\ -3 \\ 0
\end{bmatrix}, \quad \vec{b}_4=\begin{bmatrix}
1 \\ 0 \\ 5 \\ 8 
\end{bmatrix}\]
This suggests that $\dim U = 4$, but a careful inspection shows that the list is linearly independent. In fact, $\vec{b}_3=-0.5\vec{b}_2$. We can remove $\vec{b}_3$ from the list, and the smaller list will still span $U$, since any vector in $U$, written $p\vec{b}_1+q\vec{b}_2+r\vec{b}_3+t\vec{b}_4$, can also be written $p\vec{b}_1+(q-0.5r)\vec{b}_2+t\vec{b}_4$. And since there's a pivot position in every column of $\begin{bmatrix}
\vec{b}_1 & \vec{b}_2 & \vec{b}_4
\end{bmatrix}$, we know the list is independent, and thus a basis of $U$. So $\dim U= 3$. \blacksquare
\end{book}

We'll need the following lemma for the ``buy two, get one free'' theorem below.

\begin{bbox}
\textbf{Lemma.} Suppose $U$ is a subspace of $\R^c$, with $\dim U = n$, and $\mathcal{V} = \vec{v}_1,\ldots,\vec{v}_n$ is a list of $n$ vectors in $U$. Then $\mathcal{V}$ is linearly dependent if and only if $\mathcal{V}$ doesn't span $U$.
\end{bbox}

\begin{book}
\textbf{Proof.} The edge case where $n=0$, $U=\left\{ \vec{0} \right\}$, and $V = ()$ is left as exercise \ref{exercise:1.7edge-case}. So we'll assume $n >0$. Also, if all the vectors in $\mathcal{V}$ are $\vec{0}$, then $\mathcal{V}$ is certainly linearly dependent and $\spn ( \mathcal{V} ) = \left\{ \vec{0} \right\} \ne U$. So we can assume $\mathcal{V}$ has at least one non-zero vector in it.

(\rightarrow) Suppose $\mathcal{V}$ is linearly dependent. By theorem \ref{theorem:characterization-of-linear-dependence}, one of the vectors in $\mathcal{V}$ can be written as a linear combination of the others, so we can remove this vector without changing the span (theorem \ref{theorem:removing-dependent-vectors}). If the remaining list is still dependent, repeat until you have an independent list $\mathcal{W}$. (We know this process has to stop before $\mathcal{W}$ is empty, because $\mathcal{V}$ has at least one non-zero vector in it.) 

Now, $\mathcal{W}$ can't span $U$, because if it did, it would be a basis of $U$ with fewer vectors than a basis ($n = \dim U$). Thus:

 \[\spn \mathcal{V} = \spn \mathcal{W} \ne U\]

(\leftarrow) Suppose $\mathcal{V}$ doesn't span $U$. Then there is some vector $\vec{u}$ in $U$ but not in $\spn \mathcal{V}$. Now, the list $\vec{v}_1,\ldots,\vec{v}_n, \vec{u}$ contains $n+1$ vectors in it, so it's linearly dependent by theorem \ref{theorem:more-vectors-than-a-basis}. And by theorem \ref{theorem:characterization-of-linear-dependence} one of the vectors in $\vec{v}_1,\ldots,\vec{v}_n, \vec{u}$ is a linear combination of the other vectors, and it ain't $\vec{u}$, since $\vec{u}$ isn't in $\spn ( \vec{v}_1,\ldots,\vec{v}_n )$. Thus one of the vectors in $\mathcal{V}$ is a linear combination of the other vectors in $\mathcal{V}$, so $\mathcal{V}$ is linearly dependent. \blacksquare
\end{book}

\begin{notes}
\vfill
\end{notes}
\begin{theorem}[Buy two, get one free.]
\label{theorem:buy-two-get-one-free}

Suppose $\vec{v}_1, \ldots, \vec{v}_n$ is a list of $n$ vectors in a subspace $U$ of $\R^c$. If any two of the following are true, then so is the third:
\begin{enumerate}[label=(\roman*)]
\item $\dim U = n$.
\item $\vec{v}_1,\ldots,\vec{v}_n$ is linearly independent.
\item $\vec{v}_1,\ldots,\vec{v}_n$ spans $U$.
\end{enumerate}
\end{theorem}

\begin{notes}
\columnbreak
\end{notes}

\begin{book}
\textbf{Proof.}  If (ii) and (iii) are true, then $\vec{v}_1,\ldots,\vec{v}_n$ is a basis of $U$, and by definition, $\dim U = n$.

If (i) and (ii) are true, then the above lemma shows $\vec{v}_1,\ldots,\vec{v}_n$ spans $U$.

If (i) and (iii) are true, then the above lemma shows $\vec{v}_1,\ldots,\vec{v}_n$ is linearly independent. \blacksquare
\end{book}

\begin{definition}
The dimension of the column space of $A$ is called the \textbf{rank} of $A$.
\end{definition}

Recall the pivot columns of $A$ are a basis of $\col A$, and the pivot rows of an echelon matrix equivalent to $A$ are a basis of $\row A$. And the way echelon form works, you can't have more than one pivot in single row or column. Thus the number of pivot rows of $A$ equals the number of pivot columns of $A$, and $\dim(\row A) = \dim(\col A) = \rank A$.

\example Find $\rank A$ and $\dim(\nul A)$ for $A =\begin{bmatrix}
1 & 3 & -5 & 1 & 2 \\
-2 & -6 & 7 & 0 & -10
\end{bmatrix}$.

\begin{notes}
\vfill
\end{notes}

\begin{book}
\textbf{Solution.} Theorem \ref{theorem:basis-for-column-space} says that the pivot columns of $A$ form a basis for $\col A$. Let's row-reduce $A$ to find the pivot columns:
\[\begin{bmatrix}
1 & 3 & -5 & 1 & 2 \\
-2 & -6 & 7 & 0 & -10
\end{bmatrix} \sim \begin{bmatrix}
1 & 3 & -5 & 1 & 2 \\
0 & 0 & -3 & 2 & -6
\end{bmatrix}\]
So we can see the first and third columns of $A$ are pivot columns, so $\left( \bbvec{1}{-2}, \bbvec{-5}{7} \right)$ is a basis of $A$ and $\rank A = \dim (\col A) = 2$. Now let's find a basis for the null space by continuing to row-reduce the augmented matrix:
\begin{align*}
\begin{bmatrix}[ccccc|c]
1 & 3 & -5 & 1 & 2 & 0\\
-2 & -6 & 7 & 0 & -10 & 0
\end{bmatrix} &\sim \begin{bmatrix}[ccccc|c]
1 & 3 & -5 & 1 & 2 & 0\\
0 & 0 & -3 & 2 & -6 & 0
\end{bmatrix}\\ &\sim \begin{bmatrix}[ccccc|c]
1 & 3 & -5 & 1 & 2 & 0\\
0 & 0 & 1 & -2/3 & 2 & 0
\end{bmatrix}\\ &\sim \begin{bmatrix}[ccccc|c]
1 & 3 & 0 & -7/3 & 12 & 0\\
0 & 0 & 1 & -2/3 & 2 & 0
\end{bmatrix}
\end{align*}
This corresponds to the system $\begin{aligned}
x_1+3x_2-7/3x_4+12x_5 &=0 \\
x_3-2/3x_4+2x_5 &= 0
\end{aligned}$ and we can write all the variables in terms of the free variables $x_2,x_4,x_5$:
\[\begin{bmatrix}
x_1 \\ x_2 \\ x_3 \\ x_4 \\ x_5
\end{bmatrix}=\begin{bmatrix}
-3x_2+7/3x_4-12x_5 \\ x_2 \\ 2/3x_4-2x_5 \\ x_4 \\ x_5
\end{bmatrix}=x_2\begin{bmatrix}
-3 \\ 1 \\ 0 \\ 0 \\ 0
\end{bmatrix}+x_4\begin{bmatrix}
+7/3 \\ 0 \\ 2/3 \\ 1 \\ 0
\end{bmatrix}+x_5\begin{bmatrix}
-12 \\ 0 \\ -2 \\ 0 \\1
\end{bmatrix}\]
Each of these three vectors on the right has a 1 in the place the others have a 0, so this method always produces a list of linearly independent vectors that spans the null space. In this case, $\dim(\nul A) = 3$, which is the number of free variables (non-pivot columns of $A$). \blacksquare

\end{book}

The next theorem is quite useful.
Note that $\dim(\nul A)$ is sometimes called the \emph{nullity} of $A$.

\begin{theorem}[The rank-nullity theorem.]
\label{theorem:rank-nullity}

If $A$ has $c$ columns, then: \[\rank A + \dim(\nul A) = c\]

\end{theorem}

\begin{notes}
\vspace{1in} 
\hspace{2pt}
\columnbreak 
\end{notes}
\begin{book}
\textbf{Proof.} We already established that \[\rank A = (\text{\# of pivot cols of }A).\] And since 
\[(\text{\# of pivot cols of }A) + (\text{\# of non-pivot cols of }A) = c,\]
we just need to show that \[\dim(\nul A) = (\text{\# of non-pivot cols of }A).\]
Every non-pivot column of $A$ corresponds with a free variable. Using the strategy laid out in section 1.4, where you row-reduce $\begin{bmatrix}[c|c]
A & \vec{0}
\end{bmatrix}$ to reduced echelon form, every non-pivot column corresponds to a free variable, and each free variable leads to a different linearly independent vector in the basis of $\nul A$. \blacksquare
% TODO: Show these vectors are actually linearly independent.
\end{book}

\exerciseheading

\exatend (T/F?) Since all vectors in $\R^4$ have four entries, every subspace of $\R^4$ is 4-dimensional.

\exatend (T/F?) If $A$ is a $7\times 5$ matrix, then $\dim(\nul A) = 5$.

\exatend (T/F?) For any matrix $A$, $\dim(\col A) = \dim(\row A)$.

\exatend (T/F?) For any matrix $A$, $\dim(\nul A) = \dim(\row A)$.

\exatend (T/F?) Every linearly independent list of $3$ vectors in $\R^5$ is the basis of some $3$-dimensional subspace of $\R^5$.

For exercises \ref{exercise:1.7dimensions-start} through \ref{exercise:1.7dimensions-end}, find $\dim(\col A), \dim(\row A)$, and $\dim(\nul A)$.

\exatend \label{exercise:1.7dimensions-start} $A =\begin{bmatrix}
1 & 2 \\
3 & 4
\end{bmatrix}$

\exatend $A =\begin{bmatrix}
1 & 3 & -2 \\
4 & 12 & -8
\end{bmatrix}$

\exatend $A = \begin{bmatrix}
3 & 7 \\
0 & -8 \\
3 & -1 \\
1 & 0
\end{bmatrix}$

\exatend $A = \begin{bmatrix}
-1 & 4 & 8 & 34 \\
-1 & 5 & 11 & 43 \\
1 & -2 & -2 & -16
\end{bmatrix}$

\exatend $A = \begin{bmatrix}
-1 & -1 & 1 \\
4 & 5 & -2 \\
8 & 11 & -2 \\
34 & 43 & -16
\end{bmatrix}$

\exatend $A =\begin{bmatrix}
3 & -2 & -4 \\
-1 & 1 & 2 \\
6 & -2 & -4
\end{bmatrix}$

\exatend \label{exercise:1.7dimensions-end} $A = \begin{bmatrix}
3 & -2 & -4 \\
-1 & 1 & 2 \\
6 & -2 & -5
\end{bmatrix}$

\exatend If $A$ is $7 \times 5$, what is the largest possible $\dim(\row A)$?

\exatend If $A$ is $7 \times 5$, and $\dim(\nul A) = 1$, how many rows of $A$ can be removed without changing the row space?

\exatend Write a $4 \times 3$ matrix $A$ so that $\dim(\nul A) = 0$ and $\rank A=3$.

\exatend Write a $4 \times 3$ matrix $A$ so that $\dim(\nul A) = 2$ and $\rank A=1$.

\exatend Write a $4 \times 3$ matrix $A$ so that $\dim(\nul A) = 3$ and $\rank A=0$.

\exatend (T/F?) If $U$ is a subspace of $\R^c$, and the list $\vec{a}_1, \ldots, \vec{a}_n$ spans $U$, then any list of $n+1$ vectors in $U$ is linearly dependent.

\exatend (T/F?) Every 1-dimensional subspace of $\R^2$ has either $( \vec{e}_1 )$ or $( \vec{e}_2 )$ as its basis.

\exatend (T/F?) If $A$ is a \emph{square} $n \times n$ matrix, then the columns of $A$ span $\R^n$ if and only if they are linearly independent.

\exatend (T/F?) If $U$ is a subspace of $\R^n$, and $\vec{u}_1, \ldots, \vec{u}_j$ is a list of $j$ linearly independent vectors in $U$, then $\dim U \ge j$.

\exatend (T/F?) If $U$ is a subspace of $\R^n$, and $\vec{u}_1, \ldots, \vec{u}_j$ is a list of $j$ linearly dependent vectors in $U$, then $\dim U \le j$.

\exatend (T/F?) If $U$ is a subspace of $\R^n$, $\vec{u}_1, \ldots, \vec{u}_j$ is a list of $j$ vectors in $U$, and $\spn (  \vec{u}_1, \ldots, \vec{u}_j) = U$, then $\dim U \le j$.

\exatend (T/F?) If $U$ is a subspace of $\R^n$, $\vec{u}_1, \ldots, \vec{u}_j$ is a list of $j$ vectors in $U$ that \emph{doesn't} span $U$, then $\dim U \ge j$.

\exatend Frank solves a homogeneous system of 30 equations with 31 variables, and finds that every solution can be written as a multiple of a single vector. Can Frank be sure that, if he were to replace all the zeros on the right sides with any list of 30 real numbers, that the new system would always have a solution?

\exatend Recall that $\left\{ (0,0,0) \right\}$ is a subspace of $\R^3$. What is the dimension of this subspace? (Find a basis for it. You may need to look up the definitions of span and linear independence.)

\exatend \label{exercise:1.7edge-case} Prove the lemma before the ``buy two, get one free'' theorem, for the edge case where $n=0$ and $V=()$.

\exatend \label{exercise:prove-coordinate-scaling} Prove theorem \ref{theorem:properties-of-coordinates} part (i).

\exatend \label{exercise:prove-coordinate-surjective} Prove theorem \ref{theorem:properties-of-coordinates} part (iv).

\exatend Prove that we can add part (iv) onto the existence theorem below. We already proved that (i), (ii), and (iii) are equivalent, so it would be sufficient to show that (i) is equivalent to (iv).

\begin{theorem}[The existence theorem v2.]
\label{theorem:1.7existence-v2}

Suppose $A =
\begin{bmatrix}
\vec{a}_1 & \cdots & \vec{a}_c
\end{bmatrix}
$ is an $r \times c$ matrix. The following are equivalent:
\begin{enumerate}[label=(\roman*)]
\item $A$ has a pivot position in every row.
\item Every system of equations that has $A$ as its coefficient matrix is consistent.
\item $\spn(\vec{a}_1,\ldots,\vec{a}_c) = \R^r$.
\item $\dim(\nul A) = c-r$.
\end{enumerate}
\end{theorem}

\exatend Explain why we can add part (iv) to the uniqueness theorem below. We already proved that (i), (ii), and (iii) are equivalent, so it would be sufficient to show that (i) is equivalent to (iv).

\begin{theorem}[The uniqueness theorem v2.]
\label{theorem:1.7uniqueness-v2}

Suppose $A =
\begin{bmatrix}
\vec{a}_1 & \cdots & \vec{a}_c
\end{bmatrix}
$ is an $r \times c$ matrix. The following are equivalent:
\begin{enumerate}[label=(\roman*)]
\item $A$ has a pivot position in every column.
\item $\nul A = \left\{ \vec{0} \right\}$.
\item The columns $\vec{a}_1, \ldots, \vec{a}_c$ are a linearly independent list of vectors.
\item $\dim(\col A) = c$
\end{enumerate}
\end{theorem}

\exatend Exercise \ref{exercise:1.5-col-nul} in 1.5 asked you to find a $4\times4$ matrix $A$ with $\col A = \nul A$. Show that it's impossible to find a $5\times 5$ matrix $A$ with $\col A = \nul A$.

%axler
\exatend Explain why there does not exist a $2\times 4$ matrix whose null space is $\left\{ (x_1,x_2,x_3,x_4:x_1=x_3) \right\}$.

% tricky! (axler) Soln: a list of the 3 basis vectors u1,u2,u3 plus the 3 basis vectors v1,v2,v3 makes a list of 6 vectors, which is dependent since 6>5. x1u1+x2u2+x3u3+x4v1+x5v2+x6v3 = 0 means x1u1+x2u2+x3u3 = -x4v1-x5v2-x6v3, which is in both.
\exatend (\emph{challenge}) Suppose $U$ and $V$ are each $3$-dimensional subspaces of $\R^5$. Show that there is a nonzero vector that is in both $U$ and $V$.

% axler
% tricky! answer: false. V = span(u3, u1+u2) is a counterexample.
\exatend (\emph{challenge}) Prove or give a counterexample: If $\vec{u}_1, \vec{u}_2, \vec{u}_3$ are linearly independent vectors in $\R^3$, $V$ is a subspace of $\R^3$ that does not contain $\vec{u}_1$ or $\vec{u}_2$, then $V = \spn ( \vec{u}_3 )$.

\chapter{Matrices and Transformations}
\label{chap:matrices-and-transformations}


Chapter 1 was mostly about one system of equations at a time, or one linear combination at a time. But there are lots of physical phenomena in the world that can be described as a \emph{composition} of linear combinations: take a linear combination of the columns of one matrix $B$, and then use the entries from the result as weights in a linear combination of another matrix $A$. Using the notation in chapter 1 to describe this would be a big ugly mess.

We'll see in 2.1 and 2.2 how we can write this composition simply as $A(B\vec{x})$ and then $(AB)\vec{x}$. In 2.2, we'll also look at other important matrix operations like matrix addition, scaling, and \emph{transposes}. Section 2.3 is about the \emph{inverse} of $A$, which is sort of like ``matrix division.'' These operations form the basics of \emph{matrix algebra}--we can manipulate equations where the objects being added, multiplied, etc, are matrices. These simple-looking equations give us powerful techniques to understand very complicated phenomena across the sciences.

\columnbreak
\section{Matrix-Vector Products and Transformations}
\label{sec:matrix-vector-products}

\exercise (\emph{Warm-up}) \label{exercise:2.1wu1} Let $f(x) = 5x+1$.
\begin{enumerate}[label=(\alph*)]
\item Find $f(4\cdot 2)$ and $4\cdot f(2)$. Are they equal?
\begin{notes}
\vspace{5em}
\end{notes}

\item Find $f(3 - 2)$ and $f(3) + f(-2)$. Are they equal?
\end{enumerate}

\begin{notes}
\vspace{6em}
\end{notes}

\exercise (\emph{Warm-up}) Let $g(x) = 5x$. \label{exercise:2.1wu2}
\begin{enumerate}[label=(\alph*)]
\item Find $g(4\cdot 2)$ and $4\cdot g(2)$. Are they equal?
\begin{notes}
\vspace{5em}
\end{notes}

\item Find $g(3 - 2)$ and $g(3) + g(-2)$. Are they equal?
\end{enumerate}

\begin{notes}
\vfill
\end{notes}

We'll see later that $g$ in exercise \ref{exercise:2.1wu2}  is an example of a \emph{linear transformation} but $f$ in exercise \ref{exercise:2.1wu1} is not.

\begin{notes}
\columnbreak
\end{notes}

\subsection{Matrix-Vector Product}\label{subsec:matrix-vector-product}

\begin{definition} \textbf{Matrix-vector product.}

Suppose $A$ is an $r \times c$ matrix with columns $\vec{a}_1,\ldots,\vec{a}_c$ and $\vec{x}$ is a vector in $\R^c$. Then we define the product $A \vec{x}$ to be the linear combination of the columns of $A$ using the corresponding entries in $\vec{x}$ as coefficients:

\[ A \vec{x}=
\begin{bmatrix}
\vec{a}_1 & \vec{a}_2 & \cdots & \vec{a}_c
\end{bmatrix}
\begin{bmatrix}
x_1 \\ x_2 \\ \vdots \\ x_c
\end{bmatrix} = x_1 \vec{a}_1+x_2\vec{a}_2+ \cdots + x_c\vec{a}_c
\]
\end{definition}

\example Find each product:

\begin{horiparts}{2}
\item $\begin{bmatrix}
1 & 2 & 3 \\
4 & 5 & 6
\end{bmatrix} \bbbvec{-3}{1}{\frac{1}{3}}$
\item $
\begin{bmatrix}
\sqrt{e} & 5 \\
\frac{113}{17} & 6 \\
\sin \frac{\pi}{12} & 7
\end{bmatrix}\bbvec{0}{2}$
\end{horiparts}

\begin{notes}
\vfill
\end{notes}

\begin{book}
\textbf{Solution.} (a) $\begin{bmatrix}
1 & 2 & 3 \\
4 & 5 & 6
\end{bmatrix} \bbbvec{-3}{1}{\frac{1}{3}} = -3 \bbvec{1}{4}+1 \bbvec{2}{5}+\frac{1}{3}\bbvec{3}{6} = \bbvec{0}{-5}$

(b) $
\begin{bmatrix}
\sqrt{e} & 5 \\
\frac{113}{17} & 6 \\
\sin \frac{\pi}{12} & 7
\end{bmatrix}\bbvec{0}{2} = 0 \bbbvec{\sqrt{e}}{\frac{113}{7}}{\sin \frac{\pi}{12}} +2 \bbbvec{5}{6}{7} = \bbbvec{10}{12}{14}$

 \blacksquare

\columnbreak
\end{book}

\example Write the system as a vector equation and as a matrix equation in the form $A \vec{x}=\vec{y}$:
\quad $\begin{aligned}[t]
\pi x_1 + 2x_2 - x_3 &= 4 \\
-6x_2+3x_3 &= 5
\end{aligned}$

\begin{notes}
\vfill
\end{notes}

\begin{book}
\textbf{Solution.} Both equations are simultaneously true if and only if the following vectors are equal:
\[\bbvec{\pi x_1 + 2x_2-x_3}{-6x_2+3x_3} = \bbvec{4}{5} \]
The left can be written as a linear combination using the variables as weights:
\[x_1\bbvec{\pi}{0} +x_2 \bbvec{2}{-6} + x_3 \bbvec{-1}{3}  = \bbvec{4}{5} \]
The $A\vec{x}$ notation gives a compact way to write this:
\[\begin{bmatrix}
\pi & 2 & -1 \\
0 & -6 & 3
\end{bmatrix}\bbbvec{x_1}{x_2}{x_3} = \bbvec{4}{5}\]
\blacksquare
\end{book}

In general, if $A$ is an $r \times c$ matrix with columns $\vec{a}_1, \ldots, \vec{a}_c$, then the following are different ways of saying the same thing:
\begin{itemize}
\item The vector equation $ x_1 \vec{a}_1 + \cdots + x_c\vec{a}_c = \vec{y}$
\item The matrix equation $A\vec{x}=\vec{y}$
\item The system of equations with augmented matrix $\begin{bmatrix}[ccc|c]
\vec{a}_{1} & \cdots & \vec{a}_c & \vec{y}
\end{bmatrix}$
\end{itemize}

\begin{notes}
\columnbreak
\end{notes}
Much of what we learned in chapter 1 can be re-stated in terms of $A\vec{x}$. Here are some definitions and theorems that have been amended, with additions in \textcolor{red}{red}:

\begin{definition}
If $A =
\begin{bmatrix}
\vec{a}_1 & \cdots & \vec{a}_c
\end{bmatrix}
$, a linear system of equations that can be written in the form 
\begin{equation}\label{eq:2.1-1}
x_1 \vec{a}_1 + \cdots + x_c\vec{a}_c = \vec{0} \quad \textcolor{red}{\text{or} \quad A\vec{x}=\vec{0}} 
\end{equation}
is called \textbf{homogeneous}. 
The solution set of equation \ref{eq:2.1-1} is called the \textbf{null space} of $A$ and is written 
$\nul A$:
\[\nul A = \left\{ (x_1,\ldots,x_c) : x_1\vec{a}_1+\cdots+x_c\vec{a}_c=\vec{0} \right\} \textcolor{red}{\hspace{.5ex} = \left\{ \vec{x} : A\vec{x}=\vec{0} \right\}} \]

The \textbf{column space} of$A$, written $\col A$, is the span of the columns of $A$:
\[\col A = \left\{ x_1\vec{a}_1 + \cdots + x_c\vec{a}_c : x_1,\ldots,x_n \text { in } \R \right\}  
\textcolor{red}{\hspace{.5ex} = \left\{ A\vec{x} : \vec{x} \text{ in } \R^c \right\}}\]
\end{definition}

\begin{theorem}[The existence theorem, v3.]
\label{theorem:1.7existence-v3}

Suppose $A =
\begin{bmatrix}
\vec{a}_1 & \cdots & \vec{a}_c
\end{bmatrix}
$ is an $r \times c$ matrix. The following are equivalent:
\begin{enumerate}[label=(\roman*)]
\item $A$ has a pivot position in every row.
\item Every system of equations that has $A$ as its coefficient matrix is consistent.
\item $\col A = \R^r$.
\item $\dim(\nul A) = c-r$.
\item \textcolor{red}{For every $\vec{y}$ in $\R^{r}$, the equation $A\vec{x}=\vec{y}$ has a solution.}
\end{enumerate}
\end{theorem}

\begin{theorem}[The uniqueness theorem, v3.]
\label{theorem:1.7uniqueness-v3}

Suppose $A =
\begin{bmatrix}
\vec{a}_1 & \cdots & \vec{a}_c
\end{bmatrix}
$ is an $r \times c$ matrix. The following are equivalent:
\begin{enumerate}[label=(\roman*)]
\item $A$ has a pivot position in every column.
\item $\nul A = \left\{ \vec{0} \right\}$.
\item The columns $\vec{a}_1, \ldots, \vec{a}_c$ are a linearly independent list of vectors.
\item $\dim(\col A) = c$.
\item \textcolor{red}{If $A\vec{x}=\vec{0}$, then $\vec{x}=\vec{0}$.}
\end{enumerate}
\end{theorem}

\begin{notes}
\columnbreak
\end{notes}

\subsection{Linear Transformations}\label{subsec:linear-transformations}

Let's take a short detour to go over some newish terminology.

\begin{definition}
A \textbf{transformation} (sometimes called a \emph{map} or simply \emph{function}) $T$ from $\R^c$ to $\R^r$ takes any input vector $\vec{x}$ in $\R^c$ and gives a unique output vector $T(\vec{x})$ in $\R^r$.

The set of all inputs, $\R^c$, is called the \textbf{domain} of $T$, and $\R^r$, where the outputs live, is called the \textbf{codomain}.  The set of all outputs $\left\{T(\vec{x}) : \vec{x} \textrm{ is in } \R^r\right\}$ is called the \textbf{range} of $T$. So, the range is a \emph{subset} of the codomain.

The notation $T:\R^c \rightarrow \R^r$ specifies the domain and codomain.
\end{definition}

\begin{book}
\columnbreak
\end{book}
\example If $f:\R^2 \rightarrow \R$ has formula $f(x_1,x_2) = x_1^2 + x_2^2$, find the domain, codomain, and range of $f$.

\begin{book}
\textbf{Solution.} As indicated by the notation $f:\R^2 \rightarrow \R$, the domain is $\R^2$ and the codomain is $\R$. Since, for any $(x_1,x_2)$ in $\R^2$, the output $x_1^2 + x_2^2$ is non-negative, the range is $[0,\inf)$.

\begin{tikzpicture}[scale=0.7]
\begin{scope}[shift={(-6,0)}]
\fill[yellow!80,opacity=0.5] (-4,-4) rectangle (4,4);
\draw[step=1cm,gray,very thin] (-4,-4) grid (4,4);
\draw[thick] (-4,0) -- (4,0);
\draw[thick] (0,-4) -- (0,4);
\foreach \x in {-4,-3,-2,-1,1,2,3,4}
\draw[thick] (\x,0.1) -- (\x,-0.1) node[below,font=\small] {\x};
\foreach \y in {-4,-3,-2,-1,1,2,3,4}
\draw[thick] (0.1,\y) -- (-0.1,\y) node[left,font=\small] {\y};
\filldraw[black] (0,0) circle (2pt);
\node[below left,font=\small] at (0,0) {$\vec{0}$};
\filldraw[black] (-2,-1) circle (2pt);
\node[below] at (-2,-1) {$\vec{x}$};
\filldraw[black] (1,1) circle (2pt);
\node[below left] at (1,1) {$\vec{u}$};
\node at (3,3) {\large$\R^2$};
\end{scope}

\draw[->,thick,bend left=20] (-0.5,0) to node[above] {$f$} (2.5,0);
    
\begin{scope}[shift={(4,-1)}]
\fill[yellow!80,opacity=0.5] (-0.3,0) rectangle (0.3,6);
\draw[thick] (0,-3.5) -- (0,6);
\foreach \y in {-3,-2,-1,0,1,2,3,4,5} \draw[thick] (-0.1,\y) -- (0.1,\y);
\filldraw[black] (0,0) circle (2pt);
\node[right] at (0.2,0) {$0 = f(\vec{0})$};
\filldraw[black] (0,5) circle (2pt);
\node[right] at (0.2,5) {$5 = f(\vec{x})$};
\filldraw[black] (0,2) circle (2pt);
\node[right] at (0.2,2) {$2 = f(\vec{u})$};
\node at (-.7,3.5) {\large$\R$};
\end{scope}
\end{tikzpicture}

\blacksquare
\end{book}

\begin{notes}
\begin{tikzpicture}[scale=0.7]
\begin{scope}[shift={(-6,0)}]
%\fill[yellow!40,opacity=0.5] (-4,-4) rectangle (4,4);
\draw[step=1cm,gray,very thin] (-4,-4) grid (4,4);
\draw (-4,0) -- (4,0);
\draw (0,-4) -- (0,4);
\foreach \x in {-4,-3,-2,-1,1,2,3,4}
\draw (\x,0.1) -- (\x,-0.1) node[below,font=\small] {\x};
\foreach \y in {-4,-3,-2,-1,1,2,3,4}
\draw (0.1,\y) -- (-0.1,\y) node[left,font=\small] {\y};
%\filldraw[black] (0,0) circle (2pt);
%\node[below left,font=\small] at (0,0) {$\vec{0}$};
%\filldraw[black] (-2,-1) circle (2pt);
%\node[below] at (-2,-1) {$\vec{x}$};
%\filldraw[black] (1,1) circle (2pt);
%\node[below left] at (1,1) {$\vec{u}$};
\node at (3,3) {\large$\R^2$};
\end{scope}

\draw[->,thick,bend left=20] (-0.5,0) to node[above] {$f$} (2.5,0);
    
\begin{scope}[shift={(4,-1)}]
%\fill[yellow!40,opacity=0.5] (-0.3,0) rectangle (0.3,6);
\draw (0,-3.5) -- (0,6);
\foreach \y in {-3,-2,-1,0,1,2,3,4,5} \draw (-0.1,\y) -- (0.1,\y);
%\filldraw[black] (0,0) circle (2pt);
\node[right] at (0.2,0) {$0$};
%\filldraw[black] (0,5) circle (2pt);
%\node[right] at (0.2,5) {$5 = f(\vec{x})$};
%\filldraw[black] (0,2) circle (2pt);
%\node[right] at (0.2,2) {$2 = f(\vec{u})$};
\node at (-.7,3.5) {\large$\R$};
\end{scope}
\end{tikzpicture}

\vfill \hspace{2pt}
\columnbreak
\end{notes}

\example Define $h(x) = \bbvec{3x}{2x}$. Find the domain, codomain, and range of $h$.

\begin{book}
\textbf{Solution.} Based on the formula $h(x) = \bbvec{3x}{2x}$, inputs are numbers in $\R$, the domain. The outputs are vectors with two entries, so the codomain is $\R^2$. The range is the set of vectors that can be written $\bbvec{3x}{2x},$ or $x\bbvec{3}{2}$, for some real number $x$. In other words, the range is $\spn \left( \bbvec{3}{2} \right)$.

\begin{tikzpicture}[scale=0.7]
\begin{scope}[shift={(-2,-1)}]
\fill[yellow!80,opacity=0.5] (-0.3,-3) rectangle (0.3,3);
\draw (0,-3) -- (0,3);
\foreach \y in {-3,-2,-1,0,1,2,3} \draw (-0.1,\y) -- (0.1,\y);
\filldraw[black] (0,0) circle (2pt);
\node[right] at (0.2,0) {$0$};
\filldraw[black] (0,1) circle (2pt);
\node[right] at (0.2,1) {$1$};
\filldraw[black] (0,-1) circle (2pt);
\node[right] at (0.2,-1) {$-1$};
\node at (.7,3) {\large$\R$};
\end{scope}

\draw[->,thick,bend left=20] (-0.5,0) to node[above] {$h$} (2.5,0);

\begin{scope}[shift={(8,0)}]
\draw[yellow!80, line width=8pt, opacity=0.5] (-4,-2.6667) -- (4,2.6667);
\draw[black] (-4,-2.6667) -- (4,2.6667);
\draw[step=1cm,gray,very thin] (-4,-4) grid (4,4);
\draw (-4,0) -- (4,0);
\draw (0,-4) -- (0,4);
\foreach \x in {-4,-3,-2,-1,1,2,3,4}
\draw (\x,0.1) -- (\x,-0.1) node[below,font=\small] {\x};
\foreach \y in {-4,-3,-2,-1,1,2,3,4}
\draw (0.1,\y) -- (-0.1,\y) node[left,font=\small] {\y};
\filldraw[black] (0,0) circle (2pt);
\node[below right,font=\small] at (0,0) {$\vec{0}$};
\filldraw[black] (3,2) circle (2pt);
\node[below right] at (3,2) {$h(1)$};
\filldraw[black] (-3,-2) circle (2pt);
\node[below right] at (-3,-2) {$h(-1)$};
\node at (-3,3) {\large$\R^2$};
\end{scope}    
\end{tikzpicture}

\blacksquare
\end{book}

\begin{notes}
\hspace{1em}\begin{tikzpicture}[scale=0.7]
\begin{scope}[shift={(-2,-1)}]
%\fill[yellow!80,opacity=0.5] (-0.3,-3) rectangle (0.3,3);
\draw (0,-3) -- (0,3);
\foreach \y in {-3,-2,-1,0,1,2,3} \draw[thick] (-0.1,\y) -- (0.1,\y);
%\filldraw[black] (0,0) circle (2pt);
\node[right] at (0.2,0) {$0$};
%\filldraw[black] (0,1) circle (2pt);
\node[right] at (0.2,1) {$1$};
%\filldraw[black] (0,-1) circle (2pt);
\node[right] at (0.2,-1) {$-1$};
\node at (.7,3) {\large$\R$};
\end{scope}

\draw[->,thick,bend left=20] (-0.5,0) to node[above] {$h$} (2.5,0);

\begin{scope}[shift={(8,0)}]
%\draw[yellow!80, line width=8pt, opacity=0.5] (-4,-2.6667) -- (4,2.6667);
%\draw[thick, black] (-4,-2.6667) -- (4,2.6667);
\draw[step=1cm,gray,very thin] (-4,-4) grid (4,4);
\draw (-4,0) -- (4,0);
\draw (0,-4) -- (0,4);
\foreach \x in {-4,-3,-2,-1,1,2,3,4}
\draw (\x,0.1) -- (\x,-0.1) node[below,font=\small] {\x};
\foreach \y in {-4,-3,-2,-1,1,2,3,4}
\draw (0.1,\y) -- (-0.1,\y) node[left,font=\small] {\y};
%\filldraw[black] (0,0) circle (2pt);
\node[below right,font=\small] at (0,0) {$\vec{0}$};
%\filldraw[black] (3,2) circle (2pt);
%\node[below right] at (3,2) {$h(1)$};
%\filldraw[black] (-3,-2) circle (2pt);
%\node[below right] at (-3,-2) {$h(-1)$};
\node at (-3,3) {\large$\R^2$};
\end{scope}    
\end{tikzpicture}
\vfill
\end{notes}

\example Let $A =\begin{bmatrix}
2 & -\pi & 3 & \sqrt{2} \\
4 & 0 & 1 & -8
\end{bmatrix}$ and $T(\vec{x}) = A \vec{x}$. Find $T(0,1,-2,0)$. Describe the domain, codomain, range of $T$.

\begin{notes}
\vfill \vfill \hspace{2pt}
\columnbreak
\end{notes}

\begin{book}
\textbf{Solution.} \[T(0,1,-2,0) = \begin{bmatrix}
2 & -\pi & 3 & \sqrt{2} \\
4 & 0 & 1 & -8
\end{bmatrix}\begin{bmatrix}
0 \\ 1 \\ -2 \\ 0
\end{bmatrix} = 0 \bbvec{2}{4}+1 \bbvec{-\pi}{0} -2 \bbvec{3}{1} + 0 \bbvec{\sqrt{2}}{-8}= \bbvec{-\pi-6}{-2}\]
The domain is $\R^4$ because $A\vec{x}$ is only defined when the number of columns of $A$ equals the number of entries of $\vec{x}$. The codomain is $\R^2$ because the output, $A\vec{x}$, is a linear combination of the columns of $A$ which are $\R^2$ vectors. The range is the set of actual outputs that you get:
\[\text{range}(T) = \left\{ A\vec{x}:\vec{x} \text{ is in } \R^4 \right\}\]
This is the set of all linear combinations of the columns of $A$, also known as $\col A$. Row-reduction shows that there is a pivot in each row of $A$ which means $\text{range}(T) = \col A = \R^2$. \blacksquare
\end{book}

\begin{theorem}[Linearity of $\mathbf{A\accentset{\rightharpoonup}x}$.]
\label{theorem:linearity-of-Ax}

Suppose $A$ is an $r \times c$ matrix, $\vec{x}, \vec{z}$ are in $\R^c$, and $k$ is a scalar. Then:
\begin{enumerate}[label=(\roman*)]
\item $A(\vec{x}+\vec{z}) = A\vec{x}+A\vec{z}$
\item $A(k\vec{x}) = k(A\vec{x})$
\end{enumerate}
\end{theorem}

\textbf{Proof.} (i) Exercise \ref{exercise:2.1-linearity-pf1} asks you to fill in the blanks:
\begin{align*}
A ( \vec{x}+\vec{z} ) &= A \bbbvec{x_1+z_1}{\vdots}{x_c+z_c}  \\
&= \rule{3em}{0.8pt}\vec{a}_1 + \cdots + \rule{3em}{0.8pt}\vec{a}_c \quad \text{(def'n of matrix-vector product)} \\
&= \rule{2em}{0.8pt} + \rule{2em}{0.8pt} + \cdots + \rule{2em}{0.8pt}+\rule{2em}{0.8pt} \quad \text{(theorem \ref{properties-of-vector-operations} part (v))} \\
&= x_1\vec{a}_1 + \cdots + x_c\vec{a}_c + z_1\vec{a}_1 + \cdots + z_c\vec{a}_c \quad \text{(theorem \rule{6em}{0.8pt})} \\
&= \rule{3em}{0.8pt} + \rule{3em}{0.8pt} \quad \text{(def'n of matrix-vector product)}
\end{align*}

(ii) Exercise \ref{exercise:2.1-linearity-pf2} asks you to fill in the blanks:
\begin{align*}
A ( k\vec{x} ) &= A \bbbvec{kx_1}{\vdots}{kx_c}  \\
&= (\rule{3em}{0.8pt})\vec{a}_1 + \cdots + (\rule{3em}{0.8pt})\vec{a}_c \quad \quad \text{(def'n of matrix-vector product)} \\
&= k (  \rule{3em}{0.8pt} ) + \cdots +  k (  \rule{3em}{0.8pt} ) \quad \quad \text{(theorem \ref{properties-of-vector-operations} part (vii))} \\
&= k (\rule{10em}{0.8pt}) \quad \quad \text{(theorem \ref{properties-of-vector-operations} part (vi))} \\
&= \rule{4em}{0.8pt} \quad \quad \text{(def'n of matrix-vector product)} 
\end{align*}
\blacksquare

\begin{definition}
A transformation $T$ is called \textbf{linear} if for all real numbers $k$, and all vectors $\vec{x}$ and $\vec{z}$ in the domain,
\begin{enumerate}[label=(\roman*)]
\item $T ( \vec{x}+\vec{z} ) = T ( \vec{x} ) + T ( \vec{z} )$
\item $T ( k \vec{x} )= k T ( \vec{x} )$
\end{enumerate}
\end{definition}

If a transformation is of the form $T ( \vec{x} )= A\vec{x}$, theorem \ref{theorem:linearity-of-Ax} shows it's linear. We'll soon see that \emph{all} linear transformations $\R^c \rightarrow \R^r$ are of this form. 

\begin{notes}
\columnbreak
\end{notes}

\example Recall from section \ref{sec:subspaces-of-Rc} that, given any two vectors $\vec{u},\vec{v}$ in $\R^c$, the set of points $\left\{ \vec{u} + t\vec{v} : t \text{ in } \R \right\}$ is a line through the point $\vec{u}$ parallel to $\vec{v}$ (as an arrow).  Explain why linear transformations \emph{map lines onto lines}.
%Show that if $T: \R^c \rightarrow \R^r$, then $\left\{ T ( u + t \vec{v} ) : t \text{ in } \R \right\}$

\begin{book}
\textbf{Solution.} Suppose $L$ is the set of points on the through $\vec{u}$ parallel to $\vec{v}$:
\[L= \left\{ \vec{u}+t\vec{v}:t \text{ is in } \R \right\}\]
Then for any point $\vec{u}+t\vec{v}$ on this line the definition of linearity means that:
\[T(\vec{u}+t\vec{v}) = T(\vec{u}) + T(t\vec{v}) = T(\vec{u}) + tT(\vec{v})\]
This is a point on the line through $T(\vec{u})$ parallel to $T(\vec{v})$. Note that if $T(\vec{v}) = \vec{0}$, this ``line'' in the range is simply the point $T(\vec{u})$. \blacksquare
\end{book}

\begin{tikzpicture}
\begin{axis}[
x=27pt,y=27pt,
axis lines=middle,
axis line style={-},
xtick=\empty,
ytick=\empty,
ymajorgrids=false,
%minor y tick num=3,
%yminorgrids=true,
xmajorgrids=false,
%minor x tick num=3,
%xminorgrids=true,
%xlabel=$x_1$,
%ylabel=$x_2$,
xmin=-1, xmax=6.5,
ymin=-2, ymax=3
%xtick={-1,0,...,6},
%ytick={-2,-1,...,3},
]
\filldraw[black] (axis cs:2,1) circle (1pt);
\node[right] at (axis cs:2,1) {$\vec{u} + 0\vec{v}$};
\filldraw[black] (axis cs:4.5,-1.5) circle (1pt);
\node[right] at (axis cs:4.5,-1.5) {$\vec{u} + 2.5\vec{v}$};
\filldraw[black] (axis cs:1,2) circle (1pt);
\node[right] at (axis cs:1,2) {$\vec{u} + (-1)\vec{v}$};
\filldraw[black] (axis cs:3,0) circle (1pt);
\node[above right] at (axis cs:3,0) {$\vec{u} + 1\vec{v}$};
%\draw[-Stealth, very thick] (axis cs:0,0) -- (axis cs:3,1)
%node[right] {$\vec{u}$};
\draw[-Stealth, thick] (axis cs:0,0) -- (axis cs:1,-1)
node[right] {$\vec{v}$};
\draw[thick] (axis cs:0,3) -- (axis cs:6.5,-3.5);
\end{axis}
\end{tikzpicture}

\begin{notes}
\vfill \vfill
\end{notes}

\example Recall $I =
\begin{bmatrix}
\vec{e}_1 & \vec{e}_2 & \vec{e}_3
\end{bmatrix} =
\begin{bmatrix}
1 & 0 &0 \\
0 & 1 & 0 \\
0 & 0 & 1
\end{bmatrix}$ is called the $3 \times 3$ identity matrix. If $T ( \vec{x} )=I\vec{x}$, find $T(9,8,7)$.

\begin{notes}
\vfill \hspace{2pt}
\columnbreak
\end{notes}

\begin{book}
\textbf{Solution.} $T(9,8,7) = \begin{bmatrix}
1 & 0 & 0 \\
0 & 1 & 0 \\
0 & 0 & 1
\end{bmatrix}\bbbvec{9}{8}{7} = 9 \bbbvec{1}{0}{0} + 8 \bbbvec{0}{1}{0} + 7 \bbbvec{0}{0}{1} = \bbbvec{9}{8}{7}$ \blacksquare
\end{book}

\example The set of points $S = \left\{ (x_1,x_2) : 0\le x_1 \le 1, 0\le x_2 \le 1  \right\}$ is called the unit square. Define $T ( \vec{x} ) = A \vec{x}$ for $A =\begin{bmatrix}
1 & 0 \\
-2 & 1
\end{bmatrix}$. Describe the set $\left\{ T ( \vec{x} ) : \vec{x} \text{ in } S\right\}$. Note: this is called a \emph{shear transformation}.

\begin{tikzpicture}[scale=3]
\draw[-] (-0.2,0) -- (1.3,0);
\draw[-] (0,-0.2) -- (0,1.3);
\fill[blue!20] (0,0) rectangle (1,1);
\draw[thick, blue] (0,0) rectangle (1,1);
\node at (0.8, 0.8) {$S$};
\draw[black,thick] (1,0.06) -- (1,-0.06) node[below] {$1$};
\draw[black,thick] (0.06,1) -- (-0.06,1) node[left] {$1$};
\end{tikzpicture}

\begin{book}
\textbf{Solution.} Note that any vector in the unit square can be written
\[\bbvec{x_1}{x_2} = x_1 \bbvec{1}{0} + x_2 \bbvec{0}{1} = x_1\vec{e}_1+x_2\vec{e}_2 \text{ for }0\le x_1 \le 1, 0\le x_2 \le 1\] Since $T$ is linear:
\begin{align*}
T ( x_1\vec{e}_1+x_2\vec{e}_2 ) &= T (x_1\vec{e}_1) + T(x_2\vec{e}_2  ) \\
&= x_1T (\vec{e}_1) + x_2T(\vec{e}_2)  \text{ for }0\le x_1 \le 1, 0\le x_2 \le 1
\end{align*}
So, we just need to find $T ( \vec{e}_1 )$ and $T ( \vec{e}_2 )$, and see what linear combinations of those look like:
\[T ( \vec{e}_1 ) = \begin{bmatrix}
1 & 0 \\
-2 & 1
\end{bmatrix} \bbvec{1}{0} = \bbvec{1}{-2}\text{ and }T ( \vec{e}_2 ) = \begin{bmatrix}
1 & 0 \\
-2 & 1
\end{bmatrix} \bbvec{0}{1} = \bbvec{0}{1}\]

So, the set $\left\{ T ( \vec{x} ) : \vec{x} \text{ in } S\right\}$ includes these points and other points like:
\begin{align*}
T(0,0) &= 0 \bbvec{1}{-2} + 0 \bbvec{0}{1} = \bbvec{0}{0}, \quad 
T(1,1) = 1 \bbvec{1}{-2} + 1 \bbvec{0}{1} = \bbvec{1}{-1} \quad \text{ and } \\
T(0.5,.3) &= 0.5\bbvec{1}{-2} + 0.3\bbvec{0}{1} = \bbvec{0.5}{-0.7}
\end{align*}
All such linear combinations give you a parallelogram:

\begin{tikzpicture}[scale=2]
\draw[-] (-0.2,0) -- (1.3,0);
\draw[-] (0,-2.2) -- (0,1.3);
\fill[blue!20] (0,0) -- (0,1) -- (1,-1) -- (1,-2) -- cycle;
\draw[thick, blue] (0,0) -- (0,1) -- (1,-1) -- (1,-2) -- cycle;
\node at (0.5, -0.5) {$T(S)$};
\draw[black,thick] (1,0.06) -- (1,-0.06) node[below] {$1$};
\draw[black,thick] (0.06,1) -- (-0.06,1) node[left] {$1$};
\draw[black,thick] (0.06,-1) -- (-0.06,-1) node[left] {$-1$};
\draw[black,thick] (0.06,-2) -- (-0.06,-2) node[left] {$-2$};
\end{tikzpicture}

\blacksquare
\end{book}

\begin{notes}
\vfill
\end{notes}

The next theorem generalizes the idea from this example.

\begin{notes}
\columnbreak
\end{notes}
\begin{theorem}[Standard matrix of a transformation.]
\label{theorem:standard-matrix}

If $T:\R^c\rightarrow \R^r$ is a linear transformation, then:
\[ T ( \vec{x} ) =
\begin{bmatrix}
T ( \vec{e}_1 ) & \cdots & T ( \vec{e}_c )
\end{bmatrix}\vec{x} = x_1T ( \vec{e}_1 ) + \cdots + x_cT(\vec{e}_c)\]

The matrix $A = \begin{bmatrix}
T ( \vec{e}_1 ) & \cdots & T ( \vec{e}_c )
\end{bmatrix}$ is called the \textbf{standard matrix} of $T$.
\end{theorem}

\begin{book}
\textbf{Proof.} Let $T:\R^c \rightarrow \R^r$ be a linear transformation, and let $\vec{x}$ be in $\R^c$. Then we can write:
\begin{align*}
 T ( \vec{x} ) &= T \left( \bbbvec{x_1}{\vdots}{x_c} \right) \\
&= T \left( x_1 \bbbvec{1}{\vdots}{0} + \cdots + x_c \bbbvec{0}{\vdots}{1} \right)\\
&= T  ( x_1 \vec{e}_c + \cdots + x_c \vec{e}_c )\\
&= T  ( x_1 \vec{e}_c) + \cdots + T(x_c \vec{e}_c )  &(T \text{ is linear})\\
&= x_1T (\vec{e}_c) + \cdots + x_cT(\vec{e}_c ) & (T \text{ is linear}) \\
&= \begin{bmatrix}
T ( \vec{e}_1 ) & \cdots & T ( \vec{e}_c ) 
\end{bmatrix}\vec{x}& (\text{def'n of matrix-vector product})
\end{align*}
\blacksquare
\end{book}

\begin{notes}
\vfill \vfill
\end{notes}

\example $T:\R^2\rightarrow\R^2$ reflects points across the $x_2$-axis. Find the standard matrix of $T$ and check that it works with a couple points.

\bbtwocol{0.45}{0.55}{
\begin{tikzpicture}[scale=1.2]
\vectorgrid[xmin=-2.2,xmax=2.2,ymin=-1.5,ymax=1.5,step=1,labint=1,xu=-1,yu=0,xv=0,yv=1]
\node[above] at (-0.8,0){$T(\vec{e}_1)$};
\node[right] at (0,0.8){$T(\vec{e}_2)$};
\draw[thick,-Stealth] (0,0) -- (2,-1) node[below left]{$\vec{x}$};
\draw[thick,-Stealth] (0,0) -- (-2,-1) node[below right]{$T(\vec{x})$};
\end{tikzpicture}}{\textbf{Solution.} The graph shows $T$ reflects the point $(0,1)$, aka $\vec{e}_1$, across the $x_2$-axis to $(-1,0)$. And $T$ doesn't move $\vec{e}_2$, so $T(\vec{e}_2) = \vec{e}_2$. Thus the standard matrix of $T$ is:
\[\begin{bmatrix}
T(\vec{e}_1) & T(\vec{e}_2)
\end{bmatrix} = \begin{bmatrix}
-1 & 0 \\
0 & 1
\end{bmatrix}\]
Checking that it works with $\vec{x}=(2,-1)$:
\[T(\vec{x}) = \begin{bmatrix}
-1 & 0 \\
0 & 1
\end{bmatrix}\bbvec{2}{-1} = \bbvec{-2}{-1} \quad \quad \blacksquare \]
}

\begin{notes}
\vfill \hspace{2pt}
\columnbreak
\end{notes}

\example $R_{\phi}:\R^2\rightarrow\R^2$ rotates points an angle of $\phi$ counterclockwise around the origin. 
\begin{enumerate}[label=(\alph*)]
\item Find the standard matrix for $R_{\phi}$.

\begin{notes}
\begin{tikzpicture}[scale=3]
\vectorgrid[xmin=-1.2,xmax=1.2,ymin=-.5,ymax=1.2,step=1,labint=1,xu=.866,yu=0.5,xv=-.5,yv=0.866]
\end{tikzpicture}

\vfill
\end{notes}

\item Check the matrix for $R_{\frac{2\pi}{3}}$ works for $\vec{x}=(1, \sqrt{3})$.
\end{enumerate}

\begin{book}
\bbtwocol{0.55}{0.45}{
\textbf{Solution.} (a) As the graph shows, $T(\vec{e}_1)$ is still one unit from the origin, but at an angle $\phi$. Basic trigonometry says its coordinates are:
\[T(\vec{e}_1) = \bbvec{\cos \phi}{\sin \phi}\]
}{
\begin{tikzpicture}[scale=2]
\vectorgrid[xmin=-1.2,xmax=1.2,ymin=-.5,ymax=1.2,step=1,labint=1,xu=.866,yu=0.5,xv=-.5,yv=0.866]
\coordinate (Z) at (0,0);
\coordinate (X) at (1,0);
\coordinate (A) at (0.866,0.5);
\coordinate (B) at (-0.5,0.866);
\tkzMarkAngle[size=0.6](X,Z,A)
\node[right] at (0.6,0.2) {$\phi$};
\tkzMarkAngle[size=0.45](X,Z,B)
\node[above] at (0.25,0.35) {$\phi+\frac{\pi}{2}$};
\node[above] at (0.866,0.5) {$T(\vec{e}_1)$};
\node[left] at (-0.5,0.866) {$T(\vec{e}_2)$};
\end{tikzpicture}}

Likewise $T(\vec{e}_2)$ is at an angle $\phi+\frac{\pi}{2}$. Here we'll use some identities for sine and cosine of $\phi+\frac{\pi}{2}$:
\[T(\vec{e}_2) = \bbvec{\cos \left( \phi + \frac{\pi}{2} \right)}{\sin \left( \phi+\frac{\pi}{2} \right)} = \bbvec{-\sin \phi}{\cos \phi}\]
Putting these together, the standard matrix of $R_{\phi}$ is: 
\[A = \begin{bmatrix}
T(\vec{e}_1) & T(\vec{e}_2)
\end{bmatrix} =
\begin{bmatrix}
\cos \phi & -\sin \phi \\
\sin \phi & \cos \phi
\end{bmatrix}\]
\end{book}



\begin{notes}
\begin{tikzpicture}[scale=1.6]
\vectorgrid[xmin=-2.2,xmax=2.2,ymin=-1.5,ymax=2.1,step=1,labint=10,xu=-0.5,yu=0.866,xv=-0.866,yv=-0.5]
\end{tikzpicture}

\vfill \hspace{2pt}
\end{notes}

\begin{book}
\bbtwocol{0.58}{0.38}{(b) First let's find $R_{\frac{2\pi}{3}}(\vec{x})$ without the matrix so we can check that the matrix works. The point $\vec{x}$ is at a distance $\sqrt{1^2+\sqrt{3}^2} = 2$ from the origin, at an angle $\tan^{-1}\frac{\sqrt{3}}{1} = \frac{\pi}{3}$. Rotating this point by $\frac{2\pi}{3}$ take it to an angle $\frac{\pi}{3} + \frac{2\pi}{3} = \pi$, still at a distance 2 from the origin. This makes $T(\vec{x})=(-2,0)$.}{
\begin{tikzpicture}[scale=1.1]
\vectorgrid[xmin=-2.3,xmax=2.2,ymin=-1.5,ymax=2.1,step=1,labint=10,xu=-0.5,yu=0.866,xv=-0.866,yv=-0.5]
\draw[-Stealth,thick] (0,0) -- (1,1.73) node[right] {$\vec{x}$};
\draw[-Stealth,thick,blue] (0,0) -- (-2,0) node[above] {$T(\vec{x})$};
\node[blue,left] at (-.5,0.866) {$T(\vec{e}_1)$};
\node[blue,below] at (-0.866,-0.5) {$T(\vec{e}_2)$};
\end{tikzpicture}}

Now, the formula from part (a) gives 
\begin{align*}R_{\frac{2\pi}{3}}(\vec{x}) = \begin{bmatrix}
\cos \frac{2\pi}{3} & -\sin \frac{2\pi}{3} \\
\sin \frac{2\pi}{3} & \cos \frac{2\pi}{3}
\end{bmatrix}\vec{x} = \begin{bmatrix}
\frac{-1}{2} & \frac{-\sqrt{3}}{2} \\
\frac{\sqrt{3}}{2} & \frac{-1}{2}
\end{bmatrix}\bbvec{1}{\sqrt{3}} &= 1 \bbvec{\frac{-1}{2}}{\frac{\sqrt{3}}{2}} + \sqrt{3} \bbvec{\frac{-\sqrt{3}}{2}}{\frac{-1}{2} } \\
&= \bbvec{\frac{-1}{2}-\frac{3}{2}}{\frac{\sqrt{3}}{2} - \frac{\sqrt{3}}{2}} = \bbvec{-2}{0}
\end{align*}
as expected. Note the position of $\vec{x}$ relative to the black grid (defined by $\vec{e}_1$ and $\vec{e}_2$) is the same as the position of \textcolor{blue}{$T(\vec{x})$} relative to the blue grid (defined by \textcolor{blue}{$T(\vec{e}_1)$} and \textcolor{blue}{$T(\vec{e}_2)$}).
\blacksquare
\end{book}

\columnbreak

\exerciseheading

\textbf{\ref{subsec:matrix-vector-product} Matrix-Vector Product}

\exatend Multiply: $\begin{bmatrix}
7 & 3 \\
0 & 2 \\
-5 & 2
\end{bmatrix}\bbvec{2}{-1}$.

\exatend Multiply: $
\begin{bmatrix}
8 & 7 & 6 \\
5 & 4 & 3
\end{bmatrix}\bbbvec{1}{-2}{1}$.

\exatend Multiply: $\begin{bmatrix}
8 & -3 & 1 & 2 
\end{bmatrix}
\begin{bmatrix}
1 \\ 2 \\ 3 \\ 0.5
\end{bmatrix}$

\exatend Multiply (write as a single vector): $\begin{bmatrix}
a_{11} & a_{12} & a_{13} \\
a_{21} & a_{22} & a_{23}
\end{bmatrix}\bbbvec{x_1}{x_2}{x_3}$

\exatend Write as a system of linear equations (no need to solve):
$\begin{bmatrix}
8 & 2 & 0 \\
3 & \pi & 7 \\
0.5 & -1 & 3
\end{bmatrix}\bbbvec{x_1}{x_2}{x_3} = \bbbvec{3}{4}{-7}$.

\exatend Write as a vector equation (in the form $x_1\vec{a}_1+\cdots+x_c\vec{a}_c= \vec{y}$) and as a matrix equation (in the form $A\vec{x}=\vec{y}$). No need to solve:
$\begin{aligned}
  x_1 + x_{2} &=3 \\
  -x_{1} + 2x_{2} &= 0
\end{aligned}$


\exatend Suppose $A =
\begin{bmatrix}
a & b & c \\
x & y & z
\end{bmatrix}$, and $A\vec{p}$ exists.
\begin{enumerate}[label=(\alph*)]
\item Is $\vec{p}$ in $\R^2$ or $\R^{3}$?
\item Is $A\vec{p}$ in $\R^2$ or $\R^{3}$?
\end{enumerate}

%\exatend Let $\vec{x} = \bbbvec{2}{-7}{3}, \vec{y}=\bbbvec{1}{3}{5},$ and $\vec{z}=\bbbvec{-1}{36}{16}$. It can be shown that $-3 \vec{x} + 5\vec{y}-\vec{z}=0$. Use this fact to find a solution of:
%\[\begin{bmatrix}
%  2 & 1 \\
%  -7 & 3 \\
%  3 & 5 
%\end{bmatrix} \bbvec{x_1}{x_2} = \bbbvec{-1}{36}{16}\]
\exatend  Suppose $A$ is a $4 \times 3$ matrix. Show that there must be a vector $\vec{y}$ in $\R^4$ that makes $A \vec{x} = \vec{y}$ inconsistent.

\exatend Construct a $3 \times 3$ matrix $A$, with some non-zero entries, such that $\vec{x}=\bbbvec{-1}{1}{1}$ is a solution of $A\vec{x}=\vec{0}$.

Remember to justify all True/False answers. Maybe cite something in the text that shows it's true, or if it's false give a counterexample or write a corrected version of it.

\exatend (T/F?) If $A =
\begin{bmatrix}
\vec{a}_1 & \cdots & \vec{a}_c
\end{bmatrix}$ and $\vec{v}$ is in $\R^c$, then $A\vec{v}$ is in $\spn ( \vec{a}_1, \ldots, \vec{a}_c )$.

\exatend (T/F?) If $A\vec{x} = \vec{0}$ and $\vec{x} \ne \vec{0}$ then the columns of $A$ are linearly dependent.

\exatend (T/F?) If $A\vec{x} = \vec{0}$ and $\vec{x} = \vec{0}$ then the columns of $A$ are linearly independent.

\exatend (T/F?) The $i^\text{th}$ entry of $A\vec{x}$ is the dot product of row $i$ of $A$ and $\vec{x}$. Note: the dot product of $(u_1,u_2,\ldots,u_c)$ and $(v_1,v_2,\ldots,v_c)$ is defined to be $u_1v_1+u_2v_2+\cdots+u_cv_c$.

\textbf{\ref{subsec:linear-transformations} Linear Transformations}

\exatend (T/F?) The words ``range'' and ``codomain'' mean exactly the same thing.

\exatend (T/F?) The words ``transformation,'' ``map,'' and ``function'' all mean exactly the same thing.

\exatend (T/F?) Every transformation that can be written $T(\vec{x}) = A\vec{x}$ is a linear transformation.

\exatend (T/F?) If $T:\R^c\rightarrow\R^r$ is a linear transformation, then $T(a\vec{x}+b \vec{y}) = aT(\vec{x})+bT(\vec{y})$ for all $a,b$ in $\R$ and all $\vec{x},\vec{y}$ in $\R^c$.

\exatend Match each geometric description of a transformation with its standard matrix. 

$A =\begin{bmatrix}
  1 & 0 \\
  0 & -1
\end{bmatrix} \quad B =\begin{bmatrix}
  0 & 0 \\
  0 & 1
\end{bmatrix} \quad C =\begin{bmatrix}
  2 & 0 \\
  0 & 1
\end{bmatrix} \quad D =\begin{bmatrix}
  1 & -1 \\
  0 & 1
\end{bmatrix}$

$E =\begin{bmatrix}
  1 & 0 \\
  2 & 1
\end{bmatrix} \quad F =\begin{bmatrix}
  0 & 1 \\
  1 & 0
\end{bmatrix} \quad I =\begin{bmatrix}
  1 & 0 \\
  0 & 1
\end{bmatrix}$

Does nothing: \rule{2em}{0.8pt}

Horizontal stretch: \rule{2em}{0.8pt}

Reflection across the $x_1$-axis: \rule{2em}{0.8pt} 

Reflection across the line $x_1 = x_2$: \rule{2em}{0.8pt}

Projection onto the $x_2$-axis: \rule{2em}{0.8pt} 

Horizontal shear (points move horizontally depending on $x_2$): \rule{2em}{0.8pt} 

Vertical shear (points move vertically depending on $x_1$): \rule{2em}{0.8pt}

\columnbreak
\exatend Let $T:\R^2\rightarrow\R^2$ with $T ( \vec{e}_1 )$ and $T ( \vec{e}_2 )$ shown below.
\begin{enumerate}[label=(\alph*)]
\item Find $T(-2,-1)$ and $T(0,2)$.
\item Find $\vec{x}$ so that $T(\vec{x})=(3,0)$.
\item Find $\vec{x}$ so that $T(\vec{x}) = \vec{0}$. Explain why the solution is unique.
\item Is the range of $T$ a single point, a line, or all of $\R^2$?
\end{enumerate}
\begin{tikzpicture}[scale=0.6]
\vectorgrid[xmin=-4,xmax=4,ymin=-4,ymax=4,step=1,labint=2,xu=-1,yu=-1,xv=-1,yv=2]
\node[left] at (-1,-1) {$T(\vec{e}_1)$};
\node[left] at (-1,2) {$T(\vec{e}_2)$};
\end{tikzpicture}

\exatend Let $T:\R^3\rightarrow\R^2$ with $T ( \vec{e}_1 ), T ( \vec{e}_2 )$ and $T ( \vec{e}_3 )$ shown below.
\begin{enumerate}[label=(\alph*)]
\item Find $T(1,0,1)$ and $T(3,1,2)$.
\item Find a non-zero $\vec{x}$ so that $T(\vec{x}) = \vec{0}$.
\item If $T(\vec{x}) = A\vec{x}$, find a basis for $\nul A$.
\item Is the range of $T$ a single point, a line, or all of $\R^2$?
\end{enumerate}
\begin{tikzpicture}[scale=0.6]
\vectorgrid[xmin=-4,xmax=4,ymin=-4,ymax=4,step=1,labint=2,xu=-1,yu=-1,xv=-1,yv=2]
\node[left] at (-1,-1) {$T(\vec{e}_1)$};
\node[left] at (-1,2) {$T(\vec{e}_2)$};
\draw[very thick,blue,-stealth] (0,0) -- (3,0); 
\node[above] at (3,0) {$T(\vec{e}_3)$};
\end{tikzpicture}

\exatend Write your own example of a matrix $A$ that is the standard matrix for $T:\R^2\rightarrow\R^2$, with the property that the range of $T$ is a line in $\R^2$.  Is it possible to find a non-zero $\vec{x}$ so that $T(\vec{x}) = \vec{0}$?

\exatend Write your own example of a matrix $A$ that is the standard matrix for $T:\R^3\rightarrow\R^3$, with the property that the range of $T$ is all of $\R^3$.  Is it possible to find a non-zero $\vec{x}$ so that $T(\vec{x}) = \vec{0}$?

\exatend \label{exercise:2.1-linearity-pf1} Write a finished proof of theorem \ref{theorem:linearity-of-Ax} part (i) with the blanks filled in.

\exatend \label{exercise:2.1-linearity-pf2} Write a finished proof of theorem \ref{theorem:linearity-of-Ax} part (ii) with the blanks filled in.

\exatend Find the standard matrix for an $\R^2\rightarrow \R^2$ linear transformation that rotates points $90^{\circ}$ counterclockwise, and then stretches all points away from the origin by a factor of 2.

\exatend Find the standard matrix for an $\R^3\rightarrow \R^3$ linear transformation that projects points onto the nearest point on the $yz$-plane.

\exatend Find the standard matrix for an $\R^2\rightarrow \R^2$ linear transformation that reflects points across the line $x_1=x_2$.

\exatend Find the standard matrix for an $\R^3\rightarrow \R^3$ linear transformation that reflects points across the $xy$-plane and then rotates them counterclockwise by $45^{\circ}$ around the $x$-axis.

\exatend Find the standard matrix for $T:\R^2\rightarrow \R^3$ defined by $T(x_1,x_2) = (-x_2, 3x_1+x_2, 8x_1)$.

\exatend Find the standard matrix for $T:\R^4\rightarrow \R^2$ defined by $T(x_1,x_2,x_3,x_1) = (-2x_2+\pi x_4, x_1+x_2+x_3+x_4)$.

\exatend Find the standard matrix for an $\R^2\rightarrow\R^2$ that vertically shears $\vec{e}_1$ to $\vec{e}_1+0.5\vec{e}_2$ but leaves the $x_2$-axis unchanged.

\exatend (T/F?) The transformation $T:\R^2\rightarrow\R^2$ defined by $T(\vec{x}) = \bbvec{3}{-2} + \vec{x}$ is a linear transformation.

\exatend (T/F?) The transformation $T:\R^2\rightarrow\R^2$ defined by $T(x_1,x_2) = \bbvec{\sin x_1}{|x_2|}$ is a linear transformation.

\exatend (T/F?) If $A$ has a pivot position in every column, then the range of $T(\vec{x}) = A\vec{x}$ equals its codomain.

\exatend (T/F?) If $A$ has a pivot position in every column, and $T(\vec{x}) = A\vec{x}$, then the only solution of $T(\vec{x}) = \vec{0}$ is $\vec{x}=\vec{0}$.

\exatend (T/F?) If $A$ has a pivot position in every row, then the range of $T(\vec{x}) = A\vec{x}$ equals its codomain.

\exatend (T/F?) If $A$ has a pivot position in every row, and $T(\vec{x}) = A\vec{x}$, then the only solution of $T(\vec{x}) = \vec{0}$ is $\vec{x}=\vec{0}$.

\exatend Show that if $T:\R^c\rightarrow\R^r$ is a linear transformation, then $T(\vec{0}) = \vec{0}$.

\exatend  Suppose $T:\R^c\rightarrow \R^r$ is a linear transformation, and let $\vec{x}, \vec{y}, \vec{z}$ be a linearly dependent list in $\R^c$. Prove that $T(\vec{x}), T(\vec{y}), T(\vec{z})$ is linearly dependent.

\exatend Find a counterexample to the following false statement: If $T:\R^c\rightarrow \R^r$ is a linear transformation, and $\vec{x}, \vec{y}, \vec{z}$ is a linearly independent list in $\R^c$, then $T(\vec{x}), T(\vec{y}), T(\vec{z})$ is linearly independent. \hintfootnote{You need an independent list $\vec{x},\vec{y},\vec{z}$, and a linear $T$ that makes $T(\vec{x}), T(\vec{y}), T(\vec{z})$ dependent.}

\exatend Suppose $T:\R^c\rightarrow\R^r$ and $S:\R^r\rightarrow\R^n$ are linear transformations. Show that $Q:\R^c\rightarrow\R^n$ defined by $Q(\vec{x}) = S(T(\vec{x}))$ is linear.

\exatend Suppose $A$ and $B$ are $r \times c$ matrices, and that $A \vec{x} = B \vec{x}$ for all $\vec{x}$ in $\R^{c}$. Prove that $A = B$. (Note: you can't just divide by $\vec{x}$. \hintfootnote{What happens if $\vec{x} = \vec{e}_{1}$ or $\vec{e}_2$?}

\exatend Suppose $\vec{x}$ and $\vec{z}$ are two points in $\R^c$, and $T:\R^c\rightarrow\R^r$ is linear.
\begin{enumerate}[label=(\alph*)]
\item Let $\vec{w}$ be the point that is 70\% of the way from $\vec{x}$ to $\vec{z}$ along a straight line. Write a formula for $\vec{w}$ in terms of $\vec{x}$ and $\vec{z}$. \hintfootnote{Draw a picture. What does $\vec{z}-\vec{x}$ look like?}
\item Show that $T(\vec{w})$ is 70\% of the way from $T(\vec{x})$ to $T(\vec{z})$ along a straight line.
\end{enumerate}

\exatend Show that every linear transformation $T:\R \rightarrow \R$ has a formula $T(x) = kx$ for some real number $k$.

\exatend \emph{(Challenge)} Find the matrix of the transformation $T:\R^2\rightarrow\R^2$ that projects points onto the nearest point on the line $x_2=2x_1$. \hintfootnote{Let $\theta$ be the angle from the $x_1$-axis up to the line. Start by rotating points by $-\theta$, which brings the line to the $x_1$-axis. Then project points onto the $x_1$-axis, and then rotate them back by $\theta$.}

\exatend \emph{(Challenge)} Find an example of a function $f:\R^2\rightarrow\R$ such that $f(k\vec{v}) = kf(\vec{v})$ for all real numbers $k$ and all $\vec{v}$ in $\R^2$.


\columnbreak
\section{Matrix Multiplication}
\label{sec:matrix-multiplication}

\exercise (\emph{Warm-up}) Multiply:
\begin{horiparts}{2}

\item $\begin{bmatrix}
-2 & 3 \\
4 & 0
\end{bmatrix}\bbvec{1}{3}$
\item $\begin{bmatrix}
1 & 0 & 8 \\
0 & 1 & 7 
\end{bmatrix}\bbbvec{\pi}{e}{0}$
\end{horiparts}

\begin{notes}
\vfill \vfill
\end{notes}

\exercise (\emph{Warm-up}) Let $A = \begin{bmatrix}
  0 & -4 \\
  3 & -2
\end{bmatrix}$ and $B = \begin{bmatrix}
  0 & 1 & 2 \\
  1 & 3 & -0.5
\end{bmatrix}$. Multiply:

\begin{horiparts}{2}
  \item $A \left( B \bbbvec{1}{1}{1} \right)$
  \item $\begin{bmatrix}
    -4 & -12 & 2 \\
    -2 & -3 & 7
  \end{bmatrix} \bbbvec{1}{1}{1}$
\end{horiparts}

\begin{notes}
\vfill
\end{notes}

Note: It turns out that $A(B \vec{x}) = \begin{bmatrix}
    0 & 4 & 2 \\
    2 & -10 & -18
  \end{bmatrix}\vec{x}$ for all $\vec{x}$ in $\R^3$. Later in this section we'll figure out a formula for how to multiply two matrices that makes $AB = \begin{bmatrix}
    0 & 4 & 2 \\
    2 & -10 & -18
\end{bmatrix}$ and in general, $A(B \vec{x}) = (AB) \vec{x}$.



\exercise \label{exercise:2.2wu-dot-product} (\emph{Warm-up}) Recall that the \emph{dot product} of $(u_1,\ldots,u_c)$ and $(v_1,\ldots,v_c)$ is $u_1v_1 + \cdots + u_cv_c$.
\begin{enumerate}[label=(\alph*)]
\item  Multiply, and write as a single vector (even though you can't combine like terms):
\[\begin{bmatrix}
1 & 2 & 3 \\
4 & 5 & 6 \\
7 & 8 & 9
\end{bmatrix}\bbbvec{x_1}{x_2}{x_3}\]
\item (T/F?) In general, the $i^{\text{th}}$ entry in $A\vec{x}$ is the dot product of row $i$ of $A$ and $\vec{x}$.
 \end{enumerate}

\begin{notes}
\vfill \hspace{2pt}
\columnbreak
\end{notes}

Matrix scaling and addition work the same as with vectors, which are really just single-column matrices. If the matrices are different sizes, their sum is undefined. Also, we'll recycle the symbol $0$ to stand for a matrix where all the entries equal the number $0$.

\example Simplify:
\begin{enumerate}[label=(\alph*)]
\item $3I_2 + 0$

\begin{notes}
\vspace{1em}
\end{notes}

\item $\begin{bmatrix}
0 & 4 & 8 \\
12 & 16 & 20
\end{bmatrix}-2\begin{bmatrix}
\pi & e \\
7 & 6 \\
2 & 0
\end{bmatrix}$

\begin{notes}
\vspace{1em}
\end{notes}
\item $\begin{bmatrix}
10 & 11 & 12\\
13 & 14 & 15
\end{bmatrix} -
\begin{bmatrix}[c|c]
3I_2 & \vec{0}\end{bmatrix}$
\end{enumerate}

\begin{notes}
\vspace{1em}
\end{notes}

\begin{book}
\textbf{Solution.} (a) $3I_2+0 = 3\begin{bmatrix}
1 & 0 \\ 0 & 1
\end{bmatrix} + \begin{bmatrix}
0 & 0 \\ 0 & 0
\end{bmatrix} = \begin{bmatrix}
3\cdot 1 - 0 & 3\cdot 0 - 0 \\ 3\cdot 0 - 0 & 3\cdot 1 - 0
\end{bmatrix} = \begin{bmatrix}
3 & 0 \\
0 & 3
\end{bmatrix}$  

(b) This is undefined because the matrices are different sizes ($3\times 2$ and $2\times 3$).

(c) $\begin{bmatrix}
10 & 11 & 12 \\
13 & 14 & 15 
\end{bmatrix}- \begin{bmatrix}
3 & 0 & 0 \\
0 & 3 & 0
\end{bmatrix} = \begin{bmatrix}
7 & 11 & 12 \\
13 & 11 & 15
\end{bmatrix}$ \blacksquare
\end{book}

\begin{theorem}
\label{theorem:matrix-addition-scaling}
If $A,B,$ and $F$ are equal-sized matrices, and $p,q$ are real numbers, then:
\begin{enumerate}[label=(\roman*)]
\item $A+B = B+A$
\item $A + (B+F) = (A+B) +F$
\item $0 + A = A$
\item $(p+q)A = pA + qA$
\item $p(A+B) = pA+pB$
\item $p(qA) = (pq)A$
\end{enumerate}
\end{theorem}

Each of these are true because of a corresponding property of real number addition/multiplication that applies to each of the entries in the matrices. Proofs are left as exercises.

\begin{notes}
Matrix multiplication is where it gets complicated. The goal is to define $(AB)$ so that $(AB)\vec{x} = A(B\vec{x})$. Suppose $A =\begin{bmatrix}
\vec{a}_1 & \vec{a}_2 & \cdots & \vec{a}_c
\end{bmatrix}$ is a $r \times c$ matrix and $B =\begin{bmatrix}
\vec{b}_1 & \vec{b}_2 & \cdots & \vec{b}_{d}
\end{bmatrix}$ is a $c \times d$ matrix. Let's fill in the blanks to check our understanding:

$B$ defines a linear transformation that takes any input vector $\vec{x}$ in $\R^{\rule{1em}{0.8pt}}$ and gives an output vector $B\vec{x} = \rule{8em}{0.8pt}$ in $\R^{\rule{1em}{0.8pt}}$.

Since the number of entries in $B\vec{x}$ equals the number of columns of $A$, we know $A(B\vec{x})$ is defined:
\begin{align*}
A(B\vec{x}) &= A(\rule{10em}{0pt}\rule{0pt}{1.3em}) &(\text{defn of matrix-vector product}) \\
&= \rule{10em}{0pt}\rule{0pt}{1.3em} &(\text{linearity of matrix-vector product}) \\
&= \rule{10em}{0pt}\rule{0pt}{1.3em} &(\text{linearity of matrix-vector product}) \\
&= \left[\rule{10em}{0pt}\rule{0pt}{1.3em}\right]\vec{x} &(\text{defn of matrix-vector product})
\end{align*}
Thus, the following definition makes it so for any $\vec{x}$ in $\R^d$, we have $(AB)\vec{x} = A(B\vec{x})$.
\end{notes}

\begin{book}
Matrix multiplication is where it gets complicated. The goal is to define $(AB)$ so that $(AB)\vec{x} = A(B\vec{x})$. Suppose $A =\begin{bmatrix}
\vec{a}_1 & \vec{a}_2 & \cdots & \vec{a}_c
\end{bmatrix}$ is a $r \times c$ matrix and $B =\begin{bmatrix}
\vec{b}_1 & \vec{b}_2 & \cdots & \vec{b}_{d}
\end{bmatrix}$ is a $c \times d$ matrix. 

$B$ defines a linear transformation that takes any input vector $\vec{x}$ in $\R^{d}$ and gives an output vector $B\vec{x} = x_1\vec{b}_1 + x_2\vec{b}_2+\cdots+x_d\vec{b}_d$ in $\R^{c}$.

Since the number of entries in $B\vec{x}$ equals the number of columns of $A$, we know $A(B\vec{x})$ is defined:
\begin{align*}
A(B\vec{x}) &= A(x_1\vec{b}_1 + x_2\vec{b}_2+\cdots+x_d\vec{b}_d) &(\text{defn of } B\vec{x}) \\
&= A(x_1\vec{b}_1) + A(x_2\vec{b}_2)+\cdots+A(x_d\vec{b}_d) &(\text{thm \ref{theorem:linearity-of-Ax}(i)}) \\
&=  x_1(A\vec{b}_1) + x_2(A\vec{b}_2)+\cdots+x_d(A\vec{b}_d) &(\text{thm \ref{theorem:linearity-of-Ax}(ii)}) \\
&=\begin{bmatrix}
A\vec{b}_1 & A\vec{b}_2 & \cdots & A\vec{b}_d
\end{bmatrix}\vec{x}
&(\text{defn})
\end{align*}
Thus, the following definition makes it so for any $\vec{x}$ in $\R^d$, we have $(AB)\vec{x} = A(B\vec{x})$.
\end{book}

\begin{definition}
If $A$ is an $r \times c$ matrix, and $B = \begin{bmatrix}
\vec{b}_1 & \vec{b}_2 & \cdots & \vec{b}_{d}
\end{bmatrix}$ is a $c \times d$ matrix, then the product is the $r \times d$ matrix:
\[ AB  = \begin{bmatrix}
A\vec{b}_1 & A\vec{b}_2 & \cdots & A\vec{b}_{d}
\end{bmatrix}\]
\end{definition}

\example Let $A =\begin{bmatrix}
2 & 0 \\
-3 & 1
\end{bmatrix}$ and $B =\begin{bmatrix}
-1 & 4 & -2 \\
0 & 3 & 1
\end{bmatrix}$. Multiply:

\begin{horiparts}{2}
\item $AB$
\item $BA$
\end{horiparts}

\begin{notes}
\columnbreak
\end{notes}

\begin{book}
\textbf{Solution.} (a) Let's apply the definition of $AB$. You may ignore the color-coding for now, but it will be useful in a minute to help understand the row-column rule.
\begin{align*}\begin{bmatrix}
\textcolor{red}{2} & \textcolor{red}{0}\\
\textcolor{orange}{-3} & \textcolor{orange}{1}
\end{bmatrix}\begin{bmatrix}
\textcolor{purple}{-1} & \textcolor{blue}{4} & \textcolor{teal}{-2} \\
\textcolor{purple}{0} & \textcolor{blue}{3} & \textcolor{teal}{1} 
\end{bmatrix} &=\begin{bmatrix}
\begin{bmatrix}
\textcolor{red}{2} & \textcolor{red}{0}\\
\textcolor{orange}{-3} & \textcolor{orange}{1}
\end{bmatrix} \bbvec{\textcolor{purple}{-1}}{\textcolor{purple}{0}} &
\begin{bmatrix}
\textcolor{red}{2} & \textcolor{red}{0}\\
\textcolor{orange}{-3} & \textcolor{orange}{1}
\end{bmatrix} \bbvec{\textcolor{blue}{4}}{\textcolor{blue}{3}} &
\begin{bmatrix}
\textcolor{red}{2} & \textcolor{red}{0}\\
\textcolor{orange}{-3} & \textcolor{orange}{1}
\end{bmatrix} \bbvec{\textcolor{teal}{-2}}{\textcolor{teal}{1}} 
\end{bmatrix} \\
&= \begin{bmatrix}
\textcolor{red}{2}\cdot \textcolor{purple}{-1} + \textcolor{red}{0}\cdot \textcolor{purple}{0} &
\textcolor{red}{2}\cdot \textcolor{blue}{4} + \textcolor{red}{0}\cdot \textcolor{blue}{3} &
\textcolor{red}{2}\cdot \textcolor{teal}{-2} + \textcolor{red}{0}\cdot \textcolor{teal}{1} \\
\textcolor{orange}{-3}\cdot \textcolor{purple}{-1} + \textcolor{orange}{1}\cdot \textcolor{purple}{0} &
\textcolor{orange}{-3}\cdot \textcolor{blue}{4} + \textcolor{orange}{1}\cdot \textcolor{blue}{3} & 
\textcolor{orange}{-3}\cdot \textcolor{teal}{-2} + \textcolor{orange}{1}\cdot \textcolor{teal}{1}
\end{bmatrix} \\
&= \begin{bmatrix}
-2 & 8 & -4 \\
4 & -9 & 7
\end{bmatrix}
\end{align*}
(b) In the definition, the number of columns of the matrix on the left, $c$, must equal the number of rows of the matrix on the right. In sloppy notation:
\[(r\times c)(c\times d) \rightarrow (r \times d)\]
But $BA$ is $(2\times 3)(2\times 2)$, which is undefined since $3\ne 2$. \blacksquare
\end{book}

\example \label{example:2.2-multiplication-2} Let $A =\begin{bmatrix}
-1 & 0 \\
0 & 2
\end{bmatrix}$ and $B =\begin{bmatrix}
3 & 4 \\
5 & 6
\end{bmatrix}$. Multiply:

\begin{horiparts}{2}
\item $AB$
\item $BA$
\end{horiparts}

\begin{notes}
\vfill
\end{notes}

\begin{book}
\textbf{Solution.} (a) Let's apply the definition of $AB$. You may ignore the color-coding for now, but it will be useful in a minute to help understand the row-column rule.
\begin{align*}
\begin{bmatrix}
\textcolor{red}{-1} & \textcolor{red}{0} \\
\textcolor{orange}{0} & \textcolor{orange}{2}
\end{bmatrix} \begin{bmatrix}
\textcolor{blue}{3} & \textcolor{teal}{4} \\
\textcolor{blue}{5} & \textcolor{teal}{6}
\end{bmatrix} &= \begin{bmatrix}
\begin{bmatrix}
\textcolor{red}{-1} & \textcolor{red}{0} \\
\textcolor{orange}{0} & \textcolor{orange}{2}
\end{bmatrix} \bbvec{\textcolor{blue}{3}}{\textcolor{blue}{5}} & 
\begin{bmatrix}
\textcolor{red}{-1} & \textcolor{red}{0} \\
\textcolor{orange}{0} & \textcolor{orange}{2}
\end{bmatrix} \bbvec{\textcolor{teal}{4}}{\textcolor{teal}{6}} 
\end{bmatrix}\\
&= \begin{bmatrix}
\textcolor{red}{-1}\cdot \textcolor{blue}{3} + \textcolor{red}{0}\cdot \textcolor{blue}{5} &
\textcolor{red}{-1}\cdot \textcolor{teal}{4} + \textcolor{red}{0}\cdot \textcolor{teal}{6} \\
\textcolor{orange}{0}\cdot \textcolor{blue}{3} + \textcolor{orange}{2}\cdot \textcolor{blue}{5} &
\textcolor{orange}{0}\cdot \textcolor{teal}{4} + \textcolor{orange}{2}\cdot \textcolor{teal}{6} \\
\end{bmatrix} = \begin{bmatrix}
-3 & -4 \\
10 & 12
\end{bmatrix}
\end{align*}
(b) Similarly:
\begin{align*}
\begin{bmatrix}
\textcolor{red}{3} & \textcolor{red}{4} \\
\textcolor{orange}{5} & \textcolor{orange}{6}
\end{bmatrix} \begin{bmatrix}
\textcolor{blue}{-1} & \textcolor{teal}{0} \\
\textcolor{blue}{0} & \textcolor{teal}{2}
\end{bmatrix} &= \begin{bmatrix}
\begin{bmatrix}
\textcolor{red}{3} & \textcolor{red}{4} \\
\textcolor{orange}{5} & \textcolor{orange}{6}
\end{bmatrix} \bbvec{\textcolor{blue}{-1}}{\textcolor{blue}{0}} & 
\begin{bmatrix}
\textcolor{red}{3} & \textcolor{red}{4} \\
\textcolor{orange}{5} & \textcolor{orange}{6}
\end{bmatrix} \bbvec{\textcolor{teal}{0}}{\textcolor{teal}{2}} 
\end{bmatrix}\\
&= \begin{bmatrix}
\textcolor{red}{3}\cdot \textcolor{blue}{-1} + \textcolor{red}{4}\cdot \textcolor{blue}{0} &
\textcolor{red}{3}\cdot \textcolor{teal}{0} + \textcolor{red}{4}\cdot \textcolor{teal}{2} \\
\textcolor{orange}{5}\cdot \textcolor{blue}{-1} + \textcolor{orange}{6}\cdot \textcolor{blue}{0} &
\textcolor{orange}{5}\cdot \textcolor{teal}{0} + \textcolor{orange}{6}\cdot \textcolor{teal}{2} \\
\end{bmatrix} = \begin{bmatrix}
-3 & 8 \\
-5 & 12
\end{bmatrix}
\end{align*}
Note that $AB$ is not the same as $BA$. \blacksquare
\end{book}

Recall from exercise \ref{exercise:2.2wu-dot-product} that the $i^{\text{th}}$ entry in $A\vec{x}$ is the dot product of row $i$ of $A$ and $\vec{x}$. If we apply this rule to the $i^{\text{th}}$ entry in $A\vec{b}_j$ (the $j^{\text{th}}$ column in $AB$) we get the following rule:

\begin{theorem}[Row-column rule.]
\label{theorem:row-column-rule}
The entry in row $i$, column $j$ of $AB$ is the dot product of row $i$ of $A$ and column $j$ of $B$.

In other words, if we write $(AB)_{ij}$ for the entry in row $i$, column $j$ of $AB$, then:
\[(AB)_{ij} = A_{i1}B_{1j} + A_{i2}B_{2j} + \cdots + A_{ic}B_{cj}\]
\end{theorem}

The original definition of $AB$ is very useful for conceptual understanding, proofs, and computations with GPUs that can do lots of multiplications in parallel. The row-column rule is often more comfortable for hand calculations.

\example Repeat the last example, but use the row-column rule. 

\begin{notes}
\vfill \hspace{2pt}
\columnbreak
\end{notes}

\begin{book}
\textbf{Solution.} Fortunately this was already done in the solution of the last example, so we don't have to write everything again. Instead, let's make sure we understand the notation by looking at a couple example entries. For example, let's plug in $i=2$ and $j=1$ to the row-column rule:
\begin{align*}
(AB)_{21}&= \textcolor{orange}{A_{21}}\textcolor{blue}{B_{11}} + \textcolor{orange}{A_{22}}\textcolor{blue}{B_{21}} \\
&= \textcolor{orange}{0}\cdot \textcolor{blue}{3} + \textcolor{orange}{2}\cdot \textcolor{blue}{5}
\end{align*}
Each other entry in $AB$ is similarly a dot product of a row of $A$ and a column of $B$, which you should check. For $BA$, each entry is a dot product of a row of $B$ and a column of $A$. \blacksquare
\end{book}

\example True/False?
\begin{enumerate}[label=(\alph*)]
\item If $MA = MB$, and $M \ne 0$, then $A = B$.

\begin{notes}
\vfill
\end{notes}
\item If $AB = 0$ then either $A = 0$ or $B = 0$.
\end{enumerate}

\begin{book}
\textbf{Solution.} (a) Let's think about how to make it so some entries in $A$ and $B$ can be multiplied by zero when computing $MA$ and $MB$. In example \ref{example:2.2-multiplication-2} part (a), notice that row 1 entries in $B$ get multiplied by the column 1 entries in $A$. To construct our counterexample, we could make it so that column 1 of $M$ is zero, making the row 1 entries of $A$ and $B$ irrelevant. If $M=\begin{bmatrix}
0 & 1 \\
0 & 2
\end{bmatrix}$, $A = \begin{bmatrix}
3 & 4 \\
5 & 6
\end{bmatrix}$, and $B = \begin{bmatrix}
-9 & 17 \\
5 & 6
\end{bmatrix}$, then:
\[MA = \begin{bmatrix}
0 & 1 \\
0 & 2
\end{bmatrix} \begin{bmatrix}
3 & 4 \\
5 & 6
\end{bmatrix} = \begin{bmatrix}
0\cdot 3 + 1\cdot 5 & 0 \cdot4 + 1\cdot 6 \\
0\cdot 3 + 2\cdot 5 & 0 \cdot4 + 2\cdot 6 
\end{bmatrix} = \begin{bmatrix}
5 & 6 \\
10 & 12
\end{bmatrix}\]
Likewise:
\[MB = \begin{bmatrix}
0 & 1 \\
0 & 2
\end{bmatrix} \begin{bmatrix}
-9 & 17 \\
5 & 6
\end{bmatrix} = \begin{bmatrix}
0\cdot -9 + 1\cdot 5 & 0 \cdot17 + 1\cdot 6 \\
0\cdot -9 + 2\cdot 5 & 0 \cdot17 + 2\cdot 6 
\end{bmatrix} = \begin{bmatrix}
5 & 6 \\
10 & 12
\end{bmatrix}\]
So we have $MA=MB$, $M\ne 0$, and $A \ne B$ and the statement is false.

(b) We can use the same trick as part (a) to make it so all non-zero entries get multiplied by zero. If $A = \begin{bmatrix}
0 & 1 \\
0 & 2
\end{bmatrix}$ and $B=\begin{bmatrix}
5 & 6 \\
0 & 0
\end{bmatrix}$ then:
\[AB = \begin{bmatrix}
0 & 1 \\
0 & 2
\end{bmatrix}\begin{bmatrix}
5 & 6 \\
0 & 0
\end{bmatrix} = \begin{bmatrix}
0 \cdot 5 + 1\cdot 0 & 0\cdot 6 + 1\cdot 0 \\
0 \cdot 5 + 2\cdot 0 & 0\cdot 6 + 2\cdot 0 
\end{bmatrix} = 0\]
So the statement is false. \blacksquare
\end{book}

\begin{notes}
\vfill \hspace{2pt}
\columnbreak
\end{notes}

The last example shows you can't assume that properties of real numbers apply to matrices as well. But fortunately, numbers and matrices do have a lot in common.

\begin{theorem}[Properties of matrix multiplication.]
\label{theorem:properties-of-matrix-multiplication}
If $A,B,$ and $C$ have sizes for which the products and sums are defined, then:

\begin{enumerate}[label=(\roman*)]
\item $A(BC) = (AB)C$ \quad (matrix multiplication is \emph{associative})
\item If $A$ is $r\times c$, then $I_rA = AI_c = A$
\item $A(B+C) = AB+AC$
\item $(A+B)C = AC + BC$
\item If $k$ is a real number, then $k(AB) = (kA)B = A(kB)$
\end{enumerate}
\end{theorem}

Proofs of (ii) through (v) are left as exercises.

\begin{book}
\textbf{Proof.} (i) Let $C =\begin{bmatrix}
\vec{c}_1 & \vec{c}_2 & \cdots & \vec{c}_n
\end{bmatrix}$. Then:
\begin{align*}
 A(BC) &= A\begin{bmatrix}
B\vec{c}_1 & B\vec{c}_2 & \cdots & B\vec{c}_n
\end{bmatrix} \\
&= \begin{bmatrix}
A (B\vec{c}_1) & A(B\vec{c}_2) & \cdots & A( B\vec{c}_n)
\end{bmatrix} \\
&= \begin{bmatrix}
(AB)\vec{c}_1 & (AB)\vec{c}_2 & \cdots & (AB)\vec{c}_n 
\end{bmatrix} \\
&= (AB)C \quad \blacksquare
\end{align*}
\end{book}

\begin{notes}
\vfill
\end{notes}

\begin{definition}
The \textbf{transpose} of an $r\times c$ matrix $A$, written $A^T$, is the $c \times r$ matrix with rows and columns swapped. In other words, $(A^T)_{ij} = A_{ji}$.
\end{definition}

\begin{book}
\example $\begin{bmatrix}
1 & 2 & 3 & 4 \\
5 & 6 & 7 & 8
\end{bmatrix}^T = \begin{bmatrix}
1 & 5 \\
2 & 6 \\
3 & 7 \\
4 & 8 
\end{bmatrix}$ \blacksquare
\end{book}
\begin{notes}
\example $\begin{bmatrix}
1 & 2 & 3 & 4 \\
5 & 6 & 7 & 8
\end{bmatrix}^T = $

\vspace{3em} \hspace{2pt}
\columnbreak
\end{notes}

\begin{notes}
\example Suppose $A =\begin{bmatrix}
a_{11}  & \cdots & a_{1c} \\
a_{21}  & \cdots & a_{2c} \\
\vdots & \vdots & \vdots \\
a_{r1} & \cdots & a_{rc}
\end{bmatrix}$, and we write the rows of $A$ as $\vec{\alpha}_1 = \bbbvec{a_{11}}{\vdots}{a_{1c}}, \vec{\alpha}_2 = \bbbvec{a_{21}}{\vdots}{a_{2c}}, \ldots, \vec{\alpha}_r = \bbbvec{a_{r1}}{\vdots}{a_{rc}}$. Then,
\begin{align*}
\vec{\alpha}_1^T &= \\ 
\vec{\alpha}_2^T &= \\ 
&\vdots \\
\vec{\alpha}_r^T &= 
\end{align*}
So we can write $A$ in terms of its rows:
\[A = \]
\vspace{3em}
\end{notes}

\begin{book}
\example Suppose $A =\begin{bmatrix}
a_{11}  & \cdots & a_{1c} \\
a_{21}  & \cdots & a_{2c} \\
\vdots & \vdots & \vdots \\
a_{r1} & \cdots & a_{rc}
\end{bmatrix}$, and we write the rows of $A$ as $\vec{\alpha}_1 = \bbbvec{a_{11}}{\vdots}{a_{1c}}, \vec{\alpha}_2 = \bbbvec{a_{21}}{\vdots}{a_{2c}}, \ldots, \vec{\alpha}_r = \bbbvec{a_{r1}}{\vdots}{a_{rc}}$. Then,
\begin{align*}
\vec{\alpha}_1^T &= \begin{bmatrix}
a_{11} & \cdots & a_{1c} 
\end{bmatrix} \\ 
\vec{\alpha}_2^T &= \begin{bmatrix}
a_{21} & \cdots & a_{2c} 
\end{bmatrix}\\ 
&\vdots \\
\vec{\alpha}_r^T &= \begin{bmatrix}
a_{r1} & \cdots & a_{rc} 
\end{bmatrix}\\ 
\end{align*}
So we can write $A$ in terms of its rows:
\[A =\begin{bmatrix}
\vec{\alpha}_1^T \\
\vec{\alpha}_2^T \\
\vdots \\
\vec{\alpha}_r^T 
\end{bmatrix} \quad \quad \blacksquare\]
\end{book}

\begin{theorem}\label{theorem:transpose-basics}
If $A$ and $B$ are equal-sized matrices, and $k$ is a real number, then:
\begin{enumerate}[label=(\roman*)]
\item  $( A^T )^T=A$
\item $(kA)^T = kA^T$
\item $(A+B)^T = A^T + B^T$
\end{enumerate}
\end{theorem}
The proof is straightforward and left for exercises.

\begin{notes}
\columnbreak
\end{notes}

\example Suppose $A$ and $B$ have sizes that make $AB$ defined. Is $A^TB^T$ necessarily defined? Is $B^TA^T$ necessarily defined?

\begin{notes}
\vfill
\end{notes}

\begin{book}
If $AB$ is defined, then the number of columns of $A$ must equal the number of rows of $B$: we can say $A$ is $r \times c$ and $B$ is $c\times d$. This means $A^T$ is $c \times r$ and $B^T$ is $d\times c$. Since $r$ is not necessarily equal to $d$, $A^TB^T$ may be undefined. However, $B^TA^T$ is well-defined, because the $c$'s line up in the middle again. \blacksquare
\end{book}

\begin{theorem} \label{theorem:ABT-BTAT}
If $AB$ is defined, then $(AB)^T = B^TA^T$.
\end{theorem}

\begin{book}
\textbf{Proof.} The following steps show the $ij$-entry of $(AB)^T$ equals the $ij$-entry of $B^TA^T$: 
\begin{align*}
(AB)^T_{ij} &= (AB)_{ji} & \\ 
&= A_{j1}B_{1i} + A_{j2}B_{2i} + \cdots + A_{jc}B_{ci} & (\text{row-column rule}) \\
&= B_{1i}A_{j1} + B_{2i}A_{j2} + \cdots + B_{ci}A_{jc} & \\
&= B^T_{i1}A^T_{1j} + B^T_{i2}A^T_{j2} + \cdots + B^T_{ic}A^T_{cj} &(\text{defn of transpose}) \\
&= (B^TA^T)_{ij} &(\text{row-column rule})
\end{align*}
\blacksquare
\end{book}

\begin{notes}
\vfill
\end{notes}

\example Let $\vec{x}=(1,2,3)$. Multiply:

\begin{horiparts}{2}
\item $\vec{x}^T\vec{x}$
\item $\vec{x}\vec{x}^T$
\end{horiparts}

\begin{notes}
\vfill \hspace{2pt}
\end{notes}

\begin{book}
\textbf{Solution.} (a) $\vec{x}^T\vec{x}=\begin{bmatrix}
1 & 2 & 3
\end{bmatrix}\bbbvec{1}{2}{3}= 1\cdot 1 + 2\cdot 2+ 3\cdot 3 = 14$.

(b) $\vec{x}\vec{x}^T=\bbbvec{1}{2}{3}\begin{bmatrix}
1 & 2 & 3
\end{bmatrix} = \begin{bmatrix}
1\cdot 1 & 1\cdot 2 & 1\cdot 3 \\
2\cdot 1 & 2\cdot 2 & 2\cdot 3 \\
3\cdot 1 & 3\cdot 2 & 3\cdot 3 
\end{bmatrix} = \begin{bmatrix}
1 & 2 & 3 \\
2 & 4 & 6 \\
3 & 6 & 9
\end{bmatrix}$. \blacksquare
\end{book}

\columnbreak
\exerciseheading

For exercises \ref{exercise:2.2-mult-start} through \ref{exercise:2.2-mult-end}, either multiply the matrices, or if it's undefined, say so.

\exatend \label{exercise:2.2-mult-start} $\begin{bmatrix}
2 & 3 \\
4 & 5 \\
6 & 7 
\end{bmatrix}\begin{bmatrix}
1 & 0 & 0.5\\
-1 & 2 & 0
\end{bmatrix}$

\exatend $\begin{bmatrix}
1 & 0 & 0.5\\
-1 & 2 & 0
\end{bmatrix}\begin{bmatrix}
2 & 3 \\
4 & 5 \\
6 & 7
\end{bmatrix}$

\exatend $\begin{bmatrix}
1 & 0 & 0 \\
0 & 4 & 0 \\
2 & 0 & 1
\end{bmatrix}\begin{bmatrix}
1 & 0 & 0 \\
0 & 0.25 & 0 \\
-2 & 0 & 1
\end{bmatrix}$

\exatend $\begin{bmatrix}
2 \\ 0 \\ 4
\end{bmatrix}^T\begin{bmatrix}
3 \\ 5\\ -1
\end{bmatrix}$

\exatend $\begin{bmatrix}
2 \\0 \\ 4 
\end{bmatrix}\begin{bmatrix}
3 \\ 5 \\ -1
\end{bmatrix}^{T}$

\exatend $\begin{bmatrix}
3 & 2 \\
1 & 0
\end{bmatrix}\begin{bmatrix}
1 & 0 \\
0 & 1 \\
4 & \pi 
\end{bmatrix}$

\exatend $\begin{bmatrix}
1 & 0 \\
0 & 1 \\
4 & \pi 
\end{bmatrix}\begin{bmatrix}
3 & 2 \\
1 & 0
\end{bmatrix}$

\exatend $\begin{bmatrix}
7 & -3 \\
-2 & 1
\end{bmatrix}\begin{bmatrix}
1 & 3 \\
2 & 7
\end{bmatrix}$

\exatend \label{exercise:2.2-mult-end} $\bbvec{7}{-2}\begin{bmatrix}
1 & 3
\end{bmatrix} + \bbvec{-3}{1}
\begin{bmatrix}
2 & 7
\end{bmatrix}$

\exatend If $A =\begin{bmatrix}
1 & 3 \\
2 & 4
\end{bmatrix}$, simplify:
\begin{horiparts}{3}
\item $2A$ \item $A + A^T$ \item $A+10I_2$
\end{horiparts}

\exatend If $E =\begin{bmatrix}
5 & 0 & 0  \\
0 & 0 & 1 \\
0 & 1 & 0
\end{bmatrix}$, simplify:
\begin{horiparts}{3}
\item $A^2$ \item $A^2-I_3$ \item $2A-I_3$
\end{horiparts}

\exatend Suppose $A$ is $5 \times 7$, $B$ is $7\times 4$, $\vec{x}$ is in $\R^7$, $\vec{y}$ is in $\R^5$. For each of the following products, determine whether it's defined or not, and if it's defined, give the size of the result.

\begin{horiparts}{3}
\item $AB$
\item $(AB)^T$
\item $A\vec{x}$
\item $\vec{y}A$
\item $\vec{y}^TA$
\item $A \vec{y}$
\item $BA$
\item $AA^T$
\item $A^TA$
\end{horiparts}

\exatend Suppose $F$ is $6 \times 3$, $G$ is $5\times 3$, $\vec{w}$ is in $\R^3$, $\vec{z}$ is in $\R^5$. For each of the following products, determine whether it's defined or not, and if it's defined, give the size of the result.

\begin{horiparts}{3}
\item $FG$
\item $FG^T$
\item $(GF^T)^T$
\item $G\vec{w}$
\item $F\vec{w}$
\item $\vec{w}^TF^T$
\item $\vec{z}^TG$
\item $\vec{z}\vec{z}$
\item $\vec{z}^T\vec{z}$
\end{horiparts}

% basic T/F
\exatend (T/F?) The linear transformation $T(\vec{x}) = (AB)\vec{x}$ first multiplies $\vec{x}$ by $A$ and then multiplies the result by $B$.

\exatend (T/F?) If $A$ is $3\times2$ and $B$ is $2\times3$, then both $AB$ and $BA$ are defined.

\exatend (T/F?) The columns of $AB$ are linear combinations of the columns of $B$.

\exatend (T/F?) The columns of $AB$ are linear combinations of the columns of $A$.

\exatend (T/F?) If $AB = 0$, then either $A=0$ or $B=0$.

\exatend (T/F?) If $A,B,C$ are matrices that makes the products defined, then $(AB)C = A(BC)$.

\exatend (T/F?) If $A$ and $B$ have equal sizes, then $(A+B)^T = (B+A)^T$.

\exatend (T/F?) If $A,B$ are $2\times2$ matrices, and the two columns of $A$ are equal, then the two columns of $AB$ are equal.

\exatend (T/F?) If $A,B$ are $2\times2$ matrices, and the two columns of $B$ are equal, then the two columns of $AB$ are equal.

\exatend Suppose $A = \begin{bmatrix}
5 & -6 \\
4 & -5
\end{bmatrix}$ and $B = \begin{bmatrix}
-7 & k \\
-6 & 8
\end{bmatrix}$. Find $k$ that makes $AB = BA$, if possible.

\exatend Suppose $A = \begin{bmatrix}
1 & 2 \\
3 & 4
\end{bmatrix}$ and $B = \begin{bmatrix}
1 & 3 \\
2 & k
\end{bmatrix}$. Find $k$ that makes $AB = BA$, if possible.

\exatend A transformation $T:\R^2 \rightarrow \R^2$ first rotates points $30^{\circ}$ degrees counterclockwise, then reflects them across the $x_2$-axis. Find the standard matrix for $T$ two ways:
\begin{enumerate}[label=(\alph*)]
\item By finding $T(\vec{e}_1)$ and $T(\vec{e}_2)$ directly, like in section  \ref{sec:matrix-vector-products}.
\item By finding the matrix $R_{30^{\circ}}$ that rotates, and the matrix $F$ that reflects, and multiplying them in the proper order.
\end{enumerate}

\exatend A transformation $T:\R^3 \rightarrow \R^3$ first rotates points $90^{\circ}$ around the $y$-axis, in a direction that rotates the $x$-axis onto the $z$-axis. Then, it rotates points $30^{\circ}$ around the $x$ axis, in a direction that moves the $y$-axis towards the $z$-axis. Find the standard matrix for $T$ two ways:
\begin{enumerate}[label=(\alph*)]
\item By finding $T(\vec{e}_1)$, $T(\vec{e}_2)$, and $T(\vec{e}_3)$ directly, like in section \ref{sec:matrix-vector-products}.
\item By finding two matrices for the two rotations and multiplying them in the proper order.
\end{enumerate}

\exatend Find $R_{45^{\circ}}$, the matrix for the transformation that rotates points $45^{\circ}$ counterclockwise. Then find $R_{-45^{\circ}}$. Multiply them to check that $R_{45^{\circ}}R_{-45^{\circ}} = I$. \creditfootnote{Inspired by an exercise in Sergei Treil's excellent book \emph{Linear Algebra Done Wrong}.}

\exatend Note that if you rotate points around the origin by an angle of $\beta$ and then by an angle of $\alpha$, it is as if you moved them by an angle of $\alpha + \beta$. Multiply the rotation matrices $R_{\alpha}$ and $R_{\beta}$ to derive formulas for $\cos(\alpha+\beta)$ and $\sin(\alpha+\beta)$. \creditfootnote{Inspired by an exercise in Sergei Treil's excellent book \emph{Linear Algebra Done Wrong}.}


%just assign maybe one of these, they're pretty bland.
\exatend Prove theorem \ref{theorem:matrix-addition-scaling} part (i).

\exatend Prove theorem \ref{theorem:matrix-addition-scaling} part (iv).

\exatend Prove theorem \ref{theorem:matrix-addition-scaling} part (vi).

\exatend Prove theorem \ref{theorem:properties-of-matrix-multiplication} part (ii).

\exatend Prove theorem \ref{theorem:properties-of-matrix-multiplication} part (iii).

\exatend Prove theorem \ref{theorem:transpose-basics} part (i).

\exatend Prove theorem \ref{theorem:transpose-basics} part (ii).

\exatend Which of these is true about $(A^TA)^T$? (pick one.)
\begin{enumerate}[label=(\Alph*)]
\item It might be undefined.
\item It's always equal to $A^TA$.
\item It's always equal to $AA^T$.
\end{enumerate}

%harder T/F
\exatend (T/F?) If $AB$ has a column with all zeros, then the columns of $A$ are linearly dependent.

\exatend (T/F?) If $AB$ has a column with all zeros, then the columns of $B$ are linearly dependent.

\exatend Find a counterexample that shows this is false:
 If $A,B$ are square matrices, and $AB = 0$, then $BA=0$.

\exatend Show that if $A$ is a scalar multiple of the $n\times n$ identity matrix, and $B$ is any $n\times n$ matrix, then $AB = BA$.

\exatend \label{exercise:2.2-diagonal-commute} A \emph{diagonal matrix} is a square matrix where all the entries \emph{not} on the main diagonal ($a_{ij}$ with $i\ne j$) are zero. Show that if $A$ and $B$ are $n\times n$ diagonal matrices, then $AB = BA$.

%TODO: result (a) is a theorem in 2.3; let's keep this here for fun though. But let's add some examples of elementary matrices to motivate this idea.
\exatend (\emph{challenge}) \label{exercise:2.2AB-rows-are-combos-of-B-rows}
\begin{enumerate}[label=(\alph*)]
\item Show that each row of $AB$ is a linear combination of the rows of $B$. \hintfootnote{Theorem \ref{theorem:ABT-BTAT}.}
\item Show that if $AB$ has a row of all zeros, then the rows of $B$ are linearly dependent.
\end{enumerate}

%treil credit
\exatend (\emph{challenge}) Find the matrix for the $\R^2\rightarrow\R^2$ transformation that reflects points across the line $x_2=\frac{-3}{2}x_1$, by finding matrices that rotate the line onto the $x_1$-axis, reflect across the $x_1$-axis, and then rotate the $x_1$-axis back to the line, and then multiplying the matrices in the right order. \creditfootnote{Inspired by Treil.}

\columnbreak
\section{Inverses}
\label{sec:inverses}

\exercise (\emph{Warm-up}) Let's recall basic facts of ``multiplicative inverses'' of real numbers.
\begin{enumerate}[label=(\alph*)]
\item (T/F?) There exists a real number $x$ such that $5x = x5 = 1$.
\item (T/F?) There exists a real number $y$ such that $0y = y0 = 1$.
\end{enumerate}

\exercise (\emph{Warm-up}) Let $A =\begin{bmatrix}
3 & 2 \\
-2 & -2
\end{bmatrix}$. 
\begin{enumerate}[label=(\alph*)]
\item Solve the equations $A\vec{x}_1=\vec{e}_1$ and $A\vec{x}_2=\vec{e}_2$.
\item Let $X = \begin{bmatrix}
\vec{x}_1 & \vec{x}_2
\end{bmatrix}$, with your solutions from (a) in the columns, and multiply: $AX$. Are you surprised by the result?
\item Multiply $XA$. Are you surprised by the result?
\item (T/F?) $A$ is row-equivalent to $I$, and the row operations that reduced $A$ to $I$ also reduced $\vec{e}_1$ to $\vec{x}_1$, and $\vec{e}_2$ to $\vec{x}_2$.
\end{enumerate}

\begin{notes}
\vfill \hspace{2pt}
\columnbreak
\end{notes}

\begin{definition}
A square ($c\times c$) matrix $A$ is called \textbf{invertible} if there exists a $c \times c$ matrix $A^{-1}$, called the \textbf{inverse} of $A$, such that:

\[ A^{-1}A = AA^{-1} = I\]
\end{definition}

If there are two matrices $P$ and $Q$ each satisfying the definition of the inverse of $A$, then
\[P = PI = P(AQ) = (PA)Q = IQ = Q \]
and so the inverse of $A$ is unique.

\example \label{example:2.3-first-elementary} Let $E = \begin{bmatrix}
1 & 0 & 0 \\
0 & 1 & 0 \\
4 & 0 & 1
\end{bmatrix}$ and $F = \begin{bmatrix}
1 & 0 & 0 \\
0 & 1 & 0 \\
-4 & 0 & 1
\end{bmatrix}$. Show that $E = F^{-1}$ and $F = E^{-1}$.

\begin{notes}
\vfill \vfill
\end{notes}

\begin{book}
\textbf{Solution.} To show $E = F^{-1}$, we need to show that $EF = FE = I$ (take the definition of $A^{-1}$ above and replace $A$ with $F$ and $A^{-1}$ with $E$). Likewise to show $F = E^{-1}$  we need to show that $FE = EF = I$; the same thing. Note that this means $E = F^{-1}$ if and only if $F = E^{-1}$.

The reader should check, using the row-column rule, that the first two rows of $EF$ are $\begin{bmatrix}
1 & 0 & 0 \\
0 & 1 & 0
\end{bmatrix}$. Let's see what interesting things happen in row 3 of $EF$, writing the entries from $F$ in parentheses:
\begin{align*}
\text{row 3 of } EF &= \begin{bmatrix}
4(1) + 0(0) + 1(-4) & 4(0) + 0(1) + 1(0) & 4(0) + 0(0) + 1(1)
\end{bmatrix} \\
&= \begin{bmatrix}
0 & 0 & 1
\end{bmatrix}
\end{align*}

Note that this 3rd row of $EF$ is a linear combination of the rows of $F$ using the entries from row 3 of $E$ as weights:
\[\text{row 3 of } EF = 4\begin{bmatrix}
1 & 0 & 0
\end{bmatrix} + 0\begin{bmatrix}
0 & 1 & 0
\end{bmatrix} + 1\begin{bmatrix}
-4 & 0 & 1
\end{bmatrix} \]

Either way you do the multiplication, you get $EF = I$. The reader should check that $FE = I$ as well. \blacksquare
\end{book}

\example Explain why $A = \begin{bmatrix}
6 & 0 \\
7 & 0
\end{bmatrix}$ is not invertible.

\begin{notes}
\vfill \hspace{2pt}
\columnbreak
\end{notes}

\begin{book} \textbf{Solution.}
For any $2\times2$ matrix $\begin{bmatrix}
\vec{b}_1 & \vec{b}_2
\end{bmatrix}$, multiplying gives \[BA = \begin{bmatrix}
B \bbvec{6}{7} & B \bbvec{0}{0}
\end{bmatrix} = \begin{bmatrix}
B \bbvec{6}{7} & \bbvec{0}{0}
\end{bmatrix} \ne \begin{bmatrix}
1 & 0 \\
0 & 1
\end{bmatrix}\]
\blacksquare
\end{book}

\example Let $A = \begin{bmatrix}
1 & 3 & 4 \\
-1 & -4 & -2 \\
0 & 1 & -1
\end{bmatrix}$. Find $A^{-1}$ if possible. \hintfootnote{Denote $A^{-1} = \begin{bmatrix}
\vec{x}_1 & \vec{x}_2 & \vec{x}_3
\end{bmatrix}$, and think about $AA^{-1} = I$.}

\begin{notes}
\columnbreak
\end{notes}

\begin{book}
\textbf{Solution.} If $A^{-1}$ exists, then:
\[AA^{-1} = \begin{bmatrix}
A\vec{x}_1 & A\vec{x}_2 & A\vec{x}_3
\end{bmatrix} = \begin{bmatrix}
\vec{e}_1 & \vec{e}_2 & \vec{e}_3
\end{bmatrix}\]
We could solve the three equations $A\vec{x}_1=\vec{e}_1, A\vec{x}_2=\vec{e}_2, A\vec{x}_3= \vec{e}_3$ by row-reducing the augmented matrices $\begin{bmatrix}[c|c]
A & \vec{e}_{1}
\end{bmatrix}, \begin{bmatrix}[c|c]
A & \vec{e}_2
\end{bmatrix}, \begin{bmatrix}[c|c]
A & \vec{e}_3
\end{bmatrix}$ but we'd be doing the same row operations three times. So let's just reduce one big $3\times6$ augmented matrix $\begin{bmatrix}[c|ccc]
A &\vec{e}_1 & \vec{e}_2 & \vec{e}_3
\end{bmatrix}$:
\begin{align*}
\begin{bmatrix}[ccc|ccc]
1 & 3 & 4 & 1 & 0 & 0\\
-1 & -4 & -2&0& 1 & 0 \\
0 & 1 & -1 &0 & 0 & 1
\end{bmatrix}
&\sim \begin{bmatrix}[ccc|ccc]
1 & 3 & 4 & 1 & 0 & 0\\
0 & -1 & 2&1& 1 & 0 \\
0 & 1 & -1 &0 & 0 & 1
\end{bmatrix} &r_2 \leftarrow r_2+r_1  \\
&\sim \begin{bmatrix}[ccc|ccc]
1 & 3 & 4 & 1 & 0 & 0\\
0 & -1 & 2&1& 1 & 0 \\
0 & 0 & 1 &1 & 1 & 1
\end{bmatrix} &\begin{matrix}
\\ \\ r_3 \leftarrow r_3+r_2
\end{matrix} \\
&\sim \begin{bmatrix}[ccc|ccc]
1 & 3 & 0 & -3 & -4 & -4\\
0 & -1 & 0&-1& -1 & -2 \\
0 & 0 & 1 &1 & 1 & 1
\end{bmatrix} &\begin{matrix}
r_1\leftarrow r_1-4r_3\\ r_2\leftarrow r_2-2r_3\\ \hspace{2pt}
\end{matrix} \\
&\sim \begin{bmatrix}[ccc|ccc]
1 & 0 & 0 & -6 & -7 & -10\\
0 & -1 & 0&-1& -1 & -2 \\
0 & 0 & 1 &1 & 1 & 1
\end{bmatrix} &\begin{matrix}
r_1\leftarrow r_1+3r_2\\\hspace{2pt} \\\hspace{2pt}
\end{matrix} \\
&\sim \begin{bmatrix}[ccc|ccc]
1 & 0 & 0 & -6 & -7 & -4\\
0 & 1 & 0&1& 1 & 2 \\
0 & 0 & 1 &1 & 1 & 1
\end{bmatrix} &\begin{matrix}
\hspace{2pt} \\ r_2\leftarrow-r_2 \\ \hspace{2pt}
\end{matrix}
\end{align*}
This shows three solutions of the three separate equations: $\vec{x}_1=(-6,1,1)$  is a solution of $A\vec{x}_1=\vec{e}_1$, $\vec{x}_2= (-7,1,1)$ is a solution of $A\vec{x}_2=\vec{e}_2$, and $\vec{x}_3 = (-4,2,1)$ is a solution of $A\vec{x}_3=\vec{e}_3$. Thus $A \begin{bmatrix}
\vec{x}_1 & \vec{x}_2 & \vec{x}_3
\end{bmatrix} = I$, and it's left to the reader to check that $\begin{bmatrix}
\vec{x}_1 & \vec{x}_2 & \vec{x}_3
\end{bmatrix}A = I$ as well. \blacksquare
\end{book}

The following theorem generalizes the idea of the last example.

\begin{theorem}
\label{theorem:AIequivIA-inv}
If $A$ is a $c\times c$ invertible matrix, then the $c\times2c$ matrix $\begin{bmatrix}[c|c]
A & I
\end{bmatrix}$ is row-equivalent to $\begin{bmatrix}[c|c]
I & A^{-1}
\end{bmatrix}$
\end{theorem}

\begin{notes}
A proof is in the textbook.
\end{notes}

\begin{book}
\textbf{Proof.} Suppose $A$ is invertible. The columns of $A$ are linearly independent, because if $\vec{x}$ is a solution of $A\vec{x}= \vec{0}$, then $A^{-1}A\vec{x}=A^{-1}\vec{0}$ and $\vec{x}=\vec{0}$. By theorem \ref{theorem:1.7uniqueness-v3}, $A$ has a pivot position in every column. The only place you can fit $c$ pivots in a $c\times c$ matrix is in the main diagonal, so the reduced row echelon form of $A$ is $I$.

Now that we've shown $A \sim I$, we need to show the same sequence of row operations that reduces $A$ to $I$ also reduces $I$ to $A^{-1}$. Consider the equations
\begin{equation}
\label{eq:2.3-1}
A\vec{x_1}= \vec{e}_1, \ldots, A\vec{x}_c=\vec{e}_c
\end{equation}
which have augmented matrices:
\[ \begin{bmatrix}[c|c]
A & \vec{e}_1
\end{bmatrix}, \ldots, \begin{bmatrix}[c|c]
A & \vec{e}_c
\end{bmatrix}\]
Since $A \sim I$, each of these can be row reduced to
\[ \begin{bmatrix}[c|c]
I & \vec{w}_1
\end{bmatrix},\ldots, \begin{bmatrix}[c|c]
I & \vec{w}_c
\end{bmatrix}\]
which are the augmented matrices for the equations:
\[I\vec{x}_1=\vec{w}_1,\ldots, I\vec{x}_c=\vec{w}_c\]
So row reducing $A$ to $I$ also reduced each vector $\vec{e}_j$ to $\vec{w}_j$. If we write $W = \begin{bmatrix}
\vec{w}_1 & \cdots & \vec{w}_c
\end{bmatrix}$, the same row operations reduced $I$ to $W$. Since $\vec{w}_1,\ldots,\vec{w}_c$ are solutions of equations \ref{eq:2.3-1}, we have:
\[AW = \begin{bmatrix}
A\vec{w}_1 & \cdots & A\vec{w}_c
\end{bmatrix} = \begin{bmatrix}
\vec{e}_1 & \cdots & \vec{e}_c
\end{bmatrix} = I\]
It sure looks like $W = A^{-1}$ but we should be careful because we haven't shown $WA = I$ and matrix multiplication isn't generally commutative. But multiplying both sides of $AW = I$ by $A^{-1}$ gives:
\begin{align*}
A^{-1}AW &= A^{-1}I\\
IW &= A^{-1}\\
W &= A^{-1} \quad \blacksquare
\end{align*}
\end{book}

\example If $A = \begin{bmatrix}
1 & 2 \\
-2 & -4
\end{bmatrix}$, find $A^{-1}$ if possible.

\begin{book}
\textbf{Solution.} Reducing $\begin{bmatrix}[c|c]
A & I
\end{bmatrix}$:
\[\begin{bmatrix}[cc|cc]
1 & 2 & 1&0 \\
-2 & -4& 0&1
\end{bmatrix} \sim \begin{bmatrix}
1 & 2 & 1 & 0 \\
0 & 0 & 2 & 1
\end{bmatrix} \sim \begin{bmatrix}
1 & 2 & 0 & -0.5 \\
0 & 0 & 1 & 0.5
\end{bmatrix} \]
Since the reduced echelon form of any matrix is unique, this shows that $\begin{bmatrix}[c|c]
A & I
\end{bmatrix}$ is not row-equivalent to $\begin{bmatrix}[c|c]
I & A^{-1}
\end{bmatrix}$. So theorem \ref{theorem:AIequivIA-inv} shows $A$ \emph{can't} be invertible. \blacksquare
\end{book}
\begin{notes}
\vfill
\end{notes}

\bbtwocol{.7}{.25}{
Recall that the \emph{contrapositive} of a statement in the form ``If $P$, then $Q$'' is ``If not $Q$, then not $P$.'' In general, a statement and its contrapositive are equivalent, as the Venn diagram suggests. The following corollary is essentially the contrapositive of theorem \ref{theorem:AIequivIA-inv}.
}{
\begin{tikzpicture}[scale=0.8]
\draw[thick] (-2.5,-2.1) rectangle (2.1,2.1);
\draw[thick] (0,0) circle (1.8cm);
\node at (-1.5, 1.5) {$Q$};
\draw[thick, fill=lightgray!30] (0.5,0) circle (1cm);
\node at (0.5, 0) {$P$};
\end{tikzpicture}}

\begin{bbox}
\textbf{Corollary:} If $A$ is not row-equivalent to $I$, then $A$ is not invertible.
\end{bbox}

In the section \ref{sec:elementary-matrices-and-IMT} we'll prove the converse of theorem \ref{theorem:AIequivIA-inv}: that if $A$ is row-equivalent to $I$, then $A$ is invertible; and the Venn diagram looks like just one circle.

\begin{notes}
\columnbreak
\end{notes}

\begin{theorem}[$\mathbf{2 \times 2}$ inverses.]\label{theorem:2x2-inverses}
Let $A = \begin{bmatrix}
a & b \\
c & d
\end{bmatrix}$. Then:
\begin{enumerate}[label=(\roman*)]
\item $A$ is invertible if and only if $ad-bc \ne 0$. 
\item If $A$ is invertible, then 
$ A^{-1}=\frac{1}{ad-bc}\begin{bmatrix}
d & -b \\
-c & a
\end{bmatrix}$
\end{enumerate}
\end{theorem}

\begin{book}
\textbf{Proof.} For part(i) (\rightarrow), suppose $A$ is invertible. We want to show that $ad-bc \ne 0$ and the formula in part (ii) works. Assume $a\ne0$; exercise \ref{exercise:2x2-inverse-proof} asks you to prove it for the $a=0$ case. Let's row reduce $\begin{bmatrix}[c|c]
A & I
\end{bmatrix}$:
\begin{align*}
\begin{bmatrix}[cc|cc]
a & b & 1 & 0 \\
c & d & 0 & 1
\end{bmatrix} &\sim \begin{bmatrix}[cc|cc]
a & b & 1 & 0 \\
ac & ad & 0 & a
\end{bmatrix} &\begin{matrix}
\hspace{2pt} \\
r_2\leftarrow ar_2
\end{matrix} \\
 &\sim \begin{bmatrix}[cc|cc]
a & b & 1 & 0 \\
0 & ad-bc & -c & a
\end{bmatrix} &\begin{matrix}
\hspace{2pt} \\
r_2\leftarrow r_2 -cr_{1}
\end{matrix}
\end{align*}
At this stage, we can conclude $ad-bc\ne0$ because theorem \ref{theorem:AIequivIA-inv} implies $A\sim I$. So we're done with part (i)(\rightarrow). For (ii), we could keep reducing until we get to $\begin{bmatrix}[c|c]
I & A^{-1}
\end{bmatrix}$, but it would get messy. Instead, now that we know the denominator isn't zero, we can just multiply to check that:  
\begin{equation}\label{eq:AA-inv}
\left(\frac{1}{ad-bc}\begin{bmatrix}
d & -b \\
-c & a
\end{bmatrix}\right)\begin{bmatrix}
a & b \\
c & d
\end{bmatrix} = I \quad \text{ and } \quad \begin{bmatrix}
a & b \\
c & d
\end{bmatrix}\left(\frac{1}{ad-bc}\begin{bmatrix}
d & -b \\
-c & a
\end{bmatrix}\right) = I
\end{equation}
And, this also proves the (\leftarrow) direction of part (i); that if $ad-bc\ne0$, then $A$ is invertible, because we just found its inverse. \blacksquare
\end{book}

\begin{notes}
\columnbreak
\end{notes}

Let's finish up with some nice comfortable theorems about inverses.

\begin{theorem}[Properties of inverses.]
\label{theorem:properties-of-inverses}

Suppose $A$ and $B$ are invertible matrices. Then:
\begin{enumerate}[label=(\roman*)]
\item $A^{-1}$ is invertible and $( A^{-1} )^{-1}=A$.
\item $A^T$ is invertible and $(A^T)^{-1} = ( A^{-1} )^T$.
\item $AB$ is invertible and $(AB)^{-1}=B^{-1}A^{-1}$.
\end{enumerate}
\end{theorem}

\begin{notes}
\vfill
\end{notes}

\begin{book}
\textbf{Proof.} (i) and (ii) are left as exercises. Robert Beezer, in his linear algebra book, calls  part (iii) the ``socks and shoes'' theorem. Because, if you think of $I$ as your bare feet, $BI=B$ as your feet with socks on, and $AB$ as your feet with socks and shoes on, then how do you take your socks and shoes off? In reverse order, by taking shoes off, then socks off. Anyway, let's do the math:
\[ (  B^{-1}A^{-1})(AB) = B^{-1}( A^{-1}A )B = B^{-1}IB = B^{-1}B = I\]
and likewise,
\[ (AB)(B^{-1}A^{-1}) = A ( BB^{-1} )A^{-1}=AIA^{-1}=AA^{-1}=I\]
This shows that $B^{-1}A^{-1}$ is the inverse of $AB$. \blacksquare
\end{book}

\begin{theorem}
\label{theorem:2.3invertible-unique-soln}
If $A$ is an invertible $c\times c$ matrix, then for any $\vec{y}$ in $\R^c$, the equation $A\vec{x}=\vec{y}$ has the unique solution $\vec{x}=A^{-1}\vec{y}$.
\end{theorem}

\begin{notes}
\vfill \hspace{2pt}
\end{notes}

\begin{book}
\textbf{Proof.} We can see $\vec{x}=A^{-1}\vec{y}$ is a solution of the equation by substituting it:
\[ A ( A^{-1}\vec{y})=(AA^{-1})\vec{y}=\vec{y}\]
To see it's the unique solution, suppose there's another solution $\vec{w}$ such that $A\vec{w}=\vec{y}$. Then multiplying both sides by $A^{-1}$ gives:
\begin{align*}
A^{-1}A\vec{w} &= A^{-1}\vec{y} \\
I\vec{w} &= A^{-1}\vec{y} \\
\vec{w} &= A^{-1}\vec{y} 
\end{align*} 
So, any solution of the equation is actually $A^{-1}\vec{y}$ and the solution is unique. \blacksquare
\end{book}

\columnbreak
\exerciseheading

\exatend Find the inverse of $\begin{bmatrix}
1 & 7 & 0 \\
0 & 1 & 0 \\
0 & 0 & 2
\end{bmatrix}$.

\exatend Find the inverse of $\begin{bmatrix}
1 & 0 & 0 \\
2 & 1 & 0 \\
-4 & 2 & 1
\end{bmatrix}$

%TODO: A little more with transformation notation.

\exatend Give your own example of a non-invertible $2\times 2$ matrix $A$. Justify your answer.

\exatend Give your own example of a non-invertible $3\times 3$ matrix $A$. Justify your answer.

\exatend Check that $A = \begin{bmatrix}
0 & 1 \\
1 & 0
\end{bmatrix}$ is its own inverse.

\exatend (T/F?) If $A$ and $B$ are invertible, then $(AB)^{-1}=A^{-1}B^{-1}$.

\exatend (T/F?) We have proven in this section that if $A$ is row equivalent to $I$, then $A$ is invertible.

\exatend (T/F?) If $A$ has two identical rows, then $A$ is not invertible.

\exatend (T/F?) If $A$ is invertible, and $T(\vec{x})=A\vec{x}$, then the range of $T$ is equal to its codomain.

\exatend (T/F?) There exist matrices $A,P,Q$ such that $AP = PA = I$, $AQ = QA = I$, and $P \ne Q$.

\exatend If $A$ is an invertible $c \times c$ matrix, what is $\rank A$?

\exatend If $A$ is an invertible $c \times c$ matrix, then $\dim ( \nul A )$?

\exatend \label{exercise:2.3diag}Find the inverse of $A = \begin{bmatrix}
-4 & 0 & 0 \\
0 & 3 & 0 \\
0 & 0 & 2
\end{bmatrix}$ and check by multiplying.

\exatend The matrix in \ref{exercise:2.3diag} is an example of a \emph{diagonal} matrix, where $a_{ij}=0$ if $i\ne j$.
\begin{enumerate}[label=(\alph*)]
\item How can you quickly determine if a diagonal matrix is invertible or not?
\item Explain why the inverse of an invertible diagonal matrix is also diagonal.
\end{enumerate}

\exatend Let $A = \begin{bmatrix}
1 & 2 \\
3 & 4
\end{bmatrix}$.
\begin{enumerate}[label=(\alph*)]
\item Find $A^{-1}$ by row-reducing $\begin{bmatrix}[c|c]
A & I
\end{bmatrix}$ to $\begin{bmatrix}
I & A^{-1}
\end{bmatrix}$.
\item Find $A^{-1}$ using the formula in theorem \ref{theorem:2x2-inverses}.
\item Use theorem \ref{theorem:2.3invertible-unique-soln} to solve $A\vec{x}=\bbvec{1}{7}$.
\item Use theorem \ref{theorem:2.3invertible-unique-soln} to write the solution to $A\vec{x} = \bbvec{y_1}{y_2}$ in terms of $y_1$ and $y_2$.
\end{enumerate}

\exatend Let $A = \begin{bmatrix}
1 & -2 & 0 \\
2 & -3 & 2 \\
1 & -2 & 1
\end{bmatrix}$.
\begin{enumerate}[label=(\alph*)]
\item Find $A^{-1}$ by row-reducing $\begin{bmatrix}[c|c]
A & I
\end{bmatrix}$ to $\begin{bmatrix}
I & A^{-1}
\end{bmatrix}$.
\item Use theorem \ref{theorem:2.3invertible-unique-soln} to write the solution to $A\vec{x}=\bbbvec{p}{q}{r}$ in terms of $p,q$ and $r$.
\end{enumerate}

\exatend Define a transformation $T:\R^5\rightarrow \R^5$ by $T(x_1,x_2,x_3,x_4,x_5) = (x_5,x_1,x_2,x_3,x_4)$. This is an example of a ``permutation.''
\begin{enumerate}[label=(\alph*)]
\item Find $A$, the standard matrix for $T$.
\item Find $A^{-1}$ and check by multiplying.
\end{enumerate}

\exatend Suppose $B = \begin{bmatrix}
\vec{b}_1 & \cdots & \vec{b}_c
\end{bmatrix}$ is an invertible $c \times c$ matrix. 
\begin{enumerate}[label=(\alph*)]
\item Show that $\mathcal{B} = ( \vec{b}_1, \ldots, \vec{b}_c )$ is a basis of $\R^c$.
\item Show that $\left[ \vec{y} \right]_{\mathcal{B}}=B^{-1}\vec{y}$ for any $\vec{y}$ in $\R^c$.
\end{enumerate}

\exatend If $P$ is invertible, and $B = PAP^{-1}$, solve for $A$. Can you conclude that $A = B$?

\exatend Prove or give a counterexample: If $A$ and $B$ are invertible, then $A+B$ is invertible.

\exatend Prove or give a counterexample: If $A$ is invertible, and $AB = BC$, then $B=C$.

\exatend Prove or give a counterexample: If $A$ has two identical columns, then $A$ is not invertible.

\exatend Let $A = \begin{bmatrix}
1 & 2 & 3
\end{bmatrix}$. Describe all vectors $\vec{x} = \bbbvec{x_1}{x_2}{x_3}$ such that $A\vec{x}= \begin{bmatrix}
1
\end{bmatrix}$, aka $I_1$.

\exatend Give an example of a $2\times 3$ matrix $A$ and $3\times 2$ matrix $B$ such that $AB = I$. Is $BA = I$?

\exatend Suppose $A$ is a $c\times c$ matrix, and $(  \vec{b}_1, \ldots, \vec{b}_c)$ is a basis of $\R^c$. Show that $( A\vec{b}_1,\ldots,A\vec{b}_c )$ is a basis of $\R^c$.

\exatend \label{exercise:2x2-inverse-proof} Prove the following claim that was left out of the proof of theorem \ref{theorem:2x2-inverses}: If $A = \begin{bmatrix}
a & b \\
c & d
\end{bmatrix}$ is invertible, and $a \ne 0$, then $ad-bc\ne0$. \hintfootnote{Show that $b \ne 0$ and $c\ne0$.}

\exatend \label{exercise:2x2-inverse-proof-2} Perform the two multiplications below from the proof of theorem \ref{theorem:2x2-inverses} to check that you get $I$ both ways.
\[
\left(\frac{1}{ad-bc}\begin{bmatrix}
d & -b \\
-c & a
\end{bmatrix}\right)\begin{bmatrix}
a & b \\
c & d
\end{bmatrix} = I \quad \text{ and } \quad \begin{bmatrix}
a & b \\
c & d
\end{bmatrix}\left(\frac{1}{ad-bc}\begin{bmatrix}
d & -b \\
-c & a
\end{bmatrix}\right) = I
\]

\exatend Prove theorem \ref{theorem:properties-of-inverses} part (i).

\exatend Prove theorem \ref{theorem:properties-of-inverses} part (ii).

\exatend Prove that if $A^2 = 0$, then $A$ is not invertible.

\exatend Show that if $A$ is invertible and \emph{symmetric}, which means $A^{T}=A$, then $A^{-1}$ is also symmetric.

\exatend \label{exercise:PA-implies-independence} Prove theorem \ref{theorem:PA-implies-independence} below. \hintfootnote{The goal is to prove that if $A\vec{x}=\vec{0}$, then $\vec{x}=\vec{0}$.}

\begin{theorem}
\label{theorem:PA-implies-independence}
If $A$ is an $r \times c$ matrix, and there exists a matrix $P$ such that $PA = I$, then the columns of $A$ are a linearly independent list of vectors.
\end{theorem}

\exatend \label{exercise:AP-implies-span} Prove theorem \ref{theorem:AQ-implies-spanning} below. \hintfootnote{Let $\vec{y}$ be in $\R^{r}$. What can you do with $Q\vec{y}$?}

\begin{theorem}\label{theorem:AQ-implies-spanning}
If $A$ is an $r \times c$ matrix, and there exists a matrix $Q$ such that $AQ = I$, then for any $\vec{y}$ in $\R^r$, the equation $A\vec{x}=\vec{y}$ has a solution.
\end{theorem}

\exatend Suppose $A$ is $r\times c$, $F$ is $c \times r$, and $AF = FA = I$. The following parts outline a proof that $A$ must be square.
\begin{enumerate}[label=(\alph*)]
\item What does theorem \ref{theorem:PA-implies-independence} say about the pivot positions in $A$? What does this say about $r$ and $c$?
\item What does theorem \ref{theorem:AQ-implies-spanning} say about the pivot positions in $A$? What does this say about $r$ and $c$?
\end{enumerate}


\columnbreak
\section{Elementary Matrices and the Invertible Matrix Theorem}
\label{sec:elementary-matrices-and-IMT}

\exercise (\emph{Warm-up}) Review theorem \ref{theorem:buy-two-get-one-free} (buy two, get one free) and answer these true/false questions.
\begin{enumerate}[label=(\alph*)]
\item (T/F?) If $A = \begin{bmatrix}
\vec{a}_1 & \cdots & \vec{a}_c
\end{bmatrix}$ is a square matrix with linearly independent columns, then for every $\vec{y}$ in $\R^{c}$, the equation $A\vec{x}=\vec{y}$ has a solution.
\item (T/F?) If $A = \begin{bmatrix}
\vec{a}_1 & \cdots & \vec{a}_c
\end{bmatrix}$ is a square matrix such that for every $\vec{y}$ in $\R^{c}$, the equation $A\vec{x}=\vec{y}$ has a solution, then $A$ has linearly independent columns.
\item Is part (a) true if $A$ is not necessarily square?
\item Is part (b) true if $A$ is not square?
\end{enumerate}

\exercise (\emph{Warm-up}) Let $E=\begin{bmatrix}
1 & -3 \\
0 & 1
\end{bmatrix}$ and $F=\begin{bmatrix}
1 & 3 \\
0 & 1
\end{bmatrix}$. Check that $EF = FE = I$. What is the row operation that transforms $I$ to $E$? Can you transform $E$ back to $I$ with a different row operation?

\begin{notes}
\vfill
\end{notes}

\begin{definition}
An \textbf{elementary matrix} is the result of applying a \emph{single} row operation (swapping, scaling, or multaddiplacement) to an identity matrix.
\end{definition}

\example Here are some elementary matrices:

\bbthreecol{0.3}{0.3}{0.3}{\[
\begin{bmatrix}
1 & 0 \\
\pi & 1
\end{bmatrix} \]
\[\begin{bmatrix}
1 & -7 & 0 \\
0 & 1 & 0 \\
0 & 0 & 1
\end{bmatrix} \]
\[\begin{bmatrix}
1 & 0 & 0 & 0 \\
0 & 1 & 0 & -2 \\
0 & 0 & 1 & 0 \\
0 & 0 & 0 & 1
\end{bmatrix} \]}{ \[ \begin{bmatrix}
-3 & 0 \\
0 & 1
\end{bmatrix}\]
\[\begin{bmatrix}
1 & 0 & 0 \\
0 & 1 & 0 \\
0 & 0 & 5
\end{bmatrix} \]
\[\begin{bmatrix}
1 & 0 & 0 & 0 \\
0 & \sqrt{2} & 0 & 0 \\
0 & 0 & 1 & 0 \\
0 & 0 & 0 & 1
\end{bmatrix} \]}{\[\begin{bmatrix}
0 & 1 \\
1 & 0
\end{bmatrix}\]
\[\begin{bmatrix}
1 & 0 & 0 \\
0 & 0 & 1 \\
0 & 1 & 0 
\end{bmatrix} \]
\[\begin{bmatrix}
1 & 0 & 0 & 0 \\
0 & 1 & 0 & 0 \\
0 & 0 & 0 & 1 \\
0 & 0 & 1 & 0 
\end{bmatrix} \]}

\begin{notes}
\columnbreak
\end{notes}
We'll soon see that multiplication by an elementary matrix \emph{does} the corresponding row operation, but we need a little more machinery to prove it first.
Recall the definition of $AB = \begin{bmatrix}
A\vec{b}_1 & \cdots & A\vec{b}_c
\end{bmatrix}$ shows that column $j$ of $AB$ is a linear combination of the columns of $A$, using the entries from $\vec{b}_j$ as weights. The next theorem shows an analogous result about rows.

\begin{theorem}[Rows of $BA$.]\label{theorem:rows-of-BA}

Suppose $B$ is a $d \times r$ matrix, and A is an $r \times c$ matrix. Write the rows $B$ as column vectors $\vec{\beta}_1, \ldots \vec{\beta}_d$, so that $B = \bbbvec{\vec{\beta}_1^T}{\vdots}{\vec{\beta}_d^T}$. Then row $i$ of $BA$ is $\vec{\beta}_i^TA$ which is a linear combination of the rows of $A$ using the entries in $\vec{\beta}_i$ as weights.
\end{theorem}
% prove this in class too!

\begin{book}
\textbf{Proof.} Row $i$ of $BA$ is column $j$ of $(BA)^T$, which can be written:
\begin{align*}
 (BA)^T &= A^TB^T  & \text{(thm \ref{theorem:ABT-BTAT})} \\
&= A^T\begin{bmatrix}
\vec{\beta}_1 & \cdots & \vec{\beta}_d
\end{bmatrix} \\
&= \begin{bmatrix}
A^T\vec{\beta}_1 & \cdots & A^T\vec{\beta}_d
\end{bmatrix}
\end{align*}
Column $j$ of this matrix, $A^T\vec{\beta}_j$, is a linear combination of the columns of $A^T$, \emph{which are rows of $A$}, using entries in $\vec{\beta}_j$ as weights. Also, transposing this column of $(BA)^T$ back into a row of $BA$ gives: 
\[(A^T\vec{\beta}_j)^T = \vec{\beta}_j^T(A^{T})^T = \vec{\beta}_j^TA \]
\blacksquare
\end{book}

\begin{notes}
\vfill \hspace{2pt}
\columnbreak
\end{notes}
\example Let $E=\begin{bmatrix}
1 & 0 \\
\pi & 1
\end{bmatrix}$, $F=\begin{bmatrix}
-3 & 0 \\
0 & 1
\end{bmatrix}$, and  $A = \begin{bmatrix}
a & b \\
c & d
\end{bmatrix}$.
\begin{enumerate}[label=(\alph*)]
\item Multiply $EA$ and $FA$ using theorem \ref{theorem:rows-of-BA}. What do you notice?
\item Find $F^{-1}$.
\end{enumerate}

\begin{book}
\textbf{Solution.} (a) Write the rows of $E$ as vectors: $\vec{\epsilon}_1 = \bbvec{1}{0}, \vec{\epsilon}_2= \bbvec{\pi}{1}$ so that $E = \bbvec{\vec{\epsilon}_1^T}{\vec{\epsilon}_2^T}$. Using theorem \ref{theorem:rows-of-BA}:
\[
EA = \bbvec{\vec{\epsilon}_1^TA}{\vec{\epsilon}_2^TA} 
  = \bbvec{\begin{bmatrix}
1 & 0
\end{bmatrix}\begin{bmatrix}
a & b \\
c & d
\end{bmatrix}}{\begin{bmatrix}
\pi & 1
\end{bmatrix}\begin{bmatrix}
a & b \\
c & d
\end{bmatrix}}
= \bbvec{1\begin{bmatrix}
a & b
\end{bmatrix} + 0\begin{bmatrix}
c & d
\end{bmatrix}}{\pi\begin{bmatrix}
a & b
\end{bmatrix} + 1\begin{bmatrix}
c & d
\end{bmatrix}}
= \begin{bmatrix}
a & b \\
\pi a + c & \pi b + d
\end{bmatrix}
\]
Notice that $EA$ is the result of replacing row 2 of $A$ with the sum of row 2 and $\pi$ times row 1; the same row operation that was done to $I$ to get $E$. Likewise:
\[ FA
  = \bbvec{\begin{bmatrix}
-3 & 0
\end{bmatrix}\begin{bmatrix}
a & b \\
c & d
\end{bmatrix}}{\begin{bmatrix}
0 & 1
\end{bmatrix}\begin{bmatrix}
a & b \\
c & d
\end{bmatrix}}
= \bbvec{-3\begin{bmatrix}
a & b
\end{bmatrix} + 0\begin{bmatrix}
c & d
\end{bmatrix}}{0\begin{bmatrix}
a & b
\end{bmatrix} + 1\begin{bmatrix}
c & d
\end{bmatrix}}
= \begin{bmatrix}
-3a & -3b \\
 c &  d
\end{bmatrix}
\]
Again, $FA$ is the result of scaling row 1 of $A$ by $-3$; the same row operation that was done to $I$ to get $E$.

For (b), to make $FF^{-1}=I$, we want to scale row 1 of $F^{-1}$ by $-3$ to get $\begin{bmatrix}
1 & 0
\end{bmatrix}$. A good idea would be $F^{-1}=\begin{bmatrix}
\frac{-1}{3} & 0 \\
0 & 1
\end{bmatrix}$, if the multiplication $FF^{-1}$ works like it does in part (a). The reader should check that it does, and also that $F^{-1}F = I$. \blacksquare
\end{book}

\begin{notes}
\vfill \hspace{2pt}
\columnbreak
\end{notes}
\example Let $E = \begin{bmatrix}
1 & 0 & 0 \\
0 & 0 & 1 \\
0 & 1 & 0 
\end{bmatrix}$, $F = \begin{bmatrix}
1 & -7 & 0 \\
0 & 1 & 0 \\
0 & 0 & 1
\end{bmatrix}$, and $A = \begin{bmatrix}
a & b & c \\
d & e & f \\
g & h & i
\end{bmatrix}$.
\begin{enumerate}[label=(\alph*)]
\item Multiply $EA$ and $FA$ using theorem \ref{theorem:rows-of-BA}. What do you notice?
\item Find $E^{-1}$.
\end{enumerate}

\begin{book}
\textbf{Solution.} (a) Write the rows of $E$ as vectors: $\vec{\epsilon}_1 = \bbbvec{1}{0}{0}, \vec{\epsilon}_2=\bbbvec{0}{0}{1}, \vec{\epsilon}_3 = \bbbvec{0}{1}{0}$ so that $E = \bbbvec{\vec{\epsilon}_1^T}{\vec{\epsilon}_2^T}{\vec{\epsilon}_3^T}$. Using theorem \ref{theorem:rows-of-BA}:
\begin{align*}
EA = \bbbvec{\vec{\epsilon}_1^TA}{\vec{\epsilon}_2^TA}{\vec{\epsilon}_3^TA} 
  &= \bbbvec{\begin{bmatrix}
1 & 0 & 0
\end{bmatrix}\begin{bmatrix}
a & b & c \\
d & e & f \\
g & h & i
\end{bmatrix}}{\begin{bmatrix}
0 & 0 & 1
\end{bmatrix}\begin{bmatrix}
a & b & c \\
d & e & f \\
g & h & i
\end{bmatrix}}
{\begin{bmatrix}
0 & 1 & 0
\end{bmatrix}\begin{bmatrix}
a & b & c \\
d & e & f \\
h & h & i
\end{bmatrix}} \\
&= \bbbvec{1\begin{bmatrix}
a & b & c
\end{bmatrix} + 0\begin{bmatrix}
d & e & f
\end{bmatrix} + 0\begin{bmatrix}
g & h & i
\end{bmatrix}}{0\begin{bmatrix}
a & b & c
\end{bmatrix} + 0\begin{bmatrix}
d & e & f
\end{bmatrix} + 1\begin{bmatrix}
g & h & i
\end{bmatrix}}{0\begin{bmatrix}
a & b & c
\end{bmatrix} + 1\begin{bmatrix}
d & e & f
\end{bmatrix} + 0\begin{bmatrix}
g & h & i
\end{bmatrix}}
= \begin{bmatrix}
a & b & c \\
g & h & i \\
d & e & f
\end{bmatrix}
\end{align*}
Notice that $EA$ is the result of swapping rows 2 and 3 of $A$, the same row operation that turns $I$ into $E$. Likewise, using theorem \ref{theorem:rows-of-BA} with $FA$: 
\begin{align*}
FA = \bbbvec{\begin{bmatrix}
1 & -7 & 0
\end{bmatrix}\begin{bmatrix}
a & b & c \\
d & e & f \\
g & h & i
\end{bmatrix}}{\begin{bmatrix}
0 & 1 & 0
\end{bmatrix}\begin{bmatrix}
a & b & c \\
d & e & f \\
g & h & i
\end{bmatrix}}
{\begin{bmatrix}
0 & 0 & 1
\end{bmatrix}\begin{bmatrix}
a & b & c \\
d & e & f \\
h & h & i
\end{bmatrix}} 
&= \bbbvec{1\begin{bmatrix}
a & b & c
\end{bmatrix} + -7\begin{bmatrix}
d & e & f
\end{bmatrix} + 0\begin{bmatrix}
g & h & i
\end{bmatrix}}{0\begin{bmatrix}
a & b & c
\end{bmatrix} + 1\begin{bmatrix}
d & e & f
\end{bmatrix} + 0\begin{bmatrix}
g & h & i
\end{bmatrix}}{0\begin{bmatrix}
a & b & c
\end{bmatrix} + 0\begin{bmatrix}
d & e & f
\end{bmatrix} + 1\begin{bmatrix}
g & h & i
\end{bmatrix}}\\
&= \begin{bmatrix}
a-7d & b-7e & c-7f \\
d & e & f \\
g & h & i 
\end{bmatrix}
\end{align*}
Again, notice that $FA$ is the result of adding $-7$ times row 2 of $A$ to row $1$, the same multaddiplacement that turns $I$ to $F$.

For (b), we want to make $E^{-1}E = I$, with $E^{-1}$ on the left doing a row operation to $E$. The row operation that should swap rows 2 and 3 back is just the same swap. So $E^{-1} = E$ should work, and the reader should check that $EE = I$. \blacksquare
\end{book}

\begin{notes}
\vfill \hspace{2pt}
\columnbreak
\end{notes}

\textbf{Note:} The above examples should convince you that every elementary matrix $E$ has an inverse, which is defined by the inverse row operation of the one that was done to $I$ to make $E$.

The above examples also justify the following theorem. They're not formal proof, but there's nothing special about the particular row operations in the examples, and the concept works the same way in general.

\begin{theorem}[Elementary matrices do row operations.]\label{theorem:elementary-operations}

If $A$ is any $r\times c$ matrix, and $E$ is an $r\times r$ elementary matrix that is the result of doing some row operation to $I$, then $EA$ is the result of doing the same operation to $A$.
\end{theorem}

\begin{book}
\textbf{Proof.} The examples above are a good way to convince yourself this is true. If you're not convinced, maybe try similar examples with some $4x4$ elementary matrices. If you're still not convinced, here's a general proof.

First some notation. Notice that the rows of an identity matrix are equal to its columns, so we can write $I_r = \bbbvec{\vec{e}_1^T}{\vdots}{\vec{e}_r}$. And we'll write the rows of $E$ as column vectors $\vec{\epsilon}_1, \ldots, \vec{\epsilon}_r$, so that $E = \bbbvec{\vec{\epsilon}_1^T}{\vdots}{\vec{\epsilon}_r^T}$. Also, we'll refer to the row operation that turns $I$ into $E$ as simply ``the operation.''  

Now, if row $i$ is unaffected by the operation, then row $i$ in $E$ is still $\vec{\epsilon}_i^T=\vec{e}_i^T$. And by theorem \ref{theorem:rows-of-BA} (rows of BA), row $i$ of $EA$ is $\vec{e}_i^TA$ which is a linear combination of the rows of $A$ using $\vec{e}_i=(0,\ldots, 0, 1, 0, \ldots, 0)$ as weights:
\[\text{row }i\text{ of } EA = \vec{e}_i^TA = \text{row }i\text{ of }A\]
Thus the rows of $A$ that should be unchanged by the operation are indeed unchanged. What about the rows of $E$ that were changed by the operation? The gist is that these rows of $E$ make the corresponding rows of $EA$ equal to the correct linear combination of the rows of $A$. Let's check the three different types of row operations:

If the operation is scaling a row by $k$, then there's a single row $\vec{\epsilon}_i^T = \begin{bmatrix}
0 & \cdots & 0 & k & 0 & \cdots & 0
\end{bmatrix}$ with the $k$ in the $i^{\textrm{th}}$ entry. Using theorem \ref{theorem:rows-of-BA} again,
\[\text{row }i \text{ of } EA =  \vec{\epsilon}_i^TA = \begin{bmatrix}
0 & \cdots & 0 & k & 0 & \cdots & 0
\end{bmatrix}A = k(\text{row } i \text{ of } A),\]
as desired.

If the operation is a swap, then there are two rows of $E$ of the form $\vec{\epsilon}^T_i=\vec{e}^T_j$ and $\vec{\epsilon}_j^T= \vec{e}^T_i$. Using theorem \ref{theorem:rows-of-BA} again,
\[ \text{row }i \text{ of } EA =  \vec{\epsilon}_i^TA = \vec{e}_j^TA = \text{row }j \text{ of } A\]
and 
\[ \text{row }j \text{ of } EA =  \vec{\epsilon}_j^TA = \vec{e}_e^TA = \text{row }i \text{ of } A\]
which means the corresponding two rows of $EA$ are also the swapped rows of $A$ as desired.

If the operation is multaddiplacement, there's a row $i$ of $E$ that can be written $\vec{\epsilon}_i^T = \vec{e}_i^T + k\vec{e}_j^T$. Using theorem \ref{theorem:rows-of-BA} one final time,
\[ \text{row }i \text{ of } EA =  \vec{\epsilon}_i^TA = (\vec{e}_i^T + k\vec{e}_j^T)A =
\vec{e}_i^TA + k(\vec{e}_j^TA) = \text{row }i \text{ of} A + k ( \text{row }j\text{ of }A )\]
as desired. \blacksquare
\end{book}

The next theorem combines a bunch of things we already know with one new fact to prove. Combined, it's a powerful and easy-to-use tool.

\begin{theorem}[Invertible Matrix Theorem (IMT).]
\label{theorem:IMT}

Suppose $A$ is a $c\times c$ matrix. Then the following are equivalent:
\begin{enumerate}[label=(\roman*)]
\item $A$ is invertible.
\item There exists a matrix $P$ such that $PA = I_c$.
\item There exists a matrix $Q$ such that $AQ = I_c$.
\item $A$ is row-equivalent to $I_c$.
\item $A$ has $c$ pivot positions.
\item The columns of $A$ are linearly independent.
\item $\nul A = \left\{ \vec{0} \right\}$.
\item For every $\vec{y}$ in $\R^{c}$, $A\vec{x}=\vec{y}$ has a solution.
\item $\col A = \R^c $
\end{enumerate}
\end{theorem}

\begin{notes}
\columnbreak
\end{notes}
\textbf{Proof.} 

\begin{center}
\begin{tikzpicture}[
node distance=2cm,
every node/.style={minimum width=1cm},
arrow/.style={->, >=Stealth, thick},
doublearrow/.style={<->, >=Stealth, thick}]
\node (i) at (0,0) {(i)};
\node (iv) [right=1.5cm of i] {(iv)};
\node (v) [right=1.5cm of iv] {(v)};
\node (ii) [below right=1.2cm and -.4cm of i] {(ii)};
\node (vi) [below left=1.2cm and 0.2cm of v] {(vi)};
\node (vii) [right=0.8cm of vi] {(vii)};
\node (viii) [below right=1cm and 2cm of v] {(viii)};
\node (ix) [right=0.8cmof viii] {(ix)};
\node (iii) [below=2cm of i] {(iii)};
\draw[arrow] (i) to[bend left=15] (iv);
\draw[arrow,red] (iv) to[bend left=15] (i);
\draw[doublearrow] (iv) -- (v);
\draw[doublearrow] (v) -- (viii);
\draw[doublearrow] (v) -- (vi);
\draw[arrow] (i) -- (ii);
\draw[arrow] (ii) -- (vi);
\draw[doublearrow] (vi) -- (vii);
\draw[doublearrow] (viii) -- (ix);
\draw[arrow] (i) to[bend right=25] (iii);
\draw[arrow] (iii) to[bend right=15] (viii);
\end{tikzpicture}
\end{center}

In the graph, a black arrow indicates an implication that we either already know, or is easy to prove. The red arrow indicates the implication we need to prove. Notice that if you know that any one of the parts is true, you can follow arrows that imply any other part is true.
\begin{book}
For example, if $\col A = \R^c$, part (ix), is true, that implies (viii) which implies (v) which implies (i) which implies (iii). Likewise if you know that (iii) is false, then (ix) can't be true because (ix) implies (iii). The same goes for any other pair of parts.

Let's prove the easy ones first.

(i \rightarrow\hspace{1pt} iv) Theorem \ref{theorem:AIequivIA-inv} says that if $A$ is invertible, then $\begin{bmatrix}[c|c]
A & I_c
\end{bmatrix}$ is row equivalent to $\begin{bmatrix}[c|c]
I_c & A^{-1}
\end{bmatrix}$ which means $A$ is row-equivalent to $I_c$.

(iv \rightarrow\hspace{1pt} v) The identity matrix has a leading entry 1 in each of the $c$ columns.

(v \leftarrow\hspace{1pt} iv) You can't have more than one pivot position in a single row or column. So if $A$ has $c$ pivot positions, it has one in every row and every column, and each one is to the right of the one above it, so they all must be on the main diagonal. Which means the reduced row echelon form of $A$ is $I_c$.

(v \leftrightarrow\hspace{1pt}  vi \leftrightarrow\hspace{1pt} vii) Clearly $A$ has $c$ pivots if and only if it has one in every column. By the uniqueness theorem \ref{theorem:1.7uniqueness-v3}, this is equivalent to having a linearly independent list of column vectors. And linearly independent column vectors means the same thing as $\nul A = \left\{ \vec{0} \right\}$ by definition.

(v \leftrightarrow\hspace{1pt} viii \leftrightarrow\hspace{1pt} ix) Clearly $A$ has $c$ pivots if and only if it has one in every row. By the existence theorem \ref{theorem:1.7existence-v3}, this is equivalent to $A\vec{x}=\vec{y}$ being consistent for all $\vec{y}$ in $\R^c$. And by definition of $\col A$, that means the same thing as $\col A = \R^c$.

(i \rightarrow\hspace{1pt} ii) If $A$ is invertible, then $A^{-1}$ works as $P$.

(i \rightarrow\hspace{1pt} iii) If $A$ is invertible, then $A^{-1}$ works as $Q$.

(ii \rightarrow\hspace{1pt} vi) This is theorem \ref{theorem:PA-implies-independence} from the section \ref{sec:inverses} exercises.

(iii \rightarrow\hspace{1pt} viii) This is theorem \ref{theorem:AQ-implies-spanning} from the section \ref{sec:inverses} exercises.

\textcolor{red}{(i \leftarrow \hspace{1pt} iv)} Suppose $A \sim I$. Then, there's a sequence of row operations that reduce $A$ to $I$. And each of those row operations has an elementary matrix, that when you multiply by $A$ on the left, applies that operation to $A$ (theorem \ref{theorem:elementary-operations}). In other words, there exists $E_1,E_2,\ldots,E_n$ such that:
\[ A \sim E_1A \sim E_2E_1A \sim \cdots \sim E_nE_{n-1}\cdots E_2E_1A = I\]
Every elementary matrix is invertible, and by theorem \ref{theorem:properties-of-inverses} (c), the product of invertible matrices is invertible. So $( E_n\cdots E_1 )^{-1}$ exists and we can multiply both sides by it:
\begin{align*}
( E_n\cdots E_1 )^{-1} (  E_n\cdots E_1)A &= ( E_n\cdots E_1 )^{-1}I \\
I A&= ( E_n\cdots E_1 )^{-1}I \\
A &= ( E_n\cdots E_1 )^{-1}
\end{align*}
Now, theorem \ref{theorem:properties-of-inverses} (b) says that an inverse matrix is invertible, which means $A$ is invertible, proving the theorem. It's worth noting that:
\[(A)^{-1} = (( E_n\cdots E_1 )^{-1})^{-1} = E_n\cdots E_1 = E_n\cdots E_1I\]
which says the same sequence of row operations that reduce $A$ to $I$ also reduce $I$ to $A^{-1}$, which agrees with theorem \ref{theorem:AIequivIA-inv}. \emph{Credit:} This proof was inspired by ``Linear Algebra and Its Applications'' by David Lay et al. \blacksquare
\end{book}

\columnbreak
\exerciseheading

\exatend (T/F?) All elementary matrices are invertible.

\exatend (T/F?) If $\vec{a}_1, \ldots, \vec{a}_c$ is a list of $c$ linearly independent vectors in $\R^c$, then $\begin{bmatrix}
\vec{a}_1 \cdots \vec{a}_c
\end{bmatrix}$ is an invertible matrix.

\exatend (T/F?) If $A$ is $c\times c$, and $\dim(\nul A) = c$, then $A$ is invertible.

\exatend (T/F?) If $A$ is invertible, then there's a nontrivial solution to the homogeneous equation $A \vec{x}=\vec{0}$.

\exatend (T/F?) If the columns of $c\times c$ matrix $A$ do not span $\R^c$, then $A$ is not invertible.

\exatend (T/F?) If $A$ is an invertible $c \times c$ matrix, then the columns of $A$ are a basis of $\R^c$.

\exatend (T/F?) If $A$ is $c\times c$, and $A\vec{x}=\vec{0}$ has only the trivial solution, then $\row A = \R^c$.

\exatend (T/F?) If $A$ is invertible, and $B$ is row-equivalent to $A$, then $B$ is invertible.
\creditfootnote{Inspired by an exercise in David Austin's excellent \emph{Understanding Linear Algebra}.}

\exatend (T/F?) An elementary matrix is any square matrix that results from applying some sequence of row operations to $I$.

\exatend (T/F?) If $E$ is an elementary matrix, then $AE$ is the result of applying the same row operation to $A$ that transformed $I$ to $E$.
 
%doesn't use 2.4 but that's ok
\exatend Suppose $A = \begin{bmatrix}
a & b \\
c & d
\end{bmatrix}$, is invertible and let $\mathcal{A} = ( \vec{a}_1, \vec{a}_2 )$ be the basis of $\R^2$ made by the columns of $A$. 
\begin{enumerate}[label=(\alph*)]
\item Find $\left[ \vec{e}_1 \right]_{\mathcal{A}}$ and $\left[ \vec{e}_2 \right]_{\mathcal{A}}$.
\item Let $\mathcal{E} = ( \vec{e}_1,\vec{e}_2 )$. Find $\left[ \vec{a}_1 \right]_{\mathcal{E}}$ and $\left[ \vec{a}_2 \right]_{\mathcal{E}}$.
\end{enumerate}

\exatend If $A = \begin{bmatrix}
a & b \\
c & d
\end{bmatrix}$ and $EA = \begin{bmatrix}
a & b \\
c-2a & d-2b
\end{bmatrix}$, find $E$.

\exatend Let $A = \begin{bmatrix}
1 & 2 & 0 \\
0 & 1 & 2 \\
0 & 0 & 1
\end{bmatrix}$. Since $A \sim I$, there's a sequence of row operations that reduce $A$ to $I$, corresponding to elementary matrices $E_1, \ldots, E_n$. Find these elementary matrices, and multiply to check that  $E_n\cdots E_1A = I$. 

\exatend Let $A = \begin{bmatrix}
0 & 2 \\
5 & 0
\end{bmatrix}$.  Since $A \sim I$, there's a sequence of row operations that reduce $A$ to $I$, corresponding to elementary matrices $E_1, \ldots, E_n$. Find these elementary matrices, and multiply to check that  $E_n\cdots E_1A = I$. 

\exatend Let $A = \begin{bmatrix}
0 & 1 & 0 & 0 \\
0 & 0 & 1 & 0 \\
0 & 0 & 0& 1 \\
1 & 0 & 0 & 0
\end{bmatrix}$.  Since $A \sim I$, there's a sequence of row operations that reduce $A$ to $I$, corresponding to elementary matrices $E_1, \ldots, E_n$. Find these elementary matrices, and multiply to check that  $E_n\cdots E_1A = I$. 

\exatend Suppose $A,B,P$ are square matrices such that $A = PBP^{-1}$. \creditfootnote{Inspired by Austin.}
\begin{enumerate}[label=(\alph*)]
\item Show that if $A$ is invertible, then so is $B$.
\item Show that $A^2 = PB^2P^{-1}$.
\end{enumerate}

\exatend \label{exercise:2.4-ABinv-Ainv} Show that if $A$ is square, and $AB$ is invertible, then $A$ is invertible. \hintfootnote{Theorem \ref{theorem:properties-of-inverses} doesn't help you here, because you don't know that $A$ and $B$ are invertible. But the IMT might help.}

\exatend \label{exercise:2.4-ABinv-Binv} Show that if $B$ is square, and $AB$ is invertible, then $B$ is invertible. \hintfootnote{Theorem \ref{theorem:properties-of-inverses} doesn't help you here, because you don't know that $A$ and $B$ are invertible. But the IMT might help.}


%credit: david austin
\exatend Suppose that $A^{3} = AAA$ is invertible, and $(A^3)^{-1}=B$. Show that $A$ is invertible and find $A^{-1}$. \creditfootnote{Inspired by Austin.}

\exatend Let $E = \begin{bmatrix}
1 & 0 & 0 \\
0 & 0 & 1 \\
0 & 1 & 0 
\end{bmatrix}, F = \begin{bmatrix}
1 & -7 & 0 \\
0 & 1 & 0 \\
0 & 0 & 1
\end{bmatrix}$, and $A = \begin{bmatrix}
a & b & c \\
d & e & f \\
g & h & i
\end{bmatrix}$. Note that you can think of $E$ as the result of adding -7 times \emph{column} 1 of $I$ to \emph{column} 2 of $I$. Likewise $F$ as the result of swapping \emph{columns} 2 and 3. Multiply to find $AE$ and $AF$. (Note that the result is like applying the same column operation to $A$; this is true in general.)

% doesn't use 2.4 but good review.
\exatend Suppose $T:\R^2\rightarrow\R^2$ is a linear transformation such that $T(-3,5) = (4,2)$ and $T(1,-2) = (0,3)$. Find $A$ such that $T(\vec{x}) = A\vec{x}$ using the following steps:
\begin{enumerate}[label=(\alph*)]
\item Fill in the ?'s: $A\begin{bmatrix}
-3 & 1 \\
5 & -2
\end{bmatrix} = \begin{bmatrix}
? & ? \\
? & ?
\end{bmatrix}$.
\item Find the inverse of $\begin{bmatrix}
-3 & 1 \\
5 & -2
\end{bmatrix}$ and use it to find $A$.
\end{enumerate}

\chapter{Determinants}
\label{chap:determinants}
In this chapter, we'll see that every square matrix $A$ is associated with a real number, called the \emph{determinant} of $A$ or $\det(A)$, which tells you a bunch of useful things about $A$:
\begin{itemize}
\item $\det(A) \ne 0$ if and only if $A$ is invertible.
\item $\det(A)$ tells you the area of a parallelogram or volume of a parallelepiped described by $A$ is.
\item $\det(A)$ tells you how much $A$ stretches or compresses other vectors/matrices that you multiply by $A$, which makes determinants important in applications of multivariable calculus.
\end{itemize}
All of these are different ways of describing the ``how big'' $A$ is. Unfortunately, the formula for $\det(A)$ is pretty complicated, and it will take some work to figure it out. Fortunately, there's a lot to of interesting math to learn along the way.

This approach to covering determinants was pioneered by Sergei Treil in his excellent book \emph{Linear Algebra Done Wrong}.
\section{Signed Volume}
\label{sec:signed-volume}

\bbtwocol{0.6}{0.4}{
Recall that the area of a parallelogram is the length of one side (called the \emph{base}) times the distance from that side perpendicularly to the other side (the \emph{height}).
}{
\begin{tikzpicture}
\pgfmathsetmacro{\ux}{3}
\pgfmathsetmacro{\uy}{0}
\pgfmathsetmacro{\vx}{1}  
\pgfmathsetmacro{\vy}{2}  
\begin{axis}[
x=30pt,y=30pt,
axis lines=middle,
axis line style={-},
ymajorgrids=true,
%minor y tick num=3,
%yminorgrids=true,
xmajorgrids=true,
%minor x tick num=3,
%xminorgrids=true,
%xlabel=$x_1$,
%ylabel=$x_2$,
xmin=-0.5, xmax=4,
ymin=-0.6, ymax=2.5,
xtick={-1,...,4},
ytick={-1,...,4},
xticklabels={},
yticklabels={},
]
\filldraw[blue!10] (axis cs:0,0) -- (axis cs:\ux,\uy) -- (axis cs: \ux+\vx, \uy+\vy) -- (axis cs: \vx,\vy) -- (axis cs:0,0);
\draw[blue] (axis cs:0,0) -- (axis cs:\ux,\uy) -- (axis cs: \ux+\vx, \uy+\vy) -- (axis cs: \vx,\vy) -- (axis cs:0,0);
\draw[|-|,red,thick] (axis cs:0,-0.2) -- (axis cs: \ux,\uy-0.2);
\node[red] at (axis cs:1.5,-0.4) {$b$};
\draw[|-|,red,thick] (axis cs:2,0) -- (axis cs: 2,2);
\node[red] at (axis cs:2.2,1) {$h$};
\draw[-,red,thick] (axis cs:2,0.2) -- (axis cs: 1.8,0.2) -- (axis cs: 1.8, 0);
\end{axis}
\end{tikzpicture}
}

Let's describe a parallelogram in terms of vectors and matrices, which we know a lot about. If $A = \begin{bmatrix}
\vec{a}_1 & \vec{a}_2
\end{bmatrix}$, we'll say the \emph{parallelogram determined by} $A$, written $\pgram A$, is the parallelogram where two sides are the arrows $\vec{a}_1$ and $\vec{a}_2$, drawn starting at $\vec{0}$. For example, the one above could be written $\pgram \begin{bmatrix}
3 & 1 \\
0 & 2
\end{bmatrix}$.

\exercise (\emph{Warm-up}) Let $A =  \begin{bmatrix}
\vec{a}_1 & \vec{a}_2
\end{bmatrix} = \begin{bmatrix}
3 & 1 \\
0 & 2
\end{bmatrix}$. Find the area of each shaded parallelogram. Also, which pairs of vectors are linearly independent?

\begin{horiparts}{2}

\item \begin{tikzpicture}
\pgfmathsetmacro{\ux}{3}
\pgfmathsetmacro{\uy}{0}
\pgfmathsetmacro{\vx}{1}  
\pgfmathsetmacro{\vy}{2}  
\begin{axis}[
x=25pt,y=25pt, axis lines=middle, axis line style={-},
ymajorgrids=true, xmajorgrids=true,
xmin=-0.5, xmax=4.2, ymin=-0.6, ymax=2.5,
xtick={-1,...,4}, ytick={-1,...,4}, xticklabels={}, yticklabels={}
]
\filldraw[blue!10] (axis cs:0,0) -- (axis cs:\ux,\uy) -- (axis cs: \ux+\vx, \uy+\vy) -- (axis cs: \vx,\vy) -- (axis cs:0,0);
\draw[blue] (axis cs:0,0) -- (axis cs:\ux,\uy) -- (axis cs: \ux+\vx, \uy+\vy) -- (axis cs: \vx,\vy) -- (axis cs:0,0);
\draw[-Stealth, thick] (axis cs:0,0) -- (axis cs:\ux,\uy) node[below] {$\vec{a}_1$};
\draw[-Stealth, thick] (axis cs:0,0) -- (axis cs:\vx,\vy) node[left] {$\vec{a}_2$};
\node[left] at (axis cs:0,1) {$1$};
\node[below] at (axis cs:1,0) {$1$};
\end{axis}
\end{tikzpicture}

\item \begin{tikzpicture}
\pgfmathsetmacro{\ux}{3}
\pgfmathsetmacro{\uy}{0}
\pgfmathsetmacro{\vx}{0}  
\pgfmathsetmacro{\vy}{2}  
\begin{axis}[
x=25pt,y=25pt, axis lines=middle, axis line style={-},
ymajorgrids=true, xmajorgrids=true,
xmin=-0.5, xmax=4.2, ymin=-0.6, ymax=2.5,
xtick={-1,...,4}, ytick={-1,...,4}, xticklabels={}, yticklabels={}
]
\filldraw[blue!10] (axis cs:0,0) -- (axis cs:\ux,\uy) -- (axis cs: \ux+\vx, \uy+\vy) -- (axis cs: \vx,\vy) -- (axis cs:0,0);
\draw[blue] (axis cs:0,0) -- (axis cs:\ux,\uy) -- (axis cs: \ux+\vx, \uy+\vy) -- (axis cs: \vx,\vy) -- (axis cs:0,0);
\draw[-Stealth, thick] (axis cs:0,0) -- (axis cs:\ux,\uy) node[below] {$\vec{a}_1$};
\draw[-Stealth, thick] (axis cs:0,0) -- (axis cs:\vx,\vy) node[below right] {$\vec{a}_2-\frac{1}{3}\vec{a}_1$};
\node[left] at (axis cs:0,1) {$1$};
\node[below] at (axis cs:1,0) {$1$};
\end{axis}
\end{tikzpicture}

\vspace{1ex} \relax

\item \begin{tikzpicture}
\pgfmathsetmacro{\ux}{3}
\pgfmathsetmacro{\uy}{-2}
\pgfmathsetmacro{\vx}{0}  
\pgfmathsetmacro{\vy}{2}  
\begin{axis}[
x=25pt,y=25pt, axis lines=middle, axis line style={-},
ymajorgrids=true, xmajorgrids=true,
xmin=-0.5, xmax=4.6, ymin=-2.95, ymax=2.5,
xtick={-1,...,4}, ytick={-3,...,4}, xticklabels={}, yticklabels={}
]
\filldraw[blue!10] (axis cs:0,0) -- (axis cs:\ux,\uy) -- (axis cs: \ux+\vx, \uy+\vy) -- (axis cs: \vx,\vy) -- (axis cs:0,0);
\draw[blue] (axis cs:0,0) -- (axis cs:\ux,\uy) -- (axis cs: \ux+\vx, \uy+\vy) -- (axis cs: \vx,\vy) -- (axis cs:0,0);
\draw[-Stealth, thick] (axis cs:0,0) -- (axis cs:\ux,\uy) node[below] {$\vec{a}_1-( \vec{a}_2-\frac{1}{3}\vec{a}_1 )$};
\draw[-Stealth, thick] (axis cs:0,0) -- (axis cs:\vx,\vy) node[right] {$\vec{a}_2-\frac{1}{3}\vec{a}_1$};
\node[left] at (axis cs:0,1) {$1$};
%\node[below] at (axis cs:1,0) {$1$};
\end{axis}
\end{tikzpicture}

\item \begin{tikzpicture}
\pgfmathsetmacro{\ux}{3}
\pgfmathsetmacro{\uy}{-2}
\pgfmathsetmacro{\vx}{0}  
\pgfmathsetmacro{\vy}{1}  
\begin{axis}[
x=25pt,y=25pt, axis lines=middle, axis line style={-},
ymajorgrids=true, xmajorgrids=true,
xmin=-0.5, xmax=4.6, ymin=-2.95, ymax=1.9,
xtick={-1,...,4}, ytick={-3,...,4}, xticklabels={}, yticklabels={}
]
\filldraw[blue!10] (axis cs:0,0) -- (axis cs:\ux,\uy) -- (axis cs: \ux+\vx, \uy+\vy) -- (axis cs: \vx,\vy) -- (axis cs:0,0);
\draw[blue] (axis cs:0,0) -- (axis cs:\ux,\uy) -- (axis cs: \ux+\vx, \uy+\vy) -- (axis cs: \vx,\vy) -- (axis cs:0,0);
\draw[-Stealth, thick] (axis cs:0,0) -- (axis cs:\ux,\uy) node[below] {$\vec{a}_1-( \vec{a}_2-\frac{1}{3}\vec{a}_1 )$};
\draw[-Stealth, thick] (axis cs:0,0) -- (axis cs:\vx,\vy) node[above right] {$\frac{1}{2}(  \vec{a}_2-\frac{1}{3}\vec{a}_1)$};
\node[left] at (axis cs:0,1) {$1$};
%\node[below] at (axis cs:1,0) {$1$};
\end{axis}
\end{tikzpicture}

\vspace{2em}
\item \begin{tikzpicture}
\pgfmathsetmacro{\ux}{-1}
\pgfmathsetmacro{\uy}{0}
\pgfmathsetmacro{\vx}{3}  
\pgfmathsetmacro{\vy}{1}  
\begin{axis}[
x=25pt,y=25pt, axis lines=middle, axis line style={-},
ymajorgrids=true, xmajorgrids=true,
xmin=-1.5, xmax=4.2, ymin=-0.6, ymax=1.5,
xtick={-1,...,4}, ytick={-1,...,4}, xticklabels={}, yticklabels={}
]
\filldraw[blue!10] (axis cs:0,0) -- (axis cs:\ux,\uy) -- (axis cs: \ux+\vx, \uy+\vy) -- (axis cs: \vx,\vy) -- (axis cs:0,0);
\draw[blue] (axis cs:0,0) -- (axis cs:\ux,\uy) -- (axis cs: \ux+\vx, \uy+\vy) -- (axis cs: \vx,\vy) -- (axis cs:0,0);
\draw[-Stealth, thick] (axis cs:0,0) -- (axis cs:\ux,\uy) node[above] {$-\vec{e}_1$};
\draw[-Stealth, thick] (axis cs:0,0) -- (axis cs:\vx,\vy) node[below] {$3\vec{e}_1+\vec{e}_2$};
\node[left] at (axis cs:0,1) {$1$};
\node[below] at (axis cs:1,0) {$1$};
\end{axis}
\end{tikzpicture}

%WTF is going on with the spacing around c,d,e,f?
\vspace{2em}
\item \begin{tikzpicture}
\pgfmathsetmacro{\ux}{1}
\pgfmathsetmacro{\uy}{0.5}
\pgfmathsetmacro{\vx}{2}  
\pgfmathsetmacro{\vy}{1}  
\begin{axis}[
x=25pt,y=25pt, axis lines=middle, axis line style={-},
ymajorgrids=true, xmajorgrids=true,
xmin=-0.5, xmax=3.5, ymin=-0.6, ymax=2.5,
xtick={-1,...,4}, ytick={-1,...,4}, xticklabels={}, yticklabels={}
]
\filldraw[blue!10] (axis cs:0,0) -- (axis cs:\ux,\uy) -- (axis cs: \ux+\vx, \uy+\vy) -- (axis cs: \vx,\vy) -- (axis cs:0,0);
\draw[blue] (axis cs:0,0) -- (axis cs:\ux,\uy) -- (axis cs: \ux+\vx, \uy+\vy) -- (axis cs: \vx,\vy) -- (axis cs:0,0);
\draw[-Stealth, thick] (axis cs:0,0) -- (axis cs:\ux,\uy) node[above] {$\bbvec{1}{.5}$};
\draw[-Stealth, thick] (axis cs:0,0) -- (axis cs:\vx,\vy) node[below right] {$\bbvec{2}{1}$};
\node[left] at (axis cs:0,1) {$1$};
\node[below] at (axis cs:1,0) {$1$};
\end{axis}
\end{tikzpicture}
\end{horiparts}

\columnbreak

A \emph{parallelepiped} is like a 3-dimensional parallelogram, like the one below, where the six faces are all parallelograms. Its volume is the $Bh$, where $B$ is the area of one face, and $h$ is the distance perpendicularly from that face to the opposite face. Note that $B=lw$, so its volume is $lwh$. We'll also describe parallelepipeds as $\pgram \begin{bmatrix}
\vec{a}_1 & \vec{a}_2 & \vec{a}_3
\end{bmatrix}$ just like parallelograms.

\begin{tikzpicture}[scale=0.85,line join = round, line cap = round]
\pgfmathsetmacro{\ux}{-5}
\pgfmathsetmacro{\uy}{0}
\pgfmathsetmacro{\uz}{0}
\pgfmathsetmacro{\vx}{1.5}  
\pgfmathsetmacro{\vy}{0}  
\pgfmathsetmacro{\vz}{-3}  
\pgfmathsetmacro{\wx}{-.9}  
\pgfmathsetmacro{\wy}{3}  
\pgfmathsetmacro{\wz}{1.8}
\draw[blue] (0,0,0) -- (\ux,\uy,\uz);
\draw[blue] (0,0,0) -- (\vx,\vy,\vz) -- (\wx+\vx,\wy+\vy,\wz+\vz) -- (\wx,\wy,\wz) -- cycle;  
\draw[blue,dashed] (\ux,\uy,\uz) -- (\ux+\vx,\uy+\vy,\uz+\vz) -- (\vx,\vy,\vz); 
\draw[blue,dashed] (\ux+\vx,\uy+\vy,\uz+\vz) -- (\ux+\vx+\wx,\uy+\vy+\wy,\uz+\vz+\wz);
\draw[blue] (\wx,\wy,\wz) -- (\wx+\ux,\wy+\uy,\wz+\uz) -- (\ux+\vx+\wx,\uy+\vy+\wy,\uz+\vz+\wz) -- (\vx+\wx,\vy+\wy,\vz+\wz);
\draw[blue] (\ux,\uy,\uz) -- (\wx+\ux,\wy+\uy,\wz+\uz);

\draw[-stealth,thick] (0,0,0) -- (\ux,\uy,\uz);
\node[left] at (\ux,\uy,\uz) {$\vec{a}_{1}$};
\draw[-stealth,thick] (0,0,0) -- (\vx,\vy,\vz);
\node[right] at (\vx,\vy,\vz) {$\vec{a}_{2}$};
\draw[-stealth,thick] (0,0,0) -- (\wx,\wy,\wz);
\node[below left] at (\wx,\wy,\wz) {$\vec{a}_{3}$};

\draw[|-|,red,thick] (0,-.2,0) -- (\ux,\uy-0.2,\uz);
\node[red] at (\ux/2,-0.5,0) {$l$};
\draw[-,red,thick] (-2.5,0,0) -- (-2.5,0,-3);
\draw[-,red,thick] (-2.65,0,0) -- (-2.35,0,0);
\draw[-,red,thick] (-2.65,0,-3) -- (-2.35,0,-3);
\draw[-,red,thick] (-2.5,0,-.3) -- (-2.2,0,-0.3) -- (-2.2,0,0);
\node[red] at (-2.8,0,-1.5) {$w$};
\draw[|-|,red,thick] (-3.5,0,-1) -- (-3.5,3,-1);
\draw[-,red,thick] (-3.5,.3,-1) -- (-3.2,0.3,-1) -- (-3.2,0,-1);
\node[red] at (-3.7,1.5,-1) {$h$};
\end{tikzpicture}

If we have a $4\times4$ matrix $A = \begin{bmatrix}
\vec{a}_1 & \vec{a}_2 & \vec{a}_3 & \vec{a}_4
\end{bmatrix}$ (or a bigger one), we can't draw a picture of $\pgram A$, but we can talk about its ``hypervolume'' as its length times width times breadth times depth, where those distances are all perpendicular to each other. We'll refer to all of these things (\emph{area} of parallelogram, \emph{volume} of parallelepiped, etc) as the \textbf{volume} of $\pgram A$. But, geometry in $\R^4$, $\R^{10}$, etc can be very weird. How do we even know that the volume of an object there is even well-defined? Maybe the result depends on which vector you start with to be the length and which direction is the width? What does ``perpendicular'' even \emph{mean} in $\R^{17}$? Instead of answering all these questions, we'll define properties that \emph{volume} should have, that are consistent with $\R^{2}$ and $\R^{3}$. 

\begin{definition}
A \textbf{determinant} is a function whose domain is the set of all square matrices, and whose codomain is $\R$, which satisfies the properties below. Since a matrix is basically the same thing as a list of it's column vectors, we'll often write $\det(\vec{a}_1, \ldots, \vec{a}_c)$ instead of $\det\left( \begin{bmatrix}
\vec{a}_1 & \cdots & \vec{a}_c
\end{bmatrix} \right)$.
\begin{itemize}
\item \textbf{Normality:} $\det(I_c) = \det(\vec{e}_1, \ldots, \vec{e}_c)= 1$ for any $c$.
\item \textbf{Linearity (scaling):} Multiplying any one column by a scalar also scales the determinant: \[\det (\vec{a}_1, \ldots, k\vec{a}_j, \ldots, \vec{a}_c) = k \det (\vec{a}_1, \ldots, \vec{a}_j, \ldots, \vec{a}_c)\]
\item \textbf{Linearity (addition):} The determinant distributes over vector addition in any one column:
\[\det(\vec{a}_1, \ldots, \vec{a}_j+\vec{b}_j,\ldots, \vec{a}_c) = \det(\vec{a}_1, \ldots, \vec{a}_j,\ldots, \vec{a}_c) + \det(\vec{a}_1, \ldots, \vec{b}_j,\ldots, \vec{a}_c)\]
\item \textbf{Preservation under column multaddiplacement:} Replacing one column with the sum of itself and a scalar multiple of a \emph{different} column doesn't change the determinant. For $i \ne j$:
\[\det(\vec{a}_1,\ldots,\vec{a}_j+k\vec{a}_i, \ldots, \vec{a}_c) = \det(\vec{a}_1,\ldots,\vec{a}_j, \ldots, \vec{a}_c)\]
\end{itemize}
\end{definition}

At this point in our story, maybe there doesn't actually exist a function satisfying these properties, or maybe there is more than one, so technically we should be talking about ``a determinant (if it exists).'' But we'll see in the next section that there exists a unique determinant function, and in the meantime, the it will be easier to talk about ``the determinant'' so we'll do that.
Now, let's make sure the properties in the definition are consistent with volume, starting with the easiest one:

\textbf{Normality.} The volume of a unit square or cube is clearly 1. (Remember in $\R^2$, \emph{volume} here really means area.)

\hfill\begin{tikzpicture}[scale=1.1]
\pgfmathsetmacro{\ux}{1}
\pgfmathsetmacro{\uy}{0}
\pgfmathsetmacro{\vx}{0}  
\pgfmathsetmacro{\vy}{1}  
\begin{axis}[
x=30pt,y=30pt, axis lines=middle, axis line style={-},
ymajorgrids=true, xmajorgrids=true,
xmin=-.5, xmax=1.5, ymin=-0.6, ymax=1.6,
xtick={-1,...,4}, ytick={-1,...,4}, xticklabels={}, yticklabels={}
]
\filldraw[blue!10] (axis cs:0,0) -- (axis cs:\ux,\uy) -- (axis cs: \ux+\vx, \uy+\vy) -- (axis cs: \vx,\vy) -- (axis cs:0,0);
\draw[blue] (axis cs:0,0) -- (axis cs:\ux,\uy) -- (axis cs: \ux+\vx, \uy+\vy) -- (axis cs: \vx,\vy) -- (axis cs:0,0);
\draw[-Stealth, thick] (axis cs:0,0) -- (axis cs:\ux,\uy) node[above right] {$\vec{e}_1$};
\draw[-Stealth, thick] (axis cs:0,0) -- (axis cs:\vx,\vy) node[above right] {$\vec{e}_2$};
\node[left] at (axis cs:0,1) {$1$};
\node[below] at (axis cs:1,0) {$1$};
\end{axis}
\end{tikzpicture} 
\hfill
\begin{tikzpicture}[scale=1.3,line join = round, line cap = round]
\pgfmathsetmacro{\ux}{1}
\pgfmathsetmacro{\uy}{0}
\pgfmathsetmacro{\uz}{0}
\pgfmathsetmacro{\vx}{0}  
\pgfmathsetmacro{\vy}{1}  
\pgfmathsetmacro{\vz}{0}  
\pgfmathsetmacro{\wx}{0}  
\pgfmathsetmacro{\wy}{0}  
\pgfmathsetmacro{\wz}{1}
\draw[blue] (0,0,0) -- (\ux,\uy,\uz);
\draw[blue] (0,0,0) -- (\vx,\vy,\vz) -- (\wx+\vx,\wy+\vy,\wz+\vz) -- (\wx,\wy,\wz) -- cycle;  
\draw[blue] (\ux,\uy,\uz) -- (\ux+\vx,\uy+\vy,\uz+\vz) -- (\vx,\vy,\vz); 
\draw[blue] (\ux+\vx,\uy+\vy,\uz+\vz) -- (\ux+\vx+\wx,\uy+\vy+\wy,\uz+\vz+\wz);
\draw[blue] (\wx,\wy,\wz) -- (\wx+\ux,\wy+\uy,\wz+\uz) -- (\ux+\vx+\wx,\uy+\vy+\wy,\uz+\vz+\wz) -- (\vx+\wx,\vy+\wy,\vz+\wz);
\draw[blue] (\ux,\uy,\uz) -- (\wx+\ux,\wy+\uy,\wz+\uz);

\draw[-stealth,thick] (0,0,0) -- (\ux,\uy,\uz);
\node[right] at (\ux,\uy,\uz) {$\vec{e}_{2}$};
\draw[-stealth,thick] (0,0,0) -- (\vx,\vy,\vz);
\node[above] at (\vx,\vy,\vz) {$\vec{e}_{3}$};
\draw[-stealth,thick] (0,0,0) -- (\wx,\wy,\wz);
\node[below left] at (\wx,\wy,\wz) {$\vec{e}_{1}$};
\end{tikzpicture} \hfill \hspace{1pt}

\bbtwocol{.7}{.3}{
\textbf{Linearity (scaling).} When you scale one of the vectors by a factor of $k$, either the base or the height changes by that factor. For example, the volume of $\pgram \begin{bmatrix}
2\vec{e}_1 & \vec{e}_2
\end{bmatrix}$ is 2 which equals $2\det(I)$.}{
\hfill\begin{tikzpicture}
\pgfmathsetmacro{\ux}{2}
\pgfmathsetmacro{\uy}{0}
\pgfmathsetmacro{\vx}{0}  
\pgfmathsetmacro{\vy}{1}  
\begin{axis}[
x=30pt,y=30pt, axis lines=middle, axis line style={-},
ymajorgrids=true, xmajorgrids=true,
xmin=-.5, xmax=2.6, ymin=-0.6, ymax=1.6,
xtick={-1,...,4}, ytick={-1,...,4}, xticklabels={}, yticklabels={}
]
\filldraw[blue!10] (axis cs:0,0) -- (axis cs:\ux,\uy) -- (axis cs: \ux+\vx, \uy+\vy) -- (axis cs: \vx,\vy) -- (axis cs:0,0);
\draw[blue] (axis cs:0,0) -- (axis cs:\ux,\uy) -- (axis cs: \ux+\vx, \uy+\vy) -- (axis cs: \vx,\vy) -- (axis cs:0,0);
\draw[-Stealth, thick] (axis cs:0,0) -- (axis cs:\ux,\uy) node[above right] {$2\vec{e}_1$};
\draw[-Stealth, thick] (axis cs:0,0) -- (axis cs:\vx,\vy) node[above right] {$\vec{e}_2$};
\node[left] at (axis cs:0,1) {$1$};
\node[below] at (axis cs:1,0) {$1$};
\end{axis}
\end{tikzpicture}}

\bbtwocol{.6}{.38}{
Note this also allows for determinants to be negative. For example:
\[\det\begin{bmatrix}
-2\vec{e}_1 & \vec{e}_2
\end{bmatrix} = -2\det(I) = -2(1) = -2\]}{
\begin{tikzpicture}
\pgfmathsetmacro{\ux}{-2}
\pgfmathsetmacro{\uy}{0}
\pgfmathsetmacro{\vx}{0}  
\pgfmathsetmacro{\vy}{1}  
\begin{axis}[
x=30pt,y=30pt, axis lines=middle, axis line style={-},
ymajorgrids=true, xmajorgrids=true,
xmin=-2.95, xmax=1.2, ymin=-0.6, ymax=1.6,
xtick={-2,...,4}, ytick={-1,...,4}, xticklabels={}, yticklabels={}
]
\filldraw[red!10] (axis cs:0,0) -- (axis cs:\ux,\uy) -- (axis cs: \ux+\vx, \uy+\vy) -- (axis cs: \vx,\vy) -- (axis cs:0,0);
\draw[red] (axis cs:0,0) -- (axis cs:\ux,\uy) -- (axis cs: \ux+\vx, \uy+\vy) -- (axis cs: \vx,\vy) -- (axis cs:0,0);
\draw[-Stealth, thick] (axis cs:0,0) -- (axis cs:\ux,\uy) node[above left] {$-2\vec{e}_1$};
\draw[-Stealth, thick] (axis cs:0,0) -- (axis cs:\vx,\vy) node[above right] {$\vec{e}_2$};
\node[left] at (axis cs:0,1) {$1$};
\node[below] at (axis cs:1,0) {$1$};
\end{axis}
\end{tikzpicture}
}

Negative ``area'' or ``volume'' might seem like a contradiction, but it lets us have linearity which makes for elegant and powerful mathematics. Sort of like how, in calculus, you allow an integral to be ``signed area'' between the $x$-axis and a curve, and it makes important rules work, such as $\int c f(x)dx = c \int f(x)dx$.

It turns out that the sign of $2\times2$ determinants depends on whether you turn counterclockwise or clockwise to get from the first column to the second. And there's a ``right hand rule'' that describes the sign of $3\times3$ determinants, having to do with the way the three vectors are oriented.

\textbf{Linearity (addition).} When two vectors in column $j$ are added, we can think of the volume of the resulting parallepiped,
\[\det (\vec{a}_1, \ldots, \vec{a}_j + \vec{b}_j , \ldots , \vec{a}_c),\]
as the ``height'' of $\vec{a}_j + \vec{b}_j$, times the perpendicular ``base,'' which is a lower-dimensional parallelogram in the span of the other vector(s). The height of $\vec{a}_j + \vec{b}_j$ is the sum of the heights of $\vec{a}_j$ and $\vec{b}_j$, as the picture shows.

\begin{tikzpicture}
\begin{scope}[shift={(0,0)}]
\pgfmathsetmacro{\ux}{2}
\pgfmathsetmacro{\uy}{0}
\pgfmathsetmacro{\vx}{0}  
\pgfmathsetmacro{\vy}{1}  
\begin{axis}[
x=30pt,y=30pt, axis lines=middle, axis line style={-},
ymajorgrids=true, xmajorgrids=true,
xmin=-.5, xmax=2.2, ymin=-0.6, ymax=3.6,
xtick={-1,...,4}, ytick={-1,...,4}, xticklabels={}, yticklabels={}
]
\filldraw[blue!10] (axis cs:0,0) -- (axis cs:\ux,\uy) -- (axis cs: \ux+\vx, \uy+\vy) -- (axis cs: \vx,\vy) -- (axis cs:0,0);
\draw[blue] (axis cs:0,0) -- (axis cs:\ux,\uy) -- (axis cs: \ux+\vx, \uy+\vy) -- (axis cs: \vx,\vy) -- (axis cs:0,0);
\draw[-Stealth, thick] (axis cs:0,0) -- (axis cs:\ux,\uy) node[below] {$\vec{a}_1$};
\draw[-Stealth, thick] (axis cs:0,0) -- (axis cs:\vx,\vy) node[above right] {$\vec{a}_2$};
\node[left] at (axis cs:0,1) {$1$};
\node[below] at (axis cs:1,0) {$1$};
\draw[|-|,red,thick] (axis cs:1.5,0) -- (axis cs: 1.5,1);
\node[red] at (axis cs:1,.5) {$h=1$};
\end{axis}
\end{scope}
\node[red] at (3.35,1.3) {\huge$+$};
\begin{scope}[shift={(4,0)}]
\pgfmathsetmacro{\ux}{2}
\pgfmathsetmacro{\uy}{0}
\pgfmathsetmacro{\vx}{1}  
\pgfmathsetmacro{\vy}{2}  
\begin{axis}[
x=30pt,y=30pt, axis lines=middle, axis line style={-},
ymajorgrids=true, xmajorgrids=true,
xmin=-.2, xmax=3.2, ymin=-0.6, ymax=3.6,
xtick={-1,...,4}, ytick={-1,...,4}, xticklabels={}, yticklabels={}
]
\filldraw[blue!10] (axis cs:0,0) -- (axis cs:\ux,\uy) -- (axis cs: \ux+\vx, \uy+\vy) -- (axis cs: \vx,\vy) -- (axis cs:0,0);
\draw[blue] (axis cs:0,0) -- (axis cs:\ux,\uy) -- (axis cs: \ux+\vx, \uy+\vy) -- (axis cs: \vx,\vy) -- (axis cs:0,0);
\draw[-Stealth, thick] (axis cs:0,0) -- (axis cs:\ux,\uy) node[below] {$\vec{a}_1$};
\draw[-Stealth, thick] (axis cs:0,0) -- (axis cs:\vx,\vy) node[left] {$\vec{b}_2$};
\node[left] at (axis cs:0,1) {$1$};
\node[below] at (axis cs:1,0) {$1$};
\draw[|-|,red,thick] (axis cs:1.5,0) -- (axis cs: 1.5,2);
\node[red] at (axis cs:1,.8) {$h=2$};
\end{axis}
\end{scope}
\node[red] at (8.1,1.3) {\huge$=$};
\begin{scope}[shift={(8.7,0)}]
\pgfmathsetmacro{\ux}{2}
\pgfmathsetmacro{\uy}{0}
\pgfmathsetmacro{\vx}{1}  
\pgfmathsetmacro{\vy}{3}  
\begin{axis}[
x=30pt,y=30pt, axis lines=middle, axis line style={-},
ymajorgrids=true, xmajorgrids=true,
xmin=-.3, xmax=3.2, ymin=-0.6, ymax=3.6,
xtick={-1,...,4}, ytick={-1,...,4}, xticklabels={}, yticklabels={}
]
\filldraw[blue!10] (axis cs:0,0) -- (axis cs:\ux,\uy) -- (axis cs: \ux+\vx, \uy+\vy) -- (axis cs: \vx,\vy) -- (axis cs:0,0);
\draw[blue] (axis cs:0,0) -- (axis cs:\ux,\uy) -- (axis cs: \ux+\vx, \uy+\vy) -- (axis cs: \vx,\vy) -- (axis cs:0,0);
\draw[-Stealth, thick] (axis cs:0,0) -- (axis cs:\ux,\uy) node[below] {$\vec{a}_1$};
\draw[-Stealth, thick] (axis cs:0,0) -- (axis cs:\vx,\vy) node[left] {$\vec{a}_2+\vec{b}_2$};
\node[left] at (axis cs:0,1) {$1$};
\node[below] at (axis cs:1,0) {$1$};
\draw[|-|,red,thick] (axis cs:1.5,0) -- (axis cs: 1.5,3);
\node[red] at (axis cs:1,.8) {$h=3$};
\draw[dotted,thick] (axis cs:0,0) -- (axis cs: 0,1);
\draw[dotted,thick] (axis cs: 0,1) -- (axis cs:1,3);
\end{axis}
\end{scope}
\end{tikzpicture}

\textbf{Preservation under column multaddiplacement.} As the figure below shows, replacing $\vec{a}_2$ with $\vec{a}_2-\frac{1}{3}\vec{a}_1$ moves one side in a direction parallel to the opposite side, which doesn't change the area. Thus we want $d(\vec{a}_1,\vec{a}_2)$ to equal $d(\vec{a}_1,\vec{a}_2-\frac{1}{3}\vec{a}_1)$.

\begin{tikzpicture}
\begin{scope}[shift={(0,0)}]
\pgfmathsetmacro{\ux}{3}
\pgfmathsetmacro{\uy}{0}
\pgfmathsetmacro{\vx}{1}  
\pgfmathsetmacro{\vy}{2}  

\begin{axis}[
x=30pt,y=30pt,
axis lines=middle,
axis line style={-},
ymajorgrids=true,
xmajorgrids=true,
xmin=-0.2, xmax=4.2,
ymin=-0.8, ymax=2.7,
xtick={-1,...,4},
ytick={-4,-3,...,4},
xticklabels={},
yticklabels={},
]
\filldraw[blue!10] (axis cs:0,0) -- (axis cs:\ux,\uy) -- (axis cs: \ux+\vx, \uy+\vy) -- (axis cs: \vx,\vy) -- (axis cs:0,0);
\draw[blue] (axis cs:0,0) -- (axis cs:\ux,\uy) -- (axis cs: \ux+\vx, \uy+\vy) -- (axis cs: \vx,\vy) -- (axis cs:0,0);
\draw[-Stealth,thick] (axis cs:0,0) -- (axis cs:\ux,\uy)
node[above] {$\vec{a}_1$};
\draw[-Stealth,thick] (axis cs:0,0) -- (axis cs:\vx,\vy)
node[above left] {$\vec{a}_2$};
\draw[|-|,red,thick] (axis cs:0,-.2) -- (axis cs: \ux,\uy-.2);
\node[red] at (axis cs:1.5,-.4) {$b$};
\draw[|-|,red,thick] (axis cs:2,0) -- (axis cs: 2,2);
\node[red] at (axis cs:2.3,1) {$h$};
\draw[-,red,thick] (axis cs:2,0.2) -- (axis cs: 1.8,0.2) -- (axis cs: 1.8, 0);
\end{axis}
\end{scope}

\draw[-stealth,thick,bend right=10] (4.7,1.5) to (7.3,1.5);
\node at (6,2.7) {\small column};
\node at (6,2.2) {\small multaddiplacement};
\begin{scope}[shift={(7.5,0)}]
\pgfmathsetmacro{\ux}{3}
\pgfmathsetmacro{\uy}{0}
\pgfmathsetmacro{\vx}{0}  
\pgfmathsetmacro{\vy}{2}  
\begin{axis}[
x=30pt,y=30pt,
axis lines=middle,
axis line style={-},
ymajorgrids=true,
xmajorgrids=true,
xmin=-0.5, xmax=4.2,
ymin=-0.8, ymax=2.7,
xtick={-1,...,4},
ytick={-1,...,4},
xticklabels={},
yticklabels={},
]
\filldraw[blue!10] (axis cs:0,0) -- (axis cs:\ux,\uy) -- (axis cs: \ux+\vx, \uy+\vy) -- (axis cs: \vx,\vy) -- (axis cs:0,0);
\draw[blue] (axis cs:0,0) -- (axis cs:\ux,\uy) -- (axis cs: \ux+\vx, \uy+\vy) -- (axis cs: \vx,\vy) -- (axis cs:0,0);
\draw[-Stealth,thick] (axis cs:0,0) -- (axis cs:\ux,\uy)
node[above left] {$\vec{a}_1$};
\draw[-Stealth,thick,dashed] (axis cs:0,0) -- (axis cs:1,2);
\draw[-Stealth,thick,dashed] (axis cs:1,2) -- (axis cs:0,2);
\draw[-Stealth,thick] (axis cs:0,0) -- (axis cs:\vx,\vy)
node[above right] {$\vec{a}_2 - \frac{1}{3}\vec{a}_1$};
\draw[|-|,red,thick] (axis cs:0,-0.2) -- (axis cs: \ux,\uy-0.2);
\node[red] at (axis cs:1.5,-0.4) {$b$};
\draw[|-|,red,thick] (axis cs:2,0) -- (axis cs: 2,2);
\node[red] at (axis cs:2.2,1) {$h$};
\draw[-,red,thick] (axis cs:2,0.2) -- (axis cs: 1.8,0.2) -- (axis cs: 1.8, 0);
\end{axis}
\end{scope}
\end{tikzpicture}

The figure below shows a $3\times3$ example. To make the parallelepiped rectangular, $\vec{a}_3$ can be replaced with $\vec{a}_3+0.3\vec{a}_1$ without changing $h$, since you're moving the top face parallel to the bottom face. And $\vec{a}_2$ can be replaced with $\vec{a}_2+0.3\vec{a}_1$ which doesn't change $w$ because you're moving the back edge parallel to the front edge. So the volume of $\pgram\begin{bmatrix}
\vec{a}_1 & \vec{a}_2 & \vec{a}_3
\end{bmatrix}$ equals the volume of $\pgram\begin{bmatrix}
\vec{a}_1 & \vec{a}_2+0.3\vec{a}_1 & \vec{a}_3+0.6\vec{a}_2
\end{bmatrix}$, so the determinants of the matrices should be equal too.

% Probably junk, but I will hoard it
% \begin{tikzpicture}[line join = round, line cap = round]
% \pgfmathsetmacro{\ux}{-5}
% \pgfmathsetmacro{\uy}{0}
% \pgfmathsetmacro{\uz}{0}
% \pgfmathsetmacro{\vx}{1.5}  
% \pgfmathsetmacro{\vy}{0}  
% \pgfmathsetmacro{\vz}{-3}  
% \pgfmathsetmacro{\wx}{1.5}  
% \pgfmathsetmacro{\wy}{3}  
% \pgfmathsetmacro{\wz}{0}
% \draw[blue] (0,0,0) -- (\ux,\uy,\uz);
% \draw[blue] (0,0,0) -- (\vx,\vy,\vz) -- (\wx+\vx,\wy+\vy,\wz+\vz) -- (\wx,\wy,\wz) -- cycle;  
% \draw[blue,dashed] (\ux,\uy,\uz) -- (\ux+\vx,\uy+\vy,\uz+\vz) -- (\vx,\vy,\vz); 
% \draw[blue,dashed] (\ux+\vx,\uy+\vy,\uz+\vz) -- (\ux+\vx+\wx,\uy+\vy+\wy,\uz+\vz+\wz);
% \draw[blue] (\wx,\wy,\wz) -- (\wx+\ux,\wy+\uy,\wz+\uz) -- (\ux+\vx+\wx,\uy+\vy+\wy,\uz+\vz+\wz) -- (\vx+\wx,\vy+\wy,\vz+\wz);
% \draw[blue] (\ux,\uy,\uz) -- (\wx+\ux,\wy+\uy,\wz+\uz);

% \draw[-stealth,thick] (0,0,0) -- (\ux,\uy,\uz);
% \node[left] at (\ux,\uy,\uz) {$\vec{a}_{1}$};
% \draw[-stealth,thick] (0,0,0) -- (\vx,\vy,\vz);% \node[right];
% \node[right] at (\vx,\vy,\vz) {$\vec{a}_{2}$};
% \draw[-stealth,thick] (0,0,0) -- (\wx,\wy,\wz);% \node[right];
% \node[below right] at (\wx,\wy,\wz) {$\vec{a}_{3}$};

% \draw[|-|,red,thick] (0,-.2,0) -- (\ux,\uy-0.2,\uz);
% \node[red] at (\ux/2,-0.4,0) {$l$};
% \draw[-,red,thick] (-1.5,0,0) -- (-1.5,0,-3);
% \draw[-,red,thick] (-1.65,0,0) -- (-1.35,0,0);
% \draw[-,red,thick] (-1.65,0,-3) -- (-1.35,0,-3);
% \draw[-,red,thick] (-1.5,0,-.3) -- (-1.2,0,-0.3) -- (-1.2,0,0);
% \node[red] at (-1.2,0,-1.5) {$w$};
% \draw[|-|,red,thick] (-2.5,0,-1) -- (-2.5,3,-1);
% \draw[-,red,thick] (-2.5,.3,-1) -- (-2.2,0.3,-1) -- (-2.2,0,-1);
% \node[red] at (-2.7,1.5,-1) {$h$};
% \end{tikzpicture}


% \begin{tikzpicture}[line join = round, line cap = round]
% \begin{scope}[shift={(0,5.2,0)}]
% \pgfmathsetmacro{\ux}{-5}
% \pgfmathsetmacro{\uy}{0}
% \pgfmathsetmacro{\uz}{0}
% \pgfmathsetmacro{\vx}{1.5}  
% \pgfmathsetmacro{\vy}{0}  
% \pgfmathsetmacro{\vz}{-3}  
% \pgfmathsetmacro{\wx}{-.9}  
% \pgfmathsetmacro{\wy}{3}  
% \pgfmathsetmacro{\wz}{1.8}
% \draw[blue] (0,0,0) -- (\ux,\uy,\uz);
% \draw[blue] (0,0,0) -- (\vx,\vy,\vz) -- (\wx+\vx,\wy+\vy,\wz+\vz) -- (\wx,\wy,\wz) -- cycle;  
% \draw[blue,dashed] (\ux,\uy,\uz) -- (\ux+\vx,\uy+\vy,\uz+\vz) -- (\vx,\vy,\vz); 
% \draw[blue,dashed] (\ux+\vx,\uy+\vy,\uz+\vz) -- (\ux+\vx+\wx,\uy+\vy+\wy,\uz+\vz+\wz);
% \draw[blue] (\wx,\wy,\wz) -- (\wx+\ux,\wy+\uy,\wz+\uz) -- (\ux+\vx+\wx,\uy+\vy+\wy,\uz+\vz+\wz) -- (\vx+\wx,\vy+\wy,\vz+\wz);
% \draw[blue] (\ux,\uy,\uz) -- (\wx+\ux,\wy+\uy,\wz+\uz);

% \draw[-stealth,thick] (0,0,0) -- (\ux,\uy,\uz);
% \node[left] at (\ux,\uy,\uz) {$\vec{a}_{1}$};
% \draw[-stealth,thick] (0,0,0) -- (\vx,\vy,\vz);% \node[right];
% \node[right] at (\vx,\vy,\vz) {$\vec{a}_{2}$};
% \draw[-stealth,thick] (0,0,0) -- (\wx,\wy,\wz);% \node[right];
% \node[below left] at (\wx,\wy,\wz) {$\vec{a}_{3}$};

% \draw[|-|,red,thick] (0,-.2,0) -- (\ux,\uy-0.2,\uz);
% \node[red] at (\ux/2,-0.4,0) {$l$};
% \draw[-,red,thick] (-2.5,0,0) -- (-2.5,0,-3);
% \draw[-,red,thick] (-2.65,0,0) -- (-2.35,0,0);
% \draw[-,red,thick] (-2.65,0,-3) -- (-2.35,0,-3);
% \draw[-,red,thick] (-2.5,0,-.3) -- (-2.2,0,-0.3) -- (-2.2,0,0);
% \node[red] at (-2.8,0,-1.5) {$w$};
% \draw[|-|,red,thick] (-3.5,0,-1) -- (-3.5,3,-1);
% \draw[-,red,thick] (-3.5,.3,-1) -- (-3.2,0.3,-1) -- (-3.2,0,-1);
% \node[red] at (-3.7,1.5,-1) {$h$};
% \end{scope}
% \draw[-stealth, thick, bend left = 30] (2,5.5,0) to (2,3.5,0); 
% \node[right] at (2.4,4.75,0) {column};
% \node[right] at (2.4,4.3,0) {multaddiplacement};
% \pgfmathsetmacro{\ux}{-5}
% \pgfmathsetmacro{\uy}{0}
% \pgfmathsetmacro{\uz}{0}
% \pgfmathsetmacro{\vx}{0}
% \pgfmathsetmacro{\vy}{0}
% \pgfmathsetmacro{\vz}{-3}
% \pgfmathsetmacro{\wx}{0}
% \pgfmathsetmacro{\wy}{3}
% \pgfmathsetmacro{\wz}{0}
% \draw[blue] (0,0,0) -- (\ux,\uy,\uz);
% \draw[blue] (0,0,0) -- (\vx,\vy,\vz) -- (\wx+\vx,\wy+\vy,\wz+\vz) -- (\wx,\wy,\wz) -- cycle;  
% \draw[blue,dashed] (\ux,\uy,\uz) -- (\ux+\vx,\uy+\vy,\uz+\vz) -- (\vx,\vy,\vz); 
% \draw[blue,dashed] (\ux+\vx,\uy+\vy,\uz+\vz) -- (\ux+\vx+\wx,\uy+\vy+\wy,\uz+\vz+\wz);
% \draw[blue] (\wx,\wy,\wz) -- (\wx+\ux,\wy+\uy,\wz+\uz) -- (\ux+\vx+\wx,\uy+\vy+\wy,\uz+\vz+\wz) -- (\vx+\wx,\vy+\wy,\vz+\wz);
% \draw[blue] (\ux,\uy,\uz) -- (\wx+\ux,\wy+\uy,\wz+\uz);

% \draw[-stealth,thick] (0,0,0) -- (\ux,\uy,\uz);
% \node[left] at (\ux,\uy,\uz) {$\vec{a}_{1}$};
% \draw[-stealth,thick] (0,0,0) -- (\vx,\vy,\vz);% \node[right];
% \draw[-stealth, dashed] (0,0,0) -- (1.5,0,-3);
% \draw[-stealth, dashed] (1.5,0,-3) -- (0,0,-3);
% \node[above right] at (\vx,\vy,\vz) {$\vec{a}_{2}+0.3\vec{a}_{1}$};
% \draw[-stealth,thick] (0,0,0) -- (\wx,\wy,\wz);% \node[right];
% \draw[-stealth, dashed] (0,0,0) -- (-0.9,3,1.8);
% \draw[-stealth, dashed] (-0.9,3,1.8) -- (0,3,0);
% \node[above left] at (\wx,\wy,\wz) {$\vec{a}_{3}+0.6\vec{a}_2$};
% \draw[|-|,red,thick] (0,-.2,0) -- (\ux,\uy-0.2,\uz);
% \node[red] at (\ux/2,-0.4,0) {$l$};
% \draw[-,red,thick] (-2.5,0,0) -- (-2.5,0,-3);
% \draw[-,red,thick] (-2.65,0,0) -- (-2.35,0,0);
% \draw[-,red,thick] (-2.65,0,-3) -- (-2.35,0,-3);
% \draw[-,red,thick] (-2.5,0,-.3) -- (-2.2,0,-0.3) -- (-2.2,0,0);
% \node[red] at (-2.8,0,-1.5) {$w$};
% \draw[|-|,red,thick] (-3.5,0,-1) -- (-3.5,3,-1);
% \draw[-,red,thick] (-3.5,.3,-1) -- (-3.2,0.3,-1) -- (-3.2,0,-1);
% \node[red] at (-3.7,1.5,-1) {$h$};
% \end{tikzpicture}

\hfill
\begin{tikzpicture}[scale=0.95,line join = round, line cap = round]
\begin{scope}[shift={(0,5,0)}]
\pgfmathsetmacro{\ux}{-5}
\pgfmathsetmacro{\uy}{0}
\pgfmathsetmacro{\uz}{0}
\pgfmathsetmacro{\vx}{1.5}  
\pgfmathsetmacro{\vy}{0}  
\pgfmathsetmacro{\vz}{-3}  
\pgfmathsetmacro{\wx}{1.5}  
\pgfmathsetmacro{\wy}{3}  
\pgfmathsetmacro{\wz}{0}
\draw[blue] (0,0,0) -- (\ux,\uy,\uz);
\draw[blue] (0,0,0) -- (\vx,\vy,\vz) -- (\wx+\vx,\wy+\vy,\wz+\vz) -- (\wx,\wy,\wz) -- cycle;  
\draw[blue,dashed] (\ux,\uy,\uz) -- (\ux+\vx,\uy+\vy,\uz+\vz) -- (\vx,\vy,\vz); 
\draw[blue,dashed] (\ux+\vx,\uy+\vy,\uz+\vz) -- (\ux+\vx+\wx,\uy+\vy+\wy,\uz+\vz+\wz);
\draw[blue] (\wx,\wy,\wz) -- (\wx+\ux,\wy+\uy,\wz+\uz) -- (\ux+\vx+\wx,\uy+\vy+\wy,\uz+\vz+\wz) -- (\vx+\wx,\vy+\wy,\vz+\wz);
\draw[blue] (\ux,\uy,\uz) -- (\wx+\ux,\wy+\uy,\wz+\uz);

\draw[-stealth,thick] (0,0,0) -- (\ux,\uy,\uz);
\node[left] at (\ux,\uy,\uz) {$\vec{a}_{1}$};
\draw[-stealth,thick] (0,0,0) -- (\vx,\vy,\vz);% \node[right];
\node[right] at (\vx,\vy,\vz) {$\vec{a}_{2}$};
\draw[-stealth,thick] (0,0,0) -- (\wx,\wy,\wz);% \node[right];
\node[below right] at (\wx,\wy,\wz) {$\vec{a}_{3}$};

\draw[|-|,red,thick] (0,-.2,0) -- (\ux,\uy-0.2,\uz);
\node[red] at (\ux/2,-0.4,0) {$l$};
\draw[-,red,thick] (-1.5,0,0) -- (-1.5,0,-3);
\draw[-,red,thick] (-1.65,0,0) -- (-1.35,0,0);
\draw[-,red,thick] (-1.65,0,-3) -- (-1.35,0,-3);
\draw[-,red,thick] (-1.5,0,-.3) -- (-1.2,0,-0.3) -- (-1.2,0,0);
\node[red] at (-1.8,0,-1.5) {$w$};
\draw[|-|,red,thick] (-2.5,0,-1) -- (-2.5,3,-1);
\draw[-,red,thick] (-2.5,.3,-1) -- (-2.2,0.3,-1) -- (-2.2,0,-1);
\node[red] at (-2.7,1.5,-1) {$h$};
\end{scope}

\draw[-stealth, thick, bend left = 30] (2,5.5,0) to (2,3.5,0); 
\node[right] at (2.4,4.75,0) {column};
\node[right] at (2.4,4.3,0) {multaddiplacement};
\pgfmathsetmacro{\ux}{-5}
\pgfmathsetmacro{\uy}{0}
\pgfmathsetmacro{\uz}{0}
\pgfmathsetmacro{\vx}{0}
\pgfmathsetmacro{\vy}{0}
\pgfmathsetmacro{\vz}{-3}
\pgfmathsetmacro{\wx}{0}
\pgfmathsetmacro{\wy}{3}
\pgfmathsetmacro{\wz}{0}
\draw[blue] (0,0,0) -- (\ux,\uy,\uz);
\draw[blue] (0,0,0) -- (\vx,\vy,\vz) -- (\wx+\vx,\wy+\vy,\wz+\vz) -- (\wx,\wy,\wz) -- cycle;  
\draw[blue,dashed] (\ux,\uy,\uz) -- (\ux+\vx,\uy+\vy,\uz+\vz) -- (\vx,\vy,\vz); 
\draw[blue,dashed] (\ux+\vx,\uy+\vy,\uz+\vz) -- (\ux+\vx+\wx,\uy+\vy+\wy,\uz+\vz+\wz);
\draw[blue] (\wx,\wy,\wz) -- (\wx+\ux,\wy+\uy,\wz+\uz) -- (\ux+\vx+\wx,\uy+\vy+\wy,\uz+\vz+\wz) -- (\vx+\wx,\vy+\wy,\vz+\wz);
\draw[blue] (\ux,\uy,\uz) -- (\wx+\ux,\wy+\uy,\wz+\uz);

\draw[-stealth,thick] (0,0,0) -- (\ux,\uy,\uz);
\node[left] at (\ux,\uy,\uz) {$\vec{a}_{1}$};
\draw[-stealth,thick] (0,0,0) -- (\vx,\vy,\vz);% \node[right];
\draw[-stealth, dashed] (0,0,0) -- (1.5,0,-3);
\draw[-stealth, dashed] (1.5,0,-3) -- (0,0,-3);
\node[above right] at (\vx,\vy,\vz) {$\vec{a}_{2}+0.3\vec{a}_{1}$};
\draw[-stealth,thick] (0,0,0) -- (\wx,\wy,\wz);% \node[right];
\draw[-stealth, dashed] (0,0,0) -- (1.5,3,0);
\draw[-stealth, dashed] (1.5,3,0) -- (0,3,0);
\node[below left] at (\wx,\wy,\wz) {$\vec{a}_{3}+0.3\vec{a}_1$};
\draw[|-|,red,thick] (0,-.2,0) -- (\ux,\uy-0.2,\uz);
\node[red] at (\ux/2,-0.4,0) {$l$};
\draw[-,red,thick] (-1.5,0,0) -- (-1.5,0,-3);
\draw[-,red,thick] (-1.65,0,0) -- (-1.35,0,0);
\draw[-,red,thick] (-1.65,0,-3) -- (-1.35,0,-3);
\draw[-,red,thick] (-1.5,0,-.3) -- (-1.2,0,-0.3) -- (-1.2,0,0);
\node[red] at (-1.8,0,-1.5) {$w$};
\draw[|-|,red,thick] (-3.5,0,-1) -- (-3.5,3,-1);
\draw[-,red,thick] (-3.5,.3,-1) -- (-3.2,0.3,-1) -- (-3.2,0,-1);
\node[red] at (-3.7,1.5,-1) {$h$};
\end{tikzpicture}

Recall the three basic row operations: scaling, swapping, multaddiplacement. The definition of the determinant how column scaling and column multaddiplacement affect the determinant. This theorem covers column swaps:

\begin{theorem}[Antisymmetry.]\label{theorem:antisymmetry}

If you swap two column vectors, the determinant changes sign. Writing the column number underneath the vector:
\[ \det( \vec{a}_1, \ldots, \underset{i}{\vec{a}_i}, \ldots, \underset{j}{\vec{a}_j},\ldots,\vec{a}_c )
= - \det( \vec{a}_1, \ldots, \underset{i}{\vec{a}_j}, \ldots, \underset{j}{\vec{a}_i},\ldots,\vec{a}_c ) \]
\end{theorem}

\begin{notes}
\vfill
\vfill
\end{notes}

\begin{book}
\textbf{Proof.} Each of the following are equal because of preservation under column multaddiplacement:
\begin{align*}
\det( \vec{a}_1, \ldots, \underset{i}{\vec{a}_i}, \ldots, \underset{j}{\vec{a}_j},\ldots,\vec{a}_c ) 
&=  \det(\vec{a}_1, \ldots, \underset{i}{\vec{a}_i+\vec{a}_j}, \ldots, \underset{j}{\vec{a}_j}, \ldots,\vec{a}_c) \\
&=  \det(\vec{a}_1, \ldots, \underset{i}{\vec{a}_i+\vec{a}_j}, \ldots, \underset{j}{\vec{a}_j-(\vec{a}_i+\vec{a}_j)}, \ldots,\vec{a}_c) \\
&=  \det(\vec{a}_1, \ldots, \underset{i}{\vec{a}_i+\vec{a}_j}, \ldots, \underset{j}{-\vec{a}_i}, \ldots,\vec{a}_c) \\
&=  \det(\vec{a}_1, \ldots, \underset{i}{\vec{a}_i+\vec{a}_j-\vec{a}_i}, \ldots, \underset{j}{-\vec{a}_i}, \ldots,\vec{a}_c) \\
&=  \det(\vec{a}_1, \ldots, \underset{i}{\vec{a}_j}, \ldots, \underset{j}{-\vec{a}_i}, \ldots,\vec{a}_c)
\end{align*}
And finally, the scalar $-1$ can be factored out using linearity:
\[ = - \det(\vec{a}_1, \ldots, \underset{i}{\vec{a}_j}, \ldots, \underset{j}{\vec{a}_i}, \ldots,\vec{a}_c)\]
\blacksquare
\end{book}

\begin{lemma}
\label{theorem:zero-determinants}
If $A$ is a square matrix, with $\vec{0}$ in a column, then $\det(A) = 0$.
\end{lemma}

This is intuitively consistent with the determinant representing area in $\R^{2}$: if $\vec{0}$ is one of the two vectors, then the parallelogram is just a line segment with zero area. In $\R^3$, if $\vec{0}$ is one of the three vectors, then the parallelepiped is a 2D parallelogram with zero volume.

\textbf{Proof:} Using linearity (scaling):
\begin{align*}
 \det( \vec{a}_1, \ldots, \vec{0}, \ldots,\vec{a}_c ) &=  \det( \vec{a}_1, \ldots, 0\vec{0}, \ldots,\vec{a}_c ) \\
&=  0\det( \vec{a}_1, \ldots, \vec{0}, \ldots,\vec{a}_c ) = 0
\end{align*}
\blacksquare

\begin{notes}
\columnbreak
\end{notes}

\begin{definition}
The \textbf{main diagonal} of a matrix $A = \begin{bmatrix}
a_{11} & a_{12} & \cdots & a_{1c} \\
a_{21} & a_{22} & \cdots & a_{2c} \\
\vdots & \vdots & \ddots & \vdots \\
a_{r1} & a_{r2} & \cdots & a_{rc}
\end{bmatrix}$ is all the entries of the form $a_{ii}$; that is, the row number equals the column number. A \textbf{diagonal matrix} is a square matrix where all the entries \emph{not} on the main diagonal are zero. (For example, $I$ is diagonal.) An \textbf{upper triangular matrix} is a square matrix where all the entries \emph{below} the main diagonal are zero. A \textbf{lower triangular matrix} is a square matrix where all the entries \emph{above} the main diagonal are zero. A \textbf{triangular matrix} is either upper or lower triangular. 
\end{definition}

Soon we'll find an easy shortcut for calculating the determinants of triangular matrices; but first, let's practice using the properties in the definition. The idea is to use column multaddiplacements, swaps, and scaling to transform $A$ to $I$ which has determinant 1.

\example \label{example:3.1col-reduction} Find $\det(A)$ for $A =\begin{bmatrix}
1 & 2 & 3 \\
0 & 4 & 5 \\
0 & 0 & 6
\end{bmatrix}$.

\begin{notes}
\vfill \hspace{1pt}
\columnbreak
\end{notes}
\begin{book}
\textbf{Solution.} First we'll do some column multaddiplacements that don't change the determinant. Note that because $A$ is upper triangular, these operations don't change the main diagonal.
\begin{align*}
\det(A) &= \det \begin{bmatrix}
1 & 0 & 0 \\
0 & 4 & 5 \\
0 & 0 & 6
\end{bmatrix} & \begin{matrix}c_2\leftarrow c_2-2c_1 \\ c_3 \leftarrow c_3-3c_1 \end{matrix} \\
 &= \det \begin{bmatrix}
1 & 0 & 0 \\
0 & 4 & 0 \\
0 & 0 & 6
\end{bmatrix} & \begin{matrix}c_3\leftarrow c_3-\frac{5}{4}c_2 \end{matrix}
\end{align*}
Now the matrix is the result of scaling two columns of $I$ by 4 and 6, so by linearity (scaling):
\[ 
\det(A) = (4)\det \begin{bmatrix}
1 & 0 & 0 \\
0 & 1 & 0 \\
0 & 0 & 6
\end{bmatrix} = (4)(6)\det \begin{bmatrix}
1 & 0 & 0 \\
0 & 1 & 0 \\
0 & 0 & 1
\end{bmatrix} = 4(6)(1) = 24\]
\blacksquare
\end{book}

\example \label{example:3.1-between1-2} Is $A = \begin{bmatrix}
0 & 1 \\
1 & -3
\end{bmatrix}$ triangular? Find $\det(A)$.

\begin{notes}
\vfill
\end{notes}
\begin{book}
\textbf{Solution.} Below the main diagonal of $A$ is the 1 in the bottom left. Above the main diagonal is the $1$ on the top right. So $A$ is not triangular. To find $\det(A)$, first we can clear the 3 with a column multaddiplacement that doesn't change the determinant, and then do a swap:
\begin{align*}
\det(A) &= \det \begin{bmatrix}
0 & 1 \\
1 & 0
\end{bmatrix} &c_2\leftarrow c_2+3c_1 \\
 &= -\det \begin{bmatrix}
1 & 0  \\
0 & 1 
\end{bmatrix} &\text{(antisymmetry)} \\
&= -1
\end{align*}
\blacksquare
\end{book}

\example \label{example:3.1-2} Find $\det(A)$ for $A=\begin{bmatrix}
3 & 2 & -6 & 3 \\
0 & 2 & 4 & 5 \\
0 & 0 & 0 & 7 \\
0 & 0 & 0 & 1 
\end{bmatrix}$.

\begin{notes}
\vfill
\end{notes}
\begin{book}
\textbf{Solution.} Note that $A$ is triangular with a $0$ in the main diagonal. This suggests that third column is in the span of the first two columns. Or in other words the ``height'' of $\vec{a}_3$ above $\spn(\vec{a}_1,\vec{a}_2)$ is zero, which should make the volume of $\pgram A$ zero, and $\det(A)$ zero too. Let's prove it. 

We can make the third column $\vec{0}$ with some choice column multaddiplacements:
\begin{align*}
\det(A) &= \det \begin{bmatrix}
3 & 0 & 0 & 3 \\
0 & 2 & 4 & 5 \\
0 & 0 & 0 & 7 \\
0 & 0 & 0 & 1 
\end{bmatrix} & \begin{matrix}c_2\leftarrow c_2-\frac{2}{3}c_1 \\ c_3 \leftarrow c_3+2c_1 \end{matrix} \\
&= \det \begin{bmatrix}
3 & 0 & 0 & 3 \\
0 & 2 & 0 & 5 \\
0 & 0 & 0 & 7 \\
0 & 0 & 0 & 1 
\end{bmatrix} & \begin{matrix}c_3\leftarrow c_3-2c_2 \end{matrix}
\end{align*}
Now that there's a $\vec{0}$ column, lemma \ref{theorem:zero-determinants} shows $\det(A)=0$.
\end{book}

\begin{lemma} \label{theorem:detA-triangular-zero}
 If $A$ is triangular with a $0$ entry in the main diagonal, then $\det(A) = 0$.
\end{lemma}

Intuitively, the $0$ in the main diagonal suggests that $A$ is not row-equivalent to $I$, and that the columns are linearly dependent, and like the last example, there's a sequence of column multaddiplacements that turn that column to $\vec{0}$, making $\det(A) = 0$.  
\begin{notes}
The proof in the book generalizes the idea of example \ref{example:3.1-2}.

\columnbreak
\end{notes}

\begin{book}
\textbf{Proof.} Suppose $A$ is triangular with $0$ somewhere in the main diagonal. First suppose $A$ is upper triangular. If $a_{11} = 0$ then the first column of $A$ equals $\vec{0}$ and by theorem \ref{theorem:zero-determinants} $\det(A) = 0$.

Otherwise, let $a_{ii}$ be the upper-left-most 0 on the main diagonal (in case there's more than one zero). This means that $a_{11}, a_{22}, \ldots, a_{(i-1)(i-1)}$ are all non-zero. Like in the last example, adding multiples of column 1, $\begin{bmatrix}
a_{11} \\ 0 \\ \vdots \\ 0
\end{bmatrix}$ to the other columns can make all the other row 1 entries zero without changing the determinant. Now that column 2 is $\begin{bmatrix}
0 \\ a_{22}\\0 \\\vdots \\0
\end{bmatrix}$, we can add multiples of column 2 to the other columns to make all the other row 2 entries zero without changing the determinant. Continue until you've made everything right of the main diagonal in rows 1 through $i-1$ equal to zero. Since $a_{ii}$ and below was already zero, column $i$ now equals $\vec{0}$, so the determinant of this transformed matrix, which equals $det(A)$, is 0.

If $A$ is lower triangular, let $a_{ii}$ be the bottom-right-most 0 on the main diagonal. A similar process can create all zeros \emph{left} of the main diagonal in all the rows \emph{below} $i$, which makes column $i$ equal to $\vec{0}$. \blacksquare
\end{book}

\begin{theorem}[Determinants of triangular matrices.]
\label{theorem:det-of-triangular}
If $A$ is triangular, then $\det(A)$ is the product of the entries on the main diagonal.
\end{theorem}

\begin{notes}
The proof, which you can find in the book, generalizes the idea in examples \ref{example:3.1col-reduction} and \ref{example:3.1-2}: The zeros on one side of the diagonal make it so column multaddiplacements transform $A$ to a diagonal matrix, and linearity transforms that to $I$.
\end{notes}

\begin{book}
\textbf{Proof.} First, if one of the main diagonal entries is zero, then product of the diagonal will be zero; and the lemma above shows that $\det(A)=0$ too, proving the theorem if zero is on the diagonal.

Now suppose $A$ is diagonal with no zero entries on the diagonal. In this case the columns of $A$ are scalar multiples of the columns of $I$ and we can use linearity (scaling):
\begin{align*}
\det\begin{bmatrix}
a_{11} & 0 & \cdots & 0 \\
0 & a_{22} & \cdots & 0 \\
\vdots & \vdots & \ddots & \vdots \\
0 & 0 & \cdots & a_{cc}
\end{bmatrix} &= a_{11}\det\begin{bmatrix}
1 & 0 & \cdots & 0 \\
0 & a_{22} & \cdots & 0 \\
\vdots & \vdots & \ddots & \vdots \\
0 & 0 & \cdots & a_{cc}
\end{bmatrix}\\
&= a_{11}a_{22}\det\begin{bmatrix}
1 & 0 & \cdots & 0 \\
0 & 1 & \cdots & 0 \\
\vdots & \vdots & \ddots & \vdots \\
0 & 0 & \cdots & a_{cc}
\end{bmatrix}\\ 
&= \cdots \\ 
&= a_{11}a_{22} \cdots a_{cc}\det(I) \\ 
&= a_{11}a_{22} \cdots a_{cc}(1)
\end{align*}
This proves the theorem for diagonal matrices. Next we'll show that any triangular matrix can be transformed to a diagonal matrix with column multaddiplacements, which don't change the determinant.

Suppose $A$ is upper triangular, with no zeros in the diagonal. Since $a_{11}\ne 0$, we can add multiples of $\vec{a}_1= \begin{bmatrix}
a_{11} \\ 0 \\ \vdots \\ 0
\end{bmatrix}$ to the other columns to make the rest of the first row entries zero, without changing the determinant or any other rows. Then we can add multiples of the second column, which is now $\begin{bmatrix}
0 \\ a_{22} \\ 0 \\ \vdots \\ 0
\end{bmatrix}$ to make the rest of the second row entries zero, without changing the determinant or any other rows. Since every diagonal entry is non-zero, we can do this in every row and end up with a diagonal matrix with the same determinant as $A$; and we've already proved the theorem for diagonal matrices.

So, we've proved it for triangular matrices with 0 in the diagonal, and for diagonal and upper triangular matrices without zeros in the diagonal. The proof that it works for lower triangular matrices without zeros in the diagonal is exercise \ref{exercise:3.1-LT-proof}. \blacksquare
\end{book}

\example Find $\det(A)$ if $A = \begin{bmatrix}
3 & -1 & 2 \\
4 & 0 & 4 \\
-5 & 2 & -3
\end{bmatrix}$. Note that $\vec{a}_{3}=\vec{a}_1+\vec{a}_2$.

\begin{notes}
\vfill
\end{notes}
\begin{book}
\textbf{Solution.} $\vec{a}_{3}=\vec{a}_1+\vec{a}_2$ means $\vec{a}_3-\vec{a}_1-\vec{a}_2=\vec{0}$, so we can use column multaddiplacements to clear out column 3 without changing the determinant:
\begin{align*}
\det(A) &= \det \begin{bmatrix}
3 & -1 & -1 \\
4 & 0 & 0 \\
-5 & 2 & 2
\end{bmatrix} &c_3\leftarrow c_3-c_1 \\
&= \det \begin{bmatrix}
3 & -1 & 0 \\
4 & 0 & 0 \\
-5 & 2 & 0
\end{bmatrix} &c_3\leftarrow c_3-c_2  \\
&= 0 &\text{since it has } \vec{0}\text{ in a column.}
\end{align*}
\blacksquare
\end{book}

The next theorem gives a way to check if a matrix is invertible or not which will be important when we study eigenvalues.

\begin{notes}
\columnbreak
\end{notes}
\begin{theorem}[Determinants and invertibility.]
\label{theorem:determinants-and-invertibility}

Suppose $A$ is square. Then $\det(A) \ne 0$ if and only $A$ is invertible.
\end{theorem}

\begin{notes}
\vfill
\end{notes}
\begin{book}
\textbf{Proof.} Consider $A^T$, and let $B^T$ be the reduced row echelon form of $A^T$. So, there's a sequence of row operations that reduce $A^T$ to $B^T$. And every one of these row operations is the same as a column operation that transforms $A$ to $B$. The effects of these operations are a series of the following:
\begin{itemize}
\item Row multaddiplacement with $A^T$ (column multaddiplacement with $A$) doesn't change the determinant.
\item Scaling a row of $A^T$ by a non-zero scalar (scaling a column of $A$ by a non-zero scalar) multiplies $\det(A)$ by the non-zero scalar. Recall that scaling a row by zero is not a valid row operation.
\item Swapping rows of $A^T$ (swapping columns of $A$) multiplies the determinant by $-1$.
\end{itemize}
Thus:
\[(\text{product of some non-zero scalars})\det(A) = \det(B)\]% = \text{product of diagonal entries in }B
or equivalently:
\begin{equation}
\label{eq:detAtoB}
\det(A) = \frac{\det(B)}{\text{product of some non-zero scalars}}
\end{equation}
Now, the following statements are equivalent:
\begin{itemize}
\item $A$ is invertible.
\item $A^T$ is invertible (\ref{theorem:properties-of-inverses}).
\item $B^T = I$ (IMT).
\item $B = I$, since $B=( B^T )^T = I^T = I$.
\item $\det(B) \ne 0$ since $\det(I) = 1 \ne 0$, and if $B^T \ne I$, then $B^T$ must have a row of zeros because it's square in echelon form, and this means $B$ has a column of zeros making $\det(B) = 0$.
\item $\det(A) \ne 0$ by equation \ref{eq:detAtoB}.
\end{itemize}
\blacksquare
\end{book}


This section has been a lot. Keeping in mind that determinants are only defined for square matrices, let's summarize everything we'll need for later in one place:

\begin{bbox}
\textbf{Definition:}
\begin{itemize}
\item \textbf{Normality:} $\det(I) = 1$.
\item \textbf{Linearity (scaling):} Multiplying a column by a scalar also scales the determinant.

\item \textbf{Linearity (addition):} 
\[\det(\vec{a}_1, \ldots, \vec{a}_j+\vec{b}_j,\ldots, \vec{a}_c) = \det(\vec{a}_1, \ldots, \vec{a}_j,\ldots, \vec{a}_c) + \det(\vec{a}_1, \ldots, \vec{b}_j,\ldots, \vec{a}_c)\]
\item \textbf{Preservation under column multaddiplacement:} Replacing a column with the sum of itself and a different column does not change the determinant.

\end{itemize}

\textbf{Theorem \ref{theorem:antisymmetry} (antisymmetry):} Swapping two columns changes the sign of the determinant.

\textbf{Theorem \ref{theorem:det-of-triangular}:} If $A$ is triangular, then $\det(A)$ is the product of the entries on the main diagonal.

\textbf{Theorem \ref{theorem:determinants-and-invertibility}:} $\det(A) \ne 0$ if and only if $A$ is invertible.

\textbf{Theorem \ref{theorem:2x2-det}:}  $\det\begin{bmatrix}
a & b \\
c & d
\end{bmatrix} = ad-bc.$ (to be proven in exercises.)
\end{bbox}

\exerciseheading

For exercises \ref{exercise:3.1-basic-start} to \ref{exercise:3.1-basic-end} find $\det(A)$ by using column operations to transform $A$ into a triangular matrix.

\exatend \label{exercise:3.1-basic-start} $A=\begin{bmatrix}
3 & 6 \\
1 & 2
\end{bmatrix}$

\exatend $A=\begin{bmatrix}
3 & 6 \\
1 & 1
\end{bmatrix}$

\exatend $A=\begin{bmatrix}
1 & 1 \\
3 & 6 
\end{bmatrix}$

\exatend $A = \begin{bmatrix}
\frac{\sqrt{2}}{2} & \frac{-\sqrt{2}}{2} \\
\frac{\sqrt{2}}{2} & \frac{\sqrt{2}}{2} \\
\end{bmatrix}$, the matrix for a $45^{\circ}$ rotation.

\exatend $A = \begin{bmatrix}
3 & 2 & 1 \\
5 & 4 & 0 \\
6 & 0 & 0 \\
\end{bmatrix}$

\exatend $A = \begin{bmatrix}
1 & -1 & 3 \\
2 & -2 & 3 \\
0 & 1 & -3
\end{bmatrix}$

\exatend $A = \begin{bmatrix}
1 & -2 & 2 & -2 \\
3 & -5 & 6 & -6 \\
-2 & 7 & -3 & 4 \\
0 & -3 & 5 & 0
\end{bmatrix}$

\exatend 
$\begin{bmatrix}
0 & 0 & 1 & 0 & 0 & 0 \\
0 & 1 & 0 & 0 & 0 & 0 \\
1 & 0 & 0 & 0 & 0 & 0 \\
0 & 0 & 0 & 0 & 1 & 0 \\
0 & 0 & 0 & 1 & 0 & 0 \\
0 & 0 & 0 & 0 & 0 & 1 \\
\end{bmatrix}$

\exatend $\begin{bmatrix}
0 & 0 & 0 & a \\
0 & 0 & b & 0 \\
0 & c & 0 & 0 \\
d & 0 & 0 & 0
\end{bmatrix}$

\exatend \label{exercise:3.1-basic-end} $\begin{bmatrix}
0& 0 & 0 & 0 & a \\
0&0 & 0 & b & 0 \\
0&0 & c & 0 & 0 \\
0&d & 0 & 0 & 0 \\
e&0 & 0 & 0 & 0 
\end{bmatrix}$



\exatend (T/F?) Since $\vec{e}_1=(1,0)$ and $\vec{e}_2=(0,1)$ lie flat on the plane, $\pgram \begin{bmatrix}
\vec{e}_1 & \vec{e}_2
\end{bmatrix}$ has zero volume and so $\det(\vec{e}_1,\vec{e}_2) = 0$.

\exatend (T/F?) The matrix $\begin{bmatrix}
1 & 2 & 3 \\
4 & 5 & 0 \\
6 & 0 & 0
\end{bmatrix}$ is triangular.

\exatend (T/F?) If $A$ is a $3\times 3$ matrix, $\det(2A) = 2\det(A)$. \hintfootnote{Think about a rectangular box whose length, width, and height are all doubled.}

\exatend (T/F?) Triangular matrices can't be square matrices, because, like, triangles and squares are, like, different shapes.

\exatend (T/F?) If $A$ is invertible, then $\det(A) = 0$.

\exatend (T/F?) The linearity (addition) property in the definition says that $\det(A+B) =\det(A) + \det(B)$.

\exatend\begin{enumerate}[label=(\alph*)]
  \item  Draw the parallelogram with vertices at (0,0), (1,2), (4,1), and (5,3). Note that this is the parallelogram determined by vectors $\bbvec{1}{2}$ and $\bbvec{4}{1}$. Find its area.
  \item Suppose you didn't draw the parallelogram first and thought it was determined by vectors $\bbvec{1}{2}$ and $\bbvec{5}{3}$. Explain how you know this mistake doesn't matter because the parallelogram has the same area as the correct one.
\end{enumerate}


%credit: david austin
\exatend Suppose $T(\vec{x}) = A\vec{x}$ is a transformation that reflects points across some line through the origin. For some arbitrary line, draw a picture of $T(\vec{e}_1)$ and $T(\vec{e}_2)$. Thinking about area of $\pgram A$, what is $\det(A)$? \creditfootnote{Inspired by an exercise in David Austin's excellent \emph{Understanding Linear Algebra}.}

%credit: david austin
\exatend Suppose $T(\vec{x}) = A\vec{x}$ is a transformation that projects all points to the nearest point on some line through the origin. For some arbitrary line, draw a picture of $T(\vec{e}_1)$ and $T(\vec{e}_2)$. Thinking about area of $\pgram A$, what is $\det(A)$? \creditfootnote{Inspired by D. Austin.}

\exatend Use the antisymmetry theorem (but no other theorems/lemmas from this section) to show that if two columns of $A$ are equal, then $\det(A) = 0$.

\exatend Use the column replacement property of the definition, and lemma \ref{theorem:zero-determinants} (but no other theorems/lemmas from this section) to show that if two columns of $A$ are equal, then $\det(A) = 0$.

\exatend Use theorem \ref{theorem:determinants-and-invertibility} (but no other theorems/lemmas from this section) to show that if two columns of $A$ are equal, then $\det(A) = 0$.


\exatend Prove theorem \ref{theorem:2x2-det} below by following these steps:
\begin{enumerate}[label=(\alph*)]
\item Show the theorem is true if $a\ne0$. Use column multaddiplacement to turn $\begin{bmatrix}
a & b \\
c & d
\end{bmatrix}$ into a triangular matrix.
\item Show the theorem is true if $a = 0$. Use a swap to turn $\begin{bmatrix}
a & b \\
c & d
\end{bmatrix}$ into a triangular matrix. (This completes the proof.)
\item Explain how theorems \ref{theorem:2x2-inverses}, \ref{theorem:determinants-and-invertibility}, and \ref{theorem:2x2-det} agree with each other.
\end{enumerate}

\begin{theorem}\label{theorem:2x2-det}
$\det\begin{bmatrix}
a & b \\
c & d
\end{bmatrix} = ad-bc.$
\end{theorem}

\exatend Suppose $(p,y_p)$ and $(q,y_q)$ two points in $\R^2$, and you want to find $b,m$ such that the graph of $y = b+mx$ includes the two points. 
\begin{enumerate}[label=(\alph*)]
\item Explain why finding the parabola that includes the three points is equivalent to solving $A\bbvec{b}{m} = \bbvec{y_p}{y_q}$ for $(b,m)$, where $A=\begin{bmatrix}
1 & p  \\
1 & q 
\end{bmatrix}$.
\item Use column multaddiplacement to show $\det(A) = q-p$.
\item What needs to be true about the points $(p,y_p)$ and $(q,y_q)$ to guarantee there's a unique linear function that passes through the points?
\end{enumerate}

%credit: david austin
\exatend \emph{Challenge:} Suppose $(p,y_p),(q,y_q),$ and $(r,y_r)$ are three points in $\R^2$, and you want to find $a,b,c$ such that the graph of $y = c+bx+ax^2$ includes the three points. \creditfootnote{Inspired by D. Austin.}
\begin{enumerate}[label=(\alph*)]
\item Explain why finding the parabola that includes the three points is equivalent to solving $A\bbbvec{c}{b}{a} = \bbbvec{y_p}{y_q}{y_r}$ for $(c,b,a)$, where $A=\begin{bmatrix}
1 & p & p^2 \\
1 & q & q^2 \\
1 & r & r^2
\end{bmatrix}$.
\item Use column multaddiplacement to show $\det(A) = (q-p)(r-p)(r-q)$.
\item What needs to be true about the points $(p,y_p),(q,y_q),$ and $(r,y_r)$ to guarantee there's a unique quadratic function that passes through all three?
\end{enumerate}

\exatend \label{exercise:3.1-LT-proof} Show that if $A$ is lower triangular, with no zeros in the main diagonal, then $\det(A)$ is the product of the diagonal entries of $A$. This will finish the proof of theorem \ref{theorem:det-of-triangular}. You may use the fact that the theorem works with diagonal matrices, which was already proven.

\columnbreak
\section{More Determinants and the Cofactor Expansion}
\label{sec:cofactor-expansion}

\exercise (\emph{Warm-up}) For each of the following column operations, write your own example of a $3\times3$ elementary matrix that is the result doing the same operation to $I$. Find the determinant of each.
\begin{enumerate}[label=(\alph*)]
\item Scaling a column.
\item Column multaddiplacement.
\item Swapping a column.
\end{enumerate}

\begin{notes}
\vfill
\end{notes}

\exercise (\emph{Warm-Up}) Let $A = \begin{bmatrix}
2 & 4 & \pi \\
0 & 3 & 7 \\
0 & 0 & 1
\end{bmatrix}$ and $E = \begin{bmatrix}
0 & 1 & 0 \\
1 & 0 & 0 \\
0 & 0 & 1
\end{bmatrix}$. 
\begin{enumerate}[label=(\alph*)]
\item Find $AE$. What do you notice?
\item Find $\det(A), \det(E)$, and $\det(AE)$. What do you notice?
\end{enumerate}

\begin{notes}
\vfill \hspace{1pt}
\columnbreak
\end{notes}

\subsection{Some Nice Theorems About Determinants}\label{subsec:some-nice-theorems-about-determinants}

\example \label{example:3.1-AE-column-op} Let $A = \begin{bmatrix}
a & 0 & 0 \\
b & d & 0 \\
c & e & f
\end{bmatrix}$, and $E = \begin{bmatrix}
1 & 0 & 0 \\
0 & 1 & 0 \\
0 & \frac{-e}{f} & 1
\end{bmatrix}$. Note that $E$ is an elementary matrix, the result of doing the row operation $r_3\leftarrow r_3+\frac{-e}{f}r_2$ to $I$; and $E$ is also the result of doing the \emph{column} operation $c_2\leftarrow c_2+\frac{-e}{f}c_3$. 
\begin{enumerate}[label=(\alph*)]
\item Find $AE$. What do you notice?
\item Find $\det(A), \det(E)$, and $\det(AE)$. What do you notice?
\end{enumerate}

\begin{notes}
\vfill \hspace{1pt}
\columnbreak
\end{notes}
\begin{book}
\textbf{Solution.} Using the definition of matrix multiplication:
\[AE=\begin{bmatrix}
A \bbbvec{1}{0}{0} & A \bbbvec{0}{1}{\frac{-e}{f}} & A \bbbvec{0}{0}{1}
\end{bmatrix}
= \begin{bmatrix}
\vec{a}_1 & \vec{a}_2+\frac{-e}{f}\vec{a}_3 & \vec{a}_3
\end{bmatrix}
= \begin{bmatrix}
a & 0 & 0 \\
b & d & 0 \\
c & 0 & f
\end{bmatrix}\]
Notice that $AE$ is the result of doing the same column operation to $A$, that was done to get $E$ from $I$.

Since $A,E$ and $AE$ are (lower) triangular, $\det(A) = adf$, $\det(E) = 1^3=1$, and $\det(AE) = adf$. Notice that $\det(AE) = \det(A)\det(E)$.
\end{book}

\begin{lemma}\label{theorem:lemma-column-operation}
 If $A$ is any $r\times c$ matrix, and $E$ is an $c\times c$ elementary matrix that is the result of doing some \emph{column} operation to $I$, then $AE$ is the result of doing the same \emph{column} operation to $A$.
\end{lemma}

Note that this lemma is the same as theorem \ref{theorem:elementary-operations}, but about $AE$ instead of $EA$, and column instead of row operations.

\begin{notes}
\vfill
\end{notes}
\begin{book}
\textbf{Proof.} Note that:
\begin{equation}
\label{eq:3.1-lemma-proof}
AE = (E^TA^T)^T
\end{equation}
Since the rows of $E^T$ are the columns of $E$, the column operation that turns $I$ into $E$ is equivalent to the row operation that turns $I^T = I$  into $E^T$. For example, if the $I \rightarrow E$ operation is the multaddiplacement $c_3\leftarrow c_3-4c_1$, then the equivalent row operation is $r_3\leftarrow r_3-4r_1$, reflected in $E^T$. So $E^TA^T$ is the result of applying that operation to the rows of $A^T$, which are columns of $A$, and thus $(E^TA^T)^T$ is the result of applying that operation to the columns of $A$. \blacksquare
\end{book}

\begin{lemma}
\label{theorem:lemma-dAE-dAdE}

If $A$ is square and $E$ is elementary (of the same size) then: \[\det(AE) = \det(A)\det(E)\]
\end{lemma}

\begin{notes}
\vfill \hspace{2pt}
\columnbreak
\end{notes}
\begin{book}
\textbf{Proof.} Let's refer to the column operation that turns $I$ to $E$ as ``the operation.'' So, by lemma \ref{theorem:lemma-column-operation}, $AE$ is the result of applying the operation to $A$. And $\det(AE)$ is $\det(A)$ multiplied by a factor of either $-1$, if the operation is a swap, or a scalar, if the operation is scaling, or a factor of $1$ (no change), if the operation is multaddiplacement. That factor equals $\det(E)$ since that same operation changed $\det(I)=1$ to $\det(E)$. \blacksquare
\end{book}

The next theorem should be surprising, because the vectors that make $\pgram A^T$ are very different from the vectors that make $\pgram A$.

\begin{theorem}
\label{theorem:detA-detAT} If $A$ is square, then $\det(A) = \det(A^T)$.
\end{theorem}

\begin{notes}
\vfill
\end{notes}
\begin{book}
\textbf{Proof.} If $A$ is not invertible, then neither is $A^T$, and $\det(A) = 0 = \det(A^T)$. So we just need to prove it for invertible matrices. Assume $A$ is invertible. Then by the IMT, $A$ is row-equivalent to $I$. Row-operations are reversible, so there's a sequence of row operations that transform $I$ to $A$, and each of these row operations has a corresponding elementary matrix that does the operation when multiplied by $I$ on the left:
\[ A =E_n\cdots E_2E_1I = E_n\cdots E_2E_1 \] 
Taking the determinant and applying lemma \ref{theorem:lemma-dAE-dAdE} $n$ times gives:
\begin{equation}\label{eq:3.1-product-of-elem}
\det(A) = \det(E_n\cdots E_2E_1) = \det(E_n)\cdots \det(E_2)\det(E_1)
\end{equation}
Now let's do the same thing with $A^T$:
%\begin{align}\label{eq:3.1-product-of-ET}
\[A^{T} =  ( E_n\cdots E_2E_1 )^T=E_1^TE_2^T\cdots E_n^T\]% \\ \nonumber
\begin{equation}\label{eq:3.1-product-of-ET}
\det(A^T) = \det(E_1^TE_2^T\cdots E_n^T) = \det(E_1^T)\det(E_2^T)\cdots\det(E_n^T) 
\end{equation}
%\end{align}
Now, to show equations \ref{eq:3.1-product-of-elem} and \ref{eq:3.1-product-of-ET} are the same, we just have to notice that every kind of elementary matrix and its transpose have the same determinant:
\begin{itemize}
\item If $E$ results from multaddiplacing a row of $I$, then $E$ is triangular, and so is $E^T$, with 1's on the main diagonal, so $\det E = 1 = \det E^T$.
\item If $E$ results from swapping two rows of $I$, then $E$ is symmetrical: $E =E^T$. (Check with a couple examples to see why.) So, $\det(E) = \det(E^T)$. %TODO: maybe find a better way to write this...
\item If $E$ results from scaling a row of $I$, then $E$ is diagonal so $E = E^T$ and $\det(E) = \det(E^T)$.
\end{itemize}
\blacksquare
\end{book}

\begin{theorem}
\label{theorem:detAB-detAdetB}

If $A$ and $B$ are $c\times c$ matrices, then $\det(AB) = \det(A)\det(B)$.
\end{theorem}

\begin{notes}
\vfill \hspace{1pt}
\columnbreak
\end{notes}
\begin{book}
\textbf{Proof.} If $B$ is not invertible, then exercise \ref{exercise:2.4-ABinv-Ainv} in section \ref{sec:elementary-matrices-and-IMT} implies $AB$ is not invertible, which makes both sides of the equation zero. So we can assume $B$ is invertible. Then by the IMT, $B$ is row-equivalent to $I$. Row-operations are reversible, so there's a sequence of row operations that transform $I$ to $B$, and each of these row operations has a corresponding elementary matrix that does the operation when multiplied by $I$ on the left:
\[ B =E_n\cdots E_2E_1I = E_n\cdots E_2E_1 \]
Plugging this in for $B$:% and applying lemma \ref{theorem:lemma-dAE-dAdE} repeatedly in each step:
\begin{align*}
 \det(AB) &= \det(AE_n\cdots E_2E_1) \\
&= \det(A)\det(E_n)\cdots\det(E_2)\det(E_1) & (\text{lemma \ref{theorem:lemma-dAE-dAdE}, }n\text{ times}) \\
&= \det(A)\det(E_n \cdots E_2E_1) & (\text{lemma \ref{theorem:lemma-dAE-dAdE}, }n-1\text{ times}) \\
&= \det(A)\det(B)
\end{align*}
\blacksquare
\end{book}


\textbf{Remarks: Existence and Uniqueness of Determinant}

The definition of the determinant in section \ref{sec:signed-volume} gives four properties that a function should have to be called a determinant. But we haven't shown that there actually exists a function that satisfies all those properties. And, if there does exist a determinant function, we haven't shown that it's unique. What we have done throughout this chapter is \emph{implicitly assume} that we have a function, $\det()$, that satisfies the definition properties, and showed that some other things about $\det()$ must be true. For example the wording for theorem \ref{theorem:determinants-and-invertibility} should have been:

\textbf{Theorem \ref{theorem:determinants-and-invertibility} (precise version):} Suppose $\det()$ is a function which satisfies the definition of a determinant, and $A$ is square. Then $\det(A) \ne 0$ if and only if $A$ is invertible.

But instead we went with the following instead:

\textbf{Theorem \ref{theorem:determinants-and-invertibility} (readable version):} Suppose $A$ is square. Then $\det(A) \ne 0$ if and only if $A$ is invertible.

Note that we did come up with a way of calculating $\det(A)$, by column-reducing $A$ to a triangular matrix, keeping track of how the operations change the determinant, and then multiplying those factors by the diagonal entries. But it's not obvious this is a well-defined function because there are many triangular matrices equivalent to $A$. Fortunately, all is well and good, as the next two theorems show. 

You can find the proof of the theorem below in Sergei Treil's \emph{Linear Algebra Done Wrong} section 3.4 (the book that inspired this book's whole approach to determinants and is available free online). It's quite interesting and illuminating, but is long and uses some math that is slightly beyond the scope of this course..

\begin{theorem}
\label{theorem:existence-of-determinants} A function that satisfies the definition of the determinant exists.
\end{theorem}

In the next subsection we show the determinant is unique and find a formula for it.

\subsection{The Cofactor Expansion}\label{subsec:the-cofactor-expansion}

First some notation we'll use in the big theorem:
$A_{ij}$ is the result of removing row $i$ and column $j$ from $A$. For example, if $A = \begin{bmatrix}
  1 & 2 & 3 \\
  4 & 5 & 6 \\
  7 & 8 & 9
\end{bmatrix}$ then $A_{12} = \begin{bmatrix}
  4 & 6 \\
  7 & 9
\end{bmatrix}$. 

\begin{theorem}[Cofactor Expansion.] \label{theorem:cofactor-expansion}

Suppose $\det$ is a function satisfying the definition of a determinant. Define a function $d: \left\{c \times c\text{ matrices}\right\} \rightarrow \R$ by the formula:
\begin{align*}
 d(A) &= (-1)^{1+1}a_{11}\det(A_{11}) + (-1)^{2+1}a_{21}\det(A_{21}) + \\
& \hspace{8.3em} + \cdots + (-1)^{c+1}a_{c1}\det(A_{c1}) \\
&= \sum_{k=1}^c(-1)^{k+1}a_{k1}d(A_{k1})
\end{align*}
\begin{enumerate}[label=(\roman*)]
\item $\det(A) = d(A)$. (This also implies the determinant is unique.)
\item Instead of row 1, $d(A)$ can be calculated along any row or column, and the result is the same:
\begin{align*}
d(A) &= \sum_{k=1}^c(-1)^{i+k}a_{ik}d(A_{ik}) \quad \text{(along row } i)\\
    &=  \sum_{k=1}^c(-1)^{k+j}a_{kj}d(A_{kj}) \quad \text{(along column }j)
\end{align*}
\end{enumerate}
\end{theorem}

\bbtwocol{0.65}{0.35}{
First let's think about the $(-1)^{i+j}$. In row 1, column 1, we have $(-1)^{1+1} = +1$. And every time you increase the row or column index by one, the sign of $(-1)^{i+j}$ changes. This leads to the pattern on the right.
}{
\[\begin{bmatrix}
+ & - & + & \cdots \\
- & + & - & \cdots \\
+ & - & + & \cdots \\
\vdots & \vdots & \vdots & \ddots
\end{bmatrix}\]
}

\begin{book}
\textbf{Proof.} For (i), suppose $A$ is any $c\times c$ matrix. First let's define the $c \times c$ matrices $A_{1*}, A_{2*}, \ldots, A_{c*}$, which are like $A$, but with all the entries in the first column replaced with 0 except one entry:
\begin{align*}
A_{1*}&= \begin{bmatrix}
a_{11} & \text{rest of row 1 of }A \\
\vec{0} & A_{11}
\end{bmatrix}\\
A_{2*} &= \begin{bmatrix}
0 & \text{row 1 of }A_{21} \\
a_{21} & \text{rest of row 2 of }A \\
\vec{0} & \text{other rows of }A_{21}
\end{bmatrix} \\
&\vdots \\
A_{c*} &= \begin{bmatrix}
\vec{0} & A_{c1} \\
a_{c1} & \text{rest of row c of} A
\end{bmatrix}
\end{align*}
We'll find formulas for each $\det(A_{i*})$ and put them together at the end to make $d(A)$. 

To find $\det(A_{1*})$, using column operations, $A$ can be transformed to an echelon (triangular) matrix $B$ while ignoring the first row and first column. All of these column operations change the determinant by a non-zero scalar (see the proof of theorem \ref{theorem:determinants-and-invertibility} for details). So:
\begin{align}
\det(A_{1*}) &= \frac{\text{product of diagonal entries of }B}{\text{product of non-zero scalars}} \nonumber \\ 
&= \frac{a_{11}(\text{product of diagonal entries of }B_{11})}{\text{product of non-zero scalars}} \nonumber \\ 
&= a_{11}\det(A_{11}) \nonumber \\ 
&= (-1)^{1+1}a_{11}\det(A_{11}) \label{eq:3.2-A1}
%&= d(A_{1*})
\end{align}
Thus the theorem is true for matrices in the form of $A_{1*}$, where all of the first column is zero except $a_{11}$. For $\det(A_{2*})$, since $\det(A_{2*}) = \det((A_{2*})^{T})$, swapping the first two \emph{rows} changes the sign of the determinant by antisymmetry. Thus
\[\det(A_{2*}) = -\det \begin{bmatrix}
a_{21} & \text{rest of row 2 of }A_{2*} \\
\vec{0} & A_{21}
\end{bmatrix}\]
where the matrix on the right has the same form as $A_{1*}$ so the argument above shows:
\begin{align}
\det(A_{2*}) &= -a_{21}\det(A_{21}) \nonumber \\
&= (-1)^{2+1}a_{21}\det(A_{21}) \label{eq:3.2-A2}
%&= d(A_{2*}) 
\end{align}
Similarly, with $A_{i*}$, one can use row swaps to move row $i$ into row 1, leaving the order of the rest of the rows intact, by swapping row $i$ with $i-1$, then row $i-1$ with $i-2$, and so on, which multiplies the determinant by $(-1)^{i-1}$. Thus using the same argument as for $A_{1*}$:
\begin{align}
\det(A_{i*}) &= (-1)^{i-1}\det \begin{bmatrix}
a_{i1} & \text{rest of row }i\text{ of }A_{i*} \\
\vec{0} & A_{i1}
\end{bmatrix} \nonumber \\ 
&=(-1)^{i-1}a_{i1}\det(A_{i1}) \nonumber \\
&= (-1)^{i+1}a_{i1}\det(A_{i1}) \label{eq:3.2-Ai}
%&= d(A_{i*}) 
\end{align}

Now let's bring it all together. The first column of $A$ is the sum of the first columns of $A_{1*}, A_{2*}, \ldots, A_{c*}$, so using the linearity (addition) property in the definition of the determinant and equation \ref{eq:3.2-Ai},
\begin{align*}
\det(A) &= \det(A_{1*})  + \cdots + \det(A_{c*}) \\
&= (-1)^{1+1}a_{11}\det(A_{11})  + \cdots + (-1)^{c+1}a_{cn}\det(A_{c1}) \\
&= d(A)
\end{align*}
which proves part (i).

For part (ii), we know that a determinant exists from theorem \ref{theorem:existence-of-determinants}, and from part (i) we know that if a determinant exists, it's equal to $d(A)$. Together, this means that $d(A)$ is indeed a determinant function, so all of the properties of determinants apply to $d(A)$. Specifically, column swaps change the sign of $d(A)$. So, to calculate the column 2 expansion, you could swap columns 1 and 2, and use the original column 1 formula from part (i), giving:
\[ d(A) = -\sum_{k=1}^c(-1)^{k+1}a_{k2}\det(A_{k2}) = \sum_{k=1}^c(-1)^{k+2}a_{k2}\det(A_{k2}) \]
Likewise, to the general column $j$ formula, we could swap columns $j$ with $j-1$, and then swap $j-1$ with $j-2$, and so on until you swap columns 1 and 2. This is a total of $j-1$ swaps which moves column $j$ into column 1, and leaves the order of the other columns intact. Thus:
\[ d(A) = (-1)^{j-1}\sum_{k=1}^c(-1)^{k+1}a_{kj}\det(A_{kj}) = \sum_{k=1}^c(-1)^{k+j}a_{kj}\det(A_{kj}) \]
To get the formula for the row expansions, use the fact that $d(A^T)=d(A)$, and the column expansion formula on $d(A^T)$. \blacksquare
\end{book}

The reason for the term ``cofactor expansion'' is that $C_{ij} = (-1)^{i+j}\det(A_{ij})$ is called the $(i,j)$-\textbf{cofactor} of $A$. Cofactors can describe some nice formulas for $A^{-1}$ as well as solutions of $A\vec{x}=\vec{y}$ that will probably be included in a future version of this book.

\example Let  for $A=\begin{bmatrix}
2 & 4 & \sqrt{7} \\
0 & 3 & -8 \\
0 & 0 & 1
\end{bmatrix}$. Use the cofactor expansions along column 1 and row 2 to calculate $\det(A)$.

\begin{notes}
\vfill 
\end{notes}
\begin{book}
\textbf{Solution.} Along column 1, and using theorem \ref{theorem:2x2-det} for $2\times 2$ determinants,
\begin{align*}
\det(A) &= 2\det \begin{bmatrix}
3 & -8 \\
0 & 1
\end{bmatrix} - 0\det \begin{bmatrix}
4 & \sqrt{7} \\
0 & 1
\end{bmatrix}
+0 \det \begin{bmatrix}
4 & \sqrt{7} \\
3 & -8
\end{bmatrix} \\
&= 2 ( 3(1) - 0(-8) ) -0+0 = 2(3)(1) = 6
\end{align*}
which is the product of the main diagonal entries, agreeing with theorem \ref{theorem:det-of-triangular}.

Along row 2, the signs of the $(-1)^{i+j}$ factors are $-1, 1$, and $-1$, giving:
\begin{align*}
\det(A) &= -0\det\begin{bmatrix}
4 & \sqrt{7} \\
0 & 1
\end{bmatrix} + 3\det\begin{bmatrix}
2 & \sqrt{7} \\
0 & 1
\end{bmatrix} -1\det\begin{bmatrix}
2 & 4 \\
0 & 0
\end{bmatrix} \\
&=-0 +3 ( 2(1)-0(\sqrt{7}) )-1 ( 2(0)-0(4) )\\
&=-0+ 3(2)(1) -0 = 6
\end{align*}
\blacksquare
\end{book}

In practice, calculating determinants (either by hand or computer) with the cofactor expansion becomes extremely inefficient with matrices bigger than $3\times 3$. A $c\times c$ matrix requires more than $c!$ multiplications to compute the cofactor expansion, which means that doing column or row operations until you get a triangular matrix is much more efficient for larger matrices.
%TODO: Cramer's rule and A^{-1} formulas

\begin{notes}
\columnbreak
\end{notes}

\exerciseheading

%NOTE: A lot of the problems below are harder than the cofactor expansions and should be done in a different order.
\textbf{\ref{subsec:some-nice-theorems-about-determinants} Some Nice Theorems About Determinants}

\exatend (T/F?) If $A$ is $3\times3$ and $A \bbbvec{1}{2}{3} = \vec{0}$, then $\det(A) \ne 0$.

\exatend (T/F?) If $A$ and $B$ are row-equivalent, then $\det(A) = \det(B)$.

\exatend (T/F?) If $A$ and $B$ are row-equivalent, then either $\det(A)$ and $\det(B)$ both equal zero, or neither of them are zero.

\exatend (T/F?) The volume of the parallelogram/parallelepiped determined by the columns of $A$ equals the volume of the parallelogram/parallelepiped determined by the \emph{rows} of $A$.

\exatend (T/F?) If $A$ is $c\times c$ and $\det A \ne 0$, then the columns of $A$ are a basis of $\R^c$.

\exatend (T/F?) If $A$ is square, then $\det(A^TA) \ge 0$.

\exatend (T/F?) If $A$ is square, then $\det(-A) = -\det(A)$.


\exatend Suppose $\vec{b}_1$ and $\vec{b}_2$ are vectors in $\R^2$, and that $A$ is a $2\times 2$ matrix. We know that the signed area of $\pgram \begin{bmatrix}
\vec{b}_1 & \vec{b}_2
\end{bmatrix}$ is $\det(B)$. How does multiplication by $A$ affect this area? In other words, what is the area of $\pgram \begin{bmatrix}
A\vec{b}_1 & A\vec{b}_2
\end{bmatrix}$?

\exatend In chapter \ref{chap:eigenstuff}, we'll say that square matrices $A$ and $B$ are \emph{similar} if there exists an invertible $P$ such that $A = PBP^{-1}$. Show that if $A$ and $B$ are similar, then $\det(A) = \det(B)$.

%credit: ken kuttler
\exatend In chapter \ref{chap:dot-products-etc}, we'll look at \emph{orthogonal} matrices, which are square matrices that turn out to have the property $A^TA = I$. What are possible values of $\det(A)$ if $A$ is orthogonal? \creditfootnote{Inspired by an exercise in Ken Kuttler's excellent \emph{A First Course in Linear Algebra}.}

\exatend Prove that if $A$ is invertible, then $\det(A^{-1}) = \frac{1}{\det(A)}$.

\exatend $A^n$ means $AA\cdots A$ multiplied with $n$ copies of $A$. Show that $\det(A^n)=\det(A)^n$.

% credit:margalit/rabinoff
\exatend Let $A_1, A_2, \ldots, A_n$ be $c\times c$ matrices. Show that the product $A_1A_2\cdots A_n$ is invertible if and only if each matrix $A_i$ is invertible. \creditfootnote{Inspired by an example in Dan Margalit and Joseph Rabinoff's excellent \emph{Interactive Linear Algebra}.}

\exatend Explain why row operations on $A$ affect $\det(A)$ the same way the corresponding column operations do. \hintfootnote{Theorem \ref{theorem:detA-detAT}.}

\exatend Prove that $\det(AB) = \det(BA)$.

\exatend Suppose $A=\begin{bmatrix}
\vec{a}_1 & \vec{a}_2 & \vec{a}_3 & \vec{a}_4
\end{bmatrix}$ and $B$ are $4\times4$ matrices, with $\det(A) = 2$ and $\det(B)=-3$. Find:

\begin{horiparts}{2}
\item $\det(3\vec{a}_1,\vec{a}_1+\vec{a}_2,\vec{a}_4,\vec{a}_3)$
\item $\det(\vec{a}_2,\vec{a}_3,\vec{a}_4,\vec{a}_1)$
\item $\det(2A)$
\item $\det(AB)$
\item $\det(B^2)$
\item $\det(BA^{-1})$
\end{horiparts}

\textbf{\ref{subsec:the-cofactor-expansion} The Cofactor Expansion}

\exatend (T/F?) To compute a $6\times 6$ determinant with a cofactor expansion requires computing $6$ different $5\times 5$ determinants.

For exercises \ref{exercise:3.2-basic-start} through \ref{exercise:3.2-basic-end}, calculate the determinant with two cofactor expansions, along different rows or columns, and check you get the same result.

\exatend \label{exercise:3.2-basic-start} $\det \begin{bmatrix}
2 & -1 & 4 \\
3 & 0 & 1 \\
-2 & 5 & 3
\end{bmatrix}$

\exatend $\det \begin{bmatrix}
2 & 1 & 0 \\
-3 & 4 & 7 \\
-1 & 5 & 7
\end{bmatrix}$

\exatend $\det \begin{bmatrix}
1 & 2 & 3 \\
4 & 5 & 6 \\
7 & 8 & 9
\end{bmatrix}$

\exatend \label{exercise:3.2-basic-end} $\det \begin{bmatrix}
1 & 3 & 0 & -1 \\
0 & 2 & 3 & 0 \\
1 & 0 & 0 & 5 \\
2 & 1 & 0 & 3
\end{bmatrix}$


\exatend Suppose you don't yet know that $\det\begin{bmatrix}
a & b \\
c & d 
\end{bmatrix} = ad-bc$, but you do know how to calculate $1\times1$ determinants: $\det \left[ a \right] = a$. Use this and the cofactor expansion to find $\det\begin{bmatrix}
a & b \\
c & d 
\end{bmatrix}$.

\exatend Do a column 1 cofactor expansion to calculate each of the following. Do you notice a pattern?
\begin{enumerate}[label=(\alph*)]
\item $\det\begin{bmatrix}
1 & 1 \\
-1 & 1
\end{bmatrix}$
\item $\det\begin{bmatrix}
1 & 1  & 0\\
-1 & 1 & 1 \\
0 & -1 & 1
\end{bmatrix}$
\item $\det\begin{bmatrix}
1 & 1 & 0 & 0 \\
-1 & 1 & 1 & 0 \\
0 & -1 & 1 & 1 \\
0 & 0 & -1 & -1
\end{bmatrix}$
\end{enumerate}

\exatend A simple computer program calculates determinants using the cofactor expansion. The program doesn't do row operations, or any tricks to look for zero factors; it just applies the formula until it gets to a $2\times2$ matrix where it does $ad-bc$. How many multiplications does it compute for:
\begin{enumerate}[label=(\alph*)]
\item $2\times2$ determinants?
\item $3\times3$ determinants?
\item $4\times4$ determinants?
\item $5\times5$ determinants?
\end{enumerate}

\chapter{Diagonalization and Eigenstuff}
\label{chap:eigenstuff}

In this chapter, we'll study techniques for analyzing sequences of vectors of the form 
\[\vec{y}_{n+1} =A\vec{y}_n\] 
This simple-looking notation hides lots of depth. As we'll see, this kind of formula can model many situations where there are many quantities that change at a rate that depends on the other quantities--making the theory useful in differential equations, dynamical systems, probability theory, and many real-life applications of these.

\section{Review and Diagonalization}
\label{sec:review-and-diagonalization}
A solid understanding of the concepts of bases, coordinates, and dimension, and how they relate to matrices, will make this chapter much easier. So we'll start with an extra long set of warm-up exercises to review these concepts, as well as introduce a little new stuff.

\exercise \emph{(Warm-up)} Review the following definitions and theorems (unless you're confident you already have them mastered).
\begin{enumerate}[label=(\alph*)]
\item The definition of a basis in section \ref{sec:bases-and-coordinates}.
\item The definition of $\mathcal{B}$-coordinates in section \ref{sec:bases-and-coordinates}.
\item The definition of dimension in section \ref{sec:dimension}.
\item Theorem \ref{theorem:buy-two-get-one-free} (buy two, get one free) in section \ref{sec:dimension}.
\end{enumerate}

If you struggle with the rest of these warm-ups, the examples and discussions in sections \ref{sec:bases-and-coordinates} and \ref{sec:dimension} might help.

\begin{notes}
\columnbreak
\end{notes}

\exercise \emph{(Warm-up)} (T/F?) The list of vectors $\bbbvec{3}{0}{-2}, \bbbvec{0}{1}{0}$ is a basis of $\R^3$.

\begin{notes}
\vspace{2em}
\end{notes}

\exercise \emph{(Warm-up)} (T/F?) The list of vectors $\bbbvec{3}{0}{-2}, \bbbvec{0}{1}{0}$ is a basis of $\spn \left( \bbbvec{3}{0}{-2}, \bbbvec{0}{1}{0} \right)$, which is a 2-dimensional subspace of $\R^3$.

\begin{notes}
\vspace{2em}
\end{notes}

\exercise \emph{(Warm-up)} (T/F?) The list of vectors $\bbvec{1}{2}, \bbvec{3}{4}, \bbvec{5}{6}$ is a basis of $\R^2$.

\begin{notes}
\vspace{2em}
\end{notes}

\exercise \emph{(Warm-up)} (T/F?) The list of vectors $\bbvec{1}{2}, \bbvec{3}{4}, \bbvec{5}{6}$ is a basis of $\spn \left( \bbvec{1}{2}, \bbvec{3}{4}, \bbvec{5}{6} \right)$ which is a 3-dimensional subspace of $\R^2$.

\begin{notes}
\vspace{2em}
\end{notes}
\exercise \emph{(Warm-up)} The grid below shows $\mathcal{B} = ( \vec{b}_1,\vec{b}_2 )$, a basis of $\R^{2}$. Find:

\begin{horiparts}{2}
\item $ \bbvec{6}{0}_{\mathcal{B}}$
\item $\vec{y}$ if $\left[ \vec{y} \right]_{\mathcal{B}} = \bbvec{6}{0}$
\item $\bbvec{-2}{4}_{\mathcal{B}}$
\item $\vec{y}$ if $\left[ \vec{y} \right]_{\mathcal{B}} = \bbvec{-2}{4}$
\end{horiparts}

\begin{tikzpicture}[scale=0.5]
\vectorgrid[xmin=-6.2,xmax=8.2,ymin=-4.2,ymax=6.2,step=1,labint=2,xu=2,yu=2,xv=4,yv=-2]
\node[right] at (2, 2) {$\vec{b}_1$};
\node[right] at (4, -2) {$\vec{b}_2$};
\end{tikzpicture}

\begin{notes}
\columnbreak
\end{notes}

\exercise \emph{(Warm-up)} With $\vec{b}_1=\bbvec{1}{1}, \vec{b}_2=\bbvec{2}{-1}$ as in the last exercise, let $P = \begin{bmatrix}
\vec{b}_1 & \vec{b}_2
\end{bmatrix}$.
\begin{enumerate}[label=(\alph*)]
\item Explain how you know $P$ is invertible.
\item In the last exercise you found $\bbvec{6}{0}_{\mathcal{B}}$. Check that $P \bbvec{6}{0}_{\mathcal{B}} = \bbvec{6}{0}$.
\item Find $P^{-1}$.
\item In the last exercise you found $\vec{y}$ if $\left[ \vec{y} \right]_{\mathcal{B}} = \bbvec{6}{0}$. Check that $\left[ \vec{y} \right]_{\mathcal{B}} = P^{-1}\vec{y}$. 
\end{enumerate}

\begin{notes}
\vfill
\end{notes}
\exercise \emph{(Warm-up)} Using the IMT, convince yourself that the following statement is true: A list of vectors $\vec{b}_1,\ldots,\vec{b}_c$ in $\R^c$ is a basis of $\R^c$ if and only if the matrix $\begin{bmatrix}
\vec{b}_1 & \cdots & \vec{b}_c
\end{bmatrix}$ is invertible. No need to write anything; just ask for help if you don't see why it's true. (Warm-ups end here.)

\subsection{Difference Equations and Diagonalizability}\label{subsec:difference-equations-and-diagonalizability}

Now let's look at a certain type of application problem, and we'll soon see how to apply bases and coordinates to them.

\example Spock accidentally transports a population of 100 baby tribbles (cute, furry, little creatures) to an alien planet where there is plenty of food, but also predators. Every week, 80\% of the baby tribbles are eaten by predators, but the other 20\% grow to become adult tribbles. And, every week, every adult tribble gives birth to 5 baby tribbles, and at the end of the week, 60\% of the adults are eaten. Let $\vec{x}_w = (b_w,a_w)$ represent the number of baby and adult tribbles $w$ weeks later; starting with $\vec{x}_0 = (100,0)$.
\begin{enumerate}[label=(\alph*)]
\item Find $\vec{x}_1, \vec{x}_2$ and $\vec{x}_3$.
\item Find a matrix $A$ so that $\vec{x}_{w+1} = A\vec{x}_{w}$.
\item Use a computer to see what happens to the tribble population in the long run (as $w \rightarrow \infty$).
\end{enumerate}

\begin{notes}
\columnbreak
\hspace{1ex} \vfill
\end{notes}

\begin{book}
\textbf{Solution.} (a) After one week, 20\% of the beginning 100 babies survive to adulthood. There were no adults at the beginning to give birth, so: 
\[\vec{x}_1=\bbvec{b_1}{a_1} = \bbvec{0}{0.2b_0}=\bbvec{0}{0.2(100)} = \bbvec{0}{20}\]
At the end of week 2, the 20 adults have given birth to 5 babies each, and then 40\% of them remain uneaten:
\[\vec{x}_2=\bbvec{b_2}{a_2} = \bbvec{5a_{1}}{0.4a_1}=\bbvec{5(20)}{0.4(20)} = \bbvec{100}{8} \]
One week later, 20 of the babies survive to adulthood, the 8 adults have 40 babies, and 40\% of the adults survive:
\[\vec{x}_3=\bbvec{b_3}{a_3} = \bbvec{5a_{2}}{0.2b_2+0.4a_2} = \bbvec{5(8)}{0.2(100)+0.4(8)} = \bbvec{40}{23.2} \]
(b) In general, the number of babies in week $w+1$ is $b_{w+1}=5a_w$, and the number of adults is $a_{w+1}=0.2b_w+0.4a_w$. We can write this using the elegant notation of linear algebra like so:
\[ \vec{x}_{w+1}=\bbvec{b_{w+1}}{a_{w+1}} = \bbvec{5a_w}{0.2b_w+0.4a_w} = b_w\bbvec{0}{0.2} + a_w \bbvec{5}{0.4} = \begin{bmatrix}
0 & 5 \\
0.2 & 0.4
\end{bmatrix}\bbvec{b_w}{a_w}\]
Thus the matrix $A = \begin{bmatrix}
0 & 5 \\
0.2 & 0.4
\end{bmatrix}$ makes $\vec{x}_{w+1}=A\vec{x}_w$
(c) Let's get a general formula for $\vec{x}_w$. We know $\vec{x}_0=(100,0)$, and $\vec{x}_1=Ax_0$. Then:
\[\vec{x}_2=A\vec{x}_1=A ( A\vec{x}_0 ) = A^{2}\vec{x}_0\]
And likewise,
\[ \vec{x}_3=A\vec{x}_2=A(A^2\vec{x}_0) = A^3\vec{x}_0\]
So every time a week goes by, you just multiply by a factor of $A$, and in general,
\[ \vec{x}_w=A^w\vec{x}_0\]
(c) To see what happens in the long run, let's just plug some big values of $w$ into a computer:
\begin{align*}
\vec{x}_{50}&=A^{50}\vec{x}_{0} \approx \bbvec{829209}{2002294}\\
\vec{x}_{100}&=A^{100}\vec{x}_{0} \approx \bbvec{17.1 \text{ trillion}}{41.7 \text{ trillion}}\\
\vec{x}_{101}&=A^{101}\vec{x}_{0} \approx \bbvec{20.9 \text{ trillion}}{50.9 \text{ trillion}}
\end{align*}
So, it appears the tribbles are thriving on their new planet, increasing the population by about 20\% each week; in the long run, it seems $b_w$ and $a_w$ will increase without bound. \blacksquare
\end{book}

Of course, in the real world, the population of tribbles wouldn't grow to infinity for many reasons. For one, the \emph{real} Spock doesn't make mistakes. Some other factors to consider:
\begin{itemize}
\item The tribbles' food supply on the planet isn't infinite. Assuming it's a living thing, the rate it can replenish itself depends on how much of it there is, and how many tribbles are eating it.
\item The well-fed population of predators would also increase, and the more predators there are, the smaller a percentage of tribbles survive to reproduce.
\item Over thousands of years, the predators, the tribbles, and the tribbles' food will all evolve by natural selection, gradually changing their abilities to eat and reproduce.
\end{itemize}
Many of the mathematical models that take these factors into account are based on more or less complicated versions of $\vec{y}_{n+1}=A\vec{y}_n$. As we saw this is equivalent to $\vec{y}_n=A^n\vec{y}_{0}$; either way it's called a \textbf{linear difference equation}.

\begin{notes}
\columnbreak
\end{notes}

\example A new virus emerges in a population of 1000 susceptible people. Each day, 8\% of the susceptible people become infected, and 40\% of the infected people recover. Also, 1\% of the recovered population become reinfected. Let $\vec{p}_n = (s_n, i_n, r_n)$ represent the number of susceptible, infected, and recovered people, $n$ days later.
\begin{enumerate}[label=(\alph*)]
  \item Find $\vec{p}_1, \vec{p}_2$ and $\vec{p}_3$.
  \item Find $A$ so that $\vec{p}_{n+1}=A\vec{p}_n$.
  \item Use a computer to see what happens in the long run.
  \item Find $A \bbbvec{0}{24.39}{975.61}.$
\end{enumerate}

\begin{book}
\textbf{Solution.} (a) After one day, 92\% of the 1000 susceptible people are still susceptible after 8\% become infected:
\[\vec{p}_1=\bbbvec{s_1}{i_1}{r_1}=\bbbvec{0.92s_0}{.08s_0}{0} = \bbbvec{920}{80}{0} \]
The next day, there are some infected people and 40\% of them recover:
\[\vec{p}_2=\bbbvec{s_2}{i_2}{r_2}=\bbbvec{0.92s_1}{.08s_1 + .6i_1}{.4i_1} = \bbbvec{.92(920)}{.08(920)+.6(80)}{.4(80)} = \bbbvec{846.4}{121.6}{32} \]
After three days, 1\% of the recovered people become reinfected:
\begin{align*}
\vec{p}_3&=\bbbvec{s_3}{i_3}{r_3}=\bbbvec{0.92s_2}{.08s_2 + .6i_2 + .01r_2}{.4i_2+.99s_2} = \\ &=\bbbvec{.92(846.4)}{.08(846.4)+.6(121.6) + .01(32)}{.4(121.6)+.99(32)} = \bbbvec{778.688}{140.992}{80.32}
\end{align*}
(b) In general:
\begin{align*}
\vec{p}_{n+1} = \bbbvec{s_{n+1}}{i_{n+1}}{r_{n+1}} = \bbbvec{.92s_n}{.08s_n+.6i_n+.01r_n}{.4i_n+.99r_n} =
\begin{bmatrix}
.92 & 0 & 0 \\
.08 & .6 & .01 \\
0 & .4 & .99
\end{bmatrix}\bbbvec{s_n}{i_n}{r_n}
\end{align*}
(c) Like with the last example, we know $\vec{p}_0=(1000,0,0)$, and from part (b):
\begin{align*}
\vec{p}_1 &= A\vec{p}_0 \\
\vec{p}_2 &= A\vec{p}_1= A(A\vec{p}_0) = A^2\vec{p}_0 \\
\vec{p}_3 &= A\vec{p}_2= A(A^2\vec{p}_0) = A^3\vec{p}_0 \\
&\cdots \\
\vec{p}_n &= A^n\vec{p}_0
\end{align*}
So if we plug in some big values of $n$ into a computer:
\begin{align*}
\vec{p}_{50}=A^{50}(1000,0,0) \approx \bbbvec{15.5}{27.7}{956.9} \\
\vec{p}_{100}=A^{100}(1000,0,0) \approx \bbbvec{0.24}{24.44}{975.32} \\
\vec{p}_{101}=A^{101}(1000,0,0) \approx \bbbvec{0.22}{24.44}{975.34} \\
\end{align*}
It looks like it's approaching an \emph{equilibrium} where the number of infected people recovering equals the number of recovered people getting reinfected.
(d) $A \bbbvec{0}{24.39}{975.61} \approx \bbbvec{0}{24.39}{975.61}$. \blacksquare
\end{book}

\begin{notes}
\vfill \hspace{1pt}
\columnbreak

\vfill
\end{notes}

Notice that in the last example the entries in each column of $A$ add up to 1, because each column describes what happens to 100\% of each group from the previous day. This kind of matrix is very useful in statistics, and is called a \emph{stochastic matrix}, and when $A$ is stochastic, the sequence $\vec{x}_n=A^n\vec{x}_0$ is called a \emph{Markov chain}.

When calculating $A^n$ in many real world applications, the matrix $A$ might be enormous. For example, instead of just baby and adult tribbles, you may have entries in $\vec{y}_w$ for thousands of species in the ecosystem whose populations affect each other in various ways. And Google's PageRank algorithm, for assessing the quality of web pages, uses a powers of a square matrix with a row and column for \emph{every page on the web}, and there are many trillions of those. In general, if $A$ is $c\times c$, calculating $A^n$ requires $nc^3$ multiplications. We'll soon see that the concept of diagonalization makes calculation much easier, as well as helping with theoretical analysis.

\begin{definition}
A square matrix $A$ is \textbf{similar} to $B$ if there exists an invertible matrix $P$ such that $A = PBP^{-1}$. And $A$ is \textbf{diagonalizable} if $A$ is similar to a diagonal matrix $D$.
\end{definition}

\begin{notes}
\columnbreak
\end{notes}
\example Let $A=\begin{bmatrix}
-7 & 9 \\
-6 & 8
\end{bmatrix}$. It can be shown that if $P=\begin{bmatrix}
3 & 1 \\
2 & 1
\end{bmatrix}$, and $D=\begin{bmatrix}
-1 & 0 \\
0 & 2
\end{bmatrix}$, then $A=PDP^{-1}$, with $P^{-1}=\begin{bmatrix}
1 & -1 \\
-2 & 3
\end{bmatrix}$. 
\begin{enumerate}[label=(\alph*)]
\item Calculate $A^2$ and $A^3$ by hand (to see that it's tedious).
\item Calculate $D^2$ and $D^3$ by hand, and find a general formula for $D^n$.
\item Find a formula for $A^n$ in terms of $P$ and $D$.
\end{enumerate}
\begin{book}
\textbf{Solution.} (a)
\[ A^2= \begin{bmatrix}
-7 & 9 \\
-6 & 8
\end{bmatrix}\begin{bmatrix}
-7 & 9 \\
-6 & 8
\end{bmatrix} = \begin{bmatrix}
-7(-7)+9(-6) & -7(9) + 9(8) \\
-6(-7)+8(-6) & -6(9) + 8(8)
\end{bmatrix} = \begin{bmatrix}
-5 & 9 \\
-6 & 10
\end{bmatrix}
\]
\[A^{3} = \begin{bmatrix}
-7 & 9 \\
-6 & 8
\end{bmatrix}\begin{bmatrix}
-5 & 9 \\
-6 & 10
\end{bmatrix} = \begin{bmatrix}
-7(-5)+9(-6) & -7(9) + 9(10) \\
-6(-5)+8(-6) & -6(9) + 8(10)
\end{bmatrix} = \begin{bmatrix}
-19 & 27 \\
-18 & 26
\end{bmatrix}\]
(b)
\[ D^2= \begin{bmatrix}
-1 & 0 \\
0 & 2
\end{bmatrix}\begin{bmatrix}
-1 & 0 \\
0 & 2
\end{bmatrix} = \begin{bmatrix}
(-1)^2 + 0^2 & -1(0) + 0(2) \\
0(-1) + 2(0) & 0^2+2^2
\end{bmatrix} = \begin{bmatrix}
(-1)^2 & 0 \\
0 & 2^2
\end{bmatrix}\]
\[ D^3= \begin{bmatrix}
-1 & 0 \\
0 & 2
\end{bmatrix}\begin{bmatrix}
(-1)^2 & 0 \\
0 & 2^{2}
\end{bmatrix} = \begin{bmatrix}
(-1)^3 + 0^2 & -1(0) + 0(2^2) \\
0(-1)^2 + 2(0) & 0^2+2^3
\end{bmatrix} = \begin{bmatrix}
(-1)^3 & 0 \\
0 & 2^3
\end{bmatrix}\]
And in general, $D^n = \begin{bmatrix}
(-1)^n & 0 \\
0 & 2^n
\end{bmatrix}$.

(c)
\begin{align*}
A^n=( PDP^{-1} )^{n} &= ( PDP^{-1} )( PDP^{-1} )( PDP^{-1} ) \cdots( PDP^{-1} ) \\
&=PD ( P^{-1}P )D ( P^{-1}P ) D ( P^{-1} \cdots P )DP^{-1} \\
&= PDIDIDI\cdots IDP^{-1} = PD^nP^{-1}
\end{align*}
\blacksquare
\end{book}

\begin{notes}
\columnbreak
\end{notes}

In general, if $D$ is any diagonal matrix,
\[D= \begin{bmatrix}
d_1 & 0 & \cdots & 0 \\
0 & d_2 & \cdots & 0 \\
\vdots & \vdots & \ddots & \vdots \\
0 & 0 & \cdots & d_c
\end{bmatrix}\]
it is easy to calculate powers of $D$. Using the row-column rule, any entry in $D^2$ not on the main diagonal (say, row $i$, column $j$, with $i \ne j$) will be of the form:

\begin{book}
 \[0^2+\cdots+0^2+d_i(0) + 0^2+\cdots+0^2 +0(d_j)+0^2+\cdots+0^2 = 0\]
\end{book}

\begin{notes}

\vspace{2em}
\end{notes}
And any entry of $D^{2}$ on the main diagonal will be:

\begin{book}
\[ 0^2+ \cdots + 0^{2}+ d_i^2+ 0^2+\cdots+0^2=d_i^2 \]
\end{book}

\begin{notes}
\vspace{2em}
\end{notes}
This gives the following formula for $D^2$. Repeated multiplication by $D$ likewise gives the formula for $D^{n}$:

\begin{book}
\begin{equation}
\label{eq:4.1-Dn}
D^2= \begin{bmatrix}
d_1^2 & 0 & \cdots & 0 \\
0 & d_2^2 & \cdots & 0 \\
\vdots & \vdots & \ddots & \vdots \\
0 & 0 & \cdots & d_c^2
\end{bmatrix} \quad \text{and} \quad D^n= \begin{bmatrix}
d_1^n & 0 & \cdots & 0 \\
0 & d_2^n & \cdots & 0 \\
\vdots & \vdots & \ddots & \vdots \\
0 & 0 & \cdots & d_c^n
\end{bmatrix}
\end{equation}
\end{book}
\begin{notes}
\vspace{5em}
\end{notes}

Like in the last example, if $P$ is invertible, then $( PDP^{-1} )^n=PD^nP^{-1}$. The next theorem states a nice version of this in terms of coordinates.

\subsection{Diagonal Coordinates}\label{subsec:diagonal-coordinates}

\begin{theorem}[Diagonal coordinates.]
\label{theorem:diagonal-coordinates}

Suppose $A = PDP^{-1}$ with $D$ diagonal, and $P = \begin{bmatrix}
\vec{b}_1 & \cdots & \vec{b}_c
\end{bmatrix}$. Let $\mathcal{B} = ( \vec{b}_1,\ldots,\vec{b}_c )$ be the basis made of the columns of $P$. Then for any positive integer $n$ and any $\vec{y}$ in $\R^c$:
\[\left[ A^n\vec{y} \right]_{\mathcal{B}} = D^n \left[ \vec{y} \right]_{\mathcal{B}}\]
\end{theorem}

\begin{book}
\textbf{Proof.} First note that since $P$ is invertible, the IMT implies its columns are linearly independent and span $\R^c$, so $\mathcal{B} = ( \vec{b}_1,\ldots,\vec{b}_c )$ is indeed a basis. Let $\vec{x} = (x_1,\ldots, x_c)$ be $\left[ \vec{y} \right]_{\mathcal{B}}$. By definition this means:
\[x_1\vec{b}_1 + \cdots + x_c\vec{b}_c =P\vec{x} = P\left[ \vec{y} \right]_{\mathcal{B}} = \vec{y}\]
Multiplying by $P^{-1}$ on the left gives:
\begin{align}\label{eq:4.1-P-1y}
\left[ \vec{y} \right]_{\mathcal{B}} &= P^{-1}\vec{y} \nonumber \\
(x_1,\ldots,x_c) &= P^{-1}\vec{y}
\end{align}
Let $n$ be a positive integer and $\vec{y}$ be any vector in $\R^{c}$. To find $\left[ A^n\vec{y} \right]_{\mathcal{B}}$, we are looking for the scalar coefficients of $\vec{b}_1,\ldots,\vec{b}_c$ that add up to $A^n\vec{y}$. Using the same trick as the last example:
\begin{align*}
A^n\vec{y} &= ( PDP^{-1} )^{n}\vec{y} \\
&= ( PDP^{-1} )( PDP^{-1} )( PDP^{-1} ) \cdots( PDP^{-1} )\vec{y} \\
&=PD ( P^{-1}P )D ( P^{-1}P ) D ( P^{-1} \cdots P )DP^{-1}\vec{y} \\
&= PDIDIDI\cdots IDP^{-1}\vec{y} \\
&= PD^nP^{-1}\vec{y} \\
&= PD^n \left[ \vec{y} \right]_B & \text{(equation \ref{eq:4.1-P-1y})} \\
&= \begin{bmatrix}
\vec{b}_1 & \vec{b}_2 & \cdots & \vec{b}_c
\end{bmatrix} \begin{bmatrix}
d_1^n & 0 & \cdots & 0 \\
0 & d_2^n & \cdots & 0 \\
\vdots & \vdots & \ddots & \vdots \\
0 & 0 & \cdots & d_c^n
\end{bmatrix} \begin{bmatrix}
x_1 \\ x_2 \\ \vdots \\ x_{c}
\end{bmatrix} & \text{(equation \ref{eq:4.1-Dn})}\\
&= \begin{bmatrix}
\vec{b}_1 & \vec{b}_2 & \cdots & \vec{b}_c
\end{bmatrix} \begin{bmatrix}
d_1^nx_1 \\ d_2^nx_2 \\ \vdots \\ d_c^nx_c 
\end{bmatrix} \\
&= d_1^nx_1\vec{b}_1 + d_2^nx_2\vec{b}_2 + \cdots + d_c^nx_c\vec{b}_c
\end{align*}
So, by the definition of $\left[ A\vec{y} \right]_{\mathcal{B}}$, we have:
\begin{align*}
\left[ A\vec{y} \right]_{\mathcal{B}} &=  ( d_1^nx_1, d_2^nx_2, \ldots, d_c^nx_n ) \\
&= \begin{bmatrix}
d_1^n & 0 & \cdots & 0 \\
0 & d_2^n & \cdots & 0 \\
\vdots & \vdots & \ddots & \vdots \\
0 & 0 & \cdots & d_c^n
\end{bmatrix} \begin{bmatrix}
x_1 \\ x_2 \\ \vdots \\ x_c
\end{bmatrix} \\
&= D^n \left[ \vec{y} \right]_{\mathcal{B}}
\end{align*}
\blacksquare
\end{book}

\begin{notes}
\columnbreak

\vfill
\end{notes}

\example \label{example:4.1-coordinates-1} Let $A=\begin{bmatrix}
1.5 & 0.5 \\
1 & 1
\end{bmatrix}$. It can be shown that if $P = \begin{bmatrix}
1 & -1 \\
1 & 2
\end{bmatrix} = \begin{bmatrix}
\vec{b}_1 & \vec{b}_2
\end{bmatrix}$, and $D = \begin{bmatrix}
2 & 0 \\
0 & 0.5
\end{bmatrix}$, then $A = PDP^{-1}$.
\begin{enumerate}[label=(\alph*)]
\item Suppose $\vec{x}_n=A^n\vec{x}_0$, with $\vec{x}_0 = 1\vec{b}_1+1\vec{b}_2$. Use theorem \ref{theorem:diagonal-coordinates} to plot $\vec{x}_1$ and $\vec{x}_2$.
\item Repeat for some other choices of $\vec{x}_0$.
\end{enumerate}
\begin{tikzpicture}[scale=0.8]
\vectorgrid[xmin=-4.5,xmax=4.5,ymin=-5.5,ymax=5.5,step=2,labint=2,xu=1,yu=1,xv=-1,yv=2]
\node[right,blue] at (1, 1) {$\vec{b}_1$};
\node[left, blue] at (-1, 2) {$\vec{b}_2$};
\draw[blue, domain=-3:5, samples=200, smooth, variable=\t] 
    plot ({1*2^\t - 1*0.5^\t}, {1*2^\t + 2*0.5^\t});
\draw[blue, domain=-2:3, samples=100, smooth, variable=\t] 
    plot ({3*2^\t - 1*0.5^\t}, {3*2^\t + 2*0.5^\t});
\draw[blue, domain=-3:5, samples=200, smooth, variable=\t] 
    plot ({1*2^\t + 1*0.5^\t}, {1*2^\t - 2*0.5^\t});
\draw[blue,  domain=-2:3, samples=100, smooth, variable=\t] 
    plot ({3*2^\t + 1*0.5^\t}, {3*2^\t - 2*0.5^\t});
\draw[blue, domain=-3:5, samples=200, smooth, variable=\t] 
    plot ({-1*2^\t + 1*0.5^\t}, {-1*2^\t - 2*0.5^\t});
\draw[blue, domain=-3:5, samples=200, smooth, variable=\t] 
    plot ({-3*2^\t + 1*0.5^\t}, {-3*2^\t - 2*0.5^\t});
\draw[blue, domain=-3:5, samples=200, smooth, variable=\t] 
    plot ({-1*2^\t - 1*0.5^\t}, {-1*2^\t + 2*0.5^\t});
\draw[blue, domain=-3:5, samples=200, smooth, variable=\t] 
    plot ({-3*2^\t - 1*0.5^\t}, {-3*2^\t + 2*0.5^\t});
\end{tikzpicture}

\begin{book}
\textbf{Solution.} (a) $\vec{x}_0=1\vec{b}_1+1\vec{b}_2$ means $\left[ \vec{x}_0 \right]_{\mathcal{B}} = (1,1)$, which means $\vec{x}_0=(0,3)$. The notation may be confusing, but using theorem \ref{theorem:diagonal-coordinates}, the calculations to find $\left[ \vec{x}_1 \right]_{\mathcal{B}}$ and $\left[ \vec{x}_2 \right]_{\mathcal{B}}$ are simple:
\begin{align*}
\left[ \vec{x}_1 \right]_{\mathcal{B}} &= \left[ A\vec{x}_0 \right]_{\mathcal{B}} =
D\left[ \vec{x}_0 \right]_{\mathcal{B}} = \begin{bmatrix}
2 & 0 \\
0 & 0.5
\end{bmatrix}\bbvec{1}{1} = \bbvec{2}{0.5} \\
\left[ \vec{x}_2 \right]_{\mathcal{B}} &= \left[ A^2\vec{x}_0 \right]_{\mathcal{B}} =
D^2\left[ \vec{x}_0 \right]_{\mathcal{B}} = \begin{bmatrix}
4 & 0 \\
0 & 0.25
\end{bmatrix}\bbvec{1}{1} = \bbvec{4}{0.25}
\end{align*}
The $\mathcal{B}$-coordinate grid makes it easy to plot $\vec{x}_1$ and $\vec{x}_2$, as shown in black below.

(b) Let's try \textcolor{red}{$\left[ \vec{x}_0 \right]_{\mathcal{B}} = (0,-2)$ which means $\vec{x}_0=(2,-4).$ Like above:
\begin{align*}
\left[ \vec{x}_1 \right]_{\mathcal{B}} &= 
D\left[ \vec{x}_0 \right]_{\mathcal{B}} = \bbvec{2(0)}{0.5(-2)} = \bbvec{0}{-1} \\
\left[ \vec{x}_2 \right]_{\mathcal{B}} &= 
D^2\left[ \vec{x}_0 \right]_{\mathcal{B}} = \bbvec{4(0)}{0.25(-2)} = \bbvec{0}{-.5} \\
\left[ \vec{x}_3 \right]_{\mathcal{B}} &= 
D^3\left[ \vec{x}_0 \right]_{\mathcal{B}} = \bbvec{8(0)}{0.125(-2)} = \bbvec{0}{-.25} 
\end{align*}
Again, the $\mathcal{B}$-coordinate grid shows how to plot $\vec{x}_1, \vec{x}_2$, and $\vec{x}_3$.}
It's left to the reader to try other choices of $\vec{x}_0$. If you do, you'll find that multiplying by $A$ always makes the $\vec{b}_1$-coordinate double, and the $\vec{b}_2$-coordinate half. This keeps you on the hyperbola-like contour you started on. And, if $\vec{x}_0$ is a multiple of $\vec{b}_1$ (or $\vec{b}_2$), multiplying by $A$ leaves you with another multiple of $\vec{b}_1$ (or $\vec{b}_2$).

\begin{tikzpicture}[scale=0.8]
\vectorgrid[xmin=-4.5,xmax=4.5,ymin=-5.5,ymax=5.5,step=2,labint=2,xu=1,yu=1,xv=-1,yv=2]
\node[right,blue] at (1, 1) {$\vec{b}_1$};
\node[left, blue] at (-1, 2) {$\vec{b}_2$};
\draw[blue, domain=-3:5, samples=200, smooth, variable=\t] 
    plot ({1*2^\t - 1*0.5^\t}, {1*2^\t + 2*0.5^\t});
\draw[blue, domain=-2:3, samples=100, smooth, variable=\t] 
    plot ({3*2^\t - 1*0.5^\t}, {3*2^\t + 2*0.5^\t});
\draw[blue, domain=-3:5, samples=200, smooth, variable=\t] 
    plot ({1*2^\t + 1*0.5^\t}, {1*2^\t - 2*0.5^\t});
\draw[blue,  domain=-2:3, samples=100, smooth, variable=\t] 
    plot ({3*2^\t + 1*0.5^\t}, {3*2^\t - 2*0.5^\t});
\draw[blue, domain=-3:5, samples=200, smooth, variable=\t] 
    plot ({-1*2^\t + 1*0.5^\t}, {-1*2^\t - 2*0.5^\t});
\draw[blue, domain=-3:5, samples=200, smooth, variable=\t] 
    plot ({-3*2^\t + 1*0.5^\t}, {-3*2^\t - 2*0.5^\t});
\draw[blue, domain=-3:5, samples=200, smooth, variable=\t] 
    plot ({-1*2^\t - 1*0.5^\t}, {-1*2^\t + 2*0.5^\t});
\draw[blue, domain=-3:5, samples=200, smooth, variable=\t] 
    plot ({-3*2^\t - 1*0.5^\t}, {-3*2^\t + 2*0.5^\t});
\fill (0,3) circle (3pt) node[left] {$\vec{x}_0$};
\fill (1.5,3) circle (3pt) node[above] {$\vec{x}_1$};
\fill (3.75,4.5) circle (3pt) node[above] {$\vec{x}_2$};
\fill[red] (2,-4) circle (3pt) node[right] {$\vec{x}_0$};
\fill[red] (1,-2) circle (3pt) node[right] {$\vec{x}_1$};
\fill[red] (0.5,-1) circle (3pt) node[right] {$\vec{x}_2$};
\fill[red] (0.25,-.5) circle (3pt) node[right] {$\vec{x}_3$};
\end{tikzpicture} \blacksquare
\end{book}

\begin{notes}
\columnbreak

\vfill
\end{notes}

\example \label{example:4.1-coordinates-2} Let $A=\begin{bmatrix}
0.4 & 0.2 \\
-0.3 & 1.1
\end{bmatrix}$. It can be shown that if $P = \begin{bmatrix}
1 & 2 \\
3 & 1
\end{bmatrix} = \begin{bmatrix}
\vec{b}_1 & \vec{b}_2
\end{bmatrix}$, and $D = \begin{bmatrix}
1 & 0 \\
0 & 0.5
\end{bmatrix}$, then $A = PDP^{-1}$.
\begin{enumerate}[label=(\alph*)]
\item Suppose $\vec{x}_n=A^n\vec{x}_0$, with $\vec{x}_0 = 1\vec{b}_1+2\vec{b}_2$. Use theorem \ref{theorem:diagonal-coordinates} to plot $\vec{x}_1$ and $\vec{x}_2$.
\item Repeat for some other choices of $\vec{x}_0$.
\end{enumerate}
\begin{tikzpicture}[scale=0.7]
\vectorgrid[xmin=-3.5,xmax=5.5,ymin=-3.5,ymax=5.5,step=2,labint=2,xu=1,yu=3,xv=2,yv=1]
\node[right,blue] at (1, 3) {$\vec{b}_1$};
\node[right, blue] at (2, 1) {$\vec{b}_2$};
\end{tikzpicture}

\begin{book}
\textbf{Solution.} (a) $\vec{x}_0=1\vec{b}_1+2\vec{b}_2$ means $\left[ \vec{x}_0 \right]_{\mathcal{B}} = (1,2)$, which means $\vec{x}_0=(5,5)$. Like in the last example:
\begin{align*}
\left[ \vec{x}_1 \right]_{\mathcal{B}} &= \left[ A\vec{x}_0 \right]_{\mathcal{B}} =
D\left[ \vec{x}_0 \right]_{\mathcal{B}} = \begin{bmatrix}
1 & 0 \\
0 & 0.5
\end{bmatrix}\bbvec{1}{2} = \bbvec{1}{1} \\
\left[ \vec{x}_2 \right]_{\mathcal{B}} &= \left[ A^2\vec{x}_0 \right]_{\mathcal{B}} =
D^2\left[ \vec{x}_0 \right]_{\mathcal{B}} = \begin{bmatrix}
1 & 0 \\
0 & 0.25
\end{bmatrix}\bbvec{1}{2} = \bbvec{1}{0.25}
\end{align*}
The $\mathcal{B}$-coordinate grid makes it easy to plot $\vec{x}_1$ and $\vec{x}_2$, as shown in black below.

(b) Let's try \textcolor{red}{$\left[ \vec{x}_0 \right]_{\mathcal{B}} = (-1,0)$ which means $\vec{x}_0=(-1,-3).$ Like above:
\begin{align*}
\left[ \vec{x}_1 \right]_{\mathcal{B}} &= 
D\left[ \vec{x}_0 \right]_{\mathcal{B}} = \bbvec{1(-1)}{0.5(0)} = \bbvec{-1}{0} \\
\left[ \vec{x}_2 \right]_{\mathcal{B}} &= 
D^2\left[ \vec{x}_0 \right]_{\mathcal{B}} = \bbvec{1^2(-1)}{0.5^2(-0)} = \bbvec{-1}{0}
\end{align*}
So, the $1$ in $D$ makes it so the $\vec{b}_1$-coordinate doesn't change when multiplying by $A$. And if the $\vec{b}_2$-coordinate is 0, multiplying by $A$ does nothing.}
It's left to the reader to try other choices of $\vec{x}_0$. If you do, you'll find that multiplying by $A$ moves you parallel to $\vec{b}_2$ grid lines, because the $\vec{b}_1$-coordinate doesn't change.

\begin{tikzpicture}[scale=0.7]
\vectorgrid[xmin=-3.5,xmax=5.5,ymin=-3.5,ymax=5.5,step=2,labint=2,xu=1,yu=3,xv=2,yv=1]
\node[right,blue] at (1, 3) {$\vec{b}_1$};
\node[right, blue] at (2, 1) {$\vec{b}_2$};
\fill (5,5) circle (3pt) node[left] {$\vec{x}_0$};
\fill (3,4) circle (3pt) node[above] {$\vec{x}_1$};
\fill (2,3.5) circle (3pt) node[above] {$\vec{x}_2$};
\fill[red] (-1,-3) circle (3pt) node[right] {$\vec{x}_0=\vec{x}_1=\vec{x}_2$};
\end{tikzpicture}
\end{book}

\begin{notes}
\columnbreak
\end{notes}

\example \label{example:4.1-coordinates-3} Let $A=\begin{bmatrix}
-0.2 & 0.4 \\
-0.6 & 1.2
\end{bmatrix}$. It can be shown that if $P = \begin{bmatrix}
\vec{b}_1 & \vec{b}_2
\end{bmatrix}$ is the same as the last example, and $D = \begin{bmatrix}
1 & 0 \\
0 & 0
\end{bmatrix}$ (note the 0.5 was replaced with 0), then $A = PDP^{-1}$.
\begin{enumerate}[label=(\alph*)]
\item Suppose $\vec{x}_n=A^n\vec{x}_0$, with $\vec{x}_0 = 1\vec{b}_1+2\vec{b}_2$. Use theorem \ref{theorem:diagonal-coordinates} to plot $\vec{x}_1$ and $\vec{x}_2$.
\item Repeat for some other choices of $\vec{x}_0$.
\item Which of the matrices $A$ from the last three examples are invertible?
\end{enumerate}
\begin{tikzpicture}[scale=0.7]
\vectorgrid[xmin=-3.5,xmax=5.5,ymin=-3.5,ymax=5.5,step=2,labint=2,xu=1,yu=3,xv=2,yv=1]
\node[right,blue] at (1, 3) {$\vec{b}_1$};
\node[right, blue] at (2, 1) {$\vec{b}_2$};
\end{tikzpicture}

\begin{book}
\textbf{Solution.} (a) $\vec{x}_0=1\vec{b}_1+2\vec{b}_2$ means $\left[ \vec{x}_0 \right]_{\mathcal{B}} = (1,2)$, which means $\vec{x}_0=(5,5)$. Like in the last example:
\begin{align*}
\left[ \vec{x}_1 \right]_{\mathcal{B}} &= \left[ A\vec{x}_0 \right]_{\mathcal{B}} =
D\left[ \vec{x}_0 \right]_{\mathcal{B}} = \begin{bmatrix}
1 & 0 \\
0 & 0
\end{bmatrix}\bbvec{1}{2} = \bbvec{1}{0} \\
\left[ \vec{x}_2 \right]_{\mathcal{B}} &= \left[ A^2\vec{x}_0 \right]_{\mathcal{B}} =
D^2\left[ \vec{x}_0 \right]_{\mathcal{B}} = \begin{bmatrix}
1 & 0 \\
0 & 0
\end{bmatrix}\bbvec{1}{2} = \bbvec{1}{0}
\end{align*}

(b) Let's try \textcolor{red}{$\left[ \vec{x}_0 \right]_{\mathcal{B}} = (-1,0)$ which means $\vec{x}_0=(-1,-3).$ Like above:
\begin{align*}
\left[ \vec{x}_1 \right]_{\mathcal{B}} &= 
D\left[ \vec{x}_0 \right]_{\mathcal{B}} = \bbvec{1(-1)}{0.5(0)} = \bbvec{-1}{0} \\
\left[ \vec{x}_2 \right]_{\mathcal{B}} &= 
D^2\left[ \vec{x}_0 \right]_{\mathcal{B}} = \bbvec{1^2(-1)}{0.5^2(-0)} = \bbvec{-1}{0}
\end{align*}
So, the $1$ in $D$ makes it so the $\vec{b}_1$-coordinate doesn't change when multiplying by $A$. And if the $\vec{b}_2$-coordinate is 0, multiplying by $A$ does nothing.}
It's left to the reader to try other choices of $\vec{x}_0$. If you do, you'll find that multiplying by $A$ moves you parallel to $\vec{b}_2$ grid lines (because the $\vec{b}_1$-coordinate doesn't change), straight to a point in $\spn ( \vec{b}_1 )$.

\begin{tikzpicture}[scale=0.7]
\vectorgrid[xmin=-3.5,xmax=5.5,ymin=-3.5,ymax=5.5,step=2,labint=2,xu=1,yu=3,xv=2,yv=1]
\node[right,blue] at (1, 3) {$\vec{b}_1\textcolor{black}{=\vec{x}_1=\vec{x}_2}$};
\node[right, blue] at (2, 1) {$\vec{b}_2$};
\fill (5,5) circle (3pt) node[left] {$\vec{x}_0$};
\fill[red] (-1,-3) circle (3pt) node[right] {$\vec{x}_0=\vec{x}_1=\vec{x}_2$};
\end{tikzpicture}

(c) There are many ways to see the matrices $A$ in examples \ref{example:4.1-coordinates-1} and \ref{example:4.1-coordinates-2} are invertible. For one, the columns aren't multiples of each other, so they're independent. For another, you can pick any $\vec{x}_1$ in $\R^2$ and find an $\vec{x}_0$ along the same contour such that $\vec{x}_1=A\vec{x}_0$, so $\col A = \R^2$. For another, all three of $P, D$, and $P^{-1}$ are invertible, so their product must be invertible. For example \ref{example:4.1-coordinates-3}, $A$ is not invertible for similar reasons: The columns are dependent, and $\col A = \spn (\vec{b}_1) \ne \R^2$, and $D$ is not invertible.
\end{book}
\columnbreak

\exerciseheading

\textbf{\ref{subsec:difference-equations-and-diagonalizability} Difference Equations and Diagonalizability}
%NOTE: many problems here are harder than in the second subsection

\exatend Let $D = \begin{bmatrix}
0.8 & 0 \\
0 & 1
\end{bmatrix}$.
\begin{enumerate}[label=(\alph*)]
\item Find $\lim_{n \to \infty}D^n \bbvec{3}{4}$.
\item Find $\lim_{n \to \infty}D^n \bbvec{p}{q}$ in terms of $p$ and $q$.
\item Find $\vec{x}_0$ such that $\lim_{n \to \infty}D^n\vec{x}_0=\vec{0}$.
\end{enumerate}

\exatend Let $D = \begin{bmatrix}
1 & 0 & 0 \\
0 & 1.5 & 0 \\
0 & 0 & 0.7
\end{bmatrix}$.
\begin{enumerate}[label=(\alph*)]
\item Find $\lim_{n \to \infty}D^n \vec{x}_0$ for $\vec{x}_0 = (-5,0,10)$.
\item Find $\lim_{n \to \infty}D^n \vec{x}_0$ for $\vec{x}_0 = (1,1,1)$.
\item Find $\vec{x}_0$ such that $\lim_{n \to \infty}D^n \vec{x}_0 = \vec{0}$.
\end{enumerate}

\exatend A population of a species of plants is made up of either young plants or mature plants. Every year, 10\% of the young plants survive to maturity, and the other 90\% dies. And every mature plant produces seeds which become two mature plants every year. Also, 10\% of the mature plants die every year. Find $A$ so that $\bbvec{y_{n+1}}{m_{n+1}}=A \bbvec{y_n}{m_n}$ describes how the population changes over $n$ years.

\exatend A Blockbuster$^{\text{TM}}$ video rental store has three locations, $x,y$, and $z$. Every week, $90\%$ of the videos at location $x$ remain there, and $5\%$ are rented from $x$ and returned to $y$, with the other 5\% returned to $z$. And $80\%$ of the videos at location $y$ remain there; the other 20\% are rented from $y$ and returned to $z$. At location $z$, 99\% of the videos remain there every week, and the other $1\%$ is rented from $z$ and returned to $x$. Find $A$ so that $\bbbvec{x_{n+1}}{y_{n+1}}{z_{n+1}} = A \bbbvec{x_n}{y_n}{z_n}$ describes how the videos move around, where $n$ is the number of weeks.


\exatend (T/F?) Define a sequence of vectors starting with $\vec{x}_0$, and continuing with $\vec{x}_{n+1}=A\vec{x}_n$. Then $\vec{x}_n=A^n\vec{x}_0$.

\exatend (T/F?) If $A = PDP^{-1}$, then $A=D$ since the $P$ and $P^{-1}$ cancel out.

\exatend (T/F?) Every matrix that is row-equivalent to a diagonal matrix is diagonalizable.

\exatend (T/F?) If $A$ is diagonalizable, it is invertible.

\exatend Show that if $A$ is similar to $B$, then $B$ is similar to $A$.

\exatend Show that if $A$ is diagonalizable, then $A^2$ is diagonalizable.

\exatend Show that if $A^2$ is diagonalizable, then $A$ is diagonalizable.

\exatend Show that if $A$ and $B$ are similar matrices, then $A^2$ is similar to $B^2$.

%credit: beezer
\exatend Show that if $A$ and $B$ are similar matrices and $A$ is invertible, then $B$ is invertible, and $A^{-1}$ is similar to $B^{-1}$. \creditfootnote{Inspired by an exercise in Robert Beezer's excellent \emph{A First Course in Linear Algebra}.}

%credit: beezer
\exatend Show that if $A$ is invertible, then $AB$ is similar to $BA$. \creditfootnote{Inspired by R. Beezer.}

\textbf{\ref{subsec:diagonal-coordinates} Diagonal Coordinates}

\exatend The grid below shows $\mathcal{B} = ( \vec{b}_1,\vec{b}_2 )$, a basis of $\R^{2}$. Find:

\begin{enumerate}[label=(\alph*)]
\item $ \bbvec{6}{2}_{\mathcal{B}}$
\item $\vec{y}$ if $\left[ \vec{y} \right]_{\mathcal{B}} = \bbvec{6}{2}$
\item $\bbvec{-1}{5}_{\mathcal{B}}$
\item $\vec{y}$ if $\left[ \vec{y} \right]_{\mathcal{B}} = \bbvec{-1}{5}$
\item Let $P=\begin{bmatrix}
\vec{b}_1 & \vec{b}_2
\end{bmatrix}$ and multiply to check that $P \bbvec{6}{2}$ gives your answer to (b), and that $P$ times your answer to (a) gives $\bbvec{6}{2}$.
\item Use the grid to find $P^{-1}\bbvec{8}{0}$.
\item Let $D = \begin{bmatrix}
0.5 & 0 \\
0 & 2
\end{bmatrix}$, $A = PDP^{-1}$, and $\vec{x}_n=A^n\vec{x}_0$, with $\vec{x}_0=(5,-1)$. Use the method in examples \ref{example:4.1-coordinates-1} through \ref{example:4.1-coordinates-3} to find $\vec{x}_1$ and $\vec{x}_2$.
\end{enumerate}

\begin{tikzpicture}[scale=0.5]
\vectorgrid[xmin=-6.2,xmax=8.2,ymin=-4.2,ymax=6.2,step=1,labint=2,xu=-1,yu=-3,xv=2,yv=-2]
\node[left] at (-1, -3) {$\vec{b}_1$};
\node[left] at (2, -2) {$\vec{b}_2$};
\end{tikzpicture}

\exatend An electric scooter rental company has two locations where people can pick up or drop off scooters: the park, and the quarry. Let $\vec{y}_n = (p_n,q_n)$ represent the number of scooters at each location after $n$ days. They start with 100 scooters at each location, so $\vec{y}_0=(100,100)$. Each day, 70\% of scooters remain at the park, and the others are rented and return to the quarry. And 90\% of the scooters at the quarry remain at the quarry, and the rest are rented and returned to the park.
\begin{enumerate}[label=(\alph*)]
\item Find $\vec{y}_1$ and $\vec{y}_2$.
\item Find $A$ so that $\vec{y}_{n+1}= A\vec{y}_n$.
\item Before continuing, make a guess at how many scooters end up at the quarry in the long run. Do you think $\lim_{n \to \infty}q_n$ will be... 0? 200? Between 0 and 100? Between 100 and 200? 
\item Check by multiplying that $A=PDP^{-1}$, if $P = \begin{bmatrix}
-1 & 1 \\
1 & 3
\end{bmatrix}, D = \begin{bmatrix}
0.6 & 0 \\
0 & 1
\end{bmatrix}, P^{-1}= \begin{bmatrix}
-0.75 & 0.25 \\
0.25 & 0.25
\end{bmatrix}$.
\item Let $\mathcal{B} = (\vec{b}_1,\vec{b}_2 )$ be the basis made of the columns of $P$. Find $\left[ y_0 \right]_{\mathcal{B}}$.
\item Using the diagonal coordinates theorem, find $\left[ A\vec{y}_{0} \right]_{\mathcal{B}}$ and $\left[ A^2\vec{y}_{0} \right]_{\mathcal{B}}$. Multiply these by $P$ to get $A\vec{y}_0$ and $A^2\vec{y}_0$. Your answers should agree with $\vec{y}_1$ and $\vec{y}_2$ from part (a).
\item Find $\lim_{n \to \infty} A^n\vec{y}_0$.
\end{enumerate}

\exatend In a certain national park there are hawks that eat rabbits. Let $\vec{y}_m = (h_m,r_m)$ be the number of hawks and rabbits after $m$ months, starting with $\vec{y}_0=(100,500)$. The following simplified model describes how their populations affect each other:
\[ \bbvec{h_{m+1}}{r_{m+1}} = \begin{bmatrix}
.4 & .03 \\
-.07 & 1.06
\end{bmatrix}\bbvec{h_m}{r_m}\]
\begin{enumerate}[label=(\alph*)]
\item Find $\vec{y}_1$.
\item Suppose the rabbits evolve to grow their own spiked armor that protects them from hawks. Which of the four coefficient(s) do you think would change?
\item Suppose the hawks evolve to be able to eat berries and acorns in addition to rabbits, so they can more of them can survive even if the rabbit population is low. Which of the four coefficient(s) do you think would change?
\item It can be shown that $A=PDP^{-1}$, with $P=\begin{bmatrix}
9.383 & 0.046 \\
1 & 1
\end{bmatrix}, D=\begin{bmatrix}
0.403 & 0 \\
0 & 1.057
\end{bmatrix}, P^{-1}=\begin{bmatrix}
.107 & -.0049 \\
-.107 & 1.0049
\end{bmatrix}$. If $\mathcal{B}$ is the basis made from the columns of $P$, find $\left[ \vec{y}_0 \right]_{\mathcal{B}}$.
\item Find $\lim_{m \to \infty}A^m\vec{y}_0$, and describe what happens to the populations in the long run.
\end{enumerate}

\exatend Suppose $A=PDP^{-1}$ is $c\times c$, with $D$ diagonal. Explain why the diagonal coordinates theorem implies that for any $\vec{y}$ in $\R^c$, $A\vec{y}=PD \left[ \vec{y} \right]_{\mathcal{B}}$ where $\mathcal{B}$ is the basis made of columns of $P$.

\newpage
\section{Eigenstuff}
\label{sec:eigenstuff}

\exercise (\emph{Warm-up}) Let $A = \begin{bmatrix}
  -3.5 & 6 \\
  -0.75 & 1
\end{bmatrix}, \vec{x}= \bbvec{1}{0},$ and $\vec{v}=\bbvec{4}{1}$.
\begin{enumerate}[label=(\alph*)]
\item Find $A \vec{x}$  and $A \vec{v}$.
\item Sketch a graph of $\vec{x}$, $\vec{v}$, $A \vec{x}$, and $A \vec{v}$.
\item T/F? $A \vec{x}$ is a scalar multiple of $x$.
\item T/F? $A \vec{v}$ is a scalar multiple of $v$.
\end{enumerate}

\begin{tikzpicture}[line cap=round,line join=round]%,x=2.0cm,y=1.0cm]
\begin{axis}[
x=20pt,y=20pt,
axis lines=middle,
axis line style={-},
ymajorgrids=true,
xmajorgrids=true,
xmin=-8.5, xmax=5.5,
ymin=-3.5, ymax=3.5,
xtick={-8,-6,...,4},
ytick={-2,0,2},]
\end{axis}
\end{tikzpicture}

\begin{notes}
\vfill
\end{notes}
\exercise (\emph{Warm-up}) Suppose $\mathcal{B} = ( \vec{b}_1,\ldots, \vec{b}_c )$ is a basis of $\R^c$, and $P = \begin{bmatrix}
\vec{b}_1 & \cdots & \vec{b}_c
\end{bmatrix}$. 
\begin{enumerate}[label=(\alph*)]
\item Which of the following are equal to $\left[ \vec{b}_2 \right]_{\mathcal{B}}$? (Hint: Two are correct.)

$\vec{b}_2$ \quad $P\vec{b}_2$ \quad $P^{-1}\vec{b}_2$ \quad $\vec{e}_2$
\item Which of the following are equal to $\left[ \vec{e}_3 \right]_{\mathcal{B}}$? (Hint: Only one is correct.)

$\vec{e}_3$ \quad $P\vec{e}_3$ \quad $P^{-1}\vec{e}_3$ \quad $\vec{b}_{3}$
\item Which of the following are equal to $\vec{b}_3$? (Hint: Only one is correct.)

$\vec{e}_3$ \quad $P\vec{e}_3$ \quad $P^{-1}\vec{e}_3$ \quad $P^{-1}\vec{b}_{3}$
\end{enumerate}

\begin{notes}
\vfill
\end{notes}

\begin{definition}
Suppose $A$ is a $c\times c$ matrix. We say $\vec{x}$ is an \textbf{eigenvector} corresponding with an \textbf{eigenvalue} $\lambda$ (a scalar) of $A$, if:
\[A\vec{x}=\lambda \vec{x} \quad \text{and} \quad \vec{x}\ne\vec{0}\]
\end{definition}

\begin{notes}
\columnbreak
\end{notes}

\example \label{example:4.2-rotation} Let $A = \begin{bmatrix}
\frac{1}{2} & \frac{-\sqrt{3}}{2} \\
\frac{\sqrt{3}}{2} & \frac{1}{2}
\end{bmatrix}$. Explain why $A$ has no real eigenvalues.

\begin{notes}
\begin{tikzpicture}[line cap=round,line join=round]%,x=2.0cm,y=1.0cm]
\begin{axis}[
x=30pt,y=30pt,
axis lines=middle,
axis line style={-},
ymajorgrids=true,
xmajorgrids=true,
xmin=-2.5, xmax=2.5,
ymin=-2.5, ymax=2.5,
xtick={-2,-1,...,2},
ytick={-2,-1,...,2}]
\end{axis}
\end{tikzpicture}

\vfill
\end{notes}

\begin{book}
\textbf{Solution.} $A$ is the rotation matrix for a $60^{\circ}$ rotation; for any $\vec{x}$ in $\R^2$, $A\vec{x}$ is rotated $60^{\circ}$ away from the line $\spn(\vec{x})$. But any scalar multiples $\lambda \vec{x}$ are in $\spn(\vec{x})$. So $A \vec{x} \ne \lambda \vec{x}$.

\begin{tikzpicture}[scale=1]
\vectorgrid[xmin=-2.5,xmax=2.5,ymin=-2.5,ymax=2.5,step=1,labint=1,xu=0.5,yu=0.866,xv=-0.866,yv=0.5]
 \node[above] at (0.5,0.866) {$A\vec{e}_1$};
%\draw[thick,blue,->] (0,0) -- (-0.866,0.5);
 \node[left]at (-.866,0.5) {$A\vec{e}_2$};
\end{tikzpicture}
\blacksquare
\end{book}

\example If $A= \begin{bmatrix}
-3.5 & 6 \\
-0.75 & 1
\end{bmatrix}, P=\begin{bmatrix}
2 & -2 \\
0.5 & -1
\end{bmatrix}=\begin{bmatrix}
\vec{b}_1 & \vec{b}_2
\end{bmatrix}$, and $D = \begin{bmatrix}
-2 & 0 \\
0 & -0.5
\end{bmatrix}$, it can be shown that $A=PDP^{-1}$. With $\mathcal{B} = ( \vec{b}_1,\vec{b}_2 )$, show that $\vec{b}_2$ is an eigenvector of $A$ two ways:
\begin{enumerate}[label=(\alph*)]
\item By multiplying $A\vec{b}_2$ directly.
\item Using the diagonal coordinates theorem (\ref{theorem:diagonal-coordinates}).
\end{enumerate}

\begin{book}
\textbf{Solution.} (a) $A\vec{b}_2 = \begin{bmatrix}
-3.5 & 6 \\
-0.75 & 1
\end{bmatrix}\bbvec{-2}{-1} = \bbvec{7}{1.5} + \bbvec{-6}{-1} = \bbvec{1}{0.5} = -0.5\vec{b}_2$. So by definition, $\vec{b}_2$ is an eigenvector of $A$ corresponding with eigenvalue $-0.5$.

(b) Since $\vec{b}_2=0\vec{b}_1+1\vec{b}_2$, we have $\left[ \vec{b}_2 \right]_{\mathcal{B}} = (0,1)$. The diagonal coordinates theorem gives:
\[
\left[ A\vec{b}_2 \right]_{\mathcal{B}} = D \left[ \vec{b}_2 \right]_{\mathcal{B}} 
= \begin{bmatrix}
-2 & 0 \\
0 & -0.5
\end{bmatrix}\bbvec{0}{1} = \bbvec{0}{-0.5}
\]
Plugging $\bbvec{0}{-0.5}$ into the definition of $\left[ A\vec{b}_2 \right]_{B}$, we have:
\[A\vec{b}_2 = \begin{bmatrix}
\vec{b}_1 & \vec{b}_2
\end{bmatrix}\left[ \vec{b}_2 \right] = 0\vec{b}_1-0.5\vec{b}_2=-0.5\vec{b}_2\]
Thus $\vec{b}_2$ is an eigenvector of $A$ corresponding with eigenvalue $-0.5$. \blacksquare
\end{book}

\begin{notes}
\vfill\vfill\vfill
\end{notes}

Part (b) of the last example illustrates a connection between diagonalization and eigenstuff. The next theorem fleshes it out some more.

\begin{notes}
\columnbreak
\end{notes}

\begin{theorem}[Diagonalization Theorem.]
\label{theorem:diagonalization}

A $c \times c$ matrix $A$ is diagonalizable if and only if $A$ has $c$ linearly independent eigenvalues.

Specifically, $A = PDP^{-1}$ with $D$ diagonal if and only if the columns of $P$ are $c$ linearly independent eigenvectors of $A$, which correspond with eigenvalues in the diagonal of $D$ (in matching left-to-right order).
\end{theorem}

\begin{book}
\textbf{Proof.} (\rightarrow) This is a generalized version of the argument in part (b) of the last example. Suppose $A = PDP^{-1}$ with $D$ diagonal; write $P=\begin{bmatrix} \vec{b}_1 & \cdots & \vec{b}_c \end{bmatrix}, \mathcal{B}=\vec{b}_1,\ldots, \vec{b}_{c}$, and $D = \begin{bmatrix}
d_1 & \cdots & 0 \\
\vdots & \ddots & \vdots \\
0 & \cdots & d_c
\end{bmatrix}$. The goal is to show $A\vec{b}_i = d_i\vec{b}_i$ for any $\vec{b}_i$ in $\mathcal{B}$. By definition of $\mathcal{B}$-coordinates, $\left[ \vec{b}_i \right]_{\mathcal{B}}=\vec{e}_i$ since $\vec{b}_i= 0\vec{b}_1 + \cdots + 0\vec{b}_{i-1}+1\vec{b}_i+0\vec{b}_{i+1}+\cdots + 0\vec{b}_c$. Using the diagonal coordinates theorem (\ref{theorem:diagonal-coordinates}):
\begin{align*}
\left[ A\vec{b}_i \right]_{\mathcal{B}} = D \left[ \vec{b}_i \right] = D\vec{e}_i = ( 0, \ldots, 0, d_i, 0, \ldots, 0 )
\end{align*}
Again by definition of $\mathcal{B}$-coordinates, this means:
\[A\vec{b}_i = \begin{bmatrix}
\vec{b}_1 & \cdots & \vec{b}_c 
\end{bmatrix} \left[ \vec{b}_i \right]_{\mathcal{B}} = 0\vec{b}_1+\cdots+0\vec{b}_{i-1}+d_i\vec{b}_i+0\vec{b}_{i+1}+ \cdots + 0\vec{b}_c  = d_i\vec{b}_i
\]
This shows $\vec{b}_i$ is an eigenvector corresponding with eigenvalue $d_i$ of $A$.

(\leftarrow) Suppose $\vec{b}_1,\ldots,\vec{b}_c$ are linearly independent eigenvectors of $A$, corresponding respectively with eigenvalues $\lambda_1,\ldots,\lambda_c$. Since there are $c$ of them and they're linearly independent, the buy-two-get-one-free theorem (\ref{theorem:buy-two-get-one-free}) tells us that $\mathcal{B}=( \vec{b}_1,\ldots,\vec{b}_c )$ is a basis of $\R^c$, and the IMT says $P = \begin{bmatrix}
\vec{b}_1 & \cdots & \vec{b}_c
\end{bmatrix}$ is invertible. Let $D = \begin{bmatrix}
\lambda_1 & \cdots & 0 \\
\vdots & \ddots & \vdots \\
0 & \cdots & \lambda_c
\end{bmatrix}$ and let $\vec{y}$ be any vector in $\R^c$. If we can show that $A\vec{y}=PDP^{-1}\vec{y}$, we're done, since $\vec{y}$ can be any column of the identity matrix $\vec{e}_i$, and $A\vec{e}_i = PDP^{-1}\vec{e}_i$ means the $i^{\text{th}}$ columns of $A$ and $PDP^{-1}$ are equal. Let $\vec{x} = (x_1,\ldots,x_c) = \left[ \vec{y} \right]_{\mathcal{B}}$. Then:
\begin{align*}
A\vec{y} &= A ( x_1\vec{b}_1 + \cdots + x_c\vec{b}_c ) & (\text{defn of } \mathcal{B}\text{-coords})\\
&= x_1A\vec{b}_1 + \cdots + x_cA\vec{b}_c \\
&= x_1\lambda_1 \vec{b}_1 + \cdots + x_c\lambda_c \vec{b}_c & (\text{defn of eigenstuff})\\
&= \begin{bmatrix}
\vec{b}_1 & \cdots & \vec{b}_c
\end{bmatrix}\bbbvec{x_1\lambda_1}{\vdots}{x_c\lambda_c} \\
&= \begin{bmatrix}
\vec{b}_1 & \cdots & \vec{b}_c
\end{bmatrix}\begin{bmatrix}
\lambda_1 & \cdots & 0 \\
\vdots & \ddots & \vdots \\
0 & \cdots & \lambda_c
\end{bmatrix} \bbbvec{x_1}{\vdots}{x_c} \\ 
&= PD \left[ \vec{y} \right]_{\mathcal{B}} \\
&= PDP^{-1}\vec{y}
\end{align*}
\blacksquare
\end{book}
\begin{notes}
\columnbreak
\end{notes}

\example Suppose $A$ is $3\times3$, with two eigenvalues, $\lambda_1=4$ and $\lambda_2=5$. Every eigenvector corresponding with $\lambda_1 = 4$ is a multiple of $\vec{b}_1=(1,2,3)$ and every eigenvector corresponding with $\lambda_2=5$ is a multiple of $\vec{b}_2=(4,5,6)$. Is $A$ diagonalizable?

\begin{book}
\textbf{Solution.} $A$ has two linearly independent eigenvectors, $\vec{b}_1$ and $\vec{b}_2$. But it doesn't have a third. Any list of three eigenvectors must contain either two that are multiples of $\vec{b}_1$  or two that are multiples of $\vec{b}_2$, making a dependent list. Thus $A$ is not diagonalizable. \blacksquare
\end{book}
\begin{notes}
\vfill
\end{notes}

\example Suppose $A$ is $3\times3$, with two eigenvalues, $\lambda_1=4$ and $\lambda_2=5$. Every eigenvector corresponding with $\lambda_1 = 4$ is a multiple of $\vec{b}_1=(1,2,3)$ and the set of eigenvectors associated with $\lambda_2=5$ is $\spn ( \vec{b}_2, \vec{b}_3 )$, with $\vec{b}_2=(4,5,6)$ and $\vec{b}_3=(0,3,0)$. Is $A$ diagonalizable?

\begin{book}
\textbf{Solution.} Solving $\begin{bmatrix}
\vec{b}_1 & \vec{b}_2 & \vec{b}_3
\end{bmatrix}\vec{x}=\vec{0}$ gives only the trivial solution $\vec{x}=\vec{0}$, so we have a list of $3$ independent eigenvectors. So $A = PDP^{-1}$ is diagonalizable, with:
\[P=\begin{bmatrix}
\vec{b}_1 & \vec{b}_2 & \vec{b}_3
\end{bmatrix} = \begin{bmatrix}
1 & 4 & 0 \\
2 & 5 & 3 \\
3 & 6 & 0
\end{bmatrix} \quad \text{and}\quad D=\begin{bmatrix}
4 & 0 & 0 \\
0 & 5 & 0 \\
0 & 0 & 5
\end{bmatrix} \]
Other ways of ordering the eigenvectors and eigenvalues are fine, as long as the eigenvectors in $P$ correspond from left to right with the eigenvalues in $D$. For example, $A=PDP^{-1}$ with:
\[P=\begin{bmatrix}
\vec{b}_3 & \vec{b}_1 & \vec{b}_2
\end{bmatrix} = \begin{bmatrix}
0 & 1 & 4 \\
3 & 2 & 5 \\
0 & 3 & 6
\end{bmatrix} \quad \text{and}\quad D=\begin{bmatrix}
5 & 0 & 0 \\
0 & 4 & 0 \\
0 & 0 & 5
\end{bmatrix} \]
\blacksquare
\end{book}
\begin{notes}
\vfill 
\end{notes}

\example Let $A = \begin{bmatrix}
1 & -2 \\
0 & 1
\end{bmatrix}$, which is the matrix for a horizontal shear transformation as shown. Is $A$ diagonalizable?

\begin{tikzpicture}[scale=1]
\vectorgrid[xmin=-2.5,xmax=2.5,ymin=-2.5,ymax=2.5,step=1,labint=1,xu=1,yu=0,xv=-2,yv=1]
 \node[above] at (1,0) {$A\vec{e}_1$};
 \node[above]at (-2,1) {$A\vec{e}_2$};
\end{tikzpicture}

\begin{book}
\textbf{Solution.} Any vector on the $x_1$-axis, such as $\vec{e}_1$ is an eigenvalue with eigenvector 1. Any vector that is \emph{not} on the $x_1$-axis will be sheared horizontally away from its span, making $A\vec{x} \ne \lambda \vec{x}$. So it's not possible to get a list of $2$ linearly independent eigenvectors and $A$ is not diagonalizable. \blacksquare
\end{book}
\begin{notes}
\vspace{3em}
\columnbreak
\end{notes}

\example Show that a square matrix $A$ is invertible if and only iff 0 is an eigenvalue of $A$.

\begin{book}
\textbf{Solution.} The IMT says $A$ is invertible if and only if $\nul A = \left\{ \vec{0} \right\}$. In other words, if and only if the only solution to $A\vec{x}=\vec{0} = 0\vec{x}$ is $\vec{x}=\vec{0}$. So if $A$ is invertible, there's no eigenvector (which must be nonzero) to correspond with eigenvalue $0$, so 0 is not an eigenvalue. If $A$ is not invertible, there's a non-zero solution to $A\vec{x} = 0\vec{x}$, making 0 an eigenvalue. \blacksquare
\end{book}
\begin{notes}
\vfill
\end{notes}


This shows that part (xi) can be added to the IMT. Part (x) was theorem \ref{theorem:determinants-and-invertibility}.

\begin{theorem}[IMT v2.]
\label{theorem:IMT-v2}

Suppose $A$ is a $c\times c$ matrix. Then the following are equivalent:
\begin{enumerate}[label=(\roman*)]
\item $A$ is invertible.
\item There exists a matrix $P$ such that $PA = I_c$.
\item There exists a matrix $Q$ such that $AQ = I_c$.
\item $A$ is row-equivalent to $I_c$.
\item $A$ has $c$ pivot positions.
\item The columns of $A$ are linearly independent.
\item $\nul A = \left\{ \vec{0} \right\}$.
\item For every $\vec{y}$ in $\R^{c}$, $A\vec{x}=\vec{y}$ has a solution.
\item $\col A = \R^c $
\item $\det(A) \ne 0$
\item 0 is not an eigenvalue of $A$.
\end{enumerate}
\end{theorem}

\begin{notes}
\columnbreak
\end{notes}

\example Suppose $A$ is both diagonalizable and invertible. Show that $A^{-1}$ is diagonalizable, and find its eigenstuff.

\begin{book}
\textbf{Solution.} Let $A =PDP^{-1}$ with $D= \begin{bmatrix}
\lambda_1 & \cdots & 0 \\
\vdots & \ddots & \vdots \\
0 & \cdots & \lambda_c
\end{bmatrix}$. By the diagonalization theorem the columns of $P$ are eigenvectors corresponding respectively with eigenvalues $\lambda_1, \ldots, \lambda_c$. Since $A$ is invertible, we know 0 is not an eigenvalue, so $D$ is in echelon form with a pivot in every one of the diagonal entries, making $D$ invertible. In fact, $D^{-1}$ is diagonal with each diagonal entry in row $i$, column $i$ equal to $\frac{1}{\lambda_i}$ which you can check by multiplying:
\[\begin{bmatrix}
\lambda_1 & 0 & \cdots & 0 \\
0 & \lambda_2 & \cdots & 0 \\
\vdots & \vdots & \ddots & \vdots \\
0 & 0 & \cdots & \lambda_c
\end{bmatrix} \begin{bmatrix}
\frac{1}{\lambda_1} & 0 & \cdots & 0 \\
0 & \frac{1}{\lambda_2} & \cdots & 0 \\
\vdots & \vdots & \ddots & \vdots \\
0 & 0 & \cdots & \frac{1}{\lambda_c}
\end{bmatrix} = \begin{bmatrix}
1 & 0 & \cdots & 0 \\
0 & 1 & \cdots & 0 \\
\vdots & \vdots & \ddots & \vdots \\
0 & 0 & \cdots & 1
\end{bmatrix}\]

Now, since $P, D,$ and $P^{-1}$ are all invertible:
\[A^{-1} = ( PDP^{-1} )^{-1} = ( P^{-1} )^{-1}D^{-1}P^{-1} = PD^{-1}P^{-1} \]
Thus we've diagonalized $A^{-1}$. It has the same eigenvectors as $A$ (the columns of $P$), and the eigenvalues are the reciprocals of the eigenvalues of $A$.

\end{book}

\columnbreak
\exerciseheading

\exatend Construct a matrix $A$ with eigenvectors $\bbvec{3}{-5}$ and $\bbvec{2}{-4}$ which correspond respectively with eigenvalues $2$ and $3$.

\exatend Let $A = \begin{bmatrix}
-1.5 & 5 \\
-1 & 3
\end{bmatrix}, \vec{u}=\bbvec{1}{-2}, \vec{v}=\bbvec{-2}{5}$.
\begin{enumerate}[label=(\alph*)]

\item Show that $\vec{u}$ and $\vec{v}$ are eigenvectors of $A$ and find corresponding eigenvalues.
\item Diagonalize $A$ (Find $P$ and a diagonal matrix $D$ so that $A = PDP^{-1}$.)
\item Find $A^{1000000}$.
\end{enumerate}

TODO: Problem with vector grid and several powers of $A\vec{x}_0$ with different starting $\vec{x}_{0}$

TODO: More basic T/F problems.

\exatend (T/F?) If $A$ is $2\times 2$ with eigenvectors $\vec{u}$ and $\vec{v}$ that correspond respectively to eigenvalues $\mu$ and $\lambda$, then $A^n ( x_1\vec{u} +x_2\vec{v} ) = x_1\mu^n \vec{u} +x_2\lambda^n \vec{v}$.

\exatend (T/F?) If $A$ is diagonalizable, then the diagonalization is unique. In other words, if $A=P_1D_1P_1^{-1}=P_2D_2P_2^{-1}$, then $P_1=P_2$ and $D_1=D_2$.

\exatend (T/F?) If $0$ is an eigenvalue of $A$, then $A$ is invertible.

\exatend (T/F?) If $A$ is square, then $\vec{0}$ is an eigenvector of $A$ corresponding with eigenvalue $\frac{-1}{12}$ because $A\vec{0}=\frac{-1}{12}\vec{0}$.

\exatend (T/F?) If $AP = PD$ with $D$ diagonal, then the columns of $P$ are eigenvectors of $A$.

\exatend (T/F?) If $AP = PD$ with $D$ diagonal, then $P$ is invertible.

\exatend Let $A$ be the matrix that makes the transformation $T(\vec{x})=A\vec{x}$ a reflection across the $x_2$-axis. Diagonalize $A$. \hintfootnote{Drawing a picture might help you find eigenstuff.}

\exatend Let $A$ be the matrix that makes the transformation $T(\vec{x})=A\vec{x}$ a reflection across the line $x_2 =x_1$. Diagonalize $A$. \hintfootnote{Drawing a picture might help you find eigenstuff.}

\exatend Let $A$ be the $2\times 2$ matrix so that $T(\vec{x})=A\vec{x}$ projects points onto the nearest point on the $x_2$-axis. For example, $A \bbvec{3}{4}=\bbvec{0}{4}$. Diagonalize $A$. \hintfootnote{Drawing a picture will help finding eigenstuff.}

\exatend Let $A$ be the $2\times 2$ matrix so that $T(\vec{x})= A\vec{x}$ projects points onto the nearest point on the line $x_2=3x_1$. Find $A$ by finding the eigenstuff of $A$ and multiplying $PDP^{-1}$.\hintfootnote{Drawing a picture will help finding eigenstuff.}

\exatend Let $A$ be the $3\times 3$ matrix such that $T(\vec{x}) = A\vec{x}$ rotates points  $30^{\circ}$ around the $x_2$-axis. Find an eigenvector and its eigenvalue. Is $A$ diagonalizable?

\exatend \label{exercise:4.2-rotation} Let $R_{\theta}$ be the $2\times 2$ matrix such that $T(\vec{x})= R_{\theta}\vec{x}$ rotates points an angle $\theta$ around the origin. What values of $\theta$ make $R_{\theta}$ diagonalizable?

\exatend Let $A = 0$ be a square matrix with all zero entries. Is $A$ diagonalizable?

%credit: Kuttler
\exatend Let $A,B$ be square matrices such that $AB = BA$ (ie, they \emph{commute}). Suppose $\vec{x}$ is an eigenvector of $B$. Show that $A\vec{x}$ is also an eigenvector of $B$. \creditfootnote{Inspired by an exercise in Ken Kuttler's \emph{A First Course in Linear Algebra}.}

\exatend Suppose $A$ has eigenvalue $\lambda$ corresponding with eigenvector $\vec{x}$. Find an eigenvalue and corresponding eigenvector of $kA$, where $k$ is any nonzero constant.

\exatend Give an example of a matrix that is diagonalizable but not invertible.

\exatend Give an example of a matrix that is invertible but not diagonalizable.


\exatend Prove or find a counterexample: If $A$ is diagonalizable, then $A^T$ is diagonalizable.

\exatend Prove or find a counterexample: If $A$ is diagonalizable, then $A^T$ has the same eigenvalues as $A$.

\exatend Prove or find a counterexample: If $A$ is diagonalizable, then $A^T$ has the same eigenvectors as $A$.

\exatend Suppose $A$ is $3\times 3$, and  every vector in $\R^{3}$ is an eigenvector corresponding with eigenvalue 7. Show that $A = 7I$.

\exatend Show that if $A$ is $c\times c$ and has $c$ distinct eigenvalues, then $A$ is diagonalizable.

%credit: Kuttler
\exatend Let $A$ be a square matrix such that $A^n=A$ for some positive integer $n>1$. Show that if $\lambda$ is an eigenvalue of $A$, then $|\lambda|$ equals either $-1,0,$ or $1$. \creditfootnote{Inspired by K. Kuttler.}

\exatend \label{exercise:4.2-rows-sum-to-k}Suppose that $A$ is a $c\times c$ matrix such that every row of $A$ adds up to the same number $k$. Show that $k$ is an eigenvalue of $A$ by finding a corresponding eigenvector.

\newpage
\section{Finding Eigenstuff}
\label{sec:finding-eigenstuff}

\exercise (\emph{Warm-up}) If $A$ is an invertible $c \times c$ matrix...
\begin{enumerate}[label=(\alph*)]
\item How many pivot positions does $A$ have?
\item How many solutions does $A \vec{x} = \vec{0}$ have?
\item What do you know about the determinant of $A$?
\end{enumerate}

\vspace{.5in}
\exercise (\emph{Warm-up}) Write the null space of each matrix as the span of some vectors. (Review section \ref{sec:homogeneous-systems-and-linear-independence} if necessary.)
\begin{horiparts}{2}
\item $ A = \begin{bmatrix}
2 & 3 & 4 \\
0 & -1 & -\pi \\
0 & 0 & 13
\end{bmatrix}$
\item $M = \begin{bmatrix}
2 & 4 & 8 \\
0 & -4 & 2 \\
0 & 0 & 0
\end{bmatrix}$
\end{horiparts}

\begin{notes}
\vfill 
\end{notes}

We learned a lot about the usefulness and properties of eigenstuff in the last two sections, but we haven't learned how to find eigenvectors and eigenvalues if you don't already know them. Let's start by finding an eigenvector, given an eigenvalue.

\begin{notes}
\columnbreak
\end{notes}

\example Show that $2$ is an eigenvalue of $A = \begin{bmatrix}
-2 & 3 \\
-8 & 8
\end{bmatrix}$ by finding a corresponding eigenvector.

\begin{book}
\textbf{Solution.} The goal is to find $\vec{x}$ such that $A\vec{x} = 2\vec{x}$. If you didn't know any better, you might try to move the $2\vec{x}$ to the left side, factor out $A-2$, and divide to get $\vec{x}=\frac{0}{A-2}$. But sadly, $A-2$ is undefined. However, we can get around this inconvenience by noticing that $2\vec{x} = 2I\vec{x}$:
\begin{align*}
A\vec{x}&=2\vec{x} \\
A \vec{x}-2\vec{x}&=\vec{0} \\
A \vec{x}-2I\vec{x}&=\vec{0} \\
( A-2I )\vec{x}&=\vec{0}\\
\left( \begin{bmatrix}
-2 & 3 \\
-8 & 8
\end{bmatrix} - \begin{bmatrix}
2 & 0 \\
0 & 2 
\end{bmatrix} \right)\vec{x}&= \vec{0} \\
\begin{bmatrix}
-4 & 3 \\
-8 & 6 
\end{bmatrix}\vec{x} &= \vec{0} \\
\begin{bmatrix}
1 & \frac{-3}{4} \\
0 & 0 
\end{bmatrix}\vec{x}&=\vec{0}\\
\vec{x}= \bbvec{x_1}{x_2} &= \bbvec{\frac{3}{4}x_2}{x_{2}}  = x_2 \bbvec{\frac{3}{4}}{1}
\end{align*}
Thus any vector in $\spn \left( \bbvec{\frac{3}{4}}{1} \right)$, other than $\vec{0}$, is an eigenvector corresponding to eigenvalue $2$.\blacksquare
\end{book}
\begin{notes}
\columnbreak
\end{notes}

This idea works in general; the following equations are equivalent:
\begin{align*}
A\vec{x}&=\lambda \vec{x} \\
A \vec{x}- \lambda \vec{x} &= \vec{0}\\
A \vec{x} - \lambda I \vec{x} &= \vec{0}\\
( A-\lambda I ) \vec{x}&= \vec{0}
\end{align*}
This proves:

\begin{theorem}[Alternative definition of eigenstuff.]
\label{theorem:alt-defn-eigenstuff}

Suppose $A$ is a $c\times c$ matrix. Then $\vec{x}$ is an eigenvector corresponding with an eigenvalue $\lambda$ (a scalar) of $A$, if and only if:
\[( A-\lambda I )\vec{x}=\vec{0} \quad \text{and} \quad \vec{x}\ne\vec{0}\]
\end{theorem}

This means that, for any eigenvalue $\lambda$ of $A$, the set of eigenvectors corresponding with $\lambda$ equals the null space of $A-\lambda I$ (\emph{except $\vec{0}$ which is in $\nul ( A-\lambda I )$ but isn't an eigenvector}). So, the vector space $\nul ( A-\lambda I )$ is called the \textbf{eigenspace} of $A$ corresponding with $\lambda$.

\begin{book}
Now we know how to find eigenvectors given eigenvalues. But how do you find eigenvalues of $A$, if you don't know any eigenstuff at all? Recall from the IMT that $\nul(A) \ne \left\{ \vec{0} \right\}$ if and only if $\det (A) = 0$; both are equivalent to $A$ being \emph{not} invertible. Applying this to $A-\lambda I$, the following are equivalent:
\begin{itemize}
\item $\lambda$ is an eigenvalue of $A$.
%\item There's a non-trivial solution of $( A-\lambda I )\vec{x}=\vec{0}$.
\item $\nul ( A-\lambda I )\ne \left\{ \vec{0} \right\}$.
\item $\det ( A-\lambda I ) = 0$.
\end{itemize}
\end{book}

\begin{notes}
Now we know how to find eigenvectors given eigenvalues. But how do you find eigenvalues of $A$, if you don't know any eigenstuff at all? Fill in the blanks to make equivalent statements, using the IMT:
\begin{itemize}
\item $A$ is \emph{not} invertible.
\item $\nul(A) \quad \quad \left\{ \vec{0} \right\}$.
\item $\det(A) \quad \quad 0$.
\end{itemize}
 Applying this to $A-\lambda I$, the following are equivalent:
\begin{itemize}
\item $\lambda$ is an eigenvalue of $A$.
\item $\nul ( A-\lambda I ) \quad \quad \left\{ \vec{0} \right\}$.
\item $\det ( A-\lambda I ) \quad \quad 0$.
\end{itemize}
\end{notes}

This shows that solutions (values of $\lambda$) that are solutions of the \emph{characteristic equation} are eigenvalues of $A$:

\begin{definition}
The \textbf{characteristic equation} of $A$ is $\det ( A-\lambda I ) = 0$.
\end{definition}

\begin{notes}
\columnbreak
\end{notes}
\example Find the eigenstuff of the stochastic matrix $A = \begin{bmatrix}
0.9 & 0.05 \\
0.1 & 0.95
\end{bmatrix}$.

\begin{book}
\textbf{Solution.} We'll start by finding $\det ( A- \lambda I )$ and setting it to zero:
\begin{align*}
\det \left(\begin{bmatrix}
0.9 & 0.05 \\
0.1 & 0.95
\end{bmatrix} - \begin{bmatrix}
\lambda & 0 \\
0 & \lambda
\end{bmatrix} \right) &= \det \begin{bmatrix}
0.9-\lambda & 0.05 \\
0.1 & 0.95 - \lambda
\end{bmatrix} \\ 
&= (0.9-\lambda)(0.95-\lambda) - 0.1(0.05) \\
&= \lambda^2 - 1.85\lambda+0.85 = 0
\end{align*}
This characteristic equation can be solved with the quadratic formula, but there's a trick to factoring it:
\begin{align*}
(\lambda-1)(\lambda-0.85) &= 0
\end{align*}
So the eigenvalues are $\lambda_1=1$ and $\lambda_2=0.85$. Let's find eigenvectors $\vec{x}$ corresponding with $\lambda_1$:
\begin{align*}
\begin{bmatrix}
0.9-1 & 0.05 \\
0.1 & 0.95 - 1
\end{bmatrix}\vec{x} &= \vec{0} \\
\begin{bmatrix}
-0.1 & 0.05 \\
0.1 & -0.05
\end{bmatrix}\vec{x} &= \vec{0} \\
\begin{bmatrix}
1 & -0.5 \\
0 & 0
\end{bmatrix}\vec{x} &= \vec{0} \\
\vec{x}= \bbvec{x_1}{x_2} &= \bbvec{0.5x_2}{x_2} = x_2 \bbvec{0.5}{1}
\end{align*}
So the eigenspace corresponding to $\lambda_1=1$ is $\spn \left( \bbvec{0.5}{1} \right)$.
Similarly with $\lambda_2=0.85$:
\begin{align*}
\begin{bmatrix}
0.9-0.85 & 0.05 \\
0.1 & 0.95 - 0.85
\end{bmatrix}\vec{x} &= \vec{0} \\
\begin{bmatrix}
0.05 & 0.05 \\
0.1 & 0.1
\end{bmatrix}\vec{x} &= \vec{0} \\
\begin{bmatrix}
1 & 1 \\
0 & 0
\end{bmatrix}\vec{x} &= \vec{0} \\
\vec{x}= \bbvec{x_1}{x_2} &= \bbvec{-x_2}{x_2} = x_2 \bbvec{-1}{1}
\end{align*}
So the eigenspace corresponding to $\lambda_2=0.85$ is $\spn \left( \bbvec{-1}{1} \right)$.
\end{book}
\begin{notes}
\vfill
\end{notes}

\example Using the characteristic equation, find the eigenvalues of the stochastic matrix from the infectious disease example from section \ref{sec:review-and-diagonalization}. Then find an eigenvector corresponding with $\lambda=1$: 
\[A = \begin{bmatrix}
0.92 & 0 & 0 \\
0.08 & 0.6 & 0.01 \\
0 & 0.4 & 0.99
\end{bmatrix}\]

\begin{book}
\textbf{Solution.} First let's set up the characteristic equation:
\begin{align*}
\det \begin{bmatrix}
0.92 - \lambda & 0 & 0 \\
0.08 & 0.6 - \lambda & 0.01 \\
0 & 0.4 & 0.99 - \lambda
\end{bmatrix} = 0
\end{align*}
The first row of zeros makes for an easy cofactor expansion:
\begin{align*}
\det \left( (.92-\lambda)\det \begin{bmatrix}
.6-\lambda & .01 \\
.4 & .99-\lambda
\end{bmatrix}\right)-0+0 &= 0 \\
(.92-\lambda)( (.6-\lambda)(.99-\lambda) -.4(.01) ) &= 0 \\
(.92-\lambda)( \lambda^2-1.59\lambda + 0.59 ) &= 0 \\
(.92-\lambda)(\lambda-.59)(\lambda-1) &= 0 
\end{align*}
So the eigenvalues are $0.92, 0.59,$ and $1$. To find an eigenvector corresponding with $\lambda=1$, we need a solution of:
\begin{align*}
\begin{bmatrix}
0.92 - 1 & 0 & 0 \\
0.08 & 0.6 - 1 & 0.01 \\
0 & 0.4 & 0.99 - 1
\end{bmatrix}\vec{x} &= \vec{0} \\
\begin{bmatrix}
-0.08 & 0 & 0 \\
0.08 & -0.4 & 0.01 \\
0 & 0.4 & -0.01
\end{bmatrix}\vec{x} &= \vec{0} 
\end{align*}
You could row reduce this, or if you notice the second column is $-40$ times the third, you'll see that any scalar multiple of $\vec{x}=(0,1,40)$ is a solution. For extra credit, let's find an eigenvector that satisfies the conditions on the original island, where there was an unchanging population of 1000 people:
\begin{align*}
x_{1} + x_2+ x_3 &= 1000 \\
0 + x_2 + 40x_2 &= 1000 \\
x_2 &= \frac{1000}{41} \approx 24.39 \\
\vec{x}&= (0,x_2,40x_2) \approx (0,24.39,975.61)
\end{align*}
This is indeed the $\vec{x}$ from part (d) of the original problem that made $A\vec{x}=\vec{x}$. \blacksquare
\end{book}
\begin{notes}
\vfill \hspace{1pt}
\columnbreak
\end{notes}

\example In section \ref{sec:eigenstuff} example \ref{example:4.2-rotation} we saw that the matrix $A = \begin{bmatrix}
\frac{1}{2} & \frac{-\sqrt{3}}{2} \\
\frac{\sqrt{3}}{2} & \frac{1}{2}
\end{bmatrix}$ has no \emph{real} eigenvalues because multiplying by $A$ rotates all vectors (except $\vec{0}$) away from their span by $60^{\circ}$. Diagonalize $A$, by finding the \emph{complex} eigenvalues of $A$.

\begin{book}
\textbf{Solution.} The characteristic equation is:
\begin{align*}
\det \begin{bmatrix}
\frac{1}{2}-\lambda & \frac{-\sqrt{3}}{2} \\
\frac{\sqrt{3}}{2} & \frac{1}{2}-\lambda
\end{bmatrix} &= 0\\
\frac{1}{4} - \lambda + \lambda^2 + \frac{3}{4} &= 0 \\ 
\lambda^2-\lambda+1 &= 0
\end{align*}
The quadratic formula then gives:
\[\lambda = \frac{1 \pm \sqrt{1-4}}{2} = \frac{1}{2} \pm \frac{\sqrt{3}}{2}i \]
Let's find an eigenvector $\vec{x}_1$ corresponding with $\lambda_1=\frac{1}{2}+\frac{\sqrt{3}}{2}i$:
\begin{align*}
\begin{bmatrix}
\frac{1}{2}- \left( \frac{1}{2} + \frac{\sqrt{3}}{2} \right) & \frac{-\sqrt{3}}{2} \\
\frac{\sqrt{3}}{2} & \frac{1}{2}-\left( \frac{1}{2} + \frac{\sqrt{3}}{2} \right)
\end{bmatrix}\vec{x}_1 &= \vec{0}\\
\begin{bmatrix}
\frac{-\sqrt{3}}{2}i & \frac{-\sqrt{3}}{2} \\
\frac{\sqrt{3}}{2} & \frac{-\sqrt{3}}{2}i
\end{bmatrix}
\end{align*}
By inspection, $\vec{x}_1=(i,1)$ works. Similarly, an eigenvalue corresponding with  $\lambda_2=\frac{1}{2}-\frac{\sqrt{3}}{2}i$ can be found to be $\vec{x}_2=(-i,1)$. So with $P = \begin{bmatrix}
\vec{x}_1 & \vec{x}_2
\end{bmatrix} = \begin{bmatrix}
i & -i \\
1 & 1
\end{bmatrix}$, we have $P^{-1}=\begin{bmatrix}
\frac{-i}{2} & \frac{1}{2} \\
\frac{i}{2} & \frac{1}{2}
\end{bmatrix}$, and with the corresponding eigenvalues in $D = \begin{bmatrix}
\frac{1}{2}+\frac{\sqrt{3}}{2}i & 0 \\
0 & \frac{1}{2}-\frac{\sqrt{3}}{2}i
\end{bmatrix}$, we have $A = PDP^{-1}$. \blacksquare
\end{book}
\begin{notes}
\vfill
\end{notes}

In general, multiplication by complex numbers is an elegant and powerful way to represent rotation. In fact, it can be shown that $e^{at}(cos(bt) + i\sin(bt)) = e^{(a+bi)t}$, and functions of this form are solutions of many differential equations. And an important class of differential equations, called \emph{linear} differential equations, can be written in the form $\vec{y}'(t) = A\vec{y}(t)$, which are analogous to our linear difference equations, $\vec{y}_{n+1}=A\vec{y}_n$.

For this book, the above example (and a single exercise) is the only time we'll talk about complex numbers. But in a slightly more advanced linear algebra course, you'll learn that \emph{every} $c\times c$ matrix has $c$ complex eigenvalues (counting repeated factors in the characteristic equation); and there's a similar result for the analogous \emph{differential} equations.

\begin{notes}
\columnbreak
\end{notes}

The next theorem describes a nice way that eigenspaces are distinct from each other. It will be useful in chapter \ref{chap:symmetric-matrices}.

\begin{theorem}
\label{theorem:distinct-eigenvectors-indep}
If $A$ is a square matrix with eigenvectors $\vec{b}_1, \ldots, \vec{b}_n$ corresponding to \emph{distinct} eigenvalues $\lambda_1,\ldots,\lambda_n$, then $\vec{b}_1,\ldots,\vec{b}_n$ is a linearly independent list.
\end{theorem}

\begin{book}
\textbf{Proof.} We'll prove the contrapositive: that if $\vec{b}_1,\ldots,\vec{b}_n$ is a linearly \emph{dependent} list, then $\lambda_1,\ldots,\lambda_n$ are \emph{not} distinct. Suppose $\vec{b}_1,\ldots,\vec{b}_n$ is a linearly dependent list of eigenvectors. Since $\vec{b}_1$ is an eigenvector, it's not $\vec{0}$, so $( \vec{b}_1 )$ is an independent list. You can add $\vec{b}_2$ to the list, and maybe the list $\vec{b}_1,\vec{b}_2$ is still independent, or maybe not. But we can keep adding vectors, one at a time in order, and at some point the list changes from independent to dependent. Let $\vec{b}_i$ be the just before that; in other words, let $i$ be the greatest integer such that $\vec{b}_1,\ldots, \vec{b}_i$ is independent, which means $\vec{b}_1,\ldots, \vec{b}_i,\vec{b}_{i+1}$ is dependent. By theorem \ref{theorem:characterization-of-linear-dependence}, one of the vectors $\vec{b}_1,\ldots, \vec{b}_i,\vec{b}_{i+1}$ must be a linear combination of the previous vectors, and it must be $\vec{b}_{i+1}$ since the previous vectors make an independent list. So there exists $x_1,\ldots,x_i$ such that:
\begin{equation}
\label{eq:4.3-indep-1}
x_1\vec{b}_1+ \ldots + x_i\vec{b}_i = \vec{b}_{i+1}
\end{equation}
Since $\vec{b}_{i+1}$ is an eigenvector, it's not $\vec{0}$, so the $x_1,\ldots,x_i$ can't all be zero. A useful proof technique when dealing with eigenstuff is to try multiplying eigenvectors by $A$, knowing it's the same as multiplying them by the corresponding eigenvalues. Let's do that with equation \ref{eq:4.3-indep-1}:
\begin{align}
\label{eq:4.3-indep-2}
A ( x_1\vec{b}_1+ \cdots + x_i\vec{b}_i ) &= A(\vec{b}_{i+1} ) \nonumber \\
x_1A\vec{b}_1+ \cdots + x_iA\vec{b}_i &= A\vec{b}_{i+1} \nonumber \\
x_1\lambda_1\vec{b}_1+ \cdots + x_i\lambda_i\vec{b}_i &= \lambda_{i+1}\vec{b}_{i+1} 
\end{align}
Multiplying both sides of \ref{eq:4.3-indep-1} by $\lambda_{i+1}$ gives:
\begin{equation}
\label{eq:4.3-indep-3}
x_1\lambda_{i+1}\vec{b}_1+ \cdots + x_i\lambda_{i+1}\vec{b}_i = \lambda_{i+1}\vec{b}_{i+1}
\end{equation}
Equations \ref{eq:4.3-indep-2} and \ref{eq:4.3-indep-3} give expressions both equal to $\lambda_{i+1}\vec{b}_{i+1}$. Since $\lambda_{i+1}\vec{b}_{i+1} - \lambda_{i+1}\vec{b}_{i+1} = \vec{0}$, we have:
\begin{align*}
\vec{0}&=(x_1\lambda_1\vec{b}_1+ \cdots + x_i\lambda_i\vec{b}_i) - ( x_1\lambda_{i+1}\vec{b}_1+ \cdots + x_i\lambda_{i+1}\vec{b}_i ) \\
\vec{0}&= x_{1}(\lambda_1-\lambda_{i+1})\vec{b}_1 + \cdots + x_i(\lambda_i-\lambda_{i+1})\vec{b}_i 
\end{align*}
Now, since $\vec{b}_1,\cdots,\vec{b}_i$ is independent, we know all the scalars are zero:
\begin{align*}
c_1(\lambda_{i+1}-\lambda_1) &= 0 \\
&\vdots \\
c_i(\lambda_{i+1}-\lambda_i) &= 0
\end{align*}
But since not all of $x_1,\ldots,x_i$ are zero, we have some $x_k \ne 0$, with a corresponding $\lambda_{i+1}-\lambda_k=0$, so the eigenvalues $\lambda_1,\ldots,\lambda_n$ are not all distinct. \blacksquare
\end{book}

\columnbreak
\exerciseheading

\exatend Find eigenvalues and corresponding eigenvectors of $A = \begin{bmatrix}
5 & -2 \\
3 & -2
\end{bmatrix}$.

\exatend Find eigenvalues and corresponding eigenvectors of $A = \begin{bmatrix}
6 & 3 \\
-4 & -1
\end{bmatrix}$.

\exatend Find eigenvalues and corresponding eigenvectors of $A = \begin{bmatrix}
2 & 5 \\
-7 & -10
\end{bmatrix}$.

\exatend (\emph{Challenge}) Find eigenvalues and corresponding eigenvectors, which will involve complex numbers, of $A = \begin{bmatrix}
2 & -5 \\
1 & -2
\end{bmatrix}$.

\exatend Let $R_{\theta} = \begin{bmatrix}
\cos \theta & -\sin \theta \\
\sin \theta & \cos \theta
\end{bmatrix}$ be the rotation matrix that rotates points an angle $\theta$ around the origin. Find the characteristic polynomial. For what values of $\theta$ does $R_{\theta}$ have real eigenvalues? (If you did section \ref{sec:eigenstuff} exercise \ref{exercise:4.2-rotation}, check that your answers agree.)

\exatend Find eigenvalues and corresponding eigenvectors of $A = \begin{bmatrix}
0 & 2 &-2 \\
0 & 3 & 2 \\
0 & 2 & -2
\end{bmatrix}$.

\exatend Write the characteristic equation of $A = \begin{bmatrix}
1 & 2 & 3 \\
4 & 5 & 6 \\
7 & 8 & 9
\end{bmatrix}$. No need to solve or even simplify. The point of this problem is to show you that for most matrices bigger than $2\times 2$, the characteristic equation is an inefficient way to find eigenstuff. 

\exatend (T/F?) If $\lambda$ is an eigenvalue of $A$, then every vector in $\nul ( A-\lambda I )$ is an eigenvector of $A$.

\exatend (T/F?) If $A$ has eigenvalue $3$, then $A-3I$ is not invertible.

\exatend (T/F?) If $\det A = 0$, then $0$ is an eigenvalue of $A$.

\exatend (T/F?) If $A$ has eigenvalues $3$ and $6$, corresponding respectively with eigenvectors $\vec{x}$ and $\vec{y}$, then $\vec{y}=2\vec{x}$.

\exatend (T/F?) If $A$ has eigenvalues $3$ and $6$, corresponding respectively with eigenvectors $\vec{x}$ and $\vec{y}$, then $\vec{x}$ and $\vec{y}$ are linearly independent.

\exatend Suppose $A$ is a $5\times 5$ matrix with two eigenvalues. One eigenspace is 4-dimensional. Can you conclude if $A$ is diagonalizable or not?

\exatend Suppose $B$ is a $6\times 6$ matrix with $4$ eigenvalues. One eigenspace is 2-dimensional. Can you conclude if $B$ is diagonalizable or not?

\exatend Prove theorem \ref{theorem:triangular-eigenvalues} below.

\begin{theorem} \label{theorem:triangular-eigenvalues}
If $A$ is triangular, then the eigenvalues of $A$ are the entries on the main diagonal of $A$.
\end{theorem}

\exatend Find $P$ and $D$ so that $A = PDP^{-1}$ for $A = \begin{bmatrix}
3 & 0 & 0 \\
0 & 3 & 0 \\
-7 & -1 & -2
\end{bmatrix}$.

\exatend Show that $A$ is not diagonalizable: $A = \begin{bmatrix}
3 & 0 & 0 \\
7 & 3 & 0 \\
-7 & -1 & -2
\end{bmatrix}$.

%credit: treil
\exatend Find a matrix $R$ so that $R^2=A = \begin{bmatrix}
4 & \frac{3}{2} \\
0 & 1
\end{bmatrix}$. \hintfootnote{Diagonalize $A$. It's easy to find $M$ so that $M^2=D$. You can leave your answer as a product.} \creditfootnote{Inspired by an exercise in Sergei Treil's excellent \emph{Linear Algebra Done Wrong}.}

%could really be in 4.2 but here works
\exatend In section \ref{sec:matrix-multiplication} exercise \ref{exercise:2.2-diagonal-commute} you showed that if $D_1$ and $D_2$ are diagonal matrices, then $D_1D_2=D_2D_1$. Show that if $A$ and $B$ are $c\times c$ matrices that share a set of $c$ linearly independent eigenvectors, then $AB = BA$.

\exatend \label{exercise:4.3-A-AT-same-eigenvalues} Show that $\lambda$ is an eigenvalue of $A$ if and only if $\lambda$ is an eigenvalue of $A^T$. (Note $A$ might not be diagonalizable. \hintfootnote{How are $A-\lambda I$ and $A^T - \lambda I$ related?}

\exatend Use the results of exercise \ref{exercise:4.2-rows-sum-to-k} of section \ref{sec:eigenstuff} and exercise \ref{exercise:4.3-A-AT-same-eigenvalues} above to show that every stochastic matrix (whose columns all add up to 1) has 1 as an eigenvalue.

\exatend Show that if $A$ is similar to $B$, then $A$ and $B$ have the same eigenvalues. \hintfootnote{Show they have the same characteristic equation.}

\exatend Show that every $3 \times 3$ matrix $A$ has at least one real eigenvalue. 

\exatend Show that if $A$ is a non-zero diagonalizable matrix, then $A^2 \ne 0$.

\bbtwocol{0.6}{0.4}{
\exatend The rectangle illustrates $\phi$, the \emph{golden ratio}, which makes the ratio of the length to the width of the big rectangle equal to the ratio of length to width of the small rectangle. Find $\phi$.}{\quad
\begin{tikzpicture}[scale=1.6]
\draw[thick] (0,0) rectangle (1,1);
\draw[thick] (1,0) -- (1.618,0) -- (1.618,1)--(1,1);
\draw[decorate,decoration={brace,amplitude=5pt}] 
    (0,1) -- (1.618,1) node[midway,above=3pt] {$\phi$};
\node[left] at (0,0.5) {$1$};
\draw[decorate,decoration={brace,amplitude=5pt,mirror}] 
    (0,0) -- (1,0) node[midway,below=3pt] {$1$};
\draw[decorate,decoration={brace,amplitude=5pt,mirror}] 
    (1,0) -- (1.618,0) node[midway,below=3pt] {$\phi-1$};
\end{tikzpicture}}

\exatend (\emph{Challenge}) The Fibonacci sequence is defined as follows: $F_0=0, F_1=1$ and  $F_{n+2}=F_{n-1}+F_{n}$ for $n \ge 2$. This gives the sequence:
\[0,1,1,2,3,5,8,13,21,34,\ldots\]
\begin{enumerate}[label=(\alph*)]
\item Find $A$ so that $\bbvec{F_{n+2}}{F_{n+1}} = A \bbvec{F_{n+1}}{F_n}$.
\item Use the characteristic equation to find the eigenvalues of $A$. (Notice that one of them is the golden ratio $\phi$ and the other is $1-\phi$. The rest of this problem will be easier if you write them this way.)
\item Diagonalize $A$.
\item Use the diagonalization and the fact that $\bbvec{F_{n+1}}{F_n}=A^n \bbvec{F_1}{F_0} = A^n \bbvec{1}{0}$ to find a formula for $F_n$. (You'll have to find inverses and multiply some matrices. It gets pretty messy and there are many correct ways to write the formula. Check that your answer works with, say, $F_4=3$.)
\item Show $\lim_{n \to \infty}\frac{F_{n+1}}{F_n} = \phi$.
\end{enumerate}


\chapter{Orthogonality and Projections}
\label{chap:dot-products-etc}

\begin{tikzpicture}
\begin{axis}[
x=15pt,y=15pt,
axis lines=middle,
axis line style={-},
domain=0:20,
xmin=0,xmax=20,
ymin=0,ymax=6,
samples=200,]
%grid=major,]
\pgfplotstableread{
x y
3.293 4.923
3.821 4.268
4.490 3.193
5.800 1.620
6.114 1.675
6.176 0.999
6.435 1.411
6.596 1.185
6.965 1.385
7.648 1.137
9.223 3.808
9.326 3.002
9.709 3.747
11.230 4.525
14.538 1.633
15.266 0.853
15.470 1.140
15.716 0.794
18.514 3.611
19.327 4.165
}\datatable
\addplot[blue, thick] {3 + 2*sin(deg(0.7*x))};
%\addlegendentry{$y=\beta_0+\beta_1\sin(\beta_2x)$}
\addplot[only marks, mark=*, mark options={scale=0.5}, red] table {\datatable};
\draw[red] (axis cs:3.29, 4.92) -- (axis cs:3.29, {3+2*sin(deg(0.7*3.29))});
\draw[red] (axis cs:3.82, 4.27) -- (axis cs:3.82, {3+2*sin(deg(0.7*3.82))});
\draw[red] (axis cs:4.49, 3.19) -- (axis cs:4.49, {3+2*sin(deg(0.7*4.49))});
\draw[red] (axis cs:5.8, 1.62) -- (axis cs:5.8, {3+2*sin(deg(0.7*5.8))});
\draw[red] (axis cs:6.11, 1.68) -- (axis cs:6.11, {3+2*sin(deg(0.7*6.11))});
\draw[red] (axis cs:6.18, 0.999) -- (axis cs:6.18, {3+2*sin(deg(0.7*6.18))});
\draw[red] (axis cs:6.44, 1.41) -- (axis cs:6.44, {3+2*sin(deg(0.7*6.44))});
\draw[red] (axis cs:6.60, 1.19) -- (axis cs:6.60, {3+2*sin(deg(0.7*6.60))});
\draw[red] (axis cs:6.96, 1.39) -- (axis cs:6.96, {3+2*sin(deg(0.7*6.96))});
\draw[red] (axis cs:7.65, 1.14) -- (axis cs:7.65, {3+2*sin(deg(0.7*7.65))});
\draw[red] (axis cs:9.22, 3.81) -- (axis cs:9.22, {3+2*sin(deg(0.7*9.22))});
\draw[red] (axis cs:9.33, 3.00) -- (axis cs:9.33, {3+2*sin(deg(0.7*9.33))});
\draw[red] (axis cs:9.71, 3.75) -- (axis cs:9.71, {3+2*sin(deg(0.7*3.75))});
\draw[red] (axis cs:11.23, 4.53) -- (axis cs:11.23, {3+2*sin(deg(0.7*11.23))});
\draw[red] (axis cs:14.54, 1.63) -- (axis cs:14.54, {3+2*sin(deg(0.7*14.54))});
\draw[red] (axis cs:15.27, 0.85) -- (axis cs:15.27, {3+2*sin(deg(0.7*15.27))});
\draw[red] (axis cs:15.47, 1.14) -- (axis cs:15.47, {3+2*sin(deg(0.7*15.47))});
\draw[red] (axis cs:15.71, 0.79) -- (axis cs:15.71, {3+2*sin(deg(0.7*15.71))});
\draw[red] (axis cs:18.51, 3.61) -- (axis cs:18.51, {3+2*sin(deg(0.7*18.51))});
\draw[red] (axis cs:19.33, 4.17) -- (axis cs:19.33, {3+2*sin(deg(0.7*19.33))});
\end{axis}
\end{tikzpicture}

The above plot shows 20 data points $(x_1,y_1),\ldots,(x_{20},y_{20})$, taken from some real-world measurements, in red. You, a scientist, think there is a periodic or sinusoidal phenomenon here, so you want to find a function of the form $f(x)=\beta_0+\beta_1\sin(0.7x)$ to model the phenomenon. Or maybe you want this function because you're an electrical engineer doing signal processing. In any case, the goal is find $\beta_0,\beta_1$ values that make all of the differences $y_i-f(x_i)$ as small as possible. You can think of this as the distance, in $\R^{20}$, between the two points $(y_1,\ldots,y_{20})$  and $(f(x_1),\ldots,f(x_{20}))$. If you can picture it, the smallest distance from a point in a room to the nearest wall is in a direction \emph{perpendicular} to the wall. The same principle applies in $\R^{20}$, but the point in the room is $(y_1,\ldots,y_{20})$, and the ``wall'' is the set of all points of the form $(f(x_1),\ldots,f(x_{20}))$ where $f$ of the form $f(x)=\beta_0+\beta_1\sin(0.7x)$. In this chapter, we'll learn how to find a perpendicular line between them.

This concept of finding a \emph{model} that gets as close as possible to \emph{data} is fundamentally what statistics is about, making it useful for all kinds of science and engineering where data is involved. It's also the simplest kind of machine learning model.

\section{Dot Products and Orthogonality}
\label{sec:dot-products-and-orthogonality}

\exercise (\emph{Warm-up}) Let $\vec{x}=(1,0,-3) = \bbbvec{1}{0}{-3}$ and $\vec{y}=(4,7,0) = \bbbvec{4}{7}{0}$. Multiply the matrices:
\begin{enumerate}[label=(\alph*)]
\item $\vec{x}^T\vec{y}$ \quad (Note: a $1x1$ matrix is basically just a number, so you don't need to write the brackets.)
\item $\vec{x}\vec{y}^T$
\end{enumerate}

\begin{notes}
\vfill
\end{notes}
\exercise \label{exercise:5.1-wu2} (\emph{Warm-up}) Let $\vec{x}=(x_1,x_2)$ be any vector in $\R^2$.
\begin{enumerate}[label=(\alph*)]
\item Multiply $\vec{x}^T\vec{x}$.
\item (T/F?) $\vec{x}^T\vec{x} \ge 0$.
\item (T/F?) $\sqrt{\vec{x}^T\vec{x}}$ is the distance from $(0,0)$ to $(x_1,x_2)$.
\item If $\vec{x}^T\vec{x}=0$, what can you say about $\vec{x}$?
\end{enumerate}

\begin{notes}
\vfill
\end{notes}

\bbtwocol{.55}{.45}{
\exercise (\emph{Warm-up}) The graph shows $\vec{u}=(1,3)$ and $\vec{v}=(5,1)$.
\begin{enumerate}[label=(\alph*)]
\item Find $\vec{w}$ so that $\vec{u}+\vec{w}=\vec{v}$.
\item (T/F?) The distance from $(1,3)$ to $(5,1)$ is the length of the arrow $\vec{w}$.
\item (T/F?) $\vec{u}+(\vec{v}-\vec{u}) = \vec{v}$.
\item (T/F?) $\vec{v}-\vec{u} = \vec{w}$.
\end{enumerate}
}{
\quad \begin{tikzpicture}
\begin{axis}[
x=22pt,y=22pt,
axis lines=middle,
axis line style={-},
ymajorgrids=true,
xmajorgrids=true,
xmin=-.5, xmax=5.5,
ymin=-.5, ymax=3.5,
xtick={-4,-3,...,5},
ytick={-1,0,...,4},
]
\draw[-stealth, thick] (axis cs:0,0) -- (axis cs:1,3)
node[left] {$\vec{u}$};
\draw[-stealth, thick] (axis cs:1,3) -- (axis cs:5,1);
\node at (axis cs:3,2.5) {$\vec{w}$};
\draw[-stealth, thick] (axis cs:0,0) -- (axis cs:5,1)
node[below] {$\vec{v}$};
\end{axis}
\end{tikzpicture}
}

\begin{notes}
\vfill \hspace{1ex}
\columnbreak
\end{notes}

\subsection{Dot Products and Orthogonality: Basics}\label{subsec:dot-products-and-orthogonality-basics}

\begin{definition}
If $\vec{u}$ and $\vec{v}$ are vectors in $\R^c$, then their \textbf{dot product} is:
\[\vec{u}\cdot \vec{v} = \vec{u}^T\vec{v} = u_1v_1+u_2v_2+\cdots+u_cv_c\]
\end{definition}
You may have studied dot products in another class, like multivariable calculus or physics. Everything you learned there still applies here. But since we can write it as a product of matrices, we can use everything we've learned about matrix multiplication study dot products, and then use what we learn about dot products to understand some types of matrix multiplication.

\begin{theorem}[Properties of the dot product.]
\label{theorem:dot-product-properties}
If $\vec{u},\vec{v}$, and $\vec{w}$ are vectors in $\R^c$ and $k$ is a scalar, then:

\begin{enumerate}[label=(\roman*)]
\item $\vec{u}\cdot \vec{v} = \vec{v}\cdot \vec{u}$
\item $( k\vec{u} )\cdot \vec{v}=k ( \vec{u}\cdot \vec{v} )$
\item $\vec{u}\cdot( \vec{v} + \vec{w} ) = \vec{u}\cdot \vec{v} + \vec{u}\cdot\vec{w}$
\item $\vec{u}\cdot \vec{u} \ge 0$
\end{enumerate}
\end{theorem}
The proofs are left as exercises. 

\begin{definition}
The \textbf{length} (aka \emph{norm} or \emph{magnitude}) of a vector is:
\[\nm{\vec{u}} = \sqrt{\vec{u}\cdot \vec{u}} = \sqrt{u_1^2+\cdots+u_c^2}\]
A \textbf{unit vector} has length 1.
\end{definition}
Like you saw in warm-up \ref{exercise:5.1-wu2}, in $\R^2$ this definition agrees with pythagorean theorem/distance formula for the length of $\vec{u}$ as an arrow. It's not too hard to show it agrees in $\R^3$ as well; see exercise \ref{exercise:R3-length}. 

\begin{theorem}[Properties of length.] \label{theorem:properties-of-length}
If $k$ is a scalar and $\vec{u}$ is a vector, then:
\begin{enumerate}[label=(\roman*)]
\item $\nm{\vec{u}}^2=\vec{u}\cdot \vec{u}$
\item $\nm{\vec{u}} = 0$ if and only if $\vec{u}=0$
\item $\nm{k\vec{u}} = |k|\nm{\vec{u}}$
\end{enumerate}
\end{theorem}
Part (i) obviously follows from the definition, but is worth writing down here because we'll use it so much. The other proofs are straightforward exercises.

\begin{notes}
\columnbreak
\end{notes}

\example Find two unit vectors parallel to $\vec{x}=(1,-2,2)$

\begin{book}
\textbf{Solution.} First, the result of multiplying any vector by 1 over its length is always a unit vector:
\[\nm{\frac{1}{\nm{\vec{x}}}\vec{x}} = \left|\frac{1}{\nm{\vec{x}}}\right| \nm{x} = 1\]
Applying this to $\vec{x}=(1,-2,2)$ gives:
\[\frac{1}{\nm{\vec{x}}}\vec{x} = \frac{1}{\sqrt{1^2+(-2)^2+2^{2}}}\bbbvec{1}{-2}{2} = \bbbvec{1/3}{-2/3}{2/3} \]
So we have one unit vector parallel to $\vec{x}$. The other is in the opposite direction: $(-1/3,2/3,-2/3)$. \blacksquare
\end{book}

\begin{notes}
\vfill 
\vfill 
\end{notes}

\begin{definition}

\bbtwocol{0.5}{0.4}{
The \textbf{distance from} $\vec{u}$ \textbf{to} $\vec{v}$ is $||\vec{v}-\vec{u}||$.}
{\vspace{-2em}

\hspace{-2em}
\begin{tikzpicture}[scale=0.8]
\draw[-Stealth, very thick] (0,0) -- (1,3);
\node at (0.5,2.4) {$\vec{u}$};
\draw[-Stealth, very thick] (1,3) -- (3,2);
\node at (2.2,2.9) {$\vec{v}-\vec{u}$};
\draw[-Stealth, very thick] (0,0) -- (3,2);
\node at (3.8,1.3) {$\vec{v}=\vec{u}+(\vec{v}-\vec{u})$};
\node[left] at (0,0) {$\vec{0}$};
\end{tikzpicture}}
\end{definition}

\example Find the distance from $\vec{u}=(1,0,0,0,0)$ to $\vec{v}=(0,0,0,0,-1)$.

\begin{book}
\textbf{Solution.} The vector between them is $\vec{u}-\vec{v}=(1,0,0,0,1)$, and its length is $\sqrt{1^2+0^2+0^2+0^2+(-1)^2} = \sqrt{2}$. \blacksquare
\end{book}
\begin{notes}
\vfill \hspace{1ex}
\columnbreak
\end{notes}

\begin{definition}
Two vectors $\vec{u}$ and $\vec{v}$ are called \textbf{orthogonal} if $\vec{u}\cdot \vec{v} = 0$.
\end{definition}

The English word \emph{orthogonal} basically means \emph{perpendicular}. Let's look at an example to get an idea what perpendicularity has to do with $\vec{u}\cdot \vec{v}=0$ and then we'll prove it.

\example Draw each pair of vectors. Are they geometrically perpendicular? Are they orthogonal?
\begin{enumerate}[label=(\alph*)]
\item $\vec{u}=(3,2,0), \vec{v}=(0,0,1)$
\item $\vec{x}=(1,0), \vec{y}=(2,1)$
\item $\vec{w}=(1,3), \vec{z}=(-3,1)$
\end{enumerate}

\begin{book}
\textbf{Solution.} (a) Since $\vec{u}$ is on the horizontal $x_1x_2$-plane and $\vec{v}$ is on the vertical $x_3$-axis, they are perpendicular. Since $\vec{u}\cdot \vec{v}=3(0)+2(0)+0(1) = 0$, they are orthogonal.

\begin{tikzpicture}[scale=0.8, x={(-0.8cm,-0.3cm)}, y={(0.9cm,-0.1cm)}, z={(0cm,1cm)}]
% Draw axes
\draw (0,0,0) -- (3.5,0,0) node[left]{$x_1$};
\draw (0,0,0) -- (0,2.5,0) node[anchor=north west]{$x_2$};
\draw (0,0,0) -- (0,0,2) node[right]{$x_3$};
% Draw tick marks on x-axis
\foreach \x in {1,2,3} {\draw (\x,0,0) -- (\x,0.1,0) node[above left] {\small $\x$};}
% Draw tick marks on y-axis
\foreach \y in {1,2} {\draw (0,\y,0) -- (0.1,\y,0) node[below left] {\small $\y$};}
% Draw tick marks on z-axis
\foreach \z in {1} {\draw (0,0,\z) -- (0.1,0,\z) node[left] {\small $\z$};}
\draw[-Stealth,very thick] (0,0,0) -- (3,2,0) node[right] {$\vec{u}$};
\draw[-Stealth,very thick] (0,0,0) -- (0,0,1) node[right] {$\vec{v}$};
\end{tikzpicture}

(b) They are not perpendicular; the angle between them is a little less than $30^{\circ}$. And $\vec{x}\cdot \vec{y}=1(2)+0(1)=2$ so they are not orthogonal.

\begin{tikzpicture}
\begin{axis}[
x=22pt,y=22pt,
axis lines=middle,
axis line style={-},
%ymajorgrids=true,
%xmajorgrids=true,
xmin=-.5, xmax=2.5,
ymin=-.9, ymax=1.5,
xtick={-4,-3,...,5},
ytick={-1,0,...,4},
tick style={thick},
]
\draw[-Stealth, very thick] (axis cs:0,0) -- (axis cs:1,0)
node[below left] {$\vec{x}$};
\draw[-Stealth, very thick] (axis cs:0,0) -- (axis cs:2,1)
node[right] {$\vec{y}$};
\end{axis}
\end{tikzpicture}

(c) They are indeed perpendicular; one way to tell is that their slopes, $3$ and $\frac{-1}{3}$ are opposite reciprocals. And they are orthogonal since $\vec{w}\cdot \vec{z} = 1(-3)+3(1) =0$. 

\begin{tikzpicture}
\begin{axis}[
x=18pt,y=18pt,
axis lines=middle,
axis line style={-},
%ymajorgrids=true,
%xmajorgrids=true,
xmin=-3.5, xmax=1.5,
ymin=-.5, ymax=3.5,
xtick={-4,-3,...,5},
ytick={-1,0,...,4},
tick style={thick},
]
\draw[-Stealth, very thick] (axis cs:0,0) -- (axis cs:1,3)
node[left] {$\vec{w}$};
\draw[-Stealth, very thick] (axis cs:0,0) -- (axis cs:-3,1)
node[above] {$\vec{z}$};
\end{axis}
\end{tikzpicture}

\blacksquare
\end{book}
\begin{notes}
\vfill
\end{notes}

\example Suppose $\vec{u}$ and $\vec{x}$ are any vectors in $\R^c$. Show that $\vec{u}$ and $\vec{x}- \frac{\vec{x}\cdot \vec{u}}{\vec{u}\cdot \vec{u}}\vec{u}$ are orthogonal.

\begin{book}
\textbf{Solution.} Keep in mind that result of a dot product is a scalar, so the fraction $\frac{\vec{x}\cdot \vec{u}}{\vec{u}\cdot \vec{u}}$ is a scalar too.
\begin{align*}
\vec{u}\cdot \left( \vec{x}-  \frac{\vec{x}\cdot \vec{u}}{\vec{u}\cdot \vec{u}} \vec{u}\right) &=
\vec{u}\cdot \vec{x} - \vec{u} \cdot\left( \frac{\vec{x}\cdot \vec{u}}{\vec{u}\cdot \vec{u}}\vec{u}\right) & \text{(Thm \ref{theorem:dot-product-properties} (iii))}\\ 
&= \vec{u}\cdot \vec{x} - \left( \frac{\vec{x}\cdot \vec{u}}{\vec{u}\cdot \vec{u}}\vec{u}\right) \cdot \vec{u}  & \text{(Thm \ref{theorem:dot-product-properties} (i))}\\ 
&= \vec{u}\cdot \vec{x} -  \frac{\vec{x}\cdot \vec{u}}{\vec{u}\cdot \vec{u}} \left(\vec{u}\cdot\vec{u}\right) & \text{(Thm \ref{theorem:dot-product-properties} (ii))}\\ 
&= \vec{u}\cdot \vec{x} - \vec{x}\cdot \vec{u} \\
&= 0 & \text{(Thm \ref{theorem:dot-product-properties} (i))}
\end{align*}
\blacksquare
\end{book}
\begin{notes}
\vfill \hspace{1ex}
\columnbreak
\end{notes}

\begin{theorem}[Orthogonal means perpendicular.] \label{theorem:orthogonal-perpendicular}

Suppose $\vec{u}$ and $\vec{v}$ are non-zero vectors in $\R^c$. The angle between them is $90^{\circ}$ if and only if $\vec{u}\cdot \vec{v}=0$.
\end{theorem}

Note that if $\vec{u}=\vec{0}$ or $\vec{v} = \vec{0}$, then $\vec{u}\cdot \vec{v}=0$ but the ``angle between them'' doesn't really make sense.

\bbtwocol{0.58}{0.42}{
\textbf{Proof.} Exercise \ref{exercise:prove-orthogonal} asks you to fill in the gaps below.  The graph on the right shows $\vec{u}$ and $\vec{v}$ starting at the origin, labeled $A$ here, as well as copies of them to make a parallelogram $\pgram ABCD$. Note that this doesn't just apply to $\R^2$, because even in higher dimensions, $\vec{u},\vec{v},$ and $\vec{u}+\vec{v}$ all lie on the same plane.}
{\begin{tikzpicture}[scale=0.85,line cap = round]
\draw[-stealth, thick] (0,0) -- (1,3) node[midway,left] {$\vec{v}$};
\draw[-stealth, thick] (4,0) -- (5,3) node[midway,right] {$\vec{v}$};
\draw[-stealth, thick] (0,0) -- (4,0) node[midway,below] {$\vec{u}$};
\draw[-stealth, thick] (1,3) -- (5,3) node[midway,above] {$\vec{u}$};
\draw[-stealth, thick] (0,0) -- (5,3);
\node[above,rotate=31] at (3.5,2) {$\vec{u}+\vec{v}$};
\draw[-stealth, thick] (1,3) -- (4,0);
\node[below,rotate=-45] at (3.1,1) {$\vec{u}-\vec{v}$};
\node[below left] at (0,0) {$A$};
\node[above left] at (1,3) {$C$};
\node[below right] at (4,0) {$B$};
\node[above right] at (5,3) {$D$};
\coordinate (A) at (0,0);
\coordinate (B) at (4,0);
\coordinate (C) at (1,3);
\coordinate (D) at (5,3);
\tkzMarkAngle[size=0.6](B,A,C)
\node at (0.45,0.65) {$\alpha$};
\tkzMarkAngle[size=0.6](D,B,A)
\node at (3.7,0.75) {$\beta$};
%\pic[] {angle = (0,1)--(0,0)--(1,0)}
\end{tikzpicture}}

First, we'll show that:
\begin{equation}
\label{eq:5.1-parallelogram}
\nm{\vec{u}+\vec{v}}^2-\nm{\vec{u}-\vec{v}}^{2} = 4 \vec{u}\cdot \vec{v}
\end{equation}
The first step is done; fill in the missing steps:
\begin{align*}
\nm{\vec{u}+\vec{v}}^2-\nm{\vec{u}-\vec{v}}^{2} &= ( \vec{u}+\vec{v} )\cdot ( \vec{u}+\vec{v} ) - ( \vec{u}-\vec{v} )\cdot ( \vec{u}-\vec{v} ) \\
&= \ldots \\
&\vdots  \\
&= 4 \vec{u} \cdot \vec{v}
\end{align*}
Equation \ref{eq:5.1-parallelogram} shows that $\nm{\vec{u}+\vec{v}} = \nm{\vec{u}-\vec{v}}$ if and only if $\rule{4em}{0.8pt}=0$. Now we can prove the theorem:

(\rightarrow) Suppose $\alpha=90^{\circ}$. Since $\alpha$ and $\beta$ are in the same side of a parallelogram, $\alpha+\beta=180^{\circ}$ and so $\beta=\rule{2em}{0.8pt}$. This means the triangle $\bigtriangleup CAB$ is congruent to (same angles and side lengths) \rule{4em}{0.8pt}, since they have two equal sides ($\nm{\vec{u}} = \nm{\vec{u}}$ and $\nm{\vec{v}} = \nm{\vec{v}}$) and equal angles between those sides ($\alpha=\beta$). This means their hypotenuses are equal: $\nm{\vec{u}+\vec{v}} = \nm{\vec{u}-\vec{v}}$, and equation \ref{eq:5.1-parallelogram} implies $\vec{u}\cdot \vec{v}= \rule{2em}{0.8pt}$.

(\leftarrow) Now suppose $\vec{u}\cdot \vec{v}=0$. Equation \ref{eq:5.1-parallelogram} implies that $\nm{\vec{u}+\vec{v}} = \rule{5em}{0.8pt}$. The triangle $\bigtriangleup CAB$ is congruent to \rule{4em}{0.8pt}, since all three pairs of corresponding sides are equal ($\nm{\vec{u}} = \nm{\vec{u}}$, $\nm{\vec{v}} = \nm{\vec{v}}$, and the hypotenuses, as we just showed). Thus $\alpha=\beta$ and since they're on the same side of a parallelogram, $\alpha=\beta=90^{\circ}$. \blacksquare

\begin{theorem}[The Pythagorean theorem.]
\label{theorem:pythagorean}

\bbtwocol{0.58}{0.4}{
Two vectors $\vec{u}$ and $\vec{v}$ in $\R^c$ are orthogonal if and only if $\nm{\vec{u}}^2+\nm{\vec{v}}^2 = \nm{\vec{u}-\vec{v}}^{2}$
}
{\vspace{-1em}

\hspace{1.5em} \begin{tikzpicture}[scale=0.85,line cap = round]
\draw[-stealth, thick] (0,0) -- (0,3) node[midway,left] {$\vec{v}$};
\draw[-stealth, thick] (0,0) -- (4,0) node[midway,below] {$\vec{u}$};
\draw[-stealth, thick] (0,3) -- (4,0);
\node[above,rotate=-37] at (2,1.4) {$\vec{u}-\vec{v}$};
\coordinate (A) at (0,0);
\coordinate (B) at (4,0);
\coordinate (C) at (0,3);
\tkzMarkRightAngle[size=0.4](B,A,C)
\end{tikzpicture}}
\end{theorem}

Note a simple proof would be to apply theorem \ref{theorem:orthogonal-perpendicular} and use the Pythagorean theorem as you learned it years ago. But we're here to do Linear Algebra so that's what we're gonna do. The gaps in the proof are for you to fill in for exercise \ref{exercise:prove-pythagorean}.

\textbf{Proof.} Suppose $\vec{u}$ and $\vec{v}$ are any vectors in $\R^c$. Fill in the missing steps:
\begin{align*}
\nm{\vec{u}-\vec{v}}^2 &= (\vec{u}-\vec{v})\cdot(\vec{u}-\vec{v}) \\
& = \ldots \\
&\vdots  \\
\nm{\vec{u}-\vec{v}}^2 &= \nm{\vec{u}}^2 - 2\vec{u}\cdot \vec{v} + \nm{\vec{v}}^{2}
\end{align*}
Thus $\nm{\vec{u}-\vec{v}}^2 = \nm{\vec{u}}^2 + \nm{\vec{v}}^{2}$ if and only if $2\vec{u}\cdot \vec{v} = \rule{2em}{0.8pt}$ which is true if and only if $\vec{u}$ and $\vec{v}$ are \rule{8em}{0.8pt}. \blacksquare

\subsection{Orthogonal Complements}\label{subsec:orthogonal-complements}

\begin{definition}
Suppose $W$ is a subspace of $\R^c$. The \textbf{orthogonal complement} of $W$, written $W^{\perp}$, is the set of all vectors that are perpendicular to every vector in $W$. In symbols:
\[W^{\perp} = \left\{ \vec{u} \text{ in }\R^c : \vec{u}\cdot \vec{w} = 0 \text{ for all } \vec{w} \text{ in } W\right\}\]
\end{definition}

\example \label{example:5.1-four-subspaces-1} Let $\vec{a}=\bbvec{-2}{1}$. Note that $\vec{a}$ is a matrix with one column so we can write $\col \vec{a}=\spn(\vec{a})$. Describe $(\col \vec{a})^{\perp}$ and make sure your answer agrees with the graph below.

\begin{book}
\textbf{Solution.} $\vec{x}$ is in $(\col \vec{a})^{\perp}$ if and only if:
\begin{align*}
\vec{x}\cdot \vec{w} = 0 \text{ for all } \vec{w} \text{ in } (\col \vec{a}) &\text{ \quad iff \quad } \vec{x}\cdot k\vec{a} = 0 \text{ for all } k \text{ in } \R \\
&\text{ \quad iff \quad } x_1k(-2) + x_2k(1) = 0 \\
&\text{ \quad iff \quad } -2x_1 + x_2 = 0 \\
&\text{ \quad iff \quad } x_1 = \frac{1}{2}x_2 \\
&\text{ \quad iff \quad } \vec{x}=\bbvec{x_1}{x_2}= x_2 \bbvec{1/2}{1}
\end{align*}
Thus $(\col\vec{a})^{\perp}=\spn \left( \bbvec{0.5}{1} \right)$. Note that the $k$ canceled out, indicating that, as long as $\vec{x}$ is orthogonal to $(-2,1)$, it's orthogonal to any scalar multiple of $(-2,1)$. Also note that we're solving a homogeneous equation; the third equation from above means $\vec{x}$ is in $\nul(\vec{a}^T)$. The process of solving for basic variables in terms of free variables, and using it to write the solution set as a span is the same as it was in chapter 1. \blacksquare
\end{book}

\begin{tikzpicture}[line cap = round]
\begin{axis}[
x=18pt,y=18pt,
axis lines=middle,
axis line style={-},
%ymajorgrids=true,
%xmajorgrids=true,
xmin=-6.5, xmax=6.5,
ymin=-4.5, ymax=4.5,
xtick={-6,-5,...,6},
ytick={-4,-3,...,4},
tick style={thick}
]
\draw[thick, blue] (axis cs:-6,3) -- (axis cs:6,-3);
\node[above, blue] at (axis cs:-4,2.2) {$\col\vec{a}$};
\draw[-Stealth, very thick] (axis cs:0,0) -- (axis cs:-2,1)
node[above right] {$\vec{a}$};
\draw[thick, red] (axis cs:-3,-6) -- (axis cs:3,6);
\node[right, red] at (axis cs:1.25,2.5) {$(\col\vec{a})^{\perp} = \nul(\vec{a}^T)$};
\draw[-Stealth, very thick] (axis cs:0,0) -- (axis cs:0.5,1)
node[right] {$(0.5,1)$};
\draw[thin] (axis cs:.22,.44) -- (axis cs:-.22,.67) -- (axis cs:-.45,.22);
\end{axis}
\end{tikzpicture}

\example \label{example:5.1-four-subspaces-2} Let $\vec{b}=\bbbvec{1}{0}{3}$ and note $\col \vec{b}=\spn (\vec{b})$. Describe $(\col\vec{b})^{\perp}$. Hopefully the picture agrees with the answer.

\begin{book}
\textbf{Solution.} A vector $\vec{x}$ is in $\col(\vec{b})^{\perp}$ if and only iff:
\begin{align*}
\vec{x}\cdot \vec{w} = 0 \text{ for all } \vec{w} \text{ in } \col(\vec{b}) &\text{ \quad iff \quad } \vec{x}\cdot k\vec{b} = 0 \text{ for all } k \text{ in } \R \\
&\text{ \quad iff \quad } x_1k(1) + x_2k(0) + x_3k(3) = 0 \\
&\text{ \quad iff \quad } x_1+0x_2+3x_3 = 0 \\
&\text{ \quad iff \quad } \vec{x}=\bbbvec{x_1}{x_2}{x_3} = \bbbvec{-3x_3}{x_2}{x_3} = x_2 \bbbvec{0}{1}{0} + x_3 \bbbvec{-3}{0}{1}
\end{align*}
So $\col(\vec{b})^{\perp} = \spn ( \vec{u},\vec{v} )$ where $\vec{u}=(-3,0,1), \vec{v}=(0,1,0)$. Again, note that the second-to-last equation above says $\col(\vec{b})^{\perp} = \nul(A^T)$.\blacksquare
\end{book}
\begin{notes}
\vfill
\end{notes}

\begin{tikzpicture}[line cap = round, scale=0.8, x={(-0.7cm,-0.7cm)}, y={(1cm,0cm)}, z={(0cm,1cm)}]
\draw (-3.5,0,0) -- (3.5,0,0) node[left]{$x_1$};
\draw (0,0,0) -- (0,2.5,0) node[anchor=north west]{$x_2$};
\draw (0,0,0) -- (0,0,3.5) node[right]{$x_3$};
% Draw tick marks on x-axis
\foreach \x in {-3,-2,-1,1,2,3} {\draw (\x,0,0) -- (\x,0.1,0) node[right] {\scriptsize $\x$};}
% Draw tick marks on y-axis
\foreach \y in {1,2} {\draw (0,\y,0) -- (0.1,\y,0) node[below left] {\scriptsize $\y$};}
% Draw tick marks on z-axis
\foreach \z in {1,2,3} {\draw (0,0,\z) -- (0,0.1,\z) node[right] {\scriptsize $\z$};}
\draw[blue,thick] (1.5,0,4.5) -- (-1.3,0,-3.9);
\node[right,blue] at (-1,0,-3) {$\col\vec{b}$};
\fill[red!30,opacity=0.5] (3,-2,-1) -- (3,2,-1) -- (-3,2,1) -- (-3,-2,1) -- cycle;
\node[red,below left] at (-3,2.2,1) {$(\col\vec{b})^{\perp}$};
\draw[-Stealth,very thick] (0,0,0) -- (1,0,3) node[left] {$\vec{b}$};
\draw[-Stealth,very thick,red] (0,0,0) -- (-3,0,1) node[below right] {$\vec{u}$};
\draw[-Stealth,very thick,red] (0,0,0) -- (0,1,0) node[above] {$\vec{v}$};
\coordinate (A) at (0,0,0);
\coordinate (B) at (1,0,3);
\coordinate (C) at (-3,0,1);
\coordinate (D) at (0,1,0);
\tkzMarkRightAngle[size=0.3](B,A,C)
\tkzMarkRightAngle[size=0.3](B,A,D)
\end{tikzpicture}

\begin{notes}
\columnbreak
\end{notes}

The next theorem should be intuitive, if you're able to decipher the graphs in the previous examples and the next. The proof is left as an exercise.

\begin{theorem} \label{theorem:5.1-orthogonal-span}
A vector $\vec{x}$ is in $\spn ( \vec{a}_1,\ldots,\vec{a}_c ) ^{\perp}$ if and only if $\vec{x}$ is orthogonal to each vector $\vec{a}_1,\ldots,\vec{a}_c$.
\end{theorem}


\example \label{example:5.1-four-subspaces-3} Let $A = \begin{bmatrix}
\vec{a}_1 & \vec{a}_2
\end{bmatrix} = \begin{bmatrix}
1 & 1 \\
1 & 0 \\
0 & 1
\end{bmatrix}$. Describe $(\col A)^{\perp}$.

\begin{book}
\textbf{Solution.} Using theorem \ref{theorem:5.1-orthogonal-span}, $\vec{x}$ is in $(\col A)^{\perp}$ if and only if:
\begin{align*}
&\vec{x}\cdot \vec{a}_1=0 \text{ and } \vec{x}\cdot \vec{a}_2=0 \\
\text{ \quad iff \quad } &\vec{x} \text{ is a soln of the system }
\begin{matrix}
\vec{x}\cdot \vec{a}_1=0 \\
\vec{x}\cdot \vec{a}_2=0
\end{matrix}\\
\text{ \quad iff \quad } &\begin{matrix}
1x_1+1x_2+0x_3=0 \\
1x_1+0x_2+1x_3=0
\end{matrix}\\
\text{ \quad iff \quad } &\begin{bmatrix}
1 & 1 & 0 \\
1 & 0 & 1
\end{bmatrix} \vec{x} = \vec{0} \\
\text{ \quad iff \quad } &\begin{bmatrix}
1 & 0 & 1 \\
0 & 1 & -1
\end{bmatrix} \vec{x} = \vec{0} \\
\text{ \quad iff \quad } &\vec{x} = \bbbvec{x_1}{x_2}{x_3} = \bbbvec{-x_3}{x_3}{x_3} = x_3 \bbbvec{-1}{1}{1}
\end{align*}
So, $(\col A)^{\perp}=\spn ( \vec{u} )$ with $\vec{u}=(-1,1,1)$. Note that the second-to-last step above shows $\col(A)^{\perp} = \nul (A^T)$. \blacksquare
\end{book}
\begin{notes}
\vfill
\end{notes}

\begin{tikzpicture}[line cap = round, scale=1.5, x={(-0.7cm,-0.7cm)}, y={(.7cm,-.7cm)}, z={(0cm,1cm)}]
\draw (0,0,0) -- (2.5,0,0) node[left]{$x_1$};
\draw (0,0,0) -- (0,2.5,0) node[anchor=north west]{$x_2$};
\draw (0,0,0) -- (0,0,2.5) node[right]{$x_3$};
% Draw tick marks on x-axis
\foreach \x in {1,2} {\draw (\x,0,0) -- (\x,0.1,0) node[right] {\scriptsize $\x$};}
% Draw tick marks on y-axis
\foreach \y in {1,2} {\draw (0,\y,0) -- (0.1,\y,0) node[below left] {\scriptsize $\y$};}
% Draw tick marks on z-axis
\foreach \z in {1,2} {\draw (0,0,\z) -- (0,0.1,\z) node[right] {\scriptsize $\z$};}
\node[above right,blue] at (3,1.5,1.3) {$\col A$};
\fill[blue!30,opacity=0.5] (-3,-1.5,-1.5) -- (0,-1.5,1.5) -- (3,1.5,1.5) -- (0,1.5,-1.5) -- cycle;
\draw[-Stealth,very thick,blue] (0,0,0) -- (1,1,0) node[right] {$\vec{a}_1$};
\draw[-Stealth,very thick,blue] (0,0,0) -- (1,0,1) node[below] {$\vec{a}_2$};
\draw[thick, red] (1,-1,-1)--(-1.5,1.5,1.5) node[right] {$\begin{matrix}(\col A)^{\perp} \\ = \nul A^{T}\end{matrix}$};
\draw[-Stealth,very thick] (0,0,0) -- (-1,1,1) node[below] {$\vec{u}$};
\coordinate (A) at (1,1,0);
\coordinate (B) at (1,0,1);
\coordinate (C) at (0,0,0);
\coordinate (D) at (-1,1,1);
\tkzMarkRightAngle[size=0.3](A,C,D)
\tkzMarkRightAngle[size=0.3](B,C,D)
\end{tikzpicture}
\begin{notes}
\columnbreak
\end{notes}

These examples illustrate the next theorem, which is sometimes called \textsc{the fundamental theorem of linear algebra}.

\begin{theorem}[Four related subspaces.] \label{theorem:four-subspaces}
If $A$ is an $r \times c$ matrix, then:
\begin{enumerate}[label=(\roman*)]
\item $( \col A )^{\perp}=\nul  A^T$ and
\item  $(\row A )^{\perp}=\nul A$.
\end{enumerate}
\end{theorem}

\begin{notes}
\vfill
\end{notes}
\begin{book}
\textbf{Proof.} (i) Denote the columns of $A$ as $\vec{a}_1,\ldots,\vec{a}_c$, so that:
\[A = \begin{bmatrix}
\vec{a}_1 & \cdots & \vec{a}_c 
\end{bmatrix} \quad \text{and} \quad A^T= \bbbvec{\vec{a}_1^T}{\vdots}{\vec{a}_c^T}\]
The following statements are equivalent:
\begin{align*}
&\vec{x} \text{ is in } \nul (A^T)\\
\text{iff} \quad &A^T\vec{x}=\vec{0} \\
\text{iff} \quad &\bbbvec{\vec{a}_1^T}{\vdots}{\vec{a}_c^T}\vec{x} = \vec{0} \\
\text{iff} \quad &\bbbvec{\vec{a}_1\cdot \vec{x}}{\vdots}{\vec{a}_c\cdot \vec{x}} = \bbbvec{0}{0}{0} \text{ (row-column rule)} \\
\text{iff} \quad & \vec{x} \text{ is in } \spn(\vec{a}_1, \ldots, \vec{a}_c)^{\perp} = \col(A)^{\perp} \text{ (theorem \ref{theorem:5.1-orthogonal-span})}
\end{align*}
By theorem \ref{theorem:5.1-orthogonal-span} the last statement is equivalent to 

(ii) The rows of $A$ are the columns of $A^T$, so:
\begin{align*}
\row(A) &= \col(A^T) \\
&= \nul ( (A^T)^{T} ) \quad \text{(part (i))} \\
&= \nul (A)
\end{align*}
\blacksquare
\end{book}

\begin{theorem}[Properties of orthogonal complements.] \label{theorem:properties-orthogonal-complements}

If $W$ is a subspace of $\R^n$, then:
\begin{enumerate}[label=(\roman*)]
\item $W^{\perp}$ is also a subspace of $\R^n$.
\item $\dim W +\dim W^{\perp} = n$. 

\item $(W^{\perp})^{\perp}=W$. 
\end{enumerate}
\end{theorem}

\begin{notes}
The proofs of (i) and (ii) are left as exercises. The proof of (iii) is ironically the most difficult of them, and is in the textbook.

\columnbreak
\end{notes}

\begin{book}
\textbf{Proof.} (i) Exercise \ref{exercise:prove-orthogonal-properties-i}.% $W$ has a basis by thm \ref{theorem:subspaces-are-spans}; put those vectors in a matrix $A$ so that $W=\col A$. Then by thm \ref{theorem:four-subspaces}, $W^{\perp}=\nul(A^T)$ which is a subspace by theorem \ref{theorem:null-spaces-are-subspaces}. ACTUALLY IT'S BETTER TO JUST SHOW IT SATISFIES THE DEFN OF A SUBSPACE!

(ii) Exercise \ref{exercise:prove-orthogonal-properties-ii}. %Let $\dim W = c$ and put the $c$ basis vectors in a $c\times n$ matrix $A$. 
%\begin{align*}
%\dim W + \dim W^{\perp} &= \dim (\col A) + \dim (\col A)^{\perp} \\
%&= \dim(\row A) + \dim (\nul (A^{T})) & \text{ (remarks in 1.7 and thm \ref{theorem:four-subspaces})} \\
%&= \dim(\col A^{T}) + \dim (\nul (A^T))\\
%&= n &\text{rank-nullity thm (\ref{theorem:rank-nullity})}
%\end{align*}

(iii) If $\vec{w}$ is in $W$, then $\vec{w} \perp \vec{z}$ for any $\vec{z}$ in $W^{\perp}$, so $\vec{w}$ is in $(W^{\perp})^{\perp}$ and $W$ is a subset of $( W^{\perp} )^\perp$. Showing that $( W^{\perp} )^{\perp}$ is a subset of $W$ is harder. Since $W$ is a subspace, it has a basis with some number of vectors: let's say $\dim W = c$. By part (ii), $\dim(W^{\perp}) = n - c$, and $\dim((W^{\perp})^{\perp})=n-(n-c) = c$. So $W$ has the same dimension as $(W^{\perp})^{\perp}$ and the buy-two-get-one-free theorem implies that a basis of $W$ spans $(W^{\perp})^{\perp}$, since the basis is linearly independent and has the right number of vectors. Thus any vector in $(W^{\perp})^{\perp}$ is in $W$ too. \blacksquare
\end{book}

\example Suppose $\rank(A) = 1$ and $A \begin{bmatrix}
1 \\ 0 \\ -2 \\ 3 
\end{bmatrix} = \vec{0}$. Attempt to find $\dim(\nul A)$ and  $\dim(\nul A^T)$.

\begin{notes}
\vfill
\end{notes}
\begin{book}
\textbf{Solution.} Since $A$ can be multiplied by a vector with 4 entries, we know $A$ has 4 columns, and the rows of $A$ have 4 entries. Suppose $A$ has $r$ rows. Some things we know:
\begin{align*}
\dim(\col A) &= \dim (\row A) = 1 \quad \text{(def'n of rank)} \\
\dim(\row A) + \dim((\row A)^{\perp}) &= 4 \quad (\text{each is a subspace of }\R^4 )\\
(\row A)^{\perp} &= \nul A
\end{align*}
Putting these together, we get $\dim(\nul A) = \dim((\row A)^{\perp}) = 4 - 1 = 3$. To find $\dim(\nul A^T)$, these equations might help:
\begin{align*}
\dim(\col A) + \dim((\col A)^{\perp}) &= r \quad (\text{each is a subspace of }\R^r )\\
(\col A)^{\perp} &= \nul A^T
\end{align*}
Combined, this gives $\dim(\nul A^T) = \dim((\col A)^{\perp}) = r - \dim(\col A) = r-1$.
The problem is that we don't know how many rows $A$ has. In fact there are several possibilities, for example:
\[ A = \begin{bmatrix}
-1 & 0 & 1 & 1
\end{bmatrix}\quad \text{ or } \quad A = \begin{bmatrix}
-1 & 0 & 1 & 1 \\
-7 & 0 & 7 & 7 \\
0 & 0 & 0 & 0 \\
\end{bmatrix} \]
The one on the left makes $\dim((\col A)^{\perp})=1-1=0$ and the one on the right makes $\dim((\col A)^{\perp})=3-1=2$. \blacksquare
\end{book}

\example Find $W^{\perp}$ if $W = \left\{ (w_{1},w_2,w_3) : 7w_1 - w_2 = 2w_3 \right\}$.

\begin{notes}
\vfill \hspace{1ex}
\columnbreak
\end{notes}
\begin{book}
\textbf{Solution.} The equation which defines $W$ is equivalent to $7w_1-w_2-2w_3=0$, which makes $W$ equal to the null space of $\begin{bmatrix}
7 & -1 & -2
\end{bmatrix}$. Thus $W^{\perp}=\left(\nul \begin{bmatrix}
7 & -1 & -2
\end{bmatrix} \right)^{\perp}=\left(\left(\row \begin{bmatrix}
7 & -1 & -2
\end{bmatrix}\right)^{\perp}\right)^{\perp} = \row \begin{bmatrix}
7 & -1 & -2
\end{bmatrix} = \spn \bbbvec{7}{-1}{-2}$. \blacksquare
\end{book}

\columnbreak
\exerciseheading
%\begin{savenotes}

\textbf{\ref{subsec:dot-products-and-orthogonality-basics} Dot Products and Orthogonality Basics}

\exatend Find the length of each vector: 
\[\vec{x}=\bbbvec{7}{-7}{\sqrt{2}} \quad \vec{y}=\bbbvec{0}{-5}{0} \quad \vec{x}-\vec{y}\]

\exatend Which pairs of the following vectors are orthogonal to each other? Which are unit vectors?
\[\vec{a}=\bbbvec{0.6}{0.8}{0} \quad \vec{b}=\bbbvec{-0.8}{0.6}{0} \quad \vec{p}=\bbbvec{0}{0}{3} \quad \vec{q}=\bbbvec{1}{0}{0}\] 

\exatend If $\vec{u}=\bbvec{3}{5}$ and $\vec{v}=\bbvec{2}{-1}$, find $\vec{u}\cdot \vec{v}$, $\vec{v}\cdot \vec{v}$, and $\frac{\vec{u}\cdot \vec{v}}{\vec{v}\cdot \vec{v}}\vec{v}$.

\exatend Prove theorem \ref{theorem:dot-product-properties} part (i).
\exatend Prove theorem \ref{theorem:dot-product-properties} part (ii).
\exatend Prove theorem \ref{theorem:dot-product-properties} part (iii).
\exatend Prove theorem \ref{theorem:dot-product-properties} part (iv).

\exatend Prove theorem \ref{theorem:properties-of-length} part (ii).
\exatend Prove theorem \ref{theorem:properties-of-length} part (iii).

\exatend Find two unit vectors in $\R^2$ which are each parallel to $\bbvec{3}{3}$.

\exatend Find two unit vectors in $\R^2$ which are each orthogonal to $\bbvec{3}{3}$.

\exatend A triangle has three sides $(2,2),(5,6),$ and $(11.6,-5.2)$. Find the lengths of each side. Is it a right triangle?

\exatend (T/F?) $(\vec{u}\cdot \vec{u})^2=\nm{\vec{u}}$.

\exatend (T/F?) For any scalar $k$, $\nm{k \vec{u}} = k \nm{\vec{u}}$.


\exatend (T/F?) The zero vector in $\R^c$ is orthogonal to every other vector in $\R^c$.

\exatend (T/F?) If $\vec{u}\cdot \vec{u} = 0$, then $\vec{u}=\vec{0}$.

%\exatend Give an example of vectors $\vec{u}$ and $\vec{v}$ such that $\nm{\vec{u}}\nm{v}=\vec{u}\cdot \vec{v}$.



\exatend \label{exercise:R3-length} Use the Pythagorean theorem to find the distance from $(0,0,0)$ to $(u_1,u_2,u_3)$, and check that it agrees with the definition of $||(u_1,u_2,u_3)||$. \hintfootnote{The points $(0,0,0),(u_1,u_2,0)$ and $(u_1,u_2,u_3)$ form a right triangle. You can use either the version of the Pythagorean theorem that you learned years ago, or the linear algebra version, thm \ref{theorem:pythagorean}.}

\exatend Let $\vec{a}=(a_1,a_2)$ and $\vec{b}=(b_1,b_2)$ be two non-zero $\R^2$ vectors. Note that the slopes of the vectors are $\frac{a_2}{a_1}$ and $\frac{b_2}{b_1}$. You probably learned in an earlier class that two lines in $\R^2$ are perpendicular when their slopes are opposite reciprocals: $m_2= \frac{-1}{m_1}$. Show this is equivalent to $\vec{a}\cdot \vec{b}=0$. \creditfootnote{Inspired by an exercise in Gilbert Strang's excellent \emph{Linear Algebra and its Applications}.}

\exatend Give an example of two vectors that are linearly independent but not orthogonal.

\exatend Give an example of two vectors that are orthogonal but not linearly independent.

\exatend Suppose $A$ is an invertible matrix. Explain how you know row $i$ of $A$ is orthogonal to column $j$ of $A^{-1}$, if $i \ne j$.

\exatend \label{exercise:prove-orthogonal} Fill in the gaps in the proof of theorem \ref{theorem:orthogonal-perpendicular}.

\exatend \label{exercise:prove-pythagorean} Fill in the gaps in the proof of theorem \ref{theorem:pythagorean}.

\exatend (T/F?) If $\vec{u}\cdot \vec{v}=0$, then $\nm{\vec{u}}^2 + \nm{\vec{v}}^2 = \nm{\vec{u}+\vec{v}}^2$. (Hint: draw a picture.)

\exatend Show that $\vec{u}-\vec{v}$ is orthogonal to $\vec{u}+\vec{v}$ if and only if $\nm{u}=\nm{v}$. \creditfootnote{Inspired by Strang.}

\textbf{\ref{subsec:orthogonal-complements} Orthogonal Complements}

% PUT ORTHOGONAL COMPLEMENT PROBLEMS BELOW HERE!
\exatend Show that if $\vec{x}$ is in $W$ and in $W^{\perp}$, then $\vec{x}=\vec{0}$.

\exatend Recall that $\left\{ \vec{0} \right\}$, the set containing only the zero vector, is always a subspace of $\R^c$. What is $\left\{ \vec{0} \right\}^{\perp}$?

\exatend Consider $\vec{a}=\bbvec{-2}{1}$, the $2\times1$ matrix from example \ref{example:5.1-four-subspaces-1}. Describe $\row(A)$ and $\row(\vec{a})^{\perp}$.

\exatend Consider $A = \begin{bmatrix}
1 & 1 \\
1 & 0 \\
0 & 1
\end{bmatrix}$, the matrix from example \ref{example:5.1-four-subspaces-3}. Describe $\row A$ and $(\row A)^{\perp}$.

\exatend Let $A = \begin{bmatrix}
10 & 5 \\
-2 & -1
\end{bmatrix}$. Find one vector in each of:
\begin{enumerate}[label=(\alph*)]
\item $(\col(A))^{\perp}$
\item $(\nul A)^{\perp}$
\item $(\row A)^{\perp}$
\end{enumerate}

\exatend Let $A = \begin{bmatrix}
1 & 2 & 3 \\
4 & 5 & 6 \\
3 & 3 & 3
\end{bmatrix}$. Note the second row minus the first row equals the third row. And the third column is the opposite of the first column plus twice the second column. Find one vector in each of:
\begin{enumerate}[label=(\alph*)]
\item $(\col(A))^{\perp}$
\item $(\nul(A))^{\perp}$
\item $(\row(A))^{\perp}$
\end{enumerate}

\exatend Suppose $A$ is a $10 \times 15$ matrix with $\rank(A) = 6$. Find $\dim(\nul(A))$ and $\dim(\nul(A^T))$.

\exatend (T/F?) If $A$ is an $r \times c$ matrix, then $\col(A)$ is a subspace of $\R^r$ and $\row(A)$ is a subspace of $\R^c$.

\exatend (T/F?) If $A\vec{x}=\vec{0}$, then $\vec{x}$ is orthogonal to every column of $A$.

\exatend (T/F?) The vectors in $\col(A^T)$ are orthogonal to the vectors in $\nul(A)$.


\exatend (T/F?) There exists a matrix $A$ such that $(1,2,3)$ is in $\nul(A)$, and $(0,2,0)$ is in $\row(A)$.

\exatend (T/F?) The solution set of $\begin{bmatrix}
1 & 2 & 3 \\
4 & 5 &6
\end{bmatrix}\vec{x}=\vec{0}$ is orthogonal to the span of $(1,2,3)$ and $(4,5,6)$.

\exatend Find a counterexample to the false statement: If $\vec{x}\cdot \vec{y}=0$, then $\spn(\vec{x})^{\perp} = \spn(\vec{y})$.

\exatend Let $A = \begin{bmatrix}
1 & 2 & 3 \\
-1 & -2 & -3
\end{bmatrix}$. Find a basis for each subspace: $\col(A)$, $\row(A)$, $\nul(A)$, and $\nul(A^T)$.

\exatend Let $W$ be the plane in $\R^3$ with equation $2x_1-x_2+4x_3=0$. Find a matrix $A$ such that $W= \nul(A)$, and a matrix $B$ such that $W=\row(B)$. \creditfootnote{Inspired by Strang.}

% \begin{book}
% \textbf{Solution.} The second and third columns of $A$ are multiples of the first, so the first column $\bbvec{1}{-1}$ by itself is a basis of $\col(A)$. We know $\nul(A^T)=\col(A)^{\perp}$, so $\nul(A^T)$ is a dimension $2-1=1$ subspace of $\R^2$, so any vector perpendicular to $\bbvec{1}{-1}$ will work; for example $\bbvec{10}{10}$ is a basis of $\nul(A^T)$.

% Since the second row of $A$ is a multiple of the first, the first row of $A$ is a basis of $\row(A)$: $\bbbvec{1}{2}{3}$. Since $\nul(A) = \row(A)^{\perp}$, we know $\nul(A)$ is a 2-dimensional subspace of $\R^3$, so two vectors that are perpendicular to $(1,2,3)$ will make a basis of $\nul(A)$, like $(-2,1,0), (-3,0,1)$. \blacksquare
% \end{book}

\exatend Prove theorem \ref{theorem:5.1-orthogonal-span}.

\exatend Prove part (ii) of theorem \ref{theorem:four-subspaces} directly (that is, without using part (i)).

\exatend \label{exercise:prove-orthogonal-properties-i} Prove theorem \ref{theorem:properties-orthogonal-complements} part (i).

\exatend \label{exercise:prove-orthogonal-properties-ii} (\emph{Challenge}) Prove theorem \ref{theorem:properties-orthogonal-complements} part (ii). \hintfootnote{ One way to do it is with \ref{theorem:basis-for-row-space}, \ref{theorem:basis-for-column-space}, \ref{theorem:rank-nullity}, and \ref{theorem:four-subspaces}.}

% TODO maybe someday: there are more good problems in Strang 3.1, specifically 23 - 27.
%\end{savenotes}

\columnbreak
\section{Projections and Least-Squares Solutions}
\label{sec:projections}

\exercise \label{exercise:5.2-warmup} (\emph{Warm-up}) Suppose you are lost and thirsty in the desert. For a moment you get some cell reception and your maps app tells you the Derecho River is in the distance, which, oddly enough, is in the shape of a perfectly straight line, all the way to the ocean. 
\begin{enumerate}[label=(\alph*)]
\item If you walk in a straight line to the river, in a direction that makes your walk as short as possible, at what angle will your path meet the river?
\item How many points are there on the river that your straight line path could meet at that angle?
\end{enumerate}

\begin{notes}
\vfill
\end{notes}
\exercise (\emph{Warm-up}) If $\vec{v}$ is a unit vector, then $\vec{v}\cdot \vec{v} = 1$.

\begin{notes}
\vfill
\end{notes}
\exercise (\emph{Warm-up}) Let $\vec{u}=\bbvec{4}{0}$ and $\vec{v}=\bbvec{2}{3}$.
\begin{enumerate}[label=(\alph*)]
\item Is $\vec{u}$ orthogonal to $\vec{v}$?
\item Find $\vec{v}-\frac{\vec{u}\cdot \vec{v}}{\vec{u}\cdot \vec{u}}\vec{u}$. Is it orthogonal to $\vec{v}$?
\end{enumerate}

\begin{notes}
\vfill \hspace{1pt}
\columnbreak
\end{notes}

\subsection{Projections: Basic Idea} \label{subsec:projections-basic-idea}
First, let's define some notation that describes the idea of exercise \ref{exercise:5.2-warmup}, and show that your intuition there was correct in general.

\begin{theorem}[Existence and uniqueness of projections.]\label{theorem:projections}

If $W$ is a subspace of $\R^r$ and $\vec{y}$ is a vector in $\R^r$, then there exists a unique $\hat{y}$ in $W$ such that $\vec{y}-\hat{y}$ is in $W^{\perp}$.

\textbf{Definition:} The $\hat{y}$ that makes this work is called the \textbf{projection of} $\vec{y}$ \textbf{onto} $W$ and is written $\proj_W\vec{y}$.
\end{theorem}

\bbtwocol{.6}{.35}{
In the case of $\R^2$ or $\R^3$, the figures on the right and below, plus a hand-wavy appeal to intuition, might be better than the proof (which you can find in the book if you're interested). Remember any subspace $W$ of $\R^2$ or $\R^3$ is either a single point $\left\{ \vec{0} \right\}$, a line, or a plane, or all of $\R^3$. The theorem is saying that for any point $\vec{y}$, there exists exactly one point $\hat{y}$ in $W$ that makes the line from $\hat{y}$ to $\vec{y}$ perpendicular to $W$.
}{
\begin{tikzpicture}[line cap = round, scale=0.9]
\clip (-4,-1.5) rectangle (1,2);
\draw[thin, blue] (-6,3) -- (6,-3);
\node[below, blue] at (-3,2) {$W$};
%\draw[-Stealth, very thick] (0,0) -- (-2,1)
%node[above right] {$\vec{a}$};
\draw[thin, red] (-3,-6) -- (3,6);
\node[left, red] at (.75,1.5) {$W^{\perp}$};
\node[above] at (0,0) {$\vec{0}$};
%\draw[-Stealth, very thick] (0,0) -- (0.5,1)
%node[right] {$(0.5,1)$};
%\draw[thin] (axis cs:.22,.44) -- (axis cs:-.22,.67) -- (axis cs:-.45,.22);
\coordinate (A) at (0,0);
\coordinate (B) at (-3,-.5);
\coordinate (C) at (-2.2,1.1);
\tkzMarkRightAngle[size=0.4](A,C,B)
\draw[-stealth,very thick] (0,0) -- (-3,-.5) node[midway, below] {$\vec{y}$};
\draw[-stealth,very thick,blue] (0,0) -- (-2.2,1.1) node[midway,above] {$\hat{y}$};
\draw[-stealth,very thick,red] (-2.2,1.1) -- (-3,-.5) node[midway,left] {$\vec{y}-\hat{y}$};
\end{tikzpicture}
}

Note that the \emph{point} $\textcolor{red}{\vec{y}-\hat{y}}$ is in $\textcolor{red}{W^{\perp}}$, which you can see if you draw the arrow starting at $\vec{0}$, but the pictures show a copy of $\textcolor{red}{\vec{y}-\hat{y}}$ drawn tip-to-tail to illustrate $\vec{y}=\textcolor{blue}{\hat{y}}+\textcolor{red}{\vec{y}-\hat{y}}$.

\begin{center}
\begin{tikzpicture}[line cap = round, scale=0.8, x={(0.7cm,-0.7cm)}, y={(1cm,.2cm)}, z={(0cm,1cm)}]
\draw[very thick, red] (0,0,0) -- (0,0,-3);
\fill[blue!30,opacity=0.6] (-2,-2,0) -- (3,-2,0) -- (3,8,0) -- (-2,8,0) -- cycle;
\node[left] at (0,0,0) {$\vec{0}$};
\node[blue] at (-1.5,-1.5,0) {$W$};
\coordinate (A) at (0,0,0);
\coordinate (B) at (0,4,0);
\coordinate (C) at (0,4,3);
\tkzMarkRightAngle[size=0.5](A,B,C)
\draw[-stealth,very thick,blue] (A) -- (B) node[midway, below] {$\hat{y}$};
\draw[-stealth,very thick,red] (B) -- (C);
\node[right,red] at (0,4,2.2) {$\vec{y}-\hat{y}$};
\draw[-stealth,very thick] (A) -- (C) node[midway, above] {$\vec{y}$};
\draw[very thick, red] (0,0,0) -- (0,0,3) node[left] {$W^{\perp}$};
\end{tikzpicture}
\end{center}

\begin{book}
\textbf{Proof.} Since $W$ is a subspace of $\R^r$, by theorem\ref{theorem:subspaces-are-spans} (i), there exists a basis of $W$; call it $\vec{w}_1,\ldots,\vec{w}_n$. Likewise $W^{\perp}$ is also a subspace by theorem \ref{theorem:properties-orthogonal-complements}, so it has a basis; call it $\vec{z}_1,\ldots,\vec{z}_m$. This means $\dim(W) = n$ and $\dim(W^{\perp})=m$, and theorem \ref{theorem:properties-orthogonal-complements} (ii) says that $n+m=r$. This means that the combined list $(\vec{w}_1,\ldots,\vec{w}_n,\vec{z}_1,\ldots,\vec{z}_m)$ has the right number of vectors to be a basis of $\R^r$; and if we can show it's linearly independent, the buy-two-get-one-free theorem says it's a basis. Suppose $x_1,\ldots,x_{n+m}$ is a solution to:
\[x_1\vec{w}_1+\cdots+x_n\vec{w}_n+x_{n+1}\vec{z}_1+\cdots+x_{n+m}\vec{z}_m = \vec{0}\]
The goal is to show the only solution is $x_1=0, \ldots, x_{n+m}=0$. Moving the $W^{\perp}$ components to the other side gives:
\[x_1\vec{w}_1+\cdots+x_n\vec{w}_n = -x_{n+1}\vec{z}_1-\cdots-x_{n+m}\vec{z}_m \]
This shows that the vector $x_1\vec{w}_1+\cdots+x_n\vec{w}_n$ is in both $W$ and $W^{\perp}$ because it can be written as a linear combination of the basis vectors of $W^{\perp}$. The only vector in both $W$ and $W^{\perp}$ is $\vec{0}$, so we have $x_1\vec{w}_1+\cdots+x_n\vec{w}_n = \vec{0}$. Since the basis is linearly independent, this implies $x_1=0, \ldots, x_n=0$.  Likewise, $-x_{n+1}\vec{z}_1-\cdots-x_{n+m}\vec{z}_m = \vec{0}$ since it's in both $W$ and $W^{\perp}$, and $x_{n+1}=0, \ldots, x_{n+m}=0$. Thus $(\vec{w}_1,\ldots,\vec{w}_n,\vec{z}_1,\ldots,\vec{z}_m)$ is linearly independent and a basis of $\R^r$.

Now, since it spans $\R^r$, we know there exists $p_1,\ldots,p_{n+m}$ such that:
\[\vec{y}= p_1\vec{w}_1+\cdots+p_n\vec{w}_n+p_{n+1}\vec{z}_1+\cdots+p_{n+m}\vec{z}_m\]
This makes $p_1\vec{w}_1+\cdots+p_n\vec{w}_n$ work as $\hat{y}$, because 
\[\vec{y}-\hat{y} = \vec{y}- ( p_1\vec{w}_1+\cdots+p_n\vec{w}_n ) =  p_{n+1}\vec{z}_1+\cdots+p_{n+m}\vec{z}_m\]
is in $W^{\perp}$. To show that $\hat{y}$ is unique, suppose $\hat{\hat{y}}$ is another vector that makes $\vec{y}-\hat{\hat{y}}$ in $W^{\perp}$. Since $W^{\perp}$ is a vector space, this means $( \vec{y}-\hat{y} ) - ( \vec{y}-\hat{\hat{y}} ) = \hat{\hat{y}} - \hat{y}$ is in $W^{\perp}$. But $W$ is also a vector space, so $\hat{\hat{y}} - \hat{y}$ is in $W$ too, and the only vector in both $W$ and $W^{\perp}$ is $\vec{0}$, so $\hat{\hat{y}} - \hat{y} = \vec{0}$ and $\hat{\hat{y}} = \hat{y}$, so $\hat{y}=\proj_W\vec{y}$ is unique. \blacksquare
\end{book}

\begin{notes}
\columnbreak
\end{notes}

\begin{theorem}[Best approximation theorem.]\label{theorem:best-approximation}
If $W$ is a subspace of $\R^r$, $\vec{y}$ is a vector in $\R^r$, and $\hat{y}=\proj_W\vec{y}$, then for any vector $\vec{w}$ in $W$ with $\vec{w}\ne \hat{y}$:
\[\nm{\vec{y}-\hat{y}} < \nm{\vec{y}-\vec{w}}\]
\end{theorem}
As the graph below illustrates, this theorem says that the any line segment from $\vec{y}$ to $W$ that meets $W$ non-orthogonally is longer than the line segment that meets it orthogonally. The proof shows that $\vec{y}-\vec{w}$ is the hypotenuse of a right triangle with $\vec{y}-\hat{y}$ as a leg.
\begin{center}
\begin{tikzpicture}[line cap = round, scale=0.8, x={(0.7cm,-0.7cm)}, y={(1cm,.2cm)}, z={(0cm,1cm)}]
\draw[very thick, red] (0,0,0) -- (0,0,-3);
\fill[blue!30,opacity=0.6] (-2,-2,0) -- (3,-2,0) -- (3,8,0) -- (-2,8,0) -- cycle;
\node[left] at (0,0,0) {$\vec{0}$};
\node[blue] at (-1.5,-1.5,0) {$W$};
\coordinate (A) at (0,0,0);
\coordinate (B) at (0,4,0);
\coordinate (C) at (0,4,3);
\coordinate (W) at (2,6,0);
\tkzMarkRightAngle[size=0.5](A,B,C)
\draw[-stealth,very thick,blue] (A) -- (B) node[midway, above] {$\hat{y}$};
\draw[-stealth,very thick,red] (B) -- (C);
\node[right,red] at (0,2.6,1.5) {$\vec{y}-\hat{y}$};
\draw[-stealth,very thick] (A) -- (C) node[midway, above] {$\vec{y}$};
\draw[very thick, red] (0,0,0) -- (0,0,3) node[left] {$W^{\perp}$};
\draw[-stealth,very thick,blue] (A) -- (W) node[midway, below] {$\vec{w}$};
\draw[-stealth,very thick] (W) -- (C) node[midway, right] {$\vec{y}-\vec{w}$};
\end{tikzpicture}
\end{center}

\begin{book}
\textbf{Proof.} Let $\vec{w}$ be a vector in $W$ with $\vec{w}\ne \hat{y}$. Since $W$ is a vector space, $\vec{w}-\hat{y}$ is in $W$, and thus orthogonal to $\vec{y}-\hat{y}$. The Pythagorean theorem gives:
\begin{align*}
\nm{\vec{y}-\hat{y}}^2+\nm{\vec{w}-\hat{y}}^2 &= \nm{\right(\vec{y}-\hat{y}\left) - (\vec{w}-\hat{y})}^2 = \nm{ \vec{y}-\vec{w}}^2 \\
\nm{\vec{y}-\hat{y}}^2 &= \nm{ \vec{y}-\vec{w}}^2 - \nm{\vec{w}-\hat{y}}^2
\end{align*}
Since $\vec{w}\ne \hat{y}$, we know $\nm{\vec{w}-\hat{y}}^2>0$, so $\nm{\vec{y}-\hat{y}}^2<\nm{\vec{y}-\vec{w}}^2$. \blacksquare
\end{book}

\begin{notes}
\columnbreak
\end{notes}

\example \label{example:5.2-intro} Let $\vec{y}=(4,5,6)$, $W = \col \begin{bmatrix}
\vec{e}_1 & \vec{e}_2
\end{bmatrix}$, and $\vec{v}=\frac{\vec{e}_1\cdot \vec{y}}{\vec{e}_1\cdot \vec{e}_1}\vec{e}_1 + \frac{\vec{e}_2\cdot \vec{y}}{\vec{e}_2\cdot \vec{e}_2}\vec{e}_2$.
\begin{enumerate}[label=(\alph*)]
\item Find $\vec{v}$ and $\vec{y}-\vec{v}$.
\item Show that $\proj_W\vec{y} = \vec{v}$.
\end{enumerate}

\begin{notes}
\vfill
\end{notes}
\begin{book}
\textbf{Solution.} (a) $\vec{v}=\frac{4(1)+5(0)+6(0)}{1(1)+0(0)+0(0)}\bbbvec{1}{0}{0} + \frac{4(0)+5(1)+6(0)}{0(0)+1(1)+0(0)}\bbbvec{0}{1}{0} = \bbbvec{4}{5}{0}$.

$\vec{y}-\vec{v} =\bbbvec{4}{5}{6}-\bbbvec{4}{5}{0}= \bbbvec{0}{0}{6}$. \blacksquare
%Any vector in $W$ (the $x_1x_2-plane$) including $\proj_W\vec{y}$ must be of the form $(x_1,x_2,0)$. And $W^{\perp}$ is the set of vectors perpendicular to $\vec{e}_1$ and $\vec{e}_2$, which must be $\spn ( \vec{e}_3 )$. 
\end{book}

We'll soon see that the formula for $\vec{v}=\proj_W\vec{y}$ in example \ref{example:5.2-intro} turns out to work in general, when you have a basis of $W$ made of vectors that are orthogonal to each other. 

\subsection{Projections with Orthogonal Lists} \label{subsec:projections-with-orthogonal-lists}

\begin{definition}
A list of non-zero vectors $( \vec{u}_1,\ldots,\vec{u}_c )$ is called an \textbf{orthogonal list} if every vector is orthogonal to every other vector: $\vec{u}_i\cdot \vec{u}_j=0$ for $i\ne j$. If all the vectors have length one, it's an \textbf{orthonormal list}.
\end{definition}

Note that if a list contains only one vector, it's vacuously orthogonal to every other vector in the list, making the list orthogonal.

\example The standard basis vectors of $\R^{3}$, $\vec{e}_1,\vec{e}_2,\vec{e}_{3}$, are an orthonormal list.

\begin{notes}
\vspace{2em} \hspace{1pt}

\columnbreak
\end{notes}
\example \label{example:5.2-basic-orthonormal} Show that the columns of $U$ make an orthonormal list of vectors: $U = \begin{bmatrix}
\vec{u}_1 & \vec{u}_2 & \vec{u}_3
\end{bmatrix} = 
\begin{bmatrix}
1/2 & 1/2 & 1/2 \\
1/2 & 1/2 & -1/2 \\
1/2 & -1/2 & 1/2 \\
1/2 & -1/2 & -1/2
\end{bmatrix}$.

\begin{notes}
\vfill
\end{notes}
\begin{book}
\textbf{Solution.} First let's check that they're orthogonal to each other:
\begin{align*}
\vec{u}_1\cdot \vec{u}_2 &= \frac{1}{2}\cdot\frac{1}{2} + \frac{1}{2}\cdot\frac{1}{2} + \frac{1}{2}\cdot\frac{-1}{2} + \frac{1}{2}\cdot\frac{-1}{2} = 0\\
\vec{u}_1\cdot \vec{u}_3 &= \frac{1}{2}\cdot\frac{1}{2} + \frac{1}{2}\cdot\frac{-1}{2} + \frac{1}{2}\cdot\frac{1}{2} + \frac{1}{2}\cdot\frac{-1}{2} = 0 \\
\vec{u}_2\cdot \vec{u}_3 &= \frac{1}{2}\cdot\frac{1}{2} + \frac{1}{2}\cdot\frac{-1}{2} + \frac{-1}{2}\cdot\frac{1}{2} + \frac{-1}{2}\cdot\frac{-1}{2} = 0
\end{align*}
Then to see that they're unit vectors:
\begin{align*}
\nm{\vec{u}_1}^2 &=   \frac{1}{2}\cdot\frac{1}{2} + \frac{1}{2}\cdot\frac{1}{2} + \frac{1}{2}\cdot\frac{1}{2} + \frac{1}{2}\cdot\frac{1}{2} = 1\\
\nm{\vec{u}_2}^2 &=   \frac{1}{2}\cdot\frac{1}{2} + \frac{1}{2}\cdot\frac{1}{2} + \frac{-1}{2}\cdot\frac{-1}{2} + \frac{-1}{2}\cdot\frac{-1}{2} = 1\\
\nm{\vec{u}_3}^2 &=   \frac{1}{2}\cdot\frac{1}{2} + \frac{-1}{2}\cdot\frac{-1}{2} + \frac{1}{2}\cdot\frac{1}{2} + \frac{-1}{2}\cdot\frac{-1}{2} = 1
\end{align*}
\end{book}

\example With the same $U$ from example \ref{example:5.2-basic-orthonormal}, find $UU^{T}$ and $U^TU$.

\begin{notes}
\vfill \hspace{1pt} \columnbreak
\end{notes}
\begin{book}
Using the row-column rule:
\[UU^T = \begin{bmatrix}
3/4 & 1/4 & 1/4 & 1/4 \\
1/4 & 3/4 & -1/4 & -1/4 \\
1/4 & -1/4 & 3/4 & 3/4 \\
1/4 & -1/4 & 3/4 & 3/4
\end{bmatrix}\]
Notice that $UU^T$ is \emph{symmetric}, which means the entry in row $i$, column $j$ is equal to the entry in row $j$, column $i$. In other words, $( UU^T )^T = UU^T$. This will be important in the next chapter. For $U^TU$, let's show some more work, which will help justify the theorem below.
\begin{align*}
U^TU&= \begin{bmatrix}
\vec{u}_1^T \\ \vec{u}_2^T \\ \vec{u}_3^T
\end{bmatrix}\begin{bmatrix}
\vec{u}_1 & \vec{u}_2 & \vec{u}_3
\end{bmatrix} \\
&= \begin{bmatrix}
\vec{u}_1^T\vec{u}_1 & \vec{u}_1^T\vec{u}_2 & \vec{u}_1^T\vec{u}_3 \\
\vec{u}_2^T\vec{u}_1 & \vec{u}_2^T\vec{u}_2 & \vec{u}_2^T\vec{u}_3 \\
\vec{u}_3^T\vec{u}_1 & \vec{u}_3^T\vec{u}_2 & \vec{u}_3^T\vec{u}_3 
\end{bmatrix} & (\text{row-column rule}) \\
&= \begin{bmatrix}
\vec{u}_1\cdot\vec{u}_1 & \vec{u}_1\cdot\vec{u}_2 & \vec{u}_1\cdot\vec{u}_3 \\
\vec{u}_2\cdot\vec{u}_1 & \vec{u}_2\cdot\vec{u}_2 & \vec{u}_2\cdot\vec{u}_3 \\
\vec{u}_3\cdot\vec{u}_1 & \vec{u}_3\cdot\vec{u}_2 & \vec{u}_3\cdot\vec{u}_3 
\end{bmatrix}
\end{align*}
At this point, note that all the entries on the main diagonal are $\vec{u}_i\cdot \vec{u}_i = \nm{\vec{u}_i}^2 = 1$ because they're unit vectors. And because they're all orthogonal to each other, every entry not on the main diagonal is 0. So: \[U^TU = \begin{bmatrix}
1 & 0 & 0 \\
0 & 1 & 0 \\
0 & 0 & 1
\end{bmatrix}\]

\blacksquare
\end{book}

\begin{theorem}
\label{theorem:UTU-I}
If $U$ is a matrix with orthonormal columns, then $U^TU=I$.
\end{theorem}

The proof of the case where $U$ has three columns was done in the last example. If $U$ has more columns, the same principle applies: All the diagonal entries of $U^TU$ are $\vec{u}_i\cdot \vec{u}_i = \nm{\vec{u}_i}^2=1$ because they're unit vectors, and the non-diagonal entries, with $i\ne j$, are $\vec{u}_i\cdot \vec{u}_j=0$ because they're orthogonal.

\begin{theorem}[Orthogonal projection formula.]\label{theorem:orthogonal-projection-formula}

If $\vec{u}_1,\ldots,\vec{u}_c$ is an orthogonal basis of a subspace $W$ of $\R^r$, and $\vec{y}$ is in $\R^n$, then:
\[\proj_W\vec{y}=\frac{\vec{u}_1\cdot \vec{y}}{\vec{u}_1\cdot \vec{u}_1}\vec{u}_1+\cdots + \frac{\vec{u}_c\cdot \vec{y}}{\vec{u}_c\cdot \vec{u}_c}\vec{u}_c\]
\end{theorem}

Note that if $\vec{u}_1,\ldots,\vec{u}_c$ is an ortho\emph{normal} basis of $W$, then all the denominators of the coefficients equal one, making the formula simpler.

\begin{notes}
\vfill \hspace{1pt} \columnbreak
\end{notes}
\begin{book}
\textbf{Proof.} We know $\proj_W\vec{y}$ exists and is unique; and it's in $W$, so it can be written as a linear combination of the basis vectors, so let's denote it:
\[\proj_W\vec{y}=\hat{y}=x_1\vec{u}_1+\cdots+x_c\vec{u}_c\]
We want to show that each  $x_i= \frac{\vec{u}_i\cdot\vec{y}}{\vec{u}_i\cdot \vec{u}_i}$. A good strategy when trying to prove things about orthogonal lists is to make dot products that you know are zero. Let $\vec{u}_i$ be one of the basis vectors. Its dot product with $\hat{y}$ is:
\begin{align*}
\vec{u}_i\cdot \hat{y} &= \vec{u}_i\cdot ( \hat{y}=x_1\vec{u}_1+\cdots+x_n\vec{u}_n ) \\
\vec{u}_i\cdot \hat{y}&=x_1\vec{u}_i\cdot \vec{u}_1 + \cdots + x_i\vec{u}_i\cdot \vec{u}_i + \cdots + x_c\vec{u}_i\cdot \vec{u}_c \\
\vec{u}_i\cdot \hat{y}&= 0 + \cdots + 0 + c_i\vec{u}_i\cdot \vec{u}_i + 0 + \cdots + 0
\end{align*}
Dividing both sides by $\vec{u}_i\cdot \vec{u}_i$ gives $\frac{\vec{u}_i\cdot \hat{y}}{\vec{u}_i\cdot \vec{u}_i} =  x_i$. We would be done, except this has $\vec{u}_i\cdot\hat{y}$ where $\vec{u}_i\cdot\vec{y}$ is supposed to be. But they are equal because:
\begin{align*}
\vec{u}_i\cdot\vec{y} &= ( \hat{y} + ( \vec{y}-\hat{y} ) )\cdot \vec{u}_i \\
&= \hat{y}\cdot \vec{u}_i + ( \vec{y}-\hat{y} )\cdot \vec{u}_i \\
&= \hat{y}\cdot \vec{u}_i + 0
\end{align*}
The last step follows from the definition of $\hat{y}= \proj_W\vec{y}$, which says that $\vec{y}-\hat{y}$ is orthogonal to any vector in $W$, including $\vec{u}_i$. \blacksquare
\end{book}

\example Let $W=\spn(1,1,1)$ and $\vec{y}=(4,5,6)$. Find the nearest point in $W$ to $\vec{y}$, and
find the distance from $\vec{y}$ to that point.

\begin{notes}
\vfill
\end{notes}
\begin{book}
\textbf{Solution.} The best approximation theorem says that the nearest point in $W$ to $\vec{y}$ is $\proj_W\vec{y}$. With $\vec{u}=(1,1,1)$, the list of one vector, $( \vec{u} )$, is vacuously orthogonal, so the orthogonal projection formula gives:
\[\proj_W\vec{y} = \frac{\vec{u}\cdot \vec{y}}{\vec{u}\cdot \vec{u}}\vec{u}=\frac{1(4) + 1(5) + 1(6)}{1(1)+1(1)+1(1)}\vec{u}=(5,5,5)\]

The distance from $\vec{y}$ to the this point is:
\[\nm{\vec{y}-(5,5,5)} = \nm{(-1,0,1)} = \sqrt{(-1)^2+1^2} = \sqrt{2}\]
\end{book}

\example A block weighing 4 pounds is resting on an inclined plane with slope $\frac{1}{\sqrt{3}}$. Find components of its weight that are parallel and perpendicular to the slope.

\begin{center}
\begin{tikzpicture}
\draw[thin] (0,0)--(4.5*1.73,4.5) -- (4.5*1.73,0) -- (0,0);
\draw[thick,-stealth] (0,0) -- (1.73,1) node[above left] {$(\sqrt{3}/2,1/2) = \vec{u}$};
\draw[thick] (1.73*3,3) -- (1.73*3.5,3.5) -- (1.73*3.5-.5,3.5+.5*1.73) -- (1.73*3-.5,3+.5*1.73)--cycle ;
\coordinate (B) at (5.37, 3.68);
\coordinate (W) at (5.37, 0.68);
\draw[thick,-stealth] (B) -- (W) node[midway,right]{$\vec{y}=\bbvec{0}{-4}$};
\draw[thick,-stealth] (B) -- (5.37-1.73/2*1.5, 3.68-.5*1.5);
\node[above right] at (5.37-1.73/2*1.3, 3.68-.5*1.3+.1) {$\hat{y}$};
\draw[thick,-stealth] (5.37-1.73/2*1.5, 3.68-.5*1.5)--(W) node[midway, left] {$\vec{y}-\hat{y}$};
\end{tikzpicture}
\end{center}

\begin{notes}
\vfill \hspace{1pt} \columnbreak
\end{notes}
\begin{book}
\textbf{Solution.} Using the notation in the figure, they weight of the block is $\vec{y} = (0,-4)$ pounds, and $\vec{u}=( \sqrt{3}/2, 1/2 )$ is a unit vector that forms a basis for the slope of the block. Using the orthogonal projection formula for $\proj_{\vec{u}}\vec{y}$, the denominator $\vec{u}\cdot \vec{u}$ is $1$. So, the parallel component is:
\[\hat{y} = \proj_{\vec{u}}\vec{y} = (\vec{u}\cdot \vec{y})\vec{u} = \left(\frac{\sqrt{3}}{2}(0)+\frac{1}{2}(-4)\right)\bbvec{\sqrt{3}/2}{1/2} = \bbvec{-\sqrt{3}}{-1}  \]
The perpendicular component is:
$\vec{y}-\hat{y} = \bbvec{0}{-4}- \bbvec{-\sqrt{3}}{-1} = \bbvec{\sqrt{3}}{-3}$
The \emph{magnitudes} of these components, which you would probably need when doing physics or engineering, are
\[\nm{\hat{y}} = \sqrt{(-\sqrt{3})^2+(-1)^2} = 2 \text{ and } \nm{\vec{y}-\hat{y}} = \sqrt{\sqrt{3}^2+(-3)^2} = 2\sqrt{3}\]
and these add up to more than $\nm{\vec{y}} = 4$, even though the vectors $\hat{y}$ and $\vec{y}-\hat{y}$ add up to $\vec{y}$. \blacksquare
\end{book}

\subsection{Least-Squares Solutions} \label{subsec:least-squares-solutions}

Now, what if you want to find $\proj_W\vec{y}$, but you don't have an orthogonal basis of $W$? What if you don't even have a basis of $W$? Let's say that at least you have a list $\vec{a}_1,\ldots,\vec{a}_c$ (maybe linearly dependent) that spans $W$, so that $W = \col A$ with $A = \begin{bmatrix}
\vec{a}_1 & \cdots & \vec{a}_c
\end{bmatrix}$. In this case, finding $\proj_{\col A}\vec{y}$ amounts to finding a $\hat{y}$ that can be written $\hat{y}=A\hat{x}$ for some $\hat{x}$, so that $A\hat{x}$ is \emph{as close as possible} to $\vec{y}$. Here's the illustration of theorem \ref{theorem:best-approximation} in the situation $W=\col A$, where $A\vec{x}$ is some other vector in $\col A$:

\begin{center}
\begin{tikzpicture}[line cap = round, scale=0.8, x={(0.7cm,-0.7cm)}, y={(1cm,.2cm)}, z={(0cm,1cm)}]
\draw[very thick, red] (0,0,0) -- (0,0,-3);
\fill[blue!30,opacity=0.6] (-2,-2,0) -- (3,-2,0) -- (3,8,0) -- (-2,8,0) -- cycle;
\node[left] at (0,0,0) {$\vec{0}$};
\node[blue] at (-1.5,-1,0) {$W=\col A$};
\coordinate (A) at (0,0,0);
\coordinate (B) at (0,4,0);
\coordinate (C) at (0,4,3);
\coordinate (W) at (2,6,0);
\tkzMarkRightAngle[size=0.5](A,B,C)
\draw[-stealth,very thick,blue] (A) -- (B) node[midway, above] {$\hat{y} = A\hat{x}$};
\draw[-stealth,very thick,red] (B) -- (C);
\node[right,red] at (0,2.3,1.3) {$\vec{y}-A\hat{x}$};
\draw[-stealth,very thick] (A) -- (C) node[midway, above] {$\vec{y}$};
\draw[very thick, red] (0,0,0) -- (0,0,3) node[left] {$(\col A)^{\perp}$};
\draw[-stealth,very thick,blue] (A) -- (W) node[midway, below] {$A\vec{x}$};
\draw[-stealth,very thick] (W) -- (C) node[midway,right] {$\vec{y}-A\vec{x}$};
\end{tikzpicture}
\end{center}

Theorem \ref{theorem:projections} guarantees that such a $\hat{y}$ exists and is unique. The trick to finding $\hat{x}$ and $\hat{y}$ is to notice that theorem \ref{theorem:four-subspaces} says that $(\col A)^{\perp}=\nul (A^T)$. So, we're trying to find $\hat{x}$ such that $\vec{y}-A\hat{x}$ is in $\nul (A^T)$ which means:
\begin{align}\label{eq:5.1-ATAx}
A^T ( \vec{y}-A\hat{x} ) &= \vec{0} \nonumber\\
A^T\vec{y} - A^TA\hat{x} &= \vec{0} \nonumber\\
A^TA\hat{x}=A^T\vec{y}
\end{align}
Note that theorem \ref{theorem:projections} guarantees that $\hat{y}=A\hat{x}$ exists and is unique, but that doesn't mean that $\hat{x}$ is unique; there may be more that one $\hat{x}$ that satisfies equation \ref{eq:5.1-ATAx}. In any case, this proves:

\begin{theorem}[Least-squares solutions.]\label{theorem:least-squares}

Let $A$ be an $r\times c$ matrix, and $\vec{y}$ be a vector in $\R^r$. Then $A^TA\hat{x}=A^T\vec{y}$ if and only if $A\hat{x} = \proj_{\col A}\vec{y}$, which is the closest point in $\col A$ to $\vec{y}$ \,(that is, $\nm{\vec{y}-A\hat{x}}<\nm{\vec{y}-A\vec{x}}$ for any $\vec{x}\ne \hat{x}$).

\textbf{Definition.} The $\hat{x}$ described in the theorem is called the \textbf{least-squares solution} of the equation $A\vec{x}=\vec{b}$, because it minimizes $\nm{\vec{y}-A\hat{x}}$, whose formula involves a sum of squares.
\end{theorem}

\example In example \ref{example:5.2-intro} we found that with $A = \begin{bmatrix}
\vec{e}_1 & \vec{e}_2
\end{bmatrix}$ and $\vec{y}=(4,5,6)$, $\proj_{\col A} \vec{y}=(4,5,0)$. Check that theorem \ref{theorem:least-squares} gives the same result.

\begin{notes}
\vfill
\end{notes}
\begin{book}
\textbf{Solution.} Theorem \ref{theorem:least-squares} says that $\proj_{\col A} = A\hat{x}$ where:
\begin{align*}
A^TA\hat{x} &= A^T\vec{y} \\
\begin{bmatrix}
1 & 0 \\
0 & 1
\end{bmatrix}\hat{x} &= \begin{bmatrix}
1 & 0 & 0 \\
0 & 1 & 0
\end{bmatrix} \bbbvec{4}{5}{6} \\
\hat{x} &= \bbvec{4}{5}
\end{align*}
Thus $\hat{x}=\bbvec{4}{5}$ is the least-squares solution of $A\vec{x}=\vec{y}$, and $\proj_{\col A}\vec{y}=A\hat{x} = 4\vec{e}_1+5\vec{e}_2=(4,5,0)$. \blacksquare
\end{book}

\example Describe least-squares solutions to $A\vec{x}=\vec{y}$, with $A = \begin{bmatrix}
-1 & 0 & 2 \\
2 & 0 & -4 \\
1 & 2 & 3
\end{bmatrix}$ and $\vec{y}=\bbbvec{1}{3}{3}$.

\begin{notes}
\vfill \hspace{1pt} \columnbreak
\end{notes}
\begin{book}
\textbf{Solution.} First note that $A\vec{x}=\vec{y}$ doesn't have an actual solution; the first two components of $A\vec{x}$ are always going to be a multiple of $\bbvec{-1}{2}$, and so can't be equal to $\bbvec{1}{3}$. So the best we can do is get $A\hat{x}$ as close as possible to $\vec{y}$, and theorem \ref{theorem:least-squares} says we need to solve:
\begin{align*}
A^TA\hat{x}&=A^T\vec{y} \\
\begin{bmatrix}
-1 & 0 & 2 \\
2 &0 & -4 \\
1 & 2 & 3
\end{bmatrix}\begin{bmatrix} 
-1 & 2 & 1 \\
0 & 0 & 2 \\
2 & -4 & 3
\end{bmatrix}\hat{x} &= \begin{bmatrix}
-1 & 2 & 1 \\
0 & 0 & 2 \\
2 & -4 & 3
\end{bmatrix}\bbbvec{1}{3}{3} \\
\begin{bmatrix}
6 & 2 & -7 \\
2 & 4 & 6 \\
-7 & 6 & 29
\end{bmatrix}\hat{x} & = \bbbvec{8}{6}{-1} \\
\end{align*}
Row-reducing results in:
\begin{align*}
\begin{bmatrix}
1 & 0 & -2 \\
0 & 1 & 2.5 \\
0 & 0 & 0
\end{bmatrix}\hat{x} &= \bbbvec{1}{1}{0} \\
\begin{aligned}
x_1-2x_3&=1 \\
x_2+2.5x_3&=1 \\
0 &= 0
\end{aligned} \\
\hat{x} = \bbbvec{x_1}{x_2}{x_3} &= \bbbvec{1+2x_3}{1-2.5x_3}{0+x_3}
\end{align*}
So we have an infinite number of least-squares solutions, such as, with $x_3=0$, $\hat{x}=(1,1,0)$. Or, with $x_3=-4$, $\hat{x}=(-7,11,-4)$. Let's check that they give the same $A\hat{x}=\proj_{\col A} \vec{y}$:
\[\begin{bmatrix}
-1 & 0 & 2 \\
2 & 0 & -4 \\
1 & 2 & 3
\end{bmatrix} \bbbvec{1}{1}{0} = \bbbvec{-1}{2}{3}\quad\text{ and } \quad \begin{bmatrix}
-1 & 0 & 2 \\
2 & 0 & -4 \\
1 & 2 & 3
\end{bmatrix} \bbbvec{-7}{11}{-4} = \bbbvec{-1}{2}{3}\]
\blacksquare
\end{book}


\columnbreak
\exerciseheading
%\begin{savenotes}

\textbf{\ref{subsec:projections-basic-idea} Projections: Basic Idea}

%The first four problems here are good basics checks to do before orthogonal lists.
\exatend If $\vec{y}$ is in $W$, what is $\proj_W\vec{y}$? What is $\vec{y}-\proj_w\vec{y}$?

\exatend If $W = \left\{ \vec{0} \right\}$, and $\vec{y}$ is not in $W$, what is $\proj_W\vec{y}$? What is $\vec{y}-\proj_w\vec{y}$?

\exatend Let $\vec{y}=(\pi, \sqrt{2})$ and $W = \col \begin{bmatrix}
\vec{e}_1 
\end{bmatrix}$ aka the $x_1$-axis. Find $\proj_W\vec{y}$ and $\vec{y}-\proj_W\vec{y}$.
\exatend Let $\vec{y}=(-4,-2,7)$ and $W = \col \begin{bmatrix}
\vec{e}_1 & \vec{e}_3
\end{bmatrix}$, aka the $x_1x_3$-plane. Find $\proj_W\vec{y}$ and $\vec{y}-\proj_W\vec{y}$.

\exatend (T/F?) For some subspaces $W$ of $\R^r$, and some $\vec{y}$ in $\R^r$, there does not exist a $\hat{y}$ in $W$ such that $\vec{y}-\hat{y}$ is in $W^{\perp}$.

\exatend (T/F?) For some subspaces $W$ of $\R^r$, and some $\vec{y}$ in $\R^r$, there exists more than one vector $\hat{y}$ in $W$ such that $\vec{y}-\hat{y}$ is in $W^{\perp}$.

\exatend (T/F?) If $W$ is a subspace of $\R^r$ and $\vec{y}$ is in $\R^r$, then $\proj_W\vec{y}$ is closer to $\vec{y}$ than any other point in $W$.

\exatend (T/F?) If $W$ is a subspace of $\R^r$ and $\vec{y}$ is in $\R^r$, then $\nm{\proj_W\vec{y}}$ is the distance from $\vec{y}$ to the closest point in $W$.

\exatend (T/F?) Let $\vec{y}$ be in $\R^r$, and $A$ be $r\times c$. If $W = \col A$, then the equation $A\vec{x}=\proj_W\vec{y}$ is consistent.

\exatend If $\vec{y}$ is in $W^{\perp}$, what is $\proj_W\vec{y}$?

\textbf{\ref{subsec:projections-with-orthogonal-lists} Projections with Orthogonal Lists}

\exatend Let $U = \begin{bmatrix}
\vec{u}_1 & \vec{u}_2
\end{bmatrix} = \begin{bmatrix}
3 & -4 \\
4 & 3
\end{bmatrix}$. Is the list $\vec{u}_1,\vec{u}_2$ orthogonal? Is it orthonormal? Find $U^TU$.

\exatend Let $U= \begin{bmatrix}
\vec{u}_1 & \vec{u}_2
\end{bmatrix} = \begin{bmatrix}
0.6 & 0.8 \\
0.8 & 0.6
\end{bmatrix}$. Is the list $\vec{u}_1,\vec{u}_2$ orthogonal? Is it orthonormal? Find $U^TU$.

\exatend \label{exercise:5.2basic-ortho}Let $U= \begin{bmatrix}
\vec{u}_1 & \vec{u}_2 & \vec{u}_3
\end{bmatrix} = \begin{bmatrix}
-1/3 & 2/3 \\
2/3 & -1/3 \\
2/3 & 2/3 
\end{bmatrix}$. Is the list $\vec{u}_1,\vec{u}_2$ orthogonal? Is it orthonormal? Find $U^TU$ and $UU^T$.

\exatend With the same $\vec{u}_1,\vec{u}_2$ as in exercise \ref{exercise:5.2basic-ortho}, find $\vec{u}_3$ so that $\vec{u}_1,\vec{u}_2,\vec{u}_3$ is an orthonormal list. \hintfootnote{$(\col U)^{\perp}=\nul(U^T)$.}

\exatend Let $A=\begin{bmatrix}
2 & 10 & -5 \\
2 & -11 & -2 \\
1 & h & k
\end{bmatrix}$.
\begin{enumerate}[label=(\alph*)]
\item Find $h$ and $k$ so that the columns of $A$ form an orthogonal basis of $\R^3$.
\item Make the columns of $A$ into an orthonormal list by dividing each vector by its length.
\end{enumerate}

\exatend (T/F?) If $\vec{u}_i$ is one vector from an orthogonal list of vectors, then $\vec{u}_i\cdot \vec{u}_i=0$.

\exatend (T/F?) If $\vec{b}_1,\ldots,\vec{b}_c$ is a basis of $W$, then $\proj_W\vec{y} = \frac{\vec{b}_1\cdot \vec{y}}{\vec{b}_1\cdot \vec{b}_1}\vec{b}_1+\cdots + \frac{\vec{b}_c\cdot \vec{y}}{\vec{b}_c\cdot \vec{b}_c}\vec{b}_c$.

\exatend Prove theorem \ref{theorem:orthogonal-lists-are-independent} below.  \hintfootnote{One way to do this is with theorem \ref{theorem:UTU-I}. Another way is with theorems \ref{theorem:projections} and \ref{theorem:orthogonal-projection-formula}.}

\begin{theorem}[Orthogonal lists are independent.]\label{theorem:orthogonal-lists-are-independent}

If $\vec{u}_1,\ldots,\vec{u}_n$ is an orthogonal list of vectors, then it's linearly independent too.
\end{theorem}
\exatend Suppose $U$ is an $r \times c$ matrix with orthogonal columns.
\begin{enumerate}[label=(\alph*)]
\item Show that $r \ge c$. \hintfootnote{Theorem \ref{theorem:orthogonal-lists-are-independent}.}
\item What are the sizes (rows $\times$ columns) of $UU^T$ and $U^TU$?
\end{enumerate}
\exatend Suppose $\vec{u} = k\vec{v}$ with $k\ne 0$, and let $W=\spn(\vec{v})$. We know the vectors are parallel, so $W = \spn(\vec{v})=\spn(\vec{u})$. Hopefully the formula in theorem \ref{theorem:orthogonal-projection-formula} gives the same result for $\proj_W\vec{y}$ whether you use $\vec{u}$ or $\vec{v}$ as a basis of $W$. Check that this is indeed the case.
\exatend \label{exercise:UUTy}Suppose $U=\begin{bmatrix}
\vec{u}_1 & \cdots & \vec{u}_c
\end{bmatrix}$ and is a matrix with orthonormal columns, and let $W = \col U$. Show that the formula in theorem \ref{theorem:orthogonal-projection-formula} can be written $\proj_W\vec{y}= UU^T\vec{y}$. \hintfootnote{$U(U^T\vec{y})$.}

\exatend Let $U= \begin{bmatrix}
\vec{u}_1 & \vec{u}_2 & \vec{u}_3
\end{bmatrix} = \begin{bmatrix}
-1/3 & 2/3 \\
2/3 & -1/3 \\
2/3 & 2/3 
\end{bmatrix}$ and $\vec{y}=(9,0,0)$. Find $\proj_{\col U}\vec{y}$. (You showed in exercise \ref{exercise:5.2basic-ortho} that $(\vec{u}_1,\vec{u}_2)$ is an orthonormal list.)

\exatend Let $U = \begin{bmatrix}
\vec{u}_1 & \vec{u}_2 & \vec{u}_3
\end{bmatrix} = \begin{bmatrix}
-12 & -4 \\
-4 & 3 \\
3 & -12 
\end{bmatrix}$. It can be shown that $(\vec{u}_1,\vec{u}_2)$ is an orthogonal list. Find the closest point in $\col U$ to $\vec{y}=(13,13,13)$ and the distance from $\vec{y}$ to that point. %solution: proj = (16,1,9), distance = 13

\exatend It can be shown that the columns of $U = \begin{bmatrix}
\vec{u}_1 & \vec{u}_2
\end{bmatrix} = 
\begin{bmatrix}
1/2 & 1/2  \\
1/2 & 1/2  \\
1/2 & -1/2 \\
1/2 & -1/2
\end{bmatrix}$ are an orthonormal list of vectors. Find the nearest point in $\col U$ to $\vec{y}=(4,0,0,0)$ and the distance from $\vec{y}$ to that point.

\exatend Suppose $U$ is an $r\times c$ matrix with orthonormal columns, and let $\vec{x},\vec{y}$ be $\R^c$ vectors. Show that $( U\vec{x} )\cdot ( U\vec{y} )=\vec{x}\cdot \vec{y}$. \hintfootnote{The definition of $\vec{u}\cdot \vec{v}$ is $\vec{u}^T\vec{v}$.}
\exatend Suppose $U$ is an $r\times c$ matrix with orthonormal columns, and let $\vec{x}$ be in $\R^c$. Show that $\nm{U\vec{x}} = \nm{\vec{x}}$.
\exatend Suppose $\vec{u}$ is a unit vector in $\R^c$, and let $W=\spn ( \vec{u} )$. Define a function $T(\vec{y})$ that reflects $\vec{y}$ to the point on the opposite side of $W$ from $\vec{y}$, as shown in the graph below.

\bbtwocol{0.6}{0.35}{
\begin{enumerate}[label=(\alph*)]
\item Show that $T(\vec{y}) = 2\proj_W\vec{y}-\vec{y}$.
\item Use the result of exercise \ref{exercise:UUTy} to find $A$ so that $T(\vec{y}) = A\vec{y}$.
%solution: A = 2UU^T - I
\item Use the result of (b) to find $A$ so that $T(\vec{y})$ reflects $\vec{y}$ across the line $x_2=3x_1$. %solution: A=[[-.8,.6],[.6,.8]]
\end{enumerate}}{
\begin{tikzpicture}[line cap = round, scale=0.9]
\clip (-4,-1.5) rectangle (1,3);
%\draw[-Stealth, very thick] (0,0) -- (-2,1)
%node[above right] {$\vec{a}$};
%\draw[thin, red] (-3,-6) -- (3,6);
%\node[left, red] at (.75,1.5) {$W^{\perp}$};
\node[below] at (0,0) {$\vec{0}$};
%\draw[-Stealth, very thick] (0,0) -- (0.5,1)
%node[right] {$(0.5,1)$};
%\draw[thin] (axis cs:.22,.44) -- (axis cs:-.22,.67) -- (axis cs:-.45,.22);
\coordinate (A) at (0,0);
\coordinate (B) at (-3,-.5);
\coordinate (C) at (-2.2,1.1);
\tkzMarkRightAngle[size=0.4](A,C,B)
\draw[-stealth,very thick] (0,0) -- (-3,-.5) node[midway, below] {$\vec{y}$};
\draw[-stealth,very thick] (0,0) -- (-1.4,2.7) node[midway, right] {$T(\vec{y})$};
\draw[-stealth,very thick,red] (-2.2,1.1) -- (-3,-.5) node[midway,left] {$\vec{y}-\hat{y}$};
\draw[-stealth,very thick,red] (-2.2,1.1) -- (-1.4,2.7) node[midway,left] {$-\vec{y}+\hat{y}$};
\draw[thin, blue] (-6,3) -- (6,-3);
\node[below, blue] at (-3.5,2.4) {$W$};
\draw[-stealth,very thick,blue] (0,0) -- (-2.2,1.1) node[midway,above] {$\hat{y}$};
\end{tikzpicture}}

\textbf{\ref{subsec:least-squares-solutions} Least-Squares Solutions}

%% Least-squares problems start here
For problems \ref{exercise:5.2least-squares-TF-start} to \ref{exercise:5.2least-squares-TF-end}, Assume $A$ is an $r \times c$ matrix and $\vec{y}$ is in $\R^r$.

\exatend \label{exercise:5.2least-squares-TF-start} (T/F?) If $\vec{x}$ is a solution of $A\vec{x}=\vec{y}$, then $\vec{x}$ is also a least-squares solution of $A\vec{x}=\vec{y}$.

\exatend (T/F?) The equation $A^TA\hat{x}=A^T\vec{y}$ always has a solution for $\hat{x}$.

\exatend (T/F?) The equation $A^TA\hat{x}=A^T\vec{y}$ has a \emph{unique} solution for $\hat{x}$.

\exatend \label{exercise:5.2least-squares-TF-end} (T/F?) If $A^TA\hat{x}=A^T\vec{y}$ has two solutions $\hat{x}=\vec{p}$ and $\hat{x}=\vec{q}$, then $A\vec{p}=A\vec{q}$.

\exatend Let $A = \begin{bmatrix}
1 & 2 & 3 \\
4 & 5 & 6 \\
7 & 8 & 9
\end{bmatrix}$ and $\vec{y} = (-3,2,1)$. It can be shown that $\hat{x}=(3,-2,-1/3)$ is a least-squares solution of $A\hat{x}=\vec{y}$. 
\begin{enumerate}[label=(\alph*)]
\item Find $\proj_{\col A}\vec{y}$. 
\item It turns out that $\hat{x}=(10/3, -8/3,0)$ is also a least-squares solution. Do you expect this to give a different $\proj_{\col A}\vec{y}$ than your answer to (a) or the same? Explain.
\end{enumerate}

% 2 cool problems below
\exatend Let $\vec{u}=\bbbvec{3}{2}{3}$ and $\vec{y}=\bbbvec{3}{1}{0}$.
\begin{enumerate}[label=(\alph*)]
\item Use the projection formula in \ref{theorem:orthogonal-projection-formula} to find $\proj_{\col \vec{u}}\vec{y}$.
\item Solve $\vec{u}^T\vec{u}x=\vec{u}^T\vec{y}$\, to find a least-squares solution to $\vec{u}x=\vec{y}$, and find $\proj_{\col \vec{u}}\vec{y}$. Note the similarities to part (a).
\end{enumerate}

\exatend Let $U = \begin{bmatrix}
1 & 0 \\
1 & 0 \\
0 & 1 \\
0 & 1
\end{bmatrix}$, and $\vec{y}=\begin{bmatrix}
1 \\ 2 \\ 3 \\ 4
\end{bmatrix}$. Find $\proj_{\col U}\vec{y}$ two ways:
\begin{enumerate}[label=(\alph*)]
\item Notice that the columns of $U$ are orthogonal. Use the projection formula in \ref{theorem:orthogonal-projection-formula}.
\item Solve $U^TU\hat{x}=U^T\vec{y}$\, to find a least-squares solution to $U\vec{x}=\vec{y}$, and multiply it by $U$. Note the similarities to part (a).
\end{enumerate}

\exatend Let $A=\begin{bmatrix}
-1 & 0 \\
2 & 1 \\
0 & 1 
\end{bmatrix}$ and $\vec{y}=\bbbvec{0}{-11}{7}$. Find the least-squares solution of $A\vec{x}=\vec{y}$.

\exatend Let $A=\begin{bmatrix}
-1 & 0 \\
2 & 1 \\
0 & 1 
\end{bmatrix}$ and $\vec{y}=\bbbvec{3}{-4}{2}$. 
\begin{enumerate}[label=(\alph*)]
\item Find $\hat{x}$, the least-squares solution of $A\vec{x}=\vec{y}$. 
\item How far is $A\hat{x}$ from $\vec{y}$? What does this tell you about $A\vec{x}=\vec{y}$?
\end{enumerate}

\exatend Let $A = \begin{bmatrix}
8 & 3 & 5 \\
1 & 2 & -1 \\
8 & 3 & 5
\end{bmatrix}$, and $\vec{y}=(14,0,12)$. Find the least-squares solution of $A\vec{x}=\vec{y}$ and $\proj_{\col A}\vec{y}$. %soln: (1,0,1)

\exatend Prove the (i \leftrightarrow ii) part of theorem \ref{theorem:unique-least-squares-solution} below. Here's a guide to the proof:
\begin{enumerate}[label=(\alph*)]
\item (\leftarrow) Assume $A^TA$ is invertible. To show that the columns of $A$ are independent, suppose $A\vec{z}=\vec{0}$, and show this implies $\vec{z}=\vec{0}$.
\item (\rightarrow) Assume the columns of $A$ are linearly independent. Explain why, to show that $A^TA$ is invertible, it suffices to show the columns of $A^TA$ are linearly independent. To show the columns of $A^TA$ are independent, suppose $A^TA\vec{z}=\vec{0}$, and show this implies $\vec{z}=\vec{0}$. \hintfootnote{One step is to show $\vec{z}^TA^TA\vec{z} = 0$; what does this tell you about $A\vec{x}$?}
\end{enumerate}

\exatend Prove the (ii \leftrightarrow iii) part of theorem \ref{theorem:unique-least-squares-solution} below.

\begin{theorem}[Uniqueness of least-squares solutions.]\label{theorem:unique-least-squares-solution}

Suppose $A$ is an $r\times c$ matrix. The following are equivalent:
\begin{enumerate}[label=(\roman*)]
\item The columns of $A$ are linearly independent.
\item $A^TA$ is invertible.
\item For every $\vec{y}$ in $\R^r$, the equation $A\vec{x}=\vec{y}$ has a \emph{unique} least-squares solution $\hat{x}$.
\end{enumerate}
\end{theorem}

\exatend If the columns of $A$ are linearly independent, the matrix $( A^TA )^{-1}A^T$ is called the \emph{pseudo-inverse} of $A$. What does it have in common with actual inverses and how is it different? (This question is more open-ended than usual; there are many correct answers.)
%\end{savenotes} 

\columnbreak
\section{Curve Fitting}\label{sec:curve-fitting}

\exercise (\emph{Warm-up}) Suppose the size of $A$ is $3\times 5$ and $\vec{y}$ is in $\R^3$. Determine the size of each product, if it's defined:

\begin{horiparts}{4}
\item $AA^T$
\item $A^TA$
\item $A\vec{y}$
\item $A^T\vec{y}$
\end{horiparts}

\exercise (\emph{Warm-up}) Suppose $A$ is $r\times c$ and $\vec{y}$ is in $\R^r$. Does $A^TA\hat{x}=A^T\vec{y}$ always have a solution for $\hat{x}$?

A big part of any kind of science, from astronomy to biology to economics, is analyzing data, which is what statistics is all about. And a big part of data analysis is coming up with a \emph{model}, which is a way of summarizing the data in a simpler way. Some examples of models:
\begin{itemize}
\item A \emph{least-squares regression} model, which is an equation of a line (or some other curve) that gets as close as possible to a bunch of data points. For a single-variable linear model, you can describe it with only two \emph{parameters}: the slope and intercept. This is the focus of this section.
\item A \emph{weather forecasting} model, which is a mathematical representation of the atmosphere. This kind of model may store millions of parameters, for the temperature, velocity, and humidity of millions of sections of air, but it is still much simpler than the atmosphere itself, with, like, forty gazillion molecules.
\item A \emph{large language} model, which contains billions of parameters that represent the likelihood of certain phrases coming next in a piece of text, depending on all the phrases that came before. The data this kind of model is summarizing the vast amount of text on the internet that it was trained on.
\end{itemize}

All of these models can \emph{describe} the data they summarize, and also try to \emph{predict} new data. Either way, the model should get as close as possible to the data; the parameters should be chosen so that the model \emph{best fits} the data. Roughly speaking, we want the \emph{\textcolor{red}{ERROR}} to be as small as possible:
\[\textit{\textcolor{red}{ERROR }} = \textit{ DATA } - \textit{\textcolor{blue}{ MODEL}}\]
To get more specific, consider a list of data points $(x_1,y_1),(x_2,y_2),\ldots,(x_c,y_c)$. Suppose we want a model to describe/predict $y$ based on values of $x$: $\textcolor{blue}{\hat{y}=b+mx}$. So, we want $\textcolor{blue}{b,m}$ that makes all of the different errors $\textcolor{red}{y_{i}-(b+mx_{i})}$, for each $i$ between $1$ and $c$, as small as possible.

\begin{tikzpicture}
\begin{axis}[
x=32pt,y=15pt,
axis lines=middle,
axis line style={-},
domain=0:12,
xmin=0,xmax=11,
ymin=0,ymax=10,
samples=200,
xlabel=$x$,
ylabel=$y$,
ticks=none,
xticklabels={},
yticklabels={},]
%grid=major,]
\pgfplotstableread{
x y
1 1
1.5 2
2 1
2.7 1.5
2.2 1
3 3
3.3 4
3.7 4.5
3.9 7
4 3
4.1 2.5
4.4 8
4.8 5
5 7
5.2 6
6 8.5
8 6.5
}\datatable
\addplot[blue, thick] {1.167*x-0.2429};
\node[blue,rotate=29] at (axis cs:7,8.5) {$\hat{y}=b+mx$};
\node[right] at (axis cs:8,6.5) {$(x_{17},y_{17})$};
\node[right,red] at (axis cs:8,7.9) {$y_{17}-(b+mx_{17})$};
\node[left] at (axis cs:3.9,7) {$(x_{9},y_{9})$};
\node[left,red] at (axis cs:3.9,5.8) {$y_{9}-(b+mx_{9})$};
%\addlegendentry{$y=\beta_0+\beta_1\sin(\beta_2x)$}
\addplot[only marks, mark=*] table {\datatable};
\pgfplotstablegetrowsof{\datatable}
\pgfmathtruncatemacro{\rows}{\pgfplotsretval-1}
\pgfplotsinvokeforeach{0,...,\rows}{
\pgfplotstablegetelem{#1}{x}\of{\datatable}
\edef\xval{\pgfplotsretval}
\pgfplotstablegetelem{#1}{y}\of{\datatable}
\edef\yval{\pgfplotsretval}
\edef\temp{
\noexpand\draw[thin, red] (axis cs:\xval, \yval) -- (axis cs:\xval, {1.167*\xval-0.2429});
}
\temp
}
\end{axis}
\end{tikzpicture}

Fortunately, we have a good way describe how big/small a big list of numbers is--the length of a vector! The goal is to minimize:
\[\nm{\textcolor{red}{\bbbvec{y_1-(b+mx_1)}{\vdots}{y_c-(b+mx_c)}}} = \nm{\bbbvec{y_1}{\vdots}{y_c} - \textcolor{blue}{\bbbvec{b+mx_1}{\vdots}{b+mx_c}}}\]
We can use what we know about matrix-vector multiplication to help with this if we denote:
\[ \vec{y}=\bbbvec{y_1}{\vdots}{y_c}, \quad A = \begin{bmatrix}
1 & x_1 \\
1 & x_2 \\
\vdots & \vdots \\
1 & x_c
\end{bmatrix}, \quad \hat{\beta} =\bbvec{b}{m}\] 
Here $A$ is called the \textbf{design matrix}. So, we want to find $\hat{\beta}$ that minimizes the error:
\[\nm{\vec{y}-\textcolor{blue}{A\hat{\beta}}}\]
In other words, to make $\hat{y} = b + mx$ best fit the data, we want to find $\hat{\beta}$, the \emph{least-squares solution} of $A\vec{\beta}=\vec{y}$.
Theorem \ref{theorem:least-squares} tells us $\hat{\beta}$ is a solution of:
\begin{equation}\label{eq:regression}\boxed{A^TA\hat{\beta}=A^T\vec{y}}\end{equation}

\bbtwocol{0.7}{0.3}{
\example \label{example:first-best-fit} Use equation \ref{eq:regression} to find the best-fit line for the data points $(0,5),(1,3)$, and $(2,1)$. Note that the points are collinear, so it should be easy to check the answer.
}{\begin{center}
\begin{tikzpicture}
\begin{axis}[
x=18pt,y=15pt,
axis lines=middle,
axis line style={-},
ymajorgrids=true,
%minor y tick num=3,
%yminorgrids=true,
xmajorgrids=true,
%minor x tick num=3,
%xminorgrids=true,
xlabel=$x$,
ylabel=$y$,
xmin=-0.5, xmax=3.5,
ymin=-0.5, ymax=6.5,
xtick={-2,-1,...,4},
ytick={-4,-3,...,6},]
%xticklabels={},
%yticklabels={},
\pgfplotstableread[row sep=\\]{
x y\\
0.0 5.0\\
1.0 3.0\\
2.0 1.0\\
}\datatable
\addplot[only marks, mark=*] table {\datatable};
\end{axis}
\end{tikzpicture}
\end{center}
}

\begin{book}
\textbf{Solution.} Using the notation above equation \ref{eq:regression}, denote:
\[A=\begin{bmatrix}
1 & x_1 \\
1 & x_2 \\
1 & x_3
\end{bmatrix} = \begin{bmatrix}
1 & 0 \\
1 & 1 \\
1 & 2 
\end{bmatrix}, \quad \vec{y}=\bbbvec{y_1}{y_2}{y_3}= \bbbvec{5}{3}{1}\]
We want to solve equation \ref{eq:regression} for $\vec{\beta}=(b,m)$:
\begin{align*}
\begin{bmatrix}
1 & 1 & 1 \\
0 & 1 & 2
\end{bmatrix}\begin{bmatrix}
1 & 0 \\
1 & 1 \\
1 & 2
\end{bmatrix}\bbvec{b}{m} &= \begin{bmatrix}
1 & 1 & 1 \\
0 & 1 & 2
\end{bmatrix}\bbbvec{5}{3}{1}
\begin{bmatrix}
3 & 3 \\
3 & 5
\end{bmatrix} \bbvec{b}{m} = \bbvec{9}{5}
\end{align*}
We can write this as an augmented matrix and row-reduce:
\[\begin{bmatrix}[cc|c]
3 & 3 & 9 \\
3 & 5 & 5
\end{bmatrix} \sim \begin{bmatrix}[cc|c]
1 & 0 & 5 \\
0 & 1 & -2
\end{bmatrix}\]
This gives the unique solution $b=5,m=-2$, so the least-squares regression line is $\hat{y}=5-2x$, which is indeed the equation of the line that includes the data points. \blacksquare
\end{book}

The next example shows that we can use the same method for finding best-fit functions for multi-variable data.
\example \label{example:best-fit-plane} Use equation \ref{eq:regression} To find the best-fit linear function of two variables, of the form $\hat{z} =b + mx + ny$ for some constants $b,m,n$, for the following data points. Also find the \emph{least-squares error}, $\nm{\vec{z}-\hat{z}}$.
\begin{align*}
(x_1,y_1,z_1) &= (-2,-2,-4) \\
(x_2,y_2,z_2) &= (-2,2,2) \\
(x_3,y_3,z_3) &= (2,2,8) \\
(x_4,y_4,z_4) &= (2,-2,6)
\end{align*}

\begin{book}
\textbf{Solution.} Our goal is to get as close as possible to $\vec{z}=(z_1,z_2,z_3,z_4) = (-4,2,8,6)$ with a vector of the form:
\[\hat{z}= \begin{bmatrix}
b + mx_1 + ny_1 \\
b + mx_2 + ny_2 \\
b + mx_3 + ny_3 \\
b + mx_4 + ny_4 \\
\end{bmatrix} = \begin{bmatrix}
1 & x_1 & y_1 \\
1 & x_2 & y_2 \\
1 & x_3 & y_3 \\
1 & x_4 & y_4 \\
\end{bmatrix}\begin{bmatrix}
b \\ m \\ n
\end{bmatrix}\]
This describes the design matrix $A$, plugging in the given $x$ and $y$ coordinates:
\[A=\begin{bmatrix}
1 & -2 & -2 \\
1 & -2 & 2 \\
1 & 2 & 2 \\
1 & 2 & -2 \\
\end{bmatrix}\]
Now, we need to solve equation \ref{eq:regression}:
\begin{align*}
\begin{bmatrix}
1 & 1 & 1 & 1 \\
-2 & -2 & 2 & 2 \\
-2 & 2 & 2 & -2
\end{bmatrix}\begin{bmatrix}
1 & -2 & -2 \\
1 & -2 & 2 \\
1 & 2 & 2 \\
1 & 2 & -2 \\
\end{bmatrix}\bbbvec{b}{m}{n} &= \begin{bmatrix}
1 & 1 & 1 & 1 \\
-2 & -2 & 2 & 2 \\
-2 & 2 & 2 & -2
\end{bmatrix}\begin{bmatrix}
-4 \\ 2 \\ 8 \\ 6
\end{bmatrix} \\
\begin{bmatrix}
4 & 0 & 0 \\
0 & 16 & 0 \\
0 & 0 & 16
\end{bmatrix}\bbbvec{b}{m}{n} &= \bbbvec{12}{32}{16}
\end{align*}
This gives $b=3,m=2,n=1$, which means the best-fit linear function is $\hat{z}=3+2x+y$. This model gives a list of $z$-coordinates:
\[\hat{z}=A\hat{\beta}=\begin{bmatrix}
1 & -2 & -2 \\
1 & -2 & 2 \\
1 & 2 & 2 \\
1 & 2 & -2 \\
\end{bmatrix}\bbbvec{3}{2}{1} = \begin{bmatrix}
-3 \\ 1 \\ 9 \\ 5
\end{bmatrix}\]
So the least-squares error is:
\[\nm{\vec{z}-\hat{z}} = \sqrt{(-4-(-3))^2+(2-1)^2+(8-9)^2+(6-5)^2} = \sqrt{4}=2\]
 \blacksquare
\end{book}

In fact, the method of equation \ref{eq:regression} works for finding many kinds (but not all kinds) of best-fit functions.

\example \label{example:design-matrix} Describe the design matrix you would use to find the best-fit function of the form $\hat{y}=a\cos(x) + b\sin(x)$ for the data points $(x_1,y_1), (x_2,y_2), \ldots, (x_c,y_c)$.

\begin{book}
\textbf{Solution.} The goal is to find $a,b$, to get as close as possible to $\vec{y}=(y_1,y_2,\ldots,y_c)$, with a vector of the form:
\[\hat{y}=\begin{bmatrix}
a\cos(x_1)+b\sin(x_1) \\
a\cos(x_2)+b\sin(x_2) \\
\vdots\\
a\cos(x_c)+b\sin(x_c) 
\end{bmatrix} = 
\begin{bmatrix}
\cos(x_1) & \sin(x_1) \\
\cos(x_2) & \sin(x_2) \\
\vdots & \vdots \\
\cos(x_c) & \sin(x_c)
\end{bmatrix}\bbvec{a}{b}\]
So, by plugging data values into the design matrix
\[A=\begin{bmatrix}
\cos(x_1) & \sin(x_1) \\
\cos(x_2) & \sin(x_2) \\
\vdots & \vdots \\
\cos(x_c) & \sin(x_c)
\end{bmatrix}\]
we could solve $A^TA\hat{\beta}=A^T\vec{y}$ for $\hat{\beta}=(a,b)$ to find the best-fit function. \blacksquare
\end{book}

The above methods are pretty flexible and powerful, but they don't cover every kind of curve fitting problem. For example, if you want to find a best-fit function of the form $\hat{y}=\cos(ax+b)$, there is no way to construct a design matrix that works:
\[\bbbvec{\cos(ax_1+b)}{\vdots}{\cos(ax_c+b)} = \begin{bmatrix}
?? & ?? \\
\vdots & \vdots \\
?? & ??
\end{bmatrix}\bbvec{a}{b}\]
This kind of \emph{nonlinear regression} problem is beyond the scope of this course, but you might learn about algorithms for it in a numerical methods class. 

\exerciseheading

% How to use projector() from mat.py:
% projector([-1,0,1,2],[x,1],Matrix([-3,1])).rref()
% [-1,0,1,2] is the list of x-values, 
% [x,1] is the list of terms in the function to fit, eg y = b1*x + b0*1,
% Matrix([-3,1]) is the list of corresponding b1, b2 in the desired solution.
% returns an augmented matrix whose solutions are y1,y2,y3,y1.

\exatend Suppose you have data points $(x_1,y_1),\ldots, (x_{376},y_{376})$, and denote $\vec{x}=(x_1,\ldots,x_{376}), \vec{y}=(y_1,\ldots,y_{376})$. You want to find the best fit line $\hat{y}=b+mx$ for the data. Fill in the blanks with numbers: The design matrix $A$ is $\rule{2em}{0.8pt} \times \rule{2em}{0.8pt}$. The column space of $A$ is a \rule{2em}{0.8pt}-dimensional subspace of $\R^{\rule{2em}{0.8pt}}$.

\exatend Find the equation in the form $\hat{y} = b+ mx$ that best fits the $(x,y)$ data points $(-1,2), (0,4),(1,-2),(2,-6)$. %answer: y=-3x+1 

\exatend Find the equation in the form $\hat{y} =b+ mx$ that best fits the $(x,y)$ data points $(2,3), (4,2),(6,7)$. %answer: y=x

\exatend Find the equation in the form $\hat{y} = b + mx$ that best fits the $(x,y)$ data points $(2,7),(4,7),(6,4),(8,0)$. Also find the least-squares error $\nm{\vec{y}-\hat{y}}$.

\exatend (T/F?) If $\hat{y}=2x+3$ is the best-fit line for data points $(x_1,y_1),\ldots,(x_c,y_c)$, then $\sum(y_i-(2x_i+3))^2 \le \sum(y_i-(1.9x_i+2.9))^2$.

\exatend (T/F?) The methods in this section can be used to find a best-fit function of the form $\hat{y}=e^{ax} + \sin(x+b)$.

\exatend Describe the design matrix you would use to find the best-fit function of the form $\hat{y}=a_0+a_1x+a_2x^2+a_3x^3$ for data points $(x_1,y_1), (x_2,y_2), \ldots, (x_c,y_c)$. (See example \ref{example:design-matrix}.)

\exatend Describe the design matrix you would use to find the best-fit function of the form $\hat{z}=ae^{x+y}-bx+cy^2$ for data points $(x_1,y_1,z_1),(x_2,y_2,z_2),\ldots,(x_c,y_c,z_c)$. (See example \ref{example:design-matrix}.)

\exatend Describe the design matrix you would use to find the best-fit function of the form $\hat{z}=a_0x+a_1x^2+a_2y+a_3y^2 +a_4xy$ for data points $(x_1,y_1,z_1),(x_2,y_2,z_2),\ldots,(x_c,y_c,z_c)$. (See example \ref{example:design-matrix}.)

\exatend Suppose you are creating a robot outfielder to play robot baseball. When a fly ball is hit, a system of cameras estimates the ball's height at several times, giving approximate data in the form $(t_1,h_1),(t_2,h_2),\ldots,(t_c,h_c)$. The robot must find want to find $h_0,v_0,$ and $a$ for a best-fit function of the form $\hat{h}=h_0 + v_0t-\frac{1}{2}at^2$. Describe the design matrix needed to do so.

\exatend (\emph{Challenge}) Suppose you are still working on your robot outfielder. You realize that the ball is always hit from an initial height of $h_0=1$ meter, and gravity has a constant acceleration $a=9.8 \frac{\text{m}}{\text{s}^2}$, so you should only be estimating the initial velocity $v_0$. Do the methods in this section work to find a best-fit function of the form $\hat{h}=1+v_0t-\frac{9.8}{2}t^2$? Is there a way to modify the methods to make it work?
% solution: the problem is that the standard methods will give a solution (h0,v0,a) that might not have h0 = 1 and a=9.8. And forcing them to be so with extra equations (rows) in the A.T*A*beta=A.T*h equation would make it inconsistent. The trick is to make the goal function y=v0*t, with data points of the form (ti, yi) with yi=hi-1-4.9ti^2.


\exatend Cesium-137 has a half-life of 30.04 years and strontium-90 has a half-life of 28.91 years. Both are harmful components of nuclear waste. If you had a sample of nuclear waste (where these only radioactive isotopes), and measured its mass at several times, you'd have data points of the form $(t_1,m_1), (t_2,m_2), \ldots, (t_c,m_c)$. To estimate the initial amounts $s$ and $c$ of the two isotopes, you could find the best-fit function of the form $\hat{m}=b+c(0.5)^{t/30.04}+s(0.5)^{t/28.91}$. Describe the design matrix you'd need. (Note: If you could take perfectly precise measurements, you'd be able to find the exact values of $s$ and $c$ with three data points. But in real life, measurements are not perfectly precise and small errors can make a big difference in situations like this.)

\exatend This question refers to the four data points in example \ref{example:best-fit-plane}. One naive way to find a plane that gets close to the four points might be the plane that goes through three of the points, and hope it gets close enough to the fourth point. It turns out the plane that goes through $(x_1,y_1,z_1), (x_2,y_2,z_2),$ and $(x_3,y_3,z_3)$ has equation $z=2+1.5x+1.5y$. If $\vec{\gamma} = (2,1.5,1.5)$, find the error for this plane, $\nm{\vec{y}-A\vec{\gamma}}$. Which plane has smaller error: $\hat{z}=3+2x+y$ or $z=2+1.5x+1.5y$? Is this what you expected?

\exatend Suppose you want to find a best-fit function $\hat{y}=b+mx$ for the $(x,y)$ data points $(5,1), (5,3)$. Show there isn't a unique best-fit function by finding at least two that have the same least-squares error $\nm{\vec{y}-\hat{y}}$.

\exatend What must be true about data points $(x_1,y_1),\ldots,(x_c,y_c)$ to guarantee that the best-fit line $\hat{y}=b+mx$ is \emph{unique}? Explain.\hintfootnote{See theorem \ref{theorem:unique-least-squares-solution}.}

\exatend Find the best-fit function of the form $\hat{y}=c+bx+ax^2$ for the $(x,y)$ data points $(-2,7),(-1,-2),(1,2),$ and $(2,5)$. Also find the least-squares error $\nm{\vec{y}-\hat{y}}$ for your function. %solution: y=-2+2*x^2

\exatend Suppose you're finding a best-fit line $\hat{y}=b+mx$ for data points $(x_1,y_1),\ldots,(x_c,y_c)$. 
\begin{enumerate}[label=(\alph*)]
\item If you calculate the mean $\bar{x}$ of the data's $x$-values, and write a new variable with values $x^{*}_i=x_i-\bar{x}$, we say $x^{*}$ is written in \emph{mean-deviation form}, and we'll see how this saves some computation. Write the data in example \ref{example:first-best-fit} in mean-deviation form and find the best-fit line for it.
\item It can be shown that $\sum(x_i-\bar{x})=0$ for any list of numbers $x_1,\ldots,x_c$. Use this to show that if $x^{*}$ is in mean-deviation form, with corresponding design matrix $A$, then $A^TA$ is a diagonal matrix.
\end{enumerate}


\chapter{Symmetric Matrices and Quadratic Forms}
\label{chap:symmetric-matrices}

TODO: Write this chapter

\chapter{Abstract Vector Spaces}
\label{chap:abstract-vector-spaces}

Many of the concepts of linear algebra apply to mathematical objects other than $\R^c$ vectors. For example, polynomial addition and scalar multiplication works the same way:
\[5(z^2-z+3) + (2z-1) = 5z^2+3z+14\]
looks a lot like:  
\[5 \bbbvec{1}{-1}{3} + \bbbvec{0}{2}{-1}=\bbbvec{5}{3}{14}\]
Another example is that derivatives and integrals have a linearity property, which means we can think of them as \emph{linear transformations}:
\begin{align*}
\frac{d}{dz}( kf(z) + cg(z) ) &= k \frac{d}{dz}f(z) + c \frac{d}{dz}g(z) \\
\int ( kf(z)+cg(z) )dz &= k\int f(z)dz + c \int g(z)dz \\
\end{align*}
These look a lot like:
\[T(k\vec{x}+c\vec{z}) = kT(\vec{x}) + cT(\vec{z})\]
When we call objects like polynomials and other functions \emph{vectors}, we can bring the powerful theorems of linear algebra to bear when studying them. But of course, it's not just \emph{calling} objects vectors, it's the fact that the objects have certain properties in common with $\R^c$ vectors that allows us to apply the theorems. First we'll see what those properties are, and later we'll see exactly how the theorems of linear algebra apply to objects like polynomials.
\section{Vector Spaces other than $\R^c$}
\label{sec:vector-spaces-other-than-Rc}

\exercise (\emph{Warm-up}) Let $\mathbb{P}_2$ be the set of all functions of the form $f(z) = az^2 + bz + c$, for any real numbers $a, b$, and $c$. (We'll use $z$ as the variable because $x$ and $t$ are already used for other things.) So $\mathbb{P}_2$ contains functions like $3z^2 - 17.5z + \pi$, or $2z$, or just $3$. Suppose that $f(z) = az^2 + bz + c$ is in $\mathbb{P}_2$.
\begin{enumerate}[label=(\alph*)]
  \item If $k$ is a real number, is $kf(z) = kaz^2 + kbz + kc$ in $\mathbb{P}_2$?
  \item If $g(z) = pz^2 + qz + r$ is in $\mathbb{P}_2$, is $f(z) + g(z)$ in $\mathbb{P}_2$?
  \item Is $f(z) + g(z) = g(z) + f(z)$?
  \item Is $f(z) + 0 = f(z)$?
\end{enumerate}

\begin{notes}
\vfill
\end{notes}

\begin{definition}\label{definition:vector-space}
A set of objects $V$ is called a \textbf{vector space}, and the objects are called \textbf{vectors}, if it comes with addition and scalar multiplication operations for which the following properties hold:
\begin{enumerate}[label=(\roman*)]
\item \textbf{$\vec{\mathbf{0}}$ exists and works:} There exists an object, called $\vec{0}$, in $V$ such that $\vec{0}+\vec{u}=\vec{u}$ for all $\vec{u}$ in $V$.
\item \textbf{Closure under vector addition:} $\vec{u}+\vec{v}$ is in $V$ for all $\vec{u},\vec{v}$ in $V$.
\item \textbf{Closure under scalar multiplication:} $r\vec{u}$ is in $V$ for all $r$ in $\R$ and $\vec{u}$ in $V$.
\item \textbf{Vector addition is commutative:} $\vec{u}+\vec{v}=\vec{v}+\vec{u}$, for all $\vec{u},\vec{v}$ in $V$.
\item \textbf{Vector addition is associative:} $(\vec{u}+\vec{v}) + \vec{w}= \vec{u}+(\vec{v}+\vec{w})$  for all $\vec{u},\vec{v},\vec{w}$ in $V$.
\item \textbf{Scalar multiplication is associative:} $r(t\vec{u}) = (rt)\vec{u}$ for all $r,t$ in $\R$ and all $\vec{u}$ in $V$.
\item \textbf{Additive inverses exist:} For every $\vec{u}$ in $V$, there exists a vector $\vec{v}$ such that $\vec{u}+\vec{v}=\vec{0}$.
\item \textbf{1 works:} $1\vec{u}=\vec{u}$ for all $\vec{u}$ in $V$.
\item \textbf{Distributive rules work:} $r(\vec{u}+\vec{v}) = r\vec{u}+r\vec{v}$ and $(r+t)\vec{u}=r\vec{u}+t\vec{v}$ for all $r,t$ in $\R$ and all $\vec{u},\vec{v}$ in $V$.
\end{enumerate}
\end{definition}

\example $\R^c$ is a vector space, for any whole number $c$. Each of the properties of the definition is clearly familiar in $\R^c$. \blacksquare

\begin{notes}
\vfill \hspace{1pt}
\columnbreak
\end{notes}

\example We denote $\mathbb{P}$ to be the set of all polynomials, which are functions that can be written in the form $\vec{p}(z) = a_0 + a_1z + a_2z^2 + \cdots + a_nz^n$ for some non-negative integer $n$ and real numbers $a_0,\ldots,a_n$. Is $\mathbb{P}$ a vector space?

\begin{notes}
\vfill
\vfill
\end{notes}

\begin{book}
\textbf{Solution.} Let's check each of the properties of the definition.
\begin{enumerate}[label=(\roman*)]
\item Does there exist a polynomial $\vec{0}$ such that $\vec{0}+\vec{p}=\vec{p}$ for all polynomials $\vec{p}$? Yes, the zero polynomial is the function whose output is zero for all $z$: $\vec{0}(z)=0$.
\item Is $\vec{p}(z)+\vec{q}(z)$ a polynomial whenever $\vec{p}$ and $\vec{q}$ are? Yes, so the polynomials are closed under addition. For example, $3z^2+5 + (-\pi z^2+z) = (3-\pi)z^2+z+5$ is a polynomial.
\item Is $r\vec{p}(z)$ a polynomial whenever $r$ is a real number and $\vec{p}$ is a polynomial? Yes, so the polynomials are closed under scalar multiplication. For example, $\sqrt{2}(z^{3}-z^2) = \sqrt{2}z^3-\sqrt{2}z^2$ is a polynomial.
\item Is $\vec{p}(z)+\vec{q}(z) = \vec{q}(z)+\vec{p}(z)$? Yes, so polynomial addition is commutative. For example, $(3z^2+5) + (-\pi z^2+z) = (-\pi z^{2} +z) + (3z^2+5)$.
\item Is $( \vec{p}(z) + \vec{q}(z) ) + \vec{f}(z) = \vec{p}(z) + ( \vec{q}(z) + \vec{f}(z) )$? Yes, polynomial addition is associative. For example $(3z^2+(1-z^3)) + z = 3z^2 + ((1-z^3) + z)$.
\item Is $r(t\vec{p}(z)) =(rt)\vec{p}(z)$ for all real numbers $r,t$ and polynomials $\vec{p}(z)$? Yes, scalar multiplication is associative. For example, $(3\pi)(z^2+5) = 3(\pi(z^2+5))$.
\item For every polynomial $\vec{p}(z)$, does there exist a polynomial $\vec{v}(z)$ so that $\vec{p}(z)+\vec{v}(z) = 0$? Yes. For example, the additive inverse of $z^2+5$ is $-z^2-5$.
\item Is $1\vec{p}(z)=\vec{p}(z)$ for all polynomials $\vec{p}(z)$? Yes.
\item Is $r(\vec{p}(z)+\vec{q}(z)) = r\vec{p}(z) + t\vec{p}(z)$? Yes. Is $(r+t)\vec{p}(z)=r\vec{p}(z)+t\vec{p}(z)$? Yup.
\end{enumerate}
Thus $\mathbb{P}$ is a vector space. \blacksquare
\end{book}

\example Are the integers, $\mathbb{Z} = \{\ldots,-3,-2,-1,0,1,2,3,\ldots\}$, a vector space?

\begin{notes}
\vfill \hspace{1pt}
\columnbreak
\end{notes}

\begin{book}
\textbf{Solution.} Let's check each of the properties of the definition.
\begin{enumerate}[label=(\roman*)]
\item Does there exist an integer $\vec{0}$, so that $\vec{0}+\vec{u}=\vec{u}$ for all integers $\vec{u}$? Yes, it's just the number $0$.
\item Is $\vec{u}+\vec{v}$ an integer, whenever $\vec{u},\vec{v}$ are integers? Yes, so the integers are closed under addition.
\item Is $r\vec{u}$ an integer, whenever $r$ is a real number and $\vec{u}$ is an integer? No; for example, if $r=4.8$ and $\vec{u}=1$, we have $r\vec{u}=4.8$ which is not an integer. So $\mathbb{Z}$ is not a vector space. \blacksquare
\end{enumerate}
\end{book}

The following properties are familiar from $\R^c$, but it takes a little work to show that they are true for other vector spaces.

\begin{theorem}[Basic vector space properties.]\label{theorem:vector-space-properties}

If $V$ is a vector space, then:

\begin{enumerate}[label=(\roman*)]
\item There is only one vector $\vec{0}$ that satisfies property (i) of the definition.
\item $0\vec{u}=\vec{0}$ for all $\vec{u}$ in $V$.
\item $(-1)\vec{u}$ is the additive inverse of $\vec{u}$ for each $\vec{u}$ in $V$.
\item $r\vec{0}=\vec{0}$ for all $r$ in $\R$.
\end{enumerate}
\end{theorem}
Let's prove (i) and (ii) and leave the others for exercises.

\begin{notes}
\vfill
\end{notes}

\begin{book}
\textbf{Proof.} (i) Suppose that $\vec{0}_1$ and $\vec{0}_2$ are two vectors in a vector space $V$ that satisfy property (i) of the definition. That is, suppose that, for all $\vec{u}$ in $V$:
\[\vec{0}_1+\vec{u}=\vec{u} \text{ and } \vec{0}_2+\vec{u}=\vec{u}\]
Then, since $V$ is a vector space we can use the properties of the definition:
\begin{align*}
\vec{0}_1 &= \vec{0}_2+\vec{0}_1 &\text{(property (i), with } \vec{u}=\vec{0}_2)\\
&= \vec{0}_1 + \vec{0}_2 &(\text{vector addition is commutative}) \\
&= \vec{0}_2 &(\text{property (i), with } \vec{u}=\vec{0}_1)
\end{align*}
Thus if $\vec{0}_1$ and $\vec{0}_2$ each satisfy property (i) of the definition, then $\vec{0}_1=\vec{0}_2$.

(ii) Let $\vec{u}$ be in $V$. Then:
\begin{align*}
0\vec{u} &= (0+0)\vec{u} & (\text{since } 0=0+0)
\end{align*}
By property (vii), there exists an additive inverse $\vec{v}$ of $0\vec{u}$, such that $0\vec{u}+\vec{v}=\vec{0}$. Let's add it to both sides of the equation:
\begin{align*}
0\vec{u} + \vec{v} &= 0\vec{u}+0\vec{u} + \vec{v} & \\
\vec{0} &= 0\vec{u} + \vec{0} & (\text{defn of additive inverse}) \\
&= \vec{0}+0\vec{u} & (\text{vector addition is commutative}) \\
&= 0\vec{u} & (\text{definition of }\vec{0})
\end{align*}
\blacksquare
\end{book}

\example The following can each be shown to be vector spaces: 
\begin{itemize}
\item The set of all $\R\rightarrow\R$ functions is a vector space.
\item The set of all $\R\rightarrow\R$ \emph{differentiable} functions is a vector space.
\item Given any homogeneous differential equation, the set of solutions to it is a vector space.
\item The set of all $3\times 2$ matrices is a vector space.
\item The set of all $\R^2\rightarrow\R^3$ linear transformations is a vector space. (This is equivalent to the previous one).
\end{itemize}

\begin{notes}
\columnbreak
\end{notes}

Many definitions we learned in chapter \ref{chap:existence-and-uniqueness-of-solutions} can be generalized for vector spaces. Let's include them here for completeness.

\begin{definition} If $\vec{u}_1,\vec{u}_2,\ldots, \vec{u}_c$ is a list of vectors in a vector space $V$, and $x_1,x_2,\ldots,x_c$ are in $\R$, then $x_1\vec{u}_1+x_2\vec{u}_2+\cdots+x_c\vec{u}_c$ is called a \textbf{linear combination} of $\vec{u}_1,\vec{u}_2,\ldots, \vec{u}_c$. The set of all linear combinations of $\vec{u}_1,\vec{u}_2,\ldots, \vec{u}_c$ is called $\spn ( \vec{u}_1,\vec{u}_2,\ldots,\vec{u}_c )$.
\end{definition}

\example In the vector space of all $\R\rightarrow\R$ functions, describe $\spn(\cos(3z),\sin(2z))$.

\begin{notes}
\vfill
\end{notes}

\begin{book}
\textbf{Solution.} The span includes functions like:
\begin{itemize}
\item $7\cos(3z) - 5\sin(2z)$
\item $0\cos(3z) + \sqrt{2}\sin(2z) = \sqrt{2}\sin(2z)$
\item $0\cos(3z) + 0\sin(2z) = \vec{0}$
\end{itemize}
We could write $\spn ( \cos(3z),\sin(2z) )$ in set-builder notation: 
\[\left\{ r\cos(3z) + t\sin(2z) : r,t \text{ in } \R\right\}\]
Note that the scalar multiples aren't allowed to be functions of $z$. \blacksquare
\end{book}

\begin{definition}
A linear system of equations that can be written in the form 
\[ x_1 \vec{u}_1 + \cdots + x_c\vec{u}_c = \vec{0}\]
is called \textbf{homogeneous}. Homogeneous systems are always consistent since $\vec{x} = ( x_1, \ldots, x_c )=\vec{0}$ is always a solution, called the \textbf{trivial solution.} A \textbf{non-trivial} solution is any solution $\vec{x} \ne \vec{0}$. 

When a non-trivial solution exists, we say the list $\vec{u}_1, \ldots, \vec{u}_c$ is \textbf{linearly dependent}. When the only solution is the trivial solution, we say the list is \textbf{linearly independent.} The empty list $()$ is defined to be linearly independent as well.
\end{definition}

The next theorem may be so intuitive that it seems obvious, but the proof is a little tricky.
\begin{notes}
You can find it in the textbook.
\end{notes}

\begin{theorem}[Zero polynomial.]
\label{theorem:zero-polynomial} If the polynomial $\vec{p}(z) = a_0 + a_1z + a_2z^2 + \cdots + a_nz^n$ equals the $\vec{0}(z)$ function (that is, if $\vec{p}(z) = 0$ for all real numbers $z$), then all the coefficients are zero: $a_0=0, \ldots, a_n=0$.
\end{theorem}

\begin{book}
\textbf{Proof.} Suppose that $\vec{p}(z)=0$ for all $z$ in $\R$. Then, since 
\begin{align*}
\vec{p}(0) &= a_0+a_1(0) + a_2(0^2) + \cdots + a_n(0^n) = a_0 \\
\vec{p}(0) &= 0
\end{align*}
we have $a_0=0$. Now let's look at the first derivative of $\vec{p}$:
\begin{align*}
\vec{p}'(z) &= 0 + 1a_1 + 2a_2z + \cdots + na_nz^{n-1}  \\
\vec{p}'(0) &= 0 + a_1 + 2a_2(0) + \cdots + na_n(0)^{n-1}  \\
\vec{p}'(0) &= a_1
\end{align*}
Since $\vec{p}$ is a constant function, its derivative is zero (and so are its second, third, etc derivatives). Thus $a_1=0$. Now for the second derivative:
\begin{align*}
\vec{p}''(z) &= 0 + 0 + 2a_2 + 6a_3z\cdots + n(n-1)a_nz^{n-2}  \\
\vec{p}''(0) &= 0 + 0 + 2a_2 + 6a_3(0) \cdots + n(n-1)a_n(0)^{n-2}  \\
\vec{p}''(0) &= 2a_2
\end{align*}
But since $\vec{p}''(z)=0$, we have $a_2=0$. Continuing this process shows that $a_3=0,\ldots,a_n=0$.\blacksquare
% \textbf{Proof.} First let's note a simple fact. For any real numbers $x,y$:
% \begin{equation}
% \label{eq:triangle-inequality}
% |x+y| \le |x| + |y|
% \end{equation}
% This is because if $x$ and $y$ are the same sign or zero, then the left and right sides are equal; if they're opposite sign, then $|x+y|<\text{max}(|x|,|y|)<|x|+|y|$.

% Now let's prove the contrapositive of the theorem: If not all the coefficients are zero, then $\vec{p}(t)\ne 0$ for some $t$. Assume not all the coefficients are zero, and let $a_m$ be the furthest right (greatest $m$) nonzero coefficient. Let:
% \[r=\frac{|a_0|+|a_1|+\cdots+|a_{m-1}|}{|a_m|}+1\]
% We'll show that this value of $r$ makes the polynomial non-zero. Multiplying by $|a_m|$ gives:
% \begin{equation*}
% r|a_m| = |a_0|+|a_1|+\cdots+|a_{m-1}| + |a_m|
% \end{equation*}
% This implies, since $|a_m|>0$, that:
% \begin{equation}
% \label{eq:am}
% r|a_m| > |a_0|+|a_1|+\cdots+|a_{m-1}|
% \end{equation}
% Since $r\ge 1$, we have, for $j=0,1,\ldots,m-1$: 
% \begin{equation}
% \label{eq:rj}
% r^j\le r^{m-1}
% \end{equation}

% Plugging our value $r$ into the first $m-1$ terms, we get:
% \begin{align*}
% |a_0+a_1r+\cdots+a_{m-1}r^{m-1}| &\le |a_0r^{m-1} + a_1r^{m-1} + \cdots + a_{m-1}r^{m-1}| &(\text{by eqn }\ref{eq:rj})\\
% &=|a_0+a_1+\cdots+a_{m-1}||r^{m-1}| \\
% &\le(|a_0|+|a_1|+\cdots+|a_{m-1}|)|r^{m-1}| &(\text{by eqn } \ref{eq:triangle-inequality}) \\
% &<r|a_m||r^{m-1}| &(\text{by eqn } \ref{eq:am}) \\
% &=|a_mr^m| &(\text{since } r>0)
% \end{align*}
% Thus we have $|a_0+a_1r+\cdots+a_{m-1}r^{m-1}| \ne |a_mr^m|$ and thus
% \[a_0+a_1r+\cdots + a_{m-1}r^{m-1}+a_mr^m \ne 0\]
% \blacksquare
\end{book}

\begin{notes}
\columnbreak
\end{notes}


\example Is each list of polynomials linearly independent in $\mathbb{P}$?
\begin{enumerate}[label=(\alph*)]
\item $3+z^2,30+10z^2$
\item $1,z,z^2$
\item $1-2z^2,2-z-5z^2,3z+3z^2$
\end{enumerate}

\begin{book}
\textbf{Solution.} (a) The question is asking, is the only solution of the following equation $(x_1,x_2) = (0,0)$?
\[x_1(3+z^2) + x_2(30+10z^2) = \vec{0}\]
The answer is no: $x_1=-10,x_2=1$ is a non-trivial solution to this homogeneous equation, so the list is linearly dependent.

(b) Is the only solution of the homogeneous equation $x_1(1) + x_2z + x_3z^2 = \vec{0}$ the trivial solution $(x_1,x_2,x_3) = (0,0,0)$? You might be tempted to say no, that for example, $(0,-z,1)$ is a non-trivial solution. But this would be a mistake since $-z$ is not a scalar; a scalar must be a single real number, not a function of $z$. Or, you might be tempted to say that $(0,1,1)$ is a non-trivial solution because $z+z^2=0$ when $z=0$ or $z=-1$. But, the \emph{function} $z+z^2$ is not the same \emph{function} as $\vec{0}(z)=0$ so this $(0,1,1)$ is not a solution of the homogeneous equation.
One way to find solutions is to notice that $\vec{0}(z)= 0+0z+0z^{2}$, so we want to solve:
\[x_1(1)+x_2z+x_3z^2=0+0z+0z^2\]
And the only way to make these equations identical is to make $x_1=0,x_2=0,$ and $x_3=0$, so the list $1,z,z^2$ is linearly independent.

(c) Again we want to know if the only solution to the following homogeneous equation is the trivial solution $(x_1,x_2,x_3)=(0,0,0)$:
\[x_1(1-2z^2) + x_2(2-z-5z^2) + x_3(3z+3z^2) = 0 + 0z + 0z^2\]
Distributing and rearranging terms:
\[(x_1+2x_2) + (-x_2+3x_3)z + (-2x_1-5x_2+3x_3)=0 + 0z + 0z^2\]
And again we need to make the \emph{function} on the left be the same \emph{function} on the right, and that means making all three of the coefficients equal to zero:
\begin{align*}
x_1+2x_2 &= 0 \\
-x_2+3x_3&=0 \\
-2x_1-5x_2+3x_3&=0
\end{align*}
Writing this as an augmented matrix and row-reducing gives:
\[\begin{bmatrix}[ccc|c]
1 & 2 & 0 & 0 \\
0 & -1 & 3 & 0 \\
-2 & -5 & 3 & 0
\end{bmatrix} \sim \begin{bmatrix}[ccc|c]
1 & 0 & 6 & 0 \\
0 & 1 & -3 & 0\\
0 & 0 & 0 & 0
\end{bmatrix} \sim
\begin{matrix}
x_1+6x_3 = 0 \\
x_2-3x_3 = 0 \\
0 = 0
\end{matrix}
\]
Writing all the variables in terms of the free variable $x_3$:
\[\bbbvec{x_1}{x_2}{x_3} = \bbbvec{-6x_3}{3x_3}{x_3} = x_3 \bbbvec{-6}{3}{1}\]
So, some non-trivial solutions of the homogeneous equation are $(-6,3,1),(-12,6,2),$ etc. The list of polynomials is linearly dependent.
\blacksquare
\end{book}

\begin{notes}
\vfill
\end{notes}

Note that in the last example, there's a striking similarity between the equation with polynomials in $\mathbb{P}_2$
\[x_1(1-2z^2) + x_2(2-z-5z^2) + x_3(3z+3z^2) = 0 + 0z + 0z^2\]
and the $\R^3$ vector equation:
\[x_1 \bbbvec{1}{0}{-2} + x_2 \bbbvec{2}{-1}{-5} +x_3 \bbbvec{0}{3}{3} = \bbbvec{0}{0}{0} \]

\begin{notes}
\columnbreak
\end{notes}

The following theorems are analogous to theorems \ref{theorem:one-vector-is-independent}, \ref{theorem:zero-vector-is-dependent}, and \ref{theorem:characterization-of-linear-dependence}.

\begin{theorem}
\label{theorem:abstract-one-vector-is-independent}
If a list contains only a single vector and it's not $\vec{0}$, then the list is linearly independent.
\end{theorem}

\begin{theorem}
\label{theorem:abstract-zero-vector-is-dependent}
If a list of vectors contains $\vec{0}$, then it's linearly dependent.
\end{theorem}

\begin{theorem}[Characterization of linear dependence.]
\label{theorem:abstract-characterization-of-linear-dependence}

A list of vectors $W = \vec{a}_1, \ldots, \vec{a}_c $ is linearly dependent iff one of the vectors can be written as a linear combination of the others.

Specifically, if $W$ is linearly dependent, and $\vec{a}_1 \ne \vec{0}$, then some $\vec{a}_i$ with $i>1$ is a linear combination of the preceding vectors: \[x_1\vec{a}_1+\cdots+x_{i-1}\vec{a}_{i-1}=\vec{a}_i\]
\end{theorem}

The proof of theorem \ref{theorem:characterization-of-linear-dependence} did not rely on any special properties of $\R^n$; each of the steps only relies on the properties of the definition of a vector space. So it doubles as a proof of theorem \ref{theorem:abstract-characterization-of-linear-dependence}. In chapter 1, the proof of the first two theorems \ref{theorem:one-vector-is-independent} and \ref{theorem:zero-vector-is-dependent} were assigned as exercises; your answers probably likewise would stand in as proofs of theorems \ref{theorem:abstract-one-vector-is-independent}, \ref{theorem:abstract-zero-vector-is-dependent}.

\begin{definition}
A \textbf{subspace} $U$ of a vector space $V$ is a subset of $V$ for which:
\begin{enumerate}[label=(\roman*)]
\item The zero vector of $V$ is in $U$.
\item If $\vec{p}$ and $\vec{q}$ are in $U$, then $\vec{p}+\vec{q}$ is in $U$. ($U$ is \textbf{closed under vector addition}.)
\item If $\vec{p}$ is in $U$ and $c$ is in $\R$, then $c\vec{p}$ is in $U$. ($U$ is \textbf{closed under scalar multiplication}.)
\end{enumerate}
\end{definition}

\example Just like in chapter \ref{chap:existence-and-uniqueness-of-solutions}, the smallest subspace of a vector space $V$ is $\left\{\vec{0}\right\}$, and the largest subspace of $V$ is $V$ itself. \blacksquare

\example Is the set $U=\left\{ a+bz^2 : a \ge 0, b \ge 0\right\}$ a subspace of $\mathbb{P}_2$?

\begin{notes}
\vfill \hspace{1pt}
\columnbreak
\end{notes}

\begin{book}
\textbf{Solution.} (i) It does contain the zero vector, since $a,b$ can be zero. 

(ii) It is closed under addition, since if $a+bz^2$ and $c+dz^2$ are in $U$, then all of $a,b,c,d$ are greater or equal to 0, so $a+c\ge 0$ and $b+d \ge 0$, and
\[(a+bz^2) + (c+dz^2) = (a+c) + (b+d)z^2 \text{ is in } U.\]

(iii) It is not closed under scalar multiplication. For example, $3+2z^2$ is in $U$, and $-5$ is a real number, but $-5(3+2z^2) = -15 -10z^2$ is not in $U$. So $U$ is not a subspace. \blacksquare
\end{book}

Note that properties (i), (ii), and (iii) in the definition of a subspace are the same as properties (i), (ii), and (iii) of a vector space. And, if $U$ is a subspace of $V$, then any vectors $\vec{u},\vec{v},\vec{w}$ in $U$ are also vectors in $V$, which makes properties the rest of the properties, (iv) through (ix) in the definition of a vector space, also true for $U$. This proves:

\begin{theorem}[Subspaces are vector spaces.]
\label{theorem:subspaces-are-vector-spaces}
If $U$ is a subspace of a vector space $V$, then $U$ itself is a vector space.
\end{theorem}

The following theorems are just like theorems \ref{theorem:spans-are-subspaces} and \ref{theorem:subspaces-are-spans}, but with $\R^c$ replaced with an abstract vector space $V$.

\begin{theorem}[Spans are subspaces.]\label{theorem:abstract-spans-are-subspaces}

If $\vec{a}_1, \ldots, \vec{a}_c$ is a list of vectors in a vector space $V$, then $\spn (\vec{a}_1,\ldots, \vec{a}_c)$ is a subspace of $V$.
\end{theorem}

Like before, if you read the proof of theorem \ref{theorem:spans-are-subspaces}, you'll find that it only relies on properties that are true about vector spaces in general, not just $\R^c$. So theorem \ref{theorem:abstract-spans-are-subspaces} is already proven. The next definitions are also the same as in chapter 1, but with $V$ instead of $\R^n$.

\begin{definition}
Let $U$ be a subspace of a vector space $V$. (Note that since $V$ is a subspace of itself, this definition also applies to $U = V$.) A list of vectors $\mathcal{B} = ( \vec{b}_1,\ldots, \vec{b}_c )$ is called a \textbf{basis} for $U$ if:
\begin{enumerate}[label=(\roman*)]
\item $\mathcal{B}$ is linearly independent.
\item $\spn \mathcal{B} = U$.
\end{enumerate}
\end{definition}

\begin{definition}
Suppose $\mathcal{B} = (\vec{b}_1,\ldots,\vec{b}_c)$ is a basis for a subspace $U$ of a vector space $V$, and that $\vec{y}$ is a vector in $V$. If 
\[ x_1\vec{b}_1+\cdots+x_c\vec{b}_c = \vec{y}\]
then the vector $(x_1,\ldots, x_c)$ is called the \textbf{$\mathcal{B}$-coordinate vector of $\vec{y}$} and we write $(x_1,\ldots,x_c) =\left[ \vec{y} \right]_{\mathcal{B}}$.
\end{definition}

\begin{notes}
\columnbreak
\end{notes}

\example The \emph{standard basis of} $\mathbb{P}_3$ is $\mathcal{B}=( 1,z,z^2,z^3 )$.
\begin{enumerate}[label=(\alph*)]
\item Find $\left[ \pi - 7z^2 \right]_{\mathcal{B}}$.
\item If $\left[ \vec{y} \right]_{\mathcal{B}} = (5,4,3,2)$, find $\vec{y}$.
\end{enumerate}

\begin{book}
\textbf{Solution.} (a) Since we can write it as $\pi + 0z +(-7)z^2+0z^3$, we have $\left[ \pi - 7z^2 \right]_{\mathcal{B}} = (\pi,0,-7,0)$.

(b) $\vec{y} = 5 + 4z + 3z^2+2$. \blacksquare
\end{book}

\begin{notes}
\vfill
\end{notes}

Here are some more theorems that have $\R^c$ replaced with an arbitrary vector space $V$, that were really proved in chapter \ref{chap:existence-and-uniqueness-of-solutions}: 

\begin{theorem}[Unique representation theorem.]
\label{theorem:abstract-unique-representation}

Let $\mathcal{B} = ( \vec{b}_1,\ldots,\vec{b}_c )$ be a basis for a subspace $U$ of $\R^n$. Then for each $\vec{y}$ in $U$, there is only one $\left[ \vec{y} \right]_{\mathcal{B}}$.
\end{theorem}

\begin{theorem}[Properties of coordinates.]
\label{theorem:abstract-properties-of-coordinates}

Suppose $\mathcal{B} = ( \vec{b}_1, \ldots, \vec{b}_n )$ is a basis of a vector space $V$. Then, for any $\vec{u},\vec{v}$ in $V$ and any real number $k$:
\begin{enumerate}[label=(\roman*)]
\item $\left[ k\vec{u} \right]_{\mathcal{B}} = k \left[ \vec{u} \right]_{\mathcal{B}}$.
\item $\left[ \vec{u} + \vec{v} \right]_{\mathcal{B}} =  \left[ \vec{u} \right]_{\mathcal{B}}+\left[ \vec{v} \right]_{\mathcal{B}}$.
\item If $\left[ \vec{u} \right]_{\mathcal{B}} = \left[ \vec{v} \right]_{\mathcal{B}}$, then $\vec{u}=\vec{v}$. \quad (We say $f(\vec{u}) =\left[ \vec{u} \right]_{\mathcal{B}}$ is an \emph{injective} or \emph{one-to-one} function.)
\item If $\vec{x}$ is any vector in $\R^n$, then there is a $\vec{u}$ in $V$ so that $\left[ \vec{u} \right]_{\mathcal{B}}=\vec{x}$. \quad (We say $f(\vec{u}) =\left[ \vec{u} \right]_{\mathcal{B}}$ is a \emph{surjective} or \emph{onto} function.)
\end{enumerate}
\end{theorem}

To sum up theorem \ref{theorem:abstract-properties-of-coordinates}: vectors $\vec{u},\vec{v}$ in an $n$-dimensional vector space $V$ really behave just like their coordinate vectors $[\vec{u}]_{\mathcal{B}}, [\vec{v}]_{\mathcal{B}}$ in $\R^n$, at least as far as scalar multiplication and vector addition are concerned; and there's a one-to-one correspondence between them. Next we'll show that this extends to the span, linear independence, solutions of vector equations. In the next section we'll extend this to linear transformations.


\begin{notes}
\columnbreak
\end{notes}

\begin{theorem}[More properties of coordinates.]
\label{theorem:abstract-more-properties-of-coordinates}

Let $\mathcal{B} = ( \vec{b}_1, \ldots, \vec{b}_n)$ be a basis for a vector space $V$.
\begin{enumerate}[label=(\roman*)]
\item $x_1 \vec{a}_1 + \cdots + x_p \vec{a}_p = \vec{y}$ is true and only if $x_1 [\vec{a}_1]_{\mathcal B} + \cdots + x_p [\vec{a}_p]_{\mathcal B} = [\vec{y}]_{\mathcal B}$ is true.
\item A list $( \vec{a}_1, \ldots, \vec{a}_p )$ in $V$ is linearly dependent if and only if $( [\vec{a}_1]_{\mathcal{B}}, \ldots, [\vec{a}_p]_{\mathcal{B}} )$ is linearly dependent in $\R^{n}$.
\item A list $( \vec{a}_1, \ldots, \vec{a}_p )$ in $V$ spans $V$ if and only $( [\vec{a}_1]_{\mathcal{B}}, \ldots, [\vec{a}_p]_{\mathcal{B}} )$ spans $\R^{n}$.
\end{enumerate}
\end{theorem}

Let's prove (iii) and leave (i) and (ii) for exercises.

\begin{notes}
\vfill
\end{notes}

\begin{book}
\textbf{Proof.} (iii) (\rightarrow) Suppose $\mathcal{B} = ( \vec{b}_1, \ldots, \vec{b}_n)$ is a basis for a vector space $V$, and $( \vec{a}_1 , \ldots, \vec{a}_p)$ spans $V$. Let $\vec{x}$ be any vector in $\R^n$. By theorem \ref{theorem:abstract-properties-of-coordinates} part (iv), there exists $\vec{u}$ in $V$ such that $[\vec{u}]_{\mathcal{B}} = \vec{x}$; that is, $\vec{u}=x_1\vec{b}_1+\cdots + x_n\vec{b}_n$. Since $( \vec{a}_1 , \ldots, \vec{a}_p)$ spans $V$, there exist real numbers $r_1,\ldots,r_p$ such that:
\[r_1\vec{a}_1+\cdots+r_p\vec{a}_p=\vec{u}\]
Plugging both sides into the $[]_{\mathcal{B}}$ function:
\begin{align*}
\left[r_1\vec{a}_1+\cdots+r_p\vec{a}_p\right]_{\mathcal{B}}&=\left[\vec{u}\right]_{\mathcal{B}} \\
\left[r_1\vec{a}_1\right]_{\mathcal{B}}+\cdots+\left[r_p\vec{a}_p\right]_{\mathcal{B}}&= \vec{x} & \text{(theorem \ref{theorem:abstract-properties-of-coordinates} (ii))} \\ 
r_1\left[\vec{a}_1\right]_{\mathcal{B}}+\cdots+r_p\left[\vec{a}_p\right]_{\mathcal{B}}&= \vec{x} & \text{(theorem \ref{theorem:abstract-properties-of-coordinates} (i))} 
\end{align*}
Thus any vector $\vec{x}$ in $\R^n$ can be written as a linear combination of $\left[ \vec{a}_1 \right]_{\mathcal{B}},\ldots,\left[ \vec{a}_p \right]_{\mathcal{B}}$; that is, the list spans $\R^n$.

(\leftarrow) Now suppose $\left[ \vec{a}_1 \right]_{\mathcal{B}},\ldots,\left[ \vec{a}_p \right]_{\mathcal{B}}$ spans $\R^n$, and let $\vec{v}$ be any vector in $V$. Since $\left[ \vec{v} \right]_{\mathcal{B}}$ is in $\R^n$, there exist real numbers $t_1,\ldots,t_p$ such that:
\begin{align*}
t_1 \left[ \vec{a}_1 \right]_{\mathcal{B}} + \cdots + t_p \left[ \vec{a}_p \right]_{\mathcal{B}} &= \left[ \vec{v} \right]_{\mathcal{B}} \\
 \left[t_1 \vec{a}_1 \right]_{\mathcal{B}} + \cdots +  \left[t_p \vec{a}_p \right]_{\mathcal{B}} &= \left[ \vec{v} \right]_{\mathcal{B}} &(\text{theorem \ref{theorem:abstract-properties-of-coordinates} (i)})\\
 \left[t_1 \vec{a}_1  + \cdots +  t_p \vec{a}_p \right]_{\mathcal{B}} &= \left[ \vec{v} \right]_{\mathcal{B}} &(\text{theorem \ref{theorem:abstract-properties-of-coordinates} (ii)})\\
t_1 \vec{a}_1  + \cdots +  t_p \vec{a}_p  &=  \vec{v}  &(\text{theorem \ref{theorem:abstract-properties-of-coordinates} (iii)})
\end{align*}
Thus any vector $\vec{v}$ in $V$ can be written as a linear combination of $\vec{a}_1,\ldots,\vec{a}_p$, so the list spans $V$. \blacksquare
\end{book}

The next two theorems were proved in chapter \ref{chap:existence-and-uniqueness-of-solutions}.

\begin{theorem}[More vectors than a basis are dependent.]
\label{theorem:abstract-more-vectors-than-a-basis}
If a subspace $U$ of a vector space $V$ has a basis with $n$ vectors in it, then any list with more than $n$ vectors in $U$ is linearly dependent.
\end{theorem}

\begin{theorem}
\label{theorem:abstract-all-bases-same-size}
If $U$ is a subspace of a vector space $V$, then every basis of $U$ has the same number of vectors.
\end{theorem}

\begin{notes}
\columnbreak
\end{notes}

\begin{definition}
The \textbf{dimension} of a subspace $U$ of a vector space $V$, written $\dim U$, is the number of vectors in a basis of $U$. If $\dim U = n$, we say $U$ is $n$\textbf{-dimensional}. 

\emph{And something new:} If $U$ does not have a finite basis, we say that $U$ is \textbf{infinite-dimensional}.
\end{definition}

\example Show that $\mathbb{P}$, the vector space of all polynomials, is infinite dimensional.

\begin{notes}
\vfill
\end{notes}

\begin{book}
\textbf{Solution.} We'll prove that it's impossible for $\mathbb{P}$ to have a finite basis. Suppose that $\mathcal{B}$ is a finite list of vectors. Since it's finite, there's a natural number $k$ which is greater than the degree of all the polynomials in $\mathcal{B}$, and there's no way to get $x^k$ as a linear combination of the lower-degree polynomials in $\mathcal{B}$. Thus no finite list can be span $\mathbb{P}$. \blacksquare
\end{book}

\begin{theorem}[Buy two, get one free.]
\label{theorem:buy-two-get-one-free}

Suppose $\vec{v}_1, \ldots, \vec{v}_n$ is a list of $n$ vectors in a subspace $U$ of $\R^c$. If any two of the following are true, then so is the third:
\begin{enumerate}[label=(\roman*)]
\item $\dim U = n$.
\item $\vec{v}_1,\ldots,\vec{v}_n$ is linearly independent.
\item $\vec{v}_1,\ldots,\vec{v}_n$ spans $U$.
\end{enumerate}
\end{theorem}

\begin{notes}
\columnbreak
\end{notes}

\example If $( \vec{p}_1,\vec{p}_2,\vec{p}_3 ) = ( 3-z+6z^2,-2+z-2z^2, -4 + 2z-5z^2 )$, find the dimension of $\spn ( \vec{p}_1,\vec{p}_2,\vec{p}_3 )$.

\begin{book}
\textbf{Solution.} Using the standard basis $\mathcal{B} = (1,z,z^2)$, the coordinates are:
\[\left[ \vec{p}_1 \right]_{\mathcal{B}}=\bbbvec{3}{-1}{6}, \left[ \vec{p}_2 \right]_{\mathcal{B}}=\bbbvec{-2}{1}{-2},\left[ \vec{p}_3 \right]_{\mathcal{B}}=\bbbvec{-4}{2}{-5}\]
Theorems \ref{theorem:abstract-properties-of-coordinates} and \ref{theorem:abstract-more-properties-of-coordinates} essentially say that anything about $( \left[\vec{p}_1\right]_{\mathcal{B}},\left[\vec{p}_2\right]_{\mathcal{B}},\left[\vec{p}_3\right]_{\mathcal{B}} )$ is also true about $( \vec{p}_1,\vec{p}_2,\vec{p}_3 )$; so let's think about a basis of $\spn( \left[\vec{p}_1\right]_{\mathcal{B}},\left[\vec{p}_2\right]_{\mathcal{B}},\left[\vec{p}_3\right]_{\mathcal{B}} )$. If these three $\R^3$ vectors are linearly independent, then they're a basis of their span. So let's row-reduce the augmented matrix for the homogeneous equation:
\begin{align*}
x_1\left[\vec{p}_1\right]_{\mathcal{B}}+x_2\left[\vec{p}_2\right]_{\mathcal{B}}+x_3\left[\vec{p}_3\right]_{\mathcal{B}} = \vec{0} \\
\begin{bmatrix}[ccc|c]
3 & -2 & -4 &0 \\
-1 & 1 & 2 &0 \\
6 & -2 & -4 & 0
\end{bmatrix}
\sim \begin{bmatrix}[ccc|c]
1 & 0 & 0 &0\\
0 & 1 & 0 &0\\
0 & 0 & 1 &0
\end{bmatrix}
\end{align*}
This has the unique solution $(x_1,x_2,x_3) = (0,0,0)$, and by theorem \ref{theorem:abstract-more-properties-of-coordinates} (ii), so does $x_1\vec{p}_1+x_2\vec{p}_2+x_3\vec{p}_3=\vec{0}$, so the three polynomials are linearly independent, and $( \vec{p}_1,\vec{p}_2,\vec{p}_3 )$ is a basis of its span, which thus has dimension 3. \blacksquare
\end{book}

\begin{notes}
\vspace{20em}
\end{notes}

\exerciseheading

\exatend (T/F?) Every vector space is a subspace of itself.

\exatend (T/F?) The list of polynomials $z,5z^2$ is linearly dependent because $-5z$ and $1$ are scalars that make $-z(z)+1(5z^2) = \vec{0}$.

\exatend (T/F?) In $\Poly_2$, the function $\vec{f}(z) = 3z^2$ is the zero vector because $\vec{f}(0)=0$.

\exatend (T/F?) In the vector space of all $\R\rightarrow\R$ functions, $\spn(\cos(z),\sin(z))$ is the interval $\left[ -1,1 \right]$.

\exatend (T/F?) In the vector space of all $\R\rightarrow\R$ functions, $\spn(\cos(z),\sin(z))$ is the set of all functions that can be written $a\cos(z)+b\sin(z)$ for some real numbers $a,b$.

\exatend Determine if $3z$ is in $\spn ( 1+z,1-z )$.

\exatend Is the list of polynomials $( \vec{p}_1,\vec{p}_2,\vec{p}_3,\vec{p}_4 ) =$ $( 1+2z-z^2$, $-5-9z+5z^2$, $4+11z-3z^2$, $2+4z+2z^2)$ linearly independent in $\Poly_2$? Explain.

\exatend Consider the list of polynomials $( \vec{p}_1,\vec{p}_2,\vec{p}_3 ) =$ $( 1+2z+3z^2$, $4+5z+6z^2$, $7+8z+9z^2)$.
\begin{enumerate}[label=(\alph*)]
\item Is the list linearly independent? Explain.
\item Does the list span $\Poly_2$? Explain.
\end{enumerate}

\exatend It can be shown that $\mathcal{B}=(2+z$, $3+z)$ is a basis of $\Poly_1$.
\begin{enumerate}[label=(\alph*)]
\item  Find $\left[ -4+z \right]_{\mathcal{B}}$.
\item If $\left[ \vec{p} \right]= (3,5)$, find $\vec{p}$.
\end{enumerate}

\exatend It can be shown that $\mathcal{B}=( 2+z+z^2$, $4+4z+3z^2$,-$5-z-2z^2)$ is a basis of $\Poly_2$.
\begin{enumerate}[label=(\alph*)]
\item Find $\left[ z \right]_{\mathcal{B}}$.
\item If $\left[ \vec{p} \right]= (1,2,3)$, find $\vec{p}$.
\end{enumerate}


\exatend Explain why every basis of $\Poly_3$ contains at least one polynomial with degree $3$.

\exatend Prove theorem \ref{theorem:vector-space-properties} part (iii). \hintfootnote{You may use part (ii) of the theorem.}

\exatend Prove theorem \ref{theorem:vector-space-properties} part (iv). \hintfootnote{$\vec{0}=\vec{0}+\vec{0}$.}

\exatend Prove theorem \ref{theorem:abstract-more-properties-of-coordinates} part (i).

\exatend Prove theorem \ref{theorem:abstract-more-properties-of-coordinates} part (ii).

\exatend Suppose $\vec{u}_1,\ldots,\vec{u}_n$ is linearly independent in a vector space $V$ and $\vec{w}$ is in $V$. Prove that if $u_1+\vec{w}, \ldots, \vec{u}_n+\vec{w}$ is linearly dependent, then $\vec{w}$ is in $\spn(\vec{u}_1,\ldots,\vec{u}_n)$. \creditfootnote{Inspired by an exercise in Axler's excellent book \emph{Linear Algebra Done Right}.}


\exatend Let $U = \left\{ \vec{p} \text{ in } \Poly_4 : \vec{p}(6)=0 \right\}$.\creditfootnote{Inspired by Axler.}
\begin{enumerate}[label=(\alph*)]
\item Show that $U$ is a subspace of $\Poly_4$.
\item  Find a basis for $U$. 
\end{enumerate}


\exatend Let $U=\left\{ \vec{p} \text{ in } \Poly_5 : \int_{-1}^1\vec{p}(z)dz=0 \right\}$. Find a basis for $U$.\creditfootnote{Inspired by Axler.}
\begin{enumerate}[label=(\alph*)]
\item Show that $U$ is a subspace of $\Poly_5$.
\item  Find a basis for $U$. 
\end{enumerate}

\exatend Let $V$ be the set of all functions of the form $\vec{f}(z) = a\cos(5z) + b\sin(5z)$ for some $a,b$ in $\R$. Is $V$ a subspace of the vector space of all $\R \rightarrow \R$ functions? If so, what is its dimension?

\exatend Let $W$ be the set of all functions of the form $\vec{f}(z) = a\cos(kz) + b\sin(kz)$ for some $a,b,k$ in $\R$. Is $W$ a subspace of the vector space of all $\R \rightarrow \R$ functions? If so, what is its dimension?

\exatend Let $E$ be the set of all even $\R\rightarrow\R$ functions. That is, the set of functions $f:\R\rightarrow\R$ such that $f(-z) = f(z)$ for all $z$ in $\R$. Is $E$ a subspace of the vector space all $\R\rightarrow\R$ functions? Explain.

\exatend Let $U$ be the set of polynomials $\vec{p}$ in $\Poly_2$ such that $\vec{p}(3)=4$. Is $U$ a subspace of $\Poly_2$? Explain.

\exatend Let $D$ be the set $\left\{ \vec{p}(z) = az^2+bz+c \text{ in }\Poly_2:b^2-4ac=0\right\}$. Is $D$ a subspace of $\Poly_2$? Explain.

\exatend Let $U=\left\{ \vec{p} \text{ in } \Poly_4 : \vec{p}''(7)=0 \right\}$. Find a basis for $U$. \creditfootnote{Inspired by Axler.}

% This question kinda sucks; it's not really useful for our class, but might be a fun challenge.
\exatend Suppose $r$ is in $\R$ and $\vec{u}$ is in a vector space $V$. Show that if $r\vec{u}=\vec{0}$, then either $r=0$ or $\vec{u}=\vec{0}$.

\exatend (\emph{Challenge}) Prove or give a counterexample: If $(\vec{b}_1,\vec{b}_2,\vec{b}_3,\vec{b}_4)$ is a basis of a vector space $V$, and $U$ is a subspace of $V$, with $\vec{b}_1,\vec{b}_2$ in $U$, but $\vec{b}_3$ and $\vec{b}_4$ are not in $U$, then $( \vec{b}_1,\vec{b}_2 )$ is a basis of $U$. \creditfootnote{Inspired by Axler.}

%requires thm that more vectors than a basis makes a list dependent
\exatend (\emph{Challenge}) Suppose that $\vec{p}_0,\vec{p}_1,\ldots,\vec{p}_{m}$ are polynomials in $\mathbb{P}_m$, such that $\vec{p}_j(2)=0$ for each $j$. Show that $\vec{p}_0,\vec{p}_1,\ldots,\vec{p}_{m}$ is linearly dependent in $\mathbb{P}_m$. \creditfootnote{Inspired by Axler.}

\exatend Many curious students ask questions like, ``if we can expand the real numbers to include complex numbers like $i$, with $i^2=-1$, could we also expand real numbers to include $\infty$ and $-\infty$?'' Let's denote $\mathbb{I}$ to be the set that includes all real numbers, and also two distinct objects $\infty$ and $-\infty$, and define multiplication and addition in a reasonable way. Specifically, for any real number $t$:
\begin{align*}
&t+\infty=\infty+t =\infty & t+(-\infty) = (-\infty) + t =\infty, \\
&\infty+\infty=\infty  &(-\infty) + (-\infty) = (-\infty) \\
 &\infty+(-\infty) = 0 & \\
&t\infty=\left\{\begin{matrix}
-\infty \text{ if } t<0 \\ 
0 \text{ if } t=0 \\ 
\infty \text{ if } t>0 
\end{matrix}\right.
&t(-\infty)=\left\{\begin{matrix}
\infty \text{ if } t<0 \\ 
0 \text{ if } t=0 \\ 
-\infty \text{ if } t>0 
\end{matrix}\right.
\end{align*}
Is $\mathbb{I}$ a vector space? Justify your answer. \creditfootnote{Credit: Inspired by Axler.}

\exatend Suppose $\vec{p}_1,\vec{p}_2,\ldots,\vec{p}_n,\vec{p}_{n+1}$ is a list of $n+1$ polynomials in $\Poly_n$, all with the property that $\vec{p}_i(3)=0$. Show that the list is linearly dependent.

\columnbreak
\section{Abstract Linear Transformations}

In this section, we'll expand the notion of linear transformations to abstract vector spaces, and see how they are essentially the same as the $\R^c\rightarrow \R^r$ linear transformations introduced in section \ref{sec:matrix-vector-products} and studied throughout the book. First let's warm up with a review.

\exercise (\emph{Warm-up}) (T/F?) If $A$ is a $c \times r$ matrix, then $T(\vec{x}) = A\vec{x}$ is an $\R^c\rightarrow \R^r$ linear transformation.

\begin{notes}
\vspace{1em}
\end{notes}

\exercise (\emph{Warm-up}) (T/F?) If $T:\R^c\rightarrow \R^r$ is a linear transformation, then there is a $c\times r$ matrix $A$ such that $T(\vec{x})=A\vec{x}$.

\begin{notes}
\vspace{1em}
\end{notes}

\exercise (\emph{Warm-up}) (T/F?) Suppose $A$ is a matrix. Then the range of the transformation $T(\vec{x}) = A\vec{x}$ is the column space of $A$.

\begin{notes}
\vspace{1em}
\end{notes}

\exercise(\emph{Warm-up}) Let $T(\vec{x}) = A\vec{x}$ for $A = \begin{bmatrix}
1 & 2 & 3 \\
4 & 5 & 6
\end{bmatrix}$. Let $\mathcal{B}= ( \vec{e}_1,\vec{e}_2,\vec{e}_3 )$ and $\mathcal{D} = ( \vec{e}_1,\vec{e}_2 )$ be the standard bases of $\R^3$ and $\R^2$, respectively.
\begin{enumerate}[label=(\alph*)]
\item Find $T(\vec{e}_3)$.
\item (T/F?) $A = \begin{bmatrix}
T(\vec{e}_1) & T(\vec{e}_2) & T(\vec{e}_3)
\end{bmatrix}$.
\item Find $\left[T(\vec{e}_3)\right]_{\mathcal{D}}$.
\item (T/F?) For any $\vec{y}=(y_1,y_2)$ in $\R^2$, we have $\left[ \vec{y} \right]_{\mathcal{D}}=(y_1,y_2)$.
\end{enumerate}

\begin{notes}
\vfill
\end{notes}

\begin{definition}
Suppose $V$ and $W$ are vector spaces. A function $T:V\rightarrow W$ is a \textbf{linear transformation} if for all real numbers $k$ and all vectors $\vec{u},\vec{v}$ in $V$:
\begin{enumerate}[label=(\roman*)]
\item $T(\vec{u}+\vec{v}) = T(\vec{u}) + T(\vec{v})$
\item $T(k\vec{u}) = kT(\vec{u})$
\end{enumerate}
\end{definition}

\example Let $T:\Poly \rightarrow \Poly$  be differentiation; that is, $T(\vec{p}(z)) = \frac{d}{dz}\vec{p}(z)$. Is $T$ a linear transformation?

\begin{notes}
\vfill \hspace{1pt}
\columnbreak
\end{notes}

\begin{book}
\textbf{Solution.} A basic calculus fact is that if $f(z)$ and $g(z)$ are differentiable functions (including polynomials), then \[\frac{d}{dz}(f(z) + g(z)) = \frac{d}{dz}f(z) + \frac{d}{dz}g(z) \quad \text{and} \quad \frac{d}{dz}(kf(z)) = k \frac{d}{dz}f(z).\]

This shows that $T(\vec{p}(z)) = \frac{d}{dz}\vec{p}(z)$ satisfies the definition of a linear transformation. \blacksquare
\end{book}

\example  Is the function $T:\Poly\rightarrow \Poly$ defined by $T ( \vec{p}(z) ) = \vec{p}(z) + 5z$ a linear transformation?

\begin{notes}
\vfill
\end{notes}

\begin{book}
\textbf{Solution.} If there exists polynomials $\vec{p},\vec{q}$ and a real number $k$ that makes either of the following false, then $T$ is not a linear transformation:
\begin{enumerate}[label=(\roman*)]
\item Is $( \vec{p}(z)+\vec{q}(z) )+5z = ( \vec{p}(z)+5z ) + ( \vec{q}(z)+5z )$?
\item Is $( k\vec{p}(z) )+5z = k( \vec{p}(z)+5z )$?
\end{enumerate}
No, on both counts. Almost any choice of $\vec{p},\vec{q}$ and $k$ suffices as a counterexample to both, for example $\vec{p}(z)=3+z$, $\vec{q}(z) = z^2, k = 2$. \blacksquare
\end{book}

\example Suppose $T:V\rightarrow W$ is a linear transformation. Denote the zero vector of $V$ as $\vec{0}_V$ and the zero vector of $W$ as $\vec{0}_W$. Show that $T(\vec{0}_V) = \vec{0}_{W}$.

\begin{notes}
\vfill
\end{notes}

\begin{book}
\textbf{Solution.} Writing the real number zero as just $0$:
\begin{align*}
T(\vec{0}_V) &= T(0\vec{0}_V) & (\text{Thm \ref{theorem:vector-space-properties} part (ii)}) \\
&= 0T(\vec{0}_V) & (\text{Defn of linear transformation}) \\
&= \vec{0}_W & (\text{Thm \ref{theorem:vector-space-properties} part (ii)}) \quad \blacksquare
\end{align*}
\end{book}

Recall  that every $\R^c \rightarrow \R^r$ linear transformation can be written $T(\vec{x})= A\vec{x}$, with $A = \begin{bmatrix}
T ( \vec{e}_1 ) & \cdots & T ( \vec{e}_c )
\end{bmatrix}$, called the \emph{standard matrix} of $T$ (theorem \ref{theorem:standard-matrix}). The next example and theorem show how this concept can be generalized to abstract vector spaces. This will reinforce the idea that, as far as linear algebra is concerned, finite-dimensional vector spaces are essentially the same as $\R^c$. 

\begin{notes}
\columnbreak
\end{notes}

\example Let $T:\Poly_2\rightarrow \Poly_1$ be the differentiation transformation, $T(\vec{p}(z)) = \frac{d}{dz}\vec{p}(z)$, and denote the standard bases of $\Poly_2$ and $\Poly_1$ as $\mathcal{B}= \left\{ \vec{b}_1,\vec{b}_2,\vec{b}_3 \right\} = \left\{ 1,z,z^2 \right\}$ and $\mathcal{D}= \left\{ \vec{d}_1,\vec{d}_2 \right\} = \left\{ 1,z \right\}$.
\begin{enumerate}[label=(\alph*)]
\item Find each: $T(\vec{b}_1),T(\vec{b}_2)$, and $T(\vec{b}_3)$.
\item Find each: $\left[ T(\vec{b}_1) \right]_{\mathcal{D}},\left[ T(\vec{b}_2) \right]_{\mathcal{D}}$, and $\left[ T(\vec{b}_3) \right]_{\mathcal{D}}$.
\item Find $T(10-9z+7z^2)$.
\item Find $\left[10-9z+7z^2\right]_{\mathcal{B}}$.
\item Find $\left[ T(10-9z+7z^2) \right]_{\mathcal{D}}$.
\item Find $A\left[10-9z+7z^2\right]_{\mathcal{B}}$, with $A = \begin{bmatrix}
\left[ T(\vec{b}_1) \right]_{\mathcal{D}} & \left[ T(\vec{b}_2) \right]_{\mathcal{D}} & \left[ T(\vec{b}_3) \right]_{\mathcal{D}}
\end{bmatrix}$.
\item Explain why it's not a coincidence that $\left[ T(10-9z+7z^2) \right]_{\mathcal{D}}$ equals $A \left[ 10-9z+7z^2 \right]_{\mathcal{B}}$.
\end{enumerate}

\begin{notes}
\vfill \hspace{1pt}
\columnbreak
\end{notes}

\begin{book}
\textbf{Solution.} (a) $T(\vec{b}_1) = \frac{d}{dz}1=0$, $T(\vec{b}_2) = \frac{d}{dz}z=1$, and $T(\vec{b}_2) = \frac{d}{dz}z^2=2z$.

(b) Since $T(\vec{b}_1) = 0 = 0(1) + 0z = 0\vec{d}_1+0\vec{d}_2$, we have $\left[ T(\vec{b}_1) \right]_{\mathcal{D}} = \bbvec{0}{0}$.

Since $T(\vec{b}_2) = 1 = 1(1) + 0z = 1\vec{d}_1+0\vec{d}_2$, we have $\left[ T(\vec{b}_2) \right]_{\mathcal{D}} = \bbvec{1}{0}$.

Since $T(\vec{b}_3) = 2z = 0(1) + 2z = 0\vec{d}_1+2\vec{d}_2$, we have $\left[ T(\vec{b}_3) \right]_{\mathcal{D}} = \bbvec{0}{2}$.

(c) $T(10-9z+7z^2) = \frac{d}{dz}(10-9z+7z^2) = -9+14z$.

(d) Since $10-9z+7z^2=10\vec{b}_1+(-9)\vec{b}_2+7\vec{b}_3$, we have $\left[ 10-9z+7z^2 \right]_{\mathcal{B}}=\bbbvec{10}{-9}{7}$.

(e) Since $T(10-9z+7z^2) = -9+14z = -9\vec{d}_1+14\vec{d}_2$, we have $\left[ T(10-9z+7z^2) \right]_{\mathcal{d}}=\bbvec{-9}{14}$.

(f) We already found the columns of $A$ in part (b), so $A = \begin{bmatrix}
0 & 1 & 0 \\
0 & 0 & 2
\end{bmatrix}$. Putting this all together:
\[A\left[10-9z+7z^2\right]_{\mathcal{B}} = \begin{bmatrix}
0 & 1 & 0 \\
0 & 0 & 2
\end{bmatrix} \bbbvec{10}{-9}{7} = \bbvec{10(0) -9(1)+7(0) }{10(0) -9(0)+7(2)} = \bbvec{-9}{14}\]

(g) Using the linearity of $T$ and the properties of coordinates (theorem \ref{theorem:abstract-properties-of-coordinates}):
\begin{align*}
&\left[ T(10-9z+7z^2) \right]_{\mathcal{D}}\\
 &= \left[ 10T(1) -9 T(z) +7 T(z^2) \right]_{\mathcal{D}} \quad \quad \quad \text{(linearity of } T\text{)} \\
&= \left[ 10T(1) \right]_{\mathcal{D}} + \left[ -9T(z) \right]_{\mathcal{D}} + \left[ 7T(z^2) \right]_{\mathcal{D}} \quad \quad \text{(theorem \ref{theorem:abstract-properties-of-coordinates} part (ii))} \\
&= 10\left[T(1) \right]_{\mathcal{D}} -9 \left[ T(z) \right]_{\mathcal{D}} + 7\left[ T(z^2) \right]_{\mathcal{D}} \quad \quad \text{(theorem \ref{theorem:abstract-properties-of-coordinates} part (i))} \\
&= \begin{bmatrix}
\left[ T(1) \right]_{\mathcal{D}} & \left[ T(z) \right]_{\mathcal{D}} & \left[ T(z^2) \right]_{\mathcal{D}} 
\end{bmatrix}
\bbbvec{10}{-9}{7} \quad \quad \text{(defn of matrix-vector multiplication)} \\
&= \begin{bmatrix}
\left[ T(\vec{b}_1) \right]_{\mathcal{D}} & \left[ T(\vec{b}_2) \right]_{\mathcal{D}} & \left[ T(\vec{b}_3) \right]_{\mathcal{D}} 
\end{bmatrix}
\left[ 10-9z+7z^2 \right]_{\mathcal{B}}
\end{align*}
\blacksquare
\end{book}

The justification for part (g) in the above example didn't use any special properties of the differentiation transformation or $\Poly_1$ or $\Poly_2$; the proof of the theorem below is basically the same.
\begin{notes}
You can find it in the textbook.
\end{notes}

\begin{theorem}[Matrix of a transformation.] \label{theorem:abstract-matrix-of-transformation}

Suppose $V$ and $W$ are vector spaces with bases $\mathcal{B} = \left\{ \vec{b}_1,\ldots,\vec{b}_c \right\}$ and $\mathcal{D}$, respectively, and $T:V \rightarrow W$ is a linear transformation. Then for all $\vec{v}$ in $V$,
\[\left[ T(\vec{v}) \right]_{\mathcal{D}} = A \left[ \vec{v} \right]_{\mathcal{B}}\]
where $A = \begin{bmatrix}
\left[ T(\vec{b}_1) \right]_{\mathcal{D}} & \cdots  & \left[ T(\vec{b}_c) \right]_{\mathcal{D}} 
\end{bmatrix}$, which is called the \textbf{matrix of $T$ relative to $\mathcal{B}$ and $\mathcal{D}$.}
\end{theorem}

\begin{book}
\textbf{Proof.} Suppose $\vec{v}$ is in $V$. Let $\left[ \vec{v} \right]_{\mathcal{B}}=(x_1,\ldots,x_c)$; in other words, let $\vec{v}=x_1\vec{b}_1+\cdots+x_c\vec{b}_c$. Then we have:
\begin{align*}
T(\vec{v}) &= T(x_1\vec{b}_1+\cdots+x_c\vec{b}_c) \\
&= x_1T(\vec{b}_1) + \ldots + x_cT(\vec{b}_c) & \text{(linearity of }T\text{)}
\end{align*}
Taking coordinates relative to $\mathcal{D}$:
\begin{align*}
\left[ T(\vec{v}) \right]_{\mathcal{D}} &= \left[ x_1T(\vec{b}_1) + \ldots + x_cT(\vec{b}_c) \right]_{\mathcal{D}} \\
&= x_1\left[T(\vec{b}_1) \right]_{\mathcal{D}} + \ldots + x_c\left[T(\vec{b}_c) \right]_{\mathcal{D}} & \text{(thm \ref{theorem:abstract-properties-of-coordinates} (i) and (ii))} \\
&= \begin{bmatrix}
\left[ T(\vec{b}_1)\right]_{\mathcal{D}} & \cdots & \left[T(\vec{b}_c) \right]_{\mathcal{D}}
\end{bmatrix}\bbbvec{x_1}{\vdots}{x_c} & \text{(defn of matrix-vector multiplication)}\\
&= A \left[ \vec{v} \right]_{\mathcal{B}} & \blacksquare
\end{align*}
\end{book}

\example Let $T:\Poly_2\rightarrow \Poly_1$ be the differentiation transformation. It can be shown that $\mathcal{B} = (\vec{b}_1,\vec{b}_2,\vec{b}_3) = (z-1,z+1,z^2-4)$ and $\mathcal{D}= (\vec{d}_1,\vec{d}_2) = (\pi,5z)$ are bases of $\Poly_2$ and $\Poly_1$. Find the matrix of $T$ relative to $\mathcal{B}$ and $\mathcal{D}$ and check that it works with $T(2z)$.

\begin{notes}
\vfill \hspace{1pt}
\columnbreak
\end{notes}

\begin{book}
\textbf{Solution.} The formula for the matrix $A$ is:
\begin{align*}
A &= \begin{bmatrix}
\left[ T(\vec{b}_1) \right]_{\mathcal{D}} & \left[ T(\vec{b}_2) \right]_{\mathcal{D}}  & \left[ T(\vec{b}_c) \right]_{\mathcal{D}} 
\end{bmatrix} \\
&= \begin{bmatrix}
\left[ \frac{d}{dz}(z-1) \right]_{\mathcal{D}} & \left[ \frac{d}{dz}(z+1) \right]_{\mathcal{D}}  & \left[ \frac{d}{dz}(z^2-4) \right]_{\mathcal{D}} 
\end{bmatrix} \\
&= \begin{bmatrix}
\left[ 1 \right]_{\mathcal{D}} & \left[ 1 \right]_{\mathcal{D}}  & \left[ 2z \right]_{\mathcal{D}} 
\end{bmatrix} \\
&=\begin{bmatrix}
1/\pi & 1/\pi & 0 \\
0 & 0 & 2/5
\end{bmatrix}
\end{align*}
The last step follows from the fact that $1 = \frac{1}{\pi}(\pi) + 0(5z)$ and $2z = 0(\pi) + \frac{2}{5}(5z)$. Now let's check that:
\begin{equation}
\label{eq:check-matrix-7.2}
\left[ T(2z) \right]_{\mathcal{D}} = A \left[ 2z \right]_{\mathcal{D}}
\end{equation}
The left side is:
\[\left[ T(2z) \right]_{\mathcal{D}} = \left[ \frac{d}{dz}(2z) \right]_{\mathcal{D}} = \left[ 2 \right]_{\mathcal{D}} = \bbvec{2/\pi}{0}\]
The last step follows because $2 = \frac{2}{\pi}(\pi) + 0(5z)$. To compute the right side of equation \ref{eq:check-matrix-7.2}, we have $\left[ 2z \right]_{\mathcal{B}} = \bbvec{1}{1}$ since $2z = 1(z-1) + 1(z+1)$. This gives
\[A \left[ 2z \right]_{\mathcal{D}} = \begin{bmatrix}
1/\pi & 1/\pi & 0 \\
0 & 0 & 2/5
\end{bmatrix}\bbvec{1}{1} = \bbvec{1/\pi + 1/\pi}{0+0} = \bbvec{2/\pi}{0}\]
which equals the left side of equation \ref{eq:check-matrix-7.2} as expected. \blacksquare
\end{book}






\begin{definition}
Suppose $T:V \rightarrow W$ is a linear transformation. The \textbf{null space} of $T$ is the set of vectors $\vec{v}$ in $V$ such that $T(\vec{v}) = \vec{0}$. The \textbf{range} of $T$ is the set of vectors in $W$ that can be written $T(\vec{v})$ for some $\vec{v}$ in $V$. In symbols:
\begin{itemize}
\item $\nul (T) = \left\{ \vec{v} \text{ in } V: T(\vec{v}) = \vec{0}\right\}$.
\item $\range (T) = \left\{ T(\vec{v}) : \vec{v} \text{ is in } V \right\}$.
\end{itemize}
\end{definition}

If $T$ is an $\R^c\rightarrow \R^r$ linear transformation, we know $T(\vec{x})$ can be written $A\vec{x}$, and the null space of $A$ is the same thing as the null space of $T$. And since any linear combination of the columns of $A$ can be written $A\vec{x}$, the column space of $A$ is the range of $T$.

\begin{theorem}
\label{theorem:null-range-subspaces}
Suppose $T:V \rightarrow W$ is a linear transformation. Then $\nul(T)$ is a subspace of $V$ and $\range(T)$ is a subspace of $W$.
\end{theorem}

The proof is left as an exercise.

\example Let $T:\Poly_3\rightarrow \Poly_3$ be the second derivative, that is, $T(\vec{p}(z)) = \frac{d^2}{dz^2}\vec{p}(z)$, and let $\mathcal{B}$ and $\mathcal{D}$ both equal $( 1,z,z^2,z^3 )$, the standard basis of $\Poly_3$.
\begin{enumerate}[label=(\alph*)]
\item Find $A$, the matrix of $T$ relative to $\mathcal{B}$ and $\mathcal{D}$.
\item Find basis for the null space and column space of $A$.
\item Find bases for the null space and range of $T$.
\end{enumerate}

\begin{book}
\textbf{Solution.} (a) The formula in theorem \ref{theorem:abstract-matrix-of-transformation} gives:
\begin{align*}
A &= \begin{bmatrix}
\left[ T(1) \right]_D & \left[ T(z) \right]_D & \left[ T(z^2) \right]_D & \left[ T(z^3) \right]_D
\end{bmatrix} \\
&= \begin{bmatrix}
\left[ 0 \right]_D & \left[ 0 \right]_D & \left[ 2 \right]_D & \left[ 6z \right]_D 
\end{bmatrix} \\
&= \begin{bmatrix}
0 & 0 & 2 & 0 \\
0 & 0 & 0 & 6 \\
0 & 0 & 0 & 0 \\
0 & 0 & 0 & 0
\end{bmatrix}
\end{align*}
(b) For the null space of $A$, we want to solve $A\vec{z}=\vec{0}$, which is equivalent to:
$\begin{aligned}
2z_3&=0\\
6z_4&=0
\end{aligned}$ with $z_1,z_2$ free variables, giving:
\[\begin{bmatrix}
z_1 \\ z_2 \\ z_3 \\ z_4
\end{bmatrix} = \begin{bmatrix}
z_1 \\ z_2 \\ 0 \\ 0
\end{bmatrix} = z_1 \begin{bmatrix}
1 \\ 0 \\ 0 \\ 0
\end{bmatrix} + z_2 \begin{bmatrix}
0 \\ 1 \\ 0 \\ 0
\end{bmatrix}\]
Thus we have $( \vec{e}_1, \vec{e}_2 )$ as a basis for $\nul (A)$. Now, $\col (A)$ is the set of vectors that can be written as the result of $A\vec{z}$ for some $\vec{z}$:
\[\col (A) = \left\{ A\vec{z} : \vec{z} \text{ in } \R^4\right\} = \left\{ \begin{bmatrix}
2z_3 \\ 6z_4 \\ 0 \\ 0
\end{bmatrix} \right\} = \spn \left( \begin{bmatrix}
1 \\ 0 \\ 0 \\ 0
\end{bmatrix}, \begin{bmatrix}
0 \\ 1 \\ 0 \\ 0
\end{bmatrix} \right)\]
So, in this case, $\col (A)$ and $\nul (A)$ are actually equal.

(c) Let's write a formula for $T$ based on the coefficients of an arbitrary $\vec{p}(z)$:
\begin{equation*}
T(c_0+c_1z+c_2z^2+c_3z^3) = \frac{d^2}{dz^2}(c_0+c_1z+c_2z^2+c_3z^3) = 2c_2+6c_3z
\end{equation*}
So, $T(\vec{p}) = \vec{0}$ if and only if $c_2=c_3=0$, which means:
\[\nul(T) = \left\{ c_0+c_1z : c_0,c_1\text{ in } \R\right\} = \spn ( 1,z )\]
Thus $( 1,z )$ is a basis for $\nul(T)$. Note that the coordinates of $(1,z)$ are $(\vec{e}_1,\vec{e}_2)$, the basis of $\nul (A)$. The range of $T$ is:
\[\range(T) = \left\{ 2c_2+6c_3z:c_2,c_3 \text{ in } \R\right\} = \spn(1,z)\]
Thus $(1,z)$ is also a basis for $\range(T)$, and again, the coordinates of this basis are $( \vec{e}_1,\vec{e}_2 )$, the basis for $\col (T)$. \blacksquare
\end{book}

\begin{notes}
\columnbreak
\end{notes}

The next theorem generalizes the above example. 

\begin{theorem} \label{theorem:range-col-null-null}
Suppose $T:V\rightarrow W$ is a linear transformation, with $\mathcal{B} = \left\{ \vec{b}_1,\ldots,\vec{b}_c \right\}$ and $\mathcal{D}$ bases for $V,W$, respectively. Let $A$ be the basis of $T$ relative to $\mathcal{B}$ and $\mathcal{D}$. Then
\begin{enumerate}[label=(\roman*)]
\item $\dim(\range(T)) = \dim(\col(A))$ and
\item $\dim(\nul(T)) = \dim(\nul(A))$.
\end{enumerate}
\end{theorem}

\begin{notes}
The proof is relegated to the textbook, because it's really tedious, but here's the gist:
% In video: (i) Explain how the pivot columns of A are a basis of colA, and correspond with T(bk1),...,T(bkn), which is linearly independent and can be shown to span range(T).
% (ii) Explain how null(A) has a basis u1,...,um and this corresponds with vectors v1,...,vm in V, and one can show that this is a basis of null(T).

\vfill \hspace{1pt}
\newpage
\end{notes}

\begin{book}
\textbf{Proof.} (i) We know that the pivot columns of $A = \begin{bmatrix}
\left[ T(\vec{b}_1) \right]_{\mathcal{D}} & \cdots  & \left[ T(\vec{b}_c) \right]_{\mathcal{D}} 
\end{bmatrix}$ form a basis of $A$ (theorem \ref{theorem:basis-for-column-space}). Denote the pivot column numbers of $A$ as $k_1,\ldots,k_n$, so the pivot columns are $\left[T(\vec{b}_{k_1})\right]_{\mathcal{D}},\ldots, \left[T(\vec{b}_{k_n})\right]_{\mathcal{D}}$. In the previous example, we would say $k_1,k_2=3,4$ because the pivot columns of $A$ were $\left[T(\vec{b}_3)\right]_{\mathcal{D}}, \left[T(\vec{b}_4)\right]_{\mathcal{D}}$. Using this notation, $\dim(\col(A))=n$, the number of pivot columns of $A$. Since $\left[T(\vec{b}_{k_1})\right]_{\mathcal{D}},\ldots, \left[T(\vec{b}_{k_n})\right]_{\mathcal{D}}$ is a basis of $\col(A)$, it's linearly independent and spans $\col(A)$. Theorem \ref{theorem:abstract-more-properties-of-coordinates} (ii) says that $T(\vec{b}_{k_1}), \ldots, T(\vec{b}_{k_n})$ is linearly independent;  we just need to show that $T(\vec{b}_{k_1}), \ldots, T(\vec{b}_{k_n})$ spans $\range(T)$, which would mean it's a basis of $\range(T)$, and $\dim(\range(T))=n$. To do this, let $\vec{w}$ be in $\range(T)$; that is, let $\vec{w}=T(\vec{v})$ for some $\vec{v}$ in $V$. Since $\mathcal{B}$ is a basis, we can write $\vec{v}=x_1\vec{b}_1+\cdots + x_c\vec{b}_c$, which means $\vec{w}=T(x_1\vec{b}_1+\cdots + x_c\vec{b}_c) = x_1T(\vec{b}_1) + \cdots + x_cT(\vec{b}_c)$. Plugging this into $\left[  \right]_{\mathcal{D}}$ gives:
\begin{align*}
\left[ \vec{w} \right]_{\mathcal{D}} &= \left[ x_1T(\vec{b}_1) + \cdots + x_cT(\vec{b}_c) \right]_D \\
&= x_1 \left[ T(\vec{b}_1) \right]_{\mathcal{D}} + \cdots + x_c \left[ T(\vec{b}_c) \right]_{\mathcal{D}}  &\text{(thm \ref{theorem:abstract-properties-of-coordinates} (i), (ii))} \\
\end{align*}
This shows that $\left[ \vec{w} \right]_{\mathcal{D}}$ is in $\col(A)$, and since the pivot columns of $A$ span $\col(A)$, there exists $z_1,\ldots,z_n$ such that:
\begin{align*}
\left[\vec{w}\right]_{\mathcal{D}} &= z_1 \left[ T(\vec{b}_{k_1}) \right]_{\mathcal{D}} + \cdots + z_n \left[ T(\vec{b}_{k_n}) \right]_{\mathcal{D}} \\
\left[\vec{w}\right]_{\mathcal{D}} &=\left[ z_1 T(\vec{b}_{k_1}) + \cdots + z_n  T(\vec{b}_{k_n}) \right]_{\mathcal{D}} &\text{(thm \ref{theorem:abstract-properties-of-coordinates} (i), (ii))}\\
\vec{w} &= z_1T(\vec{b}_{k_1}) + \cdots + z_n T(\vec{b}_{k_n}) &\text{(thm \ref{theorem:abstract-properties-of-coordinates} (iii))}
\end{align*}
Thus $T(\vec{b}_{k_1}), \ldots, T(\vec{b}_{k_n})$ spans $\range(T)$, and as mentioned before, is linearly independent, so it's a basis of $\range(T)$, and $\dim(\range(T)) = n = \dim(\col(A))$.

(ii) We know from theorem \ref{theorem:subspaces-are-spans} that $\nul(A)$ has a basis; let's call it $\vec{u}_1,\ldots,\vec{u}_m$. Writing the components of each as $\vec{u}_i=(u_{i1},\ldots,u_{ic})$, let $\vec{v}_i=u_{i1}\vec{b}_1+\cdots+u_{ic}\vec{b}_c$ for $1\le i \le m$; in other words, let $\left[ \vec{v}_i \right]_{\mathcal{B}} = \vec{u}_i$. We'll show that $\vec{v}_1,\ldots,\vec{v}_m$ is a basis of $\nul(T)$. First, we need to show that they're \emph{in} $\nul(T)$. Theorem \ref{theorem:abstract-matrix-of-transformation} gives, for each $1\le i \le m$,
\[\left[ T(\vec{v}_i) \right]_{\mathcal{D}} = A \left[ \vec{v}_i \right]_{\mathcal{B}} = A\vec{u}_i=\vec{0} = \left[ \vec{0} \right]_{\mathcal{D}}\]
and since $\left[  \right]_{\mathcal{D}}$ is one-to-one (theorem \ref{theorem:abstract-properties-of-coordinates} (iii)), we have $T(\vec{v}_i) = \vec{0}$, and each $\vec{v}_i$ is in $\nul(T)$. Now, let $\vec{y}$ be any vector in $\nul(T)$. Then, $\left[ \vec{y} \right]_{\mathcal{B}}$ is in $\nul(A)$, since $T(\vec{y})=\vec{0}$ and $T(\vec{y})=A \left[ \vec{y} \right]_B$. Since $\vec{u}_1,\ldots,\vec{u}_m$ is a basis of $\nul(A)$, there exist $x_1,\ldots,x_m$ such that:
\begin{align*}
x_1\vec{u}_1+\cdots+x_m\vec{u}_m&=\left[ \vec{y} \right]_{B} \\
x_1 \left[ \vec{v}_1 \right]_{\mathcal{B}} + \cdots + x_m \left[ \vec{v}_m \right]_{\mathcal{B}} &= \left[ \vec{y} \right]_{\mathcal{B}} \\
\left[ x_1\vec{v}_1 + \cdots + x_m\vec{v}_m \right]_{\mathcal{B}} &= \left[ \vec{y} \right]_{\mathcal{B}} \\
x_1\vec{v}_1 + \cdots + x_m\vec{v}_m  &= \vec{y} & \text{(theorem \ref{theorem:abstract-properties-of-coordinates} (iii))} 
\end{align*}
So, any vector in $\nul(T)$ can be written as a linear combination of $\vec{v}_1,\ldots,\vec{v}_m$. And since $\vec{u}_1,\ldots,\vec{u}_m$ is a basis, it's linearly independent, and so is $\vec{v}_1,\ldots,\vec{v}_m$ (theorem \ref{theorem:abstract-more-properties-of-coordinates} part (ii)). Thus $\vec{v}_1,\ldots,\vec{v}_m$ is a basis of $\nul(T)$ and $\dim(\nul(T)) = m = \dim(\nul(A))$. \blacksquare

\columnbreak
\end{book}

\exerciseheading 

For questions \ref{exercise:LT-or-not-start} through \ref{exercise:LT-or-not-end}, determine if the given transformation is linear or not. Justify your answer.

\exatend \label{exercise:LT-or-not-start} $T:\Poly_2\rightarrow \Poly_2$, $T(\vec{p}(z)) = \vec{p}(z) + 7$.

\exatend $T:\Poly_2\rightarrow \Poly_3, T(\vec{p}(z)) = (z+3)\vec{p}(z)$.

\exatend $T:\Poly_2\rightarrow \R^3, T(\vec{p}(z)) = ( \vec{p}(0), \vec{p}(1), \vec{p}(2) )$.

\exatend $T:\Poly_2 \rightarrow \Poly_2, T(a+bz+cz^2) = cz^2$.

\exatend $T:\Poly_2\rightarrow \R^3, T(\vec{p}(z)) = ( \vec{p}(0)+8, \vec{p}(1)+8, \vec{p}(2)+8 )$.

\exatend Let $M$ be the vector space of all $4\times 4$ matrices, and $\vec{p}=(8,9,10,11)$. $T:M \rightarrow \R^4, T(\vec{A}) = \vec{A}\vec{p}$.

\exatend Let $M$ be the vector space of all $4\times 4$ matrices. $T: M \rightarrow M, T(\vec{A}) = \det(\vec{A})$.

\exatend Let $M$ be the vector space of all $4\times 4$ matrices. $T: M \rightarrow \R, T(\vec{A}) = \dim(\nul \vec{A} )$.

\exatend $T:\Poly \rightarrow \R, T(\vec{p}(z)) = \int_{-1}^1 \vec{p}(z)dz$.

\exatend $T:\Poly \rightarrow \Poly, T(\vec{p}(z)) = \int \vec{p}(z)dz$, with constant of integration always set to $0$. For example, $T(1) = z$, and $T(z+z^3) = \frac{1}{2}z^2+\frac{1}{4}z^4$.


\exatend \label{exercise:LT-or-not-end} Let $\vec{w}=(3,4,5)$, and $W = \spn(\vec{w})$. $T:\R^3 \rightarrow W, T(\vec{x}) = \proj_W\vec{x}$.

\exatend (T/F?) If $m,b$ are any real numbers, then $f:\R \rightarrow \R$ defined by $f(z) = mz + b$ is a linear transformation because the graph is a straight line.

\exatend (T/F?) If $T:V\rightarrow W$ is a linear transformation, then $S:V\rightarrow W$ defined by $S(\vec{v}) = 7T(\vec{v})$ is a linear transformation.

\exatend (T/F?) If $T:V\rightarrow W$ and $S:V \rightarrow W$ are linear transformations, then $T+S:V\rightarrow W$ defined by $(T+S)(\vec{v}) = T(\vec{v}) +S(\vec{v})$ is a linear transformation.

\exatend (T/F?) If $T:V\rightarrow W$ and $S:W \rightarrow U$ are linear transformations, then $TS:V\rightarrow U$ defined by $TS(\vec{v}) = T(S(\vec{v}))$ is a linear transformation.

\exatend Suppose $\dim(V) = 5,\dim(W) = 4$, $\mathcal{B}$ and $\mathcal{D}$ are bases of $V$ and $W$ respectively, and $T:V\rightarrow W$ is a linear transformation. Fill in the blanks: The size of the matrix for $T$ relative to $\mathcal{B}$ and $\mathcal{D}$ is $\rule{2em}{0.8pt}\times \rule{2em}{0.8pt}$.

\exatend Prove theorem \ref{theorem:null-range-subspaces}.

\exatend Prove the following theorem. \hintfootnote{Use thm \ref{theorem:rank-nullity}.}

\begin{theorem}
Suppose $T:V\rightarrow W$ is a linear transformation. Then $\dim(\range(T)) + \dim(\nul(T)) = \dim(V)$.
\end{theorem}

\exatend Suppose that $\mathcal{B} = ( \vec{b}_1,\ldots,\vec{b}_c )$ and $\mathcal{D}=( \vec{d}_1,\ldots, \vec{d}_c )$ are two bases of a vector space $V$. The matrix $P = \begin{bmatrix}
\left[ \vec{b}_1 \right]_{\mathcal{D}}  & \cdots & \left[ \vec{b}_c \right]_{\mathcal{D}}
\end{bmatrix}$ is called the \emph{change of coordinates matrix from $\mathcal{B}$ to $\mathcal{D}$}. Show that $\left[ \vec{v} \right]_{\mathcal{D}} = P \left[ \vec{v} \right]_{\mathcal{B}}$ for all $\vec{v}$ in $V$. % easy way: it's the matrix of the transformation T(v) = v.

\exatend Define a linear transformation $T:\Poly_2\rightarrow \R^4$ by $T(\vec{p}(z)) = ( \vec{p}(0), \vec{p}(1), \vec{p}(2), \vec{p}(3) )$. Let $\mathcal{B} = ( 1,z,z^2 )$ and $\mathcal{D}=( \vec{e}_1,\vec{e}_2,\vec{e}_3, \vec{e}_4 )$ be the standard bases.
\begin{enumerate}[label=(\alph*)]
\item Find $A$, the matrix of $T$ relative to $\mathcal{B}$ and $\mathcal{D}$.
\item Check that theorem \ref{theorem:abstract-matrix-of-transformation} works with $\left[ T(2+3z) \right]_{\mathcal{D}}$.
\item How many pivot positions does $A$ have? What does this say about $\dim(\nul(T))$ and $\dim(\range(T))$?
\end{enumerate}

\exatend The first four \emph{Hermite polynomials}, $\mathcal{H}=( 1,2z, 4z^2-2, 8z^3-12z )$, form a basis of $\Poly_3$, and are useful in various areas of math and physics. Let $T:\Poly_3\rightarrow \Poly_3$ be differentiation: $T(\vec{p}(z)) = \frac{d}{dz}\vec{p}(z)$. Find $A$, the matrix of $T$ relative to $\mathcal{H}$ and $\mathcal{H}$.

\exatend Define a linear transformation $T:\Poly_3\rightarrow \Poly_4$ by $T(\vec{p}(z)) = \int \vec{p}(z)dz$, with the constant of integration always set to $0$. For example, $T(2z + z^3) = z^2+\frac{1}{4}z^4$. Let $\mathcal{B}=( 1,z,z^2,z^3 )$ and $\mathcal{D}=( 1,z,z^2,z^3,z^4 )$ be the standard bases.
\begin{enumerate}[label=(\alph*)]
\item Find $A$, the matrix of $T$ relative to $\mathcal{B}$ and $\mathcal{D}$.
\item Describe $\nul(A)$ and $\nul(T)$.
\item Describe $\col(A)$ and $\range(T)$.
\end{enumerate}

\exatend Define a linear transformation $T:\Poly_2\rightarrow \Poly_3$ by $T(\vec{p}(z))=(z-5)\vec{p}(z)$. Let $\mathcal{B}=( 1,z,z^2 )$ and $\mathcal{D}=( 1,z,z^2,z^3 )$ be the standard bases.
\begin{enumerate}[label=(\alph*)]
\item Find $A$, the matrix of $T$ relative to $\mathcal{B}$ and $\mathcal{D}$.
\item Describe $\nul(A)$ and $\nul(T)$.
\item Describe $\col(A)$ and $\range(T)$.
\end{enumerate}

\exatend Let $M$ be the vector space of all $2\times 2$ matrices, with basis:
\[\mathcal{B}=\vec{B}_1,\vec{B}_2,\vec{B}_3,\vec{B}_4 = \begin{bmatrix}
1 & 0 \\
0 & 0
\end{bmatrix},\begin{bmatrix}
0 & 1 \\
0 & 0
\end{bmatrix},\begin{bmatrix}
0 & 0 \\
1 & 0
\end{bmatrix},\begin{bmatrix}
0 & 0 \\
0 & 1
\end{bmatrix}\]
Define $T:M\rightarrow \R^2$ by $T(\vec{A}) = A \bbvec{1}{2}$. 
\begin{enumerate}[label=(\alph*)]
\item Find $A$, the matrix of $T$ relative to $\mathcal{B}$ and the standard basis of $\R^2$.
\item Find a basis of $\nul(A)$.
\item Find a basis of $\nul(T)$.
\end{enumerate}

The next problem is about \emph{eigenvectors} and \emph{eigenvalues} of linear transformations, which mean the same thing as they did in chapter \ref{chap:eigenstuff}:

\begin{definition}
Suppose $T:V\rightarrow V$ is a linear transformation. We say $\vec{u}$ is an \textbf{eigenvector} corresponding with an \textbf{eigenvalue} $\lambda$ (a scalar) of $T$, if:
\[T(\vec{u})=\lambda \vec{u} \quad \text{and} \quad \vec{u}\ne\vec{0}\]
\end{definition}

\exatend Let $V$ be the set of all differentiable $\R \rightarrow \R$ functions.
\begin{enumerate}[label=(\alph*)]
\item Define $T(\vec{f}(z)) = \frac{d}{dz}\vec{f}(z)$. Find a couple eigenvectors of $T$ and their corresponding eigenvalues.
\item Define $S(\vec{f}(z)) = \frac{d^2}{dz^2}\vec{f}(z)$. Find a couple eigenvectors of $S$ that are \emph{not} eigenvectors of $T$. \hintfootnote{It's a basic kind of function that isn't polynomial, logarithmic, or exponential.}
\end{enumerate}


\appendix
\chapter{Proof-Writing Basics}

This optional section covers the basics of writing proofs of ``If $P$, then $Q$'' and ``$P$ if and only if $Q$'' statements, and should be covered before section \ref{sec:vectors-linear-combinations-span} or \ref{sec:homogeneous-systems-and-linear-independence} for students that have no experience with proof-writing.

A mathematical \emph{proof} is a written (or spoken) argument, which should convince the reader beyond any doubt that a statement (sometimes called a \emph{proposition} or a \emph{theorem}) is true, as long as the reader has the background knowledge required to understand it. Writing and reading proofs is the main way that mathematicians do mathematical research, in the same way that experiments are the main way that scientists do scientific research. When you write a proof for an assignment in a math class, it serves several purposes:
\begin{itemize}
\item It shows that you understand why the statement is true.
\item It gets you to study the definitions, theorems, and examples related to the statement you're proving.
\item It gets you to practice your reasoning and argument-writing skills, which you can apply to other areas of life, and is one of the reasons you're going to college.
\end{itemize}
Before we get into writing proofs, let's look at the concept of a \emph{definition}. In math, definitions are different than they are in, say, a dictionary. Dictionary definitions are \emph{descriptive}, which means they try to \emph{describe} the different ways people use a word. For example, when I google the word ``literally'' I get the following definition from Oxford Languages:
\begin{enumerate}
\item In a literal manner or sense; exactly. ``The driver took it literally when asked to go straight across the traffic circle.''
\item \emph{(informal)} Used for emphasis or to express strong feeling while not being literally true. ``I was literally blown away by the response I got.''
\end{enumerate}
Imagine if you needed to \emph{prove}, beyond any doubt, that something is \emph{literal} or not. It would be hopeless, because, for one, the first definition uses the word in its own definition. Even worse, the second definition is the opposite of the first! 

In math, a definition can be \emph{prescriptive}, which means that the definition is telling you exactly what the word means, every time it's used in the paper/book/course, with no room for nuance or interpretation. Sometimes one book, say, book $A$, will define a term differently than book $B$; but when you're talking about results in book $A$ you need to use the definitions of that book to avoid confusion. When the authors of $A$ and $B$ get together to talk, they need to start by agreeing on which definition they're using before using the controversial term.

Here's an example:

\begin{definition}
An integer $m$ is \textbf{even} if it can be written $m=2k$ for some integer $k$. An integer $n$ is \textbf{odd} if it can be written $n=2k+1$ for some integer $k$.
\end{definition}
From this point on, any time you see the term ``even integer'' you know that's a shorthand way of saying ``an integer that can be written as $2k$ for some integer $k$.'' And using any other meaning of the word, from everyday English, is a \textbf{bad idea}:
\begin{itemize}
\item One million is nice and even (\emph{round?}) number but 998,576 is not.
\item Zero is an odd (\emph{special, unusual?}) number because it's the only number that can be added to other numbers without changing them.
\end{itemize}
In everyday English, these statements are ambiguous, but true in some sense. In a math class that uses the above definition, they are completely false.

Now, the purpose of definitions is to make statements shorter and easier to read. For example, consider the two statements:
\begin{enumerate}
\item If $x,y$ are odd integers, then $x+y$ is an even integer.
\item If $x,y$ are integers that can be written $x=2k+1$ and $y=2c+1$ for some integers $c$ and $k$, then $x+y=2p$ for some integer $p$.
\end{enumerate}
These two statements mean exactly the same thing; the first one is easier-to-read shorthand for the second. When you prove the first statement, you need to keep in mind that you're really proving the second one.

Many students, when confronted with their first proof-writing assignments, don't know where to begin; the following strategy can help. Even experienced mathematicians follow similar strategies when they're stuck.

\begin{bbox}
\textbf{Strategy for proving statements of the form ``If $P$, then $Q$'':}
\begin{enumerate}
\item  On scratch paper, get some intuition for why the proof is true, with examples or pictures. (\emph{Optional.})
\item Start the proof by writing ``Assume $P$'' at the top, but with $P$ written in expanded terms, using the definition.
\item Write ``WTS [Want To Show] $Q$'', but with $Q$ in expanded terms, using the definition. \emph{(Optional.)}
\item Use your assumption of $P$ and any facts you know to show $Q$ is true. Also, it's nice to end the proof with a symbol like ``\blacksquare'' or ``Q.E.D.'' to let the reader know that you're done.
\end{enumerate}
\end{bbox}

Now, steps 1 and 3 are optional because a proof can be perfectly valid and convincing without them, but they're a good idea if you're not sure how to write the proof.

\example Prove that if $x,y$ are odd integers, then $x+y$ is an even integer.

\begin{notes}
\vfill \hspace{1pt}
\columnbreak
\end{notes}

\begin{book}
\textbf{Step 1.} On scratch paper, maybe think of some $x,y$ pairs of odd integers. Make sure you know they're odd using the definition. See what happens when you add them up, and think about why the result is forced to be even. It might look something like the following. (Everything in blue represents what you might write.)
\color{blue}\begin{align*}
x=5,y=13 \quad &\rightarrow \quad x+y = 18\\
x=2(2)+1, y=2(6) +1 \quad &\rightarrow \quad 2(2)+1 + 2(6)+1 = 2(8)+2 = 2(9) \\
x=-7, y=1 \quad &\rightarrow \quad x+1=-6\\
x=2(-4)+1, y=2(0) +1 \quad &\rightarrow \quad 2(-4)+1 + 2(0)+1 = 2(-4)+2 = 2(-3)
\end{align*}

\color{black}

\textbf{Steps 2, 3.} Start the proof by writing ``Assume $x,y$ are odd integers. WTS $x+y$ is even'' but written in expanded terms, using the definition, like below.

\color{blue} Proof: Assume $x=2k+1$, $y=2c+1$, for some integers $k,c$. WTS $x+y = 2p$ for some integer $p.$

\color{black}

\textbf{Step 4.} This is the hard part, but hopefully step 1 gave you some intuition, and steps 2,3 gave you a starting point. You can use the facts that $x=2k+1$ and $y=2c+1$, and any other facts you know, and the goal is to show that $x+y$ can be written $2p$ for some integer $p$. Here's what it might look like:
\color{blue}
\begin{align*}
x+y = 2k+1 + 2c+1 &= 2(k+c) + 2 \\
&= 2(k+c+1) = \text{ an even integer. } \blacksquare
\end{align*}
\color{black}
\end{book}

\begin{definition}
An integer $n$ is \textbf{divisible by} an integer $p$ if $n=pq$ for some integer $q$.
\end{definition}

\example Show that if $n$ is an odd integer, then $n^2-1$ is divisible by 8.

\begin{notes}
\vfill \hspace{1pt}
\columnbreak
\end{notes}

\begin{book}
\textbf{Solution.} \textbf{Step 1.} This one is trickier, but if you're able to figure out the trick quickly, your scratch work might look like this:
\color{blue}\btwocol{
\begin{align*}
n=7&=2(3) +1 \\
n^2-1 &= 49-1 = 8(6) \\
&= (2(3)+1)(2(3) + 1) -1 \\
&= 4(3^2) + 4(3) +1 -1 \\
&= 4(3^2 + 3) \\
&= 4(3)(3+1)
\end{align*}}{
\begin{align*}
n=13&=2(6) +1 \\
n^2-1 &= 169 -1 = 8(21) \\
&= (2(6)+1)(2(6)+1) -1 \\
&= 4(6^2) + 4(6) \\
&= 4(6)(6+1)
\end{align*}
}
\color{black}
%\includegraphics[width=.8\linewidth]{2026-01-07_181517.png}

\textbf{Step 2, 3.} The start of the proof should look something like this:

%\includegraphics[width=.8\linewidth]{2026-01-07_181610.png}
\color{blue} Proof: Let $n=2k+1$ for some integer $k$. WTS $n^2-1=8c$ for some integer $c$.
\color{black}

\textbf{Step 4.} 
\color{blue}
\begin{align*}
n^2-1 &= (2k+1)(2k+1) -1 \\
&= 4k^2+4k+1 -1 \\
&= 4(k^2+k) \\
&= 4(k)(k+1)
\end{align*}
Either $k$ or $k+1$ must be even. If $k$ is, that is if $k=2j$, then $n^2-1 = 4(2j)(k+1) = 8j(k+1)$ which is divisible by 8. If $k+1$ is even, that is if $k+1=2j$, then $n^2-1 = 4k(2j) = 8kj$ which is divisible by 8. \blacksquare
\color{black}
%\includegraphics[width=.8\linewidth]{2026-01-07_181721.png}

%\includegraphics[width=.8\linewidth]{2026-01-07_181754.png}

%\includegraphics[width=.8\linewidth]{2026-01-07_181811.png}

Note that this proof uses the fact that if $k$ is an integer, then either $k$ or $k+1$ must be even. In my class, this falls under the umbrella of ``prior math knowledge that everyone knows'' but in some classes, you may be expected to prove statements like this before using them.
\end{book}

\begin{minipage}[t]{1.0\linewidth}
The statement ``If not $Q$, then not $P$'' is called the \textbf{contrapositive} of ``If $P$, then $Q$.'' The picture below on the left should give you an idea why a statement is equivalent to its contrapositive. Sometimes it's easier to prove the contrapositive than the original statement.

The statement ``If $Q$, then $P$'' is called the \textbf{converse} of ``If $P$, then $Q$.'' The figures below in the left and middle show converses of each other.

The statement ``$P$ if and only if $Q$'' means ``If $P$, then $Q$, \emph{and} if $Q$, then $P$.'' In other words, if one of them is true, the other is true, and if one is false, the other is false.

\begin{tikzpicture}[scale=0.8,fill opacity=0.3,text opacity=1]
\draw[thick] (-2.5,-3.2) rectangle (2.1,2.1);
\node[above right,rectangle] at (-2.5,-3.2) {
\begin{tabular}{l}
If \textcolor{blue}{$P$}, then \textcolor{red}{$Q$}. \\ 
If not \textcolor{red}{$Q$}, then not \textcolor{blue}{$P$}.
\end{tabular}
};
\draw[thick,red, fill=red!30] (0,0) circle (1.8cm);
\node[red] at (-1, 1) {$Q$};
\draw[thick,blue,fill=blue!30] (0.5,0) circle (1cm);
\node[blue] at (0, 0) {$P$};
\end{tikzpicture}\quad
\begin{tikzpicture}[scale=0.8,fill opacity=0.3,text opacity=1]
\draw[thick] (-2.5,-3.2) rectangle (2.1,2.1);
\node[above right,rectangle] at (-2.5,-3.2) {
\begin{tabular}{l}
If \textcolor{red}{$Q$}, then \textcolor{blue}{$P$}. \\ 
If not \textcolor{blue}{$P$}, then not \textcolor{red}{$Q$}.
\end{tabular}
};
\draw[thick,red,fill=red!30] (0.5,0) circle (1cm);
\draw[thick,blue, fill=blue!30] (0,0) circle (1.8cm);
\node[blue] at (-1, 1) {$P$};
\node[red] at (0, 0) {$Q$};
\end{tikzpicture} \quad \begin{tikzpicture}[scale=0.8,fill opacity=0.3,text opacity=1]
\draw[thick] (-2.5,-3.2) rectangle (2.1,2.1);
\node[above right,rectangle] at (-2.5,-3.2) {
\textcolor{blue}{$P$} if and only if \textcolor{red}{$Q$}.};
\draw[thick,purple,fill=red!30] (0,0) circle (1.8cm);
\draw[thick,purple, fill=blue!30] (0,0) circle (1.8cm);
\node[blue] at (-1, 1) {$P$};
\node[red] at (-1.4, 0) {$Q$};
\end{tikzpicture}
\end{minipage}

\example Suppose $n$ is an integer. Show that $n$ is odd if and only if $n^2$ is odd.

\begin{notes}
\vfill \hspace{1pt}
\columnbreak
\end{notes}

\begin{book}
\textbf{Solution.} \textbf{Step 1.} Here's what some scratch work might look like:

\color{blue} Odd $n$: 
\begin{align*}
n=11&=2(5)+1 \\
n^2 = 121 &= (2(5)+1)(2(5)+1) \\
&= 100 + 10 + 10 + 1 = \text{ odd.}
\end{align*}
Odd $n^{2}$: $n^2=81 = 2(40)+1.$ Hmmmm...
\color{black}
%\includegraphics[width=.7\linewidth]{2026-01-08_082347.png}

\textbf{Step 2, 3, 4 (\rightarrow).} Here we want to prove that if $n$ is odd, then $n^2$ is odd. The proof is straightforward.

\color{blue} Proof (\rightarrow). Suppose $n=2k+1$. WTS $n^2=2c+1$ for some integer $c$.
\begin{align*}
n^2=(2k+1)(2k+1) &= 4k^2+4k+1 \\
&= 2(2k^2+2k) +1 = \text{ odd.}
\end{align*}
\color{black}
\textbf{Step 2, 3, 4 (\leftarrow).} Here we want to prove that if $n^2$ is odd, then $n$ is odd.
\color{blue} (\leftarrow) Suppose $n$ is an integer, and $n^2= 2p+1$ for some integer $p$. WTS $n=2q+1$ for some integer $q$.

Hmmm... I'm stuck. \color{black}

I got stuck because I don't see an easy way to break down $n^2=2p+1$ into $n$. So let's erase that and try a different approach. I know that the statement ``If $n^2$ is odd, then $n$ is odd'' is equivalent to its contrapositive ``If $n$ is not odd, then $n^2$ is not odd.'' So let's prove the contrapositive instead:

\color{blue}
Suppose $n=2k$ for some integer $k$. WTS $n^2$ must be even.
\[n^2=(2k)^2=4k^2 = 2(2k^2) = \text{ an even integer.}\] 
Thus if $n$ is even, then $n^2$ is even, so the contrapositive of this is true too: If $n^2$ is odd (not even), then $n$ is odd (not even). \blacksquare
\color{black}
\end{book}

So far, we've been showing that statements are \emph{true}. But how do you show that a statement of the form ``If $P$, then $Q$'' is \emph{false}? By finding a \textbf{counterexample}, which makes $P$ true but $Q$ false.

\example Show that it is not true that if $a-1$ is divisible by 8, then $a = n^2$ for some odd integer $n$.

\begin{notes}
\vfill \hspace{1pt}
\columnbreak
\end{notes}

\begin{book}
\textbf{Solution.} The goal is to find $a-1$ that is divisible by $8$, but $a$ cannot be written $n^2$ for any odd integer $n$. A good first guess is $a-1 = 9-1 = 8(1)$. But $a = 9 = 3^2$. Does this mean the statement is actually true? No, there are many other multiples of $8$ to try.

Let's try $a-1 = 17-1=8(2)$. Now, $17$ cannot be written $n^2$ for any odd integer $n$, so we've found a counterexample showing the statement is false. \blacksquare
\end{book}


\exerciseheading

\exatend Show that if $n,m$ are even integers, then $n+m$ is even.

\exatend Show that if $n,m$ are odd integers, then $n+m$ is even.

\exatend Show that if $a$ is even, then $3a$ is even.

\exatend Show that if $n$ is an integer, then $(n+1)^2-n^2$ is odd.

\exatend Prove or find a counterexample to the statement: If $n,m$ are integers and $n+m$ is odd, then $n$ is odd.

\exatend Suppose $a$ is an integer. Show that $a$ is even if and only if $3a$ is even.

\exatend Suppose $n$ is an integer. Show that $n$ is odd if and only if $n+5$ is odd.

\exatend Show that if $n$ is an odd integer, then there exist integers $p,q$ such that $p^2-q^2=n$.


\exatend Prove, or find a counterexample to the statement: There exist integers $p,q$ such that $p^2-q^2=n$ if and only if $n$ is odd.

\begin{definition}
An integer $a$ is \textbf{prime} if it's greater than 1, and there do not exist integers $p,q$ greater than 1 such that $a=pq$.
\end{definition}

\exatend (\emph{challenge}) Show that if $a$ is an integer greater than 1, then $a^3+1$ is not prime.

\exatend (\emph{challenge}) Show that if $n$ is a non-prime integer, then $2^n-1$ is not prime.

\exatend Suppose $n$ is an integer. Find a counterexample to the statement: ``$n$ is prime if and only if $2^n-1$ is prime.'' 


\end{multicols*}
\end{document}
